\documentclass[lang=cn,newtx,10pt,scheme=chinese,thmcnt=chapter]{elegantbook}
\usepackage{float}
\title{泛函分析}
\date{\today}
\bioinfo{邮箱}{1651486833@qq.com}
\extrainfo{不要以为抹消过去,重新来过,即可发生什么改变.—— 比企谷八幡}
\cover{cover.jpg}
\newenvironment{xiti}{\par%
\centerline{\textbf{习\quad 题\quad \thesection}}
\begin{enumerate}
}{\end{enumerate}}
\usepackage{tasks}
\usepackage{float} 
\settasks{
label = (\arabic*),
item-indent = 1.7em,
label-width = 0.5em,
label-offset = 1.2em,
column-sep=10pt,
label-align=left,
after-item-skip=0pt
}


\usepackage{tikz-cd}
\usepackage{amsfonts}

\usepackage{fancyhdr}
\usepackage{tkz-euclide}
\usetikzlibrary{shapes.callouts,arrows.meta,calc,shadings,patterns}
\renewcommand{\Re}{\operatorname{Re}}
\renewcommand{\Im}{\operatorname{Im}}
\newcommand{\ii}{\mathrm i}
\DeclareMathOperator{\ee}{\!\!\;\mathrm e}
\newcommand{\MR}{\mathbb R}
\newcommand{\MC}{\mathbb C}
\newcommand{\MF}{\mathbb F}
\newcommand{\MZ}{\mathbb Z}
\newcommand{\MN}{\mathbb N}
\newcommand{\MCF}{\mathcal F}
\newcommand{\bb}{\mathbb}
\newcommand{\dif}{\mathrm{d}}
\newcommand{\ci}{\operatorname{i}}
\newcommand{\Ker}{\operatorname{Ker}}
\newcommand{\Coker}{\operatorname{Coker}}
\author{Charlotte}
% 本文档命令
\usepackage{array}
\newcommand{\ccr}[1]{\makecell{{\color{#1}\rule{1cm}{1cm}}}}
\usepackage{algorithm}

\usepackage{algpseudocode}
% 修改标题页的橙色带
% \definecolor{customcolor}{RGB}{32,178,170}
% \colorlet{coverlinecolor}{customcolor}

\usepackage{paratype}

\usepackage{extarrows}
\usepackage{amsmath}
\usepackage{amsfonts}
\usepackage{amssymb}
\usepackage{esint}

\usepackage{graphicx}
\usepackage[export]{adjustbox}
\usepackage{mdframed}
\usepackage{booktabs,array,multirow}
\usepackage{adjustbox}
\usepackage{tabularx}
\usepackage{hyperref}
\hypersetup{colorlinks=true, linkcolor=blue, filecolor=magenta, urlcolor=cyan,}
\urlstyle{same}


\newcommand{\customfootnote}[1]{
  \let\thefootnote\relax\footnotetext{#1}
}


\cover{cover.jpg}
\author{Charlotte}
% 本文档命令
\usepackage{array}

\usepackage{algorithm}

\usepackage{algpseudocode}
\usepackage{ctex}  % 支持中文
% 修改标题页的橙色带
% \definecolor{customcolor}{RGB}{32,178,170}
% \colorlet{coverlinecolor}{customcolor}
\begin{document}
\maketitle
\tableofcontents
\mainmatter
\chapter{度量空间}
\section{度量空间}
\begin{definition}
设 $\mathscr{X}$ 是一个非空集。$\mathscr{X}$ 叫做距离空间，是指在 $\mathscr{X}$ 上定义了一个双变量的实值函数 $\rho(x, y)$，满足下列三个条件：
\begin{enumerate}
    \item $\rho(x, y) \geqslant 0$，而且 $\rho(x, y)=0$ 当且仅当 $x=y$；
    \item $\rho(x, y)=\rho(y, x)$；
    \item $\rho(x, z) \leqslant \rho(x, y)+\rho(y, z) \quad (\forall x, y, z \in \mathscr{X})$。
\end{enumerate}
这里 $\rho$ 叫做 $\mathscr{X}$ 上的一个距离；以 $\rho$ 为距离的距离空间 $\mathscr{X}$ 记做 $(\mathscr{X}, \rho)$。
\end{definition}
\begin{note}距离概念是欧氏空间中两点间距离的抽象。事实上，如果对 $\forall x=\left(x_1, x_2, \cdots, x_n\right), y=\left(y_1, y_2, \cdots, y_n\right) \in \mathbb{R}^n$，定义
\begin{equation*}
\rho(x, y)=\left[\left(x_1-y_1\right)^2+\cdots+\left(x_n-y_n\right)^2\right]^{\frac{1}{2}} ,
\end{equation*}
则容易看到条件 (1)、(2)、(3) 都满足。以后当说到欧氏空间时，我们始终用这个 $\rho$ 规定其上的距离。
\end{note}


\begin{example}
\textbf{空间 $C[a, b]$} 区间 $[a, b]$ 上的连续函数全体记为 $C[a, b]$，按距离
\begin{equation*}
\rho(x, y) \triangleq \max _{a \leqslant t \leqslant b} |x(t)-y(t)|
\end{equation*}
形成距离空间 $(C[a, b], \rho)$，以后简记作 $C[a, b]$。以后当说到连续函数空间 $C[a, b]$ 时，我们始终用上述距离 $\rho$，除非另外说明。
\end{example}

\begin{definition}
设 $(\mathscr{X}, \rho)$ 是一个距离空间。$\mathscr{X}$ 上的点列 $\left\{x_n\right\}$ 叫做收敛到 $x_0$，是指 $\rho\left(x_n, x_0\right) \rightarrow 0 \ (n \rightarrow \infty)$。这时记作 $\lim\limits _{n \rightarrow \infty} x_n=x_0$，或简单地记作 $x_n \rightarrow x_0$。

\end{definition}

\begin{note}在 $C[a, b]$ 中，点列 $\left\{x_n\right\}$ 收敛到 $x_0$ 是指 $\left\{x_n(t)\right\}$ 一致收敛到 $x_0(t)$。
\end{note}
\begin{definition}
设 $(\mathscr{X}, \rho)$ 是一个度量空间。度量空间中的一个子集 $E$ 称为闭集，指的是：对任意的点列 $\left\{x_n\right\} \subset E$，若 $x_n \rightarrow x_0$，则 $x_0 \in E$。
\end{definition}

\begin{definition}
设 $(\mathscr{X}, \rho)$ 是一个距离空间。$\mathscr{X}$ 上的点列 $\left\{x_n\right\}$ 叫做基本列，指的是 $\rho\left(x_n, x_m\right) \rightarrow 0$ 当 $n, m \rightarrow \infty$。换句话说，$\forall \varepsilon > 0, \exists N(\varepsilon)$，使得 $m, n \geqslant N(\varepsilon) \Rightarrow \rho\left(x_n, x_m\right) < \varepsilon$。如果 $\mathscr{X}$ 中所有基本列都是收敛列，则称 $\mathscr{X}$ 是完备的。
\end{definition}

\begin{example}
$(\mathbb{R}^n, \rho)$ 是完备的。
\end{example}

\begin{example}
$(C[a, b], \rho)$ 是完备的，其中 $\rho(x, y) = \max\limits_{a \leqslant t \leqslant b} |x(t) - y(t)|$。
\end{example}

\begin{proof}设 $\left\{x_n\right\}$ 是 $(C[a, b], \rho)$ 中的一个基本列，则 $\forall \varepsilon > 0$，存在 $N(\varepsilon)$ 使得对任意的 $m, n \geqslant N(\varepsilon)$ 有
\begin{equation*}
\rho(x_m, x_n) = \max_{a \leqslant t \leqslant b} |x_m(t) - x_n(t)| < \varepsilon .
\end{equation*}
因此，对任意的 $t \in [a, b]$，
\begin{equation*}
|x_m(t) - x_n(t)| < \varepsilon \quad (\forall m, n \geqslant N(\varepsilon)) .
\end{equation*}
固定 $t \in [a, b]$，可知数列 $\left\{x_n(t)\right\}$ 是基本列，因此极限 $\lim\limits_{n \rightarrow \infty} x_n(t)$ 存在。用 $x_0(t)$ 表示该极限，令 $m \rightarrow \infty$，得到
\begin{equation*}
|x_0(t) - x_n(t)| \leqslant \varepsilon \quad (\forall n \geqslant N(\varepsilon)) .
\end{equation*}
由此可见，$x_n(t)$ 一致收敛到 $x_0(t)$，从而 $x_0(t)$ 连续，因此 $x_0 \in C[a, b]$，且 $x_n \rightarrow x_0$。
\end{proof}
\begin{definition}
设 $T:(\mathscr{X}, \rho) \rightarrow (\mathscr{Y}, r)$ 是一个映射，称它是连续的，如果对于 $\mathscr{X}$ 中的任意点列 $\left\{x_n\right\}$ 和点 $x_0$，
\begin{equation*}
\rho\left(x_n, x_0\right) \rightarrow 0 \Rightarrow r\left(T x_n, T x_0\right) \rightarrow 0 \quad (n \rightarrow \infty) .
\end{equation*}
\end{definition}

\begin{proposition}
为了 $T:(\mathscr{X}, \rho) \rightarrow (\mathscr{Y}, r)$ 是连续的，必须且只须 $\forall \varepsilon>0, \forall x_0 \in \mathscr{X}$，存在 $\delta=\delta(x_0, \varepsilon)>0$ ，使得
\begin{equation*}
\rho\left(x, x_0\right) < \delta \Rightarrow r\left(T x, T x_0\right) < \varepsilon \quad (\forall x \in \mathscr{X}) .
\end{equation*}
\end{proposition}

\begin{proof}
\textbf{必要性} 若上式不成立，则必存在 $x_0 \in \mathscr{X}$ 和 $\varepsilon > 0$，使得 $\forall n \in \mathbb{N}$，存在 $x_n$ 使得 $\rho(x_n, x_0) < \frac{1}{n}$，但 $r(T x_n, T x_0) \geqslant \varepsilon$。因此，$\lim\limits_{n \rightarrow \infty} \rho(x_n, x_0) = 0$，但 $\lim\limits_{n \rightarrow \infty} r(T x_n, T x_0) \neq 0$，矛盾。\\
\textbf{充分性} 假设 上式 成立，且 $\lim\limits_{n \rightarrow \infty} \rho(x_n, x_0) = 0$，则 $\forall \varepsilon > 0$，存在 $N = N(\delta(x_0, \varepsilon))$，使得当 $n > N$ 时，$\rho(x_n, x_0) < \delta$。从而 $r(T x_n, T x_0) < \varepsilon$，即得 $\lim\limits_{n \rightarrow \infty} r(T x_n, T x_0) = 0$。
\end{proof}
\begin{definition}
称 $T:(\mathscr{X}, \rho) \rightarrow (\mathscr{X}, \rho)$ 是一个压缩映射，如果存在 $0 < \alpha < 1$，使得
\begin{equation*}
\rho(T x, T y) \leqslant \alpha \rho(x, y) \quad (\forall x, y \in \mathscr{P}) .
\end{equation*}
\end{definition}

\begin{example}
设 $\mathscr{X}= [0,1]$，$T(x)$ 是 $[0,1]$ 上的一个可微函数，满足条件：
\[
T(x) \in [0,1] \quad (\forall x \in [0,1]), \quad \text{以及} \quad \left| T^{\prime}(x) \right| \leqslant a < 1 \quad (\forall x \in [0,1]),
\]
则映射 $T: \mathscr{X} \rightarrow \mathscr{X}$ 是一个压缩映射。
\end{example}

\begin{proof}
\begin{equation*}
\begin{aligned}
& \rho(T x, T y) = |T(x) - T(y)| \\
& \quad = \left| T^{\prime}(\theta x + (1 - \theta) y) (x - y) \right| \\
& \quad \leqslant a |x - y| = a \rho(x, y) \quad (\forall x, y \in \mathscr{X}, 0 < \theta < 1) .
\end{aligned}
\end{equation*}
\end{proof}


\begin{theorem}[Banach不动点定理——压缩映象原理]
设 $(\mathscr{X}, \rho)$ 是一个完备的距离空间，$T$ 是 $(\mathscr{X}, \rho)$ 到其自身的一个压缩映射，则 $T$ 在 $\mathscr{X}$ 上存在唯一的不动点。
\end{theorem}
\begin{proof}
若 $T:(\mathscr{X}, \rho) \rightarrow (\mathscr{X}, \rho)$ 是一个压缩映射，则任取初始点 $x_0 \in \mathscr{X}$。考察迭代产生的序列
\begin{equation*}
x_{n+1} = T x_n \quad (n = 0, 1, 2, \cdots) .
\end{equation*}
我们有
\begin{equation*}
\begin{aligned}
\rho(x_{n+1}, x_n) &= \rho(T x_n, T x_{n-1}) \\
&\leqslant \alpha \rho(x_n, x_{n-1}) \leqslant \cdots \leqslant \alpha^n \rho(x_1, x_0) .
\end{aligned}
\end{equation*}
从而对 $\forall p \in \mathbb{N}$,
\begin{equation*}
\begin{aligned}
\rho(x_{n+p}, x_n) &\leqslant \sum_{i=1}^p \rho(x_{n+i}, x_{n+i-1}) \\
&\leqslant \frac{\alpha^n}{1 - \alpha} \rho(x_1, x_0) \rightarrow 0
\end{aligned}
\end{equation*}
此可见 $\{x_n\}$ 是一个基本列，设其收敛到$x_0$，由于$T$是连续映射，所以$x_0=\lim\limits_{n\to \infty}x_n=\lim\limits_{n\to \infty}Tx_{n-1}=Tx_0$。\\
若 $x^*, x^{* *}$ 都是不动点, 则
\begin{equation*}
\begin{aligned}
 \left|x^*-x^{* *}\right|  =\left|T x^*-T x^{* *}\right| < \alpha\left|x^*-x^{* *}\right|
\end{aligned}
\end{equation*}
由此推出
\begin{equation*}
x^*=x^{* *} .
\end{equation*}
\end{proof}
\begin{example}
常微分方程的初值问题的局部存在唯一性。
\begin{equation*}
x(t)=\xi+\int_0^t F(\tau, x(\tau)) \dif \tau,
\end{equation*}


\end{example}
\begin{proof}
先在 $C([-h, h])$ 上考察映射 $Tx(t)=\xi+\int_0^t F(\tau, x(\tau)) \dif \tau$，注意到
\begin{equation*}
\begin{aligned}
\rho(T x, T y) &= \max_{t \in [-h, h]} \left| \int_0^t F(\tau, x(\tau)) \, d\tau - \int_0^t F(\tau, y(\tau)) \, d\tau \right| \\
& \leqslant h \max_{t < h} |F(t, x(t)) - F(t, y(t))| .
\end{aligned}
\end{equation*}
因此，假设二元函数 $F(t, x)$ 对变元 $x$ 关于 $t$ 一致地满足局部 Lipschitz 条件：$\exists \delta > 0, L > 0$，使得当 $|t| \leqslant h, |x_1 - \xi| \leqslant \delta, |x_2 - \xi| \leqslant \delta$ 时有
\begin{equation*}
|F(t, x_1) - F(t, x_2)| \leqslant L |x_1 - x_2| ,
\end{equation*}
这时就有 $\rho(T x, T y) < L h \, \rho(x, y) \quad (\forall x, y \in \bar{B}(\xi, \delta))$，其中
\[
\bar{B}(\xi, \delta) \triangleq \{ x(t) \in C([-h, h]) \mid \max_{t \leqslant h} | x(t) - \xi | \leqslant \delta \} .
\]
在这里我们不能直接取 $C([-h, h])$ 为定理中的距离空间 $\mathscr{X}$，因为当 $L h < 1$ 时，$T$ 只是在 $C([-h, h])$ 的子集 $\bar{B}(\xi, \delta)$ 上才是压缩的。在这里我们把常数 $\xi$ 看成是 $[-h, h]$ 上恒等于 $\xi$ 的常值函数。\\
我们取 $\mathscr{X} = \bar{B}(\xi, \delta)$，为了要使 $T: \mathscr{X} \rightarrow \mathscr{X}$，再设
\begin{equation*}
M \triangleq \max \{ |F(t, x)| \mid (t, x) \in [-h, h] \times [\xi - \delta, \xi + \delta] \} ,
\end{equation*}
取 $h > 0$ 足够小，以使
\begin{equation*}
\max |(T x)(t) - \xi| = \max \left| \int_0^t F(t, x(\tau)) \, d\tau \right| \leqslant M h \leqslant \delta .
\end{equation*}
由于 $(C([-h, h]), \rho)$ 是一个完备的距离空间，而 $\mathscr{X}$ 又是它的一个闭子集，因此 $(\mathscr{X}, \rho)$ 还是一个完备的距离空间。
\end{proof}
\begin{theorem}
设函数 $F(t, x)$ 在 $[-h, h] \times[\xi-\delta, \xi+\delta]$ 上定 义、连续并满足Lipschitz条件，则当 $h<\min \{\delta / M, 1 / L\}$ 时，初值问题 (在 $[-h, h]$ 上存在唯 一解。
\end{theorem}
\begin{example}
 设 \(f(x, y) = \left(f_1(x, y), \cdots, f_m(x, y)\right) : \mathbb{R}^n \times \mathbb{R}^m \rightarrow \mathbb{R}^m, U \times V \subset \mathbb{R}^n \times \mathbb{R}^m\) 是 \( (x_0, y_0) \) 在 \( \mathbb{R}^n \times \mathbb{R}^m \) 中的一个邻域。设 \( f(x, y) \) 及
\begin{equation*}
\frac{\partial f(x, y)}{\partial y} = \left( \begin{array}{ccc}
\frac{\partial f_1}{\partial y_1}(x, y) & \cdots & \frac{\partial f_1}{\partial y_m}(x, y) \\
\vdots & \ddots & \vdots \\
\frac{\partial f_m}{\partial y_1}(x, y) & \cdots & \frac{\partial f_m}{\partial y_m}(x, y)
\end{array} \right)
\end{equation*}
在 \( U \times V \) 内连续，又设
\begin{equation*}
f(x_0, y_0) = 0, \quad \left[\operatorname{det}\left(\frac{\partial f}{\partial y}(x, y)\right)\right](x_0, y_0) \neq 0
\end{equation*}
则 \( \exists \) \( (x_0, y_0) \) 的一个邻域 \( U_0 \times V_0 \subset U \times V \) 以及唯一的连续函数 \( \varphi: U_0 \rightarrow V_0 \)，满足
\begin{equation*}
\left\{
\begin{array}{l}
f(x, \varphi(x)) = 0 \quad \text{(当 } x \in U_0), \\
\varphi(x_0) = y_0 .
\end{array}
\right.
\end{equation*}
\end{example}
\begin{proof}
用 \( \left(\frac{\partial f(x_0, y_0)}{\partial y}\right)^{-1} \cdot f(x, y)^T   = g(x, y)^T \), 可设
\begin{equation*}
\frac{\partial g}{\partial y}(x_0, y_0) = \left( \begin{array}{ccc}
1 & \cdots & 0 \\
\vdots & \ddots & \vdots \\
0 & \cdots & 1
\end{array} \right)
\end{equation*}
为恒等矩阵, 其中
\begin{equation*}
g(x, y)^T = \left( \begin{array}{c}
g_1(x, y) \\
\vdots \\
g_m(x, y)
\end{array} \right)
\end{equation*}
考虑映射 \( T: \varphi \mapsto T \varphi \),
\begin{equation*}
(T \varphi)(x) = \varphi(x) - g(x, \varphi(x))
\end{equation*}
其中 \( \varphi \in C\left(\bar{B}(x_0, r), \mathbb{R}^m\right) \), 这里 \( r > 0 \), \( C\left(\bar{B}(x_0, r), \mathbb{R}^m\right) \) 表示定义在闭球 \( \bar{B}(x_0, r) \) 上取值于 \( \mathbb{R}^m \) 上的向量值连续函数空间，其距离定义为
\begin{equation*}
\rho(\varphi, \psi) \triangleq \max_{\substack{x \in \bar{B}(x_0, r) \\ 1 \leqslant i \leqslant m}} \left|\varphi_i(x) - \psi_i(x)\right|,
\end{equation*}
其中 \( \varphi = (\varphi_1, \cdots, \varphi_m), \psi = (\psi_1, \cdots, \psi_m) \). 因为假设 \( \frac{\partial g(x, y)}{\partial y} \) 在 \( U \times V \) 上连续, 所以 \( \exists \delta > 0 \), 使得
\begin{equation*}
\begin{gathered}
\left|\delta_{ij} - \left[\frac{\partial g(x, y)}{\partial y}\right]_{ij}\right| < \frac{1}{2m} \\
i, j = 1, 2, \cdots, m, \quad x \in \bar{B}(x_0, \delta), \quad y \in \bar{B}(y_0, \delta)
\end{gathered}
\end{equation*}
其中 \( \left[ \cdot \right]_{ij} \) 表示括号内矩阵的第 \(i\) 行，第 \(j\) 列元素, \( \delta_{ij} = 0 \) 当 \( i \neq j \), \( \delta_{ij} = 1 \) 当 \( i = j \). 记 \( d_i(x) = \varphi_i(x) - \psi_i(x), i = 1, 2, \cdots, m \)，根据微分中值定理，
\begin{equation*}
\begin{aligned}
\rho(T \varphi, T \psi) & = \max_{\substack{x \in \bar{B}(x_0, r) \\ 1 \leqslant i \leqslant m}} \left|d_i(x) - g_i(x, \varphi(x)) + g_i(x, \psi(x))\right| \\
& = \max_{\substack{x \in \bar{B}(x_0, r) \\ 1 \leqslant i \leqslant m}} \left|d_i(x) - \sum_{j=1}^m \frac{\partial g_i(x, \widehat{y}(x))}{\partial y_j} d_j(x)\right|\\
&\leqslant \frac{1}{2} \max _{\substack{x \in \bar{B}\left(x_0, r\right) \\ 1 \leqslant i \leqslant m}}\left|d_i(x)\right|=\frac{1}{2} \rho(\varphi, \psi),
\end{aligned}
\end{equation*}
其中 \( \varphi, \psi \in \mathscr{X}, \widehat{y}(x) = (\widehat{y}_1(x), \widehat{y}_2(x), \cdots, \widehat{y}_m(x)) \), 这里
\begin{equation*}
\begin{aligned}
\mathscr{X} \triangleq & \left\{ \varphi \in C\left(\bar{B}(x_0, r), \mathbb{R}^m\right) \mid \varphi(x_0) = y_0, \right. \\
& \left. \varphi(x) \in \bar{B}(y_0, \delta), \, x \in \bar{B}(x_0, r) \right\} \\
\widehat{y}_i(x) = & \theta_i(x) \varphi(x) + (1 - \theta_i(x)) \psi(x) \\
& 0 < \theta_i(x) < 1, \quad x \in \bar{B}(x_0, r)
\end{aligned}
\end{equation*}
因为 \( \mathscr{X} \) 在 \( C\left(\bar{B}(x_0, r), \mathbb{R}^m\right) \) 中闭，从而是一个完备度量空间。 上式表明， \( T: \mathscr{X} \rightarrow C\left(\bar{B}(x_0, r), \mathbb{R}^m\right) \) 是压缩映射。剩下来只要再证 \( T: \mathscr{X} \rightarrow \mathscr{X} \) 就够了，因为根据压缩映射原理， \( T \) 存在唯一的不动点，这就是我们所要证的。注意到
\begin{equation*}
\begin{aligned}
\rho(T \varphi, y_0) & \leqslant \rho(T \varphi, T y_0) + \rho(T y_0, y_0) \\
& \leqslant \frac{1}{2} \rho(\varphi, y_0) + \max_{\substack{x \in \bar{B}(x_0, r) \\ 1 \leqslant i \leqslant m}} \left|f_i(x, y_0)\right|
\end{aligned}
\end{equation*}
这里把 \( y_0 \) 看成映射： \( x \mapsto y_0, x \in \bar{B}(x_0, r) \)，此时有 \( y_0 \in \mathscr{X} \)。又由 \( g \) 的连续性， \( \exists \eta > 0 \) ，使得
\begin{equation*}
\max_{\substack{x \in \bar{B}(x_0, r) \\ 1 \leqslant i \leqslant m}} \left|f_i(x, y_0)\right| < \frac{\delta}{2}
\end{equation*}
因此, 当 \( 0 < r < \min\{\eta, \delta\} \) 时, \( \rho(T y_0, y_0) < \frac{\delta}{2} \),
\begin{equation*}
\rho(T \varphi, y_0) \leqslant \frac{1}{2} \rho(\varphi, y_0) + \frac{1}{2} \delta \leqslant \frac{1}{2} \delta + \frac{1}{2} \delta = \delta, \quad \varphi \in \mathscr{X}.
\end{equation*}
此外，还有
\begin{equation*}
\begin{aligned}
(T \varphi)(x_0) & = \varphi(x_0) - g(x_0, \varphi(x_0)) \\
& = y_0 - g(x_0, y_0), \quad \varphi \in \mathscr{X}.
\end{aligned}
\end{equation*}
所以 \( T: \mathscr{X} \rightarrow \mathscr{X} \).
\end{proof}
\begin{example}
如果存在正整数 \( n \)，使得 \( T^n \) 是压缩映射，那么 \( T \) 有唯一不动点。
\end{example}

\begin{proof}
根据 Banach 不动点定理，存在 \( x_0 \) 使得 \( T^n x_0 = x_0 \)。则有
\[
T^n(T x_0) = T^{n+1} x_0 = T(T^n x_0) = T x_0.
\]
可知 \( T x_0 \) 也是 \( T^n \) 的不动点。由于压缩映射 \( T^n \) 的不动点唯一性，得出 \( T x_0 = x_0 \)。
这就证明了 \( T \) 有不动点。\\
下面再证明 \( T \) 的不动点是唯一的。我们用反证法。如果 \( x_1, x_2 \) 是 \( T \) 的两个不动点，且 \( x_1 \neq x_2 \)，即有
\[
\begin{cases}
T x_1 = x_1, \\
T x_2 = x_2,
\end{cases}
\]
那么
\[
\begin{aligned}
T^n x_1 &= T^{n-1}(T x_1) =T T^{n-1}(x_1) = \cdots = T x_1 = x_1, \\
T^n x_2 &= T^{n-1}(T x_2) =TT^{n-1}(x_2) = \cdots = T x_2 = x_2.
\end{aligned}
\]
即 \( x_1, x_2 \) 是 \( T^n \) 的两个不动点。由于 \( T^n \) 是压缩映射，它有唯一不动点，因此 \( x_1 = x_2 \)，这与假设 \( x_1 \neq x_2 \) 矛盾。\\
因此，\( T \) 的不动点是唯一的。
\end{proof}
\begin{example}
(1) 证明完备度量空间的闭子集是一个完备的子空间，而任一度量空间的完备子空间必是闭子集。\\
(2) 设 \( (\mathscr{X}, \rho) \) 是一个完备距离空间，\( M \) 是 \( X \) 的闭子集，且有 \( r_n > 0 \)，且
\[
\sum_{n=1}^{\infty} r_n < +\infty,
\]
证明：如果映射 \( T: M \to M \) 满足
\[
\rho(T^n x, T^n y) \leqslant r_n \rho(x, y), \quad x, y \in M,
\]
那么映射 \( T \) 必有唯一的不动点，并且对于任意 \( x_0 \in M \)，序列 \( T^n x_0 \) 收敛到该不动点，换句话说，该不动点可由迭代法产生。
\end{example}

\begin{proof}
(1) 设 \( \mathscr{X} \) 是完备度量空间，\( M \subset \mathscr{X} \) 是闭的。要证明 \( M \) 是一个完备的子空间。考虑
\[
\begin{aligned}
& \forall x_m, x_n \in M, \quad \left\| x_m - x_n \right\| \to 0 \quad (m, n \to \infty) \\
& \quad \Longrightarrow x_m, x_n \in \mathscr{X}, \quad \left\| x_m - x_n \right\| \to 0 \quad (m, n \to \infty),
\end{aligned}
\]
因为 \( \mathscr{X} \) 是完备度量空间，所以存在 \( x \in \mathscr{X} \)，使得 \( x_n \to x \)，从而有
\[
\begin{array}{l}
x_n \in M, x_n \to x \\
M \subset X \text{ 是闭的}
\end{array}
\Rightarrow x \in M.
\]
因此，\( \forall x_m, x_n \in M, \left\| x_m - x_n \right\| \to 0 \quad (m, n \to \infty) \)，必有存在 \( x \in M \)，使得 \( x_n \to x \)，从而证明了 \( M \) 是一个完备的子空间。\\
下设 $\mathscr{X}$ 是一度量空间, $M$ 是 $\mathscr{X}$ 的一个完备子空间. 要证 $M$ 是闭子集，即：若 $x_n \in M, x_n \rightarrow x$ ，要证 $x \in M$ 。事实上，因为收敛列是基本列，所以有
\begin{equation*}
x_n \in M, \quad\left\|x_m-x_n\right\| \rightarrow 0 \quad(m, n \rightarrow \infty) .
\end{equation*}
又 $M$ 是完备度量空间，所以 $\exists x^{\prime} \in M$ ，使得 $x_n \rightarrow x^{\prime}$ ，即有
\begin{equation*}
\left.\begin{array}{l}
x_n \rightarrow x \\
x_n \rightarrow x^{\prime}
\end{array}\right\} \Rightarrow x=x^{\prime} \in M .
\end{equation*}
(2) 设 \( (\mathscr{X}, \rho) \) 是一个完备度量空间，\( M \) 是 \( \mathscr{X} \) 的闭子集，故 \( (M, \rho) \) 也是一个完备度量空间。\\
由
\[
\sum_{n=1}^{\infty} r_n < +\infty \Longrightarrow r_n \to 0 \quad (n \to \infty) \Longrightarrow \exists n_0 \text{，使得 } r_{n_0} < 1.
\]
由
\[
\rho(T^{n_0} x, T^{n_0} y) \leqslant r_{n_0} \rho(x, y), \quad x, y \in M,
\]
即 \( T^{n_0} \) 是压缩映射，根据上例的结论，\( T \) 有唯一不动点。设该不动点为 \( x^* \)，即有
\[
T x^* = x^* = T^{n_0} x^*.
\]
进一步，对 \( \forall x_0 \in M \)，记 \( x_n = T^n x_0 \) (对于 \( n = 1, 2, \cdots \))，下面证明 \( \{ x_n = T^n x_0 \} \) 是基本列。考虑
\[
\rho(x_{n+1}, x_n) = \rho(T^{n+1} x_0, T^n x_0) = \rho(T^n(T x_0), T^n x_0) \leqslant r_n \rho(T x_0, x_0).
\]
根据三角不等式，对于任意的自然数 \( p \)，有
\[
\begin{aligned}
\rho(x_{n+p}, x_n) &\leqslant \sum_{i=1}^p \rho(x_{n+i}, x_{n+i-1}) \\
&\leqslant \rho(T x_0, x_0) \sum_{i=1}^p r_{n+i-1} \\
&= \rho(T x_0, x_0) \sum_{k=n}^{n+p-1} r_k \leqslant \rho(T x_0, x_0) \sum_{k=n}^{\infty} r_k \to 0, \quad n \to \infty.
\end{aligned}
\]
对任意的自然数 \( p \) 一致成立。因此，\( \{ x_n = T^n x_0 \} \) 是基本列。由 \( (M, \rho) \) 的完备性，必存在 \( \bar{x} \in M \)，使得 \( x_n \to \bar{x} \)。这样，一方面，因为 \( T^{n_0} \) 是压缩映射，存在唯一不动点 \( x^* \)，所以对于 \( \forall x_0 \in M \)，有
\[
T^{n_0 n} x_0 = (T^{n_0})^n x_0 \to x^* \quad (n \to \infty);
\]
另一方面，
\[
T^{n_0 n} x_0 = x_{n_0 n} \to \bar{x} \quad (n \to \infty),
\]
故有 \( \bar{x} = x^* \)。
\end{proof}

\begin{example}
设 \( M \) 是 \( \mathbb{R}^n \) 中的有界闭集，映射 \( T: M \to M \) 满足
\[
\rho(T x, T y) < \rho(x, y) \quad (\forall x, y \in M, x \neq y).
\]
求证 \( T \) 在 \( M \) 中存在唯一的不动点。
\end{example}

\begin{proof}
因为
\[
\rho(T x, T x_0) < \rho(x, x_0),
\]
所以
\[
\rho(x, x_0) \to 0 \Longrightarrow \rho(T x, T x_0) \to 0.
\]
再由三角不等式，得到
\[
\left|\rho(x, T x) - \rho(x_0, T x_0)\right| \leqslant \rho(x, x_0) + \rho(T x, T x_0).
\]
由此可见，定义 \( f(x) := \rho(x, T x) \) 在 \( M \) 上连续。因为 \( M \) 是 \( \mathbb{R}^n \) 中的有界闭集，所以存在 \( x_0 \in M \)，使得
\[
\rho(x_0, T x_0) = f(x_0) = \min_{x \in M} f(x) = \min_{x \in M} \rho(x, T x).
\]
如果 \( \rho(x_0, T x_0) = 0 \)，那么 \( x_0 \) 就是不动点。假设 \( \rho(x_0, T x_0) > 0 \)。根据假设，我们有
\[
\rho(T x_0, T^2 x_0) < \rho(x_0, T x_0) = \min_{x \in M} \rho(x, T x).
\]
但是 \( T x_0, T^2 x_0 \in M \)，这与 \( \rho(x_0, T x_0) \) 是最小值矛盾。因此，\( \rho(x_0, T x_0) = 0 \)，即存在不动点 \( x_0 \)。\\
不动点的唯一性是显然的。事实上，如果存在两个不动点 \( x_1, x_2 \)，则从
\[
\rho(x_1, x_2) = \rho(T x_1, T x_2) < \rho(x_1, x_2),
\]
即得矛盾。
\end{proof}
\begin{example}
设 \( (\mathscr{X}, \rho) \) 是一个完备距离空间，\( M \) 是 \( \mathscr{X} \) 的闭子集，且 \( 0 < \alpha < 1 \)，映射 \( T_n: M \to M \) 满足
\[
\begin{gathered}
\rho\left(T_n x, T_n y\right) \leqslant \alpha \rho(x, y) \quad (x, y \in M), \\
T x = \lim_{n \to \infty} T_n x \quad (x \in M),
\end{gathered}
\]
并且 \( x_n = T_n x_n \in M \)。证明：映射 \( T \) 有唯一的不动点，并且 \( x_n \) 收敛到该不动点。
\end{example}

\begin{proof}
(1) 由距离的连续性，对于不等式
\[
\rho\left(T_n x, T_n y\right) \leqslant \alpha \rho(x, y) \quad (x, y \in M),
\]
令 \( n \to \infty \) 取极限，得到
\[
\rho(T x, T y) \leqslant \alpha \rho(x, y) \quad (x, y \in M).
\]
由此可推出 \( T \) 是压缩映射，因此 \( T \) 有不动点，设该不动点为 \( x^* \)，即有
\[
T x^* = x^*.
\]
(2) 考虑
\[
\begin{aligned}
\rho\left(x_n, x^*\right) &= \rho\left(T_n x_n, T x^*\right) \\
&\leqslant \rho\left(T_n x_n, T_n x^*\right) + \rho\left(T_n x^*, T x^*\right) \\
&\leqslant \alpha \rho\left(x_n, x^*\right) + \rho\left(T_n x^*, T x^*\right),
\end{aligned}
\]
由此得
\[
\rho\left(x_n, x^*\right) \leqslant \frac{1}{1-\alpha} \rho\left(T_n x^*, T x^*\right).
\]
又因为 \( T x^* = \lim\limits_{n \to \infty} T_n x^* \)，所以
\[
\lim_{n \to \infty} \rho\left(x_n, x^*\right) = 0 \Longrightarrow x_n \to x^* \quad (n \to \infty).
\]
\end{proof}

\section{完备空间}
\begin{definition} 设 \( (\mathscr{X}, \rho),\left( \mathscr{X}_1, \rho_1\right) \) 是两个度量空间，如果存在映射 \( \varphi: \mathscr{X} \rightarrow \mathscr{X}_1 \) 满足
\begin{enumerate}
\item \( \varphi \) 是满射,
\item \( \rho(x, y) = \rho_1(\varphi x, \varphi y) \quad(\forall x, y \in \mathscr{X}) \),
\end{enumerate}
则称 \( (\mathscr{X}, \rho) \) 和 \( \left( \mathscr{X}_1, \rho_1\right) \) 是等距同构的，并称 \( \varphi \) 为等距同构映射，有时简称等距同构。


\end{definition}
\begin{note}
由2导出 \( \varphi \) 还是单射。\\
显然，凡是等距同构的度量空间，它们的一切与距离相联系的性质都是一样的。因此今后我们将不再区分它们。如果度量空间 \( \left( \mathscr{X}_1, \rho_1\right) \) 与另一个度量空间 \( \left( \mathscr{X}_2, \rho_2\right) \) 的子空间 \( \left( \mathscr{X}_0, \rho_2\right) \) 是等距同构的，我们就说 \( \left( \mathscr{X}_1, \rho_1\right) \) 可以嵌入 \( \left( \mathscr{X}_2, \rho_2\right) \)。在上述意义下，我们认为 \( \left( \mathscr{X}_1, \rho_1\right) \) 就是 \( \left( \mathscr{X}_2, \rho_2\right) \) 的一个子空间，并简单地记作
\begin{equation*}
\left( \mathscr{X}_1, \rho_1\right) \subset \left( \mathscr{X}_2, \rho_2\right).
\end{equation*}
\end{note}

\begin{definition}
 \( (\mathscr{X}, \rho) \) 是度量空间。集合 \( E \subset \mathscr{X} \) 叫作在 \( \mathscr{X} \) 中的稠密子集，如果 \( \forall x \in \mathscr{X} ， \forall \varepsilon>0, \exists z \in E \) ，使得 \( \rho(x, z)<\varepsilon \)。换句话说: \( \forall x \in \mathscr{X} , \exists\left\{x_n\right\} \subset E \) ，使得 \( x_n \rightarrow x(n \rightarrow \infty) \)。
\end{definition}

\begin{example}
 \([a, b]\) 上的多项式全体记为 \( P[a, b] \)。根据 Weierstrass 定理可知 \( P[a, b] \) 在 \( C[a, b] \) 中稠密。
\end{example}

\begin{definition}
包含给定度量空间 \( (\mathscr{X}, \rho) \) 的最小的完备度量空间称为 \( \mathscr{X} \) 的完备化空间，其中最小的含义是：任何一个以 \( (\mathscr{X}, \rho) \) 为子空间的完备度量空间都以此空间为子空间。
\end{definition}
\begin{proposition}
如果 \( \left( \mathscr{X}_1, \rho_1 \right) \) 是一个以 \( (\mathscr{X}, \rho) \) 为子空间的完备度量空间，\( \left. \rho_1 \right|_{ \mathscr{X} \times \mathscr{X} } = \rho \)，并且 \( \mathscr{X} \) 在 \( \mathscr{X}_1 \) 中稠密，则 \( \mathscr{X}_1 \) 是 \( \mathscr{X} \) 的完备化空间。
\end{proposition}
\begin{proof}
     事实上， \( \forall \xi \in \mathscr{X}_1, \exists x_n \in \mathscr{X} \)，使得 \( \rho_1\left(x_n, \xi\right) \rightarrow 0(n \rightarrow \infty) \)。如果存在 \( \left( \mathscr{X}_2, \rho_2 \right) \) 以 \( (\mathscr{X}, \rho) \) 为子空间，因为
\begin{equation*}
\rho_2\left(x_n, x_m\right) = \rho_1\left(x_n, x_m\right) \rightarrow 0 \quad (n, m \rightarrow \infty),
\end{equation*}
所以 \( \exists \widehat{\xi} \in \mathscr{X}_2 \)，使得 \( \rho_2\left(x_n, \widehat{\xi}\right) \rightarrow 0 \)。做映射 \( T: \mathscr{X}_1 \rightarrow \mathscr{X}_2, T \xi = \widehat{\xi} \)。我们还要证 \( T \) 是等距的。因为 \( \forall \eta \in \mathscr{X}_1 \)，又 \( \exists y_n \in \mathscr{X} \)，使得 \( \rho_1\left(y_n, \eta\right) \rightarrow 0 \)，所以
\begin{equation*}
\rho_1(\xi, \eta) = \lim_{n \rightarrow \infty} \rho_1\left(x_n, y_n\right) = \lim_{n \rightarrow \infty} \rho_2\left(x_n, y_n\right) = \rho_2(\widehat{\xi}, \widehat{\eta})
\end{equation*}
这表明 \( \left( \mathscr{X}_1, \rho_1 \right) \) 是 \( \left( \mathscr{X}_2, \rho_2 \right) \) 的一个子空间。
\end{proof}
\begin{theorem}
每一个度量空间都有一个完备化空间。
\end{theorem}

\begin{proof}
设 \( (\mathscr{X}, \rho) \) 是一个度量空间，分三步证明它有一个完备化空间。\\
（1）将 \( \mathscr{X} \) 中的基本列分类。凡是满足
\[
\lim _{n \rightarrow \infty} \rho(x_n, y_n) = 0
\]
的两个基本列 \( \{x_n\}, \{y_n\} \) 称为等价的。彼此等价的基本列归于同一类且只归一类，称为等价类。我们把一个等价类看成是一个元素，并用 \( \mathscr{X}_1 \) 表示一切这种元素（等价类）组成的集合。在 \( \mathscr{X}_1 \) 上定义距离：\( \forall \xi, \eta \in \mathscr{X}_1 \)，任取 \( \{x_n\} \in \xi, \{y_n\} \in \eta \)，令
\[
\rho_1(\xi, \eta) = \lim_{n \rightarrow \infty} \rho(x_n, y_n).
\]
容易验证上式右端的极限的确存在($\{\rho(x_n,y_m)\}$为柯西列)，并且极限值与 \( \{x_n\}, \{y_n\} \) 的选取无关。为了验证该式定义的 \( \rho_1 \) 的确是个距离，注意到定义1.1中的条件(1)和(2)是显然的，而条件(3)可由
\[
\rho(x_n, y_n) \leqslant \rho(x_n, z_n) + \rho(z_n, y_n)
\]
取极限得到，其中 \( \{x_n\}, \{y_n\}, \{z_n\} \) 是分别属于等价类 \( \xi, \eta, \zeta \) 的基本列。这样，我们就证明了 \( (\mathscr{X}_1, \rho_1) \) 是一个度量空间。\\
（2）对于每个 \( x \in \mathscr{X} \)，我们用 \( \xi_x \in \mathscr{X}_1 \) 表示包含序列 \( (x, x, \cdots, x, \cdots) \) 的等价类。这样的 \( \xi_x \) 全体记为 \( \mathscr{X}^\prime \)。显然 \( \mathscr{X}^\prime \subset \mathscr{X}_1 \)，且映射 \( T : x \mapsto \xi_x \) 作为 \( (\mathscr{X}, \rho) \rightarrow (\mathscr{X}^\prime, \rho_1) \) 的映射满足定义1.7中的条件(1)和(2)。因此，\( (\mathscr{X}, \rho) \) 和 \( (\mathscr{X}^\prime, \rho_1) \) 等距同构，即有 \( (\mathscr{X}, \rho) \subset (\mathscr{X}_1, \rho_1) \)。进一步容易验证 \( \mathscr{X} \) 在 \( \mathscr{X}_1 \) 中稠密($\forall \varepsilon>0,\exists n,N\in\mathbb{N},\forall n>N ,\rho(x_n,x_N)<\varepsilon$，取$\{x_N\}\in \mathscr{X}$)。\\
（3）证明 \( (\mathscr{X}_1, \rho_1) \) 是完备的。设 \( \{\xi^{(n)}\} \) 是 \( \mathscr{X}_1 \) 中的基本列。要证 \( \exists \xi \in \mathscr{X}_1 \)，使得
\[
\rho_1(\xi^{(n)}, \xi) \rightarrow 0 \quad (n \rightarrow \infty).
\]
\( 1^{\circ} \) 先证特殊情形，假定 \( \{\xi^{(n)}\} \subset \mathscr{X}^\prime \)。令 \( x_n = T^{-1} \xi^{(n)} \)，则 \( \{x_n\} \) 是 \( \mathscr{X} \) 中的基本列。设 \( \{x_n\} \in \xi \)，便有 \( \xi^{(n)} \rightarrow \xi \)。\\
\( 2^{\circ} \) 再证一般情形。由于 \( \mathscr{X}^\prime \) 在 \( \mathscr{X}_1 \) 中稠密，\( \forall \xi^{(n)} \in \mathscr{X}_1 \)，\( \exists \bar{\xi}^{(n)} \in \mathscr{X}^\prime \)，使得 \( \rho_1(\bar{\xi}^{(n)}, \xi^{(n)}) < \frac{1}{n} \)。由 \( 1^{\circ} \) 可设 \( \bar{\xi}^{(n)} \rightarrow \xi \in \mathscr{X}_1 \)，即可推出 \( \xi^{(n)} \rightarrow \xi \)。\\
最后综合 (1), (2), (3) 并用上个命题即得结论。
\end{proof}

\begin{example}
设 \( P[a, b] \) 表示 \( [a, b] \) 上的所有多项式。按距离
\[
\rho(x, y) = \max_{a \leqslant t \leqslant b} |x(t) - y(t)|
\]
的完备化空间是 \( C[a, b] \)。
\end{example}

\begin{example}
 \( C[a, b] \) 按照 (1.2.1) 式定义的距离 \( \rho_1 \) 完备化，完备化空间是 \( L^1[a, b] \)。
\end{example}
\begin{example}
设 \( S \) 为一切复数列的集合
\[
x = \left(\xi_1, \xi_2, \cdots, \xi_k, \cdots\right),
\]
并在 \( S \) 中定义距离为
\[
\rho(x, y) = \sum_{k=1}^{\infty} \frac{1}{2^k} \frac{\left|\xi_k - \eta_k\right|}{1 + \left|\xi_k - \eta_k\right|},
\]
其中 \( x = \left(\xi_1, \xi_2, \cdots, \xi_k, \cdots\right), y = \left(\eta_1, \eta_2, \cdots, \eta_k, \cdots\right) \)。\\
证明：\( S \) 为一个完备的距离空间。

\end{example}

\begin{proof}
首先，验证 \( \rho(x, y) \) 满足距离的正定性和对称性显然成立。接下来，验证 \( \rho(x, y) \) 满足三角不等式。\\
考虑函数
\[
f(t) = \frac{t}{1+t} = 1 - \frac{1}{1+t} \quad \text{(单调递增)}.
\]
因此，
\[
f(|a+b|) \leqslant f(|a| + |b|),
\]
即
\[
\frac{|a+b|}{1 + |a+b|} \leqslant \frac{|a| + |b|}{1 + |a| + |b|} = \frac{|a|}{1 + |a| + |b|} + \frac{|b|}{1 + |a| + |b|} \leqslant \frac{|a|}{1 + |a|} + \frac{|b|}{1 + |b|}.
\]
设 \( z = \left(\zeta_1, \zeta_2, \cdots, \zeta_k, \cdots\right) \)，则有
\[
\begin{aligned}
\rho(x, y) &= \sum_{k=1}^{\infty} \frac{1}{2^k} \frac{\left|\xi_k - \eta_k\right|}{1 + \left|\xi_k - \eta_k\right|} \\
&= \sum_{k=1}^{\infty} \frac{1}{2^k} \frac{\left|(\xi_k - \zeta_k) + (\zeta_k - \eta_k)\right|}{1 + \left|(\xi_k - \zeta_k) + (\zeta_k - \eta_k)\right|} \\
&\leqslant \sum_{k=1}^{\infty} \frac{1}{2^k} \frac{\left|\xi_k - \zeta_k\right|}{1 + \left|\xi_k - \zeta_k\right|} + \sum_{k=1}^{\infty} \frac{1}{2^k} \frac{\left|\zeta_k - \eta_k\right|}{1 + \left|\zeta_k - \eta_k\right|} \\
&= \rho(x, z) + \rho(z, y).
\end{aligned}
\]
因此，\( \rho(x, y) \) 满足三角不等式，从而 \( S \) 是一个距离空间。\\
接下来证明 \( S \) 的完备性。设 \( \left\{x^{(m)}\right\} \) 是 \( S \) 中的基本列，其中
\[
x^{(m)} = \left(x_1^{(m)}, x_2^{(m)}, \cdots, x_k^{(m)}, \cdots\right),
\]
则有
\[
\rho\left(x^{(m+p)}, x^{(m)}\right) = \sum_{i=1}^{\infty} \frac{1}{2^i} \frac{\left|x_i^{(m+p)} - x_i^{(m)}\right|}{1 + \left|x_i^{(m+p)} - x_i^{(m)}\right|} \to 0 \quad (\text{当}  m \to \infty , \text{对于所有}  p \in \mathbb{N} ).
\]
由此可得，对于每个固定的 \( k \in \mathbb{N} \)，
\[
\left|x_k^{(m+p)} - x_k^{(m)}\right| \to 0 \quad (m \to \infty, \forall p \in \mathbb{N}).
\]
实际上，对于每个固定的 \( k \in \mathbb{N} \)，对于任意 \( \varepsilon > 0 \)，存在 \( N_k \)，使得
\[
\sum_{i=1}^{\infty} \frac{1}{2^i} \frac{\left|x_i^{(m+p)} - x_i^{(m)}\right|}{1 + \left|x_i^{(m+p)} - x_i^{(m)}\right|} < \frac{\varepsilon}{2^{k+1}} \quad (\text{当}  m > N_k ,\text{对于所有}  p \in \mathbb{N} ).
\]
取级数中的第 \( k \) 项，它不会超过所有项的和，即得
\[
\begin{aligned}
\frac{\left|x_k^{(m+p)} - x_k^{(m)}\right|}{1 + \left|x_k^{(m+p)} - x_k^{(m)}\right|} < \frac{\varepsilon}{2} \quad &\Longrightarrow \left|x_k^{(m+p)} - x_k^{(m)}\right| < \frac{\frac{\varepsilon}{2}}{1 - \frac{\varepsilon}{2}} <\varepsilon(\varepsilon<1).
\end{aligned}
\]
因此，对于每个 \( k \in \mathbb{N} \)，
\[
\left|x_k^{(m+p)} - x_k^{(m)}\right| \to 0 \quad (m \to \infty, \forall p \in \mathbb{N}).
\]
这意味着对于每个 \( k \in \mathbb{N} \)，\( x^{(m)} \) 的每一个坐标序列 \( \left\{x_k^{(m)}\right\} \) 都是复数集合中的基本列。由于复数集合是完备的，每一个坐标序列 \( \left\{x_k^{(m)}\right\} \) 都收敛，并且存在 \( x_k \)，使得
\[
\left|x_k^{(m)} - x_k\right| \to 0 \quad (m \to \infty).
\]
现在令 \( x = \left(x_1, x_2, \cdots, x_k, \cdots\right) \)。\\
接下来证明 \( x^{(m)} \xrightarrow{\rho} x \quad (m \to \infty) \)。实际上，我们要证明
\[
\rho\left(x^{(m)}, x\right) = \sum_{n=1}^{\infty} \frac{1}{2^n} \frac{\left|x_n^{(m)} - x_n\right|}{1 + \left|x_n^{(m)} - x_n\right|} \to 0 \quad (m \to \infty),
\]
即要证明，对于任意 \( \varepsilon > 0 \)，存在 \( N \)，当 \( m > N \) 时，
\[
\sum_{n=1}^{\infty} \frac{1}{2^n} \frac{\left|x_n^{(m)} - x_n\right|}{1 + \left|x_n^{(m)} - x_n\right|} < \varepsilon.
\]
为了使得无穷多项之和小于 \( \frac{\varepsilon}{2} \)，只需选取 \( n_0 > 1 - \log_2 \varepsilon \)。实际上，
\[
\sum_{n=n_0+1}^{\infty} \frac{1}{2^n} \frac{\left|x_n^{(m)} - x_n\right|}{1 + \left|x_n^{(m)} - x_n\right|} < \sum_{n=n_0+1}^{\infty} \frac{1}{2^n} = \frac{1}{2^{n_0}} < \frac{\varepsilon}{2}.
\]
对于每个 \( n \leqslant n_0 \)，存在 \( N_n \)，使得当 \( m > N_n \) 时，
\[
\left|x_n^{(m)} - x_n\right| < \frac{\varepsilon}{2} \quad (n = 1, 2, \cdots, n_0).
\]
取 \( N = \max \left\{N_1, N_2, \cdots, N_{n_0}\right\} \)，当 \( m > N \) 时，
\[
\sum_{n=1}^{n_0} \frac{1}{2^n} \frac{\left|x_n^{(m)} - x_n\right|}{1 + \left|x_n^{(m)} - x_n\right|} < \sum_{n=1}^{n_0} \frac{1}{2^n} \left|x_n^{(m)} - x_n\right| < \sum_{n=1}^{n_0} \frac{1}{2^n} \cdot \frac{\varepsilon}{2} < \sum_{n=1}^{\infty} \frac{1}{2^n} \cdot \frac{\varepsilon}{2} = \frac{\varepsilon}{2}.
\]
因此，
\[
\sum_{n=1}^{\infty} \frac{1}{2^n} \frac{\left|x_n^{(m)} - x_n\right|}{1 + \left|x_n^{(m)} - x_n\right|} < \varepsilon \quad (\forall m > N).
\]
因此，\( \left\{x^{(m)}\right\} \) 按距离 \( \rho \) 收敛于 \( x \)，故 \( (S, \rho) \) 是完备的度量空间。
\end{proof}
\begin{example}
证明 \( l^\infty \) 空间是完备的。
\end{example}

\begin{proof}
设 \( \{ x^{(n)} \} \) 是一个 Cauchy 序列，假设 \( \{ x^{(n)} \} \) 是 \( l^\infty \) 空间中的 Cauchy 序列。根据 Cauchy 序列的定义，对于任意的 \( \epsilon > 0 \)，存在 \( N \in \mathbb{N} \)，使得对于所有 \( m, n \geq N \)，都有：
\[
\| x^{(m)} - x^{(n)} \|_\infty = \sup_k |x^{(m)}_k - x^{(n)}_k| < \epsilon.
\]
由于 \( \{ x^{(n)} \} \) 是 Cauchy 序列，对于任意 \( k \)，序列 \( \{ x^{(n)}_k \} \) 也是 Cauchy 序列。因此，序列 \( \{ x^{(n)}_k \} \) 收敛到某个极限 \( x_k \)，即：
\[
\lim_{n \to \infty} x^{(n)}_k = x_k.
\]
因此，我们可以定义一个极限数列 \( x = (x_1, x_2, \dots) \)，其第 \( k \) 个元素为 \( x_k \)。\\
由于 \( \{ x^{(n)} \} \) 是 Cauchy 序列，对于任意 \( \epsilon > 0 \)，存在 \( N \) 使得对于所有 \( m, n \geq N \)，有：
\[
\| x^{(m)} - x^{(n)} \|_\infty < \epsilon.
\]
我们需要证明 \( \lim_{n \to \infty}\limits \| x^{(n)} - x \|_\infty = 0 \)。对于任意 \( \epsilon > 0 \)，存在 \( N \in \mathbb{N} \)，使得对于所有 \( m, n \geq N \)，有：
\[
\| x^{(m)} - x^{(n)} \|_\infty < \epsilon.
\]
由于 \( \{ x^{(n)}_k \} \) 收敛于 \( x_k \)，对于每个 \( k \)，存在 \( N' \) 使得对于所有 \( n \geq N' \)，有：
\[
| x^{(n)}_k - x_k | \leqslant \epsilon(m\to\infty).
\]
因此：
\[
\| x^{(n)} - x \|_\infty = \sup_k | x^{(n)}_k - x_k | \leqslant \epsilon.
\]
这证明了 \( x^{(n)} \) 收敛于 \( x \)且$x\in l^\infty$。因此，\( l^\infty \) 空间是完备的。
\end{proof}
\begin{example}
 设 \( F \) 是只有有限项不为 0 的实数列全体，在 \( F \) 上引进距离
\[
\rho(x, y) = \sup_{k \geqslant 1} \left|\xi_k - \eta_k\right|.
\]
求证： \( (F, \rho) \) 不完备，并指出它的完备化空间。

\end{example}

\begin{proof}
取点列
\[
x^{(n)} = \left( 1, \frac{1}{2}, \frac{1}{3}, \cdots, \frac{1}{n}, 0, 0, \cdots \right) \in F, \quad (n = 1, 2, \cdots),
\]
则有
\[
x^{(m)} = \left( 1, \frac{1}{2}, \frac{1}{3}, \cdots, \frac{1}{m}, 0, 0, \cdots \right) \in F.
\]
当 \( m > n \) 时，\[
x^{(n)} - x^{(m)} = \left(\underbrace{0, 0, \cdots,0}_{n \text{ 个零}} , \frac{1}{n+1}, \cdots, \frac{1}{m} , 0, 0, \cdots  \right)
\]
因此，
\[
\rho\left( x^{(n)}, x^{(m)} \right) = \sup_{k \geqslant 1} \left| x_k^{(n)} - x_k^{(m)} \right| = \frac{1}{n+1} \rightarrow 0 \quad (n \rightarrow \infty).
\]
由此可见，点列 \( \{ x^{(n)} \} \) 是 \( F \) 上的基本列。 然而，对于点
\[
x = \left( 1, \frac{1}{2}, \frac{1}{3}, \cdots, \frac{1}{n}, \cdots \right),
\]
我们有
\[
\rho\left( x^{(n)}, x \right) = \sup_{k \geqslant 1} \left| x_k^{(n)} - x_k \right| = \frac{1}{n+1} \rightarrow 0 \quad (n \rightarrow \infty),
\]
即 \( \lim\limits_{n \rightarrow \infty} x^{(n)} = x \)，但 \( x \notin F \)。这就证明了 \( (F, \rho) \) 不完备。\\
\( F \) 的完备化空间应是以 0 为极限的实数列全体，即
\[
x = \left( \xi_1, \xi_2, \cdots, \xi_k, \cdots \right), \quad \xi_k \in \mathbb{R}, \quad \lim_{k \to \infty} \xi_k = 0,
\]
的全体，记此空间为 \( \widetilde{F} \)。\\
下面证明：\( \widetilde{F} \) 是完备的，并且 \( F \) 在 \( \widetilde{F} \) 中稠密。
（1）证明 \( \widetilde{F} \) 是完备的，只需证明 \( \widetilde{F} \) 是 \( l^\infty \) 中的闭集。\\
设 \( x^{(n)} = \left( \xi_1^{(n)}, \xi_2^{(n)}, \cdots, \xi_k^{(n)}, \cdots \right) \in \widetilde{F} \) (\( n = 1, 2, \cdots \))，并且 \( x^{(n)} \to x \)，即
\[
\begin{aligned}
x^{(1)} &= \left( \xi_1^{(1)}, \xi_2^{(1)}, \cdots, \xi_k^{(1)}, \cdots \right) \to 0,(k\to \infty) \\
x^{(2)} &= \left( \xi_1^{(2)}, \xi_2^{(2)}, \cdots, \xi_k^{(2)}, \cdots \right) \to 0,(k\to \infty)  \\
x^{(3)} &= \left( \xi_1^{(3)}, \xi_2^{(3)}, \cdots, \xi_k^{(3)}, \cdots \right) \to 0,(k\to \infty)  \\
& \vdots \\
x^{(n)} &= \left( \xi_1^{(n)}, \xi_2^{(n)}, \cdots, \xi_k^{(n)}, \cdots \right) \to 0,(k\to \infty) 
\end{aligned}
\]
要证明 \( x = \left( \xi_1, \xi_2, \cdots, \xi_k, \cdots \right) \in \widetilde{F} \)。由于 \( x^{(n)} \to x \)，我们有
\[
\lim_{n \to \infty} \sup_{k \geqslant 1} \left| \xi_k^{(n)} - \xi_k \right| = 0.
\]
因此，对于任意 \( \varepsilon > 0 \)，存在 \( N \)，使得
\[
\sup_{k \geqslant 1} \left| \xi_k^{(n)} - \xi_k \right| < \frac{\varepsilon}{3} \quad (n \geqslant N).
\]
特别地，
\[
\left| \xi_k^{(N)} - \xi_k \right| < \frac{\varepsilon}{3} \quad (k = 1, 2, \cdots).
\]
又，实数列 \( \{ \xi_k^{(N)} \}_{k=1}^\infty \) 是收敛列，因此是基本列，故存在 \( N_1 \)，使得
\[
\left| \xi_k^{(N)} - \xi_j^{(N)} \right| < \frac{\varepsilon}{3} \quad \text{对于} \quad k, j \geqslant N_1.
\]
且
\[
\{ \xi_j^{(N)} \}_{j=1}^\infty \in \widetilde{F} \quad \Longrightarrow \quad \lim_{j \to \infty} \left| \xi_j^{(N)} \right| = 0 \quad \Longrightarrow \quad \left| \xi_j^{(N)} \right| < \frac{\varepsilon}{3} \quad (j \geqslant N_2 \geqslant N_1).
\]
于是，当 \( k, j \geqslant N_2 \geqslant N_1 \) 时，
\[
\begin{aligned}
\left| \xi_k \right| & \leqslant \left| \xi_k - \xi_k^{(N)} \right| + \left| \xi_k^{(N)} - \xi_j^{(N)} \right| + \left| \xi_j^{(N)} \right| \\
& < \frac{\varepsilon}{3} + \frac{\varepsilon}{3} + \frac{\varepsilon}{3} = \varepsilon.
\end{aligned}
\]
这就证明了
\[
x = \left( \xi_1, \xi_2, \cdots, \xi_k, \cdots \right) \in \widetilde{F}.
\]
（2）证明 \( F \) 在 \( \widetilde{F} \) 中稠密。\\
任意给定 \( x = \left( \xi_1, \xi_2, \cdots, \xi_k, \cdots \right) \in \widetilde{F} \)，令
\[
x^{(n)} = \left( \underbrace{\xi_1, \xi_2, \cdots, \xi_n}_{n \text{个}}, 0, 0, \cdots \right) \in F \quad (n = 1, 2, \cdots),
\]
则有
\[
\rho\left( x^{(n)}, x \right) = \sup_{k \geqslant n+1} \left| \xi_k \right| \rightarrow 0 \quad (n \rightarrow \infty).
\]
这意味着，\( \widetilde{F} \) 中的任意点 \( x \) 可以由 \( F \) 中的点来逼近，即 \( F \) 在 \( \widetilde{F} \) 中稠密。
\end{proof}

\section{列紧}
\begin{definition}[列紧与自列紧]
设 \( \mathscr{X} \) 是一个度量空间，\( A \) 为其一子集。称 \( A \) 是列紧的，如果 \( A \) 中的任意点列在 \( \mathscr{X} \) 中有一个收敛子列。若这个子列还收敛到 \( A \) 中的点，则称 \( A \) 是自列紧的。如果空间 \( \mathscr{X} \) 是列紧的，那么称 \( \mathscr{X} \) 为列紧空间。
\end{definition}

\begin{proposition}
在 \( \mathbb{R}^n \) 中，任意有界集是列紧集，任意有界闭集是自列紧集。
\end{proposition}

\begin{proposition}
列紧空间内任意（闭）子集都是（自）列紧集。
\end{proposition}

\begin{proposition}
列紧空间必是完备空间。
\end{proposition}

\begin{proof}
设 \( \mathscr{X} \) 是一个列紧空间，\( \{x_n\} \) 是其中的一串基本列。不妨设其为一无穷点列。由列紧性，存在 \( \{x_n\} \) 的子列 \( \{y_n\} \) 收敛到 \( x_0 \in \mathscr{X} \)。于是可知，\( x_n \to x_0 \)（\( n \to \infty \)）。
\end{proof}

\begin{definition}[ε 网]
设 \( M \) 是 \( \mathscr{X} \) 中的一个子集，\( \varepsilon > 0 \)，\( N \subset M \)。如果对于 \( \forall x \in M \)，存在 \( y \in N \)，使得 \( \rho(x, y) < \varepsilon \)，那么称 \( N \) 是 \( M \) 的一个 \( \varepsilon \) 网。如果 \( N \) 还是一个有穷集（个数依赖于 \( \varepsilon \)），那么称 \( N \) 为 \( M \) 的一个有穷 \( \varepsilon \) 网。
\end{definition}

\begin{definition}[完全有界]
集合 \( M \) 称为完全有界的，如果 \( \forall \varepsilon > 0 \)，都存在着 \( M \) 的一个有穷 \( \varepsilon \) 网。
\end{definition}

\begin{definition}[完全有界]
集合 \( M \) 称为完全有界的，如果对于 \( \forall \varepsilon > 0 \)，都存在着 \( M \) 的一个有穷 \( \varepsilon \) 网。
\end{definition}

\begin{note}
由定义显然有
\[
M \subset \bigcup_{y \in N} B(y, \varepsilon)
\]
\end{note}

\begin{theorem}[Hausdorff]
为了(完备)度量空间 \( \mathscr{X} \) 中的集合 \( M \) 是列紧的，必须(且仅须) \( M \) 是完全有界集。
\end{theorem}

\begin{proof}
\textbf{必要性.} 用反证法，若存在 \( \varepsilon_0 > 0 \)，使得 \( M \) 中没有有穷的 \( \varepsilon_0 \) 网。任取 \( x_1 \in M \)，则存在 \( x_2 \in M \backslash B(x_1, \varepsilon_0) \)。\\
对于 \( \{x_1, x_2\} \in M \)，存在 \( x_3 \in M \backslash \left(B(x_1, \varepsilon_0) \cup B(x_2, \varepsilon_0)\right) \); \\
对 \( \{x_1, x_2, \cdots, x_n\} \in M \)，存在 \( x_{n+1} \in M \backslash \bigcup_{k=1}^n B(x_k, \varepsilon_0) \); \\
\(\cdots\)。\\这样产生的点列 \( \{x_n\} \subset M \) 显然满足 \( \rho(x_n, x_m) \geqslant \varepsilon_0 \) (\( n \neq m \))，它没有收敛的子列。这与 \( M \) 的列紧性矛盾。\\
\textbf{充分性.} 若 \( \{x_n\} \) 是 \( M \) 中的无穷点列，想找一个收敛子列。对 \( 1 \) 网，存在 \( y_1 \in M \)，使得 \( \{x_n\} \) 的子列 \( \{x_n^{(1)}\} \subset B(y_1, 1) \)。\\
对 \( 1/2 \) 网，存在 \( y_2 \in M \)，使得 \( \{x_n^{(1)}\} \) 的子列 \( \{x_n^{(2)}\} \subset B(y_2, 1/2) \); \\
对 \( 1/k \) 网，存在 \( y_k \in M \)，使得 \( \{x_n^{(k-1)}\} \) 的子列 \( \{x_n^{(k)}\} \subset B(y_k, 1/k) \); \\
\(\cdots\)。\\
最后抽出对角线子列 \( \{x_k^{(k)}\} \)，它是一个基本列。事实上，\( \forall \varepsilon > 0 \)，当 \( n > 2/\varepsilon \) 时，对 \( \forall p \in \mathbb{N} \) 有
\[
\begin{aligned}
\rho(x_{n+p}^{(n+p)}, x_n^{(n)}) &\leqslant \rho(x_{n+p}^{(n+p)}, y_n) + \rho(x_n^{(n)}, y_n) \\
&\leqslant \frac{2}{n} < \varepsilon
\end{aligned}
\]
这就证明了充分性。
\end{proof}
\begin{definition}[可分]
一个度量空间若有可数的稠密子集, 就称这个度量空间是可分的。
\end{definition}

\begin{theorem}
完全有界的度量空间是可分的。
\end{theorem}
\begin{proof}
取 \( N_n \) 为有穷的 \( \frac{1}{n} \) 网，则 \( \bigcup_{n=1}^{\infty} N_n \) 是一个可数的稠密子集。
\end{proof}

\begin{definition}[紧集]
在拓扑空间 \( \mathscr{X} \) 中，集合 \( M \) 称为是紧的，如果 \( \mathscr{X} \) 中每个覆盖 \( M \) 的开集族中有有穷个开集覆盖集合 \( M \)。
\end{definition}

\begin{theorem}
设 \( (\mathscr{X}, \rho) \) 是一个度量空间，为了 \( M \subset \mathscr{X} \) 是紧的，必须且仅须 \( M \) 是自列紧集。
\end{theorem}
\begin{proof}
\textbf{必要性.} 设 \( M \) 是紧集。先证 \( M \) 是闭集，只要证 \( M \) 的余集是开集。对任意 \( x_0 \in \mathscr{X} \backslash M \)，由于
\[
M \subset \bigcup_{x \in M} B\left(x, \frac{1}{2} \rho\left(x, x_0\right)\right),
\]
利用 \( M \) 的紧性，存在 \( x_k \in M \) (\( k = 1, 2, \dots, n \))，使得
\[
M \subset \bigcup_{k=1}^n B\left(x_k, \frac{1}{2} \rho\left(x_k, x_0\right)\right).
\]
取 \( \delta = \min\limits_{1 \leqslant k \leqslant n} \frac{1}{2} \rho\left(x_k, x_0\right) \)，则显然有 \( \delta > 0 \)，并且对任意 \( x \in B\left(x_0, \delta\right) \)，有
\[
\rho\left(x, x_k\right) \geqslant \rho\left(x_k, x_0\right) - \rho\left(x_0, x\right) > \frac{1}{2}\rho(x_k,x_0) \quad (k = 1, 2, \dots, n).
\]
因此，\( B\left(x_0, \delta\right) \cap M = \varnothing \)，从而 \( M \) 的余集是开集，得证 \( M \) 是闭集。\\
其次，证 \( M \) 是列紧集。用反证法，假若 \( M \) 中的点列 \( \{x_n\} \) 不含有收敛子列，不妨假定 \( x_n \) 是互异的。对每个 \( n \in \mathbb{N} \)，定义集合 \( S_n \triangleq \{x_1, x_2, \dots, x_{n-1}, x_{n+1}, x_{n+2}, \dots\} \)，显然每个 \( S_n \) 是闭集（因为不含收敛子列），从而每个 \( \mathscr{X} \backslash S_n \) 是开集。但有
\[
\bigcup_{n=1}^{\infty} \left( \mathscr{X} \backslash S_n \right) = \mathscr{X} \backslash \bigcap_{n=1}^{\infty} S_n = \mathscr{X} \backslash \varnothing = \mathscr{X} \supset M.
\]
由 \( M \) 的紧性，存在 \( N \in \mathbb{N} \)，使得
\[
\bigcup_{n=1}^N \left( \mathscr{X} \backslash S_n \right) \supset M,
\]
即得
\[
\mathscr{X} \backslash \left\{ x_n \right\}_{n=N+1}^{\infty} \supset M.
\]
但这是不可能的，因为 \( x_{N+1} \) 属于右端但不属于左端。此矛盾说明 \( M \) 是列紧的。\\
\textbf{充分性.} 设 \( M \) 是自列紧的，要在 \( M \) 的任一开覆盖中取出有限覆盖。用反证法，如果某个开覆盖 \( \bigcup_{\lambda \in \Lambda} G_\lambda \supset M \) 不能取出有限覆盖。由于 \( M \) 是自列紧集，\( \forall n \in \mathbb{N} \)，存在有穷的 \( \frac{1}{n} \) 网
\[
N_n = \{x_1^{(n)}, x_2^{(n)}, \dots, x_{k_{(n)}}^{(n)}\},
\]
显然
\[
\bigcup_{y \in N_n} B(y, 1/n) \supset M.
\]
因此，\( \forall n \in \mathbb{N}, \exists y_n \in N_n \)，使得 \( B(y_n, 1/n) \) 不能被有限个 \( G_\lambda \) 所覆盖。由假定 \( M \) 是自列紧集，必存在收敛子列 \( y_{n_k} \) 收敛到一点 \( y_0 \in G_{\lambda_0} \)。又 \( G_{\lambda_0} \) 是开集，所以 \( \exists \delta > 0 \)，使得 \( B(y_0, \delta) \subset G_{\lambda_0} \)。\\
对此 \( \delta > 0 \)，取 \( k \) 足够大，使得 \( n_k > 2/\delta \)，并且 \( \rho(y_{n_k}, y_0) < \delta/2 \)，则对任意 \( x \in B(y_{n_k}, 1/n_k) \)，有
\[
\rho(x, y_0) \leqslant \rho(x, y_{n_k}) + \rho(y_{n_k}, y_0) \leqslant \frac{1}{n_k} + \frac{\delta}{2} < \delta.
\]
即 \( x \in B(y_0, \delta) \)，从而 \( B(y_{n_k}, 1/n_k) \subset B(y_0, \delta) \subset G_{\lambda_0} \)。这与每个 \( B(y_n, 1/n) \) 不能被有限个 \( G_\lambda \) 所覆盖矛盾。
\end{proof}
\begin{definition}[连续函数空间]
设 \( M \) 是一个紧的度量空间，带有距离 \( \rho \)，用 \( C(M) \) 表示 \( M \to \mathbb{R} \) 的一切连续映射全体。定义
\[
d(u, v) = \max_{x \in M} |u(x) - v(x)| \quad (\forall u, v \in C(M)).
\]
\end{definition}

\begin{proposition}
\((C(M), d)\) 是一个度量空间。
\end{proposition}
\begin{proof}
其实只有定义本身的合理性是要验证的。即，对于 \( \forall u \in C(M) \)，存在着最大值 \( \max_{x \in M} |u(x)| \)。事实上，\( u(M) \) 是紧集。因为对任意点列 \( y_n \in u(M) \)，存在 \( x_n \in M \)，使得 \( u(x_n) = y_n \)。由于 \( M \) 是紧的，从而有子列 \( x_{n_k} \to x_0, k \to \infty \)，而 \( u \) 是连续的，便有 \( u(x_{n_k}) \to u(x_0) \in u(M) \)。令 \( y_0 \triangleq u(x_0) \)，即得 \( y_{n_k} = u(x_{n_k}) \to y_0 \)，从而 \( u(M) \) 是紧集。这蕴含 \( u(M) \) 是有界闭的数集。设
\[
\min u(M) = \alpha, \quad \max u(M) = \beta.
\]
由闭性推出 \( \alpha, \beta \in u(M) \)。这就证明了 \( \max_{x \in M} |u(x)| \) 的存在性。
\end{proof}

\begin{proposition}
\((C(M), d)\) 是完备的。
\end{proposition}
\begin{proof}
设$\{u_n\}$是$C(M)$中的柯西列，则由实数的完备性知$u_n$一致收敛到$u(x)$，接着只需验证$u(x)$是连续函数。$\forall x_0\in M,\forall \varepsilon>0,\exists N>0$s.t.$\max|u(x)-u_n(x)|<\frac{\varepsilon}{3}(n>N)$，由于$u_{N+1}$是连续函数，所以$\exists \delta>0$s.t.$x\in M,\rho(x,x_0)<\delta,|u_{N+1}(x)-u_{N+1}(x_0)|<\frac{\varepsilon}{3}$。
\[
|u(x)-u(x_0)|\leqslant |u(x)-u_{N+1}(x)|+|u_{N+1}(x)-u_{N+1}(x_0)|+|u_{N+1}(x_0)-u(x_0)|<\varepsilon
\]

\end{proof}
\begin{proposition}
  若 \( f: \mathscr{X} \to \mathscr{Y} \) 是一个连续函数，且 \( \mathscr{X} \) 是紧集，则 \( f \) 在 \( \mathscr{X} \) 上一致连续。  
\end{proposition}
\begin{proof}
对任意 \( \varepsilon > 0 \)，由于 \( f \) 在 \( \mathscr{X} \) 上是连续的，存在 \( U_x \subset \mathscr{X} \)，对于任意 \( x \in \mathscr{X} \)，我们可以定义
\[
U_x := \left\{ x' \in \mathscr{X} : d\left(f(x), f(x')\right) < \frac{\varepsilon}{2} \right\}.
\]
定义 \( B_x \) 为 \( U_x \) 中包含 \( x \) 的最大开球，即 \( B\left(x, \delta_x\right) \)，然后令 \( B_x' = B\left(x, \frac{\delta_x}{2}\right) \)。
考虑这族开球 \( \{ B_x' \}_{x \in \mathscr{X}} \)，它构成 \( \mathscr{X} \) 的开覆盖。由于 \( \mathscr{X} \) 是紧集，根据紧集的开覆盖定理，存在有限子覆盖
\[
\{ B_{x_i}' \}_{i=1}^n = \left\{ B\left(x_i, \frac{\delta_i}{2}\right) \right\}_{i=1}^n.
\]
因此，我们可以取 \( \delta = \frac{1}{2} \min_{1 \leqslant i \leqslant n} \delta_i \)。接下来，我们要证明 \( f \) 一致连续。\\
假设 \( z_1, z_2 \in \mathscr{X} \)，并且满足 \( d(z_1, z_2) < \delta \)。不失一般性，假设 \( z_1 \in B_1' \)，即
\[
d(z_2, x_1) \leqslant d(z_2, z_1) + d(z_1, x_1) = \delta + \frac{\delta_1}{2} \leqslant \frac{\delta_1}{2} + \frac{\delta_1}{2} = \delta_1.
\]
因此，\( z_1, z_2 \in B_1 \) 并且满足
\[
d\left(f(z_1), f(z_2)\right) \leqslant d\left(f(z_1) - f(x_1)\right) + d\left(f(z_2) - f(x_1)\right) \leqslant \frac{\varepsilon}{2} + \frac{\varepsilon}{2} = \varepsilon.
\]
这证明了 \( f \) 在 \( \mathscr{X} \) 上是一致连续的。
\end{proof}
\begin{definition}[一致有界和等度连续]
设 \( F \) 是 \( C(M) \) 的一个子集。称 \( F \) 是一致有界的，如果 \( \exists M_1 > 0 \)，使得 \( |\varphi(x)| \leqslant M_1 \quad (\forall x \in M, \forall \varphi \in F) \)；称 \( F \) 是等度连续的，如果 \( \forall \varepsilon > 0 \)，总可以找到 \( \delta(\varepsilon) > 0 \)，使得
\[
\left|\varphi(x_1) - \varphi(x_2)\right| < \varepsilon \quad (\forall x_1, x_2 \in M, \rho(x_1, x_2) < \delta, \forall \varphi \in F).
\]
\end{definition}

\begin{theorem}[Arzelà-Ascoli]
为了 \( F \subset C(M) \) 是一个列紧集，必须且仅须 \( F \) 是一致有界且等度连续的函数族。
\end{theorem}
\begin{proof}
因为 \( C(M) \) 是完备的，所以由定理 1.4，为了 \( F \) 是列紧的，必须且仅须它是完全有界的。\\
\textbf{必要性。} 因为完全有界集是有界集，所以 \( F \) 是一致有界函数族。\( \forall \varepsilon > 0 \)，要证 \( \exists \delta = \delta(\varepsilon) \)，使得 \( \forall \varphi \in F \) 有
\[
\left|\varphi(x_1) - \varphi(x_2)\right| < \varepsilon \quad (\text{当} \ \rho(x_1, x_2) < \delta).
\]
因为 \( F \) 的 \( \varepsilon/3 \) 网是一个有穷集 \( N(\varepsilon/3) = \{ \varphi_1, \varphi_2, \dots, \varphi_n \} \)，对这有穷个函数，由一致连续性，\( \exists \delta = \delta(\varepsilon/3) \)，当 \( \rho(x_1, x_2) < \delta \) 有
\[
\left|\varphi_i(x_1) - \varphi_i(x_2)\right| < \varepsilon/3 \quad (i = 1, 2, \dots, n).
\]
因为 \( \forall \varphi \in F, \exists \varphi_i \in N(\varepsilon/3) \)，使得 \( d(\varphi, \varphi_i) < \varepsilon/3 \)，所以
\[
\begin{aligned}
\left|\varphi(x) - \varphi(x')\right| &\leqslant \left|\varphi(x) - \varphi_i(x)\right| + \left|\varphi_i(x) - \varphi_i(x')\right| + \left|\varphi_i(x') - \varphi(x')\right| \\
&\leqslant 2 d(\varphi, \varphi_i) + \left|\varphi_i(x) - \varphi_i(x')\right| < \varepsilon \quad (\text{当} \ \rho(x, x') < \delta).
\end{aligned}
\]
\textbf{充分性。} 设 \( F \) 一致有界且等度连续，我们要找有穷的 \( \varepsilon \) 网。由于 \( F \) 是等度连续的，\( \exists \delta = \delta(\varepsilon/3) > 0 \)，使得当 \( \rho(x, x') < \delta \) 时，\( \left|\varphi(x) - \varphi(x')\right| < \varepsilon/3 \quad (\forall \varphi \in F) \)。就此 \( \delta \)，选取空间 \( M \) 上的有穷 \( \delta \) 网 \( N(\delta) = \{ x_1, x_2, \dots, x_n \} \)。做映射
\[
T: F \to \mathbb{R}^n, \quad T\varphi \triangleq \left( \varphi(x_1), \varphi(x_2), \dots, \varphi(x_n) \right) \quad (\forall \varphi \in F).
\]
记 \( \widetilde{F} = T(F) \)，则 \( \widetilde{F} \) 是 \( \mathbb{R}^n \) 中的有界集。事实上，设 \( |\varphi| \leqslant M_1 \) (\( \forall \varphi \in F \))，则
\[
\left( \sum_{i=1}^n \left| \varphi(x_i) \right|^2 \right)^{\frac{1}{2}} \leqslant \sqrt{n} \max_{x \in M} |\varphi(x)| \leqslant \sqrt{n} M_1 \quad (\forall \varphi \in F).
\]
从而 \( \widetilde{F} \) 是列紧集，利用定理 1.4，\( \widetilde{F} \) 有有穷的 \( \varepsilon/3 \) 网
\[
\widetilde{N}(\varepsilon/3) = \{ T\varphi_1, T\varphi_2, \dots, T\varphi_m \}.
\]
从而 \( \{ \varphi_1, \varphi_2, \dots, \varphi_m \} \) 是 \( F \) 的 \( \varepsilon \) 网，这是因为 \( \forall \varphi \in F, \exists \varphi_i \)，使得 \( \rho_n(T\varphi, T\varphi_i) < \varepsilon/3 \)，于是取定 \( x_r \in N(\delta) \)，使得 \( \rho(x, x_r) < \delta \)，有
\[
\begin{aligned}
\left| \varphi(x) - \varphi_i(x) \right| &\leqslant \left| \varphi(x) - \varphi(x_r) \right| + \left| \varphi(x_r) - \varphi_i(x_r) \right| + \left| \varphi_i(x_r) - \varphi_i(x) \right| \\
&< \frac{2}{3} \varepsilon + \rho_n(T\varphi, T\varphi_i) < \varepsilon.
\end{aligned}
\]
其中 \( \rho_n \) 表示 \( \mathbb{R}^n \) 上的距离。
\end{proof}
\begin{example}
在完备度量空间中，证明子集 \( A \) 是列紧的充分必要条件是：对任意 \( \varepsilon > 0 \)，存在 \( A \) 的列紧的 \( \varepsilon \) 网。

\end{example}

\begin{proof}
\textbf{必要性：} 显然。\\
\textbf{充分性：} 对任意 \( \varepsilon > 0 \)，设 \( N \) 是 \( A \) 的列紧的 \( \frac{\varepsilon}{2} \) 网。又设 \( N_0 \) 是 \( N \) 的有限 \( \frac{\varepsilon}{2} \) 网，则有
\[
\forall x \in A, \quad \exists \xi \in N, \quad \rho(x, \xi) < \frac{\varepsilon}{2},
\]
且
\[
\xi \in N, \quad \exists x_\varepsilon \in N_0, \quad \rho\left(\xi, x_\varepsilon\right) < \frac{\varepsilon}{2}.
\]
因此，得到
\[
\rho\left(x, x_\varepsilon\right) \leqslant \rho(x, \xi) + \rho\left(\xi, x_\varepsilon\right) < \frac{\varepsilon}{2} + \frac{\varepsilon}{2} = \varepsilon.
\]
即，\( N_0 \) 是 \( A \) 的有限 \( \varepsilon \) 网。由于 \( A \) 是完全有界的，得出结论：\( A \) 是列紧的。
\end{proof}
\begin{problem}
设 $( \mathscr{X} , \rho)$ 是度量空间； $M$ 是 $\mathscr{X}$ 中的列紧集，映射 $f: \mathscr{X} \rightarrow M$ 满足.
\begin{equation*}
\rho\left(f\left(x_1\right), f\left(x_2\right)\right)<\rho\left(x_1, x_2\right) \quad\left(\forall x_1, x_2 \in \mathscr{X} , x_1 \neq x_2\right) .
\end{equation*}
求证: $f$在 $\mathscr{X}$ 中存在唯一的不动点. 
\end{problem}
\begin{proof}
令\[
d=\inf\{\rho(x,f(x))|x\in \overline{M}\}
\]
则取$\overline{M}$中的点列$\{x_n\}$满足$\rho(x_n,f(x_n))\rightarrow d$，由于$\overline{M}$列紧且闭，所以存在子序列${x_{n_k}}$收敛于$x_0\in \overline{M}$。
\[d(x_0,f(x_0))\leqslant \rho(x_0,x_{n_k})+\rho(x_{n_k},f(x_{n_k}))+\rho(f(x_{n_k}),f(x_0))\leqslant2\rho(x_0,x_{n_k})+\rho(x_{n_k},f(x_{n_k}))\rightarrow d(k\rightarrow \infty)\]
所以$\rho(x_0,f(x_0))=d$，下面证$d=0$。否则\[
d\leqslant \rho(f(x_0),f^2(x_0))<\rho(x_0,f(x_0))=d
\]
矛盾。所以 $f$在 $\mathscr{X}$ 中存在唯一的不动点. 
\end{proof}
\section{赋泛线性空间}\begin{definition}
设 $\mathscr{X}$ 是一个非空集，$K$ 是复（或实）数域。如果下列条件满足，便称 $\mathscr{X}$ 为一复（或实）线性空间：

\begin{enumerate}
  \item $\mathscr{X}$ 是一加法交换群，即对 $\forall x, y \in \mathscr{X}$，存在 $u \in \mathscr{X}$，记作 $u = x + y$，称 $u$ 为 $x$ 和 $y$ 的和，适合
  \begin{itemize}
    \item (1.1) $x + y = y + x$；
    \item (1.2) $(x + y) + z = x + (y + z)$；
    \item (1.3) 存在唯一的 $\theta \in \mathscr{X}$，对 $\forall x \in \mathscr{X}$，$x + \theta = \theta + x$；
    \item (1.4) 对任意的 $x \in \mathscr{X}$，存在 $x' \in \mathscr{X}$，使得 $x + x' = \theta$，记此 $x'$ 为 $-x$。
  \end{itemize}
  \item 定义了数域 $K$ 中的数 $\alpha$ 与 $x \in \mathscr{X}$ 的数乘运算，即对 $\forall (\alpha, x) \in K \times \mathscr{X}$，存在 $u \in \mathscr{X}$，记作 $u = \alpha x$，称 $u$ 为 $x$ 对 $\alpha$ 的数乘，适合
  \begin{itemize}
    \item (2.1) $\alpha (\beta x) = (\alpha \beta) x$，$\forall \alpha, \beta \in K, \forall x \in \mathscr{X}$；
    \item (2.2) $1 \cdot x = x$；
    \item (2.3)$(\alpha + \beta) x = \alpha x + \beta x$，$\forall \alpha, \beta \in K, \forall x \in \mathscr{X}$；
    \item (2.4)$\alpha (x + y) = \alpha x + \alpha y$，$\forall x, y \in \mathscr{X}, \forall \alpha \in K$。
  \end{itemize}
\end{enumerate}
\end{definition}

\begin{definition}[线性同构]
设 $\mathscr{X}, \mathscr{X}_1$ 都是线性空间，$T: \mathscr{X} \rightarrow \mathscr{X}_1$ 称为是一个线性同构，如果
\begin{enumerate}
  \item 它既是单射又是满射，即它是一对一的并且是在上的；
  \item $T(\alpha x + \beta y) = \alpha T x + \beta T y$，$\forall x, y \in \mathscr{X}, \forall \alpha, \beta \in K$。
\end{enumerate}
\end{definition}

\begin{definition}[线性子空间]
设 $E \subset \mathscr{X}$，若 $E$ 依 $\mathscr{X}$ 上的加法与数乘还构成一个线性空间，则称 $E$ 是 $\mathscr{X}$ 的一个线性子空间。$\mathscr{X}$ 以及 $\{\theta\}$ 都是 $\mathscr{X}$ 的线性子空间，我们称它们为平凡的子空间，而称其他的子空间为真子空间。
\end{definition}

\begin{definition}[线性流形]
设 $E \subset \mathscr{X}$，若 $\exists x_0 \in \mathscr{X}$ 及线性子空间 $E_0 \subset \mathscr{X}$，使得 $E = E_0 + x_0 \triangleq \left\{x + x_0 \mid x \in E_0\right\}$，则称 $E$ 为线性流形。简单地说，线性流形就是子空间对某个向量的平移。
\end{definition}

\begin{definition}[线性相关]
一组向量 $x_1, x_2, \cdots, x_n \in \mathscr{X}$ 称为是线性相关的，如果存在 $\lambda_1, \lambda_2, \cdots, \lambda_n \in K$ 不全为 0，使得
\[
\lambda_1 x_1 + \lambda_2 x_2 + \cdots + \lambda_n x_n = 0;
\]
否则称为是线性无关的。
\end{definition}

\begin{definition}[线性基]
若 $A$ 是 $\mathscr{X}$ 中的一个极大线性无关向量组，即 $A$ 中的向量是线性无关的，而且任意的 $x \in \mathscr{X}$ 都是 $A$ 中的向量的线性组合，则称 $A$ 是 $\mathscr{X}$ 的一组线性基。
\end{definition}

\begin{definition}[维数]
线性空间中的线性基的元素个数（势），称为维数。
\end{definition}

\begin{definition}[线性包]
设 $\Lambda$ 是一个指标集，$\left\{x_\lambda \mid \lambda \in \Lambda\right\}$ 是 $\mathscr{X}$ 中的向量族，一切由 $\left\{x_\lambda \mid \lambda \in \Lambda\right\}$ 的有穷线性组合组成的集合
\[
\left\{y = \alpha_1 x_{\lambda_1} + \cdots + \alpha_n x_{\lambda_n} \mid \lambda_i \in \Lambda, \alpha_i \in K, i = 1, 2, \cdots, n\right\}
\]
称为 $\left\{x_\lambda \mid \lambda \in \Lambda\right\}$ 的线性包。这个线性包是一个线性子空间，不难证明它是包含 $\left\{x_\lambda \mid \lambda \in \Lambda\right\}$ 的一切线性子空间的交。因此称线性包为 $\left\{x_\lambda \mid \lambda \in \Lambda\right\}$ 张成的线性子空间，记为
\[
\operatorname{span}\left\{x_\lambda \mid \lambda \in \Lambda\right\}.
\]
\end{definition}

\begin{definition}[线性和与直和]
设 $E_1, E_2$ 是 $\mathscr{X}$ 的子空间，我们称集合 $\left\{x + y \mid x \in E_1, y \in E_2\right\}$ 为 $E_1$ 与 $E_2$ 的线性和，记为 $E_1 + E_2$。对于任意有限个子空间，定义以此类推。又若 $(E_1, E_2)$ 中的任意一对非零向量都是线性无关的，则称线性和 $E_1 + E_2$ 为直和，记作 $E_1 \oplus E_2$，这时 $E_1 \cap E_2 = \{\theta\}$。对任意 $x \in E_1 \oplus E_2$，有唯一的分解：
\[
x = x_1 + x_2 \quad (x_i \in E_i, i = 1, 2).
\]
\end{definition}
\begin{definition}
线性空间 $\mathscr{X}$ 上的准范数（准模）定义为该空间上的一个函数 $\|\cdot\|: \mathscr{X} \rightarrow \mathbb{R}$，满足以下条件：
\begin{enumerate}
  \item $\|x\| \geqslant 0 \quad (\forall x \in \mathscr{X})$，且 $\|x\| = 0 \Longleftrightarrow x = \theta$；
  \item $\|x + y\| \leqslant \|x\| + \|y\| \quad (\forall x, y \in \mathscr{X})$；
  \item $\|-x\| = \|x\| \quad (\forall x \in \mathscr{X})$；
  \item $\lim\limits_{\alpha_n \to 0} \left\|\alpha_n x\right\| = 0$，若  $\lim\limits_{\|x_n\|\to 0}\left\|\alpha x_n\right\| \to 0 \quad (\forall x \in \mathscr{X}, \forall \alpha \in K)$。
\end{enumerate}
\end{definition}


\begin{definition}
设 $\mathscr{X}$ 为一个赋准范数的线性空间。如果按照
\[
\left\| x_n - x \right\| \to 0 \quad (n \to \infty)
\]
来定义 $x_n \to x$（$n \to \infty$），那么便称其为 $F^*$ 空间。
\end{definition}


\begin{definition}
完备的 $F^*$ 空间称为 Frechet 空间，简称 $F$ 空间。$F^*$ 空间的例子很多。
\end{definition}


\begin{example}
设 $M$ 是一个紧度量空间。显然
\[
\|u\| = \max_{x \in M} |u(x)|
\]
是一个准范数。
\end{example}


\begin{example}
设 $x = (x_1, x_2, \cdots, x_n) \in \mathbb{R}^n$，定义
\[
\|x\| = \left(\sum_{i=1}^n |x_i|^2 \right)^{\frac{1}{2}}.
\]
显然它是一个准范数，$\mathbb{R}^n$ 是一个 $F$ 空间。
\end{example}
\begin{definition}
线性空间 $\mathscr{X}$ 上的范数 $\|\cdot\|$ 是一个非负值函数 $\mathscr{X} \to \mathbb{R}$，满足以下条件：
\begin{enumerate}
  \item $\|x\| \geqslant 0 \quad (\forall x \in \mathscr{X})$，且 $\|x\| = 0 \Longleftrightarrow x = \theta$（正定性）；
  \item $\|x + y\| \leqslant \|x\| + \|y\| \quad (\forall x, y \in \mathscr{X})$（三角形不等式）；
  \item $\|\alpha x\| = |\alpha| \cdot \|x\| \quad (\forall \alpha \in K, \forall x \in \mathscr{X})$（齐次性）。
\end{enumerate}
\end{definition}

\begin{definition}
当赋准范数的线性空间中的准范数是范数时，这空间叫作赋范线性空间，或称 $B^*$ 空间。完备的 $B^*$ 空间叫作 $B$ 空间或 Banach 空间。
\end{definition}
\begin{definition}[$L^p$ 范数]
设
\begin{itemize}
    \item $(\mathscr{X}, A, \mu)$ 是测度空间；
    \item $f: \mathscr{X} \to K$ 是可测函数；
    \item $p \in (0, +\infty]$。
\end{itemize}
则 $f$ 的 $L^p$ 范数（或 $p$-范数），记为 $\| f \|_p \in [0, +\infty]$，定义如下：
\begin{itemize}
    \item 若 $0 < p < \infty$，定义
    \[
    \| f \|_p = \left( \int_X |f|^p \, d\mu \right)^{1/p}
    \]
    其中的积分是 Lebesgue 积分；\\
    \item 若 $p = \infty$，定义
\[
    \| f \|_{\infty} = \inf \left\{ a \in \mathbb{R}_{\geq 0} \mid |f(x)| \leqslant a \text{ 对几乎所有 } x \in \mathscr{X} \text{ 成立} \right\}.
    \]
    当 $\mu(\mathscr{X}) > 0$ 时，这也就是 $|f|$ 的本质上确界，此时 $L^\infty$ 范数也称为一致范数。
\end{itemize}
\end{definition}

\begin{lemma}[Young不等式]
考虑正实数 $a, b, p, q$ ，且 $\frac{1}{p}+\frac{1}{q}=1$ ，则有 $a b \leqslant \frac{a^p}{p}+\frac{b^q}{q}$ 。等号成立当且仅当 $a^p=b^q$
\end{lemma}
\begin{proof}
我们知道函数 $f(x)=e^x$ 是一个凸函数，因为它的二阶导数恒为正。从而我们有:
\begin{equation*}
a b=e^{\ln (a)} e^{\ln (b)}=e^{\frac{1}{p} \ln \left(a^p\right)+\frac{1}{q} \ln \left(b^q\right)} \leqslant \frac{1}{p} e^{\ln \left(a^p\right)}+\frac{1}{q} e^{\ln \left(b^q\right)}=\frac{a^p}{p}+\frac{b^q}{q}
\end{equation*}
\end{proof}
\begin{theorem}[Hölder 不等式]
设 $1 \leqslant p \leqslant \infty, q$ 满足 $ \frac{1}{p} + \frac{1}{q} = 1$，即 $p, q$ 互为共轭指数，且 $f, g$ 是 $\mathscr{X}$ 上的可测函数，则有：
\[
\| f g \|_1 \leqslant \| f \|_p \| g \|_q.
\]
特别地，若 $f \in L^p$ 且 $g \in L^q$，则 $f g \in L^1$。等号成立的条件如下：
\begin{itemize}
    \item 如果 $1 < p, q < \infty$，则等号成立当且仅当 $|f|^p$ 和 $|g|^q$ 几乎处处成比例。
    \item 如果 $p = 1, q = \infty$，则等号成立当且仅当 $g$ 在 $f$ 的零点集以外几乎处处等于 $\| g \|_{\infty}$。
\end{itemize}
\end{theorem}

\begin{proof}
如果 $\| f \|_p$ 或 $\| g \|_q$ 中有等于 $0$ 或 $\infty$ 的情况，则不等式显然成立。下面假设不是这种情况，分两部分证明。\\
\textbf{情况 1：} 如果 $1 < p, q < \infty$，则可以假设 $\| f \|_p = \| g \|_q = 1$。此时根据 Young 不等式：
\[
|f(x) g(x)| \leqslant p^{-1} |f(x)|^p + q^{-1} |g(x)|^q.
\]
两边积分得：
\[
\| f g \|_1 \leqslant p^{-1} \int |f|^p + q^{-1} \int |g|^q = 1 = \| f \|_p \| g \|_q.
\]
因此，不等式得证。根据 Young 不等式的取等条件，得到等号成立的条件。\\
\textbf{情况 2：} 如果 $p = 1, q = \infty$，对任意 $a > 0$，设 $E_a = \{x \mid |g(x)| > a\}$。如果 $\mu(E_a) = 0$，则有：
\[
\| f g \|_1 = \int |f g| = \int_{\mathscr{X} \setminus E_a} |f g| \leqslant a \int_{\mathscr{X} \setminus E_a} |f| \leqslant a \| f \|_1.
\]
对 $a$ 取下确界，得到：
\[
\| f g \|_1 \leqslant \| f \|_1 \| g \|_{\infty}.
\]
最后，得到取等条件。
\end{proof}
\begin{theorem}[Minkowski 不等式]
如果 $1 \leqslant p \leqslant \infty, f, g \in L^p$，则
\[
\|f + g\|_p \leqslant \|f\|_p + \|g\|_p.
\]
\end{theorem}

\begin{proof}
如果 $p = 1, \infty$ 或 $f + g = 0$ 几乎处处成立，则不等式显然成立。否则，我们有
\[
|f + g|^p \leqslant (|f| + |g|) |f + g|^{p-1}.
\]
从而由Hölder 不等式得
\[
\begin{aligned}
\int |f + g|^p & \leqslant \int(|f| + |g|) |f + g|^{p-1}= \int |f| |f + g|^{p-1}+\int |g| |f + g|^{p-1}\leqslant \|f\|_p \left\| |f + g|^{p-1} \right\|_q + \|g\|_p \left\| |f + g|^{p-1} \right\|_q \\
& = \left(\|f\|_p + \|g\|_p\right) \left(\int |f + g|^p \right)^{1/q}.
\end{aligned}
\]
于是
\[
\|f + g\|_p = \left( \int |f + g|^p \right)^{1 - 1/q} \leqslant \|f\|_p + \|g\|_p.
\]
这个不等式说明 $p$-范数在 $p \geq 1$ 时确实是范数，于是 $L^p$ 空间是赋范线性空间。
\end{proof}

\begin{theorem}[Riesz-Fisher 定理]
对任意 $1 \leqslant p \leqslant \infty$，$L^p$ 都是 Banach 空间。
\end{theorem}

\begin{proof}
(1) $1 \leqslant p<\infty$ 的情况: 若函数列 $\left\{f_n\right\} \subseteq L^p(E)$ 是Cauchy列，则:
$\lim \limits_{m, n \rightarrow \infty}\left\|f_m-f_n\right\|_p=0$ ，可取子列
\[\left\{f_{n_k}\right\} \subseteq\left\{f_n\right\}, s.t. \left\|f_{n_k}-f_{n_{k+1}}\right\|_p<\frac{1}{2^k}, k \geq 1 . \forall E_1 \subseteq E, m\left(E_1\right)<\infty\]
有:
\[\int_{E_1}\left|f_{n_k}-f_{n_{k+1}}\right| d m \leqslant\left\|f_{n_k}-f_{n_{k+1}}\right\|_p\left\|\chi_{E_1}\right\|_q \leqslant \frac{m\left(E_1\right)^{\frac{1}{q}}}{2^k}\] 从而级数 \[\sum_{k=1}^{\infty} \int_{E_1}\left|f_{n_k}-f_{n_{k+1}}\right| d m \leqslant m\left(E_1\right)^{\frac{1}{q}} \cdot \sum_{k=1}^{\infty} \frac{1}{2^k}=m\left(E_1\right)^{\frac{1}{q}}<\infty\]
收敛（推出几乎处处有限），由 $E_1$ 的任意性，级数 $\sum_{k=1}^{\infty}\left|f_{n_k}(x)-f_{n_{k+1}}(x)\right|$ 在$E$上几乎处处收敛，于是：
$f_{n_k}(x)=\sum\limits_{j=1}^{k-1}\left(f_{n_{j+1}}(x)-f_{n_j}(x)\right)+f_{n_1}(x)$ 在$E$上几乎处处收敛于一可测函数 $f(x)$ ，下证 $f_n \xrightarrow{L^p} f$.\\
由Cauchy列的定义: $\forall \varepsilon>0, \exists N \in N ^*$, s.t. $n, n_k>N$ 时, $\left\|f_n-f_{n_k}\right\|_{L^p}<\varepsilon$.
对非负可测函数列 $\left\{\left|f_n-f_{n_k}\right|^p\right\}_{k=1}^{\infty}$ 使用Fatou引理：
\begin{equation*}
\begin{gathered}
\int_E\left|f_n-f\right|^p d m=\int_E \varliminf_{k\to\infty}\left|f_n-f_{n_k}\right|^p d m \\
\quad \leqslant \varliminf_{k \rightarrow \infty} \int_E\left|f_n-f_{n_k}\right\|^p d m \\
=\varliminf_{k \rightarrow \infty} \| f_n-f_{n_k}\|^p \leqslant \varepsilon^p(n>N)
\end{gathered}
\end{equation*}
因此: $f_n-f \in L^p(E) \Rightarrow f=f_n-\left(f_n-f\right) \in L^p(E)$ ，且: $f_n \xrightarrow{L^p} f$。\\
(2)当 $p=\infty$ 时: 设 $\left\{f_n\right\} \subseteq L^{\infty}(E)$ 满足: $\lim _{m, n \rightarrow \infty}\left\|f_m-f_n\right\|_{\infty}=0$.
从而 $\exists$ 零测集 $E_0$, s.t. $\forall m, n \in N ^*,\left|f_m(x)-f_n(x)\right| \leqslant\left\|f_m-f_n\right\|_{\infty}, x \in E \backslash E_0$. 故$\exists f(x)$, s.t. $f_k \xrightarrow{\text { a.e. }} f$ ，易知 $f \in L^{\infty}(E)$.
现在，任取 $\varepsilon>0, N_{\varepsilon} \in N ^*$, s.t. $\left\|f_m-f_n\right\|_{\infty}<\varepsilon, m, n>N_{\varepsilon}$.
由于当 $n>N_{\varepsilon}, x \in E \backslash E_0$ 时，有: $\left|f_n(x)-f(x)\right|=\lim _{m \rightarrow \infty}\left|f_n(x)-f_m(x)\right|<\varepsilon$ ，故当 $n>N_{\varepsilon}$ 时，有: $\left\|f_n-f\right\|_{\infty}<\varepsilon$ ，这说明 $\left\|f_n-f\right\|_{\infty} \rightarrow 0(n \rightarrow \infty)$ ，从而 $f_n \xrightarrow{L^{\infty}} f$.
\end{proof}

\begin{example}
空间 $L^p(\Omega, \mu) \quad (1 \leqslant p < \infty)$。设 $(\Omega, B, \mu)$ 是一个测度空间，$u$ 是 $\Omega$ 上的可测函数，而且 $|u(x)|^p$ 在 $\Omega$ 上是可积的。这种函数 $u$ 的全体记作 $L^p(\Omega, \mu)$，叫作 $(\Omega, B, \mu)$ 上的 $p$ 次可积函数空间。$L^p(\Omega, \mu)$ 按通常的加法与数乘规定运算，并且把几乎处处（记作 a. e.）相等的两个函数看成是同一个向量，经这样处理过的空间 $L^p(\Omega, \mu)$ 仍是一个线性空间，并且定义
\[
\|u\| = \left( \int_{\Omega} |u(x)|^p \, d\mu \right)^{\frac{1}{p}},
\]
那么 $\|\cdot\|$ 是一个范数。 条件 (2) 正是著名的 Minkowski 不等式：
\[
\left( \int_{\Omega} |u(x) + v(x)|^p \, d\mu \right)^{\frac{1}{p}} \leqslant \left( \int_{\Omega} |u(x)|^p \, d\mu \right)^{\frac{1}{p}} + \left( \int_{\Omega} |v(x)|^p \, d\mu \right)^{\frac{1}{p}}.
\]
又由 Riesz-Fisher 定理，$L^p(\Omega, \mu)$ 还是一个 $B$ 空间。
\end{example}
\begin{note}本例有两个重要的特殊情形：
\begin{enumerate}
  \item $\Omega$ 是 $\mathbb{R}^n$ 中的一个可测集，而 $d\mu$ 即是普通的 Lebesgue 测度，这时，对应的空间记作 $L^p(\Omega)$；
  \item $\Omega = \mathbb{N}$，而测度 $\mu$ 是等分布的：$\mu(\{n\}) = 1 \quad (\forall n \in \mathbb{N})$，这时空间 $L^p(\Omega, \mu)$ 由满足 $\sum_{n=1}^{\infty} |u_n|^p < \infty$ 的序列 $u = \{u_n\}_{n=1}^{\infty}$ 组成，对应的空间记作 $l^p$，其范数是
  \[
  \|u\| = \left( \sum_{n=1}^{\infty} |u_n|^p \right)^{\frac{1}{p}}.
  \]
\end{enumerate}
\end{note}
\begin{definition}
空间 $L^{\infty}(\Omega, \mu)$。设 $(\Omega, B, \mu)$ 是一个测度空间，$\mu$ 对于 $\Omega$ 是 $\sigma$-有限的，$u(x)$ 是 $\Omega$ 上的可测函数。如果 $u(x)$ 与 $\Omega$ 上的一个有界函数几乎处处相等，则称 $u(x)$ 是 $\Omega$ 上的一个本性有界可测函数。$\Omega$ 上的一切本性有界可测函数（把 a. e. 相等的两个函数视为同一个向量）的全体记作 $L^{\infty}(\Omega, \mu)$，在其上规定：
\[
\|u\| = \inf_{\substack{\mu(E_0) = 0 \\ E_0 \subset \Omega}} \left( \sup_{x \in \Omega \backslash E_0} |u(x)| \right).
\]
此式右端有时也记作 $\operatorname{ess} \sup_{x \in \Omega} |u(x)|$。

\end{definition}
\begin{proof}
显然 $L^{\infty}(\Omega, \mu)$ 是一个线性空间，以下验证 $\|\cdot\|$ 是一个范数。事实上，定义中的条件 (3) 及 $\|u\| \geqslant 0$ 都是显然的，需要验证：
\begin{enumerate}
  \item $\|u\| = 0 \Longleftrightarrow u = \theta$。充分性显然。为了证明必要性，若 $\|u\| = 0$，则对于任意 $n \in \mathbb{N}$，存在 $E_n \subset \Omega$，使得 $\mu(E_n) = 0$，并且
  \[
  \sup_{x \in \Omega \backslash E_n} |u(x)| < \frac{1}{n}.
  \]
  令
  \[
  \Omega_1 \triangleq \bigcap_{n=1}^{\infty} (\Omega \backslash E_n) = \Omega \backslash \bigcup_{n=1}^{\infty} E_n.
  \]
  因为 $u(x) = 0$ 当 $x \in \Omega_1$，且 $\mu(\bigcup_{n=1}^{\infty} E_n) = 0$，所以
  \[
  u(x) = 0 \quad (\text{a.e.} \, x \in \Omega), \quad \text{即} \, u = \theta.
  \]
  \item $\|u + v\| \leqslant \|u\| + \|v\|$。对于任意 $\varepsilon > 0$，存在 $E_0, E_1 \subset \Omega$，使得 $\mu(E_0) = \mu(E_1) = 0$，且
  \[
  \sup_{x \in \Omega \backslash E_0} |u(x)| \leqslant \|u\| + \frac{\varepsilon}{2}, \quad \sup_{x \in \Omega \backslash E_1} |v(x)| \leqslant \|v\| + \frac{\varepsilon}{2}.
  \]
  因此
  \[
  \|u + v\| \leqslant \sup_{x \in \Omega \backslash (E_0 \cup E_1)} |u(x) + v(x)| \leqslant \sup_{x \in \Omega \backslash E_0} |u(x)| + \sup_{x \in \Omega \backslash E_1} |v(x)|.
  \]
  从而
  \[
  \|u + v\| \leqslant \|u\| + \|v\| + \varepsilon.
  \]
  因为 $\varepsilon > 0$ 是任意的，所以得到了所要证的结论。
\end{enumerate}

最后来证 $L^{\infty}(\Omega, \mu)$ 是完备的。设
\[
\|u_{n+p} - u_n\| \to 0 \quad (\text{当} \, n \to \infty, \forall p \in \mathbb{N}).
\]
根据 $L^{\infty}(\Omega, \mu)$ 中范数定义，存在 $Z_{n,p} \subset \Omega, \mu(Z_{n,p}) = 0$，使得
\[
|u_{n+p}(x) - u_n(x)| \leqslant \|u_{n+p} - u_n\| + \frac{1}{2^{n+p}} \quad (\forall x \in \Omega \backslash Z_{n,p}).
\]
令 $Z = \bigcup_{n,p} Z_{n,p}$，则 $\mu(Z) = 0$，且
\[
|u_{n+p}(x) - u_n(x)| \leqslant \|u_{n+p} - u_n\| + \frac{1}{2^{n+p}} \quad (\forall x \in \Omega \backslash Z, n,p \in \mathbb{N}).
\]
, $\forall \varepsilon>0, \exists N \in N$, 使 $\forall n \geqslant N, p \in N$ 有
\begin{equation*}
\left|u_{n+p}(x)-u_n(x)\right|<\varepsilon+\frac{1}{2^{n+p}} \quad(x \in \Omega \backslash Z)
\end{equation*}
由此 $\lim\limits _{n \rightarrow \infty} u_n(x)=u(x)$ 存在, $x \in \Omega \backslash Z$, 在 上式中令 $p \rightarrow \infty$得
\begin{equation*}
\left|u(x)-u_n(x)\right| \leqslant \varepsilon \quad(x \in \Omega \backslash Z)
\end{equation*}
即
\begin{equation*}
\sup _{x \in \Omega \backslash Z}\left|u(x)-u_n(x)\right| \leqslant \varepsilon
\end{equation*}
因 $\mu(Z)=0$, 所以 $u \in L^{\infty}(\Omega, \mu)$, 且
\begin{equation*}
\left\|u_n-u\right\| \leqslant \varepsilon \quad(n \geqslant N),
\end{equation*}
即
\begin{equation*}
\left\|u_n-u\right\| \rightarrow 0 \quad(n \rightarrow \infty)
\end{equation*}
\end{proof}
\begin{note}
当 $\Omega$ 是 $\mathbb{R} ^n$ 中的一个可测集时，对应的空间记作 $L^{\infty}(\Omega)$;当 $\Omega= \mathbb{N}$ 时，对应的空间记作 $l^{\infty}$ ，它是由一切有界序列 $u=$ $\left\{u_n\right\}_{n=1}^{\infty}$ 组成的空间, 其范数是 $\|u\|=\sup _{n \geqslant 1}\left|u_n\right|$.
\end{note}


关于 \({L}^{2}\left( \Omega \right)\) 中的列紧集的刻画由Rellich给出,
\begin{theorem}
  [Rellich 紧性判别准则] 设 \(\Omega  \subset  {\mathbb{R}}^{n}\) 为带光滑边界的有界区域. 假设 \(\Omega\) 上函数族 \({\left\{  {u}_{\lambda }\right\}  }_{\lambda  \in  \Lambda }\) 满足条件: 存在 \(M > 0\) 使得

\begin{equation*}
\parallel u{\parallel }_{{L}^{2}\left( \Omega \right) } \leqslant  M,\;{\begin{Vmatrix}\frac{\partial u}{\partial {x}_{i}}\end{Vmatrix}}_{{L}^{2}\left( \Omega \right) } \leqslant  M,i = 1,2,\cdots ,n,
\end{equation*}

则函数族 \({\left\{  {u}_{\lambda }\right\}  }_{\lambda  \in  \Lambda }\) 为 \({L}^{2}\left( \Omega \right)\) 中相对紧集.
\end{theorem}

为考虑空间 \({L}^{p}\left\lbrack  {a,b}\right\rbrack  \left( {1 < p < \infty }\right)\) 中集合的紧性,我们先引进一个记号. 对任意 \(x\left( t\right)  \in  {L}^{p}\left\lbrack  {a,b}\right\rbrack\) ,定义函数 \({x}_{h}\left( t\right)\) 如下:

\begin{equation*}
{x}_{h}\left( t\right)  \mathrel{\text{ := }} \frac{1}{2h}{\int }_{t - h}^{t + h}x\left( s\right) {\dif s},\;\left( {t \in  \left\lbrack  {a,b}\right\rbrack  ,h > 0}\right) ,
\end{equation*}

其中规定当 \(t \notin  \left\lbrack  {a,b}\right\rbrack\) 时, \(x\left( t\right)  = 0\) . 显然, \({x}_{h}\left( t\right)  \in  C\left\lbrack  {a,b}\right\rbrack\) . 在 \({L}^{p}\left\lbrack  {a,b}\right\rbrack\) 上,我们定义距离

\begin{equation*}
\rho \left( {x,y}\right)  = {\left( {\int }_{a}^{b}{\left| x\left( t\right)  - y\left( t\right) \right| }^{p}\dif t\right) }^{\frac{1}{p}}.
\end{equation*}
\begin{theorem}
  空间 \({L}^{p}\left\lbrack  {a,b}\right\rbrack  \left( {1 < p < \infty }\right)\) 中的集合 \(A\) 为列紧的当且仅当下列两条件成立:

( i )存在常数 \(K\) ,使得对任意 \(x\left( t\right)  \in  A\) 有

\begin{equation*}
{\int }_{a}^{b}{\left| x\left( t\right) \right| }^{p}{\dif t} \leqslant  {K}^{p}.
\end{equation*}

(ii) 对任给的 \(\varepsilon  > 0\) ,存在 \(\delta  > 0\) 使得当 \(0 < h < \delta\) 时,对任意 \(x\left( t\right)  \in  A\) 有 \(\rho \left( {{x}_{h},x}\right)  <\)

\(\varepsilon\) ,即

\begin{equation*}
{\left( {\int }_{a}^{b}{\left| {x}_{h}\left( t\right)  - x\left( t\right) \right| }^{p}\dif t\right) }^{\frac{1}{p}} < \varepsilon ,\;\forall x\left( t\right)  \in  A.
\end{equation*}

\end{theorem}


\begin{definition}
设在线性空间 \(\mathscr{X}\) 上给定了两个范数 \(\|\cdot\|_1\) 与 \(\|\cdot\|_2\)，我们说 \(\|\cdot\|_2\) 比 \(\|\cdot\|_1\) 强，是指
\[
\|x_n\|_2 \rightarrow 0 \Longrightarrow \|x_n\|_1 \rightarrow 0 \quad (n \rightarrow \infty).
\]
如果 \(\|\cdot\|_2\) 比 \(\|\cdot\|_1\) 强，而且 \(\|\cdot\|_1\) 又比 \(\|\cdot\|_2\) 强，则称 \(\|\cdot\|_1\) 与 \(\|\cdot\|_2\) 等价。
\end{definition}

\begin{proposition}
为了 \(\|\cdot\|_2\) 比 \(\|\cdot\|_1\) 强，必须且仅须存在常数 \(C>0\)，使得
\[
\|x\|_1 \leqslant C\|x\|_2 \quad (\forall x \in \mathscr{X}).
\]
\end{proposition}

\begin{proof}
充分性是显然的，下证必要性。用反证法。若不成立，则对 \(\forall n \in \mathbb{N}\)，\(\exists x_n \in \mathscr{X}\)，使得 
\[
\|x_n\|_1 > n\|x_n\|_2.
\]
令 \(y_n \triangleq x_n / \|x_n\|_1\)。一方面 \(\|y_n\|_1 = 1\)，另一方面因为
\[
0 \leqslant \|y_n\|_2 < \frac{1}{n} \quad (\forall n \in \mathbb{N}),
\]
所以 \(\|y_n\|_2 \to 0 \ (n \to \infty)\)。又因为 \(\|\cdot\|_2\) 比 \(\|\cdot\|_1\) 强，所以 \(\|y_n\|_1 \to 0 \ (n \to \infty)\)。这显然是一个矛盾。
\end{proof}

\begin{corollary}
为了 \(\|\cdot\|_1\) 与 \(\|\cdot\|_2\) 等价，必须且仅须存在常数 \(C_1, C_2 > 0\)，使得
\[
C_1\|x\|_1 \leqslant \|x\|_2 \leqslant C_2\|x\|_1 \quad (\forall x \in \mathscr{X}).
\]
\end{corollary}
如果线性空间 $\mathscr{X}$ 的维数是有穷数 $n$, 则记 $\operatorname{dim} \mathscr{X} =n$, 否则记为 $\operatorname{dim} \mathscr{X} =\infty$. 设 $\mathscr{X}$ 是一个赋范线性空间，并且设 $\operatorname{dim} \mathscr{X} =n$ ，这时 $\mathscr{X}$ 存在一组基: $e_1, e_2, \cdots, e_n$ 。任意一个元素 $x \in \mathscr{X}$ 有下列唯一的表示：
\begin{equation*}
x=\xi_1 e_1+\xi_2 e_2+\cdots+\xi_n e_n .
\end{equation*}
每个 \(x \in \mathscr{\mathscr{X}}\) 唯一地对应着 \(\mathbb{K}^n\) 空间中的一点 \(\xi = T x \triangleq (\xi_1, \xi_2, \cdots, \xi_n)\)。自然希望建立 \(x\) 在 \(\mathscr{X}\) 中的范数 \(\|x\|\) 与 \(T x\) 在 \(\mathbb{K}^n\) 中的范数
\[
\|T x\| = \|\xi\| \triangleq \left( \sum_{j=1}^n |\xi_j|^2 \right)^{\frac{1}{2}}
\]
之间的关系。为此，考察函数
\[
p(\xi) \triangleq \left\|\sum_{j=1}^n \xi_j e_j\right\| \quad (\forall \xi \in \mathbb{K}^n).
\]
首先，\(p\) 对 \(\xi\) 是一致连续的。事实上，对 \(\forall \xi = (\xi_1, \xi_2, \cdots, \xi_n)\) 和 \(\eta = (\eta_1, \eta_2, \cdots, \eta_n) \in \mathbb{K}^n\)，由三角不等式与 Schwarz 不等式有
\[
\begin{aligned}
|p(\xi) - p(\eta)| &\leqslant p(\xi - \eta) \\
&\leqslant \sum_{i=1}^n |\xi_i - \eta_i| \|e_i\| \\
&\leqslant \left( \sum_{i=1}^n |\xi_i - \eta_i|^2 \right)^{\frac{1}{2}} \left( \sum_{i=1}^n \|e_i\|^2 \right)^{\frac{1}{2}} \\
&\leqslant \|\xi - \eta\| \left( \sum_{i=1}^n \|e_i\|^2 \right)^{\frac{1}{2}}.
\end{aligned}
\]
其次，根据范数的齐次性，对 \(\forall \xi \in \mathbb{K}^n \setminus \{\theta\}\) 有
\[
p(\xi) = \|\xi\| \left\|\sum_{j=1}^n \frac{\xi_j}{\|\xi\|} e_j\right\| = \|\xi\| p\left(\frac{\xi}{\|\xi\|}\right).
\]
注意到 \(\mathbb{K}^n\) 的单位球面 \(S_1 \triangleq \{\xi \in \mathbb{K}^n \mid |\xi| = 1\}\) 是一个紧集，因此 \(p(\xi)\) 在 \(S_1\) 上必有非负的最小值 \(C_1\) 与最大值 \(C_2\)，即有
\[
C_1 \leqslant p(\xi) \leqslant C_2 \quad (\forall \xi \in S_1).
\]
按 上 式便有
\[
C_1 \|\xi\| \leqslant p(\xi) \leqslant C_2 \|\xi\| \quad (\forall \xi \in \mathbb{K}^n).
\]
下面证明其中 \(C_1 > 0\)。用反证法，假若 \(C_1 = 0\)，那么 \(\exists \xi^* \in S_1\) 满足 \(p(\xi^*) = 0\)。设 \(\xi^* = (\xi_1^*, \xi_2^*, \cdots, \xi_n^*)\)，即有
\[
\xi_1^* e_1 + \xi_2^* e_2 + \cdots + \xi_n^* e_n = 0.
\]
因为 \(\{e_1, e_2, \cdots, e_n\}\) 是基，所以上式蕴含 \(\xi^* = \theta\)。这与 \(\xi^* \in S_1\) 矛盾。可改写为
\[
C_1 \|T x\| \leqslant \|x\| \leqslant C_2 \|T x\| \quad (\forall x \in \mathscr{X}).
\]
如果我们将 \(\|T x\|\) 看作是在 \(\mathscr{X}\) 空间中引入的另一范数 \(\|x\|_T\)，即
\[
\|x\|_T \triangleq \|T x\| \quad (\forall x \in \mathscr{X}),
\]
那么 \(\|\cdot\|_T\) 与 \(\|\cdot\|\) 是等价的。以后简称 \(\|\cdot\|_T\) 为 \(\mathscr{X}\) 的 \(\mathbb{K}^n\) 范数。于是 \(n\) 维赋范线性空间的范数与其 \(\mathbb{K}^n\) 范数等价。
\begin{theorem}
设 \(\mathscr{X}\) 是一个有穷维线性空间，若 \(\|\cdot\|_1\) 与 \(\|\cdot\|_2\) 都是 \(\mathscr{X}\) 上的范数，则必有正常数 \(C_1\) 与 \(C_2\)，使得
\[
C_1 \|x\|_1 \leqslant \|x\|_2 \leqslant C_2 \|x\|_1 \quad (\forall x \in \mathscr{X}).
\]
\end{theorem}
\begin{proof}
设 \(\operatorname{dim} \mathscr{X} = n\)。因为 \(\|\cdot\|_1\) 与 \(\|\cdot\|_2\) 都与 \(\mathbb{K}^n\) 范数等价，所以 \(\|\cdot\|_1\) 与 \(\|\cdot\|_2\) 等价。
\end{proof}
\begin{corollary}
有穷维 ${B}^*$ 空间必是 ${B}$ 空间。
\end{corollary}
\begin{proof}
    考虑$\mathbb{K}^n$范数在${B}^*$空间上的收敛，这等价于$\mathbb{K}^n$上的收敛，所以极限存在且属于该线性空间。
\end{proof}
\begin{corollary}
${B}^*$ 空间上的任意有穷维子空间必是闭子空间。
\end{corollary}

\begin{definition}
设 $P: \mathscr{X} \to \mathbb{R}$ 是线性空间 $\mathscr{X}$ 上的一个函数，若它满足以下条件：
\begin{enumerate}
    \item \(P(x + y) \leqslant P(x) + P(y) \quad (\forall x, y \in \mathscr{X})\) \hfill (次可加性),
    \item \(P(\lambda x) = \lambda P(x) \quad (\forall \lambda > 0, \forall x \in \mathscr{X})\) \hfill (正齐次性),
\end{enumerate}
则称 $P$ 为 $\mathscr{X}$ 上的一个次线性泛函。
\end{definition}
\begin{note}如果 $P$ 还满足 $P(x) \geqslant 0 \ (\forall x \in \mathscr{X})$，并且将条件 (2) 替换为齐次性
\[
P(\alpha x) = |\alpha| P(x) \quad (\forall \alpha \in \mathbb{K}, \forall x \in \mathscr{X}),
\]
则称 $P$ 是一个半范数或半模。
\end{note}



\begin{theorem}
设 $P$ 是有穷维 ${B}^*$ 空间 $\mathscr{X}$ 上的一个连续次线性泛函，如果 $P(x) \geqslant 0 \ (\forall x \in \mathscr{X})$，并且
\[
P(x) = 0 \Longleftrightarrow x = \theta,
\]
则存在正常数 $C_1, C_2$，使得
\[
C_1 \|x\| \leqslant P(x) \leqslant C_2 \|x\| \quad (\forall x \in \mathscr{X}).
\]
\end{theorem}
\begin{proof}
记 $\|\cdot\|$ 为 $\mathscr{X}$ 上的范数。由于 $\|\cdot\|$ 和 $\mathscr{X}$ 上的 $\mathbb{K} ^n$ 范数等价, 故只需证明 $P(\cdot)$ 也和 $\mathscr{X}$ 上的 $\mathbb{K} ^n$ 范数等价。设 $\left\{e_i\right\}_{i=1}^n$ 是 $\mathscr{X}$ 的一组基。定义映射 $f: \mathbb{K} ^n \rightarrow \mathbb{K} , f(\alpha)=P\left(\sum_{i=1}^n \alpha_i e_i\right)$, 其中 $\alpha=\left(\alpha_1, \alpha_2, \cdots, \alpha_n\right) 。$\\
(i) 首先证明 $f$ 连续。
对 $\forall \alpha=\left(\alpha_1, \alpha_2, \cdots, \alpha_n\right), \beta=\left(\beta_1, \beta_2, \cdots, \beta_n\right) \in \mathbb{K} ^n$,
\begin{equation*}
\begin{aligned}
|f(\alpha)-f(\beta)| & =\left|P\left(\sum_{i=1}^n \alpha_i e_i\right)-P\left(\sum_{i=1}^n \beta_i e_i\right)\right| \\
& \leqslant \max \left\{P\left(\sum_{i=1}^n\left(\alpha_i-\beta_i\right) e_i\right), P\left(\sum_{i=1}^n\left(\beta_i-\alpha_i\right) e_i\right)\right\}
\end{aligned}
\end{equation*}
由于 $P$在$\theta$处 连续, 故 $f$ 连续。\\
(ii) 其次证明 $P(\cdot)$ 和 $\mathscr{X}$ 上的 $\mathbb{K} ^n$ 范数等价。\\
记 $S$ 是 $\mathbb{K} ^n$ 中的单位球面, 则 $S$ 是紧集。而 $f$ 在 $S$ 上连续, 故 $\exists \alpha_0, \alpha_1 \in S$, 使得
\begin{equation*}
0<f\left(\alpha_0\right)=\min _{\alpha \in S} f(\alpha) \leqslant f(\alpha) \leqslant \max _{\alpha \in S} f(\alpha)=f\left(\alpha_1\right) \quad(\forall \alpha \in S)
\end{equation*}
故
\begin{equation*}
f\left(\alpha_0\right) \leqslant f\left(\frac{\alpha}{\|\alpha\|}\right) \leqslant f\left(\alpha_1\right) \quad\left(\forall \alpha \in K ^n \backslash\{\theta\}\right)
\end{equation*}
即
\begin{equation*}
f\left(\alpha_0\right)\|\alpha\| \leqslant P\left(\sum_{i=1}^n \alpha_i e_i\right) \leqslant f\left(\alpha_1\right)\|\alpha\| \quad\left(\forall \alpha \in K ^n\right)
\end{equation*}
同理可证$f(\alpha_0)\neq 0$，因此 $P(\cdot)$ 和 $\mathscr{X}$ 上的 $\mathbb{K} ^n$ 范数等价, 从而和 $\|\cdot\|$ 等价。
\end{proof}


\begin{theorem}
设 $\mathscr{X}$ 是一个 ${B}^*$ 空间。若 $e_1, e_2, \cdots, e_n$ 是 $\mathscr{X}$ 中给定的向量组，则 $\forall x \in \mathscr{X}$，存在最佳逼近系数 $\lambda_1, \lambda_2, \cdots, \lambda_n$ 使得
\[
\left\|x - \sum_{i=1}^n \lambda_i e_i\right\| = \min_{a \in \mathbb{K}^n} \left\|x - \sum_{i=1}^n a_i e_i\right\|.
\]
\end{theorem}
\begin{proof}
给定一个 $B^*$ 空间 $\mathscr{X}$，并给定 $\mathscr{X}$ 中的有穷个向量 $e_1, e_2, \cdots, e_n$。对于给定的向量 $x \in \mathscr{X}$，求一组数 $\left(\lambda_1, \lambda_2, \cdots, \lambda_n\right) \in \mathbb{K}^n$ ，使得
\begin{equation}
\left\|x-\sum_{i=1}^n \lambda_i e_i\right\|=\min _{a \in \mathbb{K}^n}\left\|x-\sum_{i=1}^n a_i e_i\right\|,
\end{equation}
其中 $a = \left(a_1, a_2, \cdots, a_n\right)$。\\
显然，我们不妨设 $e_1, e_2, \cdots, e_n$ 是线性无关的。为此，我们要求函数
\[
F(a)=\left\|x-\sum_{i=1}^n a_i e_i\right\| \quad \left(a \in \mathbb{K}^n\right)
\]
的最小值。\\
容易看出 $F$ 是 $\mathbb{K}^n$ 上的连续函数。又注意到
\[
F(a) \geqslant\left\|\sum_{i=1}^n a_i e_i\right\|-\|x\| \quad \left(\forall a \in \mathbb{K}^n\right).
\]
令
\[
P(a) \triangleq\left\|\sum_{i=1}^n a_i e_i\right\|,
\]
显然 $P(\cdot)$ 是 $\mathbb{K}^n$ 上的一个范数。而 $\mathbb{K}^n$ 是有穷维空间，根据定理 1.14，$\exists c_1 > 0$ ，使得
\[
P(a) \geqslant c_1 |a| \quad \left(\forall a \in \mathbb{K}^n\right),
\]
其中
\[
|a| \triangleq\left(\left|a_1\right|^2 + \left|a_2\right|^2 + \cdots + \left|a_n\right|^2\right)^{\frac{1}{2}} \quad \left(\forall a = \left(a_1, a_2, \cdots, a_n\right) \in \mathbb{K}^n\right).
\]
联合可得 $F(a) \to \infty$ （当 $|a| \to \infty$）。由此得出，函数 $F$ 存在最小值。
\end{proof}
\begin{definition}
设 $\mathscr{B}^*$ 空间 $(\mathscr{X}, \|\cdot\|)$，若满足以下性质，则称 $\mathscr{X}$ 是\textbf{严格凸}的：  
$\forall x, y \in \mathscr{X}, x \neq y$，必有  
\[
\|x\| = \|y\| = 1 \Longrightarrow \|\alpha x + \beta y\| < 1 \quad (\forall \alpha, \beta > 0, \alpha + \beta = 1).
\]
\end{definition}

\begin{theorem}
设 $\mathscr{X}$ 是严格凸的 ${B}^*$ 空间，$\{e_1, e_2, \cdots, e_n\}$ 是 $\mathscr{X}$ 上的一组线性无关向量，则 $\forall x \in \mathscr{X}$，存在唯一的一组最佳逼近系数 $\{\lambda_1, \lambda_2, \cdots, \lambda_n\}$ 满足
\[
\left\|x - \sum_{i=1}^n \lambda_i e_i \right\| = \min_{a \in \mathbb{K}^n} \left\|x - \sum_{i=1}^n a_i e_i \right\|.
\]
\end{theorem}
\begin{proof}
 倘若 $d \triangleq \rho(x, M)>0$, 并有 $\|x-y\|=\|x-z\|=d$,则对 $\forall \alpha, \beta>0, \alpha+\beta=1$ ，由严格凸性有
\begin{equation*}
\begin{aligned}
\frac{1}{d}\|x-(\alpha y+\beta z)\| & =\frac{1}{d}\|\alpha(x-y)+\beta(x-z)\| \\
& =\left\|\alpha\left(\frac{x-y}{d}\right)+\beta\left(\frac{x-z}{d}\right)\right\|<1
\end{aligned}
\end{equation*}
即 $\|x-(\alpha y+\beta z)\|<d$. 这显然与 $d$ 的定义矛盾. 但若 $d=0, y$ 是相应的最佳逼近元, 则必有 $\|x-y\|=0$, 即 $y=x$. 从而最佳逼近元必是唯一的.
\end{proof}
\begin{example}
空间 $L^p(\Omega, \mu)$（$1 < p < \infty$）是严格凸的。
\end{example}

\begin{proof}
因为 Minkowski 不等式 $\|u + v\| \leqslant \|u\| + \|v\|$ 等号成立的充要条件是：存在 $k_1, k_2 \geqslant 0$ 且 $k_1 + k_2 > 0$，使得 $k_1 u = k_2 v \;(\text{a.e.})$。  
这意味着 $\forall u, v \in L^p(\Omega, \mu)$，当 $\|u\| = \|v\| = 1, u \neq v$ 时，有  
\[
\|t u + (1-t) v\| < t\|u\| + (1-t)\|v\| = 1 \quad (\forall 0 < t < 1).
\]
因此，$L^p(\Omega, \mu)$ 是严格凸的。
\end{proof}
\begin{theorem}
为了 $B^*$ 空间 $\mathscr{X}$ 是有穷维的，必须且仅须 $\mathscr{X}$的单位球面是列紧的.
\end{theorem}
\begin{proof}
有穷维 $B^*$ 空间 $\mathscr{X}$ 上的单位球面
\begin{equation*}
S_1 \triangleq\{x \in \mathscr{X} \mid\|x\|=1\}
\end{equation*}
是列紧的(取$\mathbb{K}^n$范数)，\\
如果一个 $B^*$ 空间 $\mathscr{X}$ 的单位球面是列紧的，那么这空间必是有穷维的。倘若在 $S_1$ 上给定了有穷个线性无关的向量 $\left\{x_1, x_2, \cdots, x_n\right\}$ ，如果它们的线性包 $M_n$张不满 $\mathscr{X}$ ，那么 $\exists x_{n+1} \in S_1$ ，使得
\begin{equation*}
\left\|x_{n+1}-x_i\right\| \geqslant 1 \quad(i=1,2, \cdots, n)
\end{equation*}
这是因为任取 $y \in M_n$, 按定理 1.15, $\exists x \in M_n$, 使得
\begin{equation*}
\|y-x\|=d \triangleq \rho\left(y, M_n\right) .
\end{equation*}
令 $x_{n+1} \triangleq(y-x) / d$, 显然 $x_{n+1} \in S_1$, 并且
\begin{equation*}
\begin{gathered}
\left\|x_{n+1}-x_i\right\|=\frac{1}{d}\left\|y-\left(x+d x_i\right)\right\| \geqslant \frac{1}{d} d=1 \\
(i=1,2, \cdots, n) .
\end{gathered}
\end{equation*}
照此办法, 如果 $\mathscr{X}$ 是无穷维的, 我们便可以逐次在 $S_1$ 上抽选出一串 $\left\{x_n\right\}_{n=1}^{\infty}$ 适合 $\left\|x_n-x_m\right\| \geqslant 1(n, m \in N , n \neq m)$ 。这样 $S_1$ 就不是列紧的。
\end{proof}
\begin{definition}
设 ${B}^*$ 空间 $\mathscr{X}$ 上的一个子集 $A$，如果存在常数 $c > 0$，使得 $\|x\| \leqslant c \; (\forall x \in A)$，则称 $A$ 是\textbf{有界}的。
\end{definition}

\begin{corollary}
为了 ${B}^*$ 空间 $\mathscr{X}$ 是有限维的，必须且仅需其任意有界集是列紧的。
\end{corollary}
\begin{proof}
    充分性，由单位球面列紧即可知。\\
    必要性，$\mathbb{K}^n$中的有界集是列紧的。
\end{proof}

\begin{lemma}[Riesz 引理]
如果 $\mathscr{X}_0$ 是 ${B}^*$ 空间 $\mathscr{X}$ 的一个真闭子空间，那么对 $\forall 0 < \varepsilon < 1$，$\exists y \in \mathscr{X}$，使得 $\|y\| = 1$，并且
\[
\|y - x\| \geqslant 1 - \varepsilon \quad (\forall x \in \mathscr{X}_0).
\]
\end{lemma}

\begin{proof}
任取 $y_0 \in \mathscr{X} \backslash \mathscr{X}_0$。因为 $\mathscr{X}_0$ 是闭的，所以
\[
d \triangleq \inf_{x \in \mathscr{X}_0} \|y_0 - x\| > 0.
\]
因此，对任意 $\eta > 0$，存在 $x_0 \in \mathscr{X}_0$ 使得 $d \leqslant \|y_0 - x_0\| < d + \eta$。令
\[
y \triangleq \frac{y_0 - x_0}{\|y_0 - x_0\|},
\]
则 $\|y\| = 1$，并且对 $\forall x \in \mathscr{X}_0$，令 $x^{\prime} \triangleq x_0 + \|y_0 - x_0\|x \in \mathscr{X}_0$，有
\[
\|y - x\| = \frac{\|y_0 - x^{\prime}\|}{\|y_0 - x_0\|} > \frac{d}{d + \eta} = 1 - \frac{\eta}{d + \eta}.
\]
令 $\eta = d\varepsilon / (1 - \varepsilon)$，即可得 $\|y - x\| > 1 - \varepsilon$。
\end{proof}
\begin{definition}
设 $\mathscr{X}$ 是一个 $B^*$ 空间，$\mathscr{X}_0 \subset \mathscr{X}$ 是一个闭线性子空间。定义 $x' \sim x''$，如果 $x' - x'' \in \mathscr{X}_0$，这是一种等价关系。记 $[x]$ 为 $x$ 所在的等价类，所有等价类组成的空间称为关于 $\mathscr{X}_0$ 的\textbf{商空间}，记为 $\mathscr{X} / \mathscr{X}_0$。其中，加法和数乘定义如下：
\begin{enumerate}
    \item $[x] + [y] = [x + y], \quad \forall [x], [y] \in \mathscr{X} / \mathscr{X}_0$；
    \item $\lambda [x] = [\lambda x], \quad \forall \lambda \in \mathbb{K}, [x] \in \mathscr{X} / \mathscr{X}_0$。
\end{enumerate}
\end{definition}

\begin{note}容易验证，$\mathscr{X} / \mathscr{X}_0$ 关于上述运算构成一个线性空间。
\end{note}
\begin{theorem}
设 $[x] \in \mathscr{X} / \mathscr{X}_0$，定义
\[
\|[x]\|_0 = \inf_{y \in [x]} \|y\|,
\]
则 $\left(\mathscr{X} / \mathscr{X}_0, \|\cdot\|_0\right)$ 是一个 $B^*$ 空间。当 $\mathscr{X}$ 是 $B$ 空间时，$\left(\mathscr{X} / \mathscr{X}_0, \|\cdot\|_0\right)$ 也是 $B$ 空间。
\end{theorem}

\begin{proof}
根据定义，$\mathscr{X} / \mathscr{X}_0$ 中的零元素为 $[x]$，其中 $x \in \mathscr{X}_0$。以下性质显然成立：
\[
\|[x]\|_0 \geqslant 0, \quad \|\lambda [x]\|_0 = |\lambda| \|[x]\|_0, \quad \forall [x] \in \mathscr{X} / \mathscr{X}_0, \lambda \in \mathbb{K}.
\]
接下来验证三角不等式。注意：
\[
\|[x] + [y]\|_0 = \|[x + y]\|_0 = \inf_{z \in [x + y]} \|z\|.
\]
取 $x_n \in [x], y_n \in [y]$，使得
\[
\|x_n\| \to \|[x]\|_0, \quad \|y_n\| \to \|[y]\|_0, \quad (n \to \infty).
\]
此时 $x_n + y_n \in [x + y]$，因此
\[
\|[x] + [y]\|_0 \leqslant \|x_n + y_n\| \leqslant \|x_n\| + \|y_n\| \to \|[x]\|_0 + \|[y]\|_0 \quad (n \to \infty).
\]
由此，三角不等式
\[
\|[x] + [y]\|_0 \leqslant \|[x]\|_0 + \|[y]\|_0
\]
成立。\\
若 $\|[x]\|_0 = 0$，根据定义，存在 $x_n \in [x]$，使得 $\|x_n\| \to 0$ $(n \to \infty)$。由于 $\mathscr{X}_0$ 是闭的，且 $x - x_n \in \mathscr{X}_0$，可知
\[
x = \lim_{n \to \infty} (x - x_n) \in \mathscr{X}_0,
\]
即 $[x] = \theta$ 为 $\mathscr{X} / \mathscr{X}_0$ 中的零元素。因此 $\left(\mathscr{X} / \mathscr{X}_0, \|\cdot\|_0\right)$ 是一个 $B^*$ 空间。\\
再设 $\mathscr{X}$ 是 $B$ 空间，下证 $\left(\mathscr{X} / \mathscr{X}_0, \|\cdot\|_0\right)$ 完备。令 $\{\left[x_n\right]\}$ 是 Cauchy 列，则有
\[
\|\left[x_n\right] - \left[x_m\right]\|_0 = \|[x_n - x_m]\|_0 \to 0 \quad (n, m \to \infty).
\]
取子列（仍记为 $\{\left[x_{n_k}\right]\}$），使得
\[
\|[x_{n_k} - x_{n_{k+1}}]\|_0 \leqslant \frac{1}{2^{k+1}}.
\]
根据定义，存在 $y_{k, k+1} \in [x_{n_k} - x_{n_{k+1}}]$，满足
\[
\|y_{k, k+1}\| \leqslant \|[x_{n_k}\ - x_{n_{k+1}}]\|_0 + \frac{1}{2^{k+1}} \leqslant \frac{1}{2^k}.
\]
令 $y_1 = x_{n_1}, y_{k+1} = y_k - y_{k, k+1}$，并存在 $z_{k, k+1} \in \mathscr{X}_0$，使得
\[
y_{k+1} = y_k - (x_{n_k} - x_{n_{k+1}} + z_{k, k+1}),
\]
从而
\[
y_{k+1} - x_{n_{k+1}} = y_k - x_{n_k} - z_{k, k+1} \in \mathscr{X}_0.
\]
于是 $\left[y_{k+1}\right] = \left[x_{k+1}\right]$。\\
验证 $\{y_n\}$ 是 $\mathscr{X}$ 中的 Cauchy 列：
\[
\|y_n - y_{n+p}\| \leqslant \sum_{k=n}^{n+p-1} \|y_{k, k+1}\| \leqslant \sum_{k=n}^{n+p-1} \frac{1}{2^k} \leqslant \frac{1}{2^n}.
\]
因此，$\{y_n\}$ 收敛于 $y \in \mathscr{X}$，并且
\[
\|[x_n] - [y]\|_0 = \|[y_n] - [y]\|_0 = \|[y_n - y]\|_0 \leqslant \|y_n - y\| \to 0 \quad (n \to \infty).
\]
故 $\left(\mathscr{X} / \mathscr{X}_0, \|\cdot\|_0\right)$ 完备。
\end{proof}
\begin{example}
用 $C(0,1]$ 表示 $(0,1]$ 上连续且有界的函数 $x(t)$ 全体。对 $\forall x \in C(0,1]$，定义范数 $\|x\| = \sup_{0 < t \leqslant 1} |x(t)|$。证明：
\begin{enumerate}
    \item $\|\cdot\|$ 是 $C(0,1]$ 空间上的范数；
    \item $l^\infty$ 与 $C(0,1]$ 的一个子空间是等距同构的。
\end{enumerate}
\end{example}
\begin{proof}
(1) 证明 $\|\cdot\|$ 是范数的性质略。\\
(2) 对 $\forall x \in C(0,1]$，构造点列：
\[
\alpha_x = \left\{ x(1), x\left(\frac{1}{2}\right), x\left(\frac{1}{3}\right), \cdots, x\left(\frac{1}{n}\right), \cdots \right\} \in l^\infty,
\]
并有
\[
\|\alpha_x\|_\infty = \sup_{n \geqslant 1} \left| x\left( \frac{1}{n} \right) \right| \leqslant \|x\|.
\]
反之，$\forall \alpha = \left( \xi_1, \xi_2, \cdots, \xi_n, \cdots \right) \in l^\infty$，将点列 $\left(1, \xi_1\right), \left(\frac{1}{2}, \xi_2\right), \cdots, \left(\frac{1}{n}, \xi_n\right), \cdots$ 用折线连接起来，得到一个函数 $x_\alpha(t)$，并有：
\[
x_\alpha(t) \in C(0,1], \quad \|x_\alpha\| \leqslant \sup_{n \geqslant 1} |\xi_n| = \|\alpha\|_\infty.
\]
由此可得：
\[
\|\alpha_x\|_\infty \leqslant \|x\|, \quad \|x_\alpha\| \leqslant \|\alpha\|_\infty \implies \|\alpha\|_\infty = \|x\|.
\]
因此，$l^\infty$ 与 $C(0,1]$ 的一个子空间(即折线函数空间)是等距同构的。
\end{proof}

\begin{example}
在 $C^1[-1,1]$ 中，定义范数：
\[
\|f\|_1 = \left( \int_{-1}^1 \left( |f(x)|^2 + |f'(x)|^2 \right) \dif x \right)^{\frac{1}{2}}, \quad \forall f \in C^1[-1,1].
\]
证明：$\left(C^1[-1,1], \|\cdot\|_1\right)$ 不完备。


\end{example}
\begin{proof}
\textbf{分析}：
为了证明 $\left(C^1[-1,1], \|\cdot\|_1\right)$ 不完备，只需在 $\left(C^1[-1,1], \|\cdot\|_1\right)$ 中找到一个不收敛的基本列。考虑函数列
\[
f_n(x) = \sqrt{x^2 + \frac{1}{n^2}}, \quad x \in [-1, 1].
\]
易知 $\{f_n(x)\}$ 点点收敛于 $f(x) = |x|$，而 $|x| \notin C^1[-1,1]$。以下分两步验证。\\
\textbf{第一步}：证明 $\{f_n(x)\}$ 是 $\|\cdot\|_1$ 范数下的基本列。
计算 $f_n'(x) = \frac{x}{\sqrt{x^2 + \frac{1}{n^2}}}$，设 $m > n$，则
\[
\|f_m - f_n\|_1^2 = 2 \int_0^1 \left[ \left( \sqrt{x^2 + \frac{1}{m^2}} - \sqrt{x^2 + \frac{1}{n^2}} \right)^2 
+ x^2 \left( \frac{1}{\sqrt{x^2 + \frac{1}{m^2}}} - \frac{1}{\sqrt{x^2 + \frac{1}{n^2}}} \right)^2 \right] \dif x.
\]
记 $I_1$ 和 $I_2$ 分别为两部分积分，易得 $I_1 \to 0$ 和 $I_2 \to 0$，从而 $\|f_m - f_n\|_1 \to 0$。\\
\textbf{第二步}：证明 $\{f_n(x)\}$ 按范数 $\|\cdot\|_1$ 收敛的极限不在 $C^1[-1,1]$ 中。
令 $f(x) = |x|$，计算
\[
\|f_n - f\|_1^2 = \int_{-1}^1 \left( |f_n(x) - f(x)|^2 + |f_n'(x) - f'(x)|^2 \right) \dif x,
\]
进一步分析可得 $\|f_n - f\|_1 \to 0$，即 $f_n(x) \to f(x)$，但 $f(x) \notin C^1[-1,1]$。因此，$\{f_n(x)\}$ 不收敛于 $\left(C^1[-1,1], \|\cdot\|_1\right)$ 中的元素。\\
综上，$\left(C^1[-1,1], \|\cdot\|_1\right)$ 中存在不收敛的基本列，故其不完备。
\end{proof}
\begin{example}
设 $B C[0, \infty)$ 表示 $[0, \infty)$ 上连续且有界的函数 $f(x)$ 的全体，对于每个 $f \in B C[0, \infty)$ 及 $a > 0$ ，定义
\[
\|f\|_a = \left( \int_0^{\infty} e^{-a x} |f(x)|^2 \, \dif x \right)^{\frac{1}{2}}.
\]
\textbf{(1)} 证明 $\|\cdot\|_a$ 是 $B C[0, \infty)$ 上的范数。\\
\textbf{(2)} 若 $a, b > 0, a \neq b$，证明 $\|\cdot\|_a, \|\cdot\|_b$ 作为 $B C[0, \infty)$ 上的范数是不等价的。\\
\end{example}

\begin{proof}
不妨假设 $b > a > 0$，显然有 $\|f\|_b \leqslant \|f\|_a$，由此可见，为了证明不等价性，只需证明不存在 $c > 0$，使得
\[
\|f\|_a \leqslant c \|f\|_b \quad (\forall f \in B C[0, \infty)).
\]
只需证明 $\exists f_n \in B C[0, \infty)$，使得
\[
\frac{\left\|f_n\right\|_a^2}{\left\|f_n\right\|_b^2} \to \infty \quad (n \to \infty).
\]
考虑
\[
g_n(x) = \begin{cases}
e^{a x}, & 0 \leqslant x \leqslant n, \\
e^{a n}(n+1-x), & n \leqslant x \leqslant n+1, \\
0, & x \geqslant n+1,
\end{cases}
\]
\[
f_n(x)= \sqrt{g_n(x)}.
\]
则有
\[
\left\|f_n\right\|_a^2 \geqslant \int_0^n e^{-a x} \cdot e^{a x} \, \dif x = n,
\]
\[
\left\|f_n\right\|_b^2 \leqslant \int_0^{\infty} e^{-b x} \cdot e^{a x} \, \dif x = \int_0^{\infty} e^{-(b-a) x} \, \dif x = \frac{1}{b-a}.
\]
因此，
\[
\frac{\left\|f_n\right\|_a^2}{\left\|f_n\right\|_b^2} \geqslant \frac{n}{b-a} \to \infty \quad (n \to \infty).
\]
\end{proof}
\begin{example}
设 $C_0$ 表示以 0 为极限的数列全体，并在 $C_0$ 中赋以范数
\[
\|x\| = \sup_{n \geqslant 1} \left|\xi_n\right| \quad (\forall x = (\xi_1, \xi_2, \cdots, \xi_n) \in C_0),
\]
又设
\[
M = \left\{ x = \left\{ \xi_n \right\}_{n=1}^{\infty} \in C_0 \mid \sum_{n=1}^{\infty} \frac{\xi_n}{2^n} = 0 \right\}.
\]
\textbf{(1)} 证明 $M$ 是 $C_0$ 的闭线性子空间。\\
\textbf{(2)} 设 $x_0 = (2, 0, \cdots, 0, \cdots)$，求证：
\[
\inf_{z \in M} \left\| x_0 - z \right\| = 1,
\]
但对于任意 $y \in M,\|x_0-y\|\geqslant 1$。
\end{example}


\begin{proof}
（1）设 $x^{(n)} = (\xi_1^{(n)}, \xi_2^{(n)}, \cdots, \xi_k^{(n)}, \cdots) \in M$，$x = (\xi_1, \xi_2, \cdots)$，则
\[
\lim_{n \to \infty} \sup_{k \geqslant 1} \left|\xi_k^{(n)} - \xi_k\right| = 0.
\]
因此，$\forall \varepsilon > 0, \sup_{k \geqslant 1} \left|\xi_k^{(n)} - \xi_k\right| < \varepsilon \quad (n \geqslant N)$，得
\[
\sum_{k=1}^{\infty} \frac{\xi_k}{2^k} = \sum_{k=1}^{\infty} \frac{\xi_k - \xi_k^{(N)}}{2^k} + {\sum_{k=1}^{\infty} \frac{\xi_k^{(N)}}{2^k}}
\]
由于 $\sum_{k=1}^{\infty} \frac{\xi_k^{(N)}}{2^k} = 0$，因此
\[
|\sum_{k=1}^{\infty} \frac{\xi_k}{2^k}| <\varepsilon\Rightarrow |\sum_{k=1}^{\infty} \frac{\xi_k}{2^k}|=0\Rightarrow x\in M
\]
所以 $x = (\xi_1, \xi_2, \cdots, \xi_k, \cdots) \in M$。\\
（2）设 $x_0 = (2, 0, \cdots, 0, \cdots)$，对任意 $m \in \mathbb{N}$，定义
\[
x^{(m)}= \left( 1 - \frac{1}{2^{m-1}}, -1, -1, \cdots, -1, 0, 0, \cdots \right) \in M.
\]
则
\[
\rho(x_0, x^{(m)}) = 1 + \frac{1}{2^{m-1}} \quad \text{当} \quad m \to \infty \quad \Rightarrow \quad \rho(x_0, M) \leqslant 1.
\]
接下来，假设 $\exists y = (\xi_1, \xi_2, \cdots, \xi_k, \cdots) \in M$，使得 $\left\| x_0 - y \right\| \leqslant 1$，则
\[
x_0 - y = \left( 2 - \xi_1, -\xi_2, -\xi_3, \cdots, -\xi_k, \cdots \right),
\]
且
\[
\left\| x_0 - y \right\| \leqslant 1 \quad \Rightarrow \quad \left\{
\begin{array}{l}
|\xi_k| \leqslant 1, \quad k \geqslant 2, \\
2 - \xi_1 \leqslant 1.
\end{array}
\right.
\]
这说明 $\xi_1 \geqslant 1$，
\[
\frac{1}{2}\xi_1\geqslant \frac{1}{2}=\sum\limits_{k=2}^\infty\frac{1}{2^k}> -\sum\limits_{k=2}^\infty \frac{\xi_k}{2^k}
\]
因为当$k$足够大时$|\xi_k|<1$，所以$y\notin M$矛盾，因此，$\forall y \in M, \left\| x_0 - y \right\| > 1$。最终，
\[
\rho(x_0, M) = 1.
\]
\end{proof}
\begin{example}
设 $\mathscr{X}$ 是 $B^*$ 空间，$M$ 是 $\mathscr{X}$ 的有限维真子空间，求证：存在 $y \in \mathscr{X}$ 且 $\|y\| = 1$，使得对于任意 $x \in M$，有 $\|y - x\| \geqslant 1$。
\end{example}

\begin{proof}
首先，设 $y_0 \in \mathscr{X} \backslash M$，定义
\[
d =\inf_{x \in M} \left\| y_0 - x \right\|.
\]
显然 $d > 0$，因为 $y_0 \notin M$ 且 $M$ 是有限维的(闭的)，存在一个正的距离。由此可得，对于任意 $n \in \mathbb{N}$，存在 $x_n \in M$，使得
\[
d \leqslant \| y_0 - x_n \| < d + \frac{1}{n}.
\]
同时，有
\[
\| x_n \| \leqslant \| y_0 - x_n \| + \| y_0 \| \leqslant \| y_0 \| + d + 1,
\]
这说明 $\{x_n\}$ 是有界的。由于 $M$ 是有限维的，$\{x_n\}$ 必然有一个收敛子列，不妨设整个序列都收敛，且设
\[
x_n \to x_0 \in M \quad \text{当} \quad n \to \infty.
\]
因此，由于
\[
d \leqslant \| y_0 - x_n \| < d + \frac{1}{n} \quad \Rightarrow \quad \| y_0 - x_0 \| = d,
\]
我们可以定义
\[
y =\frac{y_0 - x_0}{d}.
\]
此时 $\| y \| = 1$。接下来，证明对于任意 $x \in M$，都有
\[
\| y - x \| = \left\| \frac{y_0 - x_0}{d} - x \right\| = \frac{1}{d} \| y_0 - (x_0 + d x) \| \geqslant \frac{d}{d} = 1.
\]
因此，对于所有 $x \in M$，都有 $\| y - x \| \geqslant 1$，这就证明了所需结论。
\end{proof}
\begin{example}
(1) 定义映射 $\varphi: \mathscr{X} \rightarrow \mathscr{X} / \mathscr{X}_0$ 为
\[
\varphi(x) = [x],
\]
求证 $\varphi$ 是线性连续映射。
(2) 对于任意 $[x] \in \mathscr{X} / \mathscr{X}_0$，求证：存在 $x \in \mathscr{X}$，使得 $\varphi(x) = [x]$，并且 $\|x\| \leqslant 2\|[x]\|_0$。
\end{example}

\begin{proof}
(1) 证明 $\varphi$ 是线性连续映射：\\
首先，显然映射 $\varphi(x) = [x]$ 是线性的，因为对于任意的 $x, y \in \mathscr{X}$ 和标量 $\alpha, \beta$，有：
\[
\varphi(\alpha x + \beta y) = [\alpha x + \beta y] = \alpha [x] + \beta [y] = \alpha \varphi(x) + \beta \varphi(y).
\]
接下来证明 $\varphi$ 连续。对于任意 $x \in \mathscr{X}$，有：
\[
\|\varphi(x)\| = \|[x]\| = \inf_{z \in [x]} \|z\| \leqslant \|x\|.
\]
因此，$\varphi$ 是连续的。\\
(2) 由定义，$\|[x]\| = \inf_{z \in [x]} \|z\|$，且 $\|[x]\|$ 是下确界。由于 $[x] \neq [\theta]$，我们可以找到一个 $z \in [x]$，使得 $\varphi(z) = [x]$ 且 $\|z\| \leqslant 2\|[x]\|_0([x]\neq[\theta])$。\\
\end{proof}
\begin{example}
设 $\mathscr{X} = C[0,1]$，$\mathscr{X}_0 = \{f \in \mathscr{X} \mid f(0) = 0\}$，求证：$\mathscr{X} / \mathscr{X}_0$ 与 $\mathbb{K}$ 等距同构。
\end{example}

\begin{proof}
对任意 $[f] \in \mathscr{X} / \mathscr{X}_0$，设 $f_0 = f(0) \in \mathbb{K}$。定义映射
\[
T: \mathscr{X} / \mathscr{X}_0 \rightarrow \mathbb{K}, \quad T[f] = f_0.
\]
我们需要证明 $T$ 是 $\mathscr{X} / \mathscr{X}_0$ 到 $\mathbb{K}$ 的同构映射，即证明
\[
|T[f]| = \|[f]\|.
\]
首先，注意到 $f(x) \equiv f_0 \in [f]$，因此有
\[
\|[f]\| = \inf_{f' \in [f]} \|f'\| \leqslant |f_0|.
\]
另一方面，对于任意 $\varepsilon > 0$，存在 $f_1 \in [f]$，使得
\[
\|[f]\| + \varepsilon > \|f_1\| = \max_{t \in [0,1]} |f_1(t)| \geqslant |f_1(0)| = |f_0|.
\]
当 $\varepsilon \rightarrow 0$ 时，我们得到
\[
\|[f]\| \geqslant |f_0|.
\]
综上所述，得出结论 $\|[f]\| = |f_0|$，即 $|T[f]| = |f_0|$。
\end{proof}
\section{凸集与不动点}

\begin{definition}
设 $\mathscr{X}$ 是线性空间，$E \subset \mathscr{X}$，称 $E$ 为一凸集，如果对于任意的 $x, y \in E$ 和任意的 $0 \leqslant \lambda \leqslant 1$，都有
\[
\lambda x + (1-\lambda) y \in E.
\]
\end{definition}

\begin{proposition}
若 $\{E_\lambda \mid \lambda \in \Lambda\}$ 是线性空间 $\mathscr{X}$ 中的一族凸集，则
\[
\bigcap_{\lambda \in \Lambda} E_\lambda
\]
也是凸集。
\end{proposition}


\begin{definition}
设 $\mathscr{X}$ 是线性空间，$A \subset \mathscr{X}$。若 $\{E_\lambda \mid \lambda \in \Lambda\}$ 为 $\mathscr{X}$ 中包含 $A$ 的一切凸集，那么称 $\bigcap_{\lambda \in \Lambda} E_\lambda$ 为 $A$ 的凸包，并记作 $\operatorname{co}(A)$。对于任意 $n \in \mathbb{N}$，若 $x_1, x_2, \dots, x_n \in A$，称
\[
\sum_{i=1}^n \lambda_i x_i
\]
为 $x_1, x_2, \dots, x_n$ 的凸组合，其中系数满足 $\lambda_i \geqslant 0$ 且
\[
\sum_{i=1}^n \lambda_i = 1。
\]
\end{definition}

\begin{proposition}
设 $\mathscr{X}$ 是线性空间，$A \subset \mathscr{X}$，则 $A$ 的凸包是 $A$ 中元素任意凸组合的全体，即
\[
\operatorname{co}(A) = \left\{ \sum_{i=1}^n \lambda_i x_i \mid \sum_{i=1}^n \lambda_i = 1, \lambda_i \geqslant 0, x_i \in A, i=1, 2, \dots, n, \forall n \in \mathbb{N} \right\} .
\]
\end{proposition}
\begin{proof}
 若令 $S$ 表示上式右端, 则 $A \subset S$ 而且 $S$ 是凸集，从而 $S \supset \operatorname{co}(A)$ 。反之，设 $F$ 为包含 $A$ 的任一凸集，那么 $x_i \in F(i=1,2, \cdots, n)$, 从而 $\sum_{i=1}^n \lambda_i x_i \in F$, 即得 $S \subset F$, 从而 $S \subset \operatorname{co}(A)$.
\end{proof}
\begin{definition}
设 $\mathscr{X}$ 是线性空间，$C$ 是 $\mathscr{X}$ 上含有零元素 $\theta$ 的凸子集。我们在 $\mathscr{X}$ 上定义一个取值于 $[0, \infty]$ 的函数
\[
P(x) = \inf \left\{ \lambda > 0 \mid \frac{x}{\lambda} \in C \right\} \quad (\forall x \in \mathscr{X}).
\]
与 $C$ 对应，称函数 $P$ 为 $C$ 的 Minkowski 泛函。
\end{definition}

\begin{proposition}
设 $\mathscr{X}$ 是线性空间，$C$ 是 $\mathscr{X}$ 上含有零元素 $\theta$ 的凸子集。若 $P$ 为 $C$ 的 Minkowski 泛函，则 $P$ 具有下列性质：
\begin{itemize}
    \item (1) $P(x) \in [0, \infty], P(\theta) = 0$；
    \item (2) $P(\lambda x) = \lambda P(x) \quad (\forall x \in \mathscr{X}, \forall \lambda > 0)$，(正齐次性)；
    \item (3) $P(x + y) \leqslant P(x) + P(y) \quad (\forall x, y \in \mathscr{X})$，(次可加性)。
\end{itemize}
\end{proposition}

\begin{proof}
只有 (3) 需要验证。假设 $P(x), P(y)$ 有穷。对于任意的 $\varepsilon > 0$，取 $\lambda_1 = P(x) + \frac{\varepsilon}{2}$ 和 $\lambda_2 = P(y) + \frac{\varepsilon}{2}$，则有
\[
\frac{x}{\lambda_1} \in C, \quad \frac{y}{\lambda_2} \in C.
\]
因为 $C$ 是凸集，所以
\[
\frac{x + y}{\lambda_1 + \lambda_2} = \frac{\lambda_1}{\lambda_1 + \lambda_2} \cdot \frac{x}{\lambda_1} + \frac{\lambda_2}{\lambda_1 + \lambda_2} \cdot \frac{y}{\lambda_2} \in C.
\]
这表明
\[
P(x + y) \leqslant \lambda_1 + \lambda_2 = P(x) + P(y) + \varepsilon.
\]
由于 $\varepsilon > 0$ 可以任意小，得到
\[
P(x + y) \leqslant P(x) + P(y).
\]
从而验证了 (3)。
\end{proof}

\begin{definition}
线性空间 $\mathscr{X}$ 中，含有零元素 $\theta$ 的凸集 $C$ 称为吸收的，如果对于任意的 $x \in \mathscr{X}$，存在 $\lambda > 0$ 使得 $\frac{x}{\lambda} \in C$；称 $C$ 是对称的，如果 $x \in C$ 蕴含 $-x \in C$。
\end{definition}

\begin{proposition}
为了 $C$ 是吸收凸集，必须且仅须其 Minkowski 泛函 $P(x)$ 是实值函数；为了 $C$ 是对称凸集，必须 $P(x)$ 是实齐次的，即
\[
P(\alpha x) = |\alpha| P(x) \quad (\forall \alpha \in \mathbb{R}).
\]
\end{proposition}
\begin{proposition}
设 $\mathscr{X}$ 是一个 $B^*$ 空间，$C$ 是一个含有零点 $\theta$ 的闭凸集。如果 $P(x)$ 是 $C$ 的 Minkowski 泛函，那么 $P(x)$ 下半连续，且有
\[
C = \{x \in \mathscr{X} \mid P(x) \leqslant 1\}.
\]
此外，如果 $C$ 还是有界的，那么
\[
P(x) = 0 \Longleftrightarrow x = \theta.
\]
若 $C$ 以 $\theta$ 为一内点，那么 $C$ 是吸收的，并且 $P(x)$ 还是一致连续的。
\end{proposition}

\begin{proof}
(1) 对于任意 $\alpha > 0$，若 $x \in \alpha C$，即 $\frac{x}{\alpha} \in C$，由 $P(x)$ 的定义便有 $P(x) \leqslant \alpha$。反之，若 $P(x) \leqslant \alpha$，则对于任意的 $n \in \mathbb{N}$，有
\[
\frac{x}{\alpha + \frac{1}{n}} \in C, \quad \frac{x}{\alpha + \frac{1}{n}} \rightarrow \frac{x}{\alpha} \quad (n \rightarrow \infty).
\]
又因为 $C$ 是闭集，所以 $x \in \alpha C$。这就证明了
\[
\alpha C = \{x \in \mathscr{X} \mid P(x) \leqslant \alpha\} \quad (\forall \alpha > 0).
\]
特别地，令 $\alpha = 1$ 即得
\[
C = \{x \in \mathscr{X} \mid P(x) \leqslant 1\},
\]
也即得到了命题中的式 1。又因为对于任意的 $\alpha > 0$，$\alpha C$ 是闭集，所以 $P$ 是下半连续的。\\
(2) 由于 $\theta \in C$，所以显然 $P(\theta) = 0$。在 $C$ 有界的假设下，即存在 $r > 0$，使得 $C \subset B(\theta, r)$，于是
\[
\forall x \in \mathscr{X} \setminus \{\theta\} \quad \Rightarrow \quad \frac{r x}{\|x\|} \in C \quad \Rightarrow \quad P(x) \geqslant \frac{\|x\|}{r}.
\]
由此可见，当 $P(x) = 0$ 时，$x = \theta$。\\
(3) 若 $C$ 以 $\theta$ 为内点，即存在 $r > 0$，使得 $B(\theta, r) \subset C$，那么
\[
\frac{r x}{2 \|x\|} \in C \quad (\forall x \in \mathscr{X} \setminus \{\theta\}).
\]
因此，$C$ 是吸收的，并且有
\[
P(x) \leqslant \frac{2 \|x\|}{r} \quad (\forall x \in \mathscr{X}).
\]
于是
\[
|P(x) - P(y)| \leqslant \max(P(x - y), P(y - x)) \leqslant \frac{2}{r} \|x - y\|,
\]
对于任意的 $x, y \in \mathscr{X}$。从而 $P(x)$ 是一致连续的。
\end{proof}
\begin{proposition}
若 $C$ 是 $\mathbb{R}^n$ 中的一个紧凸子集，则必存在正整数 $m \leqslant n$，使得 $C$ 同胚于 $\mathbb{R}^m$ 中的单位球。
\end{proposition}

\begin{proof}
(1) 用 $E$ 表示包含 $C$ 的最小闭线性子流形。设其维数是 $m(\leqslant n)$。于是，在 $C$ 上必有 $m+1$ 个向量 $e_1, e_2, \cdots, e_m, e_{m+1}$，使得 $e_i - e_{m+1}$ ($i = 1, 2, \cdots, m$) 是线性无关的。\\
(2) 令
\[
e_0 = \frac{1}{m+1} \sum_{i=1}^{m+1} e_i.
\]
因为 $e_0 \in C \subset E$，所以 $E - e_0$ 是一个 $m$ 维线性子空间。于是，对于任意的 $y \in E$，存在唯一的表示
\[
y = \sum_{i=1}^m \mu_i (e_i - e_0) + e_0,
\]
并在 $E - e_0$ 上可以引入一个等价范数
\[
\|z\| = \left(\sum_{i=1}^m |\mu_i|^2\right)^{\frac{1}{2}} \quad (z = y - e_0, y \in E).
\]
(3) 我们要证明：当$\|z\|$ 足够小时，蕴含上式所表示的 $y \in C$。事实上，因为
\[
\begin{aligned}
y = & \sum_{i=1}^m \mu_i e_i + \left(1 - \sum_{j=1}^m \mu_j \right) e_0 \\
= & \sum_{i=1}^m \left[\mu_i + \frac{1}{m+1}\left(1 - \sum_{j=1}^m \mu_j\right)\right] e_i  + \frac{1}{m+1}\left(1 - \sum_{j=1}^m \mu_j\right) e_{m+1},
\end{aligned}
\]
当 $\left|\mu_j\right|$ ($j=1,2,\cdots,m$) 足够小时，上式右端各项系数都是正的，并且各项系数的总和满足
\[
\sum_{i=1}^m \left[\mu_i + \frac{1}{m+1}\left(1 - \sum_{j=1}^m \mu_j\right)\right] + \frac{1}{m+1} \left(1 - \sum_{j=1}^m \mu_j\right) = 1,
\]
所以 $y \in \operatorname{co} \{e_1, e_2, \cdots, e_m, e_{m+1}\} \subset C$，由此说明$\theta$为$C-e_0$的内点。\\
(4) 在 $E - e_0$ 上，$C - e_0$ 是一个以 $\theta$ 为内点的有界闭凸集，它的 Minkowski 泛函 $P(z)$ 是 $E - e_0$ 上的一个一致连续、正齐次、次可加泛函，适合 $P(z) = 0 \Longleftrightarrow z = \theta$。应用定理 1.14，存在常数 $C_1, C_2 > 0$，使得
\[
C_1 \|z\| \leqslant P(z) \leqslant C_2 \|z\| \quad (\forall z \in E - e_0).
\]
设 $B^m(\theta, 1)$ 是 $E - e_0$ 中的单位球，若令
\[
\varphi(z) = e_0 + 
\begin{cases}
\|z\| \frac{z}{P(z)}, & z \neq \theta, \\
0, & z = \theta,
\end{cases}
\]
\[
\|\varphi(z_1)-\varphi(z_2)\|=\|k_1z_1-k_2z_2\|\leqslant |k_2|\|z_1-z_2\|+|k_1-k_2|\|z_1\|
\]
当$z_2\to z_1$时，$\|z_1-z_2\|\to 0,|k_1-k_2|\to 0$。所以 $\varphi: B^m(\theta, 1) \rightarrow C$ 是一个连续映射。\\
下证$\varphi$是单射，设$z_1\neq z_2$，若线性无关则$\varphi(z_1),\varphi(z_2)$如果相等的话那么与$z_1,z_2$线性无关矛盾，若线性相关则$\frac{z_1}{P(z_1)}=\frac{z_2}{P(z_2)}$，所以$\varphi(z_1)\neq \varphi(z_2)$。\\
最后证是满射，设$y=y'-e_0\in C-e_0$，由于$C-e_0$是闭集，所以$\frac{y}{P(y)}\in C-e_0$
，而$\frac{y}{\|y\|}\in B^m(\theta,1)$令$z=\frac{P(y)y}{\|y\|}$，则$P(z)=\frac{P(y)^2}{\|y\|},\varphi(z)=e_0+{P(y)}\frac{z}{P(z)}=e_0+y=y'
$
，所以$\varphi$是满射。\\
综上所述， $C$ 同胚于 $\mathbb{R}^m$ 中的单位球。



\end{proof}
\begin{theorem}[Brouwer]
设 $B$ 是 $\mathbb{R}^n$ 中的闭单位球，又设 $T: B \rightarrow B$ 是一个连续映射，那么 $T$ 必有一个不动点 $x \in B$。
\end{theorem}

\begin{corollary}
设 $C$ 是 $\mathbb{R}^n$ 中的一个紧凸子集，$T: C \rightarrow C$ 是连续的，则 $T$ 必有一个在 $C$ 上的不动点。
\end{corollary}

\begin{proof}
由于 $C$ 与 $\mathbb{R}^m$ ($m \leqslant n$) 中的一个单位球同胚，记此同胚为 $\varphi: B^m(\theta, 1) \rightarrow C$。考察映射
\[
T_{\varphi} = \varphi^{-1} \circ T \circ \varphi.
\]
显然 $T_{\varphi}: B^m(\theta, 1) \rightarrow B^m(\theta, 1)$。对 $T_{\varphi}$ 应用 Brouwer 不动点定理，存在 $x \in B^m(\theta, 1)$，使得 $T_{\varphi}(x) = x$。\\
由此得到 $y = \varphi(x) \in C$ 是 $T$ 的不动点。
\end{proof}
\begin{theorem}[Schauder]
设 \(C\) 是 \(B^*\) 空间 \(\mathscr{X}\) 中的一个闭凸子集，\(T: C \rightarrow C\) 连续且 \(T(C)\) 列紧，则 \(T\) 在 \(C\) 上必有一个不动点。
\end{theorem}

\begin{proof}
(1) 因为 \(T(C)\) 是列紧集，所以对 \(\forall n \in \mathbb{N}\)，存在 \(1/n\) 网 \(N_n = \{y_1, y_2, \cdots, y_{r_n}\}\)，即
\[
T(C) \subset \bigcup_{i=1}^{r_n} B\left(y_i, \frac{1}{n}\right) \quad\left(y_i \in T(C), i=1,2,\cdots,r_n\right).
\]
记 \(E_n = \operatorname{span} N_n\)，即 \(E_n\) 为由 \(N_n\) 张成的有限维线性子空间。\\
(2) 定义 \(T(C) \rightarrow \operatorname{co}(N_n)\) 的映射 \(I_n\) 如下：
\[
I_n(y) = \sum_{i=1}^{r_n} y_i \lambda_i(y) \quad (\forall y \in T(C)),
\]
其中
\[
\begin{gathered}
\lambda_i(y) = \frac{m_i(y)}{\sum_{i=1}^{r_n} m_i(y)}, \\
m_i(y) =
\begin{cases}
1 - n \|y - y_i\|, & y \in B\left(y_i, \frac{1}{n}\right), \\
0, & y \notin B\left(y_i, \frac{1}{n}\right).
\end{cases}
\end{gathered}
\]
因为 \(m_i(y) \geqslant 0\)，并且 \(\forall y \in T(C), \exists i_0 \, (1 \leqslant i_0 \leqslant r_n)\)，使得
\[
y \in B\left(y_{i_0}, \frac{1}{n}\right) \implies m_{i_0}(y) > 0,
\]
所以 \(\sum_{i=1}^{r_n} m_i(y) > 0 \, (\forall y \in T(C))\)。因此，\(\lambda_i(y) \, (1 \leqslant i \leqslant r_n)\) 有定义并满足
\[
\lambda_i(y) \geqslant 0 \quad (1 \leqslant i \leqslant r_n), \quad \sum_{i=1}^{r_n} \lambda_i(y) = 1.
\]
于是 \(I_n\) 在 \(T(C)\) 上有定义，并且公式表明 \(I_n(y)\) 是 \(N_n\) 元素的凸组合，从而 \(I_n(y) \in \operatorname{co}(N_n)\)。此外，有
\[
\begin{aligned}
\|I_n y - y\| &= \left\|\sum_{i=1}^{r_n} y_i \lambda_i(y) - \sum_{i=1}^{r_n} y \lambda_i(y)\right\| \\
&= \left\|\sum_{i=1}^{r_n} (y_i - y) \lambda_i(y)\right\| \\
&\leqslant \sum_{i=1}^{r_n} \|y_i - y\| \lambda_i(y) \\
&< \frac{1}{n}.
\end{aligned}
\]
(3) 注意到 \(T: C \rightarrow C, N_n \subset T(C)\)，而 \(C\) 是凸的，所以 \(\operatorname{co}(N_n) \subset C\)。令 \(T_n = I_n \circ T\)，那么 \(T_n: \operatorname{co}(N_n) \rightarrow \operatorname{co}(N_n)\)。又注意到 \(\operatorname{co}(N_n)\) 是 \(E_n\) 中的一个有界闭凸子集，应用推论 1.5，存在 \(x_n \in \operatorname{co}(N_n) \subset C\)，使得
\[
T_n x_n = x_n.
\]
由于 \(T(C)\) 是列紧集而 \(C\) 是闭集，所以存在子列 \(n_k\) 及 \(x \in C\)，使得
\[
T x_{n_k} \rightarrow x \quad (k \rightarrow \infty).
\]
结合以下两式：
\[
\begin{aligned}
\|x_n - x\| &= \|T_n x_n - x\| \\
&= \|I_n T x_n - T x_n + T x_n - x\| \\
&\leqslant \|I_n T x_n - T x_n\| + \|T x_n - x\| \\
&< \frac{1}{n} + \|T x_n - x\| \quad (\forall n \in \mathbb{N}),
\end{aligned}
\]
得 \(x_{n_k} \rightarrow x (k \rightarrow \infty)\)。利用 \(T\) 的连续性以及上述式子，得
\[
T x = x.
\]
\end{proof}
\begin{note}
凸组合的运算可以写为一个从参数空间到欧几里得空间的连续映射:
\begin{equation*}
f: \Delta \rightarrow \mathbb{R} ^n, \quad f\left(\lambda_1, \lambda_2, \ldots, \lambda_k\right)=\sum_{i=1}^k \lambda_i x_i
\end{equation*}
其中 $\Delta$ 表示一个标准 $k-1$ 维单纯形:
\begin{equation*}
\Delta=\left\{\left(\lambda_1, \lambda_2, \ldots, \lambda_k\right) \in \mathbb{R} ^k \mid \lambda_i \geq 0, \sum_{i=1}^k \lambda_i=1\right\} .
\end{equation*}
参数空间 $\Delta$ 是紧的，由于 $f$ 是连续的。因此， $\Delta$ 的像 $f(\Delta)=\operatorname{co}(N)$ 是一个紧集。
\end{note}
\begin{definition}
设 \(\mathscr{X}\) 是 \(B^*\) 空间，\(E\) 是 \(\mathscr{X}\) 的一个子集，称映射 \(T: E \rightarrow \mathscr{X}\) 是紧的，如果它是连续的并且把 \(E\) 中的任意有界集映为 \(\mathscr{X}\) 中的列紧集。
\end{definition}

\begin{corollary}
设 \(C\) 为 \(B^*\) 空间 \(\mathscr{X}\) 中的一个有界闭凸子集，\(T: C \rightarrow C\) 是紧的，则 \(T\) 在 \(C\) 上必有不动点。
\end{corollary}
\begin{example}
设 $k(x, y)$ 是 $[0,1] \times [0,1]$ 上的正值连续函数，
\[
(T u)(x) = \int_0^1 k(x, y) u(y) \, \dif y \quad (\forall u \in C[0,1]) .
\]
证明：存在 $\lambda > 0$ 及非负但不恒为零的连续函数 $u$，满足 $T u = \lambda u$。
\end{example}
\begin{proof}
设 $0 < m \leqslant k(x, y) \leqslant M$ 。在 $C[0,1]$ 上考查子集
\[
C=\left\{ u \in C[0,1] \mid \int_0^1 u(x) \, \dif x = 1, \, u(x) \geqslant 0 \right\} ,
\]
及映射 $S$ ：
\[
(S u)(x) = \frac{\int_0^1 k(x, y) u(y) \, \dif y}{\int_0^1 \int_0^1 k(t, y) u(y) \, \dif y \, \dif t}, \quad \forall u \in C .
\]
显然 $C$ 是闭凸集，且 $S: C \to C$ 连续。  进一步，为了证明 $S(C)$ 列紧，需要证明：\\
（1） $S(C)$ 一致有界。事实上，
\[
|S(u)(x)| \leqslant \frac{M \int_0^1 u(y) \, \dif y}{m \int_0^1 u(y) \, \dif y} = \frac{M}{m} .
\]
（2） $S(C)$ 等度连续。事实上，
\[
\int_0^1 \int_0^1 k(x, y) u(y) \, \dif y \, \dif x \geqslant \int_0^1 \int_0^1 m \cdot u(y) \,\dif y \, \dif x = m > 0 .
\]
又因为 $k(x, y)$ 是 $[0,1] \times [0,1]$ 上的连续函数，所以在 $[0,1] \times [0,1]$ 上一致连续。  
故对 $\forall \varepsilon > 0, \exists \delta > 0$，使得 $\left|k(x, y) - k(x', y)\right| < m \varepsilon$，当 $\left|x - x'\right| < \delta$ 时。  
因此，对 $\forall \varepsilon > 0, \exists \delta > 0$，使得 $\forall x, x' \in [0,1]$，只要 $\left|x - x'\right| < \delta$，便有：
\[
\begin{aligned}
|S u(x) - S u(x')| &= \left| \frac{\int_0^1 k(x, y) u(y) \, \dif y}{\int_0^1 \int_0^1 k(t, y) u(y) \, \dif y \, \dif t} - \frac{\int_0^1 k(x', y) u(y) \, \dif y}{\int_0^1 \int_0^1 k(t, y) u(y) \, \dif y \, \dif t} \right| \\
&\leqslant \frac{1}{m} \int_0^1 \left| k(x, y) - k(x', y) \right| u(y) \, \dif y < \varepsilon .
\end{aligned}
\]
根据本章Arzela-Ascoli 定理知 $S(C)$ 列紧。  从而由本节推论 1.6，$\exists u \in C$，使得 $S u = u$，即：
\[
T u = \lambda u, \quad \lambda = \int_0^1 \int_0^1 k(t, y) u(y) \, \dif y \, \dif t > 0 .
\]
\end{proof}
\begin{example}
求证: 在 $B$ 空间中, 列紧集$M$的凸包是列紧集.
\end{example}
\begin{proof}
首先证明有限点集$T=\{z_i|1\leqslant i\leqslant n\}$的凸包是自列紧集，设$\{\sum\limits_{i=1}^n \lambda_i^{(k)}z_i \}_{k=1}^\infty\subset \operatorname{co}(T)$，其中$\sum\limits_{i=1}^n\lambda_{i}^{(k)}=1,\lambda_i^{(k)}\geqslant 0,\forall k\geqslant 1$，由维尔斯特拉斯定理可知$\{\lambda_{1}^{(k)}\}_{k=1}^\infty$存在收敛子列$\{\lambda_{1}^{(k_j)}\}_{j=1}^\infty$，同样操作下去可以得到$\{\lambda_{2}^{(k_j)}\}_{j=1}^\infty$存在收敛子列$\{\lambda_{2}^{(k_{j_l})}\}_{l=1}^\infty$，以此类推得到$\{(\lambda_1^{(k)},\lambda_2^{(k)},\cdots,\lambda_n^{(k)})\}_{k=1}^\infty$存在收敛子列$\{(\lambda_1^{(k_m)},\lambda_2^{(k_m)},\cdots,\lambda_n^{(k_m)})\}_{m=1}^\infty$，于是$\{\sum\limits_{i=1}^n \lambda_i^{(k)}z_i \}_{k=1}^\infty$的子列$\{\sum\limits_{i=1}^n \lambda_i^{(k_m)}z_i \}_{m=1}^\infty$
    收敛，同时收敛到的点也是相加为$1$，那么$\operatorname{co}(T)$是自列紧集。\\
    $\forall \varepsilon>0$，存在$M$的$\frac{\varepsilon}{2}$网$T=\{z_i|1\leqslant i\leqslant n\}$，而$\operatorname{co}(T)$列紧，所以也存在$\frac{\varepsilon}{2}$网$S=\{y_i|1\leqslant i\leqslant m\}\subset \operatorname{co}(M)$，对于$\forall a=\sum\limits_{i=1}^k \lambda_ia_i\in \operatorname{co}(M)$，对于每一个$x_i$
 存在$b_i\in T$使得$\|a_i-b_i\|<\frac{\varepsilon}{2}$，令  $ b=\sum\limits_{i=1}^k \lambda_ib_i\in \operatorname{co}(T)$，存在$c\in S$使得$\|b-c\|<\frac{\varepsilon}{2}$，而$\|a-b\|\leqslant \sum\limits_{i=1}^k \lambda_i\|a_i-b_i\|<\frac{\varepsilon}{2}$，所以
 \[
 \|a-c\|\leqslant\|a-b\|+\|b-c\|<\varepsilon
 \]
 那么$S$是$\operatorname{co}(M)$的$\varepsilon$网，所以$\operatorname{co}(M)$是完全有界集，进而是列紧集。

 
\end{proof}

\section{内积空间}
\begin{definition}
线性空间 $\mathscr{X}$ 上的一个二元函数 $a(\cdot, \cdot): \mathscr{X} \times \mathscr{X} \to \mathbb{K}$，称为共轭双线性函数（sesquilinear），如果：
\begin{enumerate}
    \item $a\left(x, \alpha_1 y_1 + \alpha_2 y_2\right) = \overline{\alpha}_1 a\left(x, y_1\right) + \overline{\alpha}_2 a\left(x, y_2\right)$,
    \item $a\left(\alpha_1 x_1 + \alpha_2 x_2, y\right) = \alpha_1 a\left(x_1, y\right) + \alpha_2 a\left(x_2, y\right)$,
\end{enumerate}
其中 $\forall x, y, x_1, x_2, y_1, y_2 \in \mathscr{X}$, $\forall \alpha_1, \alpha_2 \in \mathbb{K}$。  
我们还称由
\[
q(x) \triangleq a(x, x) \quad (\forall x \in \mathscr{X})
\]
定义的函数为 $\mathscr{X}$ 上由 $a$ 诱导的二次型。
\end{definition}

\begin{proposition}
设 $a$ 是 $\mathscr{X}$ 上的共轭双线性函数，$q$ 是由 $a$ 诱导的二次型，那么：
\[
q(x) \in \mathbb{R} \ (\forall x \in \mathscr{X}) \Longleftrightarrow a(x, y) = \overline{a(y, x)} \quad (\forall x, y \in \mathscr{X}).
\]
\end{proposition}
\begin{proof}
"$\Longrightarrow$"：从
\[
q(x+y) = \overline{q(x+y)} \quad (\forall x, y \in \mathscr{X}),
\]
容易推出
\[
a(x, y) + a(y, x) = \overline{a(x, y)} + \overline{a(y, x)}.
\]
将 $y$ 替换为 $iy$，即得
\[
-a(x, y) + a(y, x) = \overline{a(x, y)} - \overline{a(y, x)}.
\]
两式相减得 $a(x, y) = \overline{a(y, x)}$。
\end{proof}
\begin{definition}
线性空间 $\mathscr{X}$ 上的一个共轭双线性函数
\[
(\cdot, \cdot): \mathscr{X} \times \mathscr{X} \to \mathbb{K}
\]
称为内积，如果它满足：
\begin{enumerate}
    \item $(x, y) = \overline{(y, x)} \quad (\forall x, y \in \mathscr{X})$ （共轭对称性）；
    \item $(x, x) \geqslant 0 \quad (\forall x \in \mathscr{X})$，且 $(x, x) = 0 \Longleftrightarrow x = \theta$ （正定性）。
\end{enumerate}
具有内积的线性空间称为内积空间，记作 $(\mathscr{X}, (\cdot, \cdot))$。
\end{definition}
\begin{note}
若在 (2) 中仅保留非负定条件 $(x, x) \geqslant 0 \ (\forall x \in \mathscr{X})$，则称 $(\cdot, \cdot)$ 为一个半内积，对应的空间称为半内积空间。
\end{note}

\begin{example}
 $\mathbb{R} ^n, \mathbb{C} ^n$ 都是内积空间, 它们的内积分别定义为
\begin{equation*}
\begin{aligned}
& (x, y)=\sum_{i=1}^n x_i y_i \quad\left(\forall x, y \in \mathbb{R}^n\right), \\
& (x, y)=\sum_{i=1}^n x_i \bar{y}_i \quad\left(\forall x, y \in \mathbb{C} ^n\right),
\end{aligned}
\end{equation*}
其中 $x=\left(x_1, x_2, \cdots, x_n\right), y=\left(y_1, y_2, \cdots, y_n\right)$.
\end{example}
\begin{example}
 $l^2$ 空间是内积空间，规定内积
\begin{equation*}
(x, y)=\sum_{i=1}^{\infty} x_i \bar{y}_i,
\end{equation*}
其中 $x=\left(x_1, x_2, \cdots, x_i, \cdots\right), y=\left(y_1, y_2, \cdots, y_i, \cdots\right) \in l^2$.
\end{example}
\begin{example}
 $L^2(\Omega, \mu)$是内积空间, 规定内积
\begin{equation*}
(u, v)=\int_{\Omega} u(x) \cdot \overline{v(x)} \dif \mu,
\end{equation*}
其中 $u, v \in L^2(\Omega, \mu)$.
\end{example}
\begin{example} $C^k(\bar{\Omega})$. 设 $\Omega$ 是 $\mathbb{R} ^n$ 中的一个有界连通开区域, $k \in N$ ，用 $C^k(\bar{\Omega})$ 表示在 $\bar{\Omega}$ 上具有直到 $k$ 阶连续偏导数的函数 $u(x)=u\left(x_1, x_2, \cdots, x_n\right)$ 的全体, 加法与数乘按自然法则定义, 再规定范数为
\begin{equation*}
\|u\|=\max _{|\alpha| \leqslant k} \max _{x \in \bar{\Omega}}\left|\partial^\alpha u(x)\right|,
\end{equation*}
其中 $\alpha=\left(\alpha_1, \alpha_2, \cdots, \alpha_n\right),|\alpha|=\alpha_1+\alpha_2+\cdots+\alpha_n$, 以及
\begin{equation*}
\partial^\alpha u(x)=\frac{\partial^{|\alpha|}}{\partial x_1^{\alpha_1} \partial x_2^{\alpha_2} \cdots \partial x_n^{\alpha_n}} u(x) .
\end{equation*}
那么$C^k(\bar{\Omega})$是$B$空间。
\end{example}
\begin{proof}
容易验证上式定义的 $\|\cdot\|$ 是一个范数，从而 $C^k(\bar{\Omega})$ 是一个赋范线性空间. 再证它是完备的，从而是 Banach 空间. 事实上，若 $\left\{u_n\right\}$ 是一个基本列, 则必存在连续函数 $v_\alpha$, 使得
\begin{equation*}
\partial^\alpha u_n(x) \rightrightarrows v_\alpha(x) \quad(n \rightarrow \infty,|\alpha| \leqslant k \text {, 对 } x \in \bar{\Omega} \text { 一致 }) .
\end{equation*}
我们只要再证明 $v_\alpha=\partial^\alpha v_\theta$ 就够了. 先证
\begin{equation*}
v_{(1,0, \cdots, 0)}=\frac{\partial}{\partial x_1} v_{(0,0, \cdots, 0)}
\end{equation*}
因为
\begin{equation*}
\begin{aligned}
& \frac{\partial}{\partial x_1} u_n(x) \rightrightarrows v_{(1,0, \cdots, 0)}(x) \quad(n \rightarrow \infty, \text { 对 } x \in \bar{\Omega} \text { 一致 }), \\
& u_n(x) \rightrightarrows v_{(0,0, \cdots, 0)}(x) \quad(n \rightarrow \infty, \text { 对 } x \in \bar{\Omega} \text { 一致 }),
\end{aligned}
\end{equation*}
以及
\begin{equation*}
u_n(x)=\int_{x_1^0}^{x_1} \frac{\partial}{\partial \xi} u_n\left(\xi, x_2, \cdots, x_n\right) \dif \xi+u_n\left(x_1^0, x_2, \cdots, x_n\right),
\end{equation*}
所以有
\begin{equation*}
\begin{aligned}
& v_{(0,0, \cdots, 0)}(x) \\
& \quad=\int_{x_1^0}^{x_1} v_{(1,0, \cdots, 0)}\left(\xi, x_2, \cdots, x_n\right) \dif \xi+v_{(0,0, \cdots, 0)}\left(x_1^0, x_2, \cdots, x_n\right)
\end{aligned}
\end{equation*}
即得 应证之式. 同理有
\begin{equation*}
v_{(\underbrace{0,0, \cdots, 0,1}_j,0, \cdots, 0)}=\frac{\partial}{\partial x_j} v_{(0,0, \cdots, 0)} \quad(j=2,3, \cdots, n) .
\end{equation*}
其余对 $|\alpha|$ 用数学归纳法类推.
\end{proof}



\begin{example}
在空间 $C^k(\bar{\Omega})$ 中, 规定内积
\begin{equation*}
(u, v)=\sum_{|\alpha| \leqslant k} \int_{\Omega} \partial^\alpha u(x) \cdot \overline{\partial^\alpha v(x)} \overline{\dif x} \quad\left(\forall u, v \in C^k(\bar{\Omega})\right) .
\end{equation*}
那么 $\left(C^k(\bar{\Omega}),(\cdot, \cdot)\right)$ 是一个内积空间.
\end{example}

\begin{proposition}[Cauchy-Schwarz 不等式]
设 $(\mathscr{X}, (\cdot, \cdot))$ 是内积空间。若令
\[
\|x\| = (x, x)^{\frac{1}{2}} \quad (\forall x \in \mathscr{X}),
\]
则有
\[
|(x, y)| \leqslant \|x\| \cdot \|y\| \quad (\forall x, y \in \mathscr{X}),
\]
且等号当且仅当 $x$ 与 $y$ 线性相关时成立。
\end{proposition}

\begin{proposition}
设 $a$ 是线性空间 $\mathscr{X}$ 上的共轭双线性函数，$q(x)$ 是由 $a$ 诱导的二次型。如果
\[
q(x) \geqslant 0 \quad (\forall x \in \mathscr{X}) \quad \text{且} \quad q(x) = 0 \Longleftrightarrow x = \theta,
\]
那么
\[
|a(x, y)| \leqslant [q(x) q(y)]^{\frac{1}{2}} \quad (\forall x, y \in \mathscr{X}),
\]
且等号当且仅当 $x$ 与 $y$ 线性相关时成立。
\end{proposition}
\begin{proof}
不妨设 $y \neq \theta$。对 $\forall \lambda \in \mathbb{K}$，考察
\[
q(x + \lambda y) = q(x) + \overline{\lambda} a(x, y) + \lambda a(y, x) + |\lambda|^2 q(y) \geqslant 0.
\]
取 $\lambda = -a(x, y) / q(y)$，由 $a(x, y) = \overline{a(y, x)}$，有
\[
q(x) - \frac{2|a(x, y)|^2}{q(y)} + \frac{|a(x, y)|^2}{q(y)} \geqslant 0.
\]
由此得证上式。且当 $x = -\lambda y (\lambda \in \mathbb{K})$ 时，等号成立。反之，若等号成立，则 $x = -\lambda y (\lambda \in \mathbb{K})$。
\end{proof}
\begin{proposition}
内积空间 $(\mathscr{X}, (\cdot, \cdot))$ 按$\|x\|=(x,x)^{\frac{1}{2}}$定义范数，是一个 $B^*$ 空间。
\end{proposition}

\begin{proof}
事实上，范数定义中的条件 (1) 和 (3) 显然成立，接下来验证 (2)。
\begin{equation*}
\begin{aligned}
\|x+y\|^2 &= (x+y, x+y) \\
&= (x, x) + (x, y) + (y, x) + (y, y) \\
&\leqslant \|x\|^2 + 2\|x\|\|y\| + \|y\|^2 \\
&= (\|x\| + \|y\|)^2 \quad (\forall x, y \in \mathscr{X}).
\end{aligned}
\end{equation*}
\end{proof}

\begin{proposition}
在内积空间 $(\mathscr{X}, (\cdot, \cdot))$ 中，内积 $(x, y)$ 是 $\mathscr{X} \times \mathscr{X}$ 上关于范数 $\|\cdot\|$ 的连续函数。
\end{proposition}

\begin{proof}
设 $x_n \rightarrow x$，$y_n \rightarrow y$。那么 $\left\|x_n\right\|$ 和 $\left\|y_n\right\|$ 有界，用 $M$ 表示它们的一个上界，则有：
\begin{equation*}
\begin{aligned}
& \left|\left(x_n, y_n\right) - (x, y)\right| \\
&\quad \leqslant \left|\left(x_n, y_n\right) - \left(x, y_n\right)\right| + \left|\left(x, y_n\right) - (x, y)\right| \\
&\quad \leqslant \|x_n - x\| \cdot \|y_n\| + \|x\| \cdot \|y_n - y\| \\
&\quad \leqslant M\|x_n - x\| + \|x\| \|y_n - y\| \rightarrow 0 \quad (n \rightarrow \infty).
\end{aligned}
\end{equation*}
\end{proof}

\begin{proposition}
内积空间 $(\mathscr{X}, (\cdot, \cdot))$ 是严格凸的 $B^*$ 空间。
\end{proposition}

\begin{proof}
对于任意 $0 < \lambda < 1$，根据Cauchy-Schwarz不等式有：
\begin{equation*}
\begin{aligned}
\|\lambda x + (1-\lambda) y\|^2 &= \lambda^2 \|x\|^2 + 2 \lambda (1-\lambda) \operatorname{Re}(x, y) + (1-\lambda)^2 \|y\|^2 \\
&< [\lambda + (1-\lambda)]^2 = 1 \quad (\text{当 } \|x\| = \|y\| = 1, x \neq y).
\end{aligned}
\end{equation*}

\end{proof}

\begin{proposition}
在 $B^*$ 空间 $(\mathscr{X}, \|\cdot\|)$ 中，为了在 $\mathscr{X}$ 上引入一个内积 $(\cdot, \cdot)$ 使其适合 $\|x\|=(x,x)^{\frac{1}{2}}$，必须且仅须范数 $\|\cdot\|$ 满足如下平行四边形等式：
\begin{equation*}
\|x + y\|^2 + \|x - y\|^2 = 2\left(\|x\|^2 + \|y\|^2\right) \quad (\forall x, y \in \mathscr{X}).
\end{equation*}
\end{proposition}

\begin{proof}
必要性可以通过直接计算得到。为了证明充分性，令
\begin{equation*}
(x, y) = \begin{cases}
\frac{1}{4}\left(\|x + y\|^2 - \|x - y\|^2\right), & \text{当 } \mathbb{K} = \mathbb{R}, \\
\frac{1}{4}\left(\|x + y\|^2 - \|x - y\|^2 + i\|x + i y\|^2 - i\|x - i y\|^2\right), & \text{当 } \mathbb{K} = \mathbb{C}.
\end{cases}
\end{equation*}
容易验证它是一个满足 上式的内积。
\end{proof}
\begin{definition}
完备的内积空间称为 Hilbert 空间。
\end{definition}
\begin{lemma}[Poincaré 不等式]设 $C_0^m(\Omega)$ 表示有界开区域 $\Omega \subset \mathbb{R} ^n$ 上一切 $m$ 次连续可微，并在边界 $\partial \Omega$ 的某邻域内为 0的函数集合，即
\begin{equation*}
C_0^m(\Omega)=\left\{u \in C^m(\bar{\Omega}) \mid u(x)=0 \text {, 当 } x \in \partial \Omega \text { 的某邻域 }\right\} .
\end{equation*}
那么 $\forall u \in C_0^m(\Omega)$ 有
\begin{equation*}
\sum_{|\alpha|<m} \int_{\Omega}\left|\partial^\alpha u(x)\right|^2 \dif x \leqslant C \sum_{|\alpha|=m} \int_{\Omega}\left|\partial^\alpha u(x)\right|^2 \dif x,
\end{equation*}
其中 $C$ 是仅依赖于区域 $\Omega$ 及 $m$ 的常数.
\end{lemma}
\begin{proof}
因为 $\Omega$ 是有界的，我们可以把 $\Omega$ 放在某个边长为 $a$ 的立方体 $\Omega_1$ 内, 适当选择坐标系, 使得
\begin{equation*}
\Omega_1=\left\{\left(x_1, x_2, \cdots, x_n\right) \in \mathbb{R} ^n \mid 0 \leqslant x_i \leqslant a(i=1,2, \cdots, n)\right\} .
\end{equation*}
在 $\Omega_1 \backslash \Omega$ 上补充定义 $u=0$ ，经补充定义后， $u(x)$ 在 $\Omega_1$ 上 $m$ 次连续可微，而且在边界上等于 $0 . \forall x \in \Omega_1$ ，
\begin{equation*}
u(x)=\int_0^{x_1} \frac{\partial u}{\partial t}\left(t, x_2, \cdots, x_n\right) \dif t .
\end{equation*}
再利用 Cauchy-Schwarz 不等式，我们有
\begin{equation*}
|u(x)|^2 \leqslant a \int_0^a\left|\frac{\partial u}{\partial x_1}\right|^2 \dif x_1 .
\end{equation*}
在 $\Omega_1$ 上积分不等式, 我们得\begin{equation*}
\begin{aligned}
\int_{\Omega}|u(x)|^2 \dif x & \leqslant a^2 \int_{\Omega}\left|\frac{\partial u}{\partial x_1}\right|^2 \dif x \\
& \leqslant a^2 \int_{\Omega}|\operatorname{grad} u(x)|^2 \dif x .
\end{aligned}
\end{equation*}
然后逐次应用不等式于 $\partial^\alpha u(x)(|\alpha|<m)$ ，即得不等式 .
\end{proof}
\begin{definition}[正交与正交补]
内积空间 $\mathscr{X}$ 上的两个元素 $x$ 与 $y$ 称为是正交的，是指
\begin{equation*}
(x, y) = 0
\end{equation*}
记作 $x \perp y$。又设 $M$ 是 $\mathscr{X}$ 的一个非空子集，若对 $\forall y \in M$ 都有 $x \perp y$，则称 $x$ 与 $M$ 正交，记作 $x \perp M$。此外，集合
\begin{equation*}
\{x \in \mathscr{X} \mid x \perp M\}
\end{equation*}
称为 $M$ 的正交补，记作 $M^{\perp}$。
\end{definition}

\begin{proposition}
设 $\mathscr{X}$ 是内积空间，$M$ 是 $\mathscr{X}$ 的一个非空子集。
\begin{enumerate}
    \item 若 $x \perp y_i \ (i=1,2)$，则
    \begin{equation*}
    x \perp \lambda_1 y_1 + \lambda_2 y_2 \quad (\forall \lambda_1, \lambda_2 \in \mathbb{K}).
    \end{equation*}
    \item 若 $x = y + z$，且 $y \perp z$，则
    \begin{equation*}
    \|x\|^2 = \|y\|^2 + \|z\|^2.
    \end{equation*}
    \item 若 $x \perp y_n \ (n \in \mathbb{N})$，且 $y_n \to y$，则 $x \perp y$。
    \item 若 $x \perp M$，则 $x \perp \operatorname{span}(M)$。
    \item $M^{\perp}$ 是 $\mathscr{X}$ 的一个闭线性子空间。
\end{enumerate}
\end{proposition}

\begin{definition}[正交集与完备正交集]
设 $\mathscr{X}$ 是一个内积空间，集合 $S = \{e_\alpha \mid \alpha \in A\}$ 是 $\mathscr{X}$ 的一个子集。称 $S$ 为正交集，是指
\begin{equation*}
e_\alpha \perp e_\beta \quad (\text{当 } \alpha \neq \beta, \ \forall \alpha, \beta \in A).
\end{equation*}
如果还有 $\|e_\alpha\| = 1 \ (\forall \alpha \in A)$，则称 $S$ 为正交规范集。\\
若 $\mathscr{X}$ 中不存在非零元与 $S$ 正交，即 $S^{\perp} = \{\theta\}$，则称 $S$ 为完备的。
\end{definition}

\begin{lemma}[Zorn 引理]
设 $\mathscr{X}$ 是一个半序集。如果它的每一个全序子集有一个上界，那么 $\mathscr{X}$ 有一个极大元。
\end{lemma}

\begin{proposition}
非 $\{\theta\}$ 内积空间 $\mathscr{X}$ 中必存在完备正交集。
\end{proposition}

\begin{proof}
因为 $\mathscr{X} \neq \{\theta\}$，所以 $\mathscr{X}$ 中的正交集依包含关系构成一个半序集类，并且每个全序子集类有一个上界，就是这些集之并集。依 Zorn 引理，这个半序集类有极大元。记此极大元为 $S$，我们证明 $S$ 是完备正交集。

若不然，则必存在 $x_0 \in S^{\perp}, x_0 \neq \theta$，令 $S_1 = \{x_0\} \cup S$，则 $S_1$ 仍为正交集，但 $S \nsubseteq S_1$，这与 $S$ 的极大性矛盾。
\end{proof}

\begin{definition}[正交规范基]
内积空间 $\mathscr{X}$ 中的正交规范集 $S = \{e_\alpha \mid \alpha \in A\}$ 称为一个基（或封闭的），是指 $\forall x \in \mathscr{X}$，有以下表示：
\begin{equation*}
x = \sum_{\alpha \in A} (x, e_\alpha) e_\alpha,
\end{equation*}
其中 $\{(x, e_\alpha) \mid \alpha \in A\}$ 称为 $x$ 关于基 $S$ 的 Fourier 系数。
\end{definition}

\begin{theorem}[Bessel 不等式]\label{thm:bessel}
设 $\mathscr{X}$ 是一个内积空间。如果 $S = \{e_\alpha \mid \alpha \in A\}$ 是 $\mathscr{X}$ 中的正交规范集，那么 $\forall x \in \mathscr{X}$，有
\begin{equation*}
\sum_{\alpha \in A}\left|\left(x, e_\alpha\right)\right|^2 \leqslant \|x\|^2. 
\end{equation*}
\end{theorem}

\begin{proof}
首先对 $A$ 的任意有限子集（设为 $1, 2, \dots, n$）证明
\begin{equation}
\sum_{i=1}^n\left|\left(x, e_i\right)\right|^2 \leqslant \|x\|^2. 
\end{equation}
因为
\begin{equation*}
\begin{aligned}
0 &\leqslant \left\|x - \sum_{i=1}^n\left(x, e_i\right)e_i\right\|^2 \\
  &= \left(x - \sum_{i=1}^n\left(x, e_i\right)e_i, x - \sum_{j=1}^n\left(x, e_j\right)e_j\right) \\
  &= \|x\|^2 - \sum_{i=1}^n\left|\left(x, e_i\right)\right|^2.
\end{aligned}
\end{equation*}
因此Bessel不等式成立。\\
由此可见，对 $\forall n \in \mathbb{N}$，使得 $\left|\left(x, e_\alpha\right)\right| > \frac{1}{n}$ 的 $\alpha \in A$ 至多有限多个，从而 $\left(x, e_\alpha\right) \neq 0$ 的 $\alpha \in A$ 至多是可数的。于是，Bessel不等式 的左端实际上是一个至多可数项求和的级数。\\
再由有限情况的结论，我们有
\begin{equation*}
\sum_{\alpha \in A_f}\left|\left(x, e_\alpha\right)\right|^2 \leqslant \|x\|^2,
\end{equation*}
其中 $A_f$ 表示 $A$ 的任意有限子集。由此得出Bessel不等式。
\end{proof}

\begin{corollary}
假设 $\mathscr{X}$ 是 Hilbert 空间，且 $S = \{e_\alpha \mid \alpha \in A\}$ 是 $\mathscr{X}$ 中的正交规范集。那么对 $\forall x \in \mathscr{X}$，有
\begin{equation*}
\sum_{\alpha \in A}\left(x, e_\alpha\right)e_\alpha \in \mathscr{X},
\end{equation*}
且
\begin{equation}
\left\|x - \sum_{\alpha \in A}\left(x, e_\alpha\right)e_\alpha\right\|^2 = \|x\|^2 - \sum_{\alpha \in A}\left|\left(x, e_\alpha\right)\right|^2. 
\end{equation}
\end{corollary}

\begin{proof}
不妨设使得 $\left(x, e_\alpha\right) \neq 0$ 的可数多个 $\alpha \in A$ 为 $1, 2, \dots, n, \dots$，那么
\begin{equation*}
\sum_{\alpha \in A}\left(x, e_\alpha\right)e_\alpha = \sum_{n=1}^\infty\left(x, e_n\right)e_n.
\end{equation*}
由 Bessel 不等式可知 $\sum\limits_{n=1}^\infty\left|\left(x, e_n\right)\right|^2$ 收敛，因此有
\begin{equation*}
\left\|\sum_{n=m}^{m+p}\left(x, e_n\right)e_n\right\|^2 = \sum_{n=m}^{m+p}\left|\left(x, e_n\right)\right|^2 \to 0 \quad (m \to \infty, \ \forall p \in \mathbb{N}).
\end{equation*}
于是，序列 $\{x_m = \sum\limits_{n=1}^m\left(x, e_n\right)e_n\}$ 是基本列，从而
\begin{equation*}
\sum_{\alpha \in A}\left(x, e_\alpha\right)e_\alpha = \sum_{n=1}^\infty\left(x, e_n\right)e_n = \lim_{m \to \infty}x_m \in \mathscr{X}.
\end{equation*}
又因为
\begin{equation*}
x - \sum_{n=1}^\infty\left(x, e_n\right)e_n \perp \sum_{n=1}^\infty\left(x, e_n\right)e_n,
\end{equation*}
所以
\begin{equation*}
\left\|x - \sum_{n=1}^\infty\left(x, e_n\right)e_n\right\|^2 = \|x\|^2 - \sum_{n=1}^\infty\left|\left(x, e_n\right)\right|^2.
\end{equation*}
\end{proof}
\begin{theorem}
设 $\mathscr{X}$ 是一个 Hilbert 空间，若 $S = \{e_\alpha \mid \alpha \in A\}$ 是 $\mathscr{X}$ 中的正交规范集，则以下三条等价：
\begin{enumerate}
    \item[(1)] $S$ 是封闭的；
    \item[(2)] $S$ 是完备的；
    \item[(3)] Parseval 等式
    \begin{equation}
    \|x\|^2 = \sum_{\alpha \in A}\left|\left(x, e_\alpha\right)\right|^2, \quad \forall x \in \mathscr{X},
    \end{equation}
    成立。
\end{enumerate}
\end{theorem}

\begin{proof}
\noindent\textbf{(1) $\Longrightarrow$ (2):}  
若 $S$ 不完备，则 $\exists x \in \mathscr{X} \setminus \{\theta\}$，使得
\begin{equation}
\left(x, e_\alpha\right) = 0, \quad \forall \alpha \in A.
\end{equation}
但由于 $S$ 是封闭的，根据定义有
\begin{equation*}
x = \sum_{\alpha \in A}\left(x, e_\alpha\right)e_\alpha = \theta,
\end{equation*}
这与 $x \neq \theta$ 矛盾，故 $S$ 是完备的。\\
\noindent\textbf{(2) $\Longrightarrow$ (3):}  \\
若 $\exists x \in \mathscr{X}$ 使 Parseval 等式  不成立，则由推论1.7知
\begin{equation}
\left\|x - \sum_{\alpha \in A}\left(x, e_\alpha\right)e_\alpha\right\|^2 = \|x\|^2 - \sum_{\alpha \in A}\left|\left(x, e_\alpha\right)\right|^2 > 0. 
\end{equation}
定义
\begin{equation*}
y \triangleq x - \sum_{\alpha \in A}\left(x, e_\alpha\right)e_\alpha \neq \theta,
\end{equation*}
显然 $y \in S^\perp$，但这与 $S$ 的完备性矛盾，因此 Parseval 等式成立。\\
\noindent\textbf{(3) $\Longrightarrow$ (1):} \\
联合 Parseval 等式和式 推论1.7，得到
\begin{equation}
\left\|x - \sum_{\alpha \in A}\left(x, e_\alpha\right)e_\alpha\right\|^2 = \|x\|^2 - \sum_{\alpha \in A}\left|\left(x, e_\alpha\right)\right|^2 = 0. 
\end{equation}
因此有
\begin{equation*}
x = \sum_{\alpha \in A}\left(x, e_\alpha\right)e_\alpha.
\end{equation*}
这表明 $S$ 是封闭的。
\end{proof}
\begin{example}
在 $L^2[0,2\pi]$ 上，
\begin{equation*}
e_n(t) = \frac{1}{\sqrt{2\pi}} e^{i n t}, \quad (n = 0, \pm 1, \pm 2, \cdots),
\end{equation*}
是一组正交规范基。对于任意 $u \in L^2[0,2\pi]$，对应的 Fourier 系数为
\begin{equation*}
\left(u, e_n\right) = \frac{1}{\sqrt{2\pi}} \int_0^{2\pi} u(t) e^{-i n t} \, dt, \quad (n = 0, \pm 1, \pm 2, \cdots).
\end{equation*}
\end{example}

\begin{example}
在 $l^2$ 空间上，
\begin{equation*}
e_n = (\underbrace{0, 0, \cdots, 0, 1}_n, 0, \cdots), \quad (n = 1, 2, 3, \cdots),
\end{equation*}
是一组正交规范基。
\end{example}

\begin{example}
设 $D$ 是复平面 $\mathbb{C}$ 中的单位开圆域，$H^2(D)$ 表示在 $D$ 内满足
\begin{equation*}
\iint_D |u(z)|^2 \, \dif x \, \dif y < \infty, \quad (z = x + i y),
\end{equation*}
的解析函数全体组成的空间，内积定义为
\begin{equation*}
(u, v) = \iint_D u(z) \overline{v(z)} \, \dif x \, \dif y.
\end{equation*}
此时函数组
\begin{equation*}
\varphi_n(z) = \sqrt{\frac{n}{\pi}} z^{n-1}, \quad (n = 1, 2, 3, \cdots),
\end{equation*}
是一组正交规范基。设 $u(z) = \sum_{k=0}^\infty b_k z^k \in H^2(D)$，其对应的 Fourier 系数为
\begin{equation*}
\begin{aligned}
\left(u, \varphi_n\right) 
&= \iint_D \sum_{k=0}^\infty b_k z^k \sqrt{\frac{n}{\pi}} \bar{z}^{n-1} \, \dif x \, \dif y \\
&= \sum_{k=0}^\infty b_k \sqrt{\frac{\pi}{k+1}} \left(\varphi_{k+1}, \varphi_n\right) \\
&= b_{n-1} \sqrt{\frac{\pi}{n}}, \quad (n = 1, 2, 3, \cdots).
\end{aligned}
\end{equation*}
\end{example}
\begin{definition}
设 $\left( \mathscr{X} _1, (\cdot, \cdot)_1 \right)$ 与 $\left( \mathscr{X} _2, (\cdot, \cdot)_2 \right)$ 是两个内积空间，如果存在 $\mathscr{X} _1 \rightarrow \mathscr{X} _2$ 的一个线性同构 $T$，满足
\begin{equation*}
(Tx, Ty)_2 = (x, y)_1 \quad (\forall x, y \in \mathscr{X}_1),
\end{equation*}
则称内积空间 $\mathscr{X}_1$ 与 $\mathscr{X}_2$ 是同构的。
\end{definition}

\begin{theorem}
为了 Hilbert 空间 $\mathscr{X}$ 是可分的，必须且仅须它有至多可数的正交规范基 $S$。又若 $S$ 的元素个数 $N < \infty$，则 $\mathscr{X}$ 同构于 $\mathbb{K}^N$；若 $N = \infty$，则 $\mathscr{X}$ 同构于 $l^2$。
\end{theorem}

\begin{proof}
必要性. 设 $\{x_n\}_{n=1}^\infty$ 是 $\mathscr{X}$ 中的可数稠密子集，那么其中必存在一个线性无关的子集 $\{y_n\}_{n=1}^N$ ($N < \infty$ 或 $N = \infty$)，使得
\begin{equation*}
\operatorname{span} \{y_n\}_{n=1}^N = \operatorname{span} \{x_n\}_{n=1}^\infty.
\end{equation*}
再对 $\{y_n\}_{n=1}^N$ 应用 Gram-Schmidt 过程，便构造出一个正交规范集 $\{e_n\}_{n=1}^N$。又因为
\begin{equation*}
\overline{\operatorname{span} \{e_n\}_{n=1}^N} = \overline{\operatorname{span} \{y_n\}_{n=1}^N} = \mathscr{X},
\end{equation*}
若存在$x\neq \theta $s.t.$(x,e_i)=0(\forall i\leqslant N )$，那么$x\notin \operatorname{span}\{e_n\}_{n=1}^N$，于是$\exists x_n=\sum\limits_{i=1}^{N}\lambda_i^{(n)}e_i$s.t.$x_n\to x$，由范数收敛可以推出依内积收敛即$0=(x,x_n)\to \|x\|^2$，进而$x=0$矛盾。所以 $\{e_n\}_{n=1}^N$ 是完备的，进而是正交规范基。\\
充分性. 设 $\{e_n\}_{n=1}^N$ ($N < \infty$ 或 $N = \infty$) 是 $\mathscr{X}$ 的正交规范基，那么集合
\begin{equation*}
\left\{x = \sum_{n=1}^N a_n e_n \;\middle|\; \operatorname{Re} a_n \text{ 与 } \operatorname{Im} a_n \text{ 皆为有理数} \right\}
\end{equation*}
是 $\mathscr{X}$ 中的稠密子集，从而 $\mathscr{X}$ 是可分的。\\
对于正交规范基 $\{e_n\}_{n=1}^N$ ($N < \infty$ 或 $N = \infty$)，定义映射
\begin{equation*}
T: x \mapsto \left\{\left(x, e_n\right)\right\}_{n=1}^N \quad (\forall x \in \mathscr{X}).
\end{equation*}
根据 Parseval 等式我们有
\begin{equation*}
\|x\|^2 = \sum_{n=1}^N \left|\left(x, e_n\right)\right|^2 \quad (\forall x \in \mathscr{X}).
\end{equation*}
因此映射 $T$ 是 $\mathscr{X} \rightarrow \mathbb{K}^N$ ($N < \infty$) 或 $\mathscr{X} \rightarrow l^2$ ($N = \infty$) 的一一对应的线性同构。此外，
\begin{equation*}
\begin{aligned}
(x, y) &= \left(\sum_{i=1}^N \left(x, e_i\right) e_i, \sum_{j=1}^N \left(y, e_j\right) e_j\right) \\
&= \sum_{i=1}^N \left(x, e_i\right) \overline{\left(y, e_i\right)} \quad (\forall x, y \in \mathscr{X}).
\end{aligned}
\end{equation*}
由此可知，映射 $T$ 还保持内积。于是，当 $N < \infty$ 时，$\mathscr{X}$ 同构于 $\mathbb{K}^N$；而当 $N = \infty$ 时，$\mathscr{X}$ 同构于 $l^2$。
\end{proof}
\begin{theorem}
如果 ${C}$ 是 Hilbert 空间 $\mathscr{X}$ 中的一个闭凸子集，那么在 ${C}$ 上存在唯一元素 $x_0$ 取到最小范数。
\end{theorem}
\begin{proof}
存在性. 若 $\theta \in {C}$，则取 $x_0 = \theta$；若 $\theta \notin {C}$，定义
\begin{equation*}
d \triangleq \inf_{z \in {C}} \|z\| > 0.
\end{equation*}
根据下确界的定义，$\forall n \in \mathbb{N}$，存在 $x_n \in {C}$，使得
\begin{equation*}
d \leqslant \|x_n\| \leqslant d + \frac{1}{n} \quad (n=1,2,\cdots).
\end{equation*}
如果 $\{x_n\}$ 有极限 $x_0$，则由 ${C}$ 的闭性可知 $x_0 \in {C}$，且由上述不等式，$\|x_0\| = d$，即 $x_0$ 取到了 ${C}$ 上元素的最小范数。\\
为了证明 $\{x_n\}$ 有极限，只需验证 $\{x_n\}$ 是一个基本列。这时可以利用平行四边形等式：
\begin{equation*}
\begin{aligned}
\|x_m - x_n\|^2 &= 2\left(\|x_m\|^2 + \|x_n\|^2\right) - 4\left\|\frac{x_m + x_n}{2}\right\|^2 \\
&\leqslant 2\left[\left(d + \frac{1}{n}\right)^2 + \left(d + \frac{1}{m}\right)^2\right] - 4d^2 \rightarrow 0 \quad (\text{当 } n, m \rightarrow \infty).
\end{aligned}
\end{equation*}
因此，$\{x_n\}$ 是基本列，从而存在极限 $x_0$。\\
唯一性. 假设存在 $x_0, \widehat{x}_0 \in {C}$，满足 $\|x_0\| = \|\widehat{x}_0\| = d$。那么有
\begin{equation*}
\begin{aligned}
\|x_0 - \widehat{x}_0\|^2 &= 2\left(\|x_0\|^2 + \|\widehat{x}_0\|^2\right) - 4\left\|\frac{x_0 + \widehat{x}_0}{2}\right\|^2 \\
&\leqslant 4d^2 - 4d^2 = 0.
\end{aligned}
\end{equation*}
因此，$\|x_0 - \widehat{x}_0\| = 0$，即 $x_0 = \widehat{x}_0$，从而唯一性得证。
\end{proof}
\begin{corollary}
若 ${C}$ 是 Hilbert 空间 $\mathscr{X}$ 中的一个闭凸子集，则对 $\forall y \in \mathscr{X}$，存在唯一的 $x_0 \in{C}$，使得
\begin{equation*}
\|y - x_0\| = \inf_{x \in {C}} \|x - y\|.
\end{equation*}
\end{corollary}

\begin{proof}
只需考察集合 ${C} - \{y\} \triangleq \{x - y \mid x \in {C}\}$，它仍是 $\mathscr{X}$ 上的一个闭凸子集。根据定理 1.24，存在唯一的 $z_0 \in {C} - \{y\}$，使得
\begin{equation*}
\|z_0\| = \inf_{z \in C  - \{y\}} \|z\|.
\end{equation*}
令 $x_0 \triangleq z_0 + y$，便得证。
\end{proof}

\begin{example}
若 $M$ 是 Hilbert 空间 $\mathscr{X}$ 上的一个闭线性子空间，则对 $\forall y \in \mathscr{X}$，存在唯一的 $x_0 \in M$，使得
\begin{equation*}
\|y - x_0\| = \inf_{x \in M} \|x - y\|.
\end{equation*}
\end{example}

\begin{theorem}
设 ${C}$ 是内积空间 $\mathscr{X}$ 中的一个闭凸子集，$\forall y \in \mathscr{X}$，为了 $x_0$ 是 $y$ 在 ${C}$ 上的最佳逼近元，必须且仅须
\begin{equation*}
\operatorname{Re}(y - x_0, x_0 - x) \geqslant 0 \quad (\forall x \in {C}).
\end{equation*}
\end{theorem}

\begin{proof}
对 $\forall x \in C$，考察函数
\begin{equation*}
\varphi_x(t) = \|y - t x - (1 - t) x_0\|^2 \quad (t \in [0, 1]).
\end{equation*}
显然，为了 $x_0$ 是 $y$ 在 $C$ 上的最佳逼近元，必须且仅须
\begin{equation*}
\varphi_x(t) \geqslant \varphi_x(0) \quad (\forall x \in C, \forall t \in [0, 1]).
\end{equation*}
下证定理中的条件成立 $\Longleftrightarrow$ 上式成立。\\
因为
\begin{equation*}
\begin{aligned}
\varphi_x(t) &= \|(y - x_0) + t (x_0 - x)\|^2 \\
&= \|y - x_0\|^2 + 2t \operatorname{Re}(y - x_0, x_0 - x) + t^2 \|x_0 - x\|^2,
\end{aligned}
\end{equation*}
所以
\begin{equation*}
\varphi_x'(0) = 2 \operatorname{Re}(y - x_0, x_0 - x).
\end{equation*}
因此，
\begin{equation*}
\text{上式成立} \Longleftrightarrow \varphi_x'(0) \geqslant 0.
\end{equation*}
\end{proof}

\begin{corollary}
设 $M$ 是 Hilbert 空间 $\mathscr{X}$ 的一个闭线性子流形。$\forall x \in \mathscr{X}$，为了 $y$ 是 $x$ 在 $M$ 上的最佳逼近元，必须且仅须
\begin{equation*}
x - y \perp M - \{y\} \triangleq \{w = z - y \mid \forall z \in M\}.
\end{equation*}
\end{corollary}

\begin{proof}
由定理 1.24，为了 $y$ 是 $x$ 在 $M$ 上的最佳逼近元，必须且仅须
\begin{equation*}
\operatorname{Re}(x - y, y - z) \geqslant 0 \quad (\forall z \in M).
\end{equation*}
因为 $M$ 是线性流形，所以 $\forall z \in M$ 可表示为
\begin{equation*}
z = y + w \quad (w \in M - \{y\}).
\end{equation*}
将其代入，得
\begin{equation*}
\operatorname{Re}(x - y, w) \leqslant 0 \quad (\forall w \in M - \{y\}).
\end{equation*}
在上式中，用 $-w$ 替换 $w$，推出
\begin{equation*}
\operatorname{Re}(x - y, w) = 0 \quad (\forall w \in M - \{y\}),
\end{equation*}
进一步用 $iw$ 替换 $w$，得
\begin{equation*}
(x - y, w) = 0 \quad (\forall w \in M - \{y\}),
\end{equation*}
即得证。
\end{proof}

\begin{example}
当 $M$ 是闭线性子空间时，$M - \{y\} = M$。因此，为了 $y$ 是 $x$ 在 $M$ 上的最佳逼近元，必须且仅须 $x - y \perp M$。
\end{example}

\begin{corollary}
设 $M$ 是 Hilbert 空间 $\mathscr{X}$ 上的一个闭线性子空间。那么 $\forall x \in \mathscr{X}$，存在下列唯一的正交分解：
\begin{equation*}
x = y + z \quad \left(y \in M, z \in M^\perp\right).
\end{equation*}
\end{corollary}

\begin{proof}
取 $y$ 为 $x$ 在 $M$ 上的最佳逼近元，而 $z = x - y$，即可满足正交分解。同时，若存在另一分解
\begin{equation*}
x = y' + z' \quad \left(y' \in M, z' \in M^\perp\right),
\end{equation*}
则
\begin{equation*}
y - y' = z - z' \quad (\in M \cap M^\perp) = \theta,
\end{equation*}
即得分解的唯一性。
\end{proof}


\begin{problem}
设 $M$ 是 Hilbert 空间 $\mathscr{X}$ 的子集，证明：
\[
\left(M^{\perp}\right)^{\perp} = \overline{\operatorname{span} M}.
\]
\end{problem}
\begin{solution}
因为：
\[
\begin{aligned}
    & x \in M^{\perp} \Longleftrightarrow x \perp M \Longleftrightarrow x \perp \operatorname{span} M \\
    & \Longleftrightarrow x \perp \overline{\operatorname{span} M} \Longleftrightarrow x \in (\overline{\operatorname{span} M})^{\perp},
\end{aligned}
\]
所以有：
\[
M^{\perp} = (\overline{\operatorname{span} M})^{\perp}.
\]
为了证明 $\left(M^{\perp}\right)^{\perp} = \overline{\operatorname{span} M}$，等价于证明 $\left[(\overline{\operatorname{span} M})^{\perp}\right]^{\perp} = \overline{\operatorname{span} M}$。问题归结为：若 $M$ 是闭的，需证明 $\left(M^{\perp}\right)^{\perp} = M$。
显然：
\[
\forall x \in M \implies x \perp M^{\perp} \implies x \in \left(M^{\perp}\right)^{\perp}.
\]
因此：
\[
M \subseteq \left(M^{\perp}\right)^{\perp}.
\]
假设反命题成立，即存在 $x \in \left(M^{\perp}\right)^{\perp} \setminus M$。由于 $M$ 是 $\left(M^{\perp}\right)^{\perp}$ 的真闭子空间，$x$ 可正交分解为：
\[
x = y + z, \quad y \in M, \, z \in M^{\perp}, \, z \neq \theta.
\]
因为 $y \in M \subseteq \left(M^{\perp}\right)^{\perp}$，所以 $z = x - y \in \left(M^{\perp}\right)^{\perp}$。于是有：
\[
z \in M^{\perp} \cap \left(M^{\perp}\right)^{\perp}.
\]
但这意味着 $z = \theta$，与 $z \neq \theta$ 矛盾。因此假设不成立，即：
\[
\left(M^{\perp}\right)^{\perp} = M.
\]
结合以上步骤，得出结论：
\[
\left(M^{\perp}\right)^{\perp} = \overline{\operatorname{span} M}.
\]
\end{solution}

\begin{problem}{8}
设 $M$ 是 Hilbert 空间 $\mathscr{X}$ 的闭线性子空间，$\{e_n\}$ 与 $\{f_n\}$ 分别是 $M$ 与 $M^{\perp}$ 的正交规范基。证明：$\{e_n\} \cup \{f_n\}$ 构成 $\mathscr{X}$ 的正交规范基。
\end{problem}

\begin{solution}
根据正交分解定理，对于 $\forall x \in \mathscr{X}$，可以将其分解为：
\[
x = y + z, \quad \text{其中 } y \in M, \; z \in M^{\perp}.
\]
依题意，$y \in M$ 和 $z \in M^{\perp}$ 可以分别展开为 $\{e_n\}$ 和 $\{f_n\}$ 的线性组合：
\[
\begin{aligned}
    y &= \sum_{n=1}^\infty (y, e_n) e_n, \\
    z &= \sum_{m=1}^\infty (z, f_m) f_m.
\end{aligned}
\]
因此，可以将 $x$ 表示为：
\[
x = y + z = \sum_{n=1}^\infty (y, e_n) e_n + \sum_{m=1}^\infty (z, f_m) f_m.
\]
由 $\{e_n\}$ 是 $M$ 的正交规范集，以及 $\{f_n\}$ 是 $M^{\perp}$ 的正交规范集，且 $M \perp M^{\perp}$ 可知：
\[
\{e_n\} \cup \{f_n\} \text{ 是 } \mathscr{X} \text{ 中的正交规范基}
\]
\end{solution}
\begin{problem}
设 $H^2(D)$ 是按例 1.6.8 定义的内积空间.\\
(1) 如果 $u(z)$ 的 Taylor 展开式是 $u(z)=\sum_{k=0}^{\infty} b_k z^k$, 求证:
\begin{equation*}
\sum_{k=0}^{\infty} \frac{\left|b_k\right|^2}{1+k}<\infty
\end{equation*}
(2) 设 $u(z), v(z) \in H^2(D)$, 并且
\begin{equation*}
u(z)=\sum_{k=0}^{\infty} a_k z^k, \quad v(z)=\sum_{k=0}^{\infty} b_k z^k
\end{equation*}
求证：
\begin{equation*}
(u, v)=\pi \sum_{k=0}^{\infty} \frac{a_k \overline{b_k}}{k+1} .
\end{equation*}
(3) 设 $u(z) \in H^2(D)$, 求证:
\begin{equation*}
|u(z)| \leqslant \frac{\|u\|}{\sqrt{\pi}(1-|z|)} \quad(\forall|z|<1) .
\end{equation*}
(4) 验证 $H^2(D)$ 是 Hilbert 空间.
\end{problem}
\begin{solution}
(1)\begin{equation*}
(u,\varphi_n)=\iint\limits_{D} \left(\sum\limits_{k=0}^{\infty} b_k z^k\right)\sqrt{\frac{n}{\pi}}\bar{z}^{n-1} \dif x\dif y=b_{n-1}\sqrt{\frac{\pi}{n}}
\end{equation*}
由Bessel不等式得
\[
\sum_{k=0}^{\infty} \frac{\left|b_k\right|^2}{1+k}=\|u\|<\infty
\]
(2)\begin{equation*}
(u,v)=\iint\limits_{D}\left(\sum\limits_{k=0}^{\infty} a_k z^k\right) \left(\sum\limits_{k=0}^{\infty} \bar{b_k} \bar{z}^k\right)=\sum\limits_{k=0}^\infty\sum\limits_{j=0}^\infty\left(
\sqrt{\frac{a_k\pi}{k+1}}\varphi_{k+1},\sqrt{\frac{b_j\pi}{j+1}}\varphi_{j+1}
\right)\\=\sum\limits_{k=0}^\infty\frac{a_k\bar{b_k}\pi}{k+1}\left(
\varphi_{k+1},\varphi_{k+1}
\right)=\pi \sum_{k=0}^{\infty} \frac{a_k \overline{b_k}}{k+1} .
\end{equation*}
(3)当$|z|<1$时，
\[
\frac{1}{1-z}=\sum\limits_{k=0}^\infty z^k,\frac{1}{(1-z)^2}=\sum\limits_{k=0}^\infty (k+1)z^k,\frac{1}{(1-|z|)^2}=\sum\limits_{k=0}^\infty (k+1)|z|^k
\]
\begin{equation*}
\frac{\|u\|^2}{(1-|z|)^2} =\pi  \sum_{k=0}^{\infty} \frac{a_k ^2}{k+1}\sum\limits_{k=0}^\infty (k+1)|z|^k\geqslant \pi\sum_{k=0}^{\infty} \frac{a_k ^2}{k+1} (k+1)|z|^k=\pi |u|^2
\end{equation*}
所以\begin{equation*}
|u(z)| \leqslant \frac{\|u\|}{\sqrt{\pi}(1-|z|)} \quad(\forall|z|<1) .
\end{equation*}
(4)设$\{u_n\}$为$H^2(D)$中的基本列，由(3)可知$\forall z\in D$，
\begin{equation*}
|u_n(z)-u_m(z)| \leqslant \frac{\|u_n-u_m\|}{\sqrt{\pi}(1-|z|)} \to 0(n,m\to \infty)
\end{equation*}
所以$\{u_n(z)\}$为$\mathbb{C}$上的基本列，设$u_n(z)\to u(z)(\forall z\in D)$，考虑$D$中的任意闭集$F,\exists 0<r<1,F\subset B(0,r),\forall z\in F$，
\begin{equation*}
|u_n(z)-u_m(z)| \leqslant \frac{\|u_n-u_m\|}{\sqrt{\pi}(1-|r|)} \to 0(n,m\to \infty)
\end{equation*}
所以$\{u_n(z)\}$在$D$上 内闭一致收敛到$u(z)$，由维尔斯特拉斯定理知$u(z)$为$D$上的解析函数，将$\{u_n\}$看做$L^2(D)$中的Cauchy列，由$L^2(D)$空间的完备性 可知$\{u_n\}$收敛到$v$，那么$v,u$几乎处处相等，所以$\{u_n\}$在$H^2(D)$中也收敛于$u$，那么$h^2(D)$完备。
\end{solution}


\chapter{线性算子与线性泛函}
\section{线性算子}
\begin{definition}
	设 ${X}, {Y}$ 是两个线性空间，$D$ 是 ${X}$ 的一个线性子空间。映射 $T: D \to {Y}$ 称为算子，$D$ 称为 $T$ 的定义域，记作 $D(T)$。  
	\[
	R(T) = \{T x \mid x \in D\}
	\]
	称为 $T$ 的值域。如果  
	\[
	T(\alpha x + \beta y) = \alpha T x + \beta T y, \quad \forall x, y \in D, \ \forall \alpha, \beta \in \mathbb{K},
	\]
	则称 $T$ 是一个\textbf{线性算子}。
\end{definition}

\begin{example}
	设 ${X} = \mathbb{R}^n, {Y} = \mathbb{R}^m, T = \left(t_{ij}\right)_{m \times n}$。如果  
	\[
	x \mapsto T x = \left( \sum_{j=1}^n t_{ij} x_j \right)_{i=1}^m, \quad \forall x = (x_1, x_2, \cdots, x_n) \in \mathbb{R}^n,
	\]
	那么 $T$ 是一个线性算子。
\end{example}

\begin{example}
	设 ${X} = {Y} = C^\infty(\bar{\Omega})$，又设微分多项式  
	\[
	P\left(\partial_x\right) = \sum_{|\alpha| \leq m} a_\alpha(x) \partial_x^\alpha, \quad a_\alpha(x) \in C^\infty(\bar{\Omega}).
	\]
	如果 $T: u(x) \mapsto P\left(\partial_x\right) u(x)$ （$\forall u \in {X}$），  
	那么 $T$ 是一个从 ${X}$ 到 ${Y}$ 的线性算子。  \\
	若 ${X} = {Y} = L^2(\Omega), D(T) = C^m(\bar{\Omega})$，则上述定义的算子 $T$ 也是线性的。
\end{example}

\begin{example}
	设 ${X} = L^1(-\infty, \infty), {Y} = L^\infty(-\infty, \infty)$，若规定  
	\[
	T: u(x) \mapsto \int_{-\infty}^\infty e^{i \xi \cdot x} u(x) \, \dif x, \quad \forall u \in {X},
	\]
	那么 $T$ 是一个从 ${X}$ 到 ${Y}$ 的线性算子。
\end{example}

\begin{definition}
	取值于实数（复数）的线性算子称为\textbf{实（复）线性泛函}，记作 $f(x)$ 或 $\langle f, x \rangle$（即线性函数）。
\end{definition}

	

	\begin{definition}
		设 ${X}, {Y}$ 是 $F^*$ 空间，$D(T) \subset {X}$。称线性算子 $T: D(T) \to {Y}$ 在 $x_0 \in D(T)$ 是连续的，如果  
		\[
		x_n \in D(T), \quad x_n \to x_0 \Longrightarrow T x_n \to T x_0.
		\]
	\end{definition}
	
	\begin{proposition}
		对于线性算子 $T$，为了它在 $D(T)$ 内处处连续，\textbf{必须且仅须}它在 $x = \theta$ 处连续。  
	\end{proposition}
	
	\begin{proof}
		若 $T$ 在 $\theta$ 处连续，则对任意 $x_n, x_0 \in D(T)$ 且 $x_n \to x_0$，有  
		\[
		x_n - x_0 \to \theta \Longrightarrow T x_n - T x_0 = T(x_n - x_0) \to T \theta = \theta.
		\]
		反之，若 $T$ 在 $D(T)$ 内处处连续，则在 $x = \theta$ 处必然连续，证毕。
	\end{proof}
	
	\begin{definition}
		设 ${X}, {Y}$ 都是 $B^*$ 空间。称线性算子 $T: {X} \to {Y}$ 是\textbf{有界}的，如果存在常数 $M \geq 0$，使得  
		\[
		\|T x\|_{{Y}} \leq M \|x\|_{{X}}, \quad \forall x \in {X}.
		\]
	\end{definition}
	
	\begin{proposition}
		设 ${X}, {Y}$ 都是 $B^*$ 空间，为了线性算子 $T$ 连续，\textbf{必须且仅须} $T$ 有界。
	\end{proposition}
	
	\begin{proof}
		充分性显然。下证必要性：  
		若不然，则存在 $x_n \in {X}$，使得  
		\[
		\|T x_n\| > n \|x_n\|.
		\]
		令 $y_n = \frac{x_n}{n \|x_n\|}$，则 $\|T y_n\| > 1$。但 $y_n \to \theta \ (n \to \infty)$，与 $T$ 的连续性矛盾，证毕。
	\end{proof}
	
	\begin{definition}
		用 ${L}({X}, {Y})$ 表示从 ${X}$ 到 ${Y}$ 的所有有界线性算子的全体，并规定  
		\[
		\|T\| = \sup_{x \in {X} \backslash \{\theta\}} \frac{\|T x\|}{\|x\|} = \sup_{\|x\| = 1} \|T x\|,
		\]
		称为 $T \in {L}({X}, {Y})$ 的范数。特别地，用 ${L}({X})$ 表示 ${L}({X}, {X})$，用 ${X}^*$ 表示 $L({X}, \mathbb{K})$，即 ${X}^*$ 表示 ${X}$ 上的有界线性泛函全体。
	\end{definition}
	
	\begin{proposition}
		设 ${X}$ 是 $B^*$ 空间，${Y}$ 是 $B$ 空间。若在 ${L}({X}, {Y})$ 上规定线性运算：
		\[
		\left(\alpha_1 T_1 + \alpha_2 T_2\right)(x) = \alpha_1 T_1 x + \alpha_2 T_2 x, \quad \forall x \in {X},
		\]
		其中 $\alpha_1, \alpha_2 \in K, \ T_1, T_2 \in {L}({X}, {Y})$，则 ${L}({X}, {Y})$ 按 $\|\cdot\|$ 构成一个 Banach 空间。
	\end{proposition}
	
	\begin{proof}
		显然 $L({X}, {Y})$ 是一个线性空间。下证 $\|\cdot\|$ 是范数：  
		\[
		\|T\| \geq 0, \quad \|T\| = 0 \Longleftrightarrow T x = \theta, \ \forall x \in {X} \Longleftrightarrow T = \theta.
		\]
		\[
		\|T_1 + T_2\| = \sup_{\|x\|=1} \|T_1 x + T_2 x\| \leq \sup_{\|x\|=1} \|T_1 x\| + \sup_{\|x\|=1} \|T_2 x\| = \|T_1\| + \|T_2\|.
		\]
		\[
		\|\alpha T\| = \sup_{\|x\|=1} \|\alpha T x\| = |\alpha| \sup_{\|x\|=1} \|T x\| = |\alpha| \|T\|.
		\]
		再证完备性：设 $\{T_n\}_{n=1}^\infty$ 是一个基本列，则 $\forall \varepsilon > 0, \exists N = N(\varepsilon)$，使得 $\forall x \in {X}$ 有  
		\[
		\|T_{n+p} x - T_n x\| \leq \varepsilon \|x\|, \quad \forall n > N, \ p \in \mathbb{N}.
		\]
		于是 $T_n x \to y \in {Y}$（$n \to \infty$）。记此 $y=T x$, 我们要证 $T \in$ ${L} ( {X} , {Y} )$. 不难看出 $T$ 是线性的, 再证其有界. 事实上, $\exists n \in N$,使得
		\begin{equation*}
			\begin{aligned}
				\|T x\|=\|y\| & \leqslant\left\|T_n x\right\|+1 \\
				& \leqslant\left(\left\|T_n\right\|+1\right)\|x\| \quad(\forall x \in {X} ,\|x\|=1)
			\end{aligned}
		\end{equation*}
		即得 $\|T\| \leqslant\left\|T_n\right\|+1$.
	\end{proof}
	
	\begin{example}
		设 $T$ 是有限维 $B^*$ 空间 ${X}$ 到 ${Y}$ 的线性映射，则 $T$ 必是连续的。
	\end{example}
	
	\begin{example}
		Hilbert 空间 ${X}$ 上的正交投影算子。设 $M$ 是 ${X}$ 的一个闭线性子空间。依正交分解定理，$\forall x \in {X}$，存在唯一分解  
		\[
		x = y + z, \quad y \in M, \ z \in M^\perp.
		\]
		对应 $x \mapsto y$ 称为从 ${X}$ 到 $M$ 的正交投影算子，记作 $P_M$。在不强调子空间 $M$ 时，省略 $M$ 记为 $P$。
	\end{example}
	
	\begin{proof}
	我们来证明 $P$ 还是一个连续线性算子, 并且如果 $M \neq\{\theta\}$, 那么 $\|P\|=1$.
	先证线性. 设 $x_i=P x_i+z_i(i=1,2)$, 其中 $z_i \in M^{\perp}$. 这时,
	\begin{equation*}
		\begin{aligned}
			\alpha_1 x_1+\alpha_2 x_2= & \left(\alpha_1 P x_1+\alpha_2 P x_2\right)+\left(\alpha_1 z_1+\alpha_2 z_2\right) \\
			& \left(\forall \alpha_1, \alpha_2 \in K \right) .
		\end{aligned}
	\end{equation*}
	因为 $\alpha_1 z_1+\alpha_2 z_2 \in M^{\perp}$, 而 $\alpha_1 P x_1+\alpha_2 P x_2 \in M$ ，所以
	\begin{equation*}
		P\left(\alpha_1 x_1+\alpha_2 x_2\right)=\alpha_1 P x_1+\alpha_2 P x_2
	\end{equation*}
	即得 $P$ 是线性算子。
	其次证连续，这是由于 $\|P x\|^2=\|x\|^2-\|z\|^2 \leqslant\|x\|^2$ 。因此， $\|P x\| \leqslant\|x\|$, 或者 $\|P\| \leqslant 1$ 。
	最后, 当 $M \neq\{\theta\}$ 时, 任取 $x \in M \backslash\{\theta\}$, 便有 $\|P x\|=\|x\|$, 从而 $\|P\|=1$.
	\end{proof}


	
	\begin{example}
		设 ${X} = L^1[a, b]$, ${Y} = C[a, b]$，定义积分算子 $T$ 为：
		\[
		(T f)(x) = \int_a^x f(t) \, dt
		\]
		求解以下问题：\( T \in L({X}, {Y}) \) 时的范数，\( T \in L({X}) \) 时的范数
	\end{example}

\begin{proof}
		 一方面，对于任意 $f \in L^1[a, b]$，有：
\[
\|T f\|_c = \max_{x \in [a, b]} |(T f)(x)| = \max_{x \in [a, b]} \left| \int_a^x f(t) \, \dif t \right| 
\leqslant \max_{x \in [a, b]}\int_a^b |f(t)| \, \dif t = \|f\|_1,
\]
因此 $\|T\| \leqslant 1$。
\\ 另一方面，取 $f_0(t) = \frac{1}{b-a}$，对于 $t \in [a, b]$，则有 $\|f_0\|_1 = 1$，并且：
\[
\begin{aligned}
	\|T\| & \geqslant \|T f_0\|_c = \max_{x \in [a, b]} |(T f_0)(x)| \\
	& = \max_{x \in [a, b]} \left|\int_a^x \frac{1}{b-a} \, \dif t\right| \\
	& = \max_{x \in [a, b]} \frac{x-a}{b-a} = 1.
\end{aligned}
\]
因此，综合以上两方面，我们得到 $\|T\| = 1$。\\
一方面，对于任意 $f \in L^1[a, b]$，有：
\[
\begin{aligned}
	\|T f\|_1 & = \int_a^b \left| \int_a^x f(t) \, \dif t \right| \dif x \\
	& \leqslant \int_a^b \int_a^x |f(t)| \, \dif t \, \dif x \\
	& \leqslant \int_a^b \|f\|_1 \, \dif x = (b - a) \|f\|_1,
\end{aligned}
\]
因此 $\|T\| \leqslant b - a$。\\另一方面，取
\[
f_{\varepsilon}(t) =
\begin{cases}
	\frac{1}{\varepsilon}, & \text{if } a \leq t \leq a + \varepsilon, \\
	0, & \text{if } a + \varepsilon < t \leq b.
\end{cases}
\]
则有 $\|f_{\varepsilon}\|_1 = 1$，并且：
\[
\begin{aligned}
	(T f_{\varepsilon})(x) &= \int_a^x f_{\varepsilon}(t) \, \dif t = 
	\begin{cases}
		\frac{x - a}{\varepsilon}, & \text{if } a \leq x \leq a + \varepsilon, \\
		1, & \text{if } a + \varepsilon < x \leq b.
	\end{cases} \\
	\|T f_{\varepsilon}\|_1 &= \int_a^b \left| (T f_{\varepsilon})(x) \right| \dif x \\
	& = \int_a^{a+\varepsilon} \frac{x - a}{\varepsilon} \, \dif x + \int_{a+\varepsilon}^b 1 \, \dif x \\
	& = b - a - \frac{1}{2} \varepsilon.
\end{aligned}
\]
因此，我们得到：
\[
\|T\| \geqslant \|T f_{\varepsilon}\|_1 = b - a - \frac{1}{2} \varepsilon.
\]
令 $\varepsilon \to 0$，从而得到 $\|T\| \geqslant b - a$。
综合以上两方面，我们确定 $\|T\| = b - a$。
\end{proof}

${{{X}}},{Y}$均为$B$空间。
\begin{problem}
求证: $T \in {L}( {X} , {Y} )$ 的充要条件是 $T$ 为线性算子并将 ${X}$ 中的有界集映为 ${Y}$ 中的有界集。

\end{problem}
\begin{proof}
必要性显然成立。\\
充分性。设$S=\{x\in {X}|\|x\|\leqslant1\}$是${X}$中的有界集，则$\exists M>0$，使得
\[
\|Tx\|\leqslant M(\forall x\in S)
\]
于是$\forall \theta\neq x\in {X}$,$\|T\frac{x}{\|x\|}\|\leqslant M$成立，则$\|T{x}\|\leqslant M{\|x\|}$成立。\\
当$x=\theta$时上式也成立，所以$\|T{x}\|\leqslant M{\|x\|}(\forall x\in{X})$成立，$T \in {L}( {X} , {Y} )$。
\end{proof}


\begin{problem}
设 $y(t) \in C[0,1]$ ，定义 $C[0,1]$ 上的泛函
\begin{equation*}
f(x)=\int_0^1 x(t) y(t) \dif t \quad(\forall x \in C[0,1]),
\end{equation*}
求 $\|f\|$.
\end{problem}
\begin{solution}
当$y(t)=0$时，$\|f\|=0=\int_0^1|y(t)|$成立，下设$y(t)$不恒为$0$。\\
\[\|f\|=\sup\limits_{\|x\|=1}\|f(x)\|=\sup\limits_{\|x\|=1}\left|\int_0^1 x(t) y(t) \dif t\right|\leqslant\sup\limits_{\|x\|=1}\left|\sup\limits_{0\leqslant t\leqslant 1}x(t)\right|\int_0^1 | y(t)| \dif t\leqslant\int_0^1 | y(t)| \dif t\]
$\forall \varepsilon>0$，由$y(t)$在$[0,1]$上的一致连续性知，$\exists n\in \mathbb{N}$将$[0,1]$$n$等分得到的区间上$\sup\limits_{t_1,t_2\in[\frac{i}{n},\frac{i+1}{n}]}|y(t_1)-y(t_2)|<\varepsilon,i=0,1\cdots n-1$，将这$n$个区间分为两类。
\begin{equation*}
    \begin{aligned}
    A=\{i|y(t)\text{在}[\frac{i}{n},\frac{i+1}{n}]\text{上不含有零点}\}\\
        B=\{i|y(t)\text{在}[\frac{i}{n},\frac{i+1}{n}]\text{上含有零点}\}\\
    \end{aligned}
\end{equation*}
取$x(t)=\begin{cases}
    \operatorname{sgn}(y(t)),t\in [\frac{i}{n},\frac{i+1}{n}],i\in A\\
\text{连续线性函数 } ,t\in [\frac{i}{n},\frac{i+1}{n}],i\in B
\end{cases}$，由于设$y(t)$不恒为$0$，所以$\|x\|=1$且$x\in C[0,1]$，因为集合$B$所对应的区间都含有零点,所以在$[\frac{i}{n},\frac{i+1}{n}],i\in B,|y(t)|<\varepsilon$。
\begin{equation*}
    \begin{aligned}
    f(x)&=\int_0^1 x(t) y(t) \dif t=\sum\limits_{i\in A}\int_{\frac{i}{n}}^{\frac{i+1}{n}} x(t) y(t) \dif t+
    \sum\limits_{i\in B}\int_{\frac{i}{n}}^{\frac{i+1}{n}} x(t)y(t) \dif t\\
&=\int_{0}^{1} |y(t)| \dif t+
    \sum\limits_{i\in B}\int_{\frac{i}{n}}^{\frac{i+1}{n}} x(t)y(t)-|y(t)| \dif t>\int_{0}^{1} |y(t)| \dif t-2\varepsilon
    \end{aligned}
\end{equation*}
由$\varepsilon$任意性可知$\|f\|\geqslant \int_{0}^{1} |y(t)| \dif t$，所以$\|f\|=\int_{0}^{1} |y(t)| \dif t$。
\end{solution}

\begin{problem}
 设 $f$ 是 ${X}$ 上的非零线性有界泛函，令 $d=\inf \{\|x\|$ $\mid f(x)=1,x\in{X}\}$, 求证: $\|f\|=1 / d$ 。
\end{problem}
\begin{proof}
$\|f\|=\sup\limits_{x\neq \theta}\frac{\|f(x)\|}{\|x\|}$\    ，所以$\frac{\|f(x)\|}{\|f\|}\leqslant \|x\|,\forall x\in {X}$,取$f(x)=1$则$d\geqslant \frac{1}{\|f\|}$成立。\\
$\forall \varepsilon >0,\exists x_\varepsilon\in {X}$使得$\frac{\|f(x_\varepsilon)\|}{\|x_\varepsilon\|}>\|f\|-\varepsilon$，而$f(\frac{x_\varepsilon}{f(x_\varepsilon)})=1$，所以$d\leqslant\frac{\|x_\varepsilon\|}{\|f(x_\varepsilon)\|}<\frac{1}{\|f\|-\varepsilon} $，由$\varepsilon$任意性知$d\leqslant \frac{1}{\|f\|}$。\\
所以$\|f\|=\frac{1}{d}$
\end{proof}
\begin{problem}
 设 $T: {X} \rightarrow {Y}$ 是线性的, 令
\begin{equation*}
N(T) \triangleq\{x \in {X} \mid T x=\theta\} .
\end{equation*}
(1)若 $T \in {L} ( {X} , {Y} )$ ，求证： $N(T)$ 是 ${X}$ 的闭线性子空间。\\
(2)问 $N(T)$ 是 ${X}$ 的闭线性子空间能否推出 $T \in {L} ( {X},{Y})$ ?\\
(3)若 $f$ 是线性泛函，求证： $f \in {X} ^* \Longleftrightarrow N(f)$ 是闭线性子空间。

\end{problem}
\begin{solution}
(1)$\forall \alpha,\beta\in {X},k\in \mathbb{K}$
\[T(\alpha+\beta)=T\alpha+T\beta=\theta+\theta=\theta,T(k\alpha)=kT\alpha=k\theta=\theta\]
所以$\alpha+\beta \in N(T),k\alpha\in N(T)$，$N(T)$是线性子空间。\\
设$\{x_n\}$是$N(T)$中的点列且收敛于$x_0$，由于$T$是连续的，所以$\theta=Tx_n\rightarrow Tx_0$
，则$Tx_0=\theta ,x\in N(T)$，所以$N(T)$是闭集。\\
(2)不能，反例如下。定义$B$空间$X$及其上的线性算子$T$。
\begin{equation*}
    \begin{aligned}
    X=\{(\xi_1,\xi_2,\cdots,\xi_n,\cdots)|\sum\limits_{n=1}^{\infty}|\xi_n|<+\infty\},\|x\|=\sup_{n\geqslant 1}|\xi_n|\\
    \forall x=(\xi_1,\xi_2,\cdots,\xi_n,\cdots)\in X,a  =(1,-1,0, \cdots) \in X ,Tx=x-a\sum\limits_{n=1}^{+\infty}\xi_n
    \end{aligned}
\end{equation*}
由$T x=\theta \Rightarrow x=a \sum\limits_{n=1}^{+\infty}\xi_n \Rightarrow \sum\limits_{n=1}^{+\infty}\xi_n=0 \Rightarrow x=a \sum\limits_{n=1}^{+\infty}\xi_n=\theta$, 即 $N(T)=\{\theta\}$. \\
下面证明 $T$ 无界. 先证 $f(x)=\sum\limits_{n=1}^{\infty} \xi_n$ 不连续. 令 $e_k=\{\overbrace{0,0, \cdots 0,1}^k, 0, \cdots\}$,
$x_n=\sum\limits_{k=1}^n e_k \in X ,\left\|x_n\right\|=1, f\left(x_n\right)=n \Rightarrow \frac{\left|f\left(x_n\right)\right|}{\| x_n \mid}=n \rightarrow \infty$, 即 $f$ 无界, 故$f$
不连续。\\
如果 $T x$ 有界,则 $x-T x$ 有界，从而
$|f(x)| \stackrel{\| a \|=1}{=}\|a\|| f(x) |=\|a f(x)\|=\|x-T x\| \leqslant M\|x\|$, 即 $f$有界，矛盾。\\
(3)必要性由(1)可知。
\\充分性。若$f$不是有界的，则
$\forall n \in N , \exists x_n\in{X},\left\|x_n\right\|=1,\left|f\left(x_n\right)\right| \geqslant n$，$y_n=\frac{x_n}{f\left(x_n\right)}-\frac{x_1}{f\left(x_1\right)} \Rightarrow f\left(y_n\right)=0 \Rightarrow y_n \in N(f)$, 但 $y_n \rightarrow-\frac{x_1}{f\left(x_1\right)} \notin N(f)$，与 $N(f)$ 闭矛盾。
\end{solution}

\begin{problem}
 设 $f$ 是 ${X}$ 上的线性泛函，记
\begin{equation*}
H_f^\lambda \triangleq\{x \in D \mid f(x)=\lambda\} \quad(\forall \lambda \in \mathbb{K} ) .
\end{equation*}
如果 $f \in {X}^*$ ，且 $\|f\|=1$ ，求证：\\
(1) $|f(x)|=\inf \left\{\|x-z\| \mid \forall z \in H_f^0\right\}(\forall x \in {X} )$;\\
(2)$\forall \lambda \in \mathbb{K}$，$H_f^\lambda$上的任意一点到$x$到$H_f^0$的距离都等于$|\lambda|$。\\
对 ${X} = \mathbb{R} ^2, \mathbb{K} = \mathbb{R} ^1$ 情形解释(1)和(2)的几何意义。
\end{problem}
\begin{proof}
    (1)由$\|f\|=1$知，$\forall x\in {X},\forall z\in H_f^0,|f(x)|=|f(x-z)|\leqslant \|x-z\|,|f(x-z)|\leqslant\inf \left\{\|x-z\| |\forall z \in H_f^0\right\}$。\\
另一方面, $\forall z \notin N(f)$, 对 $\forall x \in {X}$, 令
\begin{equation*}
	y=x-\frac{f(x)}{f(z)} z,
\end{equation*}
则有 $y \in N(f)$, 且 $f(z)(x-y)=f(x) z$, 从而
\begin{equation*}
	\frac{|f(z)|}{\|z\|}\|x-y\|=|f(x)| .
\end{equation*}
两边对 $z \neq \theta$ 取上确界，得到
\begin{equation*}
	\|x-y\| \sup _{z \neq \theta} \frac{|f(z)|}{\|z\|} \leqslant|f(x)| .
\end{equation*}
注意到 $\sup _{z \neq \theta} \frac{|f(z)|}{\|z\|}=\|f\|=1$ ，便有对 $y \in N(f)$ ，成立
\begin{equation*}
	\|x-y\| \leqslant|f(x)| \quad(\forall x \in {X} ) .
\end{equation*}
进一步对 $\forall y \in N(f)$, 上式两边取下确界, 即得
\begin{equation*}
	\rho(x, N(f))=\inf _{y \in N(f)}\|x-y\| \leqslant|f(x)|,
\end{equation*}
所以 $|f(x)|=\inf \left\{\|x-z\| \mid \forall z \in H_f^0\right\}(\forall x \in {X} )$。\\
(2)由(1)知$\forall\lambda \in \mathbb{K},\forall x\in H_f^\lambda,\rho(x,H_f^0)=|f(x)|=|\lambda|$。\\
(1)和 (2) 的几何意义\\
$\forall x=(\xi, \eta) \in \mathbb{R} ^2, f(x)=\alpha \xi+\beta \eta$ ，其中 $(\alpha, \beta)$ 是由 $f$ 确定的实数对:
$\alpha=f\left((1,0)\right), \quad \beta=f\left((0,1)\right)$ 。如果 $\|f\|=\sqrt{\alpha^2+\beta^2}=1$, 则$|f(x)|=|\alpha \xi+\beta \eta|$ 表示点 $x=(\xi, \eta)$ 到通过原点的直线
\begin{equation*}
H_f^0=\{x=(\xi, \eta) \quad \mid f(x)=\alpha \xi+\beta \eta=0\}
\end{equation*}
的距离，即 $|f(x)|=|\alpha \xi+\beta \eta|=\rho\left(x, H_f^0\right), \forall x=(\xi, \eta) \in \mathbb{R} ^2$.
\begin{equation*}
H_f^\lambda=\{x=(\xi, \eta) \mid f(x)=\alpha \xi+\beta \eta=\lambda\}
\end{equation*}
$H_f^\lambda$ 和 $H_f^0$ 是互相平行于直线， $\forall x \in H_f^\lambda, \rho\left(x, H_f^0\right)=\rho\left(\theta, H_f^\lambda\right)$ ,
\begin{equation*}
\rho\left(\theta, H_f^\lambda\right)=\left|\frac{\alpha \xi+\beta \eta-\lambda}{\sqrt{\alpha^2+\beta^2}}\right|_{(\xi, \eta)=(0,0)}=|\lambda|
\end{equation*}    
\end{proof}
\section{Riesz 表示定理}
\begin{theorem}[Riesz 表示定理（Hilbert 空间）]
	设 $f$ 是 Hilbert 空间 ${X}$ 上的一个连续线性泛函，则必存在唯一的 $y_f \in {X}$，使得
	\[
	f(x) = \left( x, y_f \right) \quad (\forall x \in {X}).
	\]
\end{theorem}
\begin{proof}
 不妨设 $f$ 不是 0 泛函，考察集合 $M \triangleq \{ x \in {X} \mid f(x) = 0 \}$。由于 $f$ 是连续线性的，则 $M$ 是一个真闭线性子空间，任取 $x_0 \perp M$ （由正交分解定理，这 $x_0$ 是存在的）。不妨设 $\|x_0\| = 1$，${X}$ 中任意元素 $x$ 可以分解如下：
\[
x = \alpha x_0 + y,
\]
其中 $y \in M, \alpha = \frac{f(x)}{f(x_0)}$。这是因为当令 $y = x - \alpha x_0$ 时，
\[
f(y) = f(x - \alpha x_0) = f(x) - \alpha f(x_0) = 0.
\]
在上式两边同时与 $x_0$ 做内积，我们得
\[
\alpha = \left( x, x_0 \right),
\]
于是
\[
f(x) = \alpha f(x_0) = \left( x, \overline{f(x_0)} x_0 \right).
\]
取 $y_f = \overline{f(x_0)} x_0$，这就是我们要求的。\\
再证唯一性。若存在 $y, y' \in {X}$ 满足
\[
f(x) = (x, y) = (x, y') \quad (\forall x \in {X}),
\]
那么
\[
(x, y - y') = 0 \quad (\forall x \in {X}).
\]
特别取 $x = y - y'$，就推得 $y = y'$。
\end{proof}
\begin{note} 这个定理的几何意义如下：连续线性泛函 $f(x)$ 的等值面都是互相平行的超平面（见习题 2.1），因此每个向量 $x$ 的泛函值 $f(x)$ 应由 $x$ 的垂直于这些等值面的分量所决定。
\\我们还知道，$\|f\| = \|y_f\|$。事实上，
\[
|f(x)| \leqslant \|y_f\| \cdot \|x\| \quad (\forall x \in {X}), \quad \text{或} \quad \|f\| \leqslant \|y_f\|.
\]
另一方面，取 $x = y_f$，可推得 $\|y_f\| \leqslant \|f\|$，即得到 $\|f\| = \|y_f\|$。
\end{note}
\begin{theorem}
	设 ${X}$ 是一个 Hilbert 空间，$a(x, y)$ 是 ${X}$ 上的共轭双线性函数，并且存在 $M > 0$，使得
	\[
	|a(x, y)| \leqslant M \|x\| \cdot \|y\| \quad (\forall x, y \in {X}),
	\]
	则存在唯一的 $A \in L({X})$，使得
	\[
	a(x, y) = (x, A y) \quad (\forall x, y \in {X}),
	\]
	且
	\[
	\|A\| = \sup_{\substack{(x, y) \in {X} \times {X} \\ x \neq 0, y \neq 0}} \frac{|a(x, y)|}{\|x\| \cdot \|y\|}.
	\]
\end{theorem}
\begin{proof}固定 $y \in {X}$，定义映射 $x \mapsto a(x, y)$，它是一个连续线性泛函。由 Riesz 表示定理，存在 $z = z(y) \in {X}$，使得
\[
a(x, y) = (x, z) \quad (\forall x \in {X}).
\]
定义映射 $A : y \mapsto z(y)$，便有 $a(x, y) = (x, A y) (\forall x, y \in {X})$。又因为
\[
\left( x, A\left( \alpha_1 y_1 + \alpha_2 y_2 \right) \right) = a\left( x, \alpha_1 y_1 + \alpha_2 y_2 \right)
= \overline{\alpha_1} a\left( x, y_1 \right) + \overline{\alpha_2} a\left( x, y_2 \right)
= \overline{\alpha_1} \left( x, A y_1 \right) + \overline{\alpha_2} \left( x, A y_2 \right)
= \left( x, \alpha_1 A y_1 + \alpha_2 A y_2 \right),
\]
因此 $A$ 是线性的，并且
\[
\|A y\| = \sup_{x \in {X} \setminus \{ 0 \}} \frac{|a(x, y)|}{\|x\|} \leqslant M \|y\|,
\]
即得 $A \in L({X})$，并满足 (2.2.3) 式。
\end{proof}
\begin{problem}
设 $f_1, f_2, \cdots, f_n$ 是 $H$ 上的一组线性有界泛函,
\begin{equation*}
\begin{gathered}
M \triangleq \bigcap_{k=1}^n N\left(f_k\right), \quad N\left(f_k\right) \triangleq\left\{x \in H \mid f_k(x)=0\right\} 
(k=1,2, \cdots, n) .
\end{gathered}
\end{equation*}
$\forall x_0 \in H$ ，记 $y_0$ 为 $x_0$ 在 $M$ 上的正交投影，求证： $\exists y_1, y_2, \cdots, y_n$ $\in H$ 及 $a_1, a_2, \cdots, a_n \in \mathbb{K}$ ，使得
\begin{equation*}
y_0=x_0-\sum_{k=1}^n a_k y_k
\end{equation*}
\end{problem}
\begin{proof}
由Riesz表示定理得
\begin{equation*}
\begin{aligned}
& \exists y_k \in H, f_k(x)=\left(x, y_k\right)(k=1,2, \cdots, n) \forall x \in M=\bigcap_{k=1}^n N\left(f_k\right) \Rightarrow\left(x, y_k\right)=0(k=1,2, \cdots, n, \forall x \in M)
\end{aligned}
\end{equation*}
\begin{itemize}
    \item 假设 $\left\{y_k\right\}_{k=1}^n$ 的极大线性无关组就是本身. 用 Gram-Schmidt过程，构造出一个正交规范集 $\left\{z_k\right\}_{k=1}^n$ s.t. \\$\operatorname{span}\left\{z_k\right\}_{k=1}^n=$  $\operatorname{span}\left\{y_k\right\}_{k=1}^n$. 则有 $z_k \perp M(k=1,2, \cdots, n)$. 令 $z_0=x_0-y_0$, 则 $z_0 \perp M$.对 $\forall x \in H$,
\begin{equation*}
    \begin{aligned}
        f_k(x-\sum_{k=1}^n\left(x, z_k\right) z_k)=\left(x-\sum_{k=1}^n\left(x, z_k\right) z_k,\sum_{k=1}^n b_k z_k
        \right)=0
    \end{aligned}
\end{equation*}
\begin{equation*}
    \begin{aligned}
 \left(x, z_0\right)=& \left(\sum_{k=1}^n\left(x, z_k\right) z_k+{x-\sum_{k=1}^n\left(x, z_k\right) z_k}, z_0\right)\\
=&\left(\sum_{k=1}^n\left(x, z_k\right) z_k, z_0\right)=\left(x, \sum_{k=1}^n \overline{\left(z_k, z_0\right)} z_k\right) 
    \end{aligned}
\end{equation*}
则$z_0= \sum\limits_{k=1}^n\overline{\left(z_k, z_0\right)} z_k,x_0-y_0\in \operatorname{span}\{y_k\}_{k=1}^n$，所以$\exists y_1, y_2, \cdots, y_n$ $\in H$ 及 $a_1, a_2, \cdots, a_n \in \mathbb{K}$ ，使得
\begin{equation*}
y_0=x_0-\sum_{k=1}^n a_k y_k
\end{equation*}
\item 因为
\[
x\in M\Longleftrightarrow (x,y_k)=0,k=1,2,\cdots,n\]
所以$M=(\operatorname{span}\{y_k\}_{k=1}^n)^\perp,M^\perp =\operatorname{span}\{y_k\}_{k=1}^n$，由正交分解定理即可证明。
\end{itemize}

\end{proof}

\begin{theorem}
设 $(\Omega, B , \mu),(\Omega, B , \nu)$ 是两个 $\sigma$-有限测度, 且 $\nu$关于 $\mu$ 绝对连续, 即
\begin{equation*}
E \in B , \quad \mu(E)=0 \Rightarrow \nu(E)=0
\end{equation*}
则存在关于 $\mu$ 的可测函数 $g$, 且 $g(x) \geqslant 0$ a.e. $\mu$, 使得

\begin{equation*}
\nu(E)=\int_E g(x) \dif \mu, \quad \forall E \in B
\end{equation*}

\end{theorem}
\begin{proof}
先假设 $\mu(\Omega)<\infty$. 考虑实 Hilbert 空间 $L^2(\Omega,(\mu+\nu)$ ),其范数为
\begin{equation*}
\|u\|^2=\int_{\Omega} u^2(x) \dif(\mu+\nu) .
\end{equation*}
令 $l(u)=\int_{\Omega} u \dif \mu$ ，显然 $l(u)$ 关于 $u$ 线性，由 Cauchy-Schwarz 不等式，
\begin{equation*}
\begin{aligned}
\|l(u)\| & \leqslant \mu(\Omega)^{\frac{1}{2}}\left(\int_{\Omega} u^2 \dif \mu\right)^{\frac{1}{2}} \\
& \leqslant \mu(\Omega)^{\frac{1}{2}}\|u\|, \quad \forall u \in L^2(\Omega,(\mu+\nu))
\end{aligned}
\end{equation*}
它还是有界的. 根据 Riesz 表示定理，存在函数 $v \in$ $L^2(\Omega,(\mu+\nu))$ ，使得
\begin{equation*}
\int_{\Omega} u \dif \mu=\int_{\Omega} u v \dif(\mu+\nu),
\end{equation*}
即
\begin{equation*}
\int_{\Omega} u(1-v) \dif \mu=\int_{\Omega} u v \dif \nu, \quad \forall u \in L^2(\Omega,(\mu+\nu)) .
\end{equation*}
我们断言，
\begin{equation*}
0<v(x) \leqslant 1 \quad \text { a.e. } \mu \text {. }
\end{equation*}
为此, 令 $F=\{x \in \Omega \mid v(x) \leqslant 0\}$, 取 $u(x)=\chi_F(x)$ 代人 (2.2.14)式，有
\begin{equation*}
\int_F(1-v) \dif \mu=\int_F v \dif \nu,
\end{equation*}
即
\begin{equation*}
\mu(F)=\int_F \dif \mu=\int_F v \dif(\mu+\nu) \leqslant 0
\end{equation*}
从而 $\mu(F)=0$.\\
同样, 令 $G=\{x \in \Omega \mid v(x)>1\}$, 取 $u(x)=\chi_G(x)$ 代入，有
\begin{equation*}
0 \geqslant \int_G(1-v) \dif \mu=\int_G v \dif \nu \geqslant \nu(G) \geqslant 0,
\end{equation*}
即
\begin{equation*}
\int_G(1-v) \dif \mu=0
\end{equation*}
因为 $1-v(x)<0, x \in G$, 所以 $\mu(G)=0$.\\
这就证明了 $0<v(x) \leqslant 1, x \in \Omega$ a.e. $\mu$. 令 $g(x)=\frac{1-v(x)}{v(x)}$,则 $g(x) \geqslant 0$, 且关于 $\mu$ 可测. 对 $E \in B$, 取 $u(x)=\frac{\chi_E(x)}{v(x)+\frac{1}{n}}$ 代入, 得
\begin{equation*}
\int_{\Omega} \chi_E(x) \frac{1-v(x)}{v(x)+\frac{1}{n}} \dif \mu=\int_{\Omega} \chi_E(x) \frac{v(x)}{v(x)+\frac{1}{n}} \dif \nu .
\end{equation*}
因为 $\nu$ 关于 $\mu$ 绝对连续, 且 $\nu>0$, a. e. $\mu$, 故 $\nu>0$, a. e. $\nu$. 令 $n \rightarrow \infty$ ，由单调收敛性定理得
\begin{equation*}
\int_E g(x) \dif \mu=\nu(E), \quad E \in B .
\end{equation*}
剩下考虑情形 $\mu(\Omega)=\infty$. 由 $\sigma$ 有限性, 取 $\Omega_n \subset \Omega_{n+1}, \Omega=$ $\bigcup_{n \geqslant 1} \Omega_n, \mu\left(\Omega_n\right)<\infty, n \geqslant 1$. 由先前结论, $E \subset \Omega$,
\begin{equation*}
\nu\left(E \cap \Omega_n\right)=\int_{E \cap \Omega_n} g_n \dif \mu,
\end{equation*}
易证: $g_n(x)=g_{n+1}(x), x \in \Omega_n$. 令 $g(x)=\lim _{n \rightarrow \infty} g_n(x)$, 由单调收敛性得
\begin{equation*}
\nu(E)=\lim _{n \rightarrow \infty} \int_{\Omega_n \cap E} \dif \nu=\lim _{n \rightarrow \infty} \int_{E \cap \Omega_n} g_n \dif \mu=\int_E g \dif \mu, \quad E \in B .
\end{equation*}
这样我们就证明了定理.

\end{proof}
\begin{problem}
设 $H$ 的元素是定义在集合 $S$ 上的复值函数。 又若 $\forall x$ $\in S$, 由
\begin{equation*}
J_x(f)=f(x) \quad(\forall f \in H),
\end{equation*}
定义的映射 $J_x: H \rightarrow \mathbb{C}$ 是 $H$ 上的线性连续泛函。求证：存在 $S \times$ $S$ 上的复值函数 $K(x, y)$ ，适合条件：\\
(1)对任意固定的 $y \in S$ ，作为 $x$ 的函数有 $K(x, y) \in H$ ;\\
(2) $f(y)=(f, K(\cdot, y))(\forall f \in H, \forall y \in S)$ 。
\end{problem}
\begin{proof}
由Riesz表示定理得，$\forall f\in H,\exists K_x\in H$s.t.$J_x(f)=f(x)=(f,K_x)$，令$K(x,y)=(K_y,K_x)$，则\[
K(x,y)=(K_y,K_x)=J_y(K_x)=K_x(y)
\]
所以条件一符合。\\
$\forall f \in H, \forall y \in S$ 
\[f(y)=J_y(f)=(f,K_y)=(f,K(\cdot,y))\]
所以条件二符合。
\end{proof}


\begin{problem}
求证: $H^2(D)$ 的再生核为
\begin{equation*}
K(z, w)=\frac{1}{\pi(1-z \bar{w})^2} \quad(z, w \in D) .
\end{equation*}

\end{problem}
\begin{solution}
由于$e_n(x)=\sqrt{\frac{n}{\pi}}x^{n-1}$为$H^2(D)$中的一组标准正交基，所以$\forall K_x,K_y\in H^2{(D)}$有
\[K_x=\sum\limits_{n=1}^\infty(K_x,e_n)e_n=\overline{e_n(x)}e_n,K(x,y)=\sum\limits_{n=1}^\infty\overline{e_n(y)}e_n(x)
\]
\[K(z,w)=\sum\limits_{n=1}^\infty(K_x,e_n)\sqrt{\frac{n}{\pi}}z^{n-1}\sqrt{\frac{n}{\pi}}\overline{w}^{n-1}=\frac{1}{\pi(1-z\overline{w})^2}\]
\end{solution}
\newpage
\section{开映射定理}
\begin{definition}
	设 $( {X} , \rho)$ 是一个度量空间，集合 $E \subset {X}$，称 $E$ 是疏集，如果 $\bar{E}$ 的内点是空的。
\end{definition}

\begin{example}
	在 $\mathbb{R}^n$ 上，有穷点集是疏集。Cantor 集是疏集。
\end{example}

\begin{proposition}
	设 $( {X} , \rho)$ 是一度量空间。为了 $E \subset {X}$ 是疏集，必须且仅须：对于任意球 $B\left(x_0, r_0\right)$，存在 $B\left(x_1, r_1\right) \subset B\left(x_0, r_0\right)$，使得
	\[
	\bar{E} \cap \bar{B}\left(x_1, r_1\right) = \varnothing.
	\]
\end{proposition}

\begin{proof}
	\textbf{必要性：} 由于 $\bar{E}$ 无内点，所以 $\bar{E}$ 不能包含任一球 $B\left(x_0, r_0\right)$。因此，存在 $x_1 \in B\left(x_0, r_0\right)$，使得 $x_1 \notin \bar{E}$。又由于 $\bar{E}$ 闭集，所以存在 $\varepsilon_1 > 0$，使得 $\bar{B}\left(x_1, \varepsilon_1\right) \cap \bar{E} = \varnothing$。取
	\[
	0 < r_1 < \min \left( \varepsilon_1, r_0 - \rho\left(x_0, x_1\right) \right)
	\]
	则有 $B\left(x_1, r_1\right) \subset B\left(x_0, r_0\right)$，且 $\bar{B}\left(x_1, r_1\right) \cap \bar{E} = \varnothing$。\\
	\textbf{充分性：} 若 $E$ 不疏，即 $\bar{E}$ 有内点，则存在 $B\left(x_0, r_0\right) \subset \bar{E}$。但由假设
	\[
	\exists B\left(x_1, r_1\right) \subset B\left(x_0, r_0\right), \text{ 使得 } \bar{B}\left(x_1, r_1\right) \cap \bar{E} = \varnothing
	\]
	一方面有 $B\left(x_1, r_1\right) \cap \bar{E} = B\left(x_1, r_1\right)$；另一方面有 $B\left(x_1, r_1\right) \cap \bar{E} = \varnothing$，即得矛盾。
\end{proof}
\begin{definition}
	在度量空间 $( {X} , \rho)$ 上，集合 $E$ 称为第一纲集，如果 $E = \bigcup_{n=1}^{\infty} E_n$，其中每个 $E_n$ 是疏集。不是第一纲的集合称为第二纲集。
\end{definition}

\begin{example}
	在 $\mathbb{R}$ 上，有理数集是第一纲集。更一般地，可数点集总是第一纲集。
\end{example}



\begin{theorem}[Baire]
	完备度量空间 $({X}, \rho)$ 是第二纲集。
\end{theorem}

\begin{proof}
	用反证法。假设 ${X}$ 是第一纲集，即存在疏集 $\{E_n\}$，使得
	\[
	{X} = \bigcup_{n=1}^{\infty} E_n.
	\]
	对任意的球 $B\left(x_0, r_0\right)$，存在 $B\left(x_1, r_1\right) \subset B\left(x_0, r_0\right)$，且 $r_1 < 1$，使得
	\[
	\bar{B}\left(x_1, r_1\right) \cap \bar{E}_1 = \varnothing.
	\]
	对球 $B\left(x_1, r_1\right)$，存在 $B\left(x_2, r_2\right) \subset B\left(x_1, r_1\right)$，且 $r_2 < \frac{1}{2}$，使得
	\[
	\bar{B}\left(x_2, r_2\right) \cap \left(\bar{E}_1 \cup \bar{E}_2\right) = \varnothing.
	\]
	如此继续下去，对球 $B\left(x_{n-1}, r_{n-1}\right)$，存在 $B\left(x_n, r_n\right) \subset B\left(x_{n-1}, r_{n-1}\right)$，且 $r_n < \frac{1}{n}$，使得
	\[
	\bar{B}\left(x_n, r_n\right) \cap \bar{E}_n = \varnothing,
	\]
	从而
	\[
	\bar{B}\left(x_n, r_n\right) \cap \left(\bigcup_{i=1}^n \bar{E}_i\right) = \varnothing \quad (\forall n \in \mathbb{N}).
	\]
	于是我们得到
	\[
	\bar{B}\left(x_0, r_0\right) \supset \bar{B}\left(x_1, r_1\right) \supset \cdots \supset \bar{B}\left(x_n, r_n\right) \supset \cdots,
	\]
	而
	\[
	\rho\left(x_{n+p}, x_n\right) \leqslant r_n < \frac{1}{n} \quad (\forall n, p \in \mathbb{N}).
	\]
	由此可见，$\{x_n\}$ 是基本列，从而存在 $x \in {X}$，使得 $\lim_{n \to \infty} x_n = x$。另一方面，在上述式子中令 $p \to \infty$ 得到 $\rho\left(x, x_n\right) \leqslant r_n$，从而
	\[
	x \in \bar{B}\left(x_n, r_n\right) \quad (\forall n \in \mathbb{N}).
	\]
	于是便有 $x \notin \bigcup_{n=1}^{\infty} E_n$，矛盾。
\end{proof}
\begin{theorem}[加强的 Baire 纲定理]完备的度量空间中,第一纲集的补集是稠密的第二纲集,可数个稠密开集的交是稠密的第二纲集.
\end{theorem}
\begin{proof}
	设 $X$ 是完备的度量空间， $\left\{E_n\right\}$ 是一列无处稠密集.
	上面实际上已证明了这一结果: 对任意的球 $B\left(x_0, r_0\right)$, 存在其中的点 $x \notin \bigcup_{n=1}^{\infty} E_n$. 故 $X \backslash \bigcup_{n=1}^{\infty} E_n$ 不止非空, 而且稠密.此外， $X \backslash \bigcup_{n=1}^{\infty} E_n$ 也是第二纲集，这是因为假如它是第一纲集， $X$ 就是两个第一纲集的并，故是第一纲集，与 Baire 纲定理矛盾. 前一结论成立。\\
	稠密的开集是无处稠密的闭集的补集，记为 $X \backslash E_n$ ，其交集是 $X \backslash \bigcup_{n=1}^{\infty} E_n$ ，故后一结论是其推论.
\end{proof}
\begin{theorem}
	在 $C[0,1]$ 中，处处不可微的函数集合 $E$ 是非空的，更确切地，$E$ 的余集是第一纲集。
\end{theorem}

\begin{proof}
	取 ${X} = C[0,1]$，设 $A_n$ 表示 ${X}$ 中这样一些元素 $f$ 的集合：对 $f$，存在 $s \in [0,1]$，使得对适合 $0 \leqslant s + h \leqslant 1$ 与 $|h| \leqslant \frac{1}{n}$ 的任何 $h$，成立
	\[
	\left|\frac{f(s+h) - f(s)}{h}\right| \leqslant n.
	\]
	若 $f$ 在某个点 $s$ 处可微，则必有正整数 $n$，使得 $f \in A_n$，于是
	\[
	{X} \setminus E \subset \bigcup_{n=1}^{\infty} A_n.
	\]
	下面我们证明每个 $A_n$ 是疏集。为此，先证 $A_n$ 是闭的。事实上，若 $f \in {X} \setminus A_n$，则 $\forall s \in [0,1]$，存在 $h_s$，使得
	\[
	\left|h_s\right| \leqslant \frac{1}{n}, \quad \text{且}\quad \left|f(s+h_s) - f(s)\right| > n |h_s|.
	\]
	又由 $f$ 的连续性，存在 $\varepsilon_s > 0$，以及 $s$ 的某个适当的邻域 $J_s$，使得对 $\forall \sigma \in J_s$，有
	\[
	\left|f(\sigma+h_s) - f(\sigma)\right| > n |h_s| + 2 \varepsilon_s.
	\]
	根据有限覆盖定理，可设 $J_{s_1}, J_{s_2}, \dots, J_{s_m}$ 覆盖 $[0,1]$，并设
	\[
	\varepsilon = \min \left\{ \varepsilon_{s_1}, \varepsilon_{s_2}, \dots, \varepsilon_{s_m} \right\}.
	\]
	今若 $g \in {X}$ 适合 $\|g - f\| < \varepsilon$，则由上述式子，对于 $\forall \sigma \in J_{s_k}(k = 1, 2, \dots, m)$，有
	\[
	\left|g(\sigma + h_{s_k}) - g(\sigma)\right| \geqslant \left|f(\sigma + h_{s_k}) - f(\sigma)\right| - 2 \varepsilon > n |h_{s_k}|.
	\]
	这证明了 ${X} \setminus A_n$ 是开集，从而 $A_n$ 是闭集。\\
	再证 $A_n$ 没有内点。对于任意 $f \in A_n$，任意 $\varepsilon > 0$，由 Weierstrass 逼近定理，存在多项式 $p$，使得
	\[
	\|f - p\| < \frac{\varepsilon}{2}.
	\]
	$p$ 的导数在 $[0,1]$ 上是有界的，因此根据中值定理，存在 $M > 0$，使得对 $\forall s \in [0,1]$ 及 $|h| < \frac{1}{n}$，成立
	\[
	|p(s+h) - p(s)| \leqslant M |h|.
	\]
	设 $g(s) \in C[0,1]$ 是一个分段线性函数，满足 $\|g\| < \varepsilon / 2$，并且各条线段斜率的绝对值都大于 $M+n$，那么
	\[
	p + g \in B(f, \varepsilon), \quad \text{而} \quad p + g \notin A_n.
	\]
	这样，我们证明了每个 $A_n$ 是疏集，从而 $\bigcup_{n=1}^{\infty} A_n$ 是第一纲集。而 ${X}$ 是完备的，由 Baire 定理${X}$ 是第二纲集，$E$ 是第一纲集的余集。
\end{proof}

\begin{theorem}[开映射定理]
	设 $X, Y$ 都是 $B$ 空间，若 $T \in {L}({X}, {Y})$ 是一个满射，则 $T$ 是开映射。
\end{theorem}

\begin{proof}
	用 $B(x_0, a)$ 和 $U(y_0, b)$ 分别表示 $X$ 和 $Y$ 中的开球。\\
	\textbf{(1) 证明 $T$ 是开映射：}\\
	为了证明 $T$ 是开映射，即 $\forall$ 开集 $W$，$T(W)$ 是开集，必须且仅须证明：$\exists \delta > 0$，使得
	\[
	T B(\theta, 1) \supset U(\theta, \delta).
	\]
	事实上，必要性是显然的。下证其充分性。由于 $T$ 的线性，上面的条件等价于
	\[
	T B(x_0, r) \supset U(T x_0, r \delta) \quad (\forall x_0 \in X, \forall r > 0).
	\]
	对于任意 $y_0 \in T(W)$，按定义，存在 $x_0 \in W$，使得 $y_0 = T x_0$。因为 $W$ 是开集，所以存在 $B(x_0, r) \subset W$。于是取 $\varepsilon = r \delta$，便有
	\[
	U(T x_0, \varepsilon) \subset T B(x_0, r) \subset T(W),
	\]
	即 $y_0 = T x_0$ 是 $T(W)$ 的内点。\\
	\textbf{(2) 证明：$\exists \delta > 0$，使得 $\overline{T B(\theta, 1)} \supset U(\theta, 3\delta)$：}\\
	这是因为
	\[
	Y = T X = \bigcup_{n=1}^{\infty} T B(\theta, n),
	\]
	而 $Y$ 是完备的，所以至少有一个 $n \in \mathbb{N}$，使得 $T B(\theta, n)$ 非疏集，即 $T B(\theta, n)$ 至少含有一个内点。因此，存在 $U(y_0, r) \subset \overline{T B(\theta, n)}$。注意到 $T B(\theta, n)$ 是一个对称凸集，便有
	\[
	U(-y_0, r) \subset \overline{T B(\theta, n)},
	\]
	从而
	\[
	U(\theta, r) \subset \frac{1}{2} U(y_0, r) + \frac{1}{2} U(-y_0, r) \subset \overline{T B(\theta, n)}.
	\]
	由 $T$ 的齐次性，取 $\delta = \frac{r}{3n}$，便有
	\[
	\overline{T B(\theta, 1)} \supset U(\theta, 3\delta).
	\]
	\textbf{(3) 证明：$T B(\theta, 1) \supset U(\theta, \delta)$：}\\
	对于任意 $y_0 \in U(\theta, \delta)$，要证 $\exists x_0 \in B(\theta, 1)$，使得 $T x_0 = y_0$，即求方程 $T x = y_0$ 在 $B(\theta, 1)$ 内的一个解 $x_0$。我们用逐次逼近法。\\
	对 $y_0 \in U(\theta, \delta)$，按 (2)，存在 $x_1 \in B(\theta, \frac{1}{3})$，使得
	\[
	\|y_0 - T x_1\| < \frac{\delta}{3}.
	\]
	对 $y_1 = y_0 - T x_1 \in U(\theta, \frac{\delta}{3})$，按 (2)，存在 $x_2 \in B(\theta, \frac{1}{3^2})$，使得
	\[
	\|y_1 - T x_2\| < \frac{\delta}{3^2}.
	\]
	对 $y_n = y_{n-1} - T x_n \in U(\theta, \frac{\delta}{3^n})$，按 (2)，存在 $x_{n+1} \in B(\theta, \frac{1}{3^{n+1}})$，使得
	\[
	\|y_n - T x_{n+1}\| < \frac{\delta}{3^{n+1}}.
	\]
	于是
	\[
	\sum_{n=1}^{\infty} \|x_n\| \leqslant \frac{1}{2},
	\]
	令 $x_0 \triangleq \sum_{n=1}^{\infty} x_n$，便有 $x_0 \in B(\theta, 1)$。而
	\[
	\begin{aligned}
		\|y_n\| &= \|y_{n-1} - T x_n\| = \cdots \\
		&= \|y_0 - T(x_1 + x_2 + \cdots + x_n)\| < \frac{\delta}{3^n} \quad (\forall n \in \mathbb{N}).
	\end{aligned}
	\]
	即得
	\[
	S_n \triangleq \sum_{i=1}^n x_i \to x_0, \quad T S_n \to y_0 \quad (n \to \infty).
	\]
	又因为 $T$ 是连续的，所以
	\[
	T x_0 = y_0,
	\]
	即得 $U(\theta, \delta) \subset T B(\theta, 1)$。
\end{proof}
\begin{theorem}[Banach]
	设 $X, Y$ 是 $B$ 空间。若 $T \in L(X, Y)$，它既是单射又是满射，那么 $T^{-1} \in L(Y, X)$。
\end{theorem}
\begin{proof}
定理 2.8的证明 依定理 2.7 证明中的第 (3) 部分, 已知
\begin{equation*}
	U(\theta, 1) \subset T B\left(\theta, \frac{1}{\delta}\right)
\end{equation*}即
\begin{equation*}
	T^{-1} U(\theta, 1) \subset B\left(\theta, \frac{1}{\delta}\right) \quad \text { 或 }\left\|T^{-1} y\right\|<\frac{1}{\delta} \quad(\forall y \in Y ,\|y\|<1) \text {. }
\end{equation*}
特别地，由范数的齐次性， $\forall y \in Y , \forall \varepsilon>0$ ，有
\begin{equation*}
	\left\|T^{-1} y\right\|<\frac{(1+\varepsilon)}{\delta}\|y\|
\end{equation*}
令 $\varepsilon \rightarrow 0$ 得
\begin{equation*}
	\left\|T^{-1} y\right\| \leqslant \frac{1}{\delta}\|y\| \quad(\forall y \in Y )
\end{equation*}
从而 $T^{-1} \in L ( Y , X )$.
\end{proof}


\begin{definition}[闭线性算子]
设 $T$ 是 $X \rightarrow Y$ 的线性算子，$D(T)$ 是其定义域。称 $T$ 是闭的，是指由 $x_n \in D(T), x_n \rightarrow x$，以及 $T x_n \rightarrow y$，就能推出 $x \in D(T)$，而且 $y = T x$。
\end{definition}

\begin{example}
在 $C[0,1]$ 上，$D(T) = C^1[0,1]$，$T = \frac{d}{dt}$ 是一个闭线性算子。
\end{example}

\begin{proof}如果 $x_n(t) \in C^1[0,1]$，并且有
\[
x_n \to x \quad \text{(在 $C[0,1]$ 中)}, \quad \frac{d x_n}{d t} \to y \quad \text{(在 $C[0,1]$ 中)},
\]
则有
\[
x_n(t) - x_n(0) \to \int_0^t y(\tau) d\tau \quad \forall t \in [0,1],
\]
并且
\[
x_n(t) - x_n(0) \to x(t) - x(0) \quad \forall t \in [0,1].
\]
因此，我们得到
\[
x(t) = x(0) + \int_0^t y(\tau) d\tau \quad \forall t \in [0,1].
\]
从而，$x \in C^1[0,1]$，并且 $\frac{d x}{d t} = y(t)$。
\end{proof}


\begin{theorem}[闭线性算子的定理]
若 $X$，$Y$ 是 $B$ 空间，$T$ 是 $X \to Y$ 的一个闭线性算子，满足 $R(T)$ 是 $Y$ 中的第二纲集，则 $R(T) = Y$ 并且 $\forall \varepsilon > 0, \exists \delta = \delta(\varepsilon) > 0$，使得 $\forall y \in Y, \|y\| < \delta$ 必有 $x \in D(T)$，且 $\|x\| < \varepsilon$，且 $y = T x$。
\end{theorem}
\begin{proof}
由与Banach定理的证明过程只需要用到闭线性算子的条件，所以我们已知对 $\varepsilon = 1$，存在 $\delta > 0$，使得
\[
U(\theta, \delta) \subset T(B(\theta, 1) \cap D(T)),
\]
其中 $U(\theta, \delta)$ 表示以 $\theta$ 为中心、半径为 $\delta$ 的邻域，$B(\theta, 1)$ 表示以 $\theta$ 为中心、半径为 1 的单位球，$D(T)$ 是算子 $T$ 的定义域。\\
\noindent 现在，考虑任意 $y \in Y$，不妨设 $y \neq \theta$（显然 $\theta \in R(T)$）。对于 $0 < \delta_1 < \delta$，根据上式，有
\[
\frac{\delta_1 y}{\|y\|} \in U(\theta, \delta) \Longrightarrow \frac{\delta_1 y}{\|y\|} \in T(B(\theta, 1) \cap D(T)).
\]
因此，存在某个 $x \in B(\theta, 1) \cap D(T)$，使得
\[
\frac{\delta_1 y}{\|y\|} = T x.
\]
由此我们得到
\[
y = T\left(\frac{\|y\|}{\delta_1} x\right),
\]
即
\[
y \in R(T).
\]
因此，$y$ 总是属于 $R(T)$，从而证明了 $R(T) = Y$。
\end{proof}
\begin{theorem}[线性算子的延拓]
设 $T$ 是 $B^*$ 空间 $X$ 到 $B$ 空间 $Y$ 的连续线性算子，那么 $T$ 能唯一地延拓到 $\overline{D(T)}$ 上成为连续线性算子 $T_1$，使得 $\left.T_1\right|_{D(T)} = T$，且 $\left\|T_1\right\| = \|T\|$。
\end{theorem}

\begin{proof}
任取 $x \in \overline{D(T)}$，则存在一个序列 $\{x_n\} \subset D(T)$，使得 $\lim\limits_{n \to \infty} x_n = x$。由于 $T$ 在 $D(T)$ 上是连续的，$T$ 是有界的，即存在常数 $M > 0$，使得
\[
\|T x\| \leqslant M \|x\| \quad (\forall x \in D(T)).
\]
因此，有
\[
\|T x_{n+p} - T x_n\| \leqslant M \|x_{n+p} - x_n\|.
\]
这表明 $\{T x_n\}$ 是 $Y$ 中的基本列。由于 $Y$ 完备，存在某个 $y \in Y$，使得 $T x_n \to y$。
不难看出，$y$ 仅依赖于 $x$，与 $D(T)$ 中的选择无关。因此，可以定义映射 $T_1: x \mapsto y$。我们可以验证以下几个性质：
\begin{itemize}
    \item 线性性：$T_1$ 是线性的。
    \item 一致性：$\left.T_1\right|_{D(T)} = T$。
    \item 有界性：对于任意 $x \in \overline{D(T)}$，有 $\|T_1 x\| \leqslant M \|x\|$。
\end{itemize}
因此，$T_1$ 是一个连续线性算子，且 $\|T_1\| = \|T\|$。唯一性由构造过程中的极限存在性和定义确定。

\end{proof}

\begin{corollary}[等价范数定理]
设线性空间 $X$ 上有两个范数 $\|\cdot\|_1$ 与 $\|\cdot\|_2$。如果 $X$ 关于这两个范数都构成 $B$ 空间，且 $\|\cdot\|_2$ 比 $\|\cdot\|_1$ 强，则 $\|\cdot\|_2$ 与 $\|\cdot\|_1$ 必等价。
\end{corollary}

\begin{proof}
考察恒同映射 $I: X \to X$，把它看成是由 $(X, \|\cdot\|_2)$ 映射到 $(X, \|\cdot\|_1)$ 的线性算子。由于假设 $\|\cdot\|_2$ 比 $\|\cdot\|_1$ 强，即存在常数 $C > 0$，使得
\[
\|I x\|_1 \leqslant C \|x\|_2 \quad (\forall x \in X).
\]
因此，$I$ 是连续的，它既是单射又是满射。依定理 2.8，$I$ 可逆且 $I^{-1}$ 连续，即存在常数 $M > 0$，使得
\[
\|I^{-1} x\|_2 \leqslant M \|x\|_1 \quad (\forall x \in X).
\]
又因为 $I^{-1} x$ 与 $x$ 是同一个元素，所以 $\|\cdot\|_1$ 与 $\|\cdot\|_2$ 等价。

\end{proof}
\begin{theorem}[闭图像定理]
设 $X, Y$ 是 $B$ 空间。若 $T$ 是 $X \to Y$ 的闭线性算子，并且 $D(T)$ 是闭的，则 $T$ 是连续的。
\end{theorem}

\begin{proof}
由于 $D(T)$ 是闭的，所以 $D(T)$ 作为 $X$ 的线性子空间可以看成是 $B$ 空间。在 $D(T)$ 上，引入另外一个范数 $\|\cdot\|_G$，定义如下：
\[
\|x\|_G = \|x\| + \|T x\| \quad (\forall x \in D(T)).
\]
现在证明 $\left(D(T), \|\cdot\|_G\right)$ 也是一个 $B$ 空间。事实上，考虑
\[
\|x_n - x_m\|_G = \|x_n - x_m\| + \|T x_n - T x_m\| \to 0 \quad \text{当} \ n, m \to \infty.
\]
由此可知，存在某个 $x^* \in X$ 和 $y^* \in Y$，使得 $x_n \to x^*$ 且 $T x_n \to y^*$。由于 $T$ 是闭算子，得出 $y^* = T x^*$，从而 $T x_n \to T x^*$。因此，有
\[
\|x_n - x^*\|_G \to 0.
\]
又显然，范数 $\|\cdot\|_G$ 比 $\|\cdot\|$ 强。根据等价范数定理，$\|\cdot\|_G$ 与 $\|\cdot\|$ 等价。因此，存在常数 $M > 0$，使得
\[
\|T x\| \leqslant \|x\|_G \leqslant M \|x\| \quad (\forall x \in D(T)).
\]
这表明，$T$ 是有界的，即 $T$ 是连续的。
\end{proof}

\begin{theorem}[共鸣定理或一致有界定理]
设 $X$ 是 $B$ 空间，$Y$ 是 $B^*$ 空间。如果 $W \subset L(X, Y)$，使得
\[
\sup_{A \in W} \|A x\| < \infty \quad (\forall x \in X),
\]
那么存在常数 $M$，使得 $\|A\| \leqslant M \quad (\forall A \in W)$。
\end{theorem}

\begin{proof}
对于任意 $x \in X$，定义新的范数
\[
\|x\|_W = \|x\| + \sup_{A \in W} \|A x\|.
\]
显然，$\|\cdot\|_W$ 是 $X$ 上的范数，且强于原范数 $\|\cdot\|$。\\
接下来证明 $\left(X, \|\cdot\|_W\right)$ 是完备的。事实上，如果
\[
\|x_m - x_n\| + \sup_{A \in W} \|A(x_m - x_n)\| \to 0 \quad \text{当} \ m, n \to \infty,
\]
由于 $X$ 是完备的，存在某个 $x \in X$，使得 $\|x_n - x\| \to 0$ （当 $n \to \infty$）。又因为对任意 $\varepsilon > 0$，存在某个 $N = N(\varepsilon)$，使得
\[
\sup_{A \in W} \|A x_m - A x_n\| < \varepsilon \quad (\forall m, n \geqslant N),
\]
所以对任意 $A \in W$ 有 $\|A x_n - A x\| \leqslant \varepsilon \quad (\forall n \geqslant N)$。于是
\[
\|x_n - x\| + \sup_{A \in W} \|A(x_n - x)\| \to 0 \quad \text{当} \ n \to \infty.
\]
即 $\|x_n - x\|_W \to 0$。这证明了 $\left(X, \|\cdot\|_W\right)$ 是完备的。\\
再根据等价范数定理，$\|\cdot\|_W$ 与原范数 $\|\cdot\|$ 等价，存在常数 $M > 0$，使得
\[
\sup_{A \in W} \|A x\| \leqslant M \|x\| \quad (\forall x \in X).
\]
因此，$\|A\| \leqslant M \quad (\forall A \in W)$，从而得证。
\end{proof}
\begin{note}
条件: $\forall x \in X , \sup _{A \in W}\|A x\|<\infty$, 意味着 $\forall x \in X , \exists M_x>0$,使得
\begin{equation*}
\|A x\| \leqslant M_x\|x\| \quad(\forall A \in W) .
\end{equation*}
而结论: $\|A\| \leqslant M(\forall A \in W)$, 则可看作是, 存在与 $x$ 无关的常数 $M$ ，使得
\begin{equation*}
\|A x\| \leqslant M\|x\| \quad(\forall A \in W) .
\end{equation*}
第一个式意味着算子族 $W$ 点点有界; 第二个式则意味着算子族 $W$ 一致有界. 因此本定理给出条件保证点点有界蕴含一致有界,故称 "一致有界" 定理. 另一方面，如果我们从反面来叙述本定理将有: $\sup _{A \in W}\|A\|=\infty \Longrightarrow \exists x_0 \in X$ ，使得
\begin{equation*}
\sup _{A \in W}\left\|A x_0\right\|=\infty
\end{equation*}
因此本定理又有 "共鸣定理"之称.
\end{note}

\begin{theorem}[Banach-Steinhaus 定理]
设 $X$ 是 $B$ 空间，$Y$ 是 $B^*$ 空间，$M$ 是 $X$ 的某个稠密子集。若 $A_n \in L(X, Y)$ 且 $A \in L(X, Y)$，则对 $\forall x \in X$ 都有
\[
\lim_{n \to \infty} A_n x = A x,
\]
的充要条件是：
\begin{itemize}
    \item[(1)] $\left\|A_n\right\|$ 有界；
    \item[(2)] 对 $\forall x \in M$ 成立 上 式。
\end{itemize}
\end{theorem}

\begin{proof}
\textbf{必要性：} 根据共鸣定理，结论显然成立。\\
\textbf{充分性：} 假定 $\left\|A_n\right\| \leqslant C \ (\forall n \in \mathbb{N})$，对 $\forall x \in X$ 和 $\forall \varepsilon > 0$，取 $y \in M$，使得
\[
\|x - y\| \leqslant \frac{\varepsilon}{4(\|A\| + C)}.
\]
则有
\[
\begin{aligned}
\|A_n x - A x\| &\leqslant \|A_n x - A_n y\| + \|A_n y - A y\| + \|A x - A y\| \\
&< \frac{\varepsilon}{2} + \|A_n y - A y\| \quad (\forall n \in \mathbb{N}).
\end{aligned}
\]
再取 $N$ 足够大，使得 $\|A_n y - A y\| < \frac{\varepsilon}{2} \ (\forall n \geqslant N)$，便有
\[
\|A_n x - A x\| < \varepsilon \quad (\forall n \geqslant N).
\]
因此，结论得证。
\end{proof}


\begin{theorem}[Lax-Milgram 定理]
设 $a(x, y)$ 是 Hilbert 空间 $X$ 上的一个共轭双线性函数，满足：
\begin{itemize}
    \item[(1)] $\exists M > 0$，使得 $|a(x, y)| \leqslant M \|x\| \cdot \|y\| \quad (\forall x, y \in X)$，
    \item[(2)] $\exists \delta > 0$，使得 $|a(x, x)| \geqslant \delta \|x\|^2 \quad (\forall x \in X)$，
\end{itemize}
那么必存在唯一的有连续逆的连续线性算子 $A \in L(X)$，满足
\[
\begin{aligned}
a(x, y) &= (x, A y) \quad (\forall x, y \in X), \\
\|A^{-1}\| &\leqslant \frac{1}{\delta}.
\end{aligned}
\]
\end{theorem}

\begin{proof}
依定理 2.4，适合上式的算子 $A \in L(X)$ 存在唯一。现证：\\
\textbf{(1) $A$ 是单射：} 若有 $y_1, y_2 \in X$，满足 $A y_1 = A y_2$，则
\[
a(x, y_1) = a(x, y_2) \quad (\forall x \in X),
\]
从而
\[
a(x, y_1 - y_2) = 0 \quad (\forall x \in X).
\]
特别取 $x = y_1 - y_2$，由上式得 $y_1 = y_2$。\\
\textbf{(2) $A$ 是满射：} 先证 $R(A)$ 是闭的。事实上，$\forall w \in \bar{R}(A)$，存在某个序列 $v_n \in X$（$n = 1, 2, \cdots$），使得
\[
w = \lim_{n \to \infty} A v_n.
\]
由条件得，
\[
\begin{aligned}
\delta \|v_{n+p} - v_n\|^2 &\leqslant |a(v_{n+p} - v_n, v_{n+p} - v_n)| \\
&= |(v_{n+p} - v_n, A(v_{n+p} - v_n))| \\
&\leqslant \|v_{n+p} - v_n\| \cdot \|A v_{n+p} - A v_n\| \quad (\forall n, p \in \mathbb{N}),
\end{aligned}
\]
从而
\[
\|v_{n+p} - v_n\| \leqslant \frac{1}{\delta} \|A v_{n+p} - A v_n\| \to 0 \quad (\text{当} \ n \to \infty, \forall p \in \mathbb{N}).
\]
因此，$\{v_n\}$ 是基本列，所以存在 $v^* \in X$，使得 $v_n \to v^*$，并由 $A$ 的连续性和上式得 $w = A v^*$，即 $w \in R(A)$。于是 $R(A)$ 闭。\\
再证 $R(A)^{\perp} = \{0\}$。假设 $w \in R(A)^{\perp}$，则
\[
(w, A v) = 0 \quad (\forall v \in X),
\]
即 $a(w, v) = 0 \quad (\forall v \in X)$。特别取 $v = w$，再利用条件有
\[
\delta \|w\|^2 \leqslant |a(w, w)| = 0,
\]
即得 $w = 0$。由此可见 $A$ 是满射。\\使用 Banach 逆算子定理： $A^{-1} \in L(X)$。因为
\[
\delta \|x\|^2 \leqslant |a(x, x)| = |(x, A x)| \leqslant \|x\| \cdot \|A x\|,
\]
所以 $\delta \|x\| \leqslant \|A x\| \quad (\forall x \in X)$，即得第二个结论。
\end{proof}
\begin{theorem}[Lax 等价定理]设 $T \in L ( X , Y )$, 其中 $X , Y$ 是 $B$空间. 给定 $y \in Y$ ，如果\begin{equation*}
\left\|T x-T_n x\right\| \rightarrow 0 \quad(n \rightarrow \infty)
\end{equation*}
对 $\forall x \in X$ 成立, 那么为了 $x_n \rightarrow x(n \rightarrow \infty)$, 其中 $x_n$ 与 $x$ 分别是\begin{equation*}
T_n x_n=y
\end{equation*}式与\begin{equation*}
T x=y .
\end{equation*}式的解, 必须且仅须 $\exists C>0$, 使得\begin{equation*}
\left\|T_n^{-1}\right\| \leqslant C \quad(\forall n \in N )
\end{equation*}式成立.
\end{theorem}

\begin{proof}
 充分性. 由条件, 我们得
\begin{equation*}
\begin{aligned}
\left\|x_n-x\right\| & =\left\|T_n^{-1} y-T_n^{-1} T_n x\right\| \\
& \leqslant\left\|T_n^{-1}\right\| \cdot\left\|T x-T_n x\right\| \\
& \leqslant C\left\|T x-T_n x\right\| \rightarrow 0 \quad(n \rightarrow \infty)
\end{aligned}
\end{equation*}
必要性. $\forall y \in Y$, 令 $x_n=T_n^{-1} y, x=T^{-1} y$, 便有 $x_n \rightarrow x(n \rightarrow \infty)$.因此,
\begin{equation*}
T_n^{-1} y \rightarrow T^{-1} y \quad(n \rightarrow \infty, \forall y \in Y )
\end{equation*}
由共鸣定理, 立得 $\left\|T_n^{-1}\right\|$ 有界.

\end{proof}

\begin{problem}
设 ${X}$ 是 $B$ 空间， ${X}_0$ 是 ${X}$ 的闭子空间。映射 $\varphi: {X} \rightarrow$ ${X} / { X }_0$ 。定义为
\begin{equation*}
\varphi: x \mapsto[x](\forall x \in {X} ),
\end{equation*}
其中 $[x]$ 表示含 $x$ 的商类。求证 $\varphi$ 是开映射。
\end{problem}
\begin{proof}
    $\forall [x]\in {X}/{X}_0,\exists x\in {X}$使得$\varphi(x)=[x] $，所以$\varphi$是满射，而${X},{X}/{X}_0$均是$B$空间，由开映射定理知$\varphi$是开映射。
\end{proof}
\begin{problem}
设 ${X} , {Y}$ 是 $B$ 空间，又设方程 $U x=y$ 对 $\forall y \in {Y}$ 有解 $x \in {X}$ ，其中 $U \in {L} ( {X}, {Y} )$ 。且 $\exists m>0$ ，使得
\begin{equation*}
\|U x\| \geqslant m\|x\| \quad(\forall x \in {X}) 
\end{equation*}
求证: $U$ 有连续逆 $U^{-1}$ ，并且 $\left\|U^{-1}\right\| \leqslant 1 / m$ 。
\end{problem}
\begin{proof}
    由题意知$U$是满射，设$x_1,x_2\in{X} ,x_1\neq x_2$，则
    \[\|Ux_1-Ux_2\|=\|U(x_1-x_2)\|\geqslant m\|x_1-x_2\|>0
    \]
    所以$U$是单射，由定理2.8知$U^{-1}\in {L}({Y},{X})$。于是$\forall\theta\neq y\in {Y},\exists\theta\neq x\in{X} $使得$Ux=y$，则
    \[\|U^{-1}\|=\sup\limits_{y\neq \theta }\frac{\|U^{-1}y\|}{\|y\|}=\sup\limits_{x\neq \theta }\frac{\|x\|}{\|Ux\|}\leqslant\sup\limits_{x\neq \theta }\frac{\|x\|}{m\|x\|}\leqslant \frac{1}{m}\]
    \end{proof}
\begin{problem}
 设 $H$ 是 Hilbert 空间， $A \in {L} (H)$ 并且 $\exists m>0$ ，使得
\begin{equation*}
|(A x, x)| \geqslant m\|x\|^2 \quad(\forall x \in H) .
\end{equation*}
求证 $\exists A^{-1} \in {L} (H)$ 。
\end{problem}

\begin{proof}
设 \( A : {H} \to {H} \) 是一个有界线性算子，满足条件：
\[
m\|x\|^2 \leqslant |(Ax, x)| \leqslant \|Ax\|\|x\|, \quad \forall x \in {H}.
\]
\textbf{1. \( A \) 是单射：}  
对任意 \( x \in {H} \)，由不等式得
\[
m\|x\|^2 \leqslant |(Ax, x)| \implies m\|x\|^2 \leqslant \|Ax\|\|x\| \implies \|Ax\| \geqslant m\|x\|.
\]
若 \( Ax = 0 \)，则 \( \|Ax\| = 0 \)，因此 \( \|x\| = 0 \)，即 \( x = \theta \)。  
故 \( A \) 是单射。\\
\textbf{2. \( A \) 是满射：}  
\begin{equation*}
\begin{aligned}
& 1^{\circ} \operatorname{Ran}(A)=\overline{\operatorname{Ran}(A)}: \\
& \qquad \begin{aligned}
y \in \overline{\operatorname{Im}(A)} & \Rightarrow \exists x_n \in H \text { s.t. } A x_n \rightarrow y \\
& \Rightarrow m\left\|x_n-x_m\right\| \leqslant\left\|A x_n-A x_m\right\| \rightarrow 0 \text { as } n, m \rightarrow \infty \\
& \Rightarrow \exists x \in H \text { s.t. } x_n \rightarrow x, A x_n \rightarrow A x \\
& \Rightarrow A x=y
\end{aligned} \\
& \begin{aligned}
2^{\circ} \operatorname{Ran}(A)^{\perp}=\{0\}, \text { 从而 } \operatorname{Im}(A)=\overline{\operatorname{Im}(A)}=H:
\end{aligned} \\
& \qquad \begin{aligned}
y \in \operatorname{Im}(A)^{\perp} & \Rightarrow\langle y, A x\rangle=0, \forall x \in H \\
& \Rightarrow 0=|\langle y, A y\rangle| \geqslant m\|y\|^2 \\
& \Rightarrow y=0
\end{aligned}
\end{aligned}
\end{equation*}
由正交分解定理, $1^{\circ}$ 表明 $H=\operatorname{Ran}(A) \oplus \operatorname{Ran}(A)^{\perp}$, 所以 $2^{\circ} \Rightarrow H=\operatorname{Ran}(A)$. 这个证明满射的思路在 Lax-Milgram定理的证明中也用到了。\\
\textbf{3. 结论：}  
根据 Banach 逆算子定理，\( A^{-1} \in L({H}) \)。
\end{proof}

\begin{problem}
设 ${X} ,{ Y}$ 是 $B^*$ 空间， $D$ 是 ${X}$ 的线性子空间且 $A$ ： $D \rightarrow {Y}$ 是线性映射。求证：\\
（1） 如果 $A$ 连续且 $D$ 是闭的，则 $A$ 是闭算子;\\
（2）如果 $A$ 连续且是闭算子，那末 ${Y}$ 完备蕴含 $D$ 闭；\\
（3）如果 $A$ 是单射的闭算子，则 $A^{-1}$ 也是闭算子；\\
（4）如果 ${X}$ 完备， $A$ 是单射的闭算子， $R(A)$ 在 ${Y}$ 中稠密且 $A^{-1}$ 连续, 那末 $R(A)= {Y}$ 。
\end{problem}
\begin{proof}
（1）设 $\left\{\begin{array}{l}x_n \in D(A), x_n \rightarrow x, \\ A x_n \rightarrow y .\end{array}\right.$ 因为 $D$ 闭, 所以 $x \in D(A)$; 又因
为 $A$ 连续, 所以 $y=A x$. 故 $A$ 是闭算子.\\
（2）因为 $A$ 连续，又 $ Y$ 完备，那么根据定理 2.13$, A$ 能唯一地延拓到 $\bar{D}$ 上成为连续线性算子 $\widetilde{A}$, 满足 $\left.\widetilde{A}\right|_D=A,\|\widetilde{A}\|=\|A\|$ 。到此, 本题还有一个条件 $A$ 是闭算子未用, 下面用它来证明 $D$ 闭. 设 $x_n \in D, x_n \rightarrow x$.要证明 $x \in D$ 。事实上, 现在有
\begin{equation*}
A x_n=\widetilde{A} x_n \rightarrow \widetilde{A} x, \quad n \rightarrow \infty .
\end{equation*}
因为 $A$ 是闭算子,所以
\begin{equation*}
\left.\begin{array}{l}
x_n \in D, x_n \rightarrow x \\
A x_n=\widetilde{A} x_n \rightarrow \widetilde{A} x
\end{array}\right\} \Rightarrow x \in D \text {, 且 } {A} x=\widetilde{A} x .
\end{equation*}
(3) 首先因为 $A$ 是单射, 所以 $A^{-1}$ 存在. 注意到 $D\left(A^{-1}\right)=R(A)$,也就是说, $A^{-1}$ 是定义在 $R(A)$ 上的线性算子。我们要证 $A^{-1}$ 是闭算子,就是要证，如果
\begin{equation*}
\left\{\begin{array}{l}
y_n \in R(A), y_n \rightarrow y, \\
x_n=A^{-1} y_n \rightarrow x,
\end{array}\right.
\end{equation*}
则有 $y \in R(A)$ 且 $x=A^{-1} y$. 事实上,
\begin{equation*}
\left\{\begin{array} { l } 
{ y _ { n } \in R ( A ) , y _ { n } \rightarrow y , } \\
{ x _ { n } = A ^ { - 1 } y _ { n } \rightarrow x }
\end{array} \Rightarrow \left\{\begin{array}{l}
x_n \in D(A), x_n \rightarrow x, \\
y_n=A x_n \rightarrow y .
\end{array}\right.\right.
\end{equation*}
因为 $A$ 是闭算子, 所以 $x \in D(A)$, 且 $y=A x$, 从而 $y \in R(A)$ 且 $x=$ $A^{-1} y$ 。\\
(4) 单先根据第 (3) 小题, 由条件 $A$ 是单射的闭算子, 即知 $A^{-1}$ 也是闭算子，
其次，注意到 $A^{-1}$ 的像空间是 $X$ ，根据第（2）小题，我们有
\begin{equation*}
\left.\begin{array}{l}
A^{-1} \text { 连续 } \\
A^{-1} \text { 是闭算子 } \\
X \text { 完备 }
\end{array}\right\} \Rightarrow R(A)=D\left(A^{-1}\right) \text { 闭. }
\end{equation*}
最后, 一方面, 因为 $R(A)$ 闭, 所以 $R(A)=\overline{R(A)}$; 另一方面, 因为 $R(A)$ 在 $Y$ 中稠密，所以 $\overline{R(A)}= Y$ 。两方面联合起来便有 $R(A)=$ $\overline{R(A)}= Y$ 。 也就是说，
$\left.\begin{array}{l}R(A) \text { 闭 } \\ R(A) \text { 在 } Y \text { 中稠密 }\end{array}\right\} \Rightarrow R(A)=\overline{R(A)}= Y$.
\end{proof}

\begin{problem}
用等价范数定理证明 $\left(C[0,1],\|\cdot\|_1\right)$ 不是 $B$ 空间,其中 $\|f\|_1=\int_0^1|f(t)| \dif t(\forall f \in C[0,1])$.
\end{problem}
\begin{proof}
假设 $\left(C[0,1],\|\cdot\|_1\right)$ 是 $B$ 空间，考虑是 $B$ 空间 $\left(C[0,1],\|\cdot\|_2\right)$ ，其中$\|f\|_2=\max\limits_{0\leqslant t\leqslant 1}|f(t)|(\forall f\in C[0,1]),\|f\|_1=\int_0^1|f(t)| \dif t\leqslant \max\limits_{0\leqslant t\leqslant 1}|f(t)|=\|f\|_2$。所以$\|\cdot\|_2$比$\|\cdot\|_1$强，由等价范数定理知$\|\cdot\|_2$与$\|\cdot\|_1$等价，则$\exists C>0$使得$\|f\|_2\leqslant C\|f\|_1$。下面考虑
\[
f(t)=\begin{cases}
    1-Ct&(0\leqslant t\leqslant \frac{1}{C})\\
    0&(\frac{1}{C}\leqslant t\leqslant1)
\end{cases}
\]
则$\|f\|_2=1\leqslant C\|f\|_1=\frac{1}{2}$矛盾，所以$\left(C[0,1],\|\cdot\|_1\right)$ 不是 $B$ 空间。
\end{proof}
\begin{problem}
设 ${X}$ 和 ${Y}$ 是 $B$ 空间，$A_n \in {L}({X}, {Y})$ (对于 $n=1, 2, \cdots$)，且对 $\forall x \in X$，序列 $\{A_n x\}$ 在 ${Y}$ 中收敛。证明：存在 $A \in {L}({X}, {Y})$，使得
\[
A_n x \rightarrow A x \quad (\forall x \in {X}) \text{，且 } \|A\| \leqslant \varliminf_{n \rightarrow \infty} \|A_n\|
\]
\end{problem}
\begin{proof}
 $\forall x \in {X}, \left\{A_n x\right\} $ 在 ${Y}$ 中收敛, $\therefore\left\{A_n x\right\}$ 在 ${Y}$中有界, 即
\begin{equation*}
\sup _{n \geqslant 1}\left\|A_n x\right\|<\infty \quad(\forall x \in {X} )
\end{equation*}
由共鸣定理知, $\exists M>0$, s. t. $\left\|A_n\right\| \leqslant M .(\forall n \geqslant 1)$。定义$Ax=\lim\limits_{n\to+\infty}A_nx$，则
\begin{equation*}
    \begin{aligned}
        A(\alpha x+\beta y)=\lim\limits_{n\to +\infty}A_n(\alpha x+\beta y)=\lim\limits_{n\to +\infty}(\alpha A_nx+\beta A_ny)=\alpha Ax+\beta Ay(\forall x,y\in{X},\alpha,\beta \in \mathbb{K})
    \end{aligned}
\end{equation*}
所以$A$是线性算子，而
\begin{equation*}
\begin{gathered}
\|A x\|=\lim_{n \rightarrow \infty}\left\|A_n x\right\| \leqslant \varliminf_{n \rightarrow \infty}\left\|A_n\right\|\|x\| \leqslant M\|x\|
\end{gathered}
\end{equation*}
所以$A\in {L}({X},{Y})\text{，且 } \|A\| \leqslant \varliminf\limits_{n \rightarrow \infty} \|A_n\|$。
\end{proof}
\begin{problem}设 $L$ 是 $B$ 空间， $p: Z \rightarrow R ^1$ 满足
(1) $p(x) \geqslant 0 \quad(\forall x \in Z )$;
(2) $p(\lambda x)=\lambda p(x)(\forall \lambda>0, \forall x \in Z )$;
(3) $p\left(x_1+x_2\right) \leqslant p\left(x_1\right)+p\left(x_2\right) \quad\left(\forall x_1, x_2 \in Z \right)$;
（4）当 $x_n \rightarrow x$ 时， $\lim _{n \rightarrow \infty} p\left(x_n\right) \geqslant p(x)$.
求证: $\exists M>0$, 使得 $p(x) \leqslant M\|x\| \quad(\forall x \in Z )$.
105
\end{problem}
\begin{problem}
设 $1 < p < \infty$ 且 $1/p + 1/q = 1$。如果序列 $\{\alpha_k\}$ 使得对 $\forall x = \{\xi_k\} \in l^p$ 保证 $\sum\limits_{k=1}^{\infty} \alpha_k \xi_k$ 收敛，证明 $\{\alpha_k\} \in l^q$。又若 $f: x \mapsto \sum\limits_{k=1}^{\infty} \alpha_k \xi_k$，证明：$f$ 作为 $l^p$ 上的线性泛函，有
\[
\|f\| = \left( \sum_{k=1}^{\infty} |\alpha_k|^q \right)^{\frac{1}{q}}.
\]
\end{problem}
\begin{proof}
     $\forall x=\left\{\xi_k\right\} \in l^p$, 令 $f(x)=\sum\limits_{k=1}^{\infty} \alpha_k \xi_k ;$ $f_n(x)=\sum\limits_{k=1}^n \alpha_k \xi_k$ 。下面证明$f_n \in {L} \left(l^p,\mathbb{ K} \right)$。由H\"older不等式得
\[|f_n(x)|=|\sum\limits_{k=1}^n \alpha_k \xi_k|\leqslant \sum\limits_{k=1}^n |\alpha_k \xi_k|\leqslant\left(\sum\limits_{k=1}^n \left|\alpha_k\right|^q \right)^{\frac{1}{q}}\left(\sum\limits_{k=1}^n \left|\xi_k\right|^p \right)^{\frac{1}{p}}\leqslant\left(\sum\limits_{k=1}^n \left|\alpha_k\right|^q \right)^{\frac{1}{q}}\|x\|_p\]
     所以$f_n \in {L} \left(l^p,\mathbb{ K} \right),\|f_n\|\leqslant\left(\sum\limits_{k=1}^n \left|\alpha_k\right|^q \right)^{\frac{1}{q}}$，取$x_n=(\xi^{(n)}_k)$如下：
    \[\xi_k^{(n)}=\begin{cases}
         \frac{|\alpha_k|^{q-1}\operatorname{sgn } \alpha_k}{\left(\sum\limits_{k=1}^n \left|\alpha_k\right|^q \right)^{\frac{1}{p}}}&k=1,2\cdots n\\
         0&k>n
     \end{cases}\]
则$x_n\in l^p,\|x_n\|_p=1,f_n(x_n)=\left(\sum\limits_{k=1}^n \left|\alpha_k\right|^q \right)^{\frac{1}{q}}$，所以$\|f_n\|=\left(\sum\limits_{k=1}^n \left|\alpha_k\right|^q \right)^{\frac{1}{q}}$， 因为 $\sum\limits_{k=1}^{\infty} \alpha_k \xi_k$ 收敛，所以$\sup\limits_{n\geqslant 1}|f_n(x)|<\infty(\forall x\in{X}
)$，由共鸣定理知
\[\sup\limits_{n\geqslant 1}\|f_n\|=\left(\sum\limits_{k=1}^\infty \left|\alpha_k\right|^q \right)^{\frac{1}{q}}<\infty\]
所以 $\{\alpha_k\} \in l^q$，由上题可知$f\in {L} \left(l^p,\mathbb{ K} \right),\|f\|\leqslant\varliminf\limits_{n\to \infty}\left(\sum\limits_{k=1}^n \left|\alpha_k\right|^q \right)^{\frac{1}{q}}=\left(\sum\limits_{k=1}^\infty \left|\alpha_k\right|^q \right)^{\frac{1}{q}} $，取$x=\{\frac{|\alpha_k|^{q-1}\operatorname{sgn}\alpha_k}{\left(\sum\limits_{k=1}^\infty \left|\alpha_k\right|^q \right)^{\frac{q-1}{q}}}\}$，则$\|x\|_p=1,|f(x)|=\left(\sum\limits_{k=1}^\infty \left|\alpha_k\right|^q \right)^{\frac{1}{q}}$，所以
有
\[
\|f\| = \left( \sum_{k=1}^{\infty} |\alpha_k|^q \right)^{\frac{1}{q}}.
\]
     
\end{proof}
\begin{problem}
如果序列 $\{a_k\}$ 使得对 $\forall x = \{\xi_k\} \in l^1$ 保证 $\sum\limits_{k=1}^{\infty} a_k \xi_k$ 收敛，证明 $\{a_k\} \in l^{\infty}$。若 $f: x \mapsto \sum\limits_{k=1}^{\infty} \alpha_k \xi_k$，作为 $l^1$ 上的线性泛函，证明：
\[
\|f\| = \sup_{k \geqslant 1} |\alpha_k|.
\]
\end{problem}
\begin{proof}
     $\forall x=\left\{\xi_k\right\} \in l^1$, 令 $f(x)=\sum\limits_{k=1}^{\infty} \alpha_k \xi_k ;$ $f_n(x)=\sum\limits_{k=1}^n \alpha_k \xi_k$ 。下面证明$f_n \in {L} \left(l^1,\mathbb{ K} \right)$。
\[|f_n(x)|=|\sum\limits_{k=1}^n \alpha_k \xi_k|\leqslant \sum\limits_{k=1}^n |\alpha_k \xi_k|\leqslant\sup\limits_{1\leqslant k\leqslant n}|\alpha_k|\left(\sum\limits_{k=1}^n \left|\xi_k\right| \right)\leqslant\sup\limits_{1\leqslant k\leqslant n}|\alpha_k|\|x\|_1\]
    所以$f_n \in {L} \left(l^1,\mathbb{ K} \right),\|f_n\|\leqslant\sup\limits_{1\leqslant k\leqslant n}|\alpha_k|$，设$\{|\alpha_k|\}_{k=1}^n$在$k=k_0$取到最大值，取$x_n=(\xi^{(n)}_k)$如下：
\[\xi_k^{(n)}=\begin{cases}
        \operatorname{sgn}\alpha_{k_0}&k=k_0\\
         0&k\neq k_0
     \end{cases}\]
则$x_n\in l^1,\|x_n\|_1=1,f_n(x_n)=\sup\limits_{1\leqslant k\leqslant n}|\alpha_k|$，所以$\|f_n\|=\sup\limits_{1\leqslant k\leqslant n}|\alpha_k|$， 因为 $\sum\limits_{k=1}^{\infty} \alpha_k \xi_k$ 收敛，所以$\sup\limits_{n\geqslant 1}|f_n(x)|<\infty(\forall x\in{X}
)$，由共鸣定理知
\[\sup\limits_{n\geqslant 1}\|f_n\|= \sup_{k \geqslant 1} |\alpha_k|<\infty\]
所以 $\{\alpha_k\} \in l^\infty$，由上上题可知$f\in {L} \left(l^1,\mathbb{ K} \right),\|f\|\leqslant\varliminf\limits_{n\to \infty}\sup\limits_{1\leqslant k\leqslant n}|\alpha_k|=\sup\limits_{1\leqslant k\leqslant n}|\alpha_k| $，设 $e _k=\{\overbrace{0, \cdots 0,1}^k,0,0 \cdots\} \in l^1$， 则 $\alpha_k=f\left( e _k\right)$ 。
\begin{equation*}
\left|\alpha_k\right|=\left|f\left( e _k\right)\right| \leqslant\|f\| \| e _k\|_1\Rightarrow |\alpha_k|\leqslant \|f\|(\forall k\geqslant 1)
\end{equation*}
所以$\|f\|\geqslant\sup\limits_{1\leqslant k}|\alpha_k| $。$
\|f\| = \sup\limits_{k \geqslant 1} |\alpha_k|
$得证。
     
\end{proof}

\begin{problem}
用 GelFand引理证明共鸣定理。
\end{problem}

\begin{problem}
设 ${X}, {Y}$ 是 $B$ 空间，$A \in {L}({X},{ Y})$ 是满射的。证明：如果在 ${Y}$ 中 $y_n \rightarrow y_0$，则存在 $C > 0$ 使得 $x_n \rightarrow x_0$，且满足 $A x_n = y_n$，且 $\|x_n\| \leqslant C \|y_n\|$。
\end{problem}
\begin{proof}
记 $N=\{x \in X: A x=0\}$, 则 $N$ 是 $X$ 的闭线性子空间, 且商空间 $X / N$ 是 Banach 空间。定义映射:
\begin{equation*}
\tilde{A}: X / N \rightarrow Y,[x] \mapsto A x
\end{equation*}
则\\
(1). $\tilde{A}$ 是单射: $\tilde{A}[x]=0 \Rightarrow A x=0 \Rightarrow x \in N \Rightarrow[x]=0$.\\
(2). $\tilde{A}$ 是满射: $y \in Y \Rightarrow \exists x \in X, A x=y \Rightarrow \tilde{A}[x]=A x=y$.\\
(3). $\tilde{A}$ 是有界线性映射： $A$ 是线性映射易知 $\tilde{A}$ 也线性. $\forall x^{\prime} \in[x]$,
\begin{equation*}
\|\tilde{A}[x]\|=\left\|A x^{\prime}\right\| \leqslant\|A\| \cdot\left\|x^{\prime}\right\|
\end{equation*}
取 $x^{\prime} \in[x]$ 满足 $\left\|x^{\prime}\right\| \leqslant 2\|[x]\|$, 则
\begin{equation*}
\|\tilde{A}[x]\| \leqslant 2\|A\| \cdot\|[x]\|
\end{equation*}
所以 $\tilde{A}$ 有界。\\
由逆算子定理可知 $\tilde{A}^{-1}$ 存在且有界, 不妨假设 $y_0=0, y_n \rightarrow 0$, 记 $\left[x_n\right]=\tilde{A}^{-1} y_n$, 则
\begin{equation*}
\left\|\left[x_n\right]\right\|=\left\|\tilde{A}^{-1} y_n\right\| \leqslant\left\|\tilde{A}^{-1} y_n\right\| \leqslant\left\|\tilde{A}^{-1}\right\| \cdot\left\|y_n\right\|
\end{equation*}
于是取 $x_n \in\left[x_n\right]$, 使得 $\left\|x_n\right\| \leqslant 2\left\|\left[x_n\right]\right\|$, 便有 $\left\|x_n\right\| \leqslant C\left\|y_n\right\|$, 其中 $C=2\left\|\tilde{A}^{-1}\right\|$.
\end{proof}
\begin{problem}
设 $P, Y$ 是 $B$ 空间，$T$ 是闭线性算子，$D(T) \subset R$，$R(T) \subset Y$，$N(T) \triangleq \{x \in X \mid T x = \theta\}$。
\begin{enumerate}
    \item 证明：$N(T)$ 是 $P$ 的闭线性子空间；
    \item 证明：$N(T) = \{\theta\}$ 和 $R(T)$ 在 $Y$ 中闭的充分且必要条件是 $\exists a > 0$，使得
    \[
    \|x\| \leqslant a\|T x\| \quad (\forall x \in D(T))；
    \]
    \item 如果用 $d(x, N(T))$ 表示点 $x \in Z$ 到集合 $N(T)$ 的距离 ($d(x, N(T)) = \inf_{x \in N(T)} \|z-x\|$)。求证：$R(T)$ 在 $Y$ 中闭的充分且必要条件是 $\exists a > 0$，使得
    \[
    d(x, N(T)) \leqslant a\|T x\| \quad (\forall x \in D(T)) .
    \]
\end{enumerate}
\end{problem}

\begin{problem}
设 $a(x, y)$ 是 Hilbert 空间 $H$ 上的一个共轭双线性泛函，满足：
\begin{enumerate}
    \item $\exists M > 0$，使得 $|a(x, y)| \leqslant M\|x\|\|y\| \quad (\forall x, y \in H)$；
    \item $\exists \delta > 0$，使得 $|a(x, x)| \geqslant \delta\|x\|^2 \quad (\forall x \in H)$。
\end{enumerate}
证明：$\forall f \in H^*$，存在唯一 $y_f \in H$，使得
\[
a\left(x, y_f\right) = f(x) \quad (\forall x \in H),
\]
而且 $y_f$ 连续地依赖于 $f$。
\end{problem}
\begin{proof}
 根据 Lax-Milgram 定理必存在唯一的有连续逆的连续线性算子 $A \in L ( H )$, s. t. $a(x, y)=(x, A y)$. 又根据Riesz表示定理，对 $\forall f \in H ^*, \quad \exists !$ $z_f \in H$, 使得 $f(x)=\left(x, z_f\right)$, 则
\begin{equation*}
A y=z_f \Rightarrow \exists ! y_f=A^{-1} z_f \Rightarrow f(x)=(x,y_f)
\end{equation*}
再证 $y_f$ 的唯一性。 设 $f(x)=a\left(x, y_f^{\prime}\right), \forall x \in H$, 则有
\begin{equation*}
0=a\left(x, y_f\right)-a\left(x, y_f^{\prime}\right)=a\left(x, y_f-y_f^{\prime}\right)
\end{equation*}
取 $x=y_f-y_f^{\prime}$, 有
\begin{equation*}
0=a\left(y_f-y_f^{\prime}, y_f-y_f^{\prime}\right){ }\geqslant \delta\left\|y_f-y_f^{\prime}\right\|^2 \Rightarrow y_f^{\prime}=y_f
\end{equation*}
\end{proof}
\begin{problem}
设 $\Omega$ 是 $\mathbb{R}^2$ 中边界光滑的有界开区域，$a: \Omega \rightarrow \mathbb{R}$ 有界可测且满足 $0 < \alpha_0 \leqslant \alpha$，$f \in {L}^2(\Omega)$。规定
\[
\begin{aligned}
& a(u, v) \triangleq \int_{\Omega} \left(\nabla u \cdot \nabla v + a u v\right) \dif x \dif y \quad (\forall u, v \in H^1(\Omega)); \\
& F(v) \triangleq \int_{\Omega} f v \, \dif x \, \dif y \quad (\forall v \in {L}^2(\Omega)).
\end{aligned}
\]
证明：存在 $u \in H^1(\Omega)$ 满足
\[
a(u, v) = F(v) \quad (\forall v \in H^1(\Omega)).
\]
\end{problem}
\section{Hahn-Banach 定理}
\begin{theorem}[实 Hahn-Banach 定理]设 \( X \) 是实线性空间，\( p \) 是定义在 \( X \) 上的次线性泛函，\( X_0 \subset X \) 是 \( X \) 的实线性子空间，\( f_0 \) 是定义在 \( X_0 \) 上的实线性泛函，满足：
\[
f_0(x) \leqslant p(x) \quad (\forall x \in X_0).
\]
那么在 \( X \) 上存在一个实线性泛函 \( f \)，满足以下两个条件：
\begin{enumerate}
    \item \( f(x) \leqslant p(x) \quad (\forall x \in X) \)（受 \( p \) 控制条件），
    \item \( f(x) = f_0(x) \quad (\forall x \in X_0) \)（延拓条件）。
\end{enumerate}

\end{theorem}

\begin{proof}
对于任意的 \( y_0 \in X \setminus X_0 \)，设  
\[
X_1 \triangleq \left\{x + \alpha y_0 \mid x \in X_0, \alpha \in \mathbb{R} \right\}.
\]
我们首先将 \( f_0 \) 延拓到 \( X_1 \)，设延拓后的线性泛函为 \( f_1 \)，满足：
\[
f_1\left(x + \alpha y_0\right) = f_0(x) + \alpha f_1(y_0) \quad (\forall x \in X_0, \forall \alpha \in \mathbb{R}).
\]
要求 \( f_1 \) 满足受 \( p \) 控制的条件：
\[
f_1\left(x + \alpha y_0\right) \leqslant p\left(x + \alpha y_0\right) \quad (\forall x \in X_0, \forall \alpha \in \mathbb{R}).
\]
除以 \( |\alpha| \) 并取极限可得：
\[
\begin{cases}
f_1\left(y_0 - z\right) \leqslant p\left(y_0 - z\right), & \forall z \in X_0, \\
f_1\left(-y_0 + y\right) \leqslant p\left(-y_0 + y\right), & \forall y \in X_0.
\end{cases}
\]
从上述不等式可得：
\[
\begin{gathered}
f_0(y) - p\left(-y_0 + y\right) \leqslant f_1\left(y_0\right) \leqslant f_0(z) + p\left(y_0 - z\right), (\forall y, z \in X_0).
\end{gathered}
\]
因此，为了能取到$f_1(y_0)$，必须且仅需：
\[
\sup_{y \in X_0} \left\{ f_0(y) - p\left(-y_0 + y\right) \right\} \leqslant \inf_{z \in X_0} \left\{ f_0(z) + p\left(y_0 - z\right) \right\}.
\]
这两个表达式总是可以保证成立。具体地，对于任意 \( y, z \in X_0 \)，有：
\[
f_0(y) - f_0(z) = f_0(y - z) \leqslant p(y - z) \leqslant p\left(y - y_0\right) + p\left(y_0 - z\right).
\]
从而，
\[
f_0(y) - p\left(-y_0 + y\right) \leqslant f_0(z) + p\left(y_0 - z\right) \quad (\forall y, z \in X_0).
\]
这表明我们可以取 \( f_1\left(y_0\right) \) 为上述两侧的中间值，从而构造 \( f_1 \)。
\\
接下来，我们逐步延拓 \( f_0 \) 到整个空间 \( X \)。为此，引入集合：
\begin{equation*}
F \triangleq\left\{\begin{array}{l|l}
\left( X _{\Delta}, f_{\Delta}\right) & \begin{array}{l} 
X _0 \subset X _{\Delta} \subset X ; \\
\forall x \in X _0 \Rightarrow f_{\Delta}(x)=f_0(x) ; \\
\forall x \in X _{\Delta} \Rightarrow f_{\Delta}(x) \leqslant p(x)
\end{array}
\end{array}\right\} .
\end{equation*}
在 $F$ 中引入序关系如下: $\left( X _{\Delta_1}, f_{\Delta_1}\right) \prec\left( X _{\Delta_2}, f_{\Delta_2}\right)$ 是指
\begin{equation*}
X _{\Delta_1} \subset X _{\Delta_2}, \quad \text { 且 } \quad f_{\Delta_1}(x)=f_{\Delta_2}(x) \quad\left(\forall x \in X _{\Delta_1}\right) .
\end{equation*}
于是 $F$ 成为半序集，又设 $M$ 是 $F$ 中的任一个全序子集，令
\begin{equation*}
X _M \triangleq \bigcup_{\left( X _{\Delta}, f_{\Delta}\right) \in M}\left\{ X _{\Delta}\right\},
\end{equation*}
及
\begin{equation*}
f_M(x)=f_{\Delta}(x) \quad\left(\forall x \in X _{\Delta},\left( X _{\Delta}, f_{\Delta}\right) \in M\right) .
\end{equation*}
由于 $M$ 是全序子集，容易验证 $X _M$ 是 $X$ 的包含 $X _0$ 的子空间，且 $f_M$ 在 $X _M$ 上是唯一确定的，满足 $f_M(x) \leqslant p(x)$ 。
根据 Zorn 引理, $F$ 存在极大元, 记为 $\left(X_A, f_A\right)$. 接下来证明 $X_A=X$. 用反证法. 若不然, 则可以选出 $y \in X \backslash X_A$. 最开始的证明告诉我们, 可以构造出
\begin{equation*}
X_B=\left\{x+\alpha y \mid x \in X_A, \alpha \in \mathbb{R} \right\}
\end{equation*}
并有对应的线性泛函 $f_B$ 使得 $\left(X_B, f_B\right) \in F$. 显然 $\left(X_A, f_A\right) \prec\left(X_B, f_B\right)$ ，且它们不相等。这与极大性矛盾. 从而 $X_A=X$ ，得到所欲证的.
\end{proof}
\begin{theorem}[复 Hahn-Banach 定理]设 $X$ 是复线性空间, $p$ 是 $X$ 上的半范数。 $X _0$ 是 $X$ 的线性子空间， $f_0$ 是 $X _0$ 上的线性泛函，并满足 $\left|f_0(x)\right| \leqslant p(x), \forall x \in X _0$ ，那么 $X$ 上必有一个线性泛函 $f$ 满足：\\
(1) $|f(x)| \leqslant p(x) \quad(\forall x \in X )$;\\
(2) $f(x)=f_0(x) \quad\left(\forall x \in X _0\right)$.
\end{theorem}
\begin{proof}
把 $X$ 看成实线性空间，相应把 $X _0$ 也看成是实线性子空间，令
\begin{equation*}
g_0(x) \triangleq \operatorname{Re} f_0(x) \quad\left(\forall x \in X _0\right),
\end{equation*}
便有 $g_0(x) \leqslant p(x)\left(\forall x \in X _0\right)$ 。从而根据实HB定理 ，必有 $X$ 上的实线性泛函 $g$ ，使得
\begin{equation*}
g(x)=g_0(x) \quad\left(\forall x \in X _0\right),
\end{equation*}
且
\begin{equation*}
g(x) \leqslant p(x) \quad(\forall x \in X )
\end{equation*}
现在，令
\begin{equation*}
f(x) \triangleq g(x)-\operatorname{i} g(\operatorname{i} x) \quad(\forall x \in X )
\end{equation*}
那么依上式, 我们有
\begin{equation*}
\begin{aligned}
f(x) & =g_0(x)-\operatorname{i} g_0(\operatorname{i} x) \\
& =\operatorname{Re} f_0(x)+\operatorname{iIm} f_0(x)=f_0(x) \quad\left(\forall x \in X _0\right)
\end{aligned}
\end{equation*}
又
\begin{equation*}
\begin{aligned}
f(\operatorname{i} x) & =g(\operatorname{i} x)-\operatorname{i} g(-x) \\
& =\operatorname{i}[g(x)-\operatorname{i} g(\operatorname{i} x)]=\operatorname{i} f(x) \quad(\forall x \in X ) .
\end{aligned}
\end{equation*}
从而 $f$ 也是复齐性的。剩下还要说明在 $X$ 上, $|f(x)|$ 受 $p(x)$ 控制. 若 $f(x)=0$ ，这是显然的。若 $f(x) \neq 0$ ，令
\begin{equation*}
\theta \triangleq \arg f(x),
\end{equation*}
那么依 上式, 有
\begin{equation*}
\begin{aligned}
|f(x)| & =e^{-i \theta} f(x)=f\left(e^{-i \theta} x\right) \\
& =g\left(e^{-i \theta} x\right) \leqslant p\left(e^{- \operatorname{i} \theta} x\right)=p(x) \quad(\forall x \in X ),
\end{aligned}
\end{equation*}
其中第三个等号是因为正数 $f\left( e ^{- \operatorname{i} \theta} x\right)=|f(x)|$ 的虚部为 0 .
\end{proof}
\begin{note}
    半范数大于等于0，$2p(x)=p(x)+p(-x)\geqslant p(0)=0$。
\end{note}


\begin{theorem}
    设 \( X \) 是 \( B^* \) 空间，\( X_0 \) 是 \( X \) 的线性子空间，\( f_0 \) 是定义在 \( X_0 \) 上的有界线性泛函，则在 \( X \) 上必有有界线性泛函 \( f \) 满足：
    \begin{enumerate}
        \item \( f(x) = f_0(x) \quad (\forall x \in X_0) \quad \) (延拓条件),
        \item \( \|f\| = \|f_0\|_0 \quad \) (保范条件),
    \end{enumerate}
    其中 \( \|f_0\|_0 \) 表示 \( f_0 \) 在 \( X_0 \) 上的范数。
\end{theorem}

\begin{proof}
    在 \( X \) 上定义半范数 \( p(x) \) 为：
    \[
    p(x) \triangleq \|f_0\|_0 \cdot \|x\|,
    \]
    那么 \( p(x) \) 是 \( X \) 上的半范数。根据HB定理，必存在 \( X \) 上的线性泛函 \( f(x) \)，满足：
    \[
    f(x) = f_0(x) \quad (\forall x \in X_0),
    \]
    以及
    \[
    |f(x)| \leq p(x) = \|f_0\|_0 \cdot \|x\| \quad (\forall x \in X).
    \]
    按泛函范数的定义，上式 蕴含 \( \|f\| \leq \|f_0\|_0 \)，又由$   f(x) = f_0(x) \quad (\forall x \in X_0),$显然有 \( \|f_0\|_0 \leq \|f\| \)。因此，得出结论：
    \[
    \|f\| = \|f_0\|_0.
    \]
\end{proof}
\begin{example} [ HBT 中延拓不唯一]
$X=\left( \mathbb{R} ^2,\|\cdot\|_1\right)$, 其中 $\left\|\left(x_1, x_2\right)\right\|_1 \stackrel{\text { def }}{=}\left|x_1\right|+\left|x_2\right|$, 取 $M= \mathbb{R} \times\{0\}, f: M \rightarrow \mathbb{R} ,(x, 0) \mapsto x$, 那么 $f$是 $M$ 上有界线性泛函, 且 $\|f\|=1$.令
\begin{equation*}
F_t: \mathbb{R} ^2 \rightarrow \mathbb{R} ,\left(x_1, x_2\right) \mapsto x_1+t x_2
\end{equation*}
那么 $\left.F_t\right|_M=f$, 而且对于 $\forall t \in(-1,1)$,
\begin{equation*}
\left|F_t\left(x_1, x_2\right)\right|=\left|x_1+t x_2\right| \leqslant\left|x_1\right|+|t|\left|x_2\right| \leqslant\left\|\left(x_1, x_2\right)\right\|_1 \Rightarrow\left\|F_t\right\| \leqslant 1
\end{equation*}
\end{example}


\begin{corollary}
    每个 \( B^* \) 空间必有足够多的连续线性泛函。
\end{corollary}

\begin{proof}
    任给 \( x_1, x_2 \in X \)，若 \( x_1 \neq x_2 \)，则 \( x_0 \triangleq x_1 - x_2 \neq \theta \)。令 \( X_0 \triangleq \left\{ \lambda x_0 \mid \lambda \in \mathbb{C} \right\} \)，并在 \( X_0 \) 上定义：
    \[
    f_0\left(\lambda x_0\right) = \lambda \|x_0\| \quad (\forall \lambda \in \mathbb{C}).
    \]
    那么 \( f_0(x_0) = \|x_0\| \) 且 \( \|f_0\|_0 = 1 \)。根据定理 2.18，存在 \( X \) 上的连续线性泛函 \( f \)，使得：
    \[
    f(x_0) = f_0(x_0) = \|x_0\|, \quad \|f\| = \|f_0\|_0 = 1.
    \]
    这个非零连续线性泛函 \( f \) 可以分辨 \( x_1 \) 和 \( x_2 \)。事实上：
    \[
    f(x_1) - f(x_2) = f(x_1 - x_2) = f(x_0) \neq 0.
    \]    
\end{proof}



\begin{corollary}
    设 \( X \) 是 \( B^* \) 空间，\( \forall x_0 \in X \setminus \{\theta\} \)，必存在 \( f \in X^* \)，使得：
\[
    f(x_0) = \|x_0\|, \quad \text{且} \quad \|f\| = 1.
    \]
\end{corollary}

\begin{note}
    本推论给出了判别 \( B^* \) 空间零元的一种方法：为了 \( x_0 = \theta \)，必须且仅须 \( \forall f \in X^* \)，蕴含 \( f(x_0) = 0 \)。
\end{note}
\begin{theorem}
    设 \( X \) 是 \( B^* \) 空间，\( M \) 是 \( X \) 的线性子空间。若 \( x_0 \in X \)，且
    \[
    d \triangleq \rho(x_0, M) > 0,
    \]
    则必存在 \( f \in X^* \) 适合下列条件：
    \begin{enumerate}
        \item \( f(x) = 0 \quad (\forall x \in M) \);
        \item \( f(x_0) = d \);
        \item \( \|f\| = 1 \).
    \end{enumerate}
\end{theorem}

\begin{proof}
    考虑 \( X_0 \triangleq \{ x = x' + \alpha x_0 \mid x' \in M, \alpha \in \mathbb{K} \} \)，对于任意 \( x \in X_0 \)，定义
    \[
    f_0(x) = \alpha d.
    \]
    显然，\( f_0 \) 满足条件 (1) 和 (2)。又若 \( x = x' + \alpha x_0 \) （其中 \( x' \in M \), \( \alpha \neq 0 \)），则有：
    \[
    \begin{aligned}
    |f_0(x)| &= |\alpha| d = |\alpha| \rho(x_0, M) \\
    &\leqslant |\alpha| \left\|\frac{x'}{\alpha} + x_0 \right\| \\
    &= \|x' + \alpha x_0\| = \|x\|.
    \end{aligned}
    \]
    因此，\( \|f_0\| \leqslant 1 \)。根据 Hahn-Banach 定理，将 \( f_0 \) 保范延拓为 \( f \in X^* \)，则有 \( f \) 满足条件 (1)、(2)，且 \( \|f\| \leqslant 1 \)。只需证明 $\|f\| \geqslant 1$.
\begin{equation*}
\begin{aligned}
d=\inf _{y \in M}\left\|x_0-y\right\| & \Rightarrow \forall n, \exists y_n \in M \text { s.t. }\left\|x_0-y_n\right\|<d+\frac{1}{n} \\
& \Rightarrow \frac{\left|f\left(x_0-y_n\right)\right|}{\left\|x_0-y_n\right\|}=\frac{\left|f\left(x_0\right)\right|}{\left\|x_0-y_n\right\|}>\frac{d}{d+\frac{1}{n}} \rightarrow 1 \text { as } n \rightarrow \infty \\
& \Rightarrow \sup _n \frac{\left|f\left(x_0-y_n\right)\right|}{\left\|x_0-y_n\right\|} \geqslant 1 \\
& \Rightarrow\|f\| \geqslant 1
\end{aligned}
\end{equation*}
\end{proof}

\begin{corollary}
    设 \( M \) 是 \( B^* \) 空间 \( X \) 的一个子集，且 \( x_0 \) 是 \( X \) 中的任一个非零元素。那么
    \[
    x_0 \in \overline{\operatorname{span} M}
    \]
    的充要条件是：对所有 \( f \in X^* \)，
    \[
    f(x) = 0 \quad (\forall x \in M) \quad \Longrightarrow \quad f(x_0) = 0.
    \]
\end{corollary}

\begin{proof}
    必要性显然。下面我们用反证法证明充分性。假设 \( x_0 \notin \overline{\operatorname{span} M} \)，那么
    \[
    d \triangleq \rho(x_0, \overline{\operatorname{span} M}) > 0.
    \]
    因此，依定理 2.19，存在 \( f \in X^* \)，使得 \( f(x) = 0 \) 对所有 \( x \in \overline{\operatorname{span} M} \)，并且 \( f(x_0) = d > 0 \)。但根据假设，对于此 \( f \)，应有 \( f(x_0) = 0 \)，这与 \( f(x_0) = d > 0 \) 矛盾。因此，结论成立。
\end{proof}
\begin{definition}
    在线性空间 \( X \) 中，\( X \) 的线性子空间 \( M \) 称为极大的，如果对于任何一个以 \( M \) 为真子集的线性子空间 \( M_1 \)，必有 \( M_1 = X \)。
\end{definition}



\begin{proposition}
    \( M \) 是极大线性子空间的充要条件是，\( M \) 是线性真子空间，并且对所有 \( x_0 \in X \setminus M \)，有
    \[
    X = \left\{ \lambda x_0 \mid \lambda \in \mathbb{R} \right\} \oplus M.
    \]
\end{proposition}

\begin{proof}
    必要性显然。为了证充分性，设 \( M_1 \) 是以 \( M \) 为真子集的线性子空间，那么必存在 \( x_0 \in M_1 \setminus M \)。于是有
    \[
    \lambda x_0 \in M_1 \quad (\forall \lambda \in \mathbb{R}),
    \]
    并且 \( M \subset M_1 \)，从而
    \[
    X = \left\{ \lambda x_0 \mid \lambda \in \mathbb{R} \right\} \oplus M \subset M_1.
    \]
    即得 \( X = M_1 \)。于是 \( M \) 是极大线性子空间。
\end{proof}



\begin{definition}
    \( X \) 的极大线性子空间 \( M \) 对向量 \( x_0 \in X \) 的平移
    \[
    L \triangleq x_0 + M
    \]
    称为极大线性流形，或简称超平面。
\end{definition}



\begin{theorem}
    为了 \( L \) 是线性 \( (B^*) \) 空间 \( X \) 上的一个（闭）超平面，必须且仅须存在非零（连续）线性泛函 \( f \) 及 \( r \in \mathbb{R} \)，使得 \( L = H_f^r \)。
\end{theorem}

\begin{proof}
 充分性：注意到 $H_f^0=\operatorname{Ker}(f)$ ，下证$H_f^0$ 是极大子空间。取 $x_0 \in X \backslash H_f^0$,
\begin{equation*}
\begin{aligned}
f\left(x-\frac{f(x)}{f\left(x_0\right)} x_0\right)=0, \forall x \in X & \Rightarrow f-\frac{f(x)}{f\left(x_0\right)} x_0 \in H_f^0, \forall x \in X \\
& \Rightarrow X=H_f^0 \oplus \operatorname{span}\left\{x_0\right\}
\end{aligned}
\end{equation*}
令 $r \stackrel{\text { def }}{=} f\left(x_0\right)$, 则
\begin{equation*}
x \in H_f^r \Leftrightarrow f\left(x-x_0\right)=f(x)-f\left(x_0\right)=0 \Leftrightarrow x-x_0 \in H_f^0 \Leftrightarrow x \in H_f^0+x_0
\end{equation*}
所以 $H_f^r$ 是极大子空间 $H_f^0$ 的平移, 为超平面，又若 \( f \) 是连续的，则 \( H_f^r \) 显然是闭的。\\
必要性: 设 $L=M+x_0, M$ 是极大(闭)子空间, 那么 $X=M \oplus \operatorname{span}\left\{x_0\right\}$, 令
\begin{equation*}
f: X \rightarrow \mathbb{R} , x=y+\lambda x_0 \mapsto \lambda
\end{equation*}
于是 $f(M)=\{0\}$ 且 $f\left(x_0\right)=1$, 进而 $M \subset H_f^0$, 由 $M$ 极大 $\Rightarrow M=H_f^0$, 因此 $L=H_f^r$ with $r=1$，由问题2.4知由$H_f^0$闭推出$f\in X^*$。
\end{proof}


\begin{definition}
所谓超平面 $L=H_f^r$ 分离集合 $E$ 与 $F$, 是指:
\begin{equation*}
\begin{aligned}
& \forall x \in E \Longrightarrow f(x) \leqslant r(\text { 或 } \geqslant r), \\
& \forall x \in F \Longrightarrow f(x) \geqslant r(\text { 或 } \leqslant r) .
\end{aligned}
\end{equation*}
如果在上面两个式子中，用 "<"与 ">" 分别代替 " $\leqslant$ " 与 " $\geqslant$"，那么就说 $H_f^r$ 严格分离 $E$ 与 $F$.
\end{definition}
\begin{theorem}[Hanh-Banach定理的几何形式]
$X$ 是实赋范空间, $C$ 是有内点的凸集, $x_0 \notin C \Rightarrow \exists f \in X^*, \exists r \in \mathbb{R}$ 满足 $H_f^r$ 分离 $x_0$ 和 $C$.
\end{theorem}
\begin{proof}
不妨设$\theta$是 $C$ 的内点, $P_C$ 是 $C$ 的 Minkowski 泛函, 由命题1.14, $P_C$ 是次线性泛函, 并且
\begin{equation*}
\bar{C}=\left\{x \in X: P_C(x) \leqslant 1\right\}
\end{equation*}
$x_0 \notin C \Rightarrow P_C\left(x_0\right) \geqslant 1,0$ 是 $C$ 的内点 $\Rightarrow \exists \varepsilon>0$ s.t. $B(0, \varepsilon) \subset C$, 于是
\begin{equation*}
\forall x \in X, x \neq 0, \varepsilon \frac{x}{\|x\|} \in \overline{B(0, \varepsilon)} \subset \bar{C}
\end{equation*}
进而
\begin{equation*}
P_C\left(\varepsilon \frac{x}{\|x\|}\right) \leqslant 1, \forall 0 \neq x \in X
\end{equation*}
也就是
\begin{equation*}
P_C(x) \leqslant \frac{1}{\varepsilon}\|x\|, \forall x \in X
\end{equation*}
令 $M=\operatorname{span}\left\{x_0\right\}$ ，定义
\begin{equation*}
f_0: M \rightarrow \mathbb{R} , x=\lambda x_0 \mapsto \lambda P_C\left(x_0\right)
\end{equation*}
则 $f_0(x) \leqslant P_C(x), \forall x \in M$, 由实 HBT 可得存在 $f: X \rightarrow \mathbb{R}$ 满足 $\left.f\right|_M=f_0$ 且 $f(x) \leqslant P_C(x), \forall x \in X$, 进而
\begin{equation*}
f\left(x_0\right)=f_0\left(x_0\right)=P_C\left(x_0\right) \geqslant 1, f(x) \leqslant P_C(x) \leqslant 1, \forall x \in C
\end{equation*}
因此 $H_f^1$ 分离 $x_0$ 和 $C$.
只剩下证明 $f \in X^*$,
\begin{equation*}
\begin{aligned}
f(x) \leqslant P_C(x) \leqslant \frac{1}{\varepsilon}\|x\|, \forall x \in X & \Rightarrow f(-x)=-f(x) \leqslant \frac{1}{\varepsilon}\|x\| \\
& \Rightarrow|f(x)| \leqslant \frac{1}{\varepsilon}\|x\|, \forall x \in X \\
& \Rightarrow f \in X^*
\end{aligned}
\end{equation*}

\end{proof}
\begin{theorem}[Hahn-Banach 凸集分离定理1]
$X$ 是实赋范空间, $A$ 是有内点的凸集, $B$ 是凸集, 若 $A \cap B=\varnothing$, 则存在 $H_f^r$ 闭, 并分离 $A, B$.
\end{theorem}
\begin{proof}
令 $C=A-B$, 即
\begin{equation*}
C=\{x-y: x \in A, y \in B\}=\bigcup_{y \in B}(A-y)
\end{equation*}
那么 $C$ 是含内点凸集，而且 $\theta\notin C, 由$ 定理 2.20， $H_f^0$ 分离 $C$ 和 $\{\theta\}$ ，即存在 $f \in X^*$ 使得
\begin{equation*}
\sup _{\delta \in C} f(\delta) \leqslant 0=f(0)
\end{equation*}
而
\begin{equation*}
\sup _{\delta \in C} f(\delta)=\sup _{x \in A, y \in B}[f(x)-f(y)]=\sup _{x \in A} f(x)-\inf _{y \in B} f(y)
\end{equation*}
因此
\begin{equation*}
\sup _{x \in A} f(x) \leqslant r \leqslant \inf _{y \in B} f(y)
\end{equation*}
这里
\begin{equation*}
r \stackrel{\text { def }}{=} \frac{1}{2}\left[\sup _{x \in A} f(x)+\inf _{y \in B} f(y)\right]
\end{equation*}
\end{proof}
\begin{theorem}[Hahn-Banach 凸集分离定理 2]
$X$ 是实赋范空间, $A$ 是闭凸集, $B$ 是紧凸集, 若 $A \cap B=\varnothing$, 则存在 $H_f^r$ 闭, 并严格分离 $A, B$.
\end{theorem}
\begin{proof}
 $A$ 闭 $B$ 紧且不交, 所以 $\operatorname{dist}(A, B)>0$,否则$A,B$相交， 令 $\varepsilon=\frac{1}{4} \operatorname{dist}(A, B)$,
\begin{equation*}
A_{\varepsilon} \stackrel{\text { def }}{=} A+B(0, \varepsilon), B_{\varepsilon} \stackrel{\text { def }}{=} B+B(0, \varepsilon)
\end{equation*}
这两个都是开凸集且不交，由定理1.4.7可知 $\exists f \in X^*, \exists r \in R$ 满足
\begin{equation*}
\sup _{x \in A_e} f(x) \leqslant r \leqslant \inf _{y \in B_e} f(y)
\end{equation*}
所以
\begin{equation*}
\begin{aligned}
f(x+\varepsilon \delta) & \leqslant r \leqslant f(y+\varepsilon \delta), \forall x \in A, \forall y \in B, \forall \delta \in B(0,1) \\
& \Rightarrow-f(\delta) \leqslant \frac{f(y)-r}{\varepsilon} \\
& \Rightarrow\|f\|=\sup _{\delta \in B(0,1)} f(-\delta) \leqslant \frac{f(y)-r}{\varepsilon} \\
& \Rightarrow r \leqslant f(y)-\varepsilon\|f\|, \forall y \in B \\
& \Rightarrow r \leqslant \inf _{y \in B} f(y)-\varepsilon\|f\|<\inf _{y \in B} f(y)
\end{aligned}
\end{equation*}
同理，
\begin{equation*}
\sup _{x \in A} f(x)<\sup _{x \in A} f(x)+\varepsilon\|f\| \leqslant r
\end{equation*}
\end{proof}
\begin{corollary}[Ascoli 定理] 设 $E$ 是实 $B^*$ 空间 $X$ 中的闭凸集，则 $\forall x_0 \in X \backslash E, \exists f \in X ^*$ 及 $\alpha \in R$ ，适合
\begin{equation*}
f(x)<\alpha<f\left(x_0\right) \quad(\forall x \in E) .
\end{equation*}
\end{corollary}
\begin{proof}
因为 $x_0 \in X \backslash E$ 及 $E$ 是闭集，所以 $\exists \delta>0$ ，使得
\begin{equation*}
B\left(x_0, \delta\right) \subset X \backslash E,
\end{equation*}
而 $B\left(x_0, \delta\right)$ 是有内点的凸集。对 $E$ 和 $B\left(x_0, \delta\right)$ 应用定理 1.15,存在非零连续线性泛函 $f$ ，适合
\begin{equation*}
\sup _{x \in E} f(x) \leqslant \inf _{y \in B\left(x_0, \delta\right)} f(y) .
\end{equation*}
进一步可以证明
\begin{equation*}
\inf _{y \in B\left(x_0, \delta\right)} f(y)<f\left(x_0\right)
\end{equation*}
事实上，倘若上式不成立，那么
\begin{equation*}
f(y) \geqslant f\left(x_0\right) \quad\left(\forall y \in B\left(x_0, \delta\right)\right) .
\end{equation*}
这表明 $f\left(x_0\right)$ 是 $f(y)$ 在 $B\left(x_0, \delta\right)$ 中的极小值, 这等价于$0$是$f(x)$在$B\left(\theta, \delta\right)$ 中的极小值，那么$f(x),f(-x)>0\Rightarrow f(\theta)>0(\forall x\in B\left(\theta, \delta\right))$，矛盾，所以
\begin{equation*}
\sup _{x \in E} f(x) <\alpha< f(x_0).
\end{equation*}
其中\begin{equation*}
r \stackrel{\text { def }}{=} \frac{1}{2}\left[\sup _{x \in E} f(x)+f(x_0)\right]
\end{equation*}

\end{proof}
\begin{corollary}[Mazur 定理]
$X$ 是实赋范空间， $C$ 是开凸集， $F$ 是线性子流形（子空间的平移），若 $C \cap F=\varnothing$ ，则存在 $H_f^r$ 闭满足 $F \subset H_f^r$ 且 $\sup _{x \in C} f(x) \leqslant r$ 。


\end{corollary}

\begin{proof}
设 $F=M+x_0, M$ 是子空间, 由分离定理
\begin{equation*}
\exists f \in X^*, \exists s \in \mathbb{R} \text { s.t. } \sup _{x \in C} f(x) \leqslant s \leqslant \inf _{y \in F} f(y)=\inf _{\delta \in M} f(\delta)+f\left(x_0\right)
\end{equation*}
进而
\begin{equation*}
\begin{aligned}
\inf _{\delta \in M} f(\delta) \geqslant s-f\left(x_0\right) & \Rightarrow f|_M=0 (\text{因为M是子空间})\\
& \Rightarrow M \subset H_f^0 \\
& \Rightarrow F \subset H_f^r \text { with } r=f\left(x_0\right)
\end{aligned}
\end{equation*}
同时,
\begin{equation*}
\sup _{x \in C} f(x) \leqslant s \leqslant f\left(x_0\right)=r
\end{equation*}
\end{proof}
\begin{note}
上述结论换句话说就是：存在 $X$ 上的非零连续线性泛函 $f$ 及 $s \in \mathbb{R}$ ，使得
\begin{equation*}
f(x) \leqslant s \quad(\forall x \in E), \quad f(x)=s \quad(\forall x \in F) .
\end{equation*}
\end{note}

\begin{definition}
超平面 $L=H_f^r$ 称为凸集 $E$ 在点 $x_0$ 的承托超平面，是指 $E$ 在 $L$ 的一侧，且 $\bar{E}$ 与 $L$ 有公共点 $x_0$ 。换句话说，
\begin{equation*}
f(x) \leqslant r=f\left(x_0\right) \quad(\forall x \in E),
\end{equation*}
或
\begin{equation*}
f(x) \geqslant r=f\left(x_0\right) \quad(\forall x \in E) .
\end{equation*}
\end{definition}
\begin{theorem}
$X$ 是实赋范空间， $C$ 是有内点的闭凸集， $\forall x_0 \in \partial C$ 均有 $C$ 的一个承托超平面。
\end{theorem}
\begin{proof}
令 $E=C^{\circ}$, 即 $C$ 的全体内点, $F=\left\{x_0\right\}$, 由 Mazur定理 可得
\begin{equation*}
\exists f \in X^*, \exists r \in \mathbb{R} \text { s.t. } \sup _{x \in E} f(x) \leqslant r \text { and }\left\{x_0\right\} \subset H_f^r
\end{equation*}
由连续性
\begin{equation*}
\sup _{x \in C} f(x) \leqslant r=f\left(x_0\right)
\end{equation*}
\end{proof}

\begin{example}
 设 $X$ 是 $B^*$ 空间， $E=\{x \in X \mid\|x\| \leqslant r\},\left\|x_0\right\|=r$ ，那么 $E$ 在 $x_0$ 有一个承托超平面.


\end{example}

\begin{proof}
根据推论 2.3, $\exists f \in X ^*$ ，使得 $f\left(x_0\right)=\left\|x_0\right\|,\|f\|=1$ 。于是 $H_f^r$ 便是 $E$ 在 $x_0$ 的承托超平面，这是因为
\begin{equation*}
f(x) \leqslant\|f\| \cdot\|x\| \leqslant r=f\left(x_0\right) \quad(\forall x \in E) .
\end{equation*}

\end{proof}
\begin{example}

设 $\sum_{k=1}^{\infty} x_k$ 是 Banach 空间 $X$ 中绝对收敛级数, $\left\{y_k\right\}_{k=1}^{\infty}$ 是 $\left\{x_k\right\}_{k=1}^{\infty}$ 的任一重排, 则
\begin{equation*}
\sum_{k=1}^{\infty} y_k=\sum_{k=1}^{\infty} x_k
\end{equation*}



\end{example}
\begin{proof}
 $\forall f \in X^*$,
\begin{equation*}
\sum_{k=1}^{\infty}\left|f\left(x_k\right)\right| \leqslant\left\|f\right\| \left\| \sum_{k=1}^{\infty}\right\| x_k \|<\infty
\end{equation*}
于是 $\sum_{k=1}^{\infty} f\left(x_k\right)$ 是 $\mathbb{K}$ 上的绝对收敛级数, 重排不变, 所以
\begin{equation*}
\sum_{k=1}^{\infty} f\left(y_k\right)=\sum_{k=1}^{\infty} f\left(x_k\right) \Rightarrow f\left(\sum_{k=1}^{\infty} y_k\right)=f\left(\sum_{k=1}^{\infty} x_k\right), \forall f \in X^* {\Rightarrow} \sum_{k=1}^{\infty} y_k=\sum_{k=1}^{\infty} x_k
\end{equation*}

\end{proof}

\begin{problem}
设 $p$ 是实线性空间 ${X}$ 上的次线性泛函, 求证:
\begin{enumerate}
    \item $p(\theta)=0$;
    \item $p(-x) \geqslant -p(x)$;
    \item 任意给定 $x_0 \in {X}$，在 ${X}$ 上必有实线性泛函 $f$ ，满足 $f\left(x_0\right)=p\left(x_0\right)$ ，以及 $f(x) \leqslant p(x) \ (\forall x \in {X})$。
\end{enumerate}
\end{problem}
\begin{proof}
(1) 由齐次性可得
\begin{equation*}
\begin{aligned}
& p(2 \theta)=2 p(\theta) \Rightarrow p(\theta)=2 p(\theta)  \Rightarrow p(\theta)=0
\end{aligned}
\end{equation*}
(2)
\begin{equation*}
\begin{aligned}
& 0=p(\theta)=p(x-x) \leqslant p(x)+p(-x) 
& \Rightarrow p(-x) \geqslant-p(x)
\end{aligned}
\end{equation*}
(3) 令${X}_0=\left\{\alpha x_0|\alpha\in\mathbb{R}\right\}$,$f_0(x)=f_0\left(\alpha x_0\right)=\alpha f_0\left(x_0\right)=\alpha p\left(x_0\right)(\forall x\in{X}_0)$。\\
当  $\alpha \geqslant 0$ ,$f_0\left(\alpha x_0\right) \leqslant p\left(\alpha x_0\right) $显然成立。\\
当 $\alpha<0$,$ p\left(\alpha x_0\right)=p\left(-|\alpha| x_0\right) \geqslant-p\left(|\alpha| x_0\right)  =-|\alpha| p\left(x_0\right)=\alpha p\left(x_0\right)=f_0\left(\alpha x_0\right)$。\\
所以$f_0(x)\leqslant p(x)(\forall x\in {X}_0)$成立，由实Hahn-Banach定理知存在 ${X}$ 上的实线性泛函 $f$ ，满足 $f\left(x_0\right)=p\left(x_0\right)$ ，以及 $f(x) \leqslant p(x) \ (\forall x \in {X})$。
\end{proof}
\begin{problem}
设 ${X}$ 是由实数列 $x =\left\{a_n\right\}$ 全体组成的实线性空间，其元素间相等和线性运算都按坐标定义，定义
\[
p(x)=\varlimsup_{n \rightarrow \infty} a_n \quad\left(\forall x=\left\{\alpha_n\right\} \in {X} \right) .
\]
求证: $p(x)$ 是 ${X}$ 上的次线性泛函。
\end{problem}

\begin{problem}
设 ${X}$ 是复线性空间，$p$ 是 ${X}$ 上的半范数。$\forall x_0 \in {X}$ ，且 $p\left(x_0\right) \neq 0$。求证：存在 ${X}$ 上的线性泛函 $f$ 满足:
\begin{enumerate}
    \item $f\left(x_0\right)=1$;
    \item $|f(x)| \leqslant \frac{p(x)}{p\left(x_0\right)} \quad (\forall x \in {X})$。
\end{enumerate}
\end{problem}
\begin{proof}
令${X}_0=\left\{\alpha x_0|\alpha\in\mathbb{C}\right\}$,$f_0(x)=f_0\left(\alpha x_0\right)=\alpha f_0\left(x_0\right)=\alpha p\left(x_0\right)(\forall x\in{X}_0)$。\\
\begin{equation*}
\begin{aligned}
 \left|f_0(x)\right|=\left|f_0\left(\alpha x_0\right)\right|=\left|\alpha f_0\left(x_0\right)\right|=|\alpha| p\left(x_0\right) 
 =p\left(\alpha x_0\right)=p(x)
\end{aligned}
\end{equation*}
于是根据 Hahn-banach 定理,
\begin{equation*}
\exists f_1(x) \in {X}^*, \text { s. t. }
\end{equation*}
(1) $\left|f_1(x)\right| \leqslant p(x) \quad(\forall x \in {X})$;\\
(2) $f_1(x)=f_0(x) \quad\left(\forall x \in {X}_0\right)$再令 $f(x)=\frac{f_1(x)}{p\left(x_0\right)}$, 即有\\
(1)
\begin{equation*}
f\left(x_0\right)=\frac{f_1\left(x_0\right)}{p\left(x_0\right)}=\frac{f_0\left(x_0\right)}{p\left(x_0\right)}=\frac{p\left(x_0\right)}{p\left(x_0\right)}=1 ;
\end{equation*}
(2) $|f(x)|=\left|\frac{f(x)}{p\left(x_0\right)}\right| \leqslant \frac{p(x)}{p\left(x_0\right)}$即得结论。
\end{proof}
\begin{problem}
设 ${X}$ 是 $B^*$ 空间，$\left\{x_n\right\}(n=1,2,3, \cdots)$ 是 ${X}$ 中的点列。如果 $\forall f \in {X}^*$，数列 $\left\{f\left(x_n\right)\right\}$ 有界，求证: $\left\{x_n\right\}$ 在 ${X}$ 内有界。
\end{problem}


\begin{proof}
    不妨设$x_n\neq \theta$，由命题2.3知${X}^*$是Banach空间，定义$\hat{x}_n\in {L}({X}^*,\mathbb{K}),\hat{x}_n(f)=f(x_n)$
    \[\|\hat{x}_n\|_{{X}^{**}}=\sup\limits_{\|f\|_{{X}^*\neq 0}}\frac{|\hat{x}_n(f)|}{\|f\|_{{X}^*}}=\sup\limits_{\|f\|_{{X}^*\neq 0}}\frac{|f(x_n)|}{\|f\|_{{X}^*}}\leqslant \sup\limits_{\|f\|_{{X}^*\neq 0}}\frac{|f(x_n)|}{\frac{|f(x_n)|}{\|x_n\|_{{X}}}}=\|x_n\|_{{X}}\]
    由推论2.3知$\exists f\in {X}^*$s.t.$|f(x_n)|=\|x_n\|_{{X}}$且$\|f\|_{{X}^*}=1$，所以$\|\hat{x}_n\|_{{X}^{**}}=\|x_n\|_{{X}}$，由共鸣定理得
    \[
\forall f\in {X}^{*},\sup\limits_{n\geqslant 1}|\hat{x}_n(f)|=\sup\limits_{n\geqslant 1}|\hat{x}_n(f)|
    =\sup\limits_{n\geqslant 1}|f(x_n)|<+\infty
    \]
    存在$M>0$使得$\|x_n\|_{{X}}=\|\hat{x}_n\|_{{X}^{**}}<M$。即 $\left\{x_n\right\}$ 在 ${X}$ 内有界。
\end{proof}
\begin{problem}
设 ${X}_0$ 是 $B^*$ 空间 ${X}$ 的闭子空间，求证：
\[
\begin{gathered}
\rho\left(x, {X}_0\right)=\sup \left\{|f(x)| \mid f \in {X}^*,  \|f\|=1,  f\left({X}_0\right)=0\right\} 
(\forall x \in {X})
\end{gathered}
\]
其中 $\rho\left(x, {X}_0\right)=\inf\limits_{y \in {X}_0}\|x-y\|$ 。
\end{problem}
\begin{proof}
由习题2.1.8知 $|f(x)|=\rho(x, H_f^0) \leqslant \rho\left(x, {X} _0\right)$$\Rightarrow \sup \left\{|f(x)| \mid f \in {X} ^*,\|f\|=1, f\left( {X} _0\right)=0\right\} \leqslant \rho\left(x, {X} _0\right)$;\\
另一方面，等式 $\sup \left\{|f(x)|\mid f \in {X} ^*,\|f\|=1, f\left( {X} _0\right)=0\right\}=\rho\left(x, {X} _0\right)$ 对 $x \in {X} _0$显然成立。\\
对 $x \notin {X} _0$, 由定理 2.19, 存在 $f \in {X} ^*,\|f\|=1, f\left( {X} _0\right)=0, f(x)=\rho\left(x, {X} _0\right)$。\\
故有 $\sup \left\{|f(x)| \mid f \in {X} ^*,\|f\|=1, f\left( {X} _0\right)=0\right\}=\rho\left(x, {X} _0\right)$。
\end{proof}
\begin{problem}
设 $B$ 是 $B^*$ 空间。给定 ${X}$ 中 $n$ 个线性无关的元 $x_1, x_2, \cdots, x_n$，与数域 $K$ 中的 $n$ 个数 $C_1, C_2, \cdots, C_n$，及 $M>0$，求证：为了 $\exists f \in {X}^*$ 适合 $f\left(x_k\right)=C_k \ (k=1,2, \cdots, n)$，以及 $\|f\| \leqslant M$，必须且仅须对任意的 $a_1, a_2, \cdots, a_n \in K$，有
\[
\left|\sum_{k=1}^n a_k C_k\right| \leqslant M \left\|\sum_{k=1}^n a_k x_k\right\| .
\]
\end{problem}

\begin{problem}
给定 $B^*$ 空间 ${X}$ 中 $n$ 个线性无关的元素 $x_1, x_2, \cdots, x_n$，求证：$\exists f_1, f_2, \cdots, f_n \in {X}^*$，使得
\[
\langle f_i, x_j \rangle = \delta_{ij} \quad (i, j = 1,2, \cdots, n).
\]
\end{problem}
\begin{proof}
令 $M_i \triangleq \operatorname{span}\limits_{j\neq i} \left\{x_j\right\}, d_i=\rho\left(x_i, M_i\right)$, 则 $d_i>0$，由推论 2.4，对 $M_i, \because d_i>0, \exists \bar{f}_i \in {X} ^*$ ，使得
\begin{equation*}
\begin{aligned}
& \left\|\bar{f}_i\right\|=1 \\
& \bar{f}_i\left(M_i\right)=0 \\
& \bar{f}_i\left(x_i\right)=d_i
\end{aligned} \stackrel{f_i=\frac{\bar{f}_i}{d_i}}{\Rightarrow} f_i\left(x_j\right)=0 \quad(j \neq i), \quad f_i\left(x_i\right)=1
\end{equation*}

\end{proof}
\begin{problem}
设 ${X}$ 是线性空间，求证：为了 $M$ 是 ${X}$ 的极大线性子空间，必须且仅须 $\operatorname{dim}({X} / M) = 1$。
\end{problem}

\begin{problem}
设 ${X}$ 是复线性空间，$E$ 是 ${X}$ 中的非空均衡集，$f$ 是 ${X}$ 上的线性泛函，求证：
\[
|f(x)| \leqslant \sup_{y \in E} \operatorname{Re} f(y) \quad (\forall x \in E).
\]
\end{problem}

\begin{problem}
设 ${X}$ 是 $B^*$ 空间，$E \subset {X}$ 是非空的均衡闭凸集，$\forall x_0 \in {X} \backslash E$，求证：$\exists f \in {X}^*$ 及 $\alpha>0$，使得
\[
|f(x)| < \alpha < \left|f\left(x_0\right)\right| \quad (\forall x \in E).
\]
\end{problem}

\begin{problem}
设 $E, F$ 是 $B^*$ 空间 ${X}$ 中的两个互不相交的非空凸集，且 $E$ 是开的和均衡的。求证：$\exists f \in {X}^*$，使得
\[
|f(x)| < \inf_{y \in F} |f(y)| \quad (\forall x \in E).
\]
\end{problem}

\begin{problem}
设 $C$ 是实 $B^*$ 空间 ${X}$ 中的一个凸集，并设 $x_0 \in \dot{C}$，$x_1 \in \partial C$，$x_2 = m(x_1 - x_0) + x_0 \ (m > 1)$。求证：$x_2 \notin \bar{C}$。
\end{problem}
\begin{problem}
设 $M$ 是 $B^*$ 空间 ${X}$ 中的闭凸集，求证：$\forall x \in {X} \backslash M$，必 $\exists f_1 \in {X}^*$，满足 $\|f_1\| = 1$，且
\[
\sup_{y \in M} f_1(y) \leqslant f_1(x) - d(x),
\]
其中 $d(x) = \inf\limits_{z \in M} \|x - z\|$。
\end{problem}

\begin{proof}
根据推论 2.6, 存在 $f \in {X} ^*, \alpha \in \mathbb{R} $, 使得 $M$ 闭凸集
\begin{equation*}
\begin{aligned}
 f(y) \leqslant \alpha \leqslant f(z) \quad(\forall y \in M, z \in B(x, d))
 \Rightarrow \sup _{y \in M} f(y) \leqslant \inf _{x \in B(x, d(x))} f(z)=\inf _{y \in B(\theta, 1)} f(x-d(x) y)  \\ =\inf_{y \in B(\theta, 1)}[f(x)-d(x) f(y)]=f(x)-d(x) \sup _{y \in B(\theta, 1)} f(y) = f(x)-d(x)\|f\| .
\end{aligned}
\end{equation*}
取 $f_1=\frac{f}{\| f\|}$ 为所求。
\end{proof}
\begin{problem}
设 $M$ 是实 $B^*$ 空间 ${X}$ 内的闭凸集，求证：
\[
\inf_{z \in M} \|x - z\| = \sup_{\substack{f \in {X}^* \\ \|f\| \leqslant 1}} \left\{ f(x) - \sup_{z \in M} f(z) \right\} \quad (\forall x \in {X}).
\]
\end{problem}
\begin{proof}
 记 $d(x)=\inf _{z \in M}\|x-z\|$，如果 $x \in M$, 左边: $d(x)=0$.
右边: 一方面对 $\forall f \in {X} ^*, f(x)-\sup _{z \in M} f(z) \leqslant 0$, 蕴含右边 $\leqslant 0$.
另一方面, 又存在 $f \in {X} ^*, f(M)=0,\|f\|=1$, 蕴含右边 $\geqslant 0$, 故右边也 $=0$.
因此当 $x \in M$ 时等式成立。\\
$ \forall x \in {X} \backslash M, \forall \varepsilon>0, \quad \exists z_{\varepsilon} \in M \text {, s.t. }\left\|x-z_{\varepsilon}\right\|<d(x)+\varepsilon . \\$
\begin{equation*}
\begin{aligned}
\sup _{\substack{f \in {X}^* \\ \|f\| \leqslant 1}}\left\{f(x)-\sup _{z \in M} f(z)\right\} \leqslant \sup _{\substack{f \in {X}^* \\ \|f\| \leqslant 1}}\left\{f(x)-f\left(z_{\varepsilon}\right)\right\} = \sup _{\substack{f \in {X}^* \\ \|f\| \leqslant 1}}\left\{f\left(x-z_{\varepsilon}\right)\right\} 
 \leqslant \sup _{\substack{f \in {X}^*\\  \|f\| \leqslant 1}}\left\{\|f\|\left\|x-z_{\varepsilon}\right\|\right\}<d(x)+\varepsilon 
\end{aligned}
\end{equation*}
根据 $\varepsilon>0$ 任意性, 有 $\sup \limits_{\substack{f \in {X}^* \\ \|f\| \leqslant 1}}\left\{f(x)-\sup\limits _{z \in M} f(z)\right\} \leqslant d(x)$。\\
另一方面，根据习题2.4.13，必 $\exists f_0 \in {X} ^*,\left\|f_0\right\|=1$ ，使得 $d(x) \leqslant f_0(x)-\sup \limits_{z \in M} f_0(z)$,\\
所以 $d(x)=\sup\limits _{\substack{f \in {X}^* \\ \|f\| \leqslant 1}}\left\{f(x)-\sup \limits_{z \in M} f(z)\right\}$. 并且右端上确界在 $f_0$ 达到。
\end{proof}

\section{共轭空间}
\begin{definition}
设 $X$ 是一个 $B^*$ 空间， $X$ 上的所有连续线性泛函全体 $X ^*$ ，按范数
\begin{equation*}
\|f\|=\sup _{\|x\|=1}|f(x)|
\end{equation*}
构成一个 $B$ 空间, 称为 $X$ 的共轭空间。
\end{definition}
\begin{theorem}[Riesz]
设 $1 \leqslant p<\infty$, 则 $\left(L^p\right)^*=L^q$ ，其中
\begin{equation*}
q= \begin{cases}\frac{p}{p-1} & , 1<p<\infty \\ \infty & , p=1\end{cases}
\end{equation*}
\end{theorem}
\begin{proof}
我们希望构造出一个线性等距同构 $J: L^q \rightarrow\left(L^p\right)^*$ ，如下：
\begin{equation*}
\begin{gathered}
\forall g \in L^q, \Lambda_g: L^p \rightarrow \mathbb{K} , f \mapsto \Lambda_g(f) \stackrel{\text { def }}{=} \int f g \\
J: L^q \rightarrow\left(L^p\right)^*, g \mapsto \Lambda_g
\end{gathered}
\end{equation*}
需要证明：\\
$1^{\circ} \Lambda_g \in\left(L^p\right)^*$ 。\\
$2^{\circ} J$ 线性。\\
$3^{\circ}\left\|\Lambda_g\right\|=\|g\|_q$ ，即 $J$ 等距。\\
$4^{\circ} \forall \Lambda \in\left(L^p\right)^*$, 存在唯一 $g \in L^{p^{\prime}}$ 使得 $\Lambda=\Lambda_g$, 即 $J$ 是双射。\\
Proof of $1^{\circ}-3^{\circ}$ : 先考虑 $1<p<\infty$,
\begin{equation*}
\forall f \in L^p,\left|\Lambda_g(f)\right|=\left|\int f g\right| \leqslant\|g\|_q \cdot\|f\|_p
\end{equation*}
于是 $\Lambda_g \in\left(L^p\right)^*$, 且 $\left\|\Lambda_g\right\| \leqslant\|g\|_q$ 。取
\begin{equation*}
\tilde{f} \stackrel{\text { def }}{=}|g|^{q-1} \operatorname{sgn}(g)
\end{equation*}
则
\begin{equation*}
\begin{aligned}
\left\{\begin{array}{l}
|\tilde{f}|^p=|g|^{(q-1) p}=|g|^q \\
\tilde{f} \cdot g=|g|^q
\end{array}\right. & \Rightarrow\left\{\begin{array}{l}
|\tilde{f}|_p^p=\|g\|_q^q \\
\Lambda_g(\tilde{f})=\|g\|_q^q
\end{array}\right. \\
& \Rightarrow \frac{\left|\Lambda_g(\tilde{f})\right|}{\|\tilde{f}\|_p}=\frac{\|g\|_q^q}{\|g\|_q^{\frac{q}{p}}}=\|g\|_q^{q\left(1-\frac{1}{p}\right)}=\|g\|_q \\
& \Rightarrow\left\|\Lambda_g\right\| \geqslant\|g\|_q
\end{aligned}
\end{equation*}
接下来考虑 $p=1$, 此时需假设 $\mu$ 是 $\sigma$-有限的, 不妨设 $\mu$ 有限, 由
\begin{equation*}
\left|\Lambda_g(f)\right| \leqslant\|g\|_{\infty}\|f\|_1
\end{equation*}
得 $\Lambda_g \in\left(L^1\right)^*$ 且 $\left\|\Lambda_g\right\| \leqslant\|g\|_{\infty}$, 令
\begin{equation*}
\begin{aligned}
& \qquad E_k \stackrel{\text { def }}{=}\left\{t \in \Omega:|g(t)|>\left\|\Lambda_g\right\|+\frac{1}{k}\right\}, f_k \stackrel{\text { def }}{=} \chi_{E_k} \cdot \operatorname{sgn}(g) \\
& \Rightarrow\left\|f_k\right\|_1=\int_{E_k}|\operatorname{sgn}(g)| \dif \mu \leqslant \mu\left(E_k\right) \\
& \Rightarrow\left\|\Lambda_g\right\| \mu\left(E_k\right) \geqslant\|\Lambda_g\|\|f_k\|_1 \geqslant\left|\Lambda_g\left(f_k\right)\right|=\int \chi_{E_k} \operatorname{sgn}(g) g \dif \mu=\int_{E_k}|g| \dif \mu \geqslant\left(\left\|\Lambda_g\right\|+\frac{1}{k}\right) \mu\left(E_k\right) \\
& \Rightarrow \mu\left(E_k\right)=0, \forall k \\
& \Rightarrow\left\{t \in \Omega:|g(t)|>\left\|\Lambda_g\right\|\right\}=\bigcup_{k=1}^{\infty} E_k \text { 零测 } \\
& \Rightarrow\|g\|_{\infty} \leqslant\left\|\Lambda_g\right\|
\end{aligned}
\end{equation*}
Proof of $4^{\circ}$ : 以下假设 $\Omega=[0,1]$,对任意的可测集 $E \subset[0,1]$, 令
\begin{equation*}
\nu(E) \triangleq \Lambda\left(\chi_E\right),
\end{equation*}
对任意不交的集合序列 $\left(E_n\right)_{n \geq 1}$ ，如果 $E=\bigcup_{n=1}^{\infty} E_n$ ，则 $\chi_E=\sum_{n=1}^{\infty} \chi_{E_n}$ ，其中级数在 $L^p$ 意义收敛:
\begin{equation*}
\left\|\chi_E-\sum_{j=1}^n \chi_{E_j}\right\|_p=\left\|\sum_{j=n+1}^{\infty} \chi_{E_j}\right\|_p=\mu\left(\bigcup_{j=n+1}^{\infty} E_j\right)^{1 / p} \rightarrow 0, \quad n \rightarrow \infty
\end{equation*}
结合 $\Lambda$ 的线性性和连续性，
\begin{equation*}
\nu(E)=\sum_{n=1}^{\infty} \Lambda\left(\chi_{E_n}\right)=\sum_{n=1}^{\infty} \nu\left(E_n\right)
\end{equation*}
这就证明了 $\nu$ 确实是一个测度. 另外, 如果 $\mu(E)=0$ ，
\begin{equation*}
\begin{aligned}
\nu\left(E\right) & =\Lambda\left(\chi_{E}\right) \leqslant\|\Lambda\| \cdot\left\|\chi_{E}\right\|_{L^p[0,1]} \\
& =\|\Lambda\|\left(\int_0^1\left|\chi_{E}\right|^p \dif \mu\right)^{\frac{1}{p}} \\
& =\|\Lambda\| \mu\left(E\right)^{\frac{1}{p}}=0 
\end{aligned}
\end{equation*}
可知 $\nu$ 关于 $\mu$ 还是绝对连续的，即由 $\mu(E)=0$ ，可以推出 $\nu(E)=0$ 。现在应用 Radon-Nikodym 定理，存在可测函数 $g$ ，使得对任意的可测集 $E$ 有
\begin{equation*}
\nu(E)=\int_E g \dif \mu
\end{equation*}
于是对于一切简单函数 $f$, 都有
\begin{equation*}
\Lambda(f)=\int_0^1 f(x) g(x) \dif \mu
\end{equation*}
进一步我们将要证明：
\begin{equation*}
\|g\|_{L^q[0,1]} \leqslant\|\Lambda\|
\end{equation*}
\begin{enumerate}
    \item  $1<p<\infty$. 对 $\forall t>0$, 记
\begin{equation*}
E_t \triangleq\{x \in[0,1]| | g(x) \mid \leqslant t\}
\end{equation*}
令 $f=\chi_{E_t}|g|^{q-2} g$ ，便有
\begin{equation*}
\begin{aligned}
\int_{E_t}|g|^q d \mu & =\int_0^1 f \cdot g \dif \mu=\Lambda(f) \\
& \leqslant\|\Lambda\| \cdot\|f\|_{L^p[0,1]}=\|\Lambda\|\left(\int_{E_t}|g|^q \dif \mu\right)^{\frac{1}{p}},
\end{aligned}
\end{equation*}
亦即
\begin{equation*}
\left(\int_{E_t}|g|^q \dif \mu\right)^{\frac{1}{q}} \leqslant\|\Lambda\| .
\end{equation*}
令 $t \rightarrow \infty$, 即得上式.
\item  $p=1$. 这时 $q=\infty$. 对 $\forall \varepsilon>0$, 令
\begin{equation*}
A \triangleq\{x \in[0,1]||g(x)|>\|\Lambda\|+\varepsilon\}
\end{equation*}
再对 $\forall t>0$, 还按上式定义 $E_t$, 并令 $f=\chi_{E_t \cap A} \operatorname{signg}$, 便有
\begin{equation*}
\|f\|_{L^1[0,1]}=\mu\left(E_t \cap A\right)
\end{equation*}
并且有
\begin{equation*}
\begin{aligned}
\mu\left(E_t \cap A\right)(\|\Lambda\|+\varepsilon) & \leqslant \int_{A \cap E_t}|g| \dif \mu \\
& =\int_0^1 f \cdot g\dif \mu \leqslant\|\Lambda\| \mu\left(E_t \cap A\right)
\end{aligned}
\end{equation*}
令 $t \rightarrow \infty$, 便得
\begin{equation*}
\mu(A)(\|\Lambda\|+\varepsilon) \leqslant\|\Lambda\| \mu(A)
\end{equation*}
由此推出 $\mu(A)=0$, 从而
\begin{equation*}
\|g\|_{L^{\infty}[0,1]} \leqslant\|\Lambda\|
\end{equation*}
这就是当 $q=\infty$ 时的 上式.

\end{enumerate}
因为简单函数集在 $L^p[0,1]$ 中是稠密的，所以对 $\forall f \in L^p[0,1]$ ，存在简单函数列 $f_n \rightarrow f\left(L^p[0,1]\right)$. 从而有
\begin{equation*}
\Lambda(f)=\lim _{n \rightarrow \infty} \Lambda\left(f_n\right),
\end{equation*}
以及
\begin{equation*}
\begin{aligned}
& \left|\int_0^1\left[f(x)-f_n(x)\right] g(x) \dif \mu\right| \\
& \quad \leqslant\left(\int_0^1\left|f(x)-f_n(x)\right|^p \dif \mu\right)^{\frac{1}{p}}\left(\int_0^1|g(x)|^q \dif \mu\right)^{\frac{1}{q}} \\
& \quad \leqslant\|\Lambda\| \cdot\left\|f-f_n\right\|_{L^p[0,1]} \rightarrow 0 \quad(n \rightarrow \infty),
\end{aligned}
\end{equation*}
亦即
\begin{equation*}
\Lambda(f)=\lim _{n \rightarrow \infty} \int_0^1 f_n(x) g(x) \dif \mu=\int_0^1 f(x) g(x) \dif \mu
\end{equation*}



\end{proof}
\begin{theorem}


\begin{equation*}
L^1 \varsubsetneqq\left(L^{\infty}\right)^*
\end{equation*}


\end{theorem}
\begin{proof}
对于 $\forall g \in L^1$,
\begin{equation*}
\left|\Lambda_g(f)\right|=\left|\int f g\right| \leqslant\|g\|_1\|f\|_{\infty} \Rightarrow \Lambda_g \in\left(L^{\infty}\right)^* \Rightarrow L^1 \subset\left(L^{\infty}\right)^*
\end{equation*}
注意到 $C[0,1]$ 是 $L^{\infty}$ 的闭子空间, 取 $f_0 \in L^{\infty} \backslash C[0,1]$, 于是 $d \stackrel{\text { def }}{=} \operatorname{dist}\left(f_0, C[0,1]\right)>0$, 由 HBT 可知存在 $\Lambda \in$$\left(L^{\infty}\right)^*,\|\Lambda\|=1$ 且
\begin{equation*}
\Lambda(C[0,1])=\{0\}, \Lambda\left(f_0\right)=d
\end{equation*}
假设存在 $g \in L^1$ 满足 $\Lambda=\Lambda_g$ ，即
\begin{equation*}
\Lambda(f)=\int f g, \forall f \in L^{\infty}
\end{equation*}
那么对于 $f \in C[0,1], \Lambda(f)=\int f g=0$, 取 $\left\{f_n\right\}_{n=1}^{\infty} \subset C[0,1]$, 满足
\begin{equation*}
\left\|f_n-\operatorname{sgn}(g)\right\|_1 \rightarrow 0 \text { as } n \rightarrow \infty
\end{equation*}
由F.Reisz定理有子列 $f_{n_k} \xrightarrow{\text { a.e. }} \operatorname{sgn}(g)$ ，取$f'_{n_k}=\max\{\min\{f_{n_k},1\},-1\},f'_{n_k} \xrightarrow{\text { a.e. }} \operatorname{sgn}(g),| f'_{n_k} g|\leqslant |g|$由 DCT 得
\begin{equation*}
\int|g|=\lim _{k \rightarrow \infty} \int f'_{n_k} g=0 \Rightarrow g \stackrel{\text { a.e. }}{=} 0 \Rightarrow \Lambda=\Lambda_g=0
\end{equation*}
这与 $\Lambda\left(f_0\right)=d>0$ 矛盾。
\end{proof}
\begin{definition}
对 $f:[a, b] \rightarrow C$ 和 $[a, b]$ 的划分 $P$, 定义
\begin{equation*}
V(f, P) \stackrel{\text { def }}{=} \sum_{k=1}^N\left|f\left(t_k\right)-f\left(t_{k-1}\right)\right|
\end{equation*}
如果 $\sup _P V(f, P)<\infty$, 则称 $f$ 是有界变差的。
\begin{equation*}
V_a^b(f) \stackrel{\text { def }}{=} \sup _P V(f, P)
\end{equation*}
称为 $f$ 在 $[a, b]$ 上的全变差。
\begin{equation*}
B V[a, b] \stackrel{\text { def }}{=}\{[a, b] \text { 上的有界变差函数全体 }\}
\end{equation*}
我们进一步给出 $B V[a, b]$ 上的范数:
\begin{equation*}
\|f\|_{B V} \stackrel{\text { def }}{=}|f(a)|+V_a^b(f)
\end{equation*}
那么 $\left(B V[a, b],\|\cdot\|_{B V}\right)$ 是 Banach 空间，定义:
\begin{equation*}
B V_0[a, b] \stackrel{\text { def }}{=}\{f \in B V[a, b]: f \text { 在 }(a, b) \text { 上右连续, } f(a)=0\}
\end{equation*}
那么 $B V_0[a, b]$ 是 $B V[a, b]$ 的闭子空间, 也是 Banach 空间。
\end{definition}
\begin{definition}[Riemann-Stieltjes 积分]
设 $f, g$ 是 $[a, b]$ 上实值函数, $I \in \mathbb{R}$, 对 $[a, b]$ 进行划分:
\begin{equation*}
\sigma(\Delta, \xi) \stackrel{\text { def }}{=} \sum_{k=1}^n f\left(\xi_k\right)\left[g\left(t_k\right)-g\left(t_{k-1}\right)\right]
\end{equation*}
其中
\begin{equation*}
\xi=\left\{\xi_k\right\}_{k=1}^{\infty}, \xi_k \in\left[t_{k-1}, t_k\right]
\end{equation*}
如果 $\sigma(\Delta, \xi) \rightarrow I$ as $\|\Delta\| \rightarrow 0$, 则记
\begin{equation*}
I=\int_a^b f \dif g
\end{equation*}
称为 $f$ 关于 $g$ 的 Riemann-Stieltjes 积分，简记为 R-S 积分。
\end{definition}
\begin{proposition}
$BV_0[0,1]$的共轭空间是$C[0,1]$。
\end{proposition}
\begin{proof}
从定义易见 $\forall g \in B V_0[0,1]$, 它对应着 $C[0,1]$ 上的一个连续线性泛函$f$
\begin{equation*}
f:\varphi \mapsto\langle f, \varphi\rangle=\int_0^1 \varphi(t) \dif g(t)
\end{equation*}
满足
\begin{equation*}
\|f\|=\sup\limits_{\theta\neq \varphi\in C[0,1]}\frac{\int_0^1 \varphi(t) \dif g(t)}{\sup\limits_{t\in [0,1]}|\varphi(t)|} \leqslant \operatorname{var}(g)=\|g\|_v
\end{equation*}
现在我们要证明反过来的结论. 这也就是说，对任意的 $f \in$ $C[0,1]^*$ ，必 $\exists \mid g \in B V[0,1]$ ，使得
\begin{equation*}
\langle f, \varphi\rangle=\int_0^1 \varphi(t) \dif g(t) \quad(\forall \varphi \in C[0,1])
\end{equation*}
并且
\begin{equation*}
\|g\|_v \leqslant\|f\| .
\end{equation*}
其中关键的步骤在于如何由 $f$ 确定出 $g$. 事实上，如果允许 $\varphi$ 可以取成间断的特征函数 $\chi_E(t)(E \subset[0,1])$ ，那么就有
\begin{equation*}
g(s)=\int_0^1 \chi_{(0, s]}(t) \dif g(t)
\end{equation*}
因此, 先把 $C[0,1]$ 看成是 $L^{\infty}[0,1]$ 的一个闭子空间, 应用 HahnBanach 定理，对给定的 $f \in C[0,1]^*, \exists \widetilde{f} \in L^{\infty}[0,1]^*$ ，使得
\begin{equation*}
\begin{gathered}
\left\langle\widetilde{f}, \chi_{\{0\}}\right\rangle=0 \\
\langle\widetilde{f}, \varphi\rangle=\langle f, \varphi\rangle \quad(\forall \varphi \in C[0,1])
\end{gathered}
\end{equation*}
并且 $\|\widetilde{f}\|=\|f\|$. 因为 $\chi_{(0, s]} \in L^{\infty}[0,1]$ ，所以可以令
\begin{equation*}
\begin{aligned}
& g(s)=\left\langle\widetilde{f}, \chi_{(0, s]}\right\rangle \quad(0<s \leqslant 1) \\
& g(0)=0
\end{aligned}
\end{equation*}
以下证明由 (2.5.13) 式定义的 $g \in B V[0,1]$, 并适合 (2.5.11)式与（2.5.12）式。
对 $[0,1]$ 的任一分割 $\Delta: 0=t_0<t_1<t_2<\cdots<t_n=1$. 记
\begin{equation*}
\begin{gathered}
\omega_k \triangleq g\left(t_{k+1}\right)-g\left(t_k\right) \quad(k=0,1,2, \cdots, n-1), \\
\lambda_k \triangleq\left\{\begin{array}{ll}
\bar{\omega}_k /\left|\omega_k\right|, & \omega_k \neq 0, \\
0, & \text { 其他 }
\end{array} \quad(k=0,1,2, \cdots, n-1),\right. \\
h_{\Delta}(t) \triangleq \sum_{k=0}^{n-1} \lambda_k \chi_{\left(t_k, t_{k+1}\right]}(t) .
\end{gathered}
\end{equation*}
我们有
\begin{equation*}
\begin{aligned}
\sum_{k=0}^{n-1}\left|\omega_k\right| & =\sum_{k=0}^{n-1} \lambda_k \omega_k=\left\langle\widetilde{f}, h_{\Delta}\right\rangle \\
& \leqslant\|\widetilde{f}\| \cdot\left\|h_{\Delta}\right\|_{L^{\infty}[0,1]} \leqslant\|\widetilde{f}\|=\|f\|
\end{aligned}
\end{equation*}
即得 $g \in B V[0,1]$, 并适合 (2.5.12) 式.
为了证明 (2.5.11) 式成立, 对 $\forall \varphi \in C[0,1]$ 及 $\forall \varepsilon>0$, 取分割 $\Delta$ 使得
\begin{equation*}
\begin{gathered}
\left|\varphi(t)-\varphi\left(t^{\prime}\right)\right|<\frac{\varepsilon}{2\|\widetilde{f}\|} \\
\left(\forall t, t^{\prime} \in\left[t_j, t_{j+1}\right], j=0,1, \cdots, n-1\right)
\end{gathered}
\end{equation*}
以及
\begin{equation*}
\left|\int_0^1 \varphi(t) d g(t)-\sum_{j=0}^{n-1} \varphi\left(t_j\right)\left(g\left(t_{j+1}\right)-g\left(t_j\right)\right)\right|<\frac{\varepsilon}{2}
\end{equation*}
令
\begin{equation*}
\varphi_{\Delta} \triangleq \sum_{j=0}^{n-1} \varphi\left(t_j\right) \chi_{\left(t_j, t_{j+1}\right]}+\varphi(0) \chi_{\{0\}}
\end{equation*}
便得
\begin{equation*}
\begin{aligned}
& \left|\langle f, \varphi\rangle-\int_0^1 \varphi(t) d g(t)\right| \\
& \quad \leqslant\left|\langle f, \varphi\rangle-\left\langle\widetilde{f}, \varphi_{\Delta}\right\rangle\right|+\left|\left\langle\widetilde{f}, \varphi_{\Delta}\right\rangle-\int_0^1 \varphi(t) d g(t)\right| \\
& \quad \leqslant\|\widetilde{f}\| \cdot\left\|\varphi-\varphi_{\Delta}\right\|_{L^{\infty}[0,1]}+ \\
& \\
& \quad\left|\sum_{j=0}^{n-1} \varphi\left(t_j\right)\left(g\left(t_{j+1}\right)-g\left(t_j\right)\right)-\int_0^1 \varphi(t) d g(t)\right| \\
& \quad<\frac{\varepsilon}{2}+\frac{\varepsilon}{2}=\varepsilon
\end{aligned}
\end{equation*}
由 $\varepsilon$ 的任意性, 即得 (2.5.11) 式.有界变差函数是单调增加函数的差, 即 $g=g_1-g_2$, 因为 $g_1, g_2$ 的右极限总是存在的，可以适当改变函数的值使得其右连续，但是不会改变积分 $\int_0^1 \varphi(t) d g_i(t), i=1,2$ ，因此我们可以假设上述得到的函数 $g$ 是右连续的。总结起来有， $g \mapsto \int_0^1 \varphi(t) d g(t)$ 是 $B V[0,1] \rightarrow C[0,1]^*$ 的一个等距同构. 换句话说
\begin{equation*}
C[0,1]^*=B V[0,1] .
\end{equation*}
\end{proof}
\begin{definition}
 因为 $B^*$ 空间 $X$ 的共轭空间 $X ^*$是一个 $B$ 空间，所以我们还可以考虑 $X ^*$ 的共轨空间，记作 $X ^{* *}$ ，称为 $X$ 的第二共轭空间. 
\end{definition}

\begin{theorem}
$B^*$ 空间 $X$ 与它的第二共轭空间 $X ^{* *}$ 的一个子空间等距同构.
\end{theorem}
\begin{proof}
注意到 $\forall x \in X$, 可以定义
\begin{equation*}
\hat{x}(f)=\langle f, x\rangle \quad\left(\forall f \in X ^*\right) .
\end{equation*}
不难验证： $\hat{x}$ 还是 $X ^*$ 上的一个线性泛函，满足
\begin{equation*}
|\hat{x}(f)| \leqslant\|f\| \cdot\|x\| .
\end{equation*}
从而 $\hat{x}$ 还是连续的, 满足
\begin{equation*}
\|\hat{x}\| \leqslant\|x\|
\end{equation*}
称映射 $T: x \mapsto \hat{x}$ 为自然映射，上式表明 $T$ 是 $X$ 到 $X ^{* *}$的连续嵌入， 注意到, 若 $\alpha, \beta \in \mathbb{C} , x, y \in X$, 记 $\hat{x}=T x, \hat{y}=T y$,则有
\begin{equation*}
\begin{aligned}
& T(\alpha x+\beta y)(f)=f(\alpha x+\beta y) \\
& \quad=\alpha f(x)+\beta f(y)=\alpha \hat{x}(f)+\beta \hat{y}(f) \\
& \quad=(\alpha\hat{x}+\beta \hat{y})(f)=(\alpha T x+\beta T y)(f) \quad\left(\forall f \in X ^*\right)
\end{aligned}
\end{equation*}
因此， $T$ 还是一个线性同构. 又应用 Hahn-Banach 定理, $\exists f \in X ^*$ ，使得
\begin{equation*}
\|f\|=1, \quad \text { 且 } \quad\langle f, x\rangle=\|x\| \text {, }
\end{equation*}
便得到
\begin{equation*}
\|x\|=\hat{x}(f) \leqslant\|\hat{x}\| \cdot\|f\|=\|\hat{x}\|
\end{equation*}
所以 $T$ 是等距的. 
\end{proof}

\begin{definition}
如果 $X$ 到 $X ^{* *}$ 的自然映射 $T$ 是满射的, 则称 $X$ 是自反的, 记作 $X = X ^{* *}$.
\end{definition}

\begin{example}
Hilbert 空间自反。
\end{example}
\begin{proof}
由 Riesz 表示定理，有等距同构 $\varphi: H \rightarrow H^*, x \mapsto\langle\cdot, x\rangle \cdot \varphi$ 诱导了 $H^*$ 上的内积 $\langle f, g\rangle=\overline{\left\langle\varphi^{-1}(f), \varphi^{-1}(g)\right\rangle}$ 。对偶空间自然完备，因此 $H^*$ 也是 Hilbert 空间，由 Riesz 表示定理，又有等距同构 $\Phi: H^* \rightarrow H^{* *}, f \mapsto$ $\overline{\left\langle\varphi^{-1}(\cdot), \varphi^{-1}(f)\right\rangle}$.
\begin{equation*}
\Phi \circ \varphi(x)(f)=\overline{\left\langle\varphi^{-1}(f), x\right\rangle}=\left\langle x, \varphi^{-1}(f)\right\rangle=f(x), \forall x \in H, f \in H^*
\end{equation*}
故 $\Phi \circ \varphi$ 就是 $H \rightarrow H^{* *}$ 的自然映射 $\phi$ ，因此 $\phi$ 也是等距同构，从而 $H$ 自反。
\end{proof}
\begin{theorem}
当 $1<p<\infty$ 时, $L^p$ 自反。

\end{theorem}
\begin{proof}
 即证明: $\forall \Lambda \in\left(L^p\right)^{* *}$, 存在 $u \in L^p$ 使得
\begin{equation*}
\Lambda(f)=f(u), \forall f \in\left(L^p\right)^*
\end{equation*}
这是因为
\begin{equation*}
i: L^p \rightarrow\left(L^p\right)^{* *} \text { 满 } \Leftrightarrow \forall \Lambda \in\left(L^p\right)^{* *}, \exists u \in L^p \text { s.t. } u^{* *}=\Lambda\left(\Lambda(f)=u^{* *}(f)=f(u)\right)
\end{equation*}
回顾定理 2.5.1,
\begin{equation*}
J: L^q \rightarrow\left(L^p\right)^*, v \mapsto f_v
\end{equation*}
是线性等距同构, 这里 $f_v(u)=\int u v$ 。取
\begin{equation*}
\varphi \stackrel{\text { def }}{=} \Lambda \circ J
\end{equation*}
则 $\varphi \in\left(L^p\right)^{**}=L^p$, 存在唯一 $u \in (L^p)^{**}$ 满足 $\varphi(v)=\int u v, \forall v \in (L^p)^*= L^q$. 那么对于 $\forall f \in\left(L^p\right)^*$, 令
\begin{equation*}
v_f \stackrel{\text { def }}{=} J^{-1}(f)
\end{equation*}
那么
\begin{equation*}
\Lambda(f)=\Lambda\left(J\left(v_f\right)\right)=\varphi\left(v_f\right)=\int v_f u=f(u)
\end{equation*}
\end{proof}
\begin{theorem}
$C[a, b]$ 不自反。


\end{theorem}
\begin{proof}
假设自反, 则 $\forall \Lambda \in C[a, b]^{* *}$, 存在 $u \in C[a, b]$ 满足
\begin{equation*}
\Lambda(f)=f(u), \forall f \in C[a, b]^*
\end{equation*}
根据 $C[a, b]^*=B V_0[a, b]$,
\begin{equation*}
\exists!v_f \in B V_0[a, b] \text { s.t. } f(u)=\int_a^b u d v_f, \forall u \in C[a, b] \text { and }\left\|v_f\right\|_{B V}=\|f\|
\end{equation*}
令 $c=\frac{a+b}{2}$, 定义
\begin{equation*}
F_c: C[a, b]^* \rightarrow R , f \mapsto v_f(c+0)-v_f(c-0)
\end{equation*}
那么
\begin{equation*}
\left|F_c(f)\right| \leqslant V_a^b\left(v_f\right)=\left\|v_f\right\|_{B V}=\|f\| \Rightarrow F_c \in C[a, b]^{* *}
\end{equation*}
根据 $\left.N ^*\right)$, 存在 $u_c \in C[a, b]$ 满足
\begin{equation*}
F_c(f)=f\left(u_c\right)=\int_a^b u_c d v_f, \forall f \in C[a, b]^*
\end{equation*}
令
\begin{equation*}
v(f) \stackrel{\text { def }}{=} \int_a^t u_c(s) d s
\end{equation*}
那么 $v \in B V_0[a, b]$, 令 $f_v \stackrel{\text { def }}{=} J(v)$, 这里
\begin{equation*}
J: B V_0[a, b] \rightarrow C[a, b]^*, v \mapsto f_v, f_v(u)=\int_a^b u d v_f
\end{equation*}
$f_v$ 对应的 $v_{f_v}=v$ 连续, 于是 $F_c\left(f_v\right)=0$, 所以
\begin{equation*}
0=F_c\left(f_v\right)=\int_a^b u_c d v=\int_a^b u_c^2 d t
\end{equation*}
所以 $u_c=0$, 进而 $F_c=0$, 矛盾。
\end{proof}

\begin{theorem}
$X^*$ 可分 $\Rightarrow X$ 可分。逆命题不成立, 例如 $L^1$ 可分但 $L^{\infty}$ 不可分。
\end{theorem}
\begin{proof}
Step1: 证明 $X^*$ 中单位球面 $S_1^*$ 可分。 $X^*$ 可分, 取 $\left\{f_n\right\}_{n=1}^{\infty}$ 为其稠密子集, 不妨 $f_n \neq 0$, 令
\begin{equation*}
g_n \stackrel{\text { def }}{=} \frac{f_n}{\left\|f_n\right\|} \in S_1^*
\end{equation*}
那么对于 $\forall g \in S_1^*$, 存在 $f_{n_k} \rightarrow g$.
\begin{equation*}
\begin{aligned}
\Rightarrow\left\|g-g_{n_k}\right\| & \leqslant\left\|g-f_{n_k}\right\|+\left\|f_{n_k}-g_{n_k}\right\| \\
& =\left\|g-f_{n_k}\right\|+\left\|\left(| | f_{n_k} \|-1\right) \frac{f_{n_k}}{\left\|f_{n_k}\right\|}\right\| \\
& =\left\|g-f_{n_k}\right\|+\left|\left\|f_{n_k} \|-1\right| \rightarrow 0 \text { as } k \rightarrow \infty\right.
\end{aligned}
\end{equation*}
Step2: 证明存在 $\left\{x_n\right\}_{n=1}^{\infty} \subset X$, 其中 $\forall\left\|x_n\right\|=1$, 且
\begin{equation*}
\overline{\operatorname{span}\left\{x_n\right\}_{n=1}^{\infty}}=X
\end{equation*}
注意到
\begin{equation*}
\left\|g_n\right\|=\sup _{x \in X,\|x\|=1}\left|g_n(x)\right|=1
\end{equation*}
所以存在 $\left\{x_n\right\}_{n=1}^{\infty} \subset X$, 其中 $\forall\left\|x_n\right\|=1$, 且 $\left|g_n\left(x_n\right)\right|>\frac{1}{2}$. 令 $M \stackrel{\text { def }}{=} \operatorname{span}\left\{x_n\right\}_{n=1}^{\infty}$, 假设 $\bar{M} \neq X$, 取 $x_0 \in X \backslash \bar{M}$,由 HBT 可得
\begin{equation*}
\exists f \in X^*,\|f\|=1 \text { s.t. } f(\bar{M})=\{0\} \text { and } f\left(x_0\right)=\operatorname{dist}\left(x_0, \bar{M}\right)>0
\end{equation*}
于是
\begin{equation*}
\left\|g_n-f\right\|=\sup _{x \in X,\|x\|=1}\left|g_n(x)-f(x)\right|>\left|g_n\left(x_n\right)-f\left(x_n\right)\right|=\left|g_n\left(x_n\right)\right|>\frac{1}{2}
\end{equation*}
这与 Step1 矛盾。\\
Step3: 证明 $\overline{\operatorname{span}^{ \mathbb{Q} }\left\{x_n\right\}_{n=1}^{\infty}}=X$.
\end{proof}

\begin{theorem}
当 $1 \leqslant p<\infty, L^p[0,1]$ 可分。

\end{theorem}
\begin{proof}
\begin{equation*}
\left\{\sum_{k=0}^{2^n-1} r_k \chi_{\left[\frac{k}{2^n}, \frac{k+1}{2^n}\right)}: r_k \in \mathbb{Q} , n \in \mathbb{N} _0\right\}
\end{equation*}
是 $L^p[0,1]$ 的可数稠密子集。
\end{proof}
\begin{theorem}
\begin{equation*}
L^{\infty} \text { 不可分。 }
\end{equation*}

\end{theorem}


\begin{proof}
假设存在稠密子集 $\left\{f_n\right\}_{n=1}^{\infty}$,
\begin{equation*}
\forall t \in(0,1), \exists f_{n_t} \in B\left(\chi_{[0, t]}, \frac{1}{3}\right)
\end{equation*}
而 $t \neq s$ 时, $\operatorname{dist}\left(\chi_{[0, t]}, \chi_{[0, s]}\right)=1$, 因此不同的 $B\left(\chi_{[0, t]}, \frac{1}{3}\right)$ 不相交, 所以 $\varphi:(0,1) \rightarrow \mathbb{N} , t \mapsto n_t$ 是单射, 于是 $(0,1)$ 可数, 矛盾。
\end{proof}

\begin{theorem}
$L^1$ 不自反。

\end{theorem}
\begin{proof}
 $\left(L^1\right)^*=L^{\infty}$, 假设 $L^1$ 自反, 则 $\left(L^{\infty}\right)^* \cong L^1, L^1$ 可分 $\Rightarrow L^{\infty}$ 可分, 矛盾。
\end{proof}

\begin{definition}[共轭算子] 设 $X , Y$ 是 $B^*$ 空间，算子 $T \in$ $L ( X , Y )$ 。算子 $T^*: Y ^* \rightarrow X ^*$ 称为是 $T$ 的共轭算子是指:
\begin{equation*}
f(T x)=\left(T^* f\right)(x) \quad\left(\forall f \in Y ^*, \forall x \in X \right) .
\end{equation*}
\end{definition}
\begin{note}
 $\forall T \in L ( X , Y ), T^*$ 是唯一存在的，并且属于 $L ( Y ^*，X ^*)$ 。事实上, 对 $\forall f \in Y ^*$ ，令
\begin{equation*}
g(x)=f(T x) \quad(\forall x \in X )
\end{equation*}
它是线性的，并且有界：
\begin{equation*}
|g(x)| \leqslant\|f\| \cdot\|T\| \cdot\|x\| \quad(\forall x \in X )
\end{equation*}
因此, $g \in X ^*$, 对应 $f \mapsto g$ 又是线性的, 正是 $T^*$ 。按定义,
\begin{equation*}
\left\|T^* f\right\|=\|g\| \leqslant\|T\| \cdot\|f\| \quad\left(\forall f \in Y ^*\right)
\end{equation*}
因此， $T^* \in L \left( Y ^*, X ^*\right) .$ 上式还蕴含
\begin{equation*}
\left\|T^*\right\| \leqslant\|T\|
\end{equation*}
因此, $T^*$ 的唯一性显然. 
\end{note}

\begin{theorem}
映射 $*: T \mapsto T^*$ 是 $L ( X , Y )$ 到 $L \left( Y ^*, X ^*\right)$子空间的等距同构。
\end{theorem}
\begin{proof}
对应 $*: T \mapsto T^*$ 是线性的$\left(\forall x \in X , \forall f \in Y ^*, \forall \alpha_1, \alpha_2 \in K \right)$:
\begin{equation*}
\begin{aligned}
& {\left[\left(\alpha_1 T_1+\alpha_2 T_2\right)^* f\right](x)=f\left[\left(\alpha_1 T_1+\alpha_2 T_2\right) x\right]} \\
& \quad=\alpha_1 f\left(T_1 x\right)+\alpha_2 f\left(T_2 x\right)=\left[\left(\alpha_1 T_1^*+\alpha_2 T_2^*\right) f\right](x)
\end{aligned}
\end{equation*}
再证等距. 只要再证 $\|T\| \leqslant\left\|T^*\right\|$ 。对 $\forall x \in X$ ，若 $T x \neq \theta$ ，由推论 2.3，必有 $f \in Y ^*$ ，使得
\begin{equation*}
f(T x)=\|T x\|, \quad \text { 且 } \quad\|f\|=1 .
\end{equation*}
从而
\begin{equation*}
\begin{aligned}
\|T x\| & =f(T x)=\left(T^* f\right)(x) \\
& \leqslant\left\|T^* f\right\| \cdot\|x\| \leqslant\left\|T^*\right\| \cdot\|x\|
\end{aligned}
\end{equation*}
即得 $\|T\| \leqslant\left\|T^*\right\|$.
\end{proof}

\begin{theorem}
设 $X , Y$ 是 $B^*$ 空间， $T \in L ( X , Y )$ ，那么 $T^{* *} \in L \left( X ^{* *}, Y { }^{* *}\right)$ 是 $T$ 在 $X ^{* *}$ 上的延拓, 并满足 $\left\|T^{* *}\right\|=\|T\|$.
\end{theorem}

\begin{proof}
对 $T^*$ 还可以再考察它的共轭算子 $T^{* *}=\left(T^*\right)^* \in$ $L \left( X ^{* *}, Y ^{* *}\right)$ 。注意到 $X \subset X ^{* *}, Y \subset Y ^{* *}$ ，并设它们的自然嵌入映射分别为 $U$ 和 $V$ ，那么
\begin{equation*}
\begin{aligned}
\left\langle T^{* *} U x, f\right\rangle & =\left\langle U x, T^* f\right\rangle \\
& =\left\langle T^* f, x\right\rangle=\langle f, T x\rangle \\
& =\langle V T x, f\rangle \quad\left(\forall f \in Y ^*, \forall x \in X \right)
\end{aligned}
\end{equation*}
从而有 $T^{* *} U x=V T x$. 即 $T^{* *}$ 是 $T$ 在 $X ^{* *}$ 上的扩张. 
\end{proof}

\begin{theorem}[pettis]
自反空间的闭子空间自反。
\end{theorem}

\begin{proof}
设 $X$ 自反， $Y$ 是其闭子空间，只需证明：
\begin{equation*}
\forall a \in Y^{* *}, \exists y \in Y \text { s.t. } a(f)=f(y), \forall f \in Y^*
\end{equation*}
定义映射
\begin{equation*}
T: X^* \rightarrow Y^*,\left.f \mapsto f\right|_Y
\end{equation*}
则 $T$ 是有界线性映射，于是取 $T^* \in L \left(Y^{* *}, X^{* *}\right)$ 满足
\begin{equation*}
\left(T^* a\right)(f)=a(T f), \forall f \in X^*
\end{equation*}
$X$ 自反，所以自然映射 $i_X$ 是满射， $T^* a \in X^{* *} \Rightarrow \exists y \in X$ s.t. $T^* a=y^{* *}$ ，所以
\begin{equation*}
\left(T^* a\right)(f)=y^{* *}(f)=f(y), \forall f \in X^*
\end{equation*}
下面证明 $y \in Y$, 假设不然, 则存在 $\tilde{f} \in X^*$ 使得 $\tilde{f}(Y)=0$,
\begin{equation*}
T(\tilde{f})=\left.\tilde{f}\right|_Y=0 \Rightarrow \tilde{f}(y)=\left(T^* a\right)(\tilde{f})=a(T(\tilde{f}))=0
\end{equation*}
这与 $\tilde{f}(y)=\operatorname{dist}(y, Y)>0$ 矛盾。\\
最后说明: $a(f)=f(y), \forall f \in Y^*$. 对于 $\forall f \in Y^*$, 由 HBT, 存在 $F \in X^*$ 使得 $f=T F$, 所以
\begin{equation*}
a(f)=a(T F)=\left(T^* a\right)(F)=F(y)=f(y)
\end{equation*}
\end{proof}

\begin{example}
设 $(\Omega, B , \mu)$ 是一个测度空间, 又设 $K(x, y)$ 是 $\Omega \times \Omega$ 上的二元平方可积函数:
\begin{equation*}
\iint_{\Omega \times \Omega}|K(x, y)|^2 \dif \mu(x) \dif \mu(y)<\infty
\end{equation*}
定义算子
\begin{equation*}
T: u \mapsto(T u)(x)=\int_{\Omega} K(x, y) u(y) \dif \mu(y) \quad\left(\forall u \in L^2(\Omega, \mu)\right)
\end{equation*}
便有 $T \in L \left(L^2(\Omega, \mu)\right)$ ，并且
\begin{equation*}
\left(T^* v\right)(x)=\int_{\Omega} K(y, x) v(y) \dif \mu(y) \quad\left(\forall v \in L^2(\Omega, \mu)\right)
\end{equation*}这是因为
\begin{equation*}
\begin{aligned}
\|T u\|^2 & =\int_{\Omega}\left|\int_{\Omega} K(x, y) u(y) \dif \mu(y)\right|^2 \dif \mu(x) \\
& \leqslant \int_{\Omega}\left(\int_{\Omega}|K(x, y)|^2 \dif \mu(y) \int_{\Omega}|u(y)|^2 \dif \mu(y)\right) \dif \mu(x) \\
& \leqslant\left(\iint_{\Omega \times \Omega}|K(x, y)|^2 \dif \mu(x) \dif \mu(y)\right)\|u\|^2 \quad\left(\forall u \in L^2(\Omega, \mu)\right)
\end{aligned}
\end{equation*}
其中 $\|\cdot\|$ 表示 $L^2(\Omega, \mu)$ 上的范数, 以及
\begin{equation*}
\begin{aligned}
\left\langle T^* v, u\right\rangle & =\langle v, T u\rangle \\
& =\int_{\Omega}\left(\int_{\Omega} K(x, y) u(y) \dif \mu(y)\right) v(x) \dif \mu(x) \\
& =\int_{\Omega \times \Omega} K(x, y) u(y) v(x) \dif \mu(x) \dif \mu(y) \\
& =\left(\int_{\Omega} K(x, y) v(x) \dif \mu(x)\right) u(y) \dif \mu(y)
\end{aligned}
\end{equation*}
$\left(\forall u, v \in L^2(\Omega, \mu)\right)$. 所以
\begin{equation*}
\left(T^* v\right)(y)=\int_{\Omega} K(x, y) v(x) \dif \mu(x) \quad\left(\forall v \in L^2(\Omega, \mu)\right)
\end{equation*}
注 理由 (1) 是因为
\begin{equation*}
\begin{aligned}
& \iint_{\Omega \times \Omega}|K(x, y)\|u(y)\| v(x)| \dif \mu(x) \dif \mu(y) \\
& \quad \leqslant\|u\| \cdot\|v\|\left(\iint_{\Omega \times \Omega}|K(x, y)|^2 \dif \mu(x) \dif \mu(y)\right)^{\frac{1}{2}}<\infty
\end{aligned}
\end{equation*}
所以可以应用 Fubini 定理.
\end{example}
\begin{theorem}[Young 不等式]设 $f \in L^p( \mathbb{R} )(1 \leqslant p \leqslant \infty)$, $K \in L^1( \mathbb{R} )$, 则
\begin{equation*}
\|K * f\|_p \leqslant\|K\|_1 \cdot\|f\|_p
\end{equation*}
其中 $\|\cdot\|_p$ 表示 $L^p( \mathbb{R} )$ 的范数 $(1 \leqslant p \leqslant \infty)$.
\end{theorem}

\begin{proof}
     当 $1<p<\infty$ 时. 由 Hölder 不等式,
\begin{equation*}
\begin{aligned}
& \left|\int_{-\infty}^{\infty} K(x-y) f(y) \dif y\right| \\
& \quad \leqslant \int_{-\infty}^{\infty}|K(x-y)|^{\frac{1}{q}} \cdot|K(x-y)|^{\frac{1}{p}}|f(y)| \dif y \\
& \quad \leqslant\left(\int_{-\infty}^{\infty}|K(x-y)| \dif y\right)^{\frac{1}{q}}\left(\int_{-\infty}^{\infty}|K(x-y)| \cdot|f(y)|^p \dif y\right)^{\frac{1}{p}}
\end{aligned}
\end{equation*}
而
\begin{equation*}
\begin{aligned}
& \int_{-\infty}^{\infty}\left|\int_{-\infty}^{\infty} K(x-y) f(y) \dif y\right|^p \dif x \\
& \quad \leqslant\|K\|_1^{\frac{p}{q}} \int_{-\infty}^{\infty} \int_{-\infty}^{\infty}|K(x-y)| \cdot|f(y)|^p \dif y \dif x \\
& \quad=\|K\|_1^{\frac{p}{q}} \cdot\|f\|_p^p \cdot\|K\|_1
\end{aligned}
\end{equation*}
由 Fubini 定理, $(K * f)(x)$ a.e. 存在有限, 而且 (2.5.27) 式成立.
当 $p=1$ 或 $\infty$ 时, 不等式 $(2.5 .27)$ 是显然的.
\end{proof}


\begin{definition}
$X$ 是赋范空间，称 $\left\{x_n\right\}_{n=1}^{\infty} \subset X$ 弱收敛于 $x_0 \in X$ 是指：

\begin{equation*}
f\left(x_n\right) \rightarrow f\left(x_0\right), \forall f \in X^*
\end{equation*}
记为 $x_n \xrightarrow{w} x_0$ 或者 $x_n \rightarrow x_0$, 称 $x_0$ 为 $\left\{x_n\right\}_{n=1}^{\infty}$ 的弱极限。\\
依范数拓扑收敛即为强收敛。


\end{definition}


\begin{proposition}
    若收敛存在必唯一，强收敛 $\Rightarrow$ 弱收敛。
\end{proposition}

\begin{proof}
若有 $x_n \rightharpoonup x, x_n \rightharpoonup y(n \rightarrow \infty)$, 由定义推得
\begin{equation*}
f(x)=\lim _{n \rightarrow \infty} f\left(x_n\right)=f(y) \quad\left(\forall f \in X ^*\right)
\end{equation*}
利用推论 2.3, 即得 $x=y$.
    \begin{equation*}
\left\|x_n-x_0\right\| \rightarrow 0 \Rightarrow\left|f\left(x_n\right)-f\left(x_0\right)\right|=\left|f\left(x_n-x_0\right)\right| \leqslant\|f\| \cdot\left\|x_n-x_0\right\| \rightarrow 0, \forall f \in X^*
\end{equation*}
\end{proof}

\begin{proposition}
$\operatorname{dim} X<\infty \Rightarrow$ 弱收敛与强收敛等价。

\end{proposition}
\begin{proof}
设 $\operatorname{dim}(X)=m$, 设 $\left\{e_1, \cdots, e_m\right\}$ 是 $X$ 的一组基，由 HBT（习题2.4.7）知存在对偶基 $f_1, \cdots, f_m \in X^*$ 满足
\begin{equation*}
f_k\left(e_j\right)=\delta_{k j}, 1 \leqslant k, j \leqslant m
\end{equation*}
设 $x_n \xrightarrow{w} x_0$ ，即
\begin{equation*}
\sum_{j=1}^m \alpha_j^{(n)} e_j \xrightarrow{w} \sum_{j=1}^m \alpha_j^{(0)} e_j
\end{equation*}
于是
\begin{equation*}
f_k\left(x_n\right) \rightarrow f_k\left(x_0\right), k=1,2, \cdots, m
\end{equation*}
因此 $\left\|x_n-x_0\right\|_{\infty} \rightarrow 0$ with $\|x\|_{\infty}=\max _{1 \leqslant k \leqslant m}\left|\alpha_i\right|$, 由有限维空间范数等价可得 $\left\|x_n-x_0\right\| \rightarrow 0$.
\end{proof}
\begin{example}[弱收敛而非强收敛]
 在 $L^2[0,1]$ 中, 设 $x_n=x_n(t)=\sin n \pi t$, 则根据 Riemann-Lebesgue 引理，显然有
\begin{equation*}
\left\langle f, x_n\right\rangle=\int_0^1 f(t) \sin n \pi t \dif t \rightarrow 0 \quad\left(\forall f \in L^2[0,1]\right)
\end{equation*}
即 $x_n \rightharpoonup \theta(n \rightarrow \infty)$ 。但 $\left\|x_n\right\|=1 / \sqrt{2}$, 不可能有 $x_n \rightarrow \theta(n \rightarrow \infty)$.
这表明弱收敛确实与强收敛不同. 
\end{example}
\begin{theorem}[Mazur]
若 $x_n \xrightarrow{w} x_0$, 则 $x_0 \in \overline{\operatorname{conv}\left(x_n: n \in \infty\right)}$.
\end{theorem}
\begin{proof}
令 $C=x_0 \in \overline{\operatorname{conv}\left(x_n: n \in \infty\right)}$, 若 $x_0 \notin C$, 则由 Ascoli 定理, 存在 $f \in X^*, \alpha \in \mathbb{R}$ 使得
\begin{equation*}
\sup _{x \in C} f(x)<\alpha<f\left(x_0\right) \Rightarrow f\left(x_n\right)<\alpha<f\left(x_0\right), \forall n \in N
\end{equation*}
因此 $f\left(x_n\right) \nrightarrow f\left(x_0\right)$, 矛盾.
\end{proof}

\begin{definition}[弱 * 收敛]
若对任意 $x \in X$ 有 $f_n(x) \rightarrow f(x)$，则称 $\{f_n: n \in \mathbb{N}\} \subset X^*$ 弱 * 收敛到 $f \in X^*$，记为 $f_n \xrightarrow{w^*} f$。
\end{definition}

\begin{proposition}[三种收敛性的关系]
在 $X^*$ 中，强收敛 $\Rightarrow$ 弱收敛 $\Rightarrow$ 弱 * 收敛。
\end{proposition}

\begin{proof}
若 $f_n \xrightarrow{w} f$，则对任意 $\Lambda \in X^{**}$ 有 $\Lambda(f_n) \rightarrow \Lambda(f)$，特别地，对任意 $x \in X$，有
\[
f_n(x) = x^{**}(f_n) \rightarrow x^{**}(f) = f(x) \Rightarrow f_n \xrightarrow{w^*} f.
\]
\end{proof}

\begin{proposition}[自反空间中的收敛等价性]
若 $X$ 自反，则 $X^*$ 中弱收敛与弱 * 收敛等价。
\end{proposition}

\begin{theorem}[弱收敛的等价刻画]
$x_n \xrightarrow{w} x_0$ 当且仅当 $\sup_n \|x_n\| < \infty$，并且存在 $X^*$ 的稠密子空间 $F$，使得 $f(x_n) \rightarrow f(x_0), \forall f \in F$。
\end{theorem}

\begin{proof}
$x_n \xrightarrow{w} x_0$ 当且仅当对任意 $f \in X^*$，有 $f(x_n) \rightarrow f(x_0)$，也即 $x_n^{**}(f) \rightarrow x_0^{**}(f)$。由于 $X^*$ 完备，因此由 Banach-Steinhaus 定理，这当且仅当
\[
\sup_n \|x_n\| = \sup_n \|x_n^{**}\| < \infty
\]
且存在稠密子空间 $F \subset X^*$，使得 $x_n^{**}(f) \rightarrow x_0^{**}(f), \forall f \in F$，即
\[
f(x_n) \rightarrow f(x_0), \forall f \in F.
\]
\end{proof}

\begin{definition}[一致收敛、强收敛、弱收敛]
设 $X , Y$ 是 Banach 空间。又设 $T_n (n=1,2,\cdots)$, $T \in L(X, Y)$。
\begin{enumerate}
    \item 若 $\|T_n - T\| \to 0$，则称 $T_n$ 一致收敛于 $T$，记作 $T_n \rightrightarrows T$。此时，$T$ 称作 $\{T_n\}$ 的一致极限。
    \item 若 $\|(T_n - T)x\| \to 0 \ (\forall x \in X)$，则称 $T_n$ 强收敛于 $T$，记作 $T_n \to T$。此时，$T$ 称作 $\{T_n\}$ 的强极限。
    \item 如果对 $\forall x \in X$ 以及 $\forall f \in Y^*$，都有
    \[
    \lim_{n \to \infty} f(T_n x) = f(T x),
    \]
    则称 $T_n$ 弱收敛于 $T$，记作 $T_n \rightharpoonup T$。此时，$T$ 称作 $\{T_n\}$ 的弱极限。
\end{enumerate}
显然，一致收敛 $\Rightarrow$ 强收敛 $\Rightarrow$ 弱收敛，且每种极限若存在必唯一，但反过来一般不成立。
\end{definition}

\begin{example}[强收敛而不一致收敛]
在空间 $l^2$ 上考察左推移算子
\[
T: x = (x_1, x_2, \cdots, x_n, \cdots) \mapsto T x = (x_2, x_3, \cdots, x_n, \cdots).
\]
令 $T_n \triangleq T^n$，则有
\[
T_n x = (x_{n+1}, x_{n+2}, \cdots) \quad (\forall x = (x_1, x_2, \cdots, x_n, \cdots) \in l^2).
\]
下面我们证明 $T_n \to 0$，但 $T_n \not\rightrightarrows 0$（即不一致收敛）。事实上，取
\[
e_n = (\underbrace{0, 0, \cdots, 0, 1}_n, 0, \cdots),
\]
则有 $T_n e_{n+1} = e_1$，并且 $\|e_n\| = 1 \ (\forall n \in \mathbb{N})$。因此
\[
\|T_n\| \geqslant \|T_n(e_{n+1})\| = 1,
\]
从而 $T_n \not\rightrightarrows 0$。但对于 $\forall x = (x_1, x_2, \cdots, x_n, \cdots) \in l^2$，有
\[
\|T_n x\| = \left(\sum_{i=1}^\infty |x_{n+i}|^2\right)^{\frac{1}{2}} \to 0 \quad (n \to \infty),
\]
即 $T_n \to 0$。
\end{example}

\begin{example}[弱收敛而不强收敛]
在空间 $l^2$ 上考察右推移算子
\[
S: x = (x_1, x_2, \cdots, x_n, \cdots) \mapsto S x = (0, x_1, x_2, \cdots, x_n, \cdots).
\]
令 $S_n \triangleq S^n$，则有
\[
S_n x = (\underbrace{0, 0, \cdots, 0}_n, x_1, x_2, \cdots) \quad (\forall x \in l^2).
\]
显然，$\|S_n x\| = \|x\| \ (\forall x \in l^2)$，从而 $S_n \nrightarrow 0$。但对于 $\forall f = (y_1, y_2, \cdots, y_n, \cdots) \in (l^2)^* = l^2$，有
\[
\begin{aligned}
\left|\langle f, S_n x \rangle\right| &= \left|\sum_{i=1}^\infty y_{i+n} x_i\right| \\
&\leqslant \left(\sum_{i=1}^\infty |y_{i+n}|^2\right)^{\frac{1}{2}} \|x\| \to 0 \quad (n \to \infty),
\end{aligned}
\]
即 $S_n \rightharpoonup 0 \ (n \to \infty)$。
\end{example}
\begin{definition}[弱列紧 \& 弱 * 列紧]
- 称 $M \subset X$ 弱列紧，若 $M$ 中任意序列都有弱收敛子列。\\
- 称 $F \subset X^*$ 弱 * 列紧，若 $F$ 中任意序列都有弱 * 收敛子列。
\end{definition}

\begin{theorem}[可分 Banach-Alaoglu 定理]
若 $X$ 可分，则 $X^*$ 中的有界集是弱 * 列紧的。
\end{theorem}

\begin{proof}
设 $\{x_n\} \subset X$ 为可数稠密子集，$\{f_n\} \subset X^*$ 有界，记 $C = \sup_n \|f_n\| < \infty$，则对任意 $m$，$\{f_n(x_m) : n \in \mathbb{N}\}$ 为有界数列，依次有收敛子列 $\{f_{n_{mj}}(x_m) : j \in \mathbb{N}\}, \forall m \in \mathbb{N}$（先取子列，再取子列的子列，以此类推），抽对角线 $\{f_{n_{mm}}(x_m) : m \in \mathbb{N}\}$，断言存在 $f \in X^*$，使得 $f_{n_{mm}} \xrightarrow{w} f$，重新编号 $f_{n_m} := f_{n_{mm}}$。\\
对任意 $x \in X$，存在 $x_m$ 使得 $\|x - x_m\| < \frac{\varepsilon}{3C}$，当 $k$ 充分大时，
\[
\begin{aligned}
\left| f_{n_{k+p}}(x) - f_{n_k}(x) \right| &\leqslant \left| f_{n_{k+p}}(x) - f_{n_{k+p}}(x_m) \right| + \left| f_{n_{k+p}}(x_m) - f_{n_k}(x_m) \right| + \left| f_{n_k}(x_m) - f_{n_k}(x) \right| \\
&\leqslant C \|x - x_m\| + \frac{\varepsilon}{3} + C \|x - x_m\| \\
&< \varepsilon
\end{aligned}
\]
因此 $\{f_{n_k}(x)\}$ 为 Cauchy 列，记 $f$ 为其逐点极限，则 $f_{n_k} \xrightarrow{w^*} f$，并且
\[
|f(x)| \leqslant \sup_n |f_n(x)| \leqslant C \|x\| \Rightarrow f \in X^*.
\]
得证。
\end{proof}

\begin{theorem}[Alaoglu 定理]
设 $X$ 为赋范空间，则 $X^*$ 中单位闭球是弱 * 紧的。
\end{theorem}

\begin{lemma}
设赋范空间 $X$ 中 $x_n \xrightarrow{w} x_0$，则
\[
\|x_0\| \leqslant \liminf_{n \to \infty} \|x_n\|.
\]
\end{lemma}

\begin{proof}
左式非负, 所以 $x_0=0$ 时成立。下设 $x_0 \neq 0$, 则存在 $f \in X^*$ s.t. $\|f\|=1$ s.t. $f\left(x_0\right)=\left\|x_0\right\|$, 从而
\begin{equation*}
\left\|x_0\right\|=\lim _{n \rightarrow \infty}\left|f\left(x_n\right)\right| \leqslant \liminf _{n \rightarrow \infty}\|f\| \cdot\left\|x_n\right\|=\liminf _{n \rightarrow \infty}\left\|x_n\right\|
\end{equation*}

\end{proof}

\begin{theorem}[Eberlein-Smulian 定理]
若 $X$ 为自反空间，则
\begin{enumerate}
  \item $X$ 中有界集是弱列紧的。
  \item $X$ 中单位闭球是弱自列紧的。
\end{enumerate}
\end{theorem}

\begin{proof}
1. 只需证明对任意 $R > 0, \overline{B(0, R)}$ 是弱列紧的。设 $\{x_n\} \subset X$, $\|x_n\| \leqslant R$，记 $Y = \overline{\operatorname{Span}\{x_n : n \in \mathbb{N}\}}$，则 $Y$ 可分，并且 $X$ 的自反性蕴含了闭子空间 $Y$ 的自反性（Petti），进而 $Y^{* *} = Y$ 可分 $\Rightarrow Y^*$ 可分，由可分 Banach-Alaoglu 定理，$Y^{* *}$ 中的有界集是弱 * 列紧的，$\|x_n^{* *}\| = \|x_n\| \leqslant R$，说明 $\{x_n^{* *}\}$ 有弱 * 收敛子列 $x_{n_k}^{* *} \xrightarrow{w^*} x_0^{* *} \in Y^{* *}$，即对任意 $f \in Y^*$ 有
\[
f(x_n) = x_n^{* *}(f) \to x_0^{* *}(f) = f(x_0).
\]
对任意 $F \in X^*, F(x_{n_k}) = F|_Y(x_{n_k}) \to F|_Y(x_0) = F(x_0)$，因此 $x_{n_k} \xrightarrow{w} x_0$，得证。\\
2. 设 $x_0$ 为其某个子列的弱收敛极限，则
\[
\|x_0\| \leqslant \liminf_{k \to \infty} \|x_{n_k}\| \leqslant 1 \Rightarrow x_0 \in \overline{B(0,1)}.
\]
\end{proof}
2.5： 1、3、8、10、11、14、15、18、19、20、22、23
\begin{problem}
  求证: $\left({l}^p\right)^* = {l}^q \quad (1 \leqslant p < \infty, \frac{1}{p} + \frac{1}{q} = 1)$.
\end{problem}
\begin{proof}
对于 $\forall x=\left\{\xi_k\right\}_{k=1}^{\infty} \in l^p, y=\left\{\eta_k\right\}_{k=1}^{\infty} \in l^q$, 由 Holder 不等式:
\begin{equation*}
\left|\sum_{k=1}^{\infty} \xi_k \eta_k\right| \leqslant\left(\sum_{k=1}^{\infty}\left|\xi_k\right|^p\right)^{\frac{1}{p}} \cdot\left(\sum_{k=1}^{\infty}\left|\eta_k\right|^q\right)^{\frac{1}{q}} \leqslant\|x\|_p \cdot\|y\|_q
\end{equation*}
定义
\begin{equation*}
T_y: l^p \rightarrow \mathbb{K} , x=\left\{\xi_k\right\}_{k=1}^{\infty} \mapsto \sum_{k=1}^{\infty} \xi_k \eta_k
\end{equation*}
则 $T_y$ 有界且 $\left\|T_y\right\| \leqslant\|y\|_q$, 下面证明映射 $\Lambda: y \mapsto T_y$ 是等距同构。
对于 $T \in\left(l^p\right)^*$ ，令 $e_k=(0, \cdots, 0,1,0, \cdots)$ ，则
\begin{equation*}
T(x)=T\left(\sum_{k=1}^{\infty} \xi_k e_k\right)=\sum_{k=1}^{\infty} \xi_k T\left(e_k\right), \forall x=\left\{\xi_k\right\}_{k=1}^{\infty} \in l^p
\end{equation*}
取 $y_T \stackrel{\text { def }}{=}\left\{T\left(e_k\right)\right\}_{k=1}^{\infty}$, 则 $\Lambda\left(y_T\right)=T$, 下面证明 $y_T \in l^p$ 且 $\left\|y_T\right\|_q \leqslant\|T\|$, 从而映射 $\Lambda$ 是满射且保距, 进而是等距同构。\\
若 $p>1$ ，设 $a_k=\left|T\left(e_k\right)\right|^{q-1} e ^{- i \cdot \arg T\left(e_k\right)}, z_n=\left(a_1, a_2, \cdots, a_n, 0,0, \cdots\right) \in l^p$ ，于是
\begin{equation*}
\begin{aligned}
\sum_{k=1}^n\left|T\left(e_k\right)\right|^q=\sum_{k=1}^n a_k \cdot T\left(e_k\right) & =T\left(z_n\right) \\
& \leqslant\|T\| \cdot\left\|z_n\right\|_p \\
& =\|T\| \cdot\left(\sum_{k=1}^n\left|a_k\right|^p\right)^{\frac{1}{p}} \\
& =\|T\| \cdot\left(\sum_{k=1}^n\left|T\left(e_k\right)\right|^{p(q-1)}\right)^{\frac{1}{p}} \\
& =\|T\| \cdot\left(\sum_{k=1}^n\left|T\left(e_k\right)\right|^q\right)^{\frac{1}{p}}
\end{aligned}
\end{equation*}
从而
\begin{equation*}
\left(\sum_{k=1}^n\left|T\left(e_k\right)\right|^q\right)^{1-\frac{1}{p}}=\left(\sum_{k=1}^n\left|T\left(e_k\right)\right|^q\right)^{\frac{1}{q}} \leqslant\|T\|, \forall n
\end{equation*}
令 $n \rightarrow \infty$ 可得 $\left\|y_T\right\|_q \leqslant\|T\|$.\\
若 $p=1, q=\infty$,
\begin{equation*}
\begin{aligned}
\left\|y_T\right\|_{\infty} & =\sup _{n \geqslant 1}\left|T\left(e_n\right)\right| \\
& =\sup _{n \geqslant 1} T\left(e_n\right) \cdot e^{-i \cdot \arg T\left(e_n\right)} \\
& =\sup _{n \geqslant 1} T\left(\left\{e^{-i \cdot \arg T\left(e_n\right)} \cdot \delta_{k n}\right\}_{k=1}^{\infty}\right) \leqslant\|T\|
\end{aligned}
\end{equation*}
最后一个不等号来自
\begin{equation*}
\|T\|=\sup _{\|x\|=1}\|T x\|
\end{equation*}
所以题目得证。
\end{proof}
\begin{problem}
  设 ${C}$ 是收敛数列的全体，赋以范数
\[
  \|\cdot\| : \left\{\xi_k\right\} \in C \mapsto \sup_{k \geqslant 1} |\xi_k|,
  \]
  求证: ${C}^* = {l}^1$.
\end{problem}
\begin{proof}
    记
    \begin{equation*}
    e_0=(1,1, \cdots), e_k=(0, \cdots, 0, \underset{k}{1}, 0, \cdots) \in C
    \end{equation*}
    对于 $x=\left\{x_n\right\} \in C, x_n \rightarrow x_0$, 设 $a=\{a_n \}\in l^1$, 定义映射
    \begin{equation*}
    T_a: C \rightarrow \mathbb{C} , x \mapsto a_1 x_0+\sum_{n=1}^{\infty} a_{n+1} x_n
    \end{equation*}
    右侧级数是收敛的，因为
    \begin{equation*}
    \sum_{n=1}^{\infty}\left|a_{n} x_n\right| \leqslant \sup _{n \geqslant 1}\left|x_n\right| \cdot \sum_{n=1}^{\infty}\left|a_{n}\right| \leqslant\|x\| \cdot\|a\|_1<+\infty
    \end{equation*}
    那么定义映射 $\Lambda: l^1 \rightarrow C^*, a \mapsto T_a$ ，只需证明 $\Lambda$ 是满射。\\
    对于任意 $T \in C^*$ ，令 $a_{k}=T\left(e_k\right), z_n=\left( e ^{- i \cdot \arg \left(a_2\right)}, \cdots, e ^{- i \cdot \arg \left(a_n\right)}, 0,0, \cdots\right) \in C$ ，于是
    \begin{equation*}
    \sum_{k=1}^n\left|a_k\right|=\sum_{k=1}^n a_k \cdot e^{-i \cdot \arg \left(a_k\right)}=T\left(z_n\right) \leqslant\|T\| \cdot\left\|z_n\right\|=\|T\|<+\infty, \forall n
    \end{equation*}
    这说明级数 $\sum_{k=1}^{\infty} a_k$ 收敛, 
    则令 $a=\left\{a_n\right\} \in l^1$, 且
    \begin{equation*}
    T_a(x)=\sum_{k=1}^{\infty} a_k x_{k}=T(x)
    \end{equation*}
    得证。
    \end{proof}
\begin{problem}
  设 ${C_0}$ 是以 0 为极限的数列全体，赋以范数
\[
  \|\cdot\| : \left\{\xi_k\right\} \in {C} \mapsto \sup_{k \geqslant 1} |\xi_k|,
  \]
  求证: ${C_0}^* ={l}^1$。
\end{problem}
\begin{proof}
\textbf{第一部分：} 对于任意的 $\eta = \{\eta_k\} \in l^1$，我们构造一个线性泛函 $f \in (C_0)^*$，设定标准基向量：
\[
e_k = (\overbrace{0, \cdots, 0,1}^{k \text{个}}, 0, 0, \cdots) \in C_0.
\]
对于任意的 $x \in C_0$，可以表示为
\[
x = \sum_{k=1}^{\infty} x_k e_k.
\]
令  
\[
f(x) = \sum_{k=1}^{\infty} \eta_k x_k.
\]
\[|f(x)|=|\sum_{k=1}^{\infty} \eta_k x_k|\leqslant \sup_{k \geqslant 1} |x_k|\left(\sum_{k=1}^{\infty}|\eta_k|\right)=\|\eta\|_{l^1}\|x\|\]
所以$f\in(l^1)^*$，且$\|f\|\leqslant\|\eta\|_{l^1}$。\\
\textbf{第二部分：} 对 $\forall f \in (C_0^*)$，令$f(e_k)=\eta_k,(k=1,2,\cdots),\eta=\{\eta_k\}$，下面证明$\eta \in l^1,f(x)=\sum\limits_{k=1}^\infty x_k\eta_k,\|f\|=\|\eta\|_{l^1}$。考虑任意的 $n \in \mathbb{N}$，构造无穷维向量 $x^{(n)}$，其坐标为
\[
x_k^{(n)} = \begin{cases} 
 e^{-i \theta_k}, & k \leqslant n, \\
0, & k > n,
\end{cases}
\]
其中 $\theta_k = \arg(\eta_k)$。由于对于每一个固定的 $n$，$x^{(n)}$ 的坐标只有有限项不为 0，故 $x^{(n)} \in l^1$。
\[
 f( x^{(n)}) = \sum_{k=1}^{\infty} f(e_k) x_k^{(n)} = \sum_{k=1}^{\infty} \eta_k  e^{-i \theta_k}.
\]
计算右边：
\[
| f( x^{(n)}) | = \sum_{k=1}^n  |\eta_k| e^{i \theta_k} \cdot e^{-i \theta_k} = \sum_{k=1}^n |\eta_k|.
\]
另一方面，利用 $f$ 的连续性得到
\[
| f(x^{(n)}) |\leqslant \|f\| \|x^{(n)}\|.
\]
计算 $x^{(n)}$ 的范数：
\[
\|x^{(n)}\| = \sup_{1\leqslant k\leqslant n} |e^{-i\theta_k}|=1
\]
因此，
\[
| f(x^{(n)}) | \leqslant \|f\| 
\]
结合上述两式，我们得到
\[
\sum_{k=1}^n |\eta_k| \leqslant \|f\| 
\]
所以，$\eta = \{\eta_1\} \in l^1$，且 $\|\eta\|_{l^1} \leqslant \|f\|$。\\
令$Tf=\eta,T:(C_0)^*\rightarrow l^1$，则$\|\eta\|_{l^1}=\|f\|$，且$T$是等距满射，所以 $(C_0)^*=l^1$。
\end{proof}
\begin{problem}
  求证：有限维 ${B}^*$ 空间必是自反的。
\end{problem}
\begin{proof}
设 $\left\{x_k\right\}_{k=1}^n$ 是 $X$ 的一组基, 则由习题 2.4.7 可知
\begin{equation*}
\exists\left\{f_k\right\}_{k=1}^n \subset X^* \text { s.t. }\left\langle f_i, x_j\right\rangle=\delta_{i j} \quad(i, j=1,2, \cdots, n)
\end{equation*}
于是对 $\forall f \in X^*, \forall x=\sum_{k=1}^n \lambda_k x_k$,
\begin{equation*}
f(x)=f\left(\sum_{k=1}^n \lambda_k x_k\right)=\sum_{k=1}^n \lambda_k f\left(x_k\right)=\sum_{k=1}^n f_k(x) f\left(x_k\right)
\end{equation*}
此即说明,
\begin{equation*}
f=\sum_{k=1}^n f\left(x_k\right) f_k
\end{equation*}
现对 $\forall y \in X^{* *}$,
\begin{equation*}
y(f)=y\left(\sum_{k=1}^n f\left(x_k\right) f_k\right)=\sum_{k=1}^n f\left(x_k\right) y\left(f_k\right)=f\left(\sum_{k=1}^n x_k y\left(f_k\right)\right)
\end{equation*}
于是 $x_0 \stackrel{\text { def }}{=} \sum_{k=1}^n x_k y\left(f_k\right) \in X$ 使得 $x_0^{* *}=y$, 这意味着从 $x$ 到 $x^{* *}$ 的自然映射是满的, 所以 $X$ 是自反的.
\end{proof}
\begin{problem}
  求证：${B}$ 空间是自反的，当且仅当它的共轭空间是自反的。
\end{problem}
\begin{proof}
设 $X$ 是 Banach 空间, 从 $X$ 到 $X^{* *}$ 的自然映射是 $T$.\\
必要性：即证明从 $X^*$ 到 $X^{* * *}$ 的自然映射是满的，考虑 $y \in X^{* * *}$ ，取
\begin{equation*}
f: X \rightarrow \mathbb{K} , x \mapsto y(T(x))=y\left(x^{* *}\right)
\end{equation*}
于是 $\forall x^{* *} \in X^{* *}$,
\begin{equation*}
y\left(x^{* *}\right)=f(x)=x^{* *}(f)=f^{**}(x^{**})
\end{equation*}
这说明 $f^{* *}=y$, 得证。\\
充分性：设 $X^*$ 自反，则由必要性知 $X^{* *}$ 自反，又因为 $X$ ，作为 $X^{* *}$ 的子空间（因为自然映射是嵌入）是闭的（因为 $X$ 是 Banach 空间），自反空间的闭子空间也自反（Pettis），从而 $X$ 自反。
\end{proof}
\begin{problem}
  设 ${X}$ 是 ${B}^*$ 空间， $T$ 是从 ${X}$ 到 ${X}^{**}$ 的自然映射，求证: $R(T)$ 是闭的充要条件是 ${X}$ 是完备的。
\end{problem}

\begin{problem}
  在 ${l}^1$ 中定义算子
  \[
  T : (x_1, x_2, \cdots, x_n, \cdots) \mapsto (0, x_1, x_2, \cdots, x_n, \cdots),
  \]
  求证: $T \in {L}({l}^1)$ 并求 $T^*$。
\end{problem}

\begin{problem}
  在 ${l}^2$ 中定义算子
  \[
  T : (x_1, x_2, \cdots, x_n, \cdots) \mapsto (x_1, \frac{x_2}{2}, \cdots, \frac{x_n}{n}, \cdots),
  \]
  求证: $T \in {L}({l}^2)$ 并求 $T^*$。
\end{problem}
\begin{proof} 对 $x=\left(x_1, x_2, \cdots, x_n, \cdots\right) \in l^2$, 有
\begin{equation*}
\|T x\|_{l^2}^2=\sum_{k=1}^{\infty}\left|\frac{x_k}{k}\right|^2 \leqslant \sum_{k=1}^{\infty}\left|x_k\right|^2=\|x\|^2
\end{equation*}
从而 $T \in L \left(l^2\right)$ 。\\
$l^2$ 是 Hilbert 空间, 由 Riesz 表示定理, $\forall f \in\left(l^2\right)^*=l^2$, 存在 $y_f=\left\{f_n\right\}_{n=1}^{\infty}$ 使得
\begin{equation*}
f=\left\langle\cdot, y_f\right\rangle
\end{equation*}
$T^*$ 满足：
\begin{equation*}
T^* f(x)=f(T(x))=\left\langle T(x), y_f\right\rangle=\sum_{n=1}^{\infty} \frac{f_n x_n}{n} \Rightarrow T^* f=\left\langle\cdot,\left\{\frac{f_n}{n}\right\}_{n=1}^{\infty}\right\rangle
\end{equation*}
所以 $T^*:\left\{f_n\right\}_{n=1}^{\infty} \mapsto\left\{f_n / n\right\}_{n=1}^{\infty}$, 也就是 $T^*=T$.
\end{proof}
\begin{problem}
  设 ${H}$ 是 Hilbert 空间， $A \in L({H})$ 并满足
  \[
  (A x, y) = (x, A y) \quad (\forall x, y \in {H}),
  \]
  求证:
  \begin{enumerate}
    \item $A^* = A$;
    \item 若 $R(A)$ 在 ${H}$ 中稠密, 则方程 $A x = y$ 对 $\forall y \in R(A)$ 存在唯一解。
  \end{enumerate}
\end{problem}
\begin{proof}
(1). 对于 $f \in H^*$, 由 Riesz 表示定理, 存在 $y_f$ 使得 $f=\left\langle\cdot, y_f\right\rangle$, 并且 $\|f\|=\left\|y_f\right\|$, 因此 $f \mapsto y_f$ 就是 $H^*$ 到 $H$的等距同构。因此： $\forall f \in H^*, \forall x \in H$,
\begin{equation*}
A^* f(x)=f(A(x)) \Rightarrow\left\langle x, y_{A^* f}\right\rangle=\left\langle A x, y_f\right\rangle=\left\langle x, A y_f\right\rangle \Rightarrow A y_f=y_{A^* f}
\end{equation*}
而 $A^*: f \mapsto A^* f$ 对应了 $y_f \mapsto y_{A^* f}$, 所以 $A^* y_f=y_{A^* f}=A y_f \Rightarrow A^*=A$.\\
(2). 对 $\forall y \in R(A)$ ，若
\begin{equation*}
A x_1=y=A x_2
\end{equation*}
则 $\forall z \in H$,
\begin{equation*}
0=\left\langle A\left(x_1-x_2\right), z\right\rangle=\left\langle x_1-x_2, A z\right\rangle
\end{equation*}
由于 $R(A)$ 在 $H$ 中稠密,
\begin{equation*}
\exists z_n \in H \text { s.t. } A z_n \rightarrow x_1-x_2
\end{equation*}
从而
\begin{equation*}
\left\|x_1-x_2\right\|^2=\lim _{n \rightarrow \infty}\left\langle x_1-x_2, A z_n\right\rangle=0 \Rightarrow x_1=x_2
\end{equation*}
\end{proof}
\begin{problem}
  设 ${X}, {Y}$ 是 ${B}^*$ 空间， $A \in {L}({X}, {Y})$，又设 $A^{-1}$ 存在且 $A^{-1} \in {L}({Y}, {X})$，求证:
  \begin{enumerate}
    \item $\left(A^*\right)^{-1}$ 存在，且 $\left(A^*\right)^{-1} \in {L}({X}^*, {Y}^*)$；
    \item $\left(A^*\right)^{-1} = \left(A^{-1}\right)^*$。
  \end{enumerate}
\end{problem}
\begin{proof}
先证 $A^*$ 是单射. 事实上, $\forall f \in {Y} ^*$,
\begin{equation*}
\begin{array}{cc}
A^*f=0\Rightarrow\langle f, y\rangle=\langle f, A x\rangle=\left\langle A^* f, x\right\rangle=0 \Longrightarrow\langle f, y\rangle=0, \forall y \in {Y}^*\Rightarrow f=0
\end{array}
\end{equation*}
再证 $A^*$ 是满射， $\forall x \in {X}$ ，
\begin{equation*}
\begin{array}{cc}
\left\langle A^*\left(A^{-1}\right)^* f, x\right\rangle =\left\langle\left(A^{-1}\right)^* f, A x\right\rangle=\left\langle f,\left(A^{-1}\right) A x\right\rangle= \langle f, x\rangle 

\end{array}
\end{equation*}
比较上述等式串的两端，即得 $A^*\left(A^{-1}\right)^* f=f$ ，再注意到 $\left(A^{-1}\right)^* f \in {Y} ^*$ ，即证得 $A^*$ 是满射。\\
进一步, 根据第一步结果, 我们有 $\left(A^*\right)^{-1}$ 存在, 从 $A^*\left(A^{-1}\right)^* f=$ $f$ ，两边作用 $\left(A^*\right)^{-1}$ ，即得
\begin{equation*}
\left(A^*\right)^{-1} f=\left(A^{-1}\right)^* f\left(\forall f \in X ^*\right) \Rightarrow\left(A^*\right)^{-1}=\left(A^{-1}\right)^* .
\end{equation*}
\end{proof}
\begin{problem}
  设 ${X}, {Y}, {Z}$ 是 ${B}^*$ 空间，而 $B \in {L}({X}, {Y})$ 以及 $A \in {L}({Y}, {Z})$，求证：$(A B)^* = B^* A^*$。
\end{problem}

\begin{problem}
  设 ${X}, {Y}$ 是 ${B}$ 空间， $T$ 是 ${X}$ 到 ${Y}$ 的线性算子，又设对 $\forall g \in {Y}^*, g(T x)$ 是 ${X}$ 上的有界线性泛函，求证: $T$ 是连续的。
\end{problem}

\begin{problem}
  设 $\{x_n\} \subset C[a, b]$, $x \in C[a, b]$ 且 $x_n \rightharpoonup x$ $(n \to \infty)$, 求证:
  \[
  \lim_{n \to \infty} x_n(t) = x(t) \quad (\forall t \in [a, b]) \quad \text{(点点收敛)}.
  \]
\end{problem}

\begin{problem}
  已知在 ${B}^*$ 空间中 $x_n \rightharpoonup x_0$ $(n \to \infty)$，求证:
  \[
  \varliminf_{n \to \infty} \|x_n\| \geqslant \|x_0\|.
  \]
\end{problem}
\begin{proof}
左式非负, 所以 $x_0=0$ 时成立。下设 $x_0 \neq 0$, 则存在 $f \in X^*$ s.t. $\|f\|=1$ s.t. $f\left(x_0\right)=\left\|x_0\right\|$, 从而
\begin{equation*}
\left\|x_0\right\|=\lim _{n \rightarrow \infty}\left|f\left(x_n\right)\right| \leqslant \liminf _{n \rightarrow \infty}\|f\| \cdot\left\|x_n\right\|=\liminf _{n \rightarrow \infty}\left\|x_n\right\|
\end{equation*}

\end{proof}
\begin{problem}
  设 ${H}$ 是 Hilbert 空间，$\{e_n\}$ 是 ${H}$ 的正交规范基，求证：在 ${H}$ 中 $x_n \rightharpoonup x_0$ $(n \to \infty)$ 的充要条件是
  \begin{enumerate}
    \item $\|x_n\|$ 有界;
    \item $(x_n, e_k) \to (x_0, e_k) \quad (n \to \infty), (k = 1, 2, \dots)$.
  \end{enumerate}
\end{problem}
\begin{proof}
充分性: 考虑$f_k=\langle\cdot,e_k\rangle$，由Resiz定定理知$\forall f\in H^*,f=\langle\cdot,y_f\rangle$,由$\{e_n\}$ 是 ${H}$ 的正交规范基知$y_f=\sum\limits_{k=1}^\infty \langle y_f,e_k\rangle e_k,g_n=\sum\limits_{k=1}^n \langle y_f,e_k\rangle f_k\to f$，所以
 $\operatorname{span}\left\{f_k\right\}_{k=1}^\infty$ 在 $H^*$中稠密，根据本节定理15 $ x_n \rightharpoonup 0$.。\\
必要性：令 $f=\left\langle\cdot, e_k\right\rangle \in H^*$ ，则 $2^{\circ}$ 得证；对于 $\forall n$ ，考虑 $x_n^{* *} \in H^{* *}$ ，则 $\left\|x_n^{* *}\right\|=\left\|x_n\right\|$ ，收敛列必有界，所以
\begin{equation*}
\sup _n\left|x_n^{* *}(f)\right|=\sup _n\left|f\left(x_n\right)\right|<+\infty
\end{equation*}
由 UBP 可得 $x_n^{* *}$ 一致有界, 进而 $1^{\circ}$ 得证。
\end{proof}
\begin{problem}
设 $S_n$ 是 $L^p( R )(1 \leqslant p<\infty)$ 到自身的算子:
\begin{equation*}
\left(S_n u\right)(x)= \begin{cases}u(x), & |x| \leqslant n \\ 0, & |x|>n\end{cases}
\end{equation*}
其中 $u \in L^p( R )$ 是任意的，求证： $\left\{S_n\right\}$ 强收敛于恒同算子 $I$ ，但不一致收敛到 $I$ 。
\end{problem}
\begin{problem}
设 $H$ 是 Hilbert 空间, 在 $H$ 中 $x_n \rightharpoonup x_0(h \rightarrow \infty)$, 而且 $y_n \rightarrow y_0(h \rightarrow \infty)$, 求证: $\left(x_n, y_n\right) \mapsto\left(x_0, y_0\right)(n \rightarrow \infty)$.

\end{problem}
\begin{problem}
设 $\left\{e_n\right\}$ 是 Hilbert 空间 $H$ 中的正交规范集, 求证: 在 $H$ 中 $e_n \rightharpoonup \theta(n \rightarrow \infty)$, 但 $e_n \nrightarrow \theta(n \rightarrow \infty)$.
\end{problem}
\begin{proof}
    由 Bessel 不等式，\(\forall x \in H\)，有
\begin{equation*}
\sum_{n=1}^{\infty} \left| \left(x, e_n\right) \right|^2 \leqslant \|x\|^2,
\end{equation*}
即 \(\sum_{n=1}^{\infty} \left| \left(x, e_n\right) \right|^2\) 收敛。由级数收敛的必要条件，可得 \(\left(x, e_n\right) \to 0\)，即 \(e_n \rightharpoonup 0\)。
但由于 \(e_n \perp e_m\)，对 \(n \neq m\)，有
\begin{equation*}
\|e_n - e_m\|^2 = \|e_n\|^2 + \|e_m\|^2 = 2 \nrightarrow 0,
\end{equation*}
因此 \(e_n \nrightarrow 0\)。
\end{proof}


\begin{problem}
  设 ${H}$ 是 Hilbert 空间，求证: 在 ${H}$ 中 $x_n \to x$ $(n \to \infty)$ 的充要条件是
  \begin{enumerate}
    \item $\|x_n\| \to \|x\|$ $(n \to \infty)$;
    \item $x_n \rightharpoonup x$ $(n \to \infty)$.
  \end{enumerate}
\end{problem}
\begin{proof}
必要性显然。下证充分性。\\
由内积定义，有
\begin{equation*}
\begin{aligned}
\|x_n - x\|^2 & = (x_n - x, x_n - x) \\
& = \|x_n\|^2 - (x, x_n) - (x_n, x) + \|x\|^2,
\end{aligned}
\end{equation*}
又当 \(n \to \infty\) 时，
\begin{equation*}
\left\{
\begin{aligned}
\|x_n\|^2 & \to \|x\|^2, \\
(x, x_n) & \to \|x\|^2, \\
(x_n, x) & \to \|x\|^2,
\end{aligned}
\right.
\end{equation*}
故有 \(\|x_n - x\|^2 \to 0 \ (n \to \infty)\)，即 \(x_n \to x \ (n \to \infty)\)。
\end{proof}
\begin{problem}
  求证: 在自反的 ${B}$ 空间中, 集合的弱列紧性与有界性是等价的。
\end{problem}
\begin{proof}
 有界性 \(\Longrightarrow\) 弱列紧性：由本节定理 20（Eberlein-Smulian 定理）推出。\\
弱列紧性 \(\Longrightarrow\) 有界性：用反证法。
设 \(\{x_n\}\) 弱列紧但是无界，则存在 \(\{x_n\}\) 的子列 \(\{y_n\}\) 满足 \(\|y_n\| > n\)。因为 \(\{y_n\}\) 弱列紧，所以 \(\exists\{y_{n_k}\}\) 弱收敛，即对 \(\forall f \in X^*\)，有 \(\{f(y_{n_k})\}\) 收敛。
根据共鸣定理，\(\{y_{n_k}\}\) 有界，这与 \(\|y_{n_k}\| > n_k\) 矛盾。
因此，弱列紧性 \(\Longrightarrow\) 有界性成立。
\end{proof}
\begin{problem}
  求证: ${B}^*$ 空间中的闭凸集是弱闭的, 即若 ${M}$ 是闭凸集， $\{x_n\} \subset {M}$ ，且 $x_n \to x_0$ $(n \to \infty)$ ，则 $x_0 \in {M}$。
\end{problem}

\begin{problem}
  设 ${X}$ 是自反的 ${B}$ 空间， ${M}$ 是 ${X}$ 中的有界闭凸集， $\forall f \in {X}^*$, 求证: $f$ 在 ${M}$ 上达到最大值和最小值。
\end{problem}
\begin{proof}
    设 \(b = \sup_{x \in M} f(x)\)，则对 \(\forall n \in \mathbb{N}\)，存在 \(x_n \in M\)，使得
\[
b - \frac{1}{n} < f(x_n) \leqslant b.
\]
因为 \(M\) 有界，根据本节定理 20，\(M\) 弱列紧，所以 \(\{x_n\}\) 有弱收敛子列 \(x_{n_k} \rightharpoonup x_b \in M\)，故对 \(\forall f \in X^*\)，有 \(f(x_{n_k}) \rightarrow f(x_b)\)。于是
\[
b - \frac{1}{n_k} < f(x_{n_k}) \leqslant b.
\]
令 \(k \to \infty\)，即得 \(f(x_b) = b\)。\\
同理，存在 \(x_a \in M\)，使得 \(f(x_a) = a = \inf_{x \in M} f(x)\)。
\end{proof}
\begin{problem}
  设 ${X}$ 是自反的 ${B}$ 空间， ${M}$ 是 ${X}$ 中的非空闭凸集，求证: 存在 $x_0 \in {M}$, 使得
  \[
  \|x_0\| = \inf \{\|x\| \mid x \in {M}\}.
  \]
\end{problem}
\begin{proof}
    设 \(d = \inf_{x \in M} \{\|x\|\}\)，则对 \(\forall n \in \mathbb{N}\)，存在 \(x_n \in M\)，使得
\[
d \leqslant \|x_n\| < d + \frac{1}{n} \leqslant d + 1.
\]
因为 \(\{x_n\}\) 有界，根据本节定理 20，存在子列 \(x_{n_k} \rightharpoonup x_0\)。对此 \(x_0\)，根据 §4 Hahn-Banach 定理推论 4，存在 \(f \in X^*\)，使得 \(\|f\| = 1, f(x_0) = \|x_0\|\)。于是有，一方面，
\[
x_0 \in M \Rightarrow \|x_0\| \geqslant d;
\]
另一方面，
\[
\|x_0\| = f(x_0) = \lim_{k \to \infty} f(x_{n_k}) \leqslant \lim_{k \to \infty} \|f\| \|x_{n_k}\| = d.
\]
故有 \(\|x_0\| = d = \inf_{x \in M} \{\|x\|\}\)。
\end{proof}
\section{谱}
\begin{definition}
设 $X$ 是一个复 Banach 空间，在 $L(X)$ 上引入乘法：
\begin{equation*}
(A B) x \stackrel{\text{def}}{=} A(B x)
\end{equation*}
则满足以下性质：
\begin{enumerate}
    \item 结合律：$(A B) C = A(B C)$.
    \item 分配律：$(A + B) C = A C + B C$，$A(B + C) = A B + A C$.
    \item 数量乘法：$\lambda(A B) = (\lambda A) B = A(\lambda B)$.
    \item 单位元：$A I = I A = A$.
    \item 不等式：$\|A B\| \leqslant \|A\| \cdot \|B\|$.
\end{enumerate}
由此可得 $L(X)$ 是一个 Banach 代数。
\end{definition}

\begin{definition}
称 $A \in L(X)$ 可逆，若存在 $B \in L(X)$，使得：
\begin{equation*}
A B = B A = I
\end{equation*}
\end{definition}
\begin{note}
由于 $X$ 为 Banach 空间, 因此由逆算子定理, 只要 $A \in L (X)$ 为双射则 $A^{-1} \in L (X)$.
\end{note}
\begin{definition}
对于 $A \in L(X)$，定义$A$ 的预解集（resolvent set）为：
\begin{equation*}
\rho(A) \stackrel{\text{def}}{=}\{\lambda \in \mathbb{C} : (\lambda I - A )^{-1}\in L(X)\} = \mathbb{C} \backslash \sigma(A)
\end{equation*}
其中，$\rho(A)$ 中的元素称为正则值。\\
 $A$ 的谱 $\sigma(A)$ 为：
\begin{equation*}
\sigma(A) \stackrel{\text{def}}{=}\mathbb{C}\backslash \rho(A)
\end{equation*}
其中，$\sigma(A)$ 中的元素称为谱点。
\end{definition}

\begin{definition}
若存在 $0 \neq x \in X$，使得 $A x = \lambda x$，则称 $\lambda$ 为 $A$ 的特征值，记为 $\sigma_p(A)$，即：
\begin{equation*}
\sigma_p(A) \stackrel{\text{def}}{=}\{\lambda \in \mathbb{C} : \operatorname{Ker}(\lambda I - A) \neq \{0\} \}
\end{equation*}
\end{definition}

\begin{example}
考虑有限维线性空间上的线性映射 $A \in L(\mathbb{C}^n)$，则 $\sigma(A) = \sigma_p(A) \neq \varnothing$。
\end{example}

\begin{example}
设 $A: C[0,1] \rightarrow C[0,1], u(t) \mapsto t \cdot u(t)$。我们来讨论 $A$ 是否有特征值。


\end{example}

\begin{solution}
首先，考察以下方程：
\begin{equation*}
(\lambda I - A) u = 0 \Leftrightarrow \lambda u(t) - t u(t) = 0, \quad \forall t \in [0,1]
\end{equation*}
这意味着对于所有 $t \in [0,1]$，我们必须有 $u(t) = 0$。因此，$A$ 没有非零的特征值，也即 $A$ 没有特征值。
\end{solution}

\begin{definition}
\begin{enumerate}
    \item $(\lambda I-A)^{-1}$ 不存在. 这相当于 $\lambda$ 是特征值,$\lambda \in \sigma_p(A)$.
    \item  $(\lambda I-A)^{-1}$ 存在，且值域 $R(\lambda I-A) \triangleq(\lambda I-A) D(A)= X$ 。这相当于 $\lambda$ 是正则值 (Banach 逆算子定理 ).
    \item  $(\lambda I-A)^{-1}$ 存在, $R(\lambda I-A) \neq X$ ，但 $\overline{R(\lambda I-A)}= X$ 。对于这部分 $\lambda$ ，我们称其为 $A$ 的连续谱 $\sigma_c(A)$。
    \item  $(\lambda I-A)^{-1}$ 存在，且 $\overline{R(\lambda I-A)} \neq X$ ，这部分 $\lambda$ 称为 $A$ 的剩余谱 $\sigma_r(A)$.

    
\end{enumerate}
因此有：
\begin{equation*}
\sigma(A)=\sigma_p(A) \cup \sigma_c(A) \cup \sigma_r(A)
\end{equation*}
\end{definition}

\begin{example}
设 $A: C[0,1] \rightarrow C[0,1]$，$u(t) \mapsto t \cdot u(t)$，则 $\sigma(A) = \sigma_r(A) = [0,1]$。


\end{example}
\begin{proof}
\textbf{(1) 证明：}$\mathbb{C} \backslash [0,1] \subset \rho(A)$。\\
设 $\lambda \in \mathbb{C} \backslash [0,1]$，令：
\[
T: C[0,1] \rightarrow C[0,1], \quad u(t) \mapsto \frac{1}{\lambda - t} u(t)
\]
于是，$(\lambda I - A) T = I = T(\lambda I - A)$，且：
\[
\|T u\|_{\infty} \leqslant \left[\max_{t \in [0,1]} \frac{1}{|\lambda - t|}\right] \|u\|_{\infty}
\]
因此 $(\lambda I - A)^{-1} = T \in L(X)$，从而 $\lambda \in \rho(A)$。\\
\textbf{(2) 证明：}$[0,1] \subset \sigma_r(A)$。\\
设 $\lambda \in [0,1]$，对于任意 $v \in \operatorname{Ran}(\lambda I - A)$，存在 $u \in C[0,1]$ 满足：
\[
(\lambda - t) u(t) = v(t), \quad t \in [0,1]
\]
由此可得：$v(\lambda) = 0$，即 $1 \notin \overline{\operatorname{Ran}(\lambda I - A)}$，因此 $\overline{\operatorname{Ran}(\lambda I - A)} \neq X$。
最后，由 $[0,1] \subset \sigma_r(A) \subset \sigma(A) \subset [0,1]$，得到 $\sigma(A) = \sigma_r(A) = [0,1]$。
\end{proof}
\begin{example}
设 $A: L^2[0,1] \rightarrow L^2[0,1]$，$u(t) \mapsto t \cdot u(t)$，则 $\sigma(A) = \sigma_c(A) = [0,1]$。
\end{example}

\begin{proof}
首先与上例同理，证明 $\mathbb{C} \backslash [0,1] \subset \rho(A)$。\\
接下来，证明：对于任意 $\lambda \in [0,1]$，$\operatorname{Ran}(\lambda I - A) \neq X$。假设 $1 \in \operatorname{Ran}(\lambda I - A)$，则存在 $u \in L^2[0,1]$，使得：
\[
1 = (\lambda - t) u(t), \quad t \in [0,1]
\]
从而：
\[
\frac{1}{\lambda - t} \in u(t) \in L^2[0,1]
\]
这显然是矛盾的，因此 $1 \notin \operatorname{Ran}(\lambda I - A)$，从而 $\operatorname{Ran}(\lambda I - A) \neq X$。\\
接着，证明：对于任意 $\lambda \in [0,1]$，$\overline{\operatorname{Ran}(\lambda I - A) }=X$。对于任意 $v \in L^2$ 和任意 $\varepsilon > 0$，定义：
\[
u_{\varepsilon}(t) \stackrel{\text{def}}{=} \frac{1}{\lambda - t} v(t) \cdot \chi_{[0,1] \backslash (\lambda - \varepsilon, \lambda + \varepsilon)}(t)
\]
于是 $u_{\varepsilon} \in L^2$，且：
\[
(\lambda I - A) u_{\varepsilon} = \chi_{[0,1] \backslash (\lambda - \varepsilon, \lambda + \varepsilon)} v \xrightarrow{L^2} v \quad \text{as} \quad \varepsilon \to 0^+
\]
得到 $v \in \overline{\operatorname{Ran}(\lambda I - A)}$。

\end{proof}

\begin{definition}
算子值函数：设 $A$ 是一个线性算子，定义如下：
\[
R_\lambda(A): \rho(A) \rightarrow L(X), \quad \lambda \mapsto (\lambda I - A)^{-1}
\]
称为 $A$ 的预解式 (resolvent)。
\end{definition}

\begin{lemma}
设 $T \in L(X)$ 且 $\|T\| < 1$，则有：
\begin{enumerate}
    \item $(I - T)^{-1} \in L(X)$。
    \item $(I - T)^{-1} = \sum_{n=0}^{\infty} T^n$。（Von Neumann 级数）
    \item $\| (I - T)^{-1} \| \leqslant (1 - \|T\|)^{-1}$。
\end{enumerate}
\end{lemma}

\begin{proof}
\textbf{(1)} 令
\[
S_n \stackrel{\text{def}}{=} \sum_{k=0}^n T^k
\]
则
\[
\| S_{n+p} - S_n \| = \left\| \sum_{k=n+1}^{n+p} T^k \right\| \leqslant \sum_{k=n+1}^{n+p} \| T \|^k < \frac{\| T \|^{k+1}}{1 - \| T \|}
\]
由于 $L(X)$ 完备性，存在 $S \in L(X)$ 使得 $\| S_n - S \| \to 0$。现在我们要证明：
\[
S (I - T) = (I - T) S = I
\]
\textbf{证明：} 注意到
\[
\| S_n (I - T) - I \| = \| I - T^{n+1} - I \| \leqslant \| T \|^{n+1} \to 0 \text{ as } n \to \infty
\]
因此，$\| S (I - T) - I \| \leqslant \| S - S_n \| \cdot \| I - T \| + \| S_n (I - T) - I \| \to 0$。从而 $S (I - T) = I$，同理 $(I - T) S = I$。\\
\textbf{(2)} 由上面的结论可得，$(I - T)^{-1} = \sum_{n=0}^{\infty} T^n$。\\
\textbf{(3)} 由于 $\| S \| \leqslant \sup_n \| S_n \|$，可以得出 $\| (I - T)^{-1} \| \leqslant (1 - \| T \|)^{-1}$。
\end{proof}

\begin{theorem}
$\rho(A)$ 是开集，进而 $\sigma(A)$ 是闭集。
\end{theorem}

\begin{proof}
设 $\lambda_0 \in \rho(A)$，则：
\[
\lambda I - A = \lambda_0 I - A + (\lambda - \lambda_0) I = (\lambda_0 I - A) \left[ I + (\lambda - \lambda_0) (\lambda_0 I - A)^{-1} \right]
\]
由引理 2.2，当 $|\lambda - \lambda_0| < \| (\lambda_0 I - A)^{-1} \|^{-1}$ 时，定义：
\[
B \stackrel{\text{def}}{=} \left[ I + (\lambda - \lambda_0) (\lambda_0 I - A)^{-1} \right]^{-1} \in L(X)
\]
于是：
\[
(\lambda I - A)^{-1} = B \cdot R_{\lambda_0}(A) \in L(X)
\]
因此 $\lambda \in \rho(A)$，从而得到 $\left( D \left( \lambda_0, \frac{1}{\| (\lambda_0 I - A)^{-1} \|} \right) \subset \rho(A) \right)$。
\end{proof}
\begin{proposition}
\begin{equation*}
A \in L (X) \Rightarrow \sigma(A) \subset \overline{ B(0,\|A\|)}
\end{equation*}
\end{proposition}
\begin{proof}
即证明: $\forall \lambda \in \mathbb{C}$ with $|\lambda|>\|A\|,(\lambda I-A)^{-1} \in L (X)$,
\begin{equation*}
\begin{aligned}
|\lambda|>\|A\|  \Rightarrow\left\|\frac{A}{\lambda}\right\|<1 
\stackrel{\text { Lem2.2 }}{\Rightarrow}  \left(I-\frac{A}{\lambda}\right)^{-1} \in L (X)  \Leftrightarrow(\lambda I-A)^{-1} \in L (X)
\end{aligned}
\end{equation*}
\end{proof}
\begin{corollary}
$\sigma(A)$ 是 $\mathbb{C}$ 中的紧集。
\end{corollary}

\begin{definition}
设 $X$ 是复 Banach 空间，$\Omega$ 是 $\mathbb{C}$ 上的开集。称算子值函数
\[
T: \Omega \rightarrow L(X), \quad \lambda \mapsto T_\lambda
\]
在 $\lambda_0 \in \Omega$ 处全纯是指：存在 $\lambda_0$ 的邻域 $U$，使得对任意 $\lambda \in U$，存在 $S_\lambda \in L(X)$，使得：
\[
\left\| \frac{T_{\lambda + \delta} - T_\lambda}{\delta} - S_\lambda \right\| \to 0 \text{ as } \delta \to 0
\]
\end{definition}
\begin{theorem}
设 $A$ 是闭线性算子, 则 $\rho(A)$ 是开集。
\end{theorem}

\begin{proof}
设 $\lambda_0 \in \rho(A)$, 则有：
\[
\begin{aligned}
\lambda I - A &= (\lambda - \lambda_0) I + (\lambda_0 I - A) \\
              &= (\lambda_0 I - A) \left[ I + (\lambda - \lambda_0) (\lambda_0 I - A)^{-1} \right]
\end{aligned}
\]
当 $|\lambda - \lambda_0| < \| (\lambda_0 I - A)^{-1} \|^{-1}$ 时，
\[
B \triangleq \left[ I + (\lambda - \lambda_0) (\lambda_0 I - A)^{-1} \right]^{-1} \in L(X)
\]
从而
\[
(\lambda I - A)^{-1} = B \cdot R_{\lambda_0}(A) \in L(X),
\]
即得 $\lambda \in \rho(A)$。
\end{proof}

\begin{lemma}[第一预解公式]设 $\lambda, \mu \in \rho(A)$, 则有
\[
R_\lambda(A) - R_\mu(A) = (\mu - \lambda) R_\lambda(A) R_\mu(A)
\]
\end{lemma}

\begin{proof}
直接计算：
\[
\begin{aligned}
(\lambda I - A)^{-1} &= (\lambda I - A)^{-1} (\mu I - A) (\mu I - A)^{-1} \\
                     &= (\lambda I - A)^{-1} \left[ (\mu - \lambda) I + \lambda I - A \right] (\mu I - A)^{-1} \\
                     &= (\mu - \lambda)(\lambda I - A)^{-1} (\mu I - A)^{-1} + (\mu I - A)^{-1}
\end{aligned}
\]
\end{proof}

\begin{theorem}预解式 $R_\lambda(A)$ 在 $\rho(A)$ 内是算子值解析函数。
\end{theorem}
\begin{proof}
Step1: 连续性。对于 $\forall \lambda_0 \in \rho(A)$,
\begin{equation*}
\lambda I-A=\left(\lambda_0 I-A\right)\left[I+\left(\lambda-\lambda_0\right)\left(\lambda_0 I-A\right)^{-1}\right]
\end{equation*}
当 $\left|\lambda-\lambda_0\right|<\|\left(\lambda_0 I-A\right)^{-1}| |^{-1}$ 时,
\begin{equation*}
R_\lambda(A)=\left[I+\left(\lambda-\lambda_0\right) R_{\lambda_0}(A)\right]^{-1} R_{\lambda_0}(A)
\end{equation*}
于是当 $\left|\lambda-\lambda_0\right|<\left(2 \| R_{\lambda_0}(A)| |\right)^{-1}$ 时,
\begin{equation*}
\left\|R_\lambda(A)\right\| \leqslant\left\|\left[I+\left(\lambda-\lambda_0\right) R_{\lambda_0}(A)\right]^{-1}\right\| \cdot\left\|R_{\lambda_0}(A)\right\| \leqslant \frac{1}{1-\frac{1}{2}}\left\|R_{\lambda_0}(A)\right\|=2\left\|R_{\lambda_0}(A)\right\|
\end{equation*}
由引理 1.6.2可知,
\begin{equation*}
\left\|R_\lambda(A)-R_{\lambda_0}(A)\right\| \leqslant\left|\lambda-\lambda_0\right| \cdot\left\|R_\lambda(A)\right\| \cdot\left\|R_{\lambda_0}(A)\right\| \leqslant 2\left\|R_{\lambda_0}(A)\right\|^2 \cdot\left|\lambda-\lambda_0\right|
\end{equation*}
Step2：全纯性。
\begin{equation*}
\begin{aligned}
&\left\|\frac{R_\lambda(A)-R_{\lambda_0}(A)}{\lambda-\lambda_0}+R_{\lambda_0}(A)^2\right\| \stackrel{\text { R.I. }}{=}\left\|-R_\lambda(A) R_{\lambda_0}(A)+R_{\lambda_0}(A)^2\right\| \\
& \leqslant\left\|R_{\lambda_0}(A)\right\| \cdot\left\|R_\lambda(A)-R_{\lambda_0}(A)\right\| \rightarrow 0 \text { as } \lambda \rightarrow \lambda_0
\end{aligned}
\end{equation*}

\end{proof}

\begin{theorem}
设 $A$ 是有界线性算子，则 $\sigma(A) \neq \varnothing$。
\end{theorem}

\begin{proof}
假设 $\sigma(A)=\varnothing$, 则 $\rho(A)= \mathbb{C}$, 说明 $\lambda \mapsto R_\lambda(A)$ 时算子值整函数, 于是
\begin{equation*}
\forall f \in L (X)^*, u_f(\lambda) \stackrel{\text { def }}{=} f\left(R_\lambda(A)\right), \lambda \in \mathbb{C}
\end{equation*}
是整函数，因为
\begin{equation*}
\left|\frac{u_f(\lambda)-u_f\left(\lambda_0\right)}{\lambda-\lambda_0}+f\left(R_{\lambda_0}(A)^2\right)\right| \leqslant\|f\| \cdot\left\|\frac{R_\lambda(A)-R_{\lambda_0}(A)}{\lambda-\lambda_0}+R_{\lambda_0}(A)^2\right\| \rightarrow 0 \text { as } \lambda \rightarrow \lambda_0
\end{equation*}
另一方面，当 $|\lambda|>2| | A \|$ 时，
\begin{equation*}
\left\|R_\lambda(A)\right\| \leqslant \frac{1}{|\lambda|} \frac{1}{1-\left\|\frac{A}{\lambda}\right\|}=\frac{1}{|\lambda|-\|A\|} \leqslant \frac{1}{\|A\|}
\end{equation*}
而 $\lambda \mapsto R_\lambda(A)$ 连续, 在 $\overline{ B(0,2\|A\|)}$ 上有界, 于是存在 $C>0$ 使得 $\left\|R_\lambda(A)\right\| \leqslant C, \forall \lambda \in \mathbb{C}$,
\begin{equation*}
\left|u_f(\lambda)\right| \leqslant\|f\| \cdot\left\|R_\lambda(A)\right\| \leqslant C\|f\|, \forall \lambda \in \mathbb{C}
\end{equation*}
由 Liouvillle 定理， $u_f$ 是常函数。从而
\begin{equation*}
f\left(R_\lambda(A)\right)=f\left(R_\mu(A)\right), \forall \lambda, \mu \in \mathbb{C} , \forall f \in L (X)^*
\end{equation*}
由 HBT, $R_\lambda(A)=R_\mu(A)$, 这与引理2.3矛盾。
\end{proof}

\begin{definition}
设 $A \in L(X)$, 称数
\[
r_\sigma(A) \triangleq \sup \{ |\lambda| : \lambda \in \sigma(A) \}
\]
为 $A$ 的谱半径。
\end{definition}

\begin{theorem}[Gelfand定理]
设 \( X \) 是 Banach 空间，且 \( A \in L(X) \)，则有：
\[
r_\sigma(A) = \lim_{n \to \infty} \|A^n\|^{\frac{1}{n}}.
\]
\end{theorem}
\begin{proof}
Step1: 先证明右式极限存在, 令 $r \stackrel{\text { def }}{=} \inf _n\left\|A^n\right\|^{\frac{1}{n}}$, 则
\begin{equation*}
\liminf _{n \rightarrow \infty}\left\|A^n\right\|^{\frac{1}{n}} \geqslant r
\end{equation*}
另一方面, $\forall \varepsilon>0$, 存在 $m$ 使得
\begin{equation*}
\left\|A^m\right\|^{\frac{1}{m}}<r+\varepsilon
\end{equation*}
所以对于 $\forall n \in N$, 有唯一分解 $n=p_n m+q_n$ with $0 \leqslant q_n<m$, 所以
\begin{equation*}
\left\|A^n\right\|^{\frac{1}{n}} \leqslant\left\|A^{p_n m}\right\|^{\frac{1}{n}}\left\|A^{q_n}\right\|^{\frac{1}{n}} \leqslant\left\|A^m\right\|^{\frac{q_n}{n}}\|A\|^{\frac{q_n}{n}}<(r+\varepsilon)^{\frac{p_n m}{n}}\|A\|^{\frac{q_n}{n}}
\end{equation*}
当 $n \rightarrow \infty$ 时, $\frac{q_n}{n} \rightarrow 0, \frac{p_n m}{n} \rightarrow 1$, 所以
\begin{equation*}
\limsup _{n \rightarrow \infty}\left\|A^n\right\|^{\frac{1}{n}} \leqslant r+\varepsilon
\end{equation*}
令 $\varepsilon \rightarrow 0$ 则
\begin{equation*}
\limsup _{n \rightarrow \infty}\left\|A^n\right\|^{\frac{1}{n}} \leqslant r
\end{equation*}
Step2: 证明 $r_\sigma(A) \leqslant \lim _{n \rightarrow \infty}\left\|A^n\right\|^{\frac{1}{n}}$. 我们知道幂级数
\begin{equation*}
\sum_{n=0}^{\infty}\left\|A^n\right\| z^n
\end{equation*}
的收敛半径为
\begin{equation*}
r=\frac{1}{\lim _{n \rightarrow \infty}\left\|A^n\right\|^{\frac{1}{n}}}
\end{equation*}
令 $z=\frac{1}{\lambda}$, 可知当 $|\lambda|>\lim _{n \rightarrow \infty}\left\|A^n\right\|^{\frac{1}{n}}$ 时（收敛圆内绝对收敛）
\begin{equation*}
\sum_{n=0}^{\infty}\left\|\frac{A^n}{\lambda^{n+1}}\right\|<\infty
\end{equation*}
$L (X)$ 完备, 根据引理 1.2, 级数
\begin{equation*}
\sum_{n=0}^{\infty}\frac{A^n}{\lambda^{n+1}}
\end{equation*}
也收敛。另一方面，
\begin{equation*}
\left\|\left(\sum_{n=0}^N \frac{A^n}{\lambda^{n+1}}\right)(\lambda I-A)-I\right\|=\left\|\frac{A^{N+1}}{\lambda^{N+1}}\right\| \rightarrow 0
\end{equation*}
所以
\begin{equation*}
\sum_{n=0}^{\infty} \frac{A^n}{\lambda^{n+1}}=(\lambda I-A)^{-1}=R_\lambda(A)
\end{equation*}
从而 $R_\lambda(A) \in L (X) \Rightarrow \lambda \in \rho(A) \rightarrow r_\sigma(A) \leqslant \lim _{n \rightarrow \infty}\left\|A^n\right\|^{\frac{1}{n}}$.\\
Step3: 证明 $r_\sigma(A) \geqslant \lim _{n \rightarrow \infty}\left\|A^n\right\|^{\frac{1}{n}}$. 设 $|\lambda|>r_\sigma(A)$, 则 $\lambda \in \rho(A) \Rightarrow \forall f \in L (X)^*, f\left(R_\lambda(A)\right)$ 在 $\lambda$ 全纯, 从而$f\left(R_\lambda(A)\right)$ 在圆环 $|\lambda|>r_\sigma(A)$ 内全纯, 故可展为收敛的 Laurent 级数。另一方面, 由 Step2, 当 $|\lambda|>\lim _{n \rightarrow \infty}\left\|A^n\right\|^{\frac{1}{n}}$时,
\begin{equation*}
R_\lambda(A)=\sum_{n=0}^{\infty} \frac{A^n}{\lambda^{n+1}} \Rightarrow f\left(R_\lambda(A)\right)=\sum_{n=0}^{\infty} \frac{f\left(A^n\right)}{\lambda^{n+1}}
\end{equation*}
Laurent 展式唯一, 所以这一展式在 $|\lambda|>r_\sigma(A)$ 上也成立。在内部绝对收敛, 所以
\begin{equation*}
\forall \varepsilon>0, \sum_{n=1}^{\infty} \frac{\left|f\left(A^n\right)\right|}{\left(r_\sigma(A)+\varepsilon\right)^{n+1}}<\infty
\end{equation*}
记
\begin{equation*}
T_n \stackrel{\text { def }}{=} \frac{A^n}{\left(r_\sigma(A)+\varepsilon\right)^{n+1}}
\end{equation*}
收敛级数通项有界，所以
\begin{equation*}
\sup _n\left|f\left(T_n\right)\right|<\infty, \forall f \in L (X)^*
\end{equation*}
由 UBP, $C \stackrel{\text { def }}{=} \sup _n\left\|T_n\right\|<\infty$ ，从而
\begin{equation*}
\left\|A^n\right\| \leqslant C\left(r_\sigma(A)+\varepsilon\right)^{n+1} \Rightarrow \lim _{n \rightarrow \infty}\left\|A^n\right\|^{\frac{1}{n}} \leqslant r_\sigma(A)+\varepsilon
\end{equation*}
令 $\varepsilon \rightarrow 0$ 得证。
\end{proof}
\begin{example}
考虑空间 \( l^2 \) 上的右推移算子 \( A \)，定义为：
\[
A: x = \left(x_1, x_2, \cdots, x_n, \cdots\right) \mapsto \left(0, x_1, x_2, \cdots, x_n, \cdots\right).
\]
我们需要证明：
\[
\begin{aligned}
\sigma_p(A) &= \varnothing, \\
\sigma_r(A) &= \{\lambda \in \mathbb{C} \mid |\lambda| < 1\}, \\
\sigma_c(A) &= \{\lambda \in \mathbb{C} \mid |\lambda| = 1\}.
\end{aligned}
\]

\end{example}
\begin{proof}
由Gelfand定理知 $\|A\|=1 \Rightarrow \sigma(A)\subset\overline{ D }$。\\
先证明: $\sigma_p(A)=\varnothing$, 否则 $\exists \lambda \in \mathbb{C} , \exists 0 \neq x \in l^2$ 使得
\begin{equation*}
\left(0, x_1, x_2, \cdots\right)=A x=\lambda x=\left(\lambda x_1, \lambda x_2, \cdots\right)
\end{equation*}
$\lambda=0 \Rightarrow x=0, \lambda \neq 0 \Rightarrow x_1=0 \Rightarrow x_2=0 \Rightarrow \cdots \Rightarrow x=0$. 矛盾。\\
再证明: $D \subset \sigma_r(A)$, 设 $\lambda \in D$, Claim: $\overline{\operatorname{Ran}(\lambda I-A)} \neq l^2$ 。这等价于 $\operatorname{Ran}(\lambda I-A)^{\perp} \neq\{0\}$. 令 $z=\left(1, \bar{\lambda}, \bar{\lambda}^2, \cdots\right)$,则
\begin{equation*}
\begin{aligned}
\langle(\lambda I-A) x, z\rangle & =\left\langle\left(\lambda x_1, \lambda x_2-x_2, \lambda x_3-x_2, \cdots\right),\left(1, \bar{\lambda}, \bar{\lambda}^2, \cdots\right)\right\rangle \\
& =\lambda x_1+\lambda^2 x_2-\lambda x_1+\lambda^3 x_3-\lambda^2 x_2+\cdots=0 \\
\Rightarrow 0 \neq z & \in \operatorname{Ran}(\lambda I-A)^{\perp}
\end{aligned}
\end{equation*}
然后证明: $\partial D \subset \sigma_c(A)$. 设 $\lambda \in \partial D$。 
\\$1^{\circ}$ 证明 $\operatorname{Ran}(\lambda I-A) \neq l^2$ 。
\begin{equation*}
\begin{aligned}
\operatorname{Ran}(\lambda I-A) \ni y=(\lambda I-A) x & \Rightarrow y_1=\lambda x_1, y_k=\lambda x_k-x_{k-1}, k \geqslant 2 \\
& \Rightarrow y_1=\lambda x_1, \lambda^{k-1} y_k=\lambda^k x_k-\lambda^{k-1} x_{k-1}, k \geqslant 2 \\
& \Rightarrow \sum_{k=1}^n \lambda^{k-1} y_k=\lambda^n x_n
\end{aligned}
\end{equation*}
假设 $\operatorname{Ran}(\lambda I-A)=l^2$, 令 $y=e_1$,
\begin{equation*}
\begin{aligned}
\exists x \in l^2 \text { s.t. } e_1=(\lambda I-A) x & \Rightarrow \lambda^n x_n=1, n=1,2, \cdots \\
& \Rightarrow x=\left(\frac{1}{\lambda}, \frac{1}{\lambda^2}, \cdots\right)
\end{aligned}
\end{equation*}
$|\lambda|=1$, 不收敛, 这与 $x \in l^2$ 矛盾。\\
$2^{\circ}$ 证明 $\overline{\operatorname{Ran}(\lambda I-A)}=l^2$. 只需证明 $\operatorname{Ran}(\lambda I-A)^{\perp}=\{0\}$. 对于 $\forall z \in \operatorname{Ran}(\lambda I-A)^{\perp}$,
\begin{equation*}
0=\left\langle z,(\lambda I-A) e_n\right\rangle=\bar{\lambda} z_n-z_{n+1}, \forall n
\end{equation*}
所以 $z_{n+1}=\bar{\lambda} z_n \Rightarrow\left|z_{n+1}\right|=\left|z_n\right| \Rightarrow z=0$.
最后: $\overline{ D } \subset \sigma_c(A) \cup \sigma_r(A) \subset \sigma(A) \subset \overline{ D }$, 由前面结论可得 $\sigma_c(A)=\partial D , \sigma_r(A)= D$.
\end{proof}


\begin{example}
设 \( (X, (\cdot, \cdot)) \) 是一个 Hilbert 空间，且 \( A \in L(X) \) 为对称算子，即
\[
(A x, y) = (x, A y) \quad \forall x, y \in X.
\]
若 \( A \) 是对称算子，则 \( \sigma(A) \subset \mathbb{R} \)，且 \( \sigma_r(A) = \varnothing \)。
\end{example}
\begin{proof}
我们有 \( (A x, x) \in \mathbb{R} \) 对于任意 \( x \in X \)。设 \( \lambda = a + ib \) 且 \( b \neq 0 \)，我们需要证明 \( \lambda I - A \) 存在有界逆。\\
首先，计算
\[
\|(\lambda I - A) x\|^2 = \|(a I - A) x\|^2 + b^2 \|x\|^2,
\]
因此 \( N(\lambda I - A) = \{0\} \)，即 \( \lambda I - A \) 为单射。\\
接下来，我们证明 \( \lambda I - A \) 是满射，即 \( R(\lambda I - A) = X \)。若成立，则由 Banach 逆算子定理，\( \lambda I - A \) 存在有界逆，且 \( \lambda \notin \sigma(A) \)。为了证明 \( R(\lambda I - A) \) 是闭的，设 \( x_n \in X \)，且
\[
(\lambda I - A) x_n = y_n \to y \quad (n \to \infty).
\]
由上式得
\[
\left\|(\lambda I - A)(x_n - x_m)\right\|^2 \geqslant b^2 \|x_n - x_m\|^2,
\]
即 \( \{x_n\} \) 是 \( X \) 中的 Cauchy 列，故存在 \( x_n \to x \)（\( n \to \infty \)）。这时有
\[
y_n = (\lambda I - A) x_n \to (\lambda I - A) x = y \quad (n \to \infty),
\]
所以
\[
R(\lambda I - A) = \overline{R(\lambda I - A)}.
\]
接下来证明 \( R(\lambda I - A)^\perp = \{0\} \)。设 \( y \in R(\lambda I - A)^\perp \)，即
\[
((\lambda I - A) x, y) = 0 \quad \forall x \in X.
\]
由于 \( A \) 是对称的，得到
\[
(x, (\bar{\lambda} I - A) y) = 0 \quad \forall x \in X,
\]
因此 \( (\bar{\lambda} I - A) y = 0 \)，即 \( y \in N(\bar{\lambda} I - A) \)。通过（2.6.13）式中用 \( \bar{\lambda} \) 代替 \( \lambda \) 得到 \( N(\bar{\lambda} I - A) = \{0\} \)，所以必有 \( y = 0 \)，即
\[
R(\lambda I - A)^\perp = \{0\}.
\]
设 \( \lambda \in \sigma(A) \)，但 \( \lambda \) 不是点谱，即 \( N(\lambda I - A) = \{0\} \)。此时
\[
R(\lambda I - A)^\perp = N(\bar{\lambda} I - A^*) = N(\lambda I - A) = \{0\},
\]
即 \( R(\lambda I - A) \) 在 \( X \) 中稠密，故 \( \lambda \in \sigma_c(A) \)，即 \( \sigma_r(A) = \varnothing \)。\\
结论得证。
\end{proof}
% 例 2.6.15 酉算子
\begin{example}
设 \( (X, (\cdot, \cdot)) \) 是 Hilbert 空间，且 \( U \in L(X) \) 为酉算子，即
\[
(U x, U y) = (x, y), \quad \forall x, y \in X, \quad R(U) = X.
\]
若 \( U \) 是酉算子，则 \( \sigma(U) \subset S^1 = \left\{ e^{i\theta} \mid \theta \in [0, 2\pi] \right\} \)，且 \( \sigma_r(U) = \varnothing \)。


\end{example}


\begin{proof}
由定理 2.6.14，得
\[
\|U x\|^2 = \|x\|^2,
\]
因此 \( U \) 为单射，并且
\[
\|U\| = 1, \quad \sigma(U) \subset \{\lambda \in \mathbb{C} \mid |\lambda| \leqslant 1\}.
\]
又由于 \( U \) 是满射，且 \( \|U^{-1}\| = 1 \)，根据 Banach 逆算子定理，\( U^{-1} \) 存在且有界。\\
设 \( |\lambda| < 1 \)，则
\[
\lambda I - U = -U \left( I - \lambda U^{-1} \right).
\]
因为 \( \| \lambda U^{-1} \| = |\lambda| \| U^{-1} \| = |\lambda| < 1 \)，所以 \( I - \lambda U^{-1} \) 存在有界逆，且
\[
(\lambda I - U)^{-1} = \left( I - \lambda U^{-1} \right)^{-1} U^{-1},
\]
因此 \( \lambda \notin \sigma(U) \)。\\
结论得证。
\end{proof}
\begin{problem}
设 $X$ 是 $B$ 空间，求证：$L(X)$ 中可逆（存在有界逆）算子集是开的。
\end{problem}
\begin{proof}
注意到
\begin{equation*}
A + \lambda I = A \left(I + \lambda A^{-1}\right) = \left(I + \lambda A^{-1}\right) A,
\end{equation*}
根据本节引理 2.2，当 $\left\|\lambda A^{-1}\right\| < 1$ 时，有
\[
\left(I + \lambda A^{-1}\right)^{-1} \in L(X).
\]
因此，根据（1）式，当 $\lambda < \frac{1}{\left\|A^{-1}\right\|}$ 时，我们有
\begin{equation*}
(A + \lambda I)\left(I + \lambda A^{-1}\right)^{-1} A^{-1} = \left(I + \lambda A^{-1}\right)^{-1} A^{-1}(A + \lambda I) = I,
\end{equation*}
再根据例 1，当 $\lambda < \frac{1}{\left\|A^{-1}\right\|}$ 时，
\begin{equation*}
(A + \lambda I)^{-1} = \left(I + \lambda A^{-1}\right)^{-1} A^{-1} \in L(X).
\end{equation*}
\end{proof}
\begin{problem}
设 $A$ 是闭线性算子，$\lambda_1, \lambda_2, \cdots, \lambda_n \in \sigma_p(A)$ 两两互异，又设 $x_i$ 是对应于 $\lambda_i$ 的特征元（$i = 1, 2, \cdots, n$）。求证：$\{x_1, x_2, \cdots, x_n\}$ 是线性无关的。
\end{problem}
\begin{proof}
 用反证法. 令 $x_m$ 为第一个可由它的前面 $m-1$ 个向量线性表出的向量, 即
\begin{equation*}
x_m=\sum_{k=1}^{m-1} \alpha_k x_k,
\end{equation*}
且 $x_1, x_2, \cdots, x_{m-1}$ 线性无关, 对 (1) 式两边施以 $\lambda_m I-A$ 得到
\begin{equation*}
0=\left(\lambda_m I-A\right) x_m=\sum_{k=1}^{m-1} \alpha_k\left(\lambda_m I-A\right) x_k=\sum_{k=1}^{m-1} \alpha_k\left(\lambda_m-\lambda_k\right) x_k
\end{equation*}
因为 $x_1, x_2, \cdots, x_{m-1}$ 线性无关, 所以
\begin{equation*}
\alpha_k\left(\lambda_m-\lambda_k\right)=0 \quad(k=1,2, \cdots, m-1) .
\end{equation*}
注意到 $\lambda_m \neq \lambda_k$, 故有
\begin{equation*}
\alpha_k=0 \quad(k=1,2, \cdots, m-1) .
\end{equation*}
于是 (1) 式蕴含 $x_m=0$, 这与 $x_m$ 为特征向量矛盾.
\end{proof}
\begin{proof}
在双边 $l^2$ 空间上，右推移算子 $A$ 定义为：
\[
x = (\cdots, \xi_{-1}, \xi_0, \xi_1, \cdots) \in l^2, \quad y = A x = (\cdots, \eta_{-1}, \eta_0, \eta_1, \cdots),
\]
其中 $\eta_m = \xi_{m-1} \; (m \in \mathbb{Z})$。我们证明：
\[
\sigma(A) = \sigma_c(A) = \{|\lambda| = 1\}.
\]

\textbf{1. $\sigma_p(A) = \varnothing$}

由 $A x$ 的定义知 $\|A x\| = \|x\|$，故 $\|A\| = 1$。结合
\[
\|A^n x\| = \|x\| \implies \|A^n\| = 1 \implies r_\sigma(A) = \lim_{n \to \infty} \|A^n\| = 1.
\]

(1) 当 $\lambda = 0$：
\[
(\lambda I - A)x = 0 \implies A x = 0 \implies x = 0.
\]

(2) 当 $\lambda \neq 0$：
\[
(\lambda I - A)x = 0 \implies \lambda \xi_k - \xi_{k-1} = 0, \quad \forall k \in \mathbb{Z}.
\]
递推得到：
\[
\xi_n = \frac{1}{\lambda^n} \xi_0, \quad \xi_{-n} = \lambda^n \xi_0.
\]
由于 $x \in l^2$，有：
\[
\sum_{n=-\infty}^\infty |\xi_n|^2 < \infty.
\]
但这与 $|\lambda| \neq 1$ 的假设矛盾，因此 $x = 0$，即 $\sigma_p(A) = \varnothing$。

\textbf{2. $\sigma_r(A) = \varnothing$}

根据定义，需证 $\overline{R(\lambda I - A)} = l^2$，即 $R(\lambda I - A)^\perp = \{\theta\}$。

设 $z \perp R(\lambda I - A)$，则：
\[
((\lambda I - A)x, z) = 0, \quad \forall x \in l^2.
\]
取特殊序列 $x^{(n)} = (0, \cdots, 0, 1, 0, \cdots) \in l^2$，计算得：
\[
\lambda z_n - z_{n+1} = 0, \quad \forall n \in \mathbb{Z}.
\]
递推与之前类似，得 $z = \theta$。故 $\sigma_r(A) = \varnothing$。

\textbf{3. $\sigma_c(A) = \{|\lambda| = 1\}$}

$R(\lambda I-A)\neq l^2$，考虑$e_0$，若$(\lambda I- A)x=e_0$，则
\begin{equation*}
    \begin{aligned}
        \lambda x_0-x_1=1,\lambda x_n-x_{n+1}=0,\lambda x_{-n}-x_{-n+1}=0,n\in \mathbb{N}
        \\
        x_n=\lambda^{n-1} x_1,x_{-n}={\lambda}^{-n}x_0,\|x\|^2=\sum\limits_{n=1}^\infty|x_1|^2+\sum\limits_{n=0}^\infty|x_0|^2< \infty
    \end{aligned}
\end{equation*}

所以$x_1=x_0=0$与$\lambda x_0-x_1=1$矛盾。

下证$R(\lambda I-A)^\perp =\{\theta\}$。若$z\in R(\lambda I-A)^\perp$，则
\begin{equation*}
    \begin{aligned}
        \langle (\lambda I-A)e_n,z\rangle=\lambda z_n-z_{n+1}=0\Rightarrow \|z\|^2=\sum\limits_{n=-\infty}^\infty |z_0|^2<\infty 
    \end{aligned}
\end{equation*}
所以$z=\theta$。

综上，$\sigma(A) = \sigma_c(A) = \{|\lambda| = 1\}$。
\end{proof}
\begin{problem}
在 $l^2$ 空间上，考察左推移算子
\[
A: \left(\xi_1, \xi_2, \cdots\right) \mapsto \left(\xi_2, \xi_3, \cdots\right).
\]
求证：$\sigma_p(A) = \{\lambda \in \mathbb{C} \mid |\lambda| < 1\}$,  
$\sigma_c(A) = \{\lambda \in \mathbb{C} \mid |\lambda| = 1\}$, 并且
\[
\sigma(A) = \sigma_p(A) \cup \sigma_c(A).
\]
\end{problem}

\begin{proof}
\textbf{1. 证明 $A$ 是线性有界算子且 $\|A\| = 1$}

我们设
\[
x = (x_1, x_2, \dots) \in l^2, \quad y = A x = (x_2, x_3, \dots).
\]

可以验证：
\[
\|A x\|^2 = \sum_{n=2}^\infty |x_n|^2 \leq \sum_{n=1}^\infty |x_n|^2 = \|x\|^2.
\]

因此：
\[
\|A x\| \leq \|x\|.
\]

另外，取 $x' = (0, 1, 0, 0, \dots)$，可以验证：
\[
\|A x'\| = \|x'\| = 1.
\]

因此算子 $A$ 是线性有界算子且其算子范数为：
\[
\|A\| = \sup_{x \neq 0} \frac{\|A x\|}{\|x\|} = 1.
\]

\textbf{2. 证明  $\{\lambda \in \mathbb{C} \mid |\lambda| > 1\} \subset \rho(A)$}

根据谱的定义，当 $|\lambda| > 1$ 时：
\[
|\lambda| > \|A\| = 1 \implies \lambda \in \rho(A),
\]
其中 $\rho(A)$ 表示 $A$ 的 resolvent 子集。

因此：
\[
\{\lambda \in \mathbb{C} \mid |\lambda| > 1\} \subset \rho(A).
\]

\textbf{3. 证明  $\{\lambda \in \mathbb{C} \mid |\lambda| < 1\} = \sigma_p(A)$}

定义：
\[
D = \{\lambda \in \mathbb{C} \mid |\lambda| < 1\}.
\]
对于任意 $\lambda \in D$，考虑数列 $\{\lambda^n\}_{n=0}^\infty \in l^2$，因为对于 $|\lambda| < 1$，
\[
\sum_{n=0}^\infty |\lambda^n|^2 < \infty,
\]
所以：
\[
A (1, \lambda, \lambda^2, \dots) = (\lambda, \lambda^2, \dots) = \lambda (1, \lambda, \lambda^2, \dots).
\]
从而，这表明 $\lambda \in \sigma_p(A)$，其中 $\sigma_p(A)$ 是 $A$ 的点谱。

反之，如果 $A x = \lambda x, \, x \neq 0, \, x \in l^2$，则：
\[
\lambda^n x = A^n x = (x_2, x_3, \dots) \to 0 \text{ 当 } n \to \infty,
\]
这说明 $|\lambda| < 1$。

因此：
\[
\sigma_p(A) = \{\lambda \in \mathbb{C} \mid |\lambda| < 1\}.
\]

\textbf{4. 证明  $\{\lambda \in \mathbb{C} \mid |\lambda| = 1\} = \sigma_c(A)$}

\begin{equation*}
\begin{aligned}
& \operatorname{Ran}(\lambda I-A) \neq \ell^2: \\
& \qquad \begin{aligned}
y \in \operatorname{Ran}(\lambda I-A) & \Rightarrow \exists 0 \neq x \in l^2, y=(\lambda I-A) x=\left(\lambda x_1-x_2, \lambda x_2-x_3, \cdots\right)=\left\{\lambda x_n-x_{n+1}\right\}_{n=1}^{\infty} \\
& \Rightarrow \sum_{n=1}^k \lambda^{-n} y_n=x_1-\lambda^{-k} x_{k+1} \rightarrow x_1 \text { as } k \rightarrow \infty
\end{aligned}
\end{aligned}
\end{equation*}

所以
\begin{equation*}
\sum_{n=1}^{\infty} \lambda^{-n} y_n=x_1
\end{equation*}

但是取

\begin{equation*}
y=\left\{\lambda^n \cdot \frac{1}{n}\right\}_{n=1}^{\infty} \in l^2
\end{equation*}

则级数发散，说明满足 $y=(\lambda I-A) x$ 的 $x$ 不存在。因此 $\operatorname{Ran}(\lambda I-A) \neq l^2$ ．

下证$R(\lambda I-A)^\perp =\{\theta\}$。若$z\in R(\lambda I-A)^\perp$，则
\begin{equation*}
    \begin{aligned}
 \langle (\lambda I-A)e_1,z\rangle=\lambda z_1=0,        \langle (\lambda I-A)e_n,z\rangle=\lambda z_n-z_{n-1}=0\Rightarrow z_n=0 
    \end{aligned}
\end{equation*}
所以$z=\theta$。

\textbf{5. 证明  $\sigma_r(A) = \varnothing$}

根据定义，要求 $\sigma_r(A)$：
\[
\sigma_r(A) = \emptyset.
\]
这是因为不存在 $\lambda$ 使得 $(\lambda I - A)^{-1}$ 不存在，且没有非零解满足性质。

\textbf{结论}

综上所述：
\[
\sigma_p(A) = \{\lambda \in \mathbb{C} \mid |\lambda| < 1\},
\]
\[
\sigma_c(A) = \{\lambda \in \mathbb{C} \mid |\lambda| = 1\},
\]
\[
\sigma_r(A) = \varnothing.
\]
\end{proof}
\chapter{广义函数与Sobolev空间}

在数学分析中,我们对函数的概念有非常明确的描述: “若存在对应法则 \(f\) 使得对于集合 \(X\) 中的任一实数 \(x\) ,存在唯一的实数 \(y\) 按照法则 \(f\) 与之对应,则称 \(f\) 是定义在 \(X\) 上的 \(x\) 的一个函数”.这一概念被广泛应用到科学技术的各个领域,但当我们在力学、物理学及技术科学中表示某些集中量的分布密度或密度分布时,这一基本概念就不够用了. 例如,在三维Euclid空间 \({\mathbb{R}}^{3}\) 中的点 \({x}^{0} = \left( {{x}_{1}^{0},{x}_{2}^{0},{x}_{3}^{0}}\right)\) 处放置单位正电荷, 若用
\begin{equation*}
\delta \left( {x - {x}^{0}}\right)  \mathrel{\text{ := }} \delta \left( {{x}_{1} - {x}_{1}^{0},{x}_{2} - {x}_{2}^{0},{x}_{3} - {x}_{3}^{0}}\right)
\end{equation*}

表示电荷的空间分布密度, 那么由密度的通常定义形式上应该有
\begin{equation*}
\left\{  \begin{array}{l} \delta \left( {x - {x}^{0}}\right)  = \mathop{\lim }\limits_{{\bigtriangleup V \rightarrow  0}}\frac{\bigtriangleup Q}{\bigtriangleup V}, \\   {\iiint }_{{\mathbb{R}}^{3}}\delta \left( {x - {x}^{0}}\right) {\dif x} = 1. \end{array}\right.
\end{equation*}

其中 \(\bigtriangleup Q\) 表示含点 \(x\) 的体积元 \(\bigtriangleup V\) 中所含电荷的总量. 通过极限计算不难得到
\begin{equation*}
\left\{  \begin{array}{l} \delta \left( {x - {x}^{0}}\right)  = \left\{  \begin{matrix} \infty , & x = {x}^{0}, \\  0, & x \neq  {x}^{0}, \end{matrix}\right. \\   {\iiint }_{{\mathbb{R}}^{3}}\delta \left( {x - {x}^{0}}\right) {\dif x} = 1. \end{array}\right.
\end{equation*}

显然,这个 \(\delta\) 函数不是前面经典意义下的函数,因为在任何积分理论中都得不出 \(\delta \left( {x - {x}^{0}}\right)\) 在 \({\mathbb{R}}^{3}\) 的积分值为 1 的结果. 但该 \(\delta\) 函数确实反映或近似反映了单位集中量的密度分布, 并给现实的不连续量提供了一种数学表述的可能性. 因此, 我们有必要拓展 “函数” 这一基本概念.

此外, 数学中的许多运算往往都是在一定限制条件下才成立的, 如求导运算等,数学自身的发展有时也需要突破这些限制. 上世纪三十年代,苏联数学家Sobolev为了确定偏微分方程解的存在性和唯一性, 发现如果限制在古典解要求的光滑函数类中, 解可能不存在, 同时也限制了很多近代数学工具的使用. 当我们扩大微分方程的解所在的函数空间时, 首先必须推广经典的导数概念. Sobolev引进了广义导数的概念, 提出了广义函数的原型.

在讲述广义函数的定义之前,我们先引进一些记号. 记多重指标 \(\alpha  = \left( {{\alpha }_{1},{\alpha }_{2}}\right.\) , \(\left. {\cdots ,{\alpha }_{n}}\right)\) ,其中 \({\alpha }_{1},{\alpha }_{2},\cdots ,{\alpha }_{n} \geq  0\) 是整数,
\begin{equation*}
\left| \alpha \right|  = \mathop{\sum }\limits_{{i = 1}}^{n}{\alpha }_{i},\;\alpha ! = {\alpha }_{1}!{\alpha }_{2}!\cdots {\alpha }_{n}!
\end{equation*}
\begin{equation*}
{x}^{\alpha } = {x}_{1}^{{\alpha }_{1}}{x}_{2}^{{\alpha }_{2}}\cdots {x}_{n}^{{\alpha }_{n}}\;\left( {\forall x = \left( {{x}_{1},{x}_{2},\cdots ,{x}_{n}}\right)  \in  {\mathbb{R}}^{n}}\right) ,
\end{equation*}
\begin{equation*}
{\partial }^{\alpha } = {\partial }_{{x}_{1}}^{{\alpha }_{1}}{\partial }_{{x}_{2}}^{{\alpha }_{2}}\cdots {\partial }_{{x}_{n}}^{{\alpha }_{n}}
\end{equation*}

而且如果 \(\beta  \leqslant  \alpha\) (即 \({\beta }_{j} \leqslant  {\alpha }_{j},j = 1,2,\cdots ,n\) ),则
\begin{equation*}
\left( \begin{array}{l} \alpha \\  \beta  \end{array}\right)  = \frac{\alpha !}{\beta !\left( {\alpha  - \beta }\right) !} = \left( \begin{array}{l} {\alpha }_{1} \\  {\beta }_{1} \end{array}\right) \left( \begin{array}{l} {\alpha }_{2} \\  {\beta }_{2} \end{array}\right) \cdots \left( \begin{array}{l} {\alpha }_{n} \\  {\beta }_{n} \end{array}\right) .
\end{equation*}

\section{广义函数的概念}
在给出广义函数的定义之前, 我们先来看看前面描述的集中量密度分布的 \(\delta\) 函数. 对于Euclid空间 \({\mathbb{R}}^{n}\) 中的任一固定点
\begin{equation*}
{x}^{0} = \left( {{x}_{1}^{0},\cdots ,{x}_{n}^{0}}\right) ,
\end{equation*}

对于在点 \({x}^{0}\) 放置集中量 \(q = f\left( {x}^{0}\right)\) 的任一连续延拓 \(f\left( x\right)\) ,均应有
\begin{equation}
  {\int }_{{\mathbb{R}}^{n}}\delta \left( {x - {x}^{0}}\right) f\left( x\right) {\dif x} = f\left( {x}^{0}\right)  = q. \label{3.1}
\end{equation}

\ref{3.1}表明,尽管我们不能将 \(\delta\) 函数视为经典意义下的函数,但却可以由它对一些函数的作用值 \(\left\langle  {\delta \left( {x - {x}^{0}}\right) ,f\left( x\right) }\right\rangle   = f\left( {x}^{0}\right)\) 来获取这个 \(\delta\) 函数的信息. 从前面关于泛函的定义可知,我们可以把 \(\delta\) 函数看成某个函数空间 \(\Phi\) 上的连续线性泛函. 这样一来,我们自然会问,应该如何选择空间 \(\Phi\) 才会使得 \(\delta\) 函数具有更方便的运算性质呢? 习惯上,我们取 \(\Phi\) 为由所有在 \({\mathbb{R}}^{n}\) 中无穷次可微而分别在不同的有界域外恒等于零的函数 \(\varphi \left( x\right)\) 组成的空间 \({C}_{0}^{\infty }\left( {\mathbb{R}}^{n}\right) ,{C}_{0}^{\infty }\left( {\mathbb{R}}^{n}\right)\) 中的函数通常称为检验函数或基本函数, 相应地称 \({C}_{0}^{\infty }\left( {\mathbb{R}}^{n}\right)\) 为检验函数空间或基本空间. 当然我们也可取其他的基本空间, 但 \({C}_{0}^{\infty }\left( {\mathbb{R}}^{n}\right)\) 是最重要的基本空间.

\subsection{基本空间 \(\mathcal{D}\)}
设 \(\Omega  \subset  {\mathbb{R}}^{n}\) 是一个开集, \(\varphi  \in  {C}^{\infty }\left( \bar{\Omega }\right)\) ,称集合 \(\{ x \in  \Omega  \mid  \varphi \left( x\right)  \neq  0\}\) 关于 \(\Omega\) 的闭包为 \(\varphi\) 关于 \(\Omega\) 的支集,记为 \(\operatorname{supp}\varphi\) . 通常称 \({C}_{0}^{\infty }\left( {\mathbb{R}}^{n}\right)\) 中的函数为紧支集函数,实际上, 在 \(\operatorname{supp}\varphi\) 之外, \(\varphi \left( x\right)  \equiv  0\) .
\begin{example}
  Dirichlet 函数
\begin{equation*}
D\left( x\right)  = \left\{  \begin{array}{ll} 1, & \text{ 当 }x\text{ 为有理数, } \\  0, & \text{ 当 }x\text{ 为无理数. } \end{array}\right.
\end{equation*}

的支集为整个实数轴.

\end{example}

一般地,对于整数 \(k \geq  0\) (可以是 \(\infty\) ),记
\begin{equation*}
{C}_{0}^{k}\left( \Omega \right)  = \left\{  {\varphi  \in  {C}^{k}\left( \bar{\Omega }\right)  \mid  \operatorname{supp}\varphi \text{ 为 }\Omega \text{ 中紧集 }}\right\}  .
\end{equation*}

于是,
\begin{equation*}
{C}_{0}^{\infty }\left( \Omega \right)  \subset  \cdots  \subset  {C}_{0}^{k + 1}\left( \Omega \right)  \subset  {C}_{0}^{k}\left( \Omega \right)  \subset  \cdots  \subset  {C}_{0}^{0}\left( \Omega \right) .
\end{equation*}

\begin{example}
\begin{equation*}
\varphi \left( {x,a,b,c}\right)  = \left\{  \begin{array}{ll} c\exp \frac{-{b}^{2}}{{a}^{2} - {\left| x\right| }^{2}}, & {\left| x\right| }^{2} = {x}_{1}^{2} + \cdots  + {x}_{n}^{2} < {a}^{2}, \\  0, & {\left| x\right| }^{2} \geq  {a}^{2} \end{array}\right.
\end{equation*}

其中 \(a,b\) 为任意正数, \(c\) 为任意实数或复数.

\end{example}

\begin{example}
  定义
\begin{equation}
  {j}_{\delta }\left( x\right)  \mathrel{\text{ := }} \varphi \left( {x,\delta ,\delta ,{C}_{\delta }}\right)  \in  {C}_{0}^{\infty }\left( {\mathbb{R}}^{n}\right) , \label{3.2}
\end{equation}

\end{example}

其中 \(\varphi \left( {x,a,b,c}\right)\) 为例3.1.2中函数且 \({C}_{\delta } = {\left( {\int }_{\left| x\right|  \leqslant  \delta }\exp \frac{-{\delta }^{2}}{{\delta }^{2} - {\left| x\right| }^{2}}\dif x\right) }^{-1}\) . 则有 \({C}_{0}^{\infty }\left( {\mathbb{R}}^{n}\right)\) 非空. 我们习惯称 \({j}_{\delta }\left( x\right)\) 为磨光核.

实际上,从 \({j}_{\delta }\left( x\right)\) 出发,我们可以得到许多 \({C}_{0}^{\infty }\left( {\mathbb{R}}^{n}\right)\) 的函数.
\begin{theorem}
  设 \(\Omega  \subset  {\mathbb{R}}^{n}\) 为开集,则 \({C}_{0}^{\infty }\left( \Omega \right)\) 在 \({C}_{0}^{0}\left( \Omega \right)\) 中稠密,即 \({C}_{0}^{0}\left( \Omega \right)\) 中任意函数 \(u\left( x\right)\) 都可用 \({C}_{0}^{\infty }\left( \Omega \right)\) 中函数列一致逼近.

\end{theorem}
\begin{proof}
   对任意 \(u\left( x\right)  \in  {C}_{0}^{0}\left( \Omega \right)\) ,在 \(\Omega\) 外通过零延拓把 \(u\left( x\right)\) 的定义域延拓到全空间 \({\mathbb{R}}^{n}\) 上. 因此,下面只需证明 \({C}_{0}^{0}\left( {\mathbb{R}}^{n}\right)\) 中任意函数 \(u\left( x\right)\) 都可用 \({C}_{0}^{\infty }\left( {\mathbb{R}}^{n}\right)\) 中函数列一致逼近即可.

对于 \(u\left( x\right)  \in  {C}_{0}^{0}\left( {\mathbb{R}}^{n}\right)\) ,构造其光滑的磨光函数或正则化
\begin{equation}
  {u}_{\delta }\left( x\right)  \mathrel{\text{ := }} {\int }_{\left| {x - y}\right|  \leqslant  \delta }u\left( y\right) {j}_{\delta }\left( {x - y}\right) {\dif y}. \label{3.3}
\end{equation}

由\ref{3.3}知 \(\operatorname{supp}{u}_{\delta }\) 包含在 \(\operatorname{supp}u\) 的 \(\delta\) 邻域中,故 \({u}_{\delta }\) 具有紧支集. 下证 \({u}_{\delta } \in  {C}_{0}^{\infty }\left( {\mathbb{R}}^{n}\right)\) .

实际上, 我们有
\begin{equation}
  {\partial }^{\alpha }{u}_{\delta }\left( x\right)  = {\int }_{\left| {x - y}\right|  < \delta }u\left( y\right) {\partial }^{\alpha }{j}_{\delta }\left( {x - y}\right) {\dif y},\;x \in  {\mathbb{R}}^{n}. \label{3.4}
\end{equation}


先考虑指标 \({\alpha }_{0} = \left( {1,0,\cdots ,0}\right)\) ,
\begin{equation*}
{\partial }^{{\alpha }_{0}}{u}_{\delta }\left( x\right)  = \mathop{\lim }\limits_{{h \rightarrow  0}}{\int }_{\left| {x - y}\right|  < \delta }\frac{1}{h}\left\lbrack  {{j}_{\delta }\left( {x + h{e}_{1} - y}\right)  - {j}_{\delta }\left( {x - y}\right) }\right\rbrack  u\left( y\right) {\dif y}
\end{equation*}
\begin{equation*}
= \mathop{\lim }\limits_{{h \rightarrow  0}}{\int }_{\left| {x - y}\right|  < \delta }{\partial }^{{\alpha }_{0}}{j}_{\delta }\left( {x + {\theta h}{e}_{1} - y}\right) u\left( y\right) {\dif y},
\end{equation*}

其中 \(\theta  = \theta \left( {x,y}\right)  \in  \left( {0,1}\right) ,{e}_{1} = \left( {1,0,\cdots ,0}\right)  \in  {\mathbb{R}}^{n}\) . 由 \({j}_{\delta }\) 的连续可微性知,存在常数 \({M}_{{\alpha }_{0}}\) 使得
\begin{equation*}
\left| {{\partial }^{\alpha }{j}_{\delta }\left( x\right) }\right|  \leqslant  {M}_{{\alpha }_{0}},\;x \in  {\mathbb{R}}^{n}.
\end{equation*}

应用Lebesgue控制收敛定理可得
\begin{equation*}
{\partial }^{{\alpha }_{0}}{u}_{\delta }\left( x\right)  = {\int }_{\left| {x - y}\right|  < \delta }\mathop{\lim }\limits_{{h \rightarrow  0}}{\partial }^{{\alpha }_{0}}{j}_{\delta }\left( {x + {\theta h}{e}_{1} - y}\right) u\left( y\right) {\dif y}
\end{equation*}
\begin{equation*}
= {\int }_{\left| {x - y}\right|  < \delta }{\partial }^{{\alpha }_{0}}{j}_{\delta }\left( {x - y}\right) u\left( y\right) {\dif y}.
\end{equation*}

重复上述步骤,可得任意指标 \(\alpha\) 的等式 (3.4). 这意味着 \({u}_{\delta } \in  {C}_{0}^{\infty }\left( {\mathbb{R}}^{n}\right)\) .

下面证当 \(\delta  \rightarrow  0\) 时, \({u}_{\delta }\) 一致逼近 \(u\) . 为此,考虑
\begin{equation*}
u\left( x\right)  - {u}_{\delta }\left( x\right)  = {\int }_{\left| {x - y}\right|  \leqslant  \delta }u\left( x\right) {j}_{\delta }\left( {x - y}\right) {\dif y} - {\int }_{\left| {x - y}\right|  \leqslant  \delta }u\left( y\right) {j}_{\delta }\left( {x - y}\right) {\dif y}
\end{equation*}
\begin{equation*}
= {\int }_{\left| {x - y}\right|  \leqslant  \delta }\left\lbrack  {u\left( x\right)  - u\left( y\right) }\right\rbrack  {j}_{\delta }\left( {x - y}\right) {\dif y}.
\end{equation*}

由 \(u\) 一致连续,故对任给的 \(\varepsilon  > 0\) ,存在 \({\delta }_{0} > 0\) 使得当 \(\left| {x - y}\right|  < {\delta }_{0}\) 时有 \(\left| {u\left( x\right)  - u\left( y\right) }\right|  < \varepsilon\) . 于是由上式可得当 \(\delta  < {\delta }_{0}\) 时有
\begin{equation*}
\left| {u\left( x\right)  - {u}_{\delta }\left( x\right) }\right|  \leqslant  \varepsilon {\int }_{\left| {x - y}\right|  \leqslant  \delta }{j}_{\delta }\left( {x - y}\right) {\dif y} = \varepsilon .
\end{equation*}

这表明当 \(\delta  \rightarrow  0\) 时, \({u}_{\delta }\) 一致逼近于 \(u\) .

\end{proof}
\begin{note}
  实际上,若 \(u\left( x\right)\) 为可积函数且在 \(\Omega\) 内有紧子集,则(3.3)中定义的函数 \({u}_{\delta }\left( x\right)\)  \(\in  {C}_{0}^{\infty }\left( \Omega \right)\) . 进而可以证明若 \(u \in  {C}_{0}^{k}\left( \Omega \right)\) ,则

\begin{equation*}
{\begin{Vmatrix}{u}_{\delta }\left( x\right)  - u\left( x\right) \end{Vmatrix}}_{{C}^{k}\left( \Omega \right) } \rightarrow  0\;\left( {\delta  \rightarrow  0}\right) .
\end{equation*}

\end{note}


 \({C}_{0}^{\infty }\left( \Omega \right)\) 显然以实数域作系数域组成线性空间. 而且可进一步在 \({C}_{0}^{\infty }\left( \Omega \right)\) 上定义收敛性的概念.
\begin{definition}
  称 \({C}_{0}^{\infty }\left( \Omega \right)\) 中序列 \(\left\{  {{\varphi }_{j}\left( x\right) }\right\}\) 收敛到 \({\varphi }_{0}\) ,如果

(1)存在 \(\Omega\) 中紧集 \(K\) 使得对任意的 \(j = 0,1,2,\cdots\) ,都有 \(\operatorname{supp}{\varphi }_{j} \subset  K\) ;

(2)对于任意指标 \(\alpha  = \left( {{\alpha }_{1},\cdots ,{\alpha }_{n}}\right)\) 都有
\begin{equation*}
\mathop{\max }\limits_{{x \in  K}}\left| {{\partial }^{\alpha }{\varphi }_{j}\left( x\right)  - {\partial }^{\alpha }{\varphi }_{0}\left( x\right) }\right|  \rightarrow  0,\;\left( {j \rightarrow  \infty }\right) .
\end{equation*}

带上述收敛性的线性空间 \({C}_{0}^{\infty }\left( \Omega \right)\) ,称为基本空间 \(\mathcal{D}\left( \Omega \right)\) .
\end{definition}
\begin{note}
  我们只在基本空间 \(\mathcal{D}\left( \Omega \right)\) 上引进了收敛性,并由此给定了拓扑,但上述收敛性并不能由任何模给出,即 \(\mathcal{D}\left( \Omega \right)\) 不是赋范线性空间.


\end{note}

\begin{note}
  基本空间 \(\mathcal{D}\left( \Omega \right)\) 为完备空间,即每个Cauchy序列都有极限,参见习题3.1.3.
\end{note}


\subsection{广义函数的定义及其收敛性}

现在我们可以引进广义函数的定义如下
\begin{definition}
  称 \(\mathcal{D}\left( \Omega \right)\) 上的连续线性泛函 \(f\) 为 \(\mathcal{D}\left( \Omega \right)\) 上的广义函数 (或分布) \(f\) , 即 \(f : \mathcal{D}\left( \Omega \right)  \rightarrow  \mathbb{R}\) 满足

(1) 线性 \(f\left( {{a\varphi } + {b\psi }}\right)  = {af}\left( \varphi \right)  + {bf}\left( \psi \right) ,\forall a,b \in  \mathbb{R},\forall \varphi ,\psi  \in  \mathcal{D}\left( \Omega \right)\) .

(2) 连续性 当 \({\varphi }_{j} \rightarrow  \varphi \left( {\mathcal{D}\left( \Omega \right) }\right)\) 时,有 \(f\left( {\varphi }_{j}\right)  \rightarrow  f\left( \varphi \right)\) .

\(\mathcal{D}\left( \Omega \right)\) 上广义函数的全体记为 \({\mathcal{D}}^{\prime }\left( \Omega \right)\) ,称为广义函数空间 (或分布空间).
\end{definition}


\begin{example}
前面的 \(\delta\) 函数可以看成 \(\mathcal{D}\left( \Omega \right)\) 上如下定义的广义函数
\begin{equation*}
\left\langle  {\delta \left( {x - {x}^{0}}\right) ,\varphi \left( x\right) }\right\rangle   = \varphi \left( {x}^{0}\right) ,\;\forall \varphi  \in  \mathcal{D}\left( \Omega \right) ,
\end{equation*}

请读者自己证明 \(\delta \left( {x - {x}^{0}}\right)  \in  {\mathcal{D}}^{\prime }\left( \Omega \right)\) . \(\delta\) 函数通常又称为Dirac分布或测度,在技术科学中则称为Dirac脉冲符号或Dirac脉冲函数.
\end{example}
\begin{proof}
    线性性质显然成立，下面证连续性.

    若${\varphi }_{j} \rightarrow  \varphi \left( {\mathcal{D}\left( \Omega \right) }\right)$,那么存在紧集 $K$,若$x_0\in K$,则
\[\left\langle  {\delta \left( {x - {x}^{0}}\right) ,\varphi_j \left( x\right) }\right\rangle   = \varphi_j \left( {x}^{0}\right)\to \varphi(x_0) =\left\langle  {\delta \left( {x - {x}^{0}}\right) ,\varphi \left( x\right) }\right\rangle \]

若$x_0\notin K$ ,则
\[\left\langle  {\delta \left( {x - {x}^{0}}\right) ,\varphi_j \left( x\right) }\right\rangle   = \varphi_j \left( {x}^{0}\right)=0\to 0= \varphi(x_0) =\left\langle  {\delta \left( {x - {x}^{0}}\right) ,\varphi \left( x\right) }\right\rangle \]
\end{proof}
\begin{example}
  记 \({L}_{loc}^{1}\left( \Omega \right)\) 为在 \(\Omega\) 的任一有界集上Lebesgue可积的函数(简称为 \(\Omega\) 上局部可积函数)的全体. 则每个 \(f\left( x\right)  \in  {L}_{loc}^{1}\left( \Omega \right)\) 都对应一个广义函数 \(f\)
\begin{equation}
\langle f,\varphi \rangle  = {\int }_{\Omega }f\left( x\right) {\varphi \dif x},\;\forall \varphi  \in  \mathcal{D}\left( \Omega \right) . \label{3.5}
\end{equation}

\end{example}
\begin{proof}
   线性性是显然的,下面验证连续性. 由于 \(f\left( x\right)  \in  {L}_{loc}^{1}\left( \Omega \right)\) ,故对任意相对于 \(\Omega\) 的紧集 \(K\) ,有
\begin{equation*}
{\int }_{K}\left| {f\left( x\right) }\right| {\dif x} < \infty .
\end{equation*}

若 \({\varphi }_{j} \rightarrow  {\varphi }_{0}\left( {\mathcal{D}\left( \Omega \right) }\right)\) ,则存在紧集 \({K}_{0} \subset  \Omega\) ,使得
\begin{equation*}
\text{ supp }{\varphi }_{j} \subset  {K}_{0}\text{ ,且 }\mathop{\max }\limits_{{x \in  {K}_{0}}}\left| {{\varphi }_{j}\left( x\right)  - {\varphi }_{0}\left( x\right) }\right|  \rightarrow  0\left( {j \rightarrow  \infty }\right) \text{ . }
\end{equation*}

从而由Lebesgue控制收敛定理知
\begin{equation*}
\left| {\left\langle  {f,{\varphi }_{j}}\right\rangle   - \left\langle  {f,{\varphi }_{0}}\right\rangle  }\right|  \leqslant  {\int }_{{K}_{0}}\left| {f\left( x\right) }\right| \left| {{\varphi }_{j}\left( x\right)  - {\varphi }_{0}\left( x\right) }\right| {\dif x} \rightarrow  0,\;j \rightarrow  \infty .
\end{equation*}

由定义知 \(f\) 的连续性.
\end{proof}
\begin{note}
例3.1.5表明,广义函数是局部可积函数的推广,即广义函数空间 \({\mathcal{D}}^{\prime }\left( \Omega \right)\) 包含所有的局部Lebesgue可积函数,故 \({\mathcal{D}}^{\prime }\left( \Omega \right)\) 非空. 但并不是所有的函数都是广义函数. 事实上, 普通的不可测函数并不能看成广义函数.
\end{note}

 \({\mathcal{D}}^{\prime }\left( \Omega \right)\) 除了包含经典函数之外还包含 \(\delta\) 函数这种“非经典”的广义函数,所以 \({\mathcal{D}}^{\prime }\left( \Omega \right)\) 确实扩大了经典的函数空间. 实际上, \(\delta\) 函数不能用任何一个局部可积函数通过例3.1.5的方式来给出. 若不然,则存在某个 \(\Omega\) 上局部Lebesgue可积函数 \(f\left( x\right)\) 使得
\begin{equation*}
\langle \delta ,\varphi \rangle  = {\int }_{\Omega }f\left( x\right) \varphi \left( x\right) {\dif x} = \varphi \left( 0\right) ,\;\forall \varphi  \in  \mathcal{D}\left( \Omega \right) .
\end{equation*}

特别对例3.1.2中函数 \(\varphi \left( {x,a,a,1}\right)  \in  \mathcal{D}\left( \Omega \right)\) 有
\begin{equation*}
{\int }_{\left| x\right|  \leqslant  a}f\left( x\right) \varphi \left( {x,a,a,1}\right) {\dif x} = \varphi \left( {0,a,a,1}\right)  = \frac{1}{e},
\end{equation*}

在上式中令 \(a \rightarrow  0\) ,注意上式右边为与 \(a\) 无关的常数,而左边则因积分域缩为一点而趋于零,从而得到矛盾. 从 \(\delta\) 函数可以引出定义
\begin{definition}
  由局部Lebesgue可积函数 \(f\left( x\right)\) 给出的广义函数称为函数型广义函数或正规广义函数;非正规广义函数统称为奇型广义函数.

\end{definition}
\begin{note}
\(\delta\) 函数是奇型广义函数,但不是仅有的奇型广义函数,后面我们还会遇到其他的奇型广义函数.
\end{note}

在应用中, 我们往往还需要考虑用速降函数作为检验函数来定义广义函数.记速降函数空间(L. Schwartz S空间) \(\mathcal{Y}\) 为
\begin{equation}
\mathcal{S}\left( {\mathbb{R}}^{n}\right)  \mathrel{\text{ := }} \left\{  {\varphi \left( x\right) \left| \mathop{\sup }\limits_{{x \in  {\mathbb{R}}^{n}}}\right| {x}^{\alpha }{\partial }^{\beta }\varphi \left( x\right)  \mid   < \infty ,\;\forall \left| \alpha \right| ,\left| \beta \right|  = 0,1,2,\cdots }\right\}   \label{3.6}
\end{equation}

即当 \(\left| x\right|  \rightarrow  \infty\) 时 \(\mathcal{Y}\) 中函数 \(\varphi\) 比 \({\left| x\right| }^{-1}\) 的任何正幂更快地趋于零, \(\varphi\) 的各阶导数 \({\partial }^{\alpha }\varphi\) 也具有该性质. 故 \(\mathcal{Y}\) 称为速降函数空间, \(\mathcal{Y}\) 中元素 \(\varphi\) 称为速降函数. \(\mathcal{Y}\) 中函数列 \(\left\{  {{\varphi }_{j}\left( x\right) }\right\}\) 的收敛性定义为
\begin{equation*}
{\varphi }_{j}\left( x\right)  \rightarrow  0\left( \mathcal{S}\right)  \Leftrightarrow  \mathop{\lim }\limits_{{j \rightarrow  \infty }}\mathop{\sup }\limits_{{x \in  {\mathbb{R}}^{n}}}\left| {{x}^{\alpha }{\partial }^{\beta }{\varphi }_{j}\left( x\right) }\right|  = 0,\forall \alpha ,\beta  \in  {\mathbb{N}}^{n}.
\end{equation*}

显然, \(\mathcal{Y}\) 在上述收敛意义下为完备空间,参见习题3.1.3. 记 \(\mathcal{Y}\) 上的连续线性泛函空间为 \({\mathcal{S}}^{\prime }\) . 由于当 \(\left| x\right|  \rightarrow  \infty\) 时 \(\mathcal{Y}\) 上广义函数仅仅能具幂增率 \(O\left( {\left| x\right| }^{k}\right)\) ,故 \({\mathcal{S}}^{\prime }\) 称为缓增广义函数空间. 显然有
\begin{equation*}
\mathcal{D} \subset  \mathcal{S},\;{\mathcal{I}}^{\prime } \subset  {\mathcal{D}}^{\prime }.
\end{equation*}

\(\mathcal{Y}\) 中不属于 \(\mathcal{D}\) 的函数有 \({e}^{-{\left| x\right| }^{2}}\) . 实际上,对任意 \(k \in  \mathbb{N}\) 及多重指标 \(\alpha ,{\left| x\right| }^{k}{D}^{\alpha }{e}^{-{\left| x\right| }^{2}}\) 都是形如 \({C}_{1}{x}^{l}{e}^{-{\left| x\right| }^{2}}\) 项之和,故当 \(\left| x\right|  \rightarrow  \infty\) 时, \({\left| x\right| }^{l}{e}^{-{\left| x\right| }^{2}} \rightarrow  0\) ,从而 \({e}^{-{\left| x\right| }^{2}} \in  \mathcal{S}\) . \({\mathcal{D}}^{\prime }\) 中不属于 \({\mathcal{S}}^{\prime }\) 的广义函数有 \({e}^{{\left| x\right| }^{2}}\) .

此外,我们也可用 \(\mathcal{E} = {C}^{\infty }\left( \Omega \right)\) 作为基本空间来定义广义函数 \({\mathcal{E}}^{\prime }\left( \Omega \right)\) . 而且
\begin{equation*}
\mathcal{D} \subset  \mathcal{S} \subset  \mathcal{E},{\mathcal{E}}^{\prime } \subset  {\mathcal{S}}^{\prime } \subset  {\mathcal{D}}^{\prime }.
\end{equation*}

为讨论广义函数空间的结构, 我们还有如下重要估计.
\begin{proposition}
      \(\mathcal{E}\left( \Omega \right)\) 的线性泛函 \(f\) 连续当且仅当存在相对于 \(\Omega\) 的紧集 \(K\) ,存在常数 \(C\) 及非负整数 \(m\) 使得对任意 \(\varphi  \in  \mathcal{E}\left( \Omega \right)\) ,有
\begin{equation}
\left| {\langle f,\varphi \rangle }\right|  \leqslant  C\mathop{\sum }\limits_{{\left| \alpha \right|  \leqslant  m}}\mathop{\sup }\limits_{{x \in  K}}\left| {{\partial }^{\alpha }\varphi \left( x\right) }\right| . 
\end{equation}
\end{proposition}

\begin{proof}
  充分性是显然的,下证必要性.

取一列紧集序列:$K_1\subset{K_2\subset{\cdots \subset{\Omega}}}$,使得$\bigcup\limits_{i=1}^{\infty}K_i=\Omega$.
  
  用反证法,若不然,则存在\(\{\varphi_j(x)\}\subset{C^\infty(\Omega)}\) 使得对任意 \(j \in  \mathbb{N}\) ,存在 \(C_j\) 使得 \(C_j\to \infty,j\to \infty\) ,有
\begin{equation*}
\left| {\langle f,\varphi_j \rangle }\right|  >  C_j\mathop{\sup }\limits_{{x \in  K_j,|\alpha|\leqslant j}}\left| {{\partial }^{\alpha }\varphi_j \left( x\right) }\right|
\end{equation*}
  $\{\varphi_j\}$满足下列两种情况,

  (1)当$j$充分大时,$C_j\mathop{\sup }\limits_{{x \in  K_j,|\alpha|\leqslant j}}\left| {{\partial }^{\alpha }\varphi_j \left( x\right) }\right|>0$(因为递增),那么令\[
\psi_j(x)=\frac{\varphi_j(x)}{C_j\mathop{\sup }\limits_{{x \in  K_j,|\alpha|\leqslant j}}\left| {{\partial }^{\alpha }\varphi_j \left( x\right) }\right|}
  \]
从而成立$|\langle f,\psi_j\rangle|>1$,考虑任意紧集$K,\exists j_0\in \mathbb{N},K\subset{ K_j},j>j_0$,从而有
\[
\lim\limits_{j\to \infty}\mathop{\sup }\limits_{{x \in  K,|\alpha|\leqslant j}}\left| {{\partial }^{\alpha }\psi_j \left( x\right) }\right|<\lim\limits_{j\to \infty}\frac{1}{C_j}=0
\]
所以$\psi_j\to 0(\mathcal{E}(\Omega))$,这与$|\langle f,\psi_j\rangle|>1$矛盾.

(2)当$j$充分大时,$C_j\mathop{\sup }\limits_{{x \in  K_j,|\alpha|\leqslant j}}\left| {{\partial }^{\alpha }\varphi_j \left( x\right) }\right|=0$,那么令$\lambda_j=|\langle f,\varphi_j\rangle|>0$,令\[
\psi_j(x)=\frac{\varphi_j(x)}{\lambda_j}
  \]
  考虑任意紧集$K,\exists j_0\in \mathbb{N},K\subset{ K_j},j>j_0$,从而有
\[
\lim\limits_{j\to \infty}\mathop{\sup }\limits_{{x \in  K,|\alpha|\leqslant j}}\left| {{\partial }^{\alpha }\psi_j \left( x\right) }\right|=0
\]
所以$\psi_j\to 0(\mathcal{E}(\Omega))$,这与$|\langle f,\psi_j\rangle|=1$矛盾.
\end{proof}
\begin{theorem}
  \(\mathcal{D}\left( \Omega \right)\) 的线性泛函 \(f\) 连续当且仅当对任意相对于 \(\Omega\) 的紧集 \(K\) ,存在常数 \(C\) 及非负整数 \(m\) 使得对任意满足 \(\operatorname{supp}\varphi  \subset  K\) 的 \(\varphi  \in  \mathcal{D}\left( \Omega \right)\) ,有
\begin{equation}
\left| {\langle f,\varphi \rangle }\right|  \leqslant  C\mathop{\sum }\limits_{{\left| \alpha \right|  \leqslant  m}}\mathop{\sup }\limits_{{x \in  K}}\left| {{\partial }^{\alpha }\varphi \left( x\right) }\right| . \label{3.7}
\end{equation}
\end{theorem}
\begin{proof}
  充分性是显然的,下证必要性. 用反证法,若不然,则存在紧集 \(K\) 使得对任意 \(j \in  \mathbb{N}\) ,存在 \({\varphi }_{j} \in  \mathcal{D}\left( \Omega \right)\) 使得 \(\operatorname{supp}{\varphi }_{j} \subset  K\) ,有
\begin{equation*}
\left| {\langle f,\varphi_j \rangle }\right|  >  C\mathop{\sum }\limits_{{\left| \alpha \right|  \leqslant  j}}\mathop{\sup }\limits_{{x \in  K}}\left| {{\partial }^{\alpha }\varphi_j \left( x\right) }\right| (\forall C\in \mathbb{R}). 
\end{equation*}
  
  进而取$C=1=\frac{\langle f,\varphi_j \rangle}{\langle f,\varphi_j \rangle}$,那么有
\begin{equation*}
\mathop{\sup }\limits_{{x \in  K}}\left| {{\partial }^{\alpha }{\varphi }_{j}}\right|  \leqslant  \frac{1}{j},\;\left( {\left| \alpha \right|  \leqslant  j}\right) .
\end{equation*}

因此, \(\left\{  {\varphi }_{j}\right\}\) 在 \(\mathcal{D}\left( \Omega \right)\) 中收敛到 0, 从而
\begin{equation*}
\left\langle  {f,{\varphi }_{j}}\right\rangle   \rightarrow  0,\;j \rightarrow  \infty .
\end{equation*}

矛盾, 故定理得证.
\end{proof}

实际上, \({L}_{loc}^{1}\left( \Omega \right)\) 中的每个元素有多个代表元,从函数型广义函数的例子可以看出, 这些不同的代表元都给出同一个广义函数.因此, 在研究广义函数的性质时, 我们往往需要用到两个广义函数相等的概念.
\begin{definition}
  设 \(\omega  \subset  \Omega\) 为 \(\Omega\) 的子集. 称 \(f \in  {\mathcal{D}}^{\prime }\left( \Omega \right)\) 在 \(\omega\) 内为零广义函数,如果
\begin{equation*}
\langle f,\varphi \rangle  = 0,\;\forall \varphi  \in  \mathcal{D}\left( \omega \right) .
\end{equation*}

称 \(f,g \in  {\mathcal{D}}^{\prime }\left( \Omega \right)\) 在 \(\omega\) 中相等,如果 \(f - g\) 为 \(\omega\) 内的零广义函数; 称 \(f\) 与 \(g\) 在 \(\Omega\) 中相等, 如果 \(f\) 与 \(g\) 在 \(\Omega\) 的任意子集 \(\omega\) 中相等. 称使得广义函数 \(f\) 为零广义函数的最大开子集 \(\omega  \subset  \Omega\) 的余集为广义函数 \(f\) 的支集或台,记为 \(\operatorname{supp}f\) .

\end{definition}
\begin{definition}
  称 \({\mathcal{D}}^{\prime }\left( \Omega \right)\) 中广义函数列 \(\left\{  {f}_{n}\right\}\) 弱收敛到 \(f\) ,记为 \({f}_{n} \rightharpoonup  f\) ,如果
\begin{equation*}
\mathop{\lim }\limits_{{n \rightarrow  \infty }}\left\langle  {{f}_{n},\varphi }\right\rangle   = \langle f,\varphi \rangle ,\;\forall \varphi  \in  \mathcal{D}\left( \Omega \right) .
\end{equation*}

习惯上也称 \(f\) 为 \(\left\{  {f}_{n}\right\}\) 的弱极限.
\end{definition}
\begin{example}
  设
\begin{equation*}
{\delta }_{h}\left( x\right)  = \left\{  \begin{array}{ll} \frac{1}{h}, & \left| x\right|  < \frac{h}{2} \\  0, & \text{ 其他, } \end{array}\right.
\end{equation*}

则当 \(h \rightarrow  0\) 时, \({\delta }_{h} \rightarrow  \delta \left( {{\mathcal{D}}^{\prime }\left( \mathbb{R}\right) }\right)\) .

\end{example}
\begin{proof}
  对任意 \(\varphi  \in  {C}_{0}^{\infty }\left( \mathbb{R}\right)\) ,有

\begin{equation*}
{\int }_{\mathbb{R}}{\delta }_{h}\left( x\right) \varphi \left( x\right) {\dif x} = \frac{1}{h}{\int }_{-\frac{h}{2}}^{\frac{h}{2}}\varphi \left( x\right) {\dif x} = \varphi \left( \xi \right) ,\;\xi  \in  \left( {-\frac{h}{2},\frac{h}{2}}\right) .
\end{equation*}

当 \(h \rightarrow  0\) 时, \(\varphi \left( \xi \right)  \rightarrow  \varphi \left( 0\right)\) . 故 \(\mathop{\lim }\limits_{{h \rightarrow  0}}\left\langle  {{\delta }_{h},\varphi }\right\rangle   = \varphi \left( 0\right)  = \langle \delta ,\varphi \rangle\) .

\end{proof}
\begin{note}
  \(\delta\) 函数可视为函数列 \(\left\{  {{\delta }_{h}\left( x\right) }\right\}\) 在 \({\mathcal{D}}^{\prime }\left( \mathbb{R}\right)\) 中的极限. 这说明正规广义函数列可能收敛到奇型广义函数. 实际上, 奇型广义函数是正规广义函数的极限. 
\end{note}
\begin{example}
  设 \({f}_{v}\left( x\right)  = \frac{1}{\pi }\frac{\sin {vx}}{x}\) ,则当 \(v \rightarrow  \infty\) 时有 \({f}_{v} \rightarrow  \delta \left( {{\mathcal{D}}^{\prime }\left( \mathbb{R}\right) }\right)\) .

\end{example}

\begin{proof}
因为有\( \mathop{\lim }\limits_{{T \rightarrow  \infty }}\frac{1}{\pi }{\int }_{-T}^{T}\frac{\sin {jx}}{x}\mathrm{\;d}x = 1, \)所以 \( \forall \varphi  \in  \mathcal{D}\left( \mathbb{R}\right) \) ,存在 \( {T}_{0} > 0 \) ,使得 \( \operatorname{supp}\left( \varphi \right)  \subset  \left\lbrack  {-{T}_{0},{T}_{0}}\right\rbrack \) .

一方面,
当 \( T > {T}_{0} \) 时,
\( {\int }_{-\infty }^{\infty }{f}_{j}\left( x\right)  \cdot  \varphi \left( x\right) \mathrm{d}x = {\int }_{-T}^{T}{f}_{j}\left( x\right)  \cdot  \varphi \left( x\right) \mathrm{d}x \)

另一方面, \( \forall \varepsilon  > 0 \) ,取 \( {T}_{1} \) 足够大,以致 \( T > {T}_{1} \) 时,
\( \left| {\frac{1}{\pi }{\int }_{-T}^{T}\frac{\sin {jx}}{x}\mathrm{\;d}x - 1}\right|  < \frac{\varepsilon }{2}. \)从而当 \( T > \max \left\{  {{T}_{0},{T}_{1}}\right\} \) 时,
\( \left| {\left\langle  {{f}_{j},\varphi }\right\rangle   - \varphi \left( 0\right) }\right|  \leqslant  \left| {\frac{1}{\pi }{\int }_{-T}^{T}\frac{\sin {jx}}{x}\left\lbrack  {\varphi \left( x\right)  - \varphi \left( 0\right) }\right\rbrack  \mathrm{d}x}\right|  + \frac{\varepsilon }{2}\left| {\varphi \left( 0\right) }\right| \)
\( = \frac{1}{\pi }\left| {{\int }_{0}^{T}\sin {jx}\frac{\varphi \left( x\right)  + \varphi \left( {-x}\right)  - {2\varphi }\left( 0\right) }{x}\mathrm{\;d}x}\right|  + \frac{\varepsilon }{2}\left| {\varphi \left( 0\right) }\right| . \)

固定 \( T \) ,由 Riemann-Lebesgue 引理,存在正整数 \( {n}_{0} \) ,当 \( j > {n}_{0} \) 时,
\( \frac{1}{\pi }\left| {{\int }_{0}^{T}\sin {jx}\frac{\varphi \left( x\right)  + \varphi \left( {-x}\right)  - {2\varphi }\left( 0\right) }{x}\mathrm{\;d}x}\right|  < \frac{\varepsilon }{2}, \)
于是得 \( \left\langle  {{f}_{j},\varphi }\right\rangle   \rightarrow  \varphi \left( 0\right)  = \langle \delta ,\varphi \rangle \left( {j \rightarrow  \infty }\right) \left( {\forall \varphi  \in  \mathcal{D}\left( \mathbb{R}\right) }\right) \) .
\end{proof}


实际上,上述弱收敛在 \({\mathcal{D}}^{\prime }\left( \Omega \right)\) 上定义了一个弱拓扑. 我们自然会问, \({\mathcal{D}}^{\prime }\left( \Omega \right)\) 关于上述弱拓扑是否完备? 对此, 我们有
\begin{theorem}
  设 \(\Omega  \subset  {\mathbb{R}}^{n}\) 为开集. \({\mathcal{D}}^{\prime }\left( \Omega \right)\) 为弱完备空间. 即 \({\mathcal{D}}^{\prime }\left( \Omega \right)\) 中每个基本列 \(\left\{  {f}_{n}\right\}\) 有弱极限 \(f \in  {\mathcal{D}}^{\prime }\left( \Omega \right)\) .

\end{theorem}
\begin{proof}
  要证明 \({\mathcal{D}}^{\prime }\left( \Omega \right)\) 弱完备,只需证明 \({\mathcal{D}}^{\prime }\left( \Omega \right)\) 中任意基本列都有弱极限. 设 \(\left\{  {f}_{j}\right\}\) 为 \({\mathcal{D}}^{\prime }\left( \Omega \right)\) 中基本列. 则由基本列的定义可知,对任给的正数 \(\varepsilon  > 0\) 及 \(\varphi  \in\)  \(\mathcal{D}\left( \Omega \right)\) ,存在 \(N\left( {\varepsilon ,\varphi }\right)\) 使得当 \(j,k \geq  N\left( {\varepsilon ,\varphi }\right)\) 时有
\begin{equation*}
\left| \left\langle  {{f}_{j} - {f}_{k},\varphi }\right\rangle  \right|  < \varepsilon
\end{equation*}

因此,对每个 \(\varphi  \in  \mathcal{D}\left( \Omega \right) ,\left\{  \left\langle  {{f}_{j},\varphi }\right\rangle  \right\}\) 为一基本列. 从而 \(\left\{  {f}_{j}\right\}\) 有弱极限存在,记为 \(f\) . 即
\begin{equation*}
\langle f,\varphi \rangle  \mathrel{\text{ := }} \mathop{\lim }\limits_{{j \rightarrow  \infty }}\left\langle  {{f}_{j},\varphi }\right\rangle  .
\end{equation*}

显然,上述弱极限 \(f\) 为线性的. 下证 \(f\) 连续. 用反证法,若不然,即 \(f\) 不连续. 则由定理 3.1.2 知,存在相对于 \(\Omega\) 的紧集 \(K\) 使得对任意 \(j \in  \mathbb{N}\) ,存在 \({\varphi }_{j} \in  \mathcal{D}\left( \Omega \right)\) 满足 \(\operatorname{supp}{\varphi }_{j} \subset  K\) ,但
\begin{equation*}
\left| \left\langle  {f,{\varphi }_{j}}\right\rangle  \right|  > {4}^{j}\mathop{\sum }\limits_{{\left| \alpha \right|  \leqslant  j}}\mathop{\sup }\limits_{{x \in  K}}\left| {{\partial }^{\alpha }{\varphi }_{j}\left( x\right) }\right| .
\end{equation*}

由齐次性,不妨设 \(\left| \left\langle  {f,{\varphi }_{j}}\right\rangle  \right|  = 1\) . 从而有
\begin{equation*}
\left| {{\partial }^{\alpha }{\varphi }_{j}\left( x\right) }\right|  \leqslant  \frac{1}{{4}^{j}},\;\left| \alpha \right|  < j,x \in  \Omega .
\end{equation*}

因此, \({\varphi }_{j} \rightarrow  0\left( {\mathcal{D}\left( \Omega \right) }\right)\) .

取 \({\psi }_{j} = {2}^{j}{\varphi }_{j} \in  \mathcal{D}\left( \Omega \right)\) ,显然有 \({\psi }_{j} \rightarrow  0\left( {\mathcal{D}\left( \Omega \right) }\right)\) ,且 \(\left| \left\langle  {f,{\psi }_{j}}\right\rangle  \right|  \rightarrow  \infty\) . 选 \(\left\{  {\psi }_{j}\right\}\) 的子列 \(\left\{  {\psi }_{1j}\right\}\) 及 \(\left\{  {f}_{j}\right\}\) 的子列 \(\left\{  {f}_{1j}\right\}\) 如下:

首先,由 \(\left| \left\langle  {f,{\psi }_{j}}\right\rangle  \right|  \rightarrow  \infty\) 知,存在 \(\left\{  {\psi }_{j}\right\}\) 的子列 \(\left\{  {\psi }_{1j}\right\}\) 使得 \(\left| \left\langle  {f,{\psi }_{1j}}\right\rangle  \right|  > 1\) . 因此, 若取 \({f}_{11} = {f}_{1}\) ,则由 \({\psi }_{1j} \rightarrow  0\left( {j \rightarrow  \infty }\right)\) 知,当 \(j \rightarrow  \infty\) 时,有 \(\left\langle  {{f}_{11},{\psi }_{1j}}\right\rangle   \rightarrow  0\) ,. 从而可选取到 \(\left\{  {\psi }_{1j}\right\}\) 中的元素,不妨记为 \({\psi }_{11}\) ,使得 \(\left| \left\langle  {{f}_{11},{\psi }_{11}}\right\rangle  \right|  < 1\) . 至此,我们可取 \({\psi }_{11}\) 和 \({f}_{11}\) 使得 \(\left| {\left\langle  {f,{\psi }_{11}}\right\rangle   > 1\text{ 及 }\left| {\left\langle  {{f}_{11},{\psi }_{11}}\right\rangle   < 1\text{ . }}\right\rangle  }\right\rangle\)

进一步,若 \({\psi }_{1,j - 1}\) 及 \({f}_{1,j - 1}\) 已选定,则可取 \({\psi }_{1j}\) ,使得
\begin{equation}
\left\{  \begin{array}{ll} \left| \left\langle  {{f}_{1\nu },{\psi }_{1j}}\right\rangle  \right|  < \frac{1}{{2}^{j - \nu }}, & \nu  = 1,\cdots ,j - 1, \\  \left| \left\langle  {f,{\psi }_{1j}}\right\rangle  \right|  > \mathop{\sum }\limits_{{\nu  = 1}}^{{j - 1}}\left| \left\langle  {f,{\psi }_{1\nu }}\right\rangle  \right|  + j. &  \end{array}\right.  \label{3.8}
\end{equation}

因为 \({\psi }_{\nu } \rightarrow  0\left( {\mathcal{D}\left( \Omega \right) }\right)\) ,故对固定的 \({f}_{1j}\) ,有 \(\left\langle  {{f}_{1j},{\psi }_{\nu }}\right\rangle   \rightarrow  0\left( {\nu  \rightarrow  \infty }\right)\) . 这意味着满足 (3.8) 的 \({\psi }_{1j}\) 完全可以取到. 为选 \({f}_{1j}\) ,注意对每个 \(\varphi  \in  \mathcal{D}\left( \Omega \right)\) ,有 \(\left\langle  {{f}_{\nu },\varphi }\right\rangle   \rightarrow\)  \(\langle f,\varphi \rangle \left( {\nu  \rightarrow  \infty }\right)\) . 故由\ref{3.8}中第二式知,存在 \({f}_{1j}\) 使得
\begin{equation*}
\left| \left\langle  {{f}_{1j},{\psi }_{1j}}\right\rangle  \right|  > \mathop{\sum }\limits_{{\nu  = 1}}^{{j - 1}}\left| \left\langle  {{f}_{1j},{\psi }_{1\nu }}\right\rangle  \right|  + j.
\end{equation*}

选定子列 \(\left\{  {\psi }_{1j}\right\}\) 及 \(\left\{  {f}_{1j}\right\}\) 后,令
\begin{equation*}
\psi  = \mathop{\sum }\limits_{{j = 1}}^{\infty }{\psi }_{1j}
\end{equation*}

则 \(\psi  \in  \mathcal{D}\left( \Omega \right)\) . 实际上,考虑余项
\begin{equation*}
{R}_{N} = \mathop{\sum }\limits_{{i = N}}^{\infty }{\psi }_{1i}
\end{equation*}

则有
\begin{equation*}
\left| {R}_{N}\right|  \leqslant  \mathop{\sum }\limits_{{i = N}}^{\infty }\left| {\psi }_{1i}\right|  \leqslant  \mathop{\sum }\limits_{{i = N}}^{\infty }\left| {\psi }_{i}\right|  \leqslant  \mathop{\sum }\limits_{{i = N}}^{\infty }\frac{1}{{2}^{i}}.
\end{equation*}

故 \({R}_{N} \rightarrow  0\left( {N \rightarrow  \infty }\right)\) . 因此, \(\psi\) 为绝对且一致收敛的级数的和函数. 同理可证对任意 \(\alpha\) ,级数 \(\mathop{\sum }\limits_{{i = 1}}^{\infty }{\partial }^{\alpha }{\psi }_{1j}\) 绝对且一致收敛.

由\ref{3.8}知
\begin{equation*}
\left| {\mathop{\sum }\limits_{{\nu  = j + 1}}^{\infty }\left\langle  {{f}_{1j},{\psi }_{1\nu }}\right\rangle  }\right|  < \mathop{\sum }\limits_{{\nu  = j + 1}}^{\infty }\frac{1}{{2}^{\nu  - j}} = 1.
\end{equation*}

由于
\begin{equation*}
\left\langle  {{f}_{1j},\psi }\right\rangle   = \mathop{\sum }\limits_{{\nu  = 1}}^{{j - 1}}\left\langle  {{f}_{1j},{\psi }_{1\nu }}\right\rangle   + \left\langle  {{f}_{1j},{\psi }_{1j}}\right\rangle   + \mathop{\sum }\limits_{{\nu  = j + 1}}^{\infty }\left\langle  {{f}_{1j},{\psi }_{1\nu }}\right\rangle  ,
\end{equation*}

故有
\begin{equation*}
\left| \left\langle  {{f}_{1j},\psi }\right\rangle  \right|  \geq  \left| \left\langle  {{f}_{1j},{\psi }_{1j}}\right\rangle  \right|  - \left| {\mathop{\sum }\limits_{{\nu  = 1}}^{{j - 1}}\left\langle  {{f}_{1j},{\psi }_{1\nu }}\right\rangle  }\right|  - \left| {\mathop{\sum }\limits_{{\nu  = j + 1}}^{\infty }\left\langle  {{f}_{1j},{\psi }_{1\nu }}\right\rangle  }\right|  > j - 1.
\end{equation*}

这表明当 \(j \rightarrow  \infty\) 时, \(\left| \left\langle  {{f}_{1j},\psi }\right\rangle  \right|  \rightarrow  \infty\) ,从而与假定 \(\left\langle  {{f}_{j},\psi }\right\rangle\) 的收敛性矛盾. 故 \(f\) 为连续线性泛函. 至此, 定理得证.
\end{proof}
\begin{theorem}[单位分解定理]
    若 \( A \subset  {\mathbb{R}}^{n} \) ,而且 \( \mathcal{O} \) 是 \( A \) 的一个开覆盖,那么必
有一族 \( {C}^{\infty } \) 函数 \( \mathcal{F} \) ,使得 \( \forall \varphi  \in  \mathcal{F} \) 定义在包含 \( A \) 的一个开集上,
具有下列性质:

(1) \( 0 \leq  \varphi \left( x\right)  \leq  1\;\left( {\forall x \in  A}\right) \) ;

(2) \( \forall x \in  A,\exists \) 含 \( x \) 的开集 \( V \) ,使得只有有穷多个 \( \varphi  \in  \mathcal{F} \) 在其
上非零;

(3) \( \mathop{\sum }\limits_{{\varphi  \in  \mathcal{F}}}\varphi \left( x\right)  \equiv  1\;\left( {\forall x \in  A}\right) \) ;

(4) \( \forall \varphi  \in  \mathcal{F},\exists U \in  \mathcal{O} \) ,使得 \( \operatorname{supp}\left( \varphi \right)  \subset  U \) .
\end{theorem}

\subsection{习 题}
\begin{problemset}
\item  求证：$\delta$ 函数不是局部可积函数．
\begin{proof}
    用反证法．如果 $\exists f(x) \in L_{\mathrm{loc}}^1$ ，使得
$$
\langle\delta, \varphi\rangle=\int_{\Omega} f(x) \varphi(x) \mathrm{d} x \quad(\forall \varphi \in \mathcal{D}(\Omega)),
$$

取 $\varphi_k(x)=j(k x)(\forall k \in \mathbb{N})$ ，其中
$$
j(x)= \begin{cases}C_n \mathrm{e}^{-\frac{1}{1-|x|^2}}, & |x|<1 \\ 0, & |x| \geqslant 1\end{cases}
$$

则有：一方面，根据积分的绝对连续性
$$
\left|\int_{B(\theta, 1 / k)} f(x) \varphi_k(x) \mathrm{d} x\right| \leqslant C_n \int_{B(\theta, 1 / k)}|f(x)| \mathrm{d} x \rightarrow 0 \quad(k \rightarrow \infty),
$$


故有
$$
\int_{\Omega} f(x) \varphi_k(x) \mathrm{d} x=\int_{B(\theta, 1 / k)} f(x) \varphi_k(x) \mathrm{d} x \rightarrow 0 \quad(k \rightarrow \infty)
$$

另一方面，
$$
\left\langle\delta, \varphi_k\right\rangle=\varphi_k(0)=C_n \mathrm{e}^{-1} \rightarrow C_n \mathrm{e}^{-1} \quad(k \rightarrow \infty)
$$

于是有 $0=C_n \mathrm{e}^{-1} \neq 0$ ，矛盾．
\end{proof}
\item 设 \(1 \leqslant  p < \infty\) ,证明 \({C}_{0}^{\infty }\left( \Omega \right)\) 在 \({L}^{p}\left( \Omega \right)\) 中稠密.
\begin{proof}
   由Lusin定理知$C_0^0(\Omega)$在$L^p(\Omega)$稠密,由定理3.1.1知$C_0^\infty$在$C^0_0(\Omega)$中稠密,进而 \({C}_{0}^{\infty }\left( \Omega \right)\) 在 \({L}^{p}\left( \Omega \right)\) 中稠密.
\end{proof}
\item 已知 \(\varphi \left( x\right)  \in  \mathcal{D}\left( {\mathbb{R}}^{n}\right)\) . 试问下列序列

(1) \(\frac{1}{k}\varphi \left( x\right) ,k = 1,2,\cdots\) ;

(2) \(\frac{1}{k}\varphi \left( {kx}\right) ,k = 1,2,\cdots\) ;

(3) \(\frac{1}{k}\varphi \left( \frac{x}{k}\right) ,k = 1,2,\cdots\)

是否在 \(\mathcal{D}\left( {\mathbb{R}}^{n}\right)\) 中收敛到零元.
\begin{solution}
(1)
\[    \lim_{k\to \infty}\sup\limits_{x\in \operatorname{supp}\varphi}|\partial^\alpha \frac{1}{k}\varphi(x)|=0,\forall |\alpha|\in \mathbb{N}\]
(2)
\[    \lim_{k\to \infty}\sup\limits_{x\in \operatorname{supp}\varphi}|\partial^\alpha \frac{1}{k}\varphi(kx)|=\infty,\forall |\alpha|>1\]
(3)紧支集不存在.
\end{solution}
\item 证明 \(\delta\) 分布为 \(\mathcal{D}\) 上的连续线性泛函.
\item 设 \(\Omega  \subset  {\mathbb{R}}^{n}\) 为开集. 证明基本空间 \(\mathcal{D}\left( \Omega \right) ,\mathcal{S}\left( {\mathbb{R}}^{n}\right)\) 为完备空间.
\begin{proof}
设 $\left\{\phi_j\right\}$ 为 $\mathcal{D}(\Omega)$ 中的柯西序列。由 $\mathcal{D}(\Omega)$ 的拓扑性质，存在紧集 $K \subset \Omega$ 使得所有 $\phi_j$ 的支撑包含于 $K$ ，且对任意多指标 $\alpha$ ，序列 $\left\{\partial^\alpha \phi_j\right\}$ 在 $K$ 上为一致柯西序列。

对任意多指标 $\alpha$ ，由于 $\left\{\partial^\alpha \phi_j\right\}$ 在 $K$ 上为一致柯西序列，存在连续函数 $\psi_\alpha$ 在 $K$ 上使得 $\partial^\alpha \phi_j \rightarrow \psi_\alpha$ 一致于 $K$ 。定义 $\phi(x)=\psi_0(x)$ 对于 $x \in K$ ，并延拓 $\phi$ 为 0 于 $\Omega \backslash K$ ，则 $\phi$ 的支撑包含于 $K$ 。

需证明 $\phi$ 光滑且 $\partial^\alpha \phi=\psi_\alpha$ 。由于 $\phi_j \rightarrow \phi$ —致于 $K$ 且 $\partial^\alpha \phi_j \rightarrow \psi_\alpha$ 一致于 $K$ ，由导数与极限交换性质，可得 $\partial^\alpha \phi=\psi_\alpha$ 。具体地，对一阶导数，有：
$$
\phi_j(x)-\phi_j\left(x_0\right)=\int_{x_0}^x \partial_i \phi_j(t) d t \rightarrow \int_{x_0}^x \psi_i(t) d t
$$

左边收敛于 $\phi(x)-\phi\left(x_0\right)$ ，故 $\phi(x)-\phi\left(x_0\right)=\int_{x_0}^x \psi_i(t) d t$ ，因此 $\partial_i \phi=\psi_i$ 。通过归纳可得对任意 $\alpha$ ，有 $\partial^\alpha \phi=\psi_\alpha$ ，故 $\phi$ 光滑。

由构造，所有 $\phi_j$ 的支撑包含于 $K$ ，且对任意多指标 $\alpha$ ，有 $\partial^\alpha \phi_j \rightarrow \partial^\alpha \phi$ 一致于 $K$ ，故 $\phi_j \rightarrow \phi$ 在 $\mathcal{D}(\Omega)$ 中。

因此， $\mathcal{D}(\Omega)$ 是序列完备的。


\end{proof}
\item 证明若 \(f\left( x\right)  \in  {L}^{p}\left( {\mathbb{R}}^{n}\right)\) ,则 \(f\left( x\right)\) 对应一个 \({\mathcal{S}}^{\prime }\left( {\mathbb{R}}^{n}\right)\) 广义函数.
\item 证明

( 1 )广义函数列 \(\{ \delta \left( {x - k}\right) {\} }_{k = 1}^{\infty }\) 在 \({\mathcal{D}}^{\prime }\left( \mathbb{R}\right)\) 中收敛到零元.

( 2 )级数 \(\mathop{\sum }\limits_{{k = 0}}^{\infty }{a}_{k}\delta \left( {x - k}\right)\) 当 \(\left| {a}_{k}\right|  \leqslant  c{\left( 1 + k\right) }^{m}\) 时在 \({\mathcal{S}}^{\prime }\left( \mathbb{R}\right)\) 中收敛,其中 \(m\) 为固定正整数, \(c > 0\) 为固定常数.
\begin{proof}
    (1)$\forall \varphi\in \mathcal{D}(\Omega)$
    \[
    \lim\limits_{k\to \infty }\langle\delta(x-k), \varphi\rangle=\lim\limits_{k\to \infty} \varphi(k)=0=\langle 0,\varphi\rangle
    \]
    (2)$\forall \varphi\in \mathcal{S}(\mathbb{R})$
    \[
  \left|  \langle\mathop{\sum }\limits_{{k = 0}}^{N}{a}_{k}\delta \left( {x - k}\right),\varphi\rangle\right|=\left| \mathop{\sum }\limits_{{k = 0}}^{N }{a}_{k}\varphi(k)\right|\leqslant \mathop{\sum }\limits_{{k = 0}}^{N }c(1+k)^m\left|\varphi(k)\right|
    \]

而\[
\sup\limits_{x\in\mathbb{R}}|(1+x)^{m+2}\varphi(x)|<+\infty\Rightarrow |\varphi (k)|<\frac{C}{(1+k)^{m+2}}
\]
所以
    \[
  \left|  \langle\mathop{\sum }\limits_{{k = 0}}^{\infty}{a}_{k}\delta \left( {x - k}\right),\varphi\rangle\right|\leqslant \mathop{\sum }\limits_{{k = 0}}^{\infty}c(1+k)^m\left|\varphi(k)\right|<\mathop{\sum }\limits_{{k = 0}}^{\infty}cC(1+k)^{-2}<\infty
    \]
    所以收敛.
\end{proof}
\item 设 \(\Omega  \subset  {\mathbb{R}}^{n}\) 是一个开集,又设 \(K\) 是 \(\Omega\) 的一个紧子集. 证明存在函数 \(\varphi  \in  {C}_{0}^{\infty }\left( \Omega \right)\) 使得 \(0 \leqslant  \varphi \left( x\right)  \leqslant  1\) ,且 \(\varphi \left( x\right)\) 在 \({K}\) 的一个邻域内恒为 1 .
\begin{proof}
作开集 \( {\Omega }_{1} \) ,使得 \( K \subset  {\Omega }_{1},{\bar{\Omega }}_{1} \subset  \Omega \) ,记 \( \chi \left( x\right) \) 为 \( {\Omega }_{1} \) 的特征函数,令
\( \delta  = \frac{1}{3}\min \left\{  {\rho \left( {K,\partial {\Omega }_{1}}\right) ,\rho \left( {{\Omega }_{1},\partial \Omega }\right) }\right\}  , \)
\( \varphi \left( x\right)  = {\int }_{\Omega }\chi \left( y\right) {j}_{\delta }\left( {x - y}\right) \mathrm{d}y = {\int }_{{\Omega }_{1}}{j}_{\delta }\left( {x - y}\right) \mathrm{d}y, \)
则有 \( \varphi  \in  {C}_{0}^{\infty }\left( \Omega \right) \) ,且 \( \varphi \left( x\right)  = {\int }_{{\Omega }_{1}}{j}_{\delta }\left( {x - y}\right) \mathrm{d}y \) ,以及
\( 0 \leqslant  \varphi \left( x\right)  \leqslant  1 \) .

又令\( {K}_{\delta } = \{ x \in  \Omega  \mid  \rho \left( {x,K}\right)  \leqslant  \delta \} . \)\( \forall x \in  {K}_{\delta } \) ,注意到
\begin{equation*}
    \begin{aligned}
        \varphi \left( x\right) =
 {\int }_{\Omega }\chi \left( y\right) {j}_{\delta }\left( {x - y}\right) \mathrm{d}y = {\int }_{{\Omega }_{1}}{j}_{\delta }\left( {x - y}\right) \mathrm{d}y= {L}_{1} + {L}_{2}\\= {\int }_{{\Omega }_{1}\left( {\left| {y - x}\right|  \leqslant  \delta }\right) }{j}_{\delta }\left( {x - y}\right) \mathrm{d}y ={\int }_{\left| {y - x}\right|  \leqslant  \delta }{j}_{\delta }\left( {x - y}\right) \mathrm{d}y = {\int }_{\left| t\right|  \leqslant  \delta }{j}_{\delta }\left( t\right) \mathrm{d}y\;=1
    \end{aligned}
\end{equation*}  其中 \( {L}_{1} = {\int }_{{\Omega }_{1}\left( {\left| {y - x}\right|  \leqslant  \delta }\right) }{j}_{\delta }\left( {x - y}\right) \mathrm{d}y,{L}_{2} = {\int }_{{\Omega }_{1}\left( {\left| {y - x}\right|  > \delta }\right) }{j}_{\delta }\left( {x - y}\right) \mathrm{d}y \),易知 \( \varphi \left( x\right) \) 在 \( K \) 的 \( \delta \) 邻域 \( {K}_{\delta } \) 内恒等于 1 .
\end{proof}
\item 设 \(\mathcal{E}\left( {\mathbb{R}}^{n}\right)  = \left\{  {\varphi  \mid  \varphi }\right.\) 为 \({\mathbb{R}}^{n}\) 上光滑函数$\}$. 称 \(\mathcal{E}\left( {\mathbb{R}}^{n}\right)\) 中点列 \({\left\{  {\varphi }_{n}\right\}  }_{n = 1}^{\infty }\) 收敛到 0,如果对任意整数 \(m\) 和多重指标 \(\alpha  \in  {\mathbb{N}}^{n}\) 都有
\begin{equation*}
\mathop{\lim }\limits_{{n \rightarrow  \infty }}\mathop{\sup }\limits_{{\left| \alpha\right|  \leqslant  m}}\left| {\left( {{\partial }^{\alpha }{\varphi }_{n}}\right) \left( x\right) }\right|  = 0.
\end{equation*}

称 \(\mathcal{E}\left( {\mathbb{R}}^{n}\right)\) 中点列 \({\left\{  {\varphi }_{n}\right\}  }_{n = 1}^{\infty }\) 收敛到 \(\varphi\) ,如果点列 \({\left\{  {\varphi }_{n} - \varphi \right\}  }_{n = 1}^{\infty }\) 收敛到 0 .

(1)证明 \(\mathcal{E}\left( {\mathbb{R}}^{n}\right)\) 关于上述收敛性为完备的.

(2)证明自然包含 \(\iota  : \mathcal{D}\left( {\mathbb{R}}^{n}\right)  \hookrightarrow  \mathcal{E}\left( {\mathbb{R}}^{n}\right)\) 及其对偶映射 \({\iota }^{\prime } : {\mathcal{E}}^{\prime }\left( {\mathbb{R}}^{n}\right)  \hookrightarrow  {\mathcal{D}}^{\prime }\left( {\mathbb{R}}^{n}\right)\) 均为连续的.

(3)证明 \(f \in  {\mathcal{D}}^{\prime }\left( {\mathbb{R}}^{n}\right)\) 属于 \({\mathcal{E}}^{\prime }\left( {\mathbb{R}}^{n}\right)\) 当且仅当 \(f\) 具有紧支集.
\begin{proof}
    (1)设$\{\varphi_\nu(x)\}$是$\mathcal{E}(\mathbb{R})^n$中的一个基本列,那么设$\{\partial^\alpha\varphi(x)\}$一致收敛到$\varphi^\alpha(x)$,由于一致收敛的序列的导数等于极限的导数,所以$\partial^\alpha\varphi_0(x)=\varphi^\alpha(x)$,所以 \(\mathcal{E}\left( {\mathbb{R}}^{n}\right)\) 关于上述收敛性为完备的.

    (2)设$\{\varphi_\nu(x)\}$是$\mathcal{D}(\mathbb{R})^n$中的一个基本列收敛到$\varphi_0(x),\exists$ 紧集$K\subset \mathbb{R}^n,\operatorname{supp}\varphi_i\subset{ K},i=0,1,\cdots$,对于任意指标$\alpha$,有
\[
\lim\limits_{i\to \infty}\sup\limits_{x\in K}|\partial^\alpha\varphi_i(x)-\partial^\alpha\varphi_0(x)|=0\Rightarrow \lim\limits_{i\to \infty}\sup\limits_{x\in \mathbb{R}^n,|\alpha|\leqslant m}|\partial^\alpha\varphi_i(x)-\partial^\alpha\varphi_0(x)|=\lim\limits_{i\to \infty}\sup\limits_{x\in K,|\alpha|\leqslant m}|\partial^\alpha\varphi_i(x)-\partial^\alpha\varphi_0(x)|=0
    \]
    所以自然包含连续.

    设$\{\varphi_i'(x)\}$收敛到$\varphi_0'(x)$
\[
    \langle \iota^\prime \varphi'_i,\psi\rangle=\langle \varphi'_i,\iota\psi\rangle=\langle \varphi'_i,\psi\rangle\to \langle \varphi'_0,\psi\rangle=\langle \iota^\prime \varphi'_0,\psi\rangle,\forall \psi\in \mathcal{D}(\mathbb{R}^n)
    \]
    所以$\iota'$连续.

    (3)设$f$的紧支集为$K\subset{\mathbb{R}^n},K^c=\omega$,
    \[
    \langle f, \varphi\rangle=
    \]
\end{proof}
\item 证明自然包含映射 \(\mathcal{D} \hookrightarrow  \mathcal{S}\) 及其对偶 \({\mathcal{S}}^{\prime } \hookrightarrow  {\mathcal{D}}^{\prime }\) 均连续. 证明 \({\mathcal{E}}^{\prime } \subset  {\mathcal{S}}^{\prime }\) .
\begin{proof}
    设$\varphi_j\to \varphi_0(\mathcal{D})$,那么存在紧集$K$使得$\operatorname{supp}\varphi_j\subset K$,设$K\subset B_R$,下面证明\[
    \sup\limits_{x\in \mathbb{R}^n}\left|x^\alpha (\partial^\beta \varphi_j-\partial^\beta\varphi_0)\right|\leqslant    R^{|\alpha|}\sup\limits_{x\in K}\left|\partial^\beta \varphi_j-\partial^\beta\varphi_0\right|\to 0,\forall \alpha,\beta \in \mathbb{N}^n 
    \]
    所以自然包含是连续的.

    设$\iota'$为对偶映射,设$f_n\to f_0(\mathcal{S}^\prime)$,证明$\iota'f_n\to \iota'f_0$.只需证明\[\left\langle\iota'f_n,\varphi\right\rangle\to \left\langle\iota'f_0,\varphi\right\rangle,\forall \varphi \in \mathcal{D} \]
  由$f_n\to f_0(\mathcal{S}^\prime)$知
\[\left\langle f_n,\iota\varphi\right\rangle\to \left\langle f_0,\iota \varphi\right\rangle,\iota \varphi \in \mathcal{S}\]
\end{proof}
\item 定义 \({\mathcal{D}}_{{L}^{\infty }}\left( {\mathbb{R}}^{n}\right)  = \left\{  {\varphi  \mid  \varphi }\right.\) 为 \({\mathbb{R}}^{n}\) 上光滑函数且其各阶导数都属于 \(\left. {{L}^{\infty }\left( {\mathbb{R}}^{n}\right) }\right\}\) . 在 \({\mathcal{D}}_{{L}^{\infty }}\left( {\mathbb{R}}^{n}\right)\) 上引进半范族 \(\left\{  {\parallel  \cdot  {\parallel }_{\alpha ,\infty }}\right\}\) ,其中 \(\alpha  \in  {\mathbb{N}}^{n},\parallel \varphi {\parallel }_{\alpha ,\infty } = {\begin{Vmatrix}{\partial }^{\alpha }\varphi \end{Vmatrix}}_{\infty }\) .

(1)证明 \({\mathcal{D}}_{{L}^{\infty }}\left( {\mathbb{R}}^{n}\right)\) 关于上述半范族完备.

(2) 证明 \({C}_{0}^{\infty }\left( {\mathbb{R}}^{n}\right)\) 在 \({\mathcal{D}}_{{L}^{\infty }}\left( {\mathbb{R}}^{n}\right)\) 中不闭. 试求其闭包.
\begin{proof}
    
\end{proof}
\item 设
\begin{equation*}
{f}_{j}\left( x\right)  = \left\{  \begin{array}{ll} 0, & \left| x\right|  \geq  \frac{1}{j}, \\  \frac{j}{2}, & \left| x\right|  < \frac{1}{j}, \end{array}\right.
\end{equation*}

证明函数列 \({\left\{  {f}_{j}\left( x\right) \right\}  }_{j = 1}^{\infty }\) 在 \({\mathcal{D}}^{\prime }\left( \mathbb{R}\right)\) 中弱收敛到 \(\delta \left( x\right)\) .

\begin{proof}
    \[
    \langle f_j(x),\varphi(x)\rangle=\int_\mathbb{R}f_j(x)\varphi(x)\dif x=\varphi(\xi)\to \varphi(0)=\langle\delta,\varphi(x)\rangle,\forall \varphi(x)\in \mathcal{D}(\mathbb{R})
    \]
    
\end{proof}
\item 设

\begin{equation*}
{f}_{j}\left( x\right)  = {\left( 1 + \frac{x}{j}\right) }^{j},\;j = 1,2,\cdots ,
\end{equation*}

证明 \({f}_{j}\left( x\right)  \rightarrow  {e}^{x}\left( {{\mathcal{D}}^{\prime }\left( \mathbb{R}\right) }\right)\) .
\begin{proof}      \begin{equation*}
    \begin{aligned}
        \langle f_j(x),\varphi(x)\rangle=\int_\mathbb{R}f_j(x)\varphi(x)\dif x=\int_K\left(\left(1+\frac{x}{j}\right)^j-e^x\right)\varphi(x)\dif x+\langle e^x,\varphi(x)\rangle\\=\int_K\left(e^{x+\xi}-e^x\right)\varphi(x)\dif x+\langle e^x,\varphi(x)\rangle\to\langle e^x,\varphi(x)\rangle,\forall \varphi(x)\in \mathcal{D}(\mathbb{R})
    \end{aligned}
\end{equation*}
    
\end{proof}
\item 设 \(\Omega  \subset  {\mathbb{R}}^{n}\) 为开集,证明 \({C}^{\infty }\left( \Omega \right)\) 为 \({\mathcal{D}}^{\prime }\left( \Omega \right)\) 中弱稠密集,即对任意 \(u \in  {\mathcal{D}}^{\prime }\left( \Omega \right)\) , 有 \({C}^{\infty }\left( \Omega \right)\) 中函数列 \(\left\{  {{u}_{j}\left( x\right) }\right\}\) 弱收敛到 \(u\) .

\item 在 \({\mathcal{D}}^{\prime }\) 中,证明:
\begin{equation*}
\frac{1}{2\pi }\frac{1 - {r}^{2}}{1 - {2r}\cos \theta  + {r}^{2}} \rightharpoonup  \delta \left( \theta \right) ,\;r \rightarrow  1,
\end{equation*}
\begin{equation*}
\frac{1}{{2a}\sqrt{\pi t}}{\exp }^{-\frac{{x}^{2}}{4{a}^{2}t}} \rightharpoonup  \delta \left( x\right) ,\;t \rightarrow  {0}^{ + },
\end{equation*}
\begin{equation*}
\frac{1}{\pi }\frac{\varepsilon }{{x}^{2} + {\varepsilon }^{2}} \rightharpoonup  \delta \left( x\right) ,\;\varepsilon  \rightarrow  {0}^{ + }.
\end{equation*}
\begin{proof}
    (1)
    \begin{equation*}
        \begin{aligned}
          & \left| \langle \frac{1}{2\pi }\frac{1 - {r}^{2}}{1 - {2r}\cos \theta  + {r}^{2}},\varphi(x)\rangle-\varphi(0)\right|=\left|\int_{-\pi}^\pi \frac{1}{2\pi }\frac{1 - {r}^{2}}{1 - {2r}\cos \theta  + {r}^{2}}(\varphi(x)-\varphi(0))\dif x \right|\\
           \leqslant& \left|\int_{-\xi}^\xi \frac{1}{2\pi }\frac{1 - {r}^{2}}{1 - {2r}\cos \theta  + {r}^{2}}(\varphi(x)-\varphi(0))\dif x \right|+ \left|\int_{|x|>\xi} \frac{1}{2\pi }\frac{1 - {r}^{2}}{1 - {2r}\cos \theta  + {r}^{2}}(\varphi(x)-\varphi(0))\dif x \right|
           \\           \leqslant &\left|\int_{-\xi}^\xi \frac{1}{2\pi }\frac{1 - {r}^{2}}{1 - {2r}\cos \theta  + {r}^{2}}\varepsilon\dif x \right|+ 2M\times\frac{\varepsilon}{M}<2\varepsilon
        \end{aligned}
    \end{equation*}

    (2)
$$
G_t(x)=\frac{1}{2 a \sqrt{\pi t}} \exp \left(-\frac{x^2}{4 a^2 t}\right), \quad t>0,
$$


要证 $G_t \rightharpoonup \delta$ 当 $t \rightarrow 0^{+}$。

对任意 $\varphi \in C_c^{\infty}(\mathbb{R})$ ，考虑
$$
I_t:=\int_{\mathbb{R}} G_t(x) \varphi(x) d x
$$


作变换 $y=\frac{x}{2 a \sqrt{t}}$ ，则 $x=2 a \sqrt{t} y$ ，且
$$
I_t=\int_{\mathbb{R}} \frac{1}{\sqrt{\pi}} e^{-y^2} \varphi(2 a \sqrt{t} y) d y
$$


注意对每个固定 $y$ ，当 $t \rightarrow 0$ 有 $\varphi(2 a \sqrt{t} y) \rightarrow \varphi(0)$ 。并且由于 $\varphi$ 连续且有紧支撑，存在常数 $M$ 使得 $|\varphi(2 a \sqrt{t} y)| \leq M$ 对所有 $t, y$ 成立；而 $(1 / \sqrt{\pi}) e^{-y^2}$ 可积。因此可用主导收敛定理将极限交换：
$$
\lim _{t \rightarrow 0^{+}} I_t=\int_{\mathbb{R}} \frac{1}{\sqrt{\pi}} e^{-y^2} \varphi(0) d y=\varphi(0) \cdot \int_{\mathbb{R}} \frac{1}{\sqrt{\pi}} e^{-y^2} d y=\varphi(0)
$$


于是 $G_t \rightharpoonup \delta$ 。

    (3)    \begin{equation*}
        \begin{aligned}
          & \left| \langle \frac{1}{\pi }\frac{\varepsilon }{{x}^{2} + {\varepsilon }^{2}},\varphi(x)\rangle-\varphi(0)\right|=\left|\int_{-\infty}^\infty \frac{1}{\pi }\frac{\varepsilon }{{x}^{2} + {\varepsilon }^{2}} (\varphi(x)-\varphi(0))\dif x \right|\\
           \leqslant& \left|\int_{-\xi}^\xi\frac{1}{\pi }\frac{\varepsilon }{{x}^{2} + {\varepsilon }^{2}} (\varphi(x)-\varphi(0))\dif x \right|+ \left|\int_{|x|>\xi} \frac{1}{\pi }\frac{\varepsilon }{{x}^{2} + {\varepsilon }^{2}} (\varphi(x)-\varphi(0))\dif x \right|
           \\           \leqslant &\left|\int_{-\xi}^\xi \frac{1}{\pi }\frac{\varepsilon }{{x}^{2} + {\varepsilon }^{2}} \zeta\dif x \right|+ 2M\times\frac{\zeta}{M}<2\zeta
        \end{aligned}
    \end{equation*}
\end{proof}
\item 设 \(\Omega  \subset  {\mathbb{R}}^{n}\) 为开集. 若 \(f,g \in  {\mathcal{D}}^{\prime }\left( \Omega \right)\) 满足对任意 \(x \in  \Omega\) ,存在 \(x\) 的邻域 \(U\) 使得对任意 \(\varphi  \in  \mathcal{D}\left( U\right)\) 都有 \(f\left( \varphi \right)  = g\left( \varphi \right)\) ,则 \(f = g\) . 特别地,若广义函数 \(f \in  {\mathcal{D}}^{\prime }\left( \Omega \right)\) 满足条件: 对任意 \(x \in  \Omega\) ,存在 \(x\) 的邻域 \(U\) 使得在 \(U\) 上成立 \(f = 0\) ,则在 \(\Omega\) 上有 \(f = 0\) .
\begin{proof}
    设紧集$K=\operatorname{supp}\varphi$,存在有限覆盖$K=\bigcup\limits_{i=1}^nU_{x_i}$,由单位分解定理知存在函数族$\{\psi_i\}_{i\in I}\subset C^\infty(\Omega),\operatorname{supp}\psi_i\subset U_{x_i},0\leqslant \psi_i\leqslant 1,\sum\limits_{i\in I}\psi_i=1$,那么
    \[f(\varphi)=f\left(\sum\limits_{i\in I}\varphi\psi_i\right)=\sum\limits_{i\in  I}f(\varphi\psi_i)=\sum\limits_{i\in  I}g(\varphi\psi_i)=g(\varphi)\]
    
    
\end{proof}
\item 证明若 \(\ell  \in  {\mathcal{D}}^{\prime }\left( {\mathbb{R}}^{n}\right)\) 具有紧支集,则存在整数 \(M\) 及连续函数 \({g}_{\alpha },\alpha  \in  {\mathbb{N}}^{n},\left| \alpha \right|  \leqslant  M\) 使得

\begin{equation*}
\ell  = \mathop{\sum }\limits_{{\left| \alpha \right|  \leqslant  M}}{\partial }^{\alpha }{g}_{\alpha }
\end{equation*}

\item 设 \(\Omega  \subset  {\mathbb{R}}^{n}\) 为开集. 称 \(f \in  {\mathcal{D}}^{\prime }\left( \Omega \right)\) 为正广义函数,如果对任意满足 \(\varphi  \geq  0\) 的 \(\varphi  \in\)  \(\mathcal{D}\left( \Omega \right)\) 都有 \(f\left( \varphi \right)  \geq  0\) . 证明正广义函数为一测度.



\end{problemset}

\section{广义函数的运算}

根据广义函数的定义, 我们可以将经典函数的一些运算法则推广到广义函数空间 \({\mathcal{D}}^{\prime }\left( \Omega \right)\) .

 \begin{definition}
  对任意两个广义函数 \(f,g \in  {\mathcal{D}}^{\prime }\left( \Omega \right)\) ,定义 \(f + g \in  {\mathcal{D}}^{\prime }\left( \Omega \right)\) 如下  
\begin{equation*}
  \langle f + g,\varphi \rangle  = \langle f,\varphi \rangle  + \langle g,\varphi \rangle ,\;\forall \varphi  \in  \mathcal{D}\left( \Omega \right) .
  \end{equation*}
  
  对任意广义函数 \(f \in  {\mathcal{D}}^{\prime }\left( \Omega \right)\) 及任意无穷次可微函数 \(a\left( x\right)  \in  \mathcal{E} = {C}^{\infty }\left( \Omega \right)\) ,定义乘积 \({af} \in  {\mathcal{D}}^{\prime }\left( \Omega \right)\) 如下  \begin{equation*}
  \langle {af},\varphi \rangle  = \langle f,{a\varphi }\rangle ,\;\forall \varphi  \in  {\mathcal{D}}^{\prime }\left( \Omega \right) .
  \end{equation*}
\end{definition}
\begin{note}
不难验证上述运算的合理性. 特别,当 \(f,g\) 为经典的局部Lebesgue可积函数时, 上述运算可归为经典函数的相应运算.
\end{note}
 \begin{example}
  对任意光滑函数 \(a\left( x\right)  \in  {C}^{\infty }\left( {\mathbb{R}}^{n}\right) ,a\left( x\right) \delta \left( x\right)  = a\left( 0\right) \delta \left( x\right)\) .
  
  L. Schwartz [Sch] 曾证明在 \({\mathcal{D}}^{\prime }\left( \Omega \right)\) 上不存在具有结合性的乘法运算. 例如, \(\delta  \cdot  \delta\) 一般是没有意义的(近年来,物理学界和数学界都以多种方式定义这类乘积). 实际上,容易看出, \({\mathcal{D}}^{\prime }\left( \Omega \right)\) 上具结合性的乘法运算会导致矛盾.
  
\end{example}
\begin{example}
  定义广义函数 \({PV}\left( \frac{1}{x}\right)  \in  {\mathcal{D}}^{\prime }\left( \mathbb{R}\right)\) 如下:  
\begin{equation*}
  \left\langle  {{PV}\left( \frac{1}{x}\right) ,\varphi }\right\rangle   \mathrel{\text{ := }} \mathop{\lim }\limits_{{\varepsilon  \rightarrow  {0}^{ + }}}\left( {{\int }_{-\infty }^{-\varepsilon } + {\int }_{\varepsilon }^{\infty }}\right) \frac{\varphi \left( x\right) }{x}{\dif x},\;\forall \varphi  \in  \mathcal{D}\left( \mathbb{R}\right) .
  \end{equation*}
  
  即 \({PV}\left( \frac{1}{x}\right)\) 是由右边极限定义的非局部可积函数 \(\frac{1}{x}\) 的 Cauchy 主值. 容易看出,
\begin{equation*}
  \left\lbrack  {{PV}\left( \frac{1}{x}\right)  \cdot  x}\right\rbrack   \cdot  \delta  = 1 \cdot  \delta  = \delta ,
  \end{equation*}
  \begin{equation*}
  {PV}\left( \frac{1}{x}\right)  \cdot  \left( {x \cdot  \delta }\right)  = {PV}\left( \frac{1}{x}\right)  \cdot  0 = 0.
  \end{equation*}
  
\end{example}\begin{definition}
   对任意 \({x}_{0} \in  {\mathbb{R}}^{n}\) ,定义平移算子 \({\tau }_{{x}_{0}} : \mathcal{D}\left( {\mathbb{R}}^{n}\right)  \rightarrow  \mathcal{D}\left( {\mathbb{R}}^{n}\right)\) 为\begin{equation*}
  \left( {{\tau }_{{x}_{0}}\varphi }\right) \left( x\right)  \mathrel{\text{ := }} \varphi \left( {x - {x}_{0}}\right) ,\;\forall \varphi  \in  \mathcal{D}\left( {\mathbb{R}}^{n}\right) .
  \end{equation*}
  
  对任意 \({x}_{0} \in  {\mathbb{R}}^{n}\) ,定义 \({\widetilde{\tau }}_{{x}_{0}} \mathrel{\text{ := }} {\left( {\tau }_{-{x}_{0}}\right) }^{ * }\) . 即对任意 \(f \in  {\mathcal{D}}^{\prime }\left( {\mathbb{R}}^{n}\right)\) ,有
\begin{equation*}
  \left\langle  {{\widetilde{\tau }}_{{x}_{0}}f,\varphi }\right\rangle   = \left\langle  {f,{\tau }_{-{x}_{0}}\varphi }\right\rangle   = \left\langle  {f,\varphi \left( {x + {x}_{0}}\right) }\right\rangle  ,\;\forall \varphi  \in  \mathcal{D}\left( {\mathbb{R}}^{n}\right) .
  \end{equation*}
  
\end{definition}
\begin{note}
   \({\tau }_{{x}_{0}} : \mathcal{D}\left( {\mathbb{R}}^{n}\right)  \rightarrow  \mathcal{D}\left( {\mathbb{R}}^{n}\right)\) 连续且 \({\widetilde{\tau }}_{{x}_{0}}\) 是平移算子的推广. 事实上,若 \(f\left( x\right)  \in\)  \({L}_{loc}^{1}\left( {\mathbb{R}}^{n}\right)\) ,则对任意 \({x}_{0} \in  {\mathbb{R}}^{n}\) ,有
\begin{equation*}
\left\langle  {{\tau }_{{x}_{0}}f,\varphi }\right\rangle   = {\int }_{{\mathbb{R}}^{n}}f\left( {x - {x}_{0}}\right) \varphi \left( x\right) {\dif x} = {\int }_{{\mathbb{R}}^{n}}f\left( x\right) \varphi \left( {x + {x}_{0}}\right) {\dif x}
= \left\langle  {f,{\tau }_{-{x}_{0}}\varphi }\right\rangle   = \left\langle  {{\widetilde{\tau }}_{{x}_{0}}f,\varphi }\right\rangle  ,\;\forall \varphi  \in  \mathcal{D}\left( {\mathbb{R}}^{n}\right) ,
\end{equation*}

即得 \({\widetilde{\tau }}_{{x}_{0}}f = {\tau }_{{x}_{0}}f = f\left( {x - {x}_{0}}\right)\) .
\end{note}
\begin{definition}
  定义反射算子 \(\sigma  : \mathcal{D}\left( {\mathbb{R}}^{n}\right)  \rightarrow  \mathcal{D}\left( {\mathbb{R}}^{n}\right)\) 如下  \begin{equation*}
  \left( {\sigma \varphi }\right) \left( x\right)  \mathrel{\text{ := }} \varphi \left( {-x}\right) ,\;\forall \varphi  \in  \mathcal{D}\left( {\mathbb{R}}^{n}\right) .
  \end{equation*}
  
  定义 \(\widetilde{\sigma } \mathrel{\text{ := }} {\sigma }^{ * }\) . 即对任意 \(f \in  {\mathcal{D}}^{\prime }\left( {\mathbb{R}}^{n}\right)\) ,有
\begin{equation*}
  \langle \widetilde{\sigma }f,\varphi \rangle  = \left\langle  {{\sigma }^{ * }f,\varphi }\right\rangle   = \langle f,{\sigma \varphi }\rangle ,\;\forall \varphi  \in  \mathcal{D}\left( {\mathbb{R}}^{n}\right) .
  \end{equation*}
  
\end{definition}
\begin{note}
\(\sigma  : \mathcal{D}\left( {\mathbb{R}}^{n}\right)  \rightarrow  \mathcal{D}\left( {\mathbb{R}}^{n}\right)\) 连续且 \(\widetilde{\sigma }\) 是反射算子的推广. 事实上,若 \(f\left( x\right)  \in\)  \({L}_{loc}^{1}\left( {\mathbb{R}}^{n}\right)\) ,则有
\begin{equation*}
\langle {\sigma f},\varphi \rangle  = {\int }_{{\mathbb{R}}^{n}}f\left( {-x}\right) \varphi \left( x\right) {\dif x} = {\int }_{{\mathbb{R}}^{n}}f\left( x\right) \varphi \left( {-x}\right) {\dif x}
= \langle f,{\sigma \varphi }\rangle  = \langle \widetilde{\sigma }f,\varphi \rangle ,\;\forall \varphi  \in  \mathcal{D}\left( {\mathbb{R}}^{n}\right) ,
\end{equation*}

即得 \(\widetilde{\sigma }f = {\sigma f} = f\left( {-x}\right)\) .
\end{note} 
在 \( {\mathcal{D}}^{\prime }\left( \Omega \right) \) 上定义算子 \( {A}^{ * } \) 如下:
\[\left\langle  {{A}^{ * }f,\varphi }\right\rangle   = \langle f,{A\varphi }\rangle \;\left( {\forall \varphi  \in  \mathcal{D}\left( \Omega \right) ,\forall f \in  {\mathcal{D}}^{\prime }\left( \Omega \right) }\right) . \]

显然, \( {A}^{ * } : {\mathcal{D}}^{\prime }\left( \Omega \right)  \rightarrow  {\mathcal{D}}^{\prime }\left( \Omega \right) \) 并且是连续的. 事实上,如果 \( {f}_{j} \rightarrow  f \)
\( \left( {{\mathcal{D}}^{\prime }\left( \Omega \right) }\right) ,j \rightarrow  \infty \) ,那么对 \( \forall \varphi  \in  \mathcal{D}\left( \Omega \right) \) ,有
\( \left\langle  {{A}^{ * }{f}_{j},\varphi }\right\rangle   = \left\langle  {{f}_{j},{A\varphi }}\right\rangle   \rightarrow  \langle f,{A\varphi }\rangle  = \left\langle  {{A}^{ * }f,\varphi }\right\rangle  \;\left( {j \rightarrow  \infty }\right) , \)
即得 \( {A}^{ * }{f}_{j} \rightarrow  {A}^{ * }f,j \rightarrow  \infty \) .
按照这种方式我们来定义广义函数的各种运算.
\begin{definition}
  定义广义函数 \(f \in  {\mathcal{D}}^{\prime }\left( \Omega \right)\) 的导数 (广义导数或弱导数) \(\widetilde{\partial } ^{\alpha }f\) 为由下式确定的广义函数
\begin{equation*}
  \left\langle  {\widetilde{\partial }^{\alpha }f,\varphi }\right\rangle   = {\left( -1\right) }^{\left| \alpha \right| }\left\langle  {f,{\partial }^{\alpha }\varphi }\right\rangle  ,\;\forall \varphi  \in  \mathcal{D},
  \end{equation*}
  
  其中 \(\alpha  \in  {\mathbb{N}}^{n}\) 为 \(n\) 重指标, \(\left| \alpha \right|  = {\alpha }_{1} + \cdots  + {\alpha }_{n}\) .
  
\end{definition}
\begin{note}
     若 \( f\left( x\right)  \in  {C}^{1}\left( \Omega \right) \) ,依自然对应:
\( \langle f,\varphi \rangle  = {\int }_{\Omega }f\left( x\right) \varphi \left( x\right) \mathrm{d}x\;\left( {\forall \varphi  \in  \mathcal{D}\left( \Omega \right) }\right) \)
产生的广义函数仍然记作 \( f,f \) 有广义微商 \( {\widetilde{\partial }}_{{x}_{i}}f \) (即 \( {\widetilde{\partial }}^{\alpha }f,\alpha  = \)
\( \left( {\underbrace{0,\cdots ,0,1},0,\cdots ,0}\right) ) \) ,同时 \( f\left( x\right) \) 有普通的微商 \( {\partial }_{{x}_{i}}f \) ,我们说 \( {\widetilde{\partial }}_{{x}_{i}}f \)
正是 \( {\partial }_{{x}_{i}}f \) 对应的广义函数. 这是因为
\( \left\langle  {{\widetilde{\partial }}_{{x}_{i}}f,\varphi }\right\rangle   =  - \left\langle  {f,{\partial }_{{x}_{i}}\varphi }\right\rangle   =  - {\int }_{\Omega }f\left( x\right) {\partial }_{{x}_{i}}\varphi \left( x\right) \mathrm{d}x \)
\( = {\int }_{\Omega }{\partial }_{{x}_{i}}f\left( x\right) \varphi \left( x\right) \mathrm{d}x\;\left( {\forall \varphi  \in  \mathcal{D}\left( \Omega \right) }\right) . \)

广义函数对微商运算是封闭的,即任意广义函数 \( f \in \)
\( {\mathcal{D}}^{\prime }\left( \Omega \right) \) 都可以做任意次广义微商. 当然,即使是局部可积函数,也
未必能做普通微商, 但当把它看作广义函数时, 总可以做广义微
商, 做广义微商后所得的是广义函数, 未必是普通函数.

若 \( \alpha ,\beta \) 是任意两个多重指标,则
\( {\widetilde{\partial }}^{\alpha } \cdot  {\widetilde{\partial }}^{\beta } = {\widetilde{\partial }}^{\alpha  + \beta } = {\widetilde{\partial }}^{\beta } \cdot  {\widetilde{\partial }}^{\alpha }. \)

由 \( {\widetilde{\partial }}^{\alpha } \) 的连续性,
\( {f}_{j} \rightarrow  {f}_{0}\;\left( {j \rightarrow  \infty }\right)  \Rightarrow  {\widetilde{\partial }}^{\alpha }{f}_{j} \rightarrow  {\widetilde{\partial }}^{\alpha }{f}_{0}\;\left( {j \rightarrow  \infty }\right) , \)
可见广义微商与极限总是可交换的.

在后文中我们不再区分$\widetilde{\partial}$和$\partial$.
\end{note}


对无穷可微函数 \(a\left( x\right)\) , Leibniz公式成立:
\begin{equation*}
{\widetilde{\partial } }^{\alpha }\left( {af}\right)  = \mathop{\sum }\limits_{{\left| \beta \right|  = 0}}^{\left| \alpha \right| }\left( \begin{array}{l} \alpha \\  \beta  \end{array}\right) {\widetilde{\partial } }^{\beta }a{\widetilde{\partial } }^{\alpha  - \beta }f,\;\forall f \in  {\mathcal{D}}^{\prime }\left( \Omega \right) .
\end{equation*}
\begin{example}
  计算Heaviside函数
\begin{equation*}
  H\left( {x - a}\right)  = \left\{  {\begin{array}{ll} 1, & x \geq  a \\  0, & x < a, \end{array}\;x \in  {\mathbb{R}}^{1}}\right.
  \end{equation*}
  
  的导数.
\end{example}
\begin{solution}根据定义有
\begin{equation*}
\left\langle  {{H}^{\prime },\varphi }\right\rangle   =  - \left\langle  {H,{\varphi }^{\prime }}\right\rangle   =  - {\int }_{a}^{\infty }{\varphi }^{\prime }\left( x\right) {\dif x} = \varphi \left( a\right)  = \langle \delta \left( {x - a}\right) ,\varphi \rangle ,\;\forall \varphi  \in  \mathcal{D}\left( \Omega \right) .
\end{equation*}

因此,
\begin{equation*}
{H}^{\prime }\left( {x - a}\right)  = \delta \left( {x - a}\right) .
\end{equation*}

特别地,当 \(a = 0\) 时有 \({H}^{\prime }\left( x\right)  = \delta \left( x\right)\) .

若引进 \(n\) 维Heaviside函数
\begin{equation*}
{H}_{n}\left( x\right)  = \left\{  \begin{array}{ll} 1, & {x}_{j} \geq  0,j = 1,\cdots ,n, \\  0, & \text{ 其它情形 } \end{array}\right.
\end{equation*}

则类似地可得
\begin{equation*}
\langle \delta ,\varphi \rangle  = {\left( -1\right) }^{n}\left\langle  {{H}_{n},\frac{{\partial }^{n}}{\partial {x}_{1}\cdots \partial {x}_{n}}\varphi }\right\rangle  .
\end{equation*}

故 \(n\) 维 \(\delta\) 分布为
\begin{equation*}
\delta \left( x\right)  = \frac{{\partial }^{n}}{\partial {x}_{1}\cdots \partial {x}_{n}}{H}_{n}\left( x\right) ,\;x \in  {\mathbb{R}}^{n}.
\end{equation*}\end{solution}
\begin{note}
从上例可以看出, 函数型广义函数的弱导数不一定还是函数型广义函数.
\end{note}
\begin{example}
   计算 \(\delta \left( {x - {x}^{0}}\right)\) 的导数.
\end{example}
\begin{solution}
   根据定义有
\begin{equation*}
\left\langle  {{\delta }^{\prime }\left( {x - {x}^{0}}\right) ,\varphi }\right\rangle   =  - \left\langle  {\delta \left( {x - {x}^{0}}\right) ,{\varphi }^{\prime }}\right\rangle   =  - {\varphi }^{\prime }\left( {x}^{0}\right) ,\;x \in  {\mathbb{R}}^{1}.
\end{equation*}

一般地可得
\begin{equation*}
\left\langle  {{\delta }^{\left( k\right) }\left( {x - {x}^{0}}\right) ,\varphi }\right\rangle   = {\left( -1\right) }^{k}{\varphi }^{\left( k\right) }\left( {x}^{0}\right) ,\;x \in  {\mathbb{R}}^{1}.
\end{equation*}
\end{solution}
\begin{example}
  设 \(\Delta  = \frac{{\partial }^{2}}{\partial {x}_{1}^{2}} + \cdots  + \frac{{\partial }^{2}}{\partial {x}_{n}^{2}}\) ,则在 \({\mathcal{D}}^{\prime }\left( {\mathbb{R}}^{n}\right)\) 中有  \begin{equation*}
  \Delta {\left| x\right| }^{2 - n} = \left( {2 - n}\right) {\omega }_{n}\delta \left( x\right) ,\;n \geq  3
  \end{equation*}
\begin{equation*}
  \Delta \ln \left| x\right|  = {2\pi \delta }\left( x\right) ,\;n = 2,
  \end{equation*}
  
  其中 \({\omega }_{n}\) 是 \({\mathbb{R}}^{n}\) 中单位球面的面积.  
\end{example}
\begin{proof}
  (1) 当 \(n \geq  3\) 时, \(\forall \varphi  \in  \mathcal{D}\left( {\mathbb{R}}^{n}\right)\) ,据定义有  
\begin{equation*}
\begin{aligned}
      \left\langle  {\Delta {\left| x\right| }^{2 - n},\varphi }\right\rangle  & = \left\langle  {{\left| x\right| }^{2 - n},{\Delta \varphi }}\right\rangle   = \mathop{\lim }\limits_{{\varepsilon  \rightarrow  0}}{\int }_{\left| x\right|  \geq  \varepsilon }{\left| x\right| }^{2 - n}{\Delta \varphi }\left( x\right) {\dif x}\\  &= \mathop{\lim }\limits_{{\varepsilon  \rightarrow  0}}\left\{  {{\int }_{\left| x\right|  \geq  \varepsilon }\Delta {\left| x\right| }^{2 - n}\varphi \left( x\right) {\dif x}}\right.
  + {\int }_{\left| x\right|  = \varepsilon }\left\lbrack  {\varphi \frac{\partial {\left| x\right| }^{2 - n}}{\partial r} - {r}^{2 - n}\frac{\partial \varphi }{\partial r}}\right\rbrack  {\dif\sigma },
\end{aligned}
  \end{equation*}
  
  其中 \({\dif\sigma }\) 是球面 \(\left\{  {x \in  {\mathbb{R}}^{n}\left| \right| x \mid   = \varepsilon }\right\}\) 上的面积元, \(r = \left| x\right|\) ,最后一个等号利用了Green公式,实际上可以取一充分大的球 \(B\left( {0,R}\right)  = \left\{  {x \in  {\mathbb{R}}^{n}\left| \right| x \mid   < R}\right\}\) 使得 \(\operatorname{supp}\varphi  \subset\)  \(B\left( {0,R}\right)\) ,然后再应用Green公式. 另外,注意到
  \begin{equation*}
  \Delta {\left| x\right| }^{2 - n} = 0\text{ ,当 }\left| x\right|  \neq  0,
  \end{equation*}\begin{equation*}
  {\varepsilon }^{2 - n}{\int }_{\left| x\right|  = \varepsilon }\frac{\partial \varphi }{\partial r}{\dif\sigma } = O\left( \varepsilon \right) \text{ ,当 }\varepsilon  \rightarrow  0,
  \end{equation*}
  
  而且
\begin{equation*}
  {\varepsilon }^{1 - n}{\int }_{\left| x\right|  = \varepsilon }\varphi \left( x\right) {\dif\sigma } \rightarrow  \varphi \left( 0\right) {\omega }_{n},\text{ 当 }\varepsilon  \rightarrow  0,
  \end{equation*}
  
  综上可得
  \begin{equation*}
  \left\langle  {\Delta {\left| x\right| }^{2 - n},\varphi }\right\rangle   = \left( {2 - n}\right) \varphi \left( 0\right) {\omega }_{n} = \left( {2 - n}\right) {\omega }_{n}\langle \delta ,\varphi \rangle ,\;\forall \varphi  \in  \mathcal{D}\left( {\mathbb{R}}^{n}\right) .
  \end{equation*}
  
  (2) 同理可证得
  \begin{equation*}
  \langle \Delta \ln \left| x\right| ,\varphi \rangle  = {2\pi }\langle \delta ,\varphi \rangle ,\;\forall \varphi  \in  \mathcal{D}\left( {\mathbb{R}}^{n}\right) .
  \end{equation*}
  
\end{proof}
\begin{example}
  验证下列公式
\begin{equation*}
  {x}^{n}{\partial }^{m}\delta \left( x\right)  = \left\{  \begin{array}{ll} {\left( -1\right) }^{n}\frac{m!}{\left( {m - n}\right) !}{\delta }^{\left( m - n\right) }\left( x\right) , & m \geq  n \\  0, & m < n \end{array}\right.
  \end{equation*}
\end{example}



\begin{proof}
  \begin{equation*}
  \begin{aligned}
        \left\langle  {{x}^{n}{\partial }^{m}\delta \left( x\right) ,\varphi \left( x\right) }\right\rangle &  = {\left( -1\right) }^{m}\left\langle  {\delta \left( x\right) ,{\partial }^{m}\left( {{x}^{n}\varphi \left( x\right) }\right) }\right\rangle  = {\left( -1\right) }^{m}\mathop{\sum }\limits_{{r = 0}}^{m}\left( \begin{matrix} m \\  r \end{matrix}\right) \left( {{\partial }^{r}{x}^{n}}\right) \left( {{\partial }^{m - r}\varphi \left( x\right) }\right) {|}_{x = 0}  \\&= \left\{  \begin{array}{ll} \frac{{\left( -1\right) }^{m}m!}{n!\left( {m - n}\right) !}n!\left( {{\partial }^{m - n}\varphi }\right) \left( 0\right) , & m \geq  n \\  0, & m < n \end{array}\right.  = \left\{  \begin{array}{ll} \frac{{\left( -1\right) }^{n}m!}{\left( {m - n}\right) !}\left\langle  {{\delta }^{\left( m - n\right) },\varphi }\right\rangle  , & m \geq  n \\  0, & m < n. \end{array}\right.
  \end{aligned}
  \end{equation*}


\end{proof}

引进了广义函数的上述运算后,就可以理解微分方程在广义函数空间 \({\mathcal{D}}^{\prime }\left( \Omega \right)\) 中的解(弱解)的含义. 例如, 对于线性偏微分算子
\begin{equation*}
P\left( {x,\partial }\right)  \equiv  \mathop{\sum }\limits_{{\left| \alpha \right|  \leqslant  m}}{a}_{\alpha }\left( x\right) {\partial }^{\alpha }
\equiv  \mathop{\sum }\limits_{{\left| \alpha \right|  \leqslant  m}}{a}_{{\alpha }_{1}\cdots {\alpha }_{n}}\left( x\right) \frac{{\partial }^{\left| \alpha \right| }}{{\partial }^{{\alpha }_{1}}{x}_{1}\cdots {\partial }^{{\alpha }_{n}}{x}_{n}},x \in  {\mathbb{R}}^{n},
\end{equation*}
\begin{definition}
  若对给定的广义函数 \(f \in  {\mathcal{D}}^{\prime }\left( \Omega \right)\) ,存在广义函数 \(u \in  {\mathcal{D}}^{\prime }\left( \omega \right) ,\omega  \subset  \Omega\) ,使得  
  \begin{equation*}
  \left\langle  {u,{P}^{ * }\varphi }\right\rangle   = \langle f,\varphi \rangle ,\;\forall \varphi  \in  \mathcal{D}\left( \Omega \right) ,
  \end{equation*}
  
  其中 \({P}^{ * }\) 称为 \(P\) 的伴随算子,则称 \(u\) 为方程 \({Pu} = f\) 在 \({\mathcal{D}}^{\prime }\left( \Omega \right)\) 中的弱解或解.
  
\end{definition}
\subsection{习 题}

\begin{problemset}
  \item 证明广义函数的混合导数与微分顺序无关. 即对任意 \(f \in  \mathcal{D}\left( \Omega \right) ,\alpha ,\beta  \in  {\mathbb{N}}^{n}\) ,有
\begin{equation*}
  {\partial }^{\alpha }{\partial }^{\beta }f = {\partial }^{\beta }{\partial }^{\alpha }f.
  \end{equation*}
  \item 证明
\begin{equation*}
  \delta \left( {{x}_{1},{x}_{2},\cdots ,{x}_{n}}\right)  = \delta \left( {x}_{1}\right)  \circ  \delta \left( {x}_{2}\right)  \circ  \cdots  \circ  \delta \left( {x}_{n}\right),
  \end{equation*}
  其中符号 \(\circ\) 表示算子的复合.

  \item 设
  \begin{equation*}
  f\left( x\right)  = \left\{  
  \begin{array}{ll} 
  {x}^{2}, & x \leqslant  1, \\  
  x, & x > 1.
  \end{array}\right.
  \end{equation*}
  试求 \(\frac{\dif f}{\dif x},\frac{{\dif}^{2}f}{\dif {x}^{2}},\frac{{\dif}^{3}f}{\dif{x}^{3}}\) .
\begin{solution}
    \begin{equation*}
        \begin{aligned}
          \langle  \frac{\dif f}{\dif x},\varphi(x)\rangle&=-\langle f, \frac{\dif \varphi}{\dif x}\rangle=-\int_{-\infty}^\infty f(x)\frac{\dif \varphi}{\dif x}\dif{x}=-\int_{-\infty}^1 x^2\frac{\dif \varphi}{\dif x}\dif{x}-\int_{1}^\infty x\frac{\dif \varphi}{\dif x}\dif{x}\\
          &= -\varphi(1)+\int_{-\infty}^1 2x\varphi(x)\dif{x}+\varphi(1)+\int_{1}^\infty \varphi(x)\dif x=\langle H(x-1)+2xH(1-x),\varphi(x)\rangle
          \\
        \frac{{\dif}^{2}f}{\dif {x}^{2}}  &=\delta(x-1)+2H(1-x)-2x\delta(x-1) =2H(1-x)-\delta(x-1)         \\
        \frac{{\dif}^{3}f}{\dif {x}^{3}}  &=-\delta'(x-1)-2\delta(x-1)
        \end{aligned}
    \end{equation*}
\end{solution}
  \item 验证公式\begin{equation*}
  \Delta \frac{1}{r} =  - {4\pi \delta }\left( x\right),\;{r}^{2} = {x}_{1}^{2} + {x}_{2}^{2} + {x}_{3}^{2}.
  \end{equation*}
\begin{proof}
    \begin{equation*}
\begin{aligned}
\Delta\frac{1}{r}=-4\pi \delta(x)
\end{aligned}
  \end{equation*}
  
\end{proof}
  \item 证明
  \begin{equation*}
  {\delta }^{\left( m\right) }\left( {{ax} + b}\right)  = \frac{1}{{a}^{m}\left| a\right| }{\delta }^{\left( m\right) }\left( {x + \frac{b}{a}}\right),\;a \neq  0.
  \end{equation*}
\begin{proof}
    \[
    \left\langle   {\delta }^{\left( m\right) }\left( {{ax} + b}\right) ,\varphi(x)\right\rangle=(-1)^m\int_{-\infty}^\infty {\delta }\left( {{ax} + b}\right) \varphi^{(m)}(x)\dif x=(-1)^m\frac{1}{{a}^{m}|a|}\int_{-\infty}^\infty {\delta }\left(y\right) \varphi^{(m)}(\frac{1}{a}y-\frac{b}{a})\dif y
    \]
\end{proof}
  \item 证明
  \begin{equation*}
  {PV}\frac{1}{x} = \frac{\dif}{\dif {x}} {\ln \left| x\right|   }  \text{ (在 }{\mathcal{D}}^{\prime }\text{ 中) .}
  \end{equation*}
\begin{proof}
    \begin{equation*}
        \begin{aligned}
             \left\langle\frac{\dif}{\dif {x}} {\ln \left| x\right|   } ,\varphi(x)\right\rangle&=-\lim\limits_{\varepsilon\to 0^+}\int_{|x|\geqslant\varepsilon}\ln|x|\varphi^\prime(x)\dif x\\&=-\lim\limits_{\varepsilon\to 0^+}\left(\varepsilon\ln|\varepsilon|\frac{\varphi(-\varepsilon)-\varphi(\varepsilon)}{\varepsilon}-\int_{|x|\geqslant\varepsilon}\frac{1}{x}\varphi(x)\dif x\right)\\&=\left\langle PV\frac{1}{x},\varphi(x)\right\rangle
        \end{aligned}
    \end{equation*}
\end{proof}
  \item 设分段可微函数 \(f\left( x\right)\) 在点 \({x}_{1},{x}_{2},\cdots\) 分别有有限的间断值 \({h}_{1},{h}_{2},\cdots\) ,在这些间断点外有连续的导数 \({f}^{\prime }\left( x\right)\) ,证明作为广义函数, \(f\) 的弱导数
  \begin{equation*}
  {f}^{\prime } = {f}^{\prime }\left( x\right)  + \mathop{\sum }\limits_{i}{h}_{i}\delta \left( {x - {x}_{i}}\right).
  \end{equation*}

  \item 证明每个广义函数 \(f \in  {\mathcal{D}}^{\prime }\left( {\mathbb{R}}^{n}\right)\) 都可表示为
  \begin{equation*}
  f = \mathop{\sum }\limits_{\alpha }{\partial }^{\alpha }{g}_{\alpha }\left( x\right),
  \end{equation*}
  其中右边的无限和在任何有界域中只有有限项不为零,且 \({g}_{\alpha } \in  {L}^{2}\left( {\mathbb{R}}^{n}\right)\) 的支集随 \(\left| \alpha \right|\) 的增大而无限远移,但含于 \(\operatorname{supp}f\) 的任意邻域中.

  \item 设 \(f \in  {\mathcal{D}}^{\prime }\left( {\mathbb{R}}^{n}\right) ,\varphi  \in  \mathcal{D}\left( {\mathbb{R}}^{n}\right)\). 定义 \(f\) 与 \(\varphi\) 的卷积为
\begin{equation*}
  f * \varphi \left( x\right)  = \langle f\left( {x - y}\right) ,\varphi \left( y\right) \rangle  = \langle f\left( y\right) ,\varphi \left( {x + y}\right) \rangle. 
  \end{equation*}
  试证明 \(f * \varphi  \in  \mathcal{E}\left( {\mathbb{R}}^{n}\right)\) 且  \begin{equation*}
  {\partial }^{\alpha }\left( {f * \varphi }\right)  = \left( {{\partial }^{\alpha }f}\right)  * \varphi  = f * \left( {{\partial }^{\alpha }\varphi }\right),\;\forall \alpha  \in  {\mathbb{N}}^{n},
  \end{equation*}  \begin{equation*}
  \operatorname{supp}\left( {f * \varphi }\right)  \subset  \operatorname{supp}f + \operatorname{supp}\varphi  \mathrel{\text{ := }} \{ x + y \mid  x \in  \operatorname{supp}f,y \in  \operatorname{supp}\varphi \}.
  \end{equation*}

  \item 设 \(f \in  \mathcal{D}\left( {\mathbb{R}}^{n}\right) ,\varphi ,\psi  \in  \mathcal{S}\left( {\mathbb{R}}^{n}\right)\), 证明
  \begin{equation*}
  \left( {f * \varphi }\right)  * \psi  = f * \left( {\varphi  * \psi }\right).
  \end{equation*}

  \item 设 \(f,g \in  {\mathcal{D}}^{\prime }\left( {\mathbb{R}}^{n}\right) ,\varphi  \in  \mathcal{D}\left( {\mathbb{R}}^{n}\right)\). 若 \(f * \varphi  \in  \mathcal{D}\left( {\mathbb{R}}^{n}\right) ,g * \varphi  \in  \mathcal{D}\left( {\mathbb{R}}^{n}\right)\), 定义卷积 \(f * g\) 如下
  \begin{equation*}
  \langle f * g,\varphi \rangle  = \langle f\left( x\right) ,\langle g\left( y\right) ,\varphi \left( {x + y}\right) \rangle \rangle.
  \end{equation*}
  证明 \(f * g = g * f\).
\end{problemset}

\section{广义函数的Fourier变换}

\subsection{广义函数的Fourier变换}

从偏微分方程的学习中, 我们已经看到, Fourier变换无论在理论上还是在应用上都是重要的工具. 从前一节可以看出, 广义函数的运算的定义一般是通过对偶方式由检验函数上的相应运算来定义.为了定义广义函数的Fourier变换, 我们先回顾一下检验函数的Fourier变换及其性质. 本节中,将记 \(\mathcal{D} = \mathcal{D}\left( {\mathbb{R}}^{n}\right) ,\mathcal{S} = \mathcal{S}\left( {\mathbb{R}}^{n}\right)\) .
\begin{definition}
  对任意函数 \(\varphi \left( x\right)  \in  \mathcal{D}\) ,定义其Fourier变换为
\begin{equation}
\mathcal{F}\left\lbrack  \varphi \right\rbrack  \left( \xi \right)  = \widehat{\varphi }\left( \xi \right)  \mathrel{\text{ := }} \frac{1}{\sqrt{{\left( 2\pi \right) }^{n}}}{\int }_{{\mathbb{R}}^{n}}\varphi \left( x\right) {e}^{-{\ci x} \cdot  \xi }{\dif x}, \label{3.10}
\end{equation}

其中 \(x \cdot  \xi  = {x}_{1}{\xi }_{1} + \cdots  + {x}_{n}{\xi }_{n}\) .

\end{definition}
\begin{note}
本书中,将这种变换运算本身和变换后的像函数 \(\widehat{\varphi }\left( \xi \right)\) 统称为Fourier变换, 这并不会引起混淆.
\end{note}
\begin{proposition}
  若 \(\varphi \left( x\right)  \in  \mathcal{D}\) ,则有 \(\widehat{\varphi }\left( \xi \right)  \in  \mathcal{S}\) .

\end{proposition}
\begin{proof}
为简单起见,我们以一维为例来证明. 由于 \(\varphi\) 及其各阶导数在某有限区间 \(\left\lbrack  {-a,a}\right\rbrack\) 外恒等于零,故函数 \(\widehat{\varphi }\left( \xi \right)\) 的定义域 \(\mathbb{R}\) 可延拓到复平面 \(\mathbb{C} = \{ \zeta  \mid  \zeta  =\)  \(\xi  + {\ci\eta },\xi ,\eta  \in  \mathbb{R}\}\) :
\begin{equation}
\widehat{\varphi }\left( \zeta \right)  = \frac{1}{\sqrt{2\pi }}{\int }_{-a}^{a}\varphi \left( x\right) {e}^{-{\ci x\zeta }}{\dif x} = \frac{1}{\sqrt{2\pi }}{\int }_{-a}^{a}\varphi \left( x\right) {e}^{-{\ci x\xi }}{e}^{\eta x}{\dif x}. \label{3.11}
\end{equation}

该变换通常称为复数域上的Fourier变换或Fourier-Laplace变换. \ref{3.11} 显然容许在积分号内求微分运算,且 \(\zeta\) 可为 \(\mathbb{C}\) 中任意有限点. 故 \(\widehat{\varphi }\left( \zeta \right)\) 为 \(\zeta\) 的整解析函数. 不难验证
\begin{equation*}
\mathcal{F}\left\lbrack  {\varphi }^{\left( k\right) }\right\rbrack  \left( \zeta \right)  = {\left( \ci\zeta \right) }^{k}\widehat{\varphi }\left( \zeta \right) ,\;k = 0,1,\cdots .
\end{equation*}

从而有估计
\begin{equation}
{\left| \zeta \right| }^{k}\left| {\widehat{\varphi }\left( \zeta \right) }\right|  = \left| {{\int }_{-a}^{a}{\varphi }^{\left( k\right) }\left( x\right) {e}^{-{\ci x\xi } + {x\eta }}{\dif x}}\right|  \leqslant  C\left( k\right) {e}^{a\left| \eta \right| },\;k = 0,1,2,\cdots , \label{3.12}
\end{equation}

其中
\begin{equation*}
C\left( k\right)  = {\int }_{-a}^{a}\left| {{\varphi }^{\left( k\right) }\left( x\right) }\right| {\dif x}.
\end{equation*}

这表明 \(\widehat{\varphi }\left( \zeta \right)\) 对有界的 \(\left| \eta \right|\) 关于 \(\left| \zeta \right|\) 速降. 由\ref{3.11}用同样方法可证明 \(\widehat{\varphi }\left( \zeta \right)\) 的各阶导数都具有这种性质. 特别有 \(\widehat{\varphi }\left( \xi \right)  \in  \mathcal{S}\) .

\end{proof}

在Fourier变换的应用中, 我们往往需要使用逆Fourier变换. 而由上可以看出, Fourier变换将 \(\mathcal{D}\) 中函数映为 \(\mathcal{S}\) . 这给Fourier变换的应用带来不便. 为此,我们有必要将Fourier变换的定义拓展到速降函数空间 \(\mathcal{A}\) .
\begin{definition}
   设 \(\varphi  \in  \mathcal{S}\left( {\mathbb{R}}^{n}\right)\) . 定义 \(\varphi\) 的Fourier变换为
\begin{equation}
\widehat{\varphi } = \frac{1}{\sqrt{{\left( 2\pi \right) }^{n}}}{\int }_{{\mathbb{R}}^{n}}{e}^{-{\ci x} \cdot  \xi }\varphi \left( x\right) {\dif x}, \label{3.13}
\end{equation}

其中 \(x \cdot  \xi  = \mathop{\sum }\limits_{{i = 1}}^{n}{x}_{i}{\xi }_{i}\) . 定义 \(\varphi\) 的Fourier逆变换为\begin{equation*}
\check{\varphi }\left( x\right)  = \frac{1}{\sqrt{{\left( 2\pi \right) }^{n}}}{\int }_{{\mathbb{R}}^{n}}{e}^{{\ci x} \cdot  \xi }\varphi \left( \xi \right) {\dif\xi }.
\end{equation*}

有时也记 \(\widehat{\varphi } = \mathcal{F}\left\lbrack  \varphi \right\rbrack\) 和 \(\check{\varphi } = \overline{\mathcal{F}}\left\lbrack  \varphi \right\rbrack\) .
\end{definition}



由于当 \(\varphi  \in  \mathcal{S}\) 时\ref{3.13}的积分区域为全空间 \({\mathbb{R}}^{n}\) , \(\widehat{\varphi }\left( \xi \right)\) 一般不能延拓为复变量 \(\xi  + {\ci\eta }\) 的函数(这时因子 \({e}^{x\eta }\) 的增长速度可能超过 \(\varphi \left( x\right)\) 的速降速度). 尽管如此, 但我们依然能证明
\begin{lemma}
  Fourier变换 \(\mathcal{F}\) 及Fourier逆变换 \(\overline{\mathcal{F}}\) 都是从 \(\mathcal{S}\left( {\mathbb{R}}^{n}\right)\) 到 \(\mathcal{S}\left( {\mathbb{R}}^{n}\right)\) 的连续线性变换. 而且对任意多重指标 \(\alpha ,\beta\) ,我们有
\begin{equation}
\left( {{\left( \ci\xi \right) }^{\alpha }{\partial }^{\beta }\widehat{\varphi }}\right) \left( \xi \right)  = \mathcal{F}\left\lbrack  {{\partial }^{\alpha }\left( {{\left( -\ci x\right) }^{\beta }\varphi \left( x\right) }\right) }\right\rbrack  . \label{3.14}
\end{equation}
\end{lemma}
\begin{proof}
  我们只证Fourier变换的情形, 类似可证Fourier逆变换的结论. 首先, 线性性是显然的. 在证明连续性之前, 考虑
\begin{equation}
\begin{aligned}
    \left( {{\xi }^{\alpha }{\partial }^{\beta }\widehat{\varphi }}\right) \left( \xi \right)  &= \frac{1}{\sqrt{{\left( 2\pi \right) }^{n}}}{\int }_{{\mathbb{R}}^{n}}{\xi }^{\alpha }{\left( -\ci x\right) }^{\beta }{e}^{-{\ci\xi } \cdot  x}\varphi \left( x\right) {\dif x}
\\&= \frac{1}{\sqrt{{\left( 2\pi \right) }^{n}}}{\int }_{{\mathbb{R}}^{n}}\frac{1}{{\left( -\ci\right) }^{\left| \alpha \right| }}\left( {{\partial }_{x}^{\alpha }{e}^{-{\ci\xi } \cdot  x}}\right) {\left( -\ci x\right) }^{\beta }\varphi \left( x\right) {\dif x}
\\&= \frac{{\left( -\ci\right) }^{\left| \alpha \right| }}{\sqrt{{\left( 2\pi \right) }^{n}}}{\int }_{{\mathbb{R}}^{n}}{e}^{-{\ci\xi } \cdot  x}{\partial }_{x}^{\alpha }\left( {{\left( -\ci x\right) }^{\beta }\varphi \left( x\right) }\right) {\dif x}(\text{用}\int_{\mathbb{R}^n}\partial^\alpha_xu\cdot v\dif x=(-1)^{|\alpha|}\int_{\mathbb{R}^n}\partial^\alpha_xv\cdot u\dif x). \label{3.15}
\end{aligned}
\end{equation}

由\ref{3.15}可得
\begin{equation*}
\mathop{\sup }\limits_{\xi }\left| {{\xi }^{\alpha }\left( {{\partial }^{\beta }\widehat{\varphi }}\right) \left( \xi \right) }\right|  \leqslant  \frac{1}{\sqrt{{\left( 2\pi \right) }^{n}}}{\int }_{{\mathbb{R}}^{n}}\left| {{\partial }_{x}^{\alpha }\left( {{x}^{\beta }\varphi }\right) }\right| {\dif x} < \infty .
\end{equation*}

这意味着Fourier变换 \(\mathcal{F}\) 映 \(\mathcal{S}\left( {\mathbb{R}}^{n}\right)\) 到 \(\mathcal{S}\left( {\mathbb{R}}^{n}\right)\) ,且\ref{3.14}成立.

此外,若 \(k\) 足够大,则有 \({\int }_{{\mathbb{R}}^{n}}{\left( 1 + {x}^{2}\right) }^{-k}{\dif x} < \infty\) . 从而有
\begin{equation*}
\mathop{\sup }\limits_{\xi }\left| {{\xi }^{\alpha }\left( {{\partial }^{\beta }\widehat{\varphi }}\right) \left( \xi \right) }\right|  \leqslant  \frac{1}{\sqrt{{\left( 2\pi \right) }^{n}}}{\int }_{{\mathbb{R}}^{n}}\frac{{\left( 1 + {x}^{2}\right) }^{-k}}{{\left( 1 + {x}^{2}\right) }^{-k}}\left| {{\partial }_{x}^{\alpha }{\left( -\ci x\right) }^{\beta }\varphi \left( x\right) }\right| {\dif x}
\end{equation*}
\begin{equation*}
\leqslant  \frac{1}{\sqrt{{\left( 2\pi \right) }^{n}}}\left( {{\int }_{{\mathbb{R}}^{n}}{\left( 1 + {x}^{2}\right) }^{-k}{\dif x}}\right) \mathop{\sup }\limits_{x}\left\{  {{\left( 1 + {x}^{2}\right) }^{k}\left| {{\partial }_{x}^{\alpha }{\left( -\ci x\right) }^{\beta }\varphi \left( x\right) }\right| }\right\}  .
\end{equation*}

由莱布尼兹法则可知,存在有限个多重指标 \({\alpha }_{j},{\beta }_{j}\) 及常数 \({c}_{j}\) 使得
\begin{equation*}
\mathop{\sup }\limits_{\xi }\left| {{\xi }^{\alpha }\left( {{\partial }^{\beta }\widehat{\varphi }}\right) \left( \xi \right) }\right|  \leqslant  \mathop{\sum }\limits_{j}{c}_{j}\mathop{\sup }\limits_{x}\left\{  \left| {{x}^{{\alpha }_{j}}{\partial }^{{\beta }_{j}}\varphi \left( x\right) }\right| \right\}  .
\end{equation*}

这意味着Fourier变换 \(\mathcal{F}\) 的连续性.

\end{proof}
\begin{proposition}
 \( \mathcal{D}\left( {\mathbb{R}}^{n}\right) \) 在 \( \mathcal{S}\left( {\mathbb{R}}^{n}\right) \) 中稠密.
\end{proposition}
\begin{proof}
对任意 \( f \in  \mathcal{S}\left( {\mathbb{R}}^{n}\right) \) ,我们定义
    \[\beta_R(x)=j_1(x)* \chi_{B(0,R)}(x),{f}_{\nu }\left( x\right)  = {\beta }_{\nu }\left( x\right) f\left( x\right)  \in  \mathcal{D}\left( {\mathbb{R}}^{n}\right) \]
下面证明,对任意的 \( \alpha,\beta  \in  {\mathbb{N}}^{n}, \) ,当 \( \nu  \rightarrow  \infty \) 时成立
\( \mathop{\sup }\limits_{{x \in  {\mathbb{R}}^{n}}}\left| {{x }^{\alpha}{\partial }^{\beta }\left( {{f}_{\nu }\left( x\right)  - f\left( x\right) }\right) }\right|  \rightarrow  0. \)

实际上由于 \( {f}_{\nu }\left( x\right)  - f\left( x\right)  = \left( {{\beta }_{\nu }\left( x\right)  - 1}\right) f\left( x\right) \) ,则
\( {\partial }^{\beta }\left( {{f}_{\nu }\left( x\right)  - f\left( x\right) }\right)  = \mathop{\sum }\limits_{{\beta  = {\alpha }_{1} + {\alpha }_{2}}}\frac{\alpha !}{{\alpha }_{1}!{\alpha }_{2}!}{\partial }^{{\alpha }_{1}}\left( {{\beta }_{\nu }\left( x\right)  - 1}\right) {\partial }^{{\alpha }_{2}}f\left( x\right) . \)
注意由 \( f \in  \mathcal{S}\left( {\mathbb{R}}^{n}\right) \) 知 \( \left| {{x }^{\alpha}{\partial }^{\beta }f\left( x\right) }\right|  \leqslant  C{\left(  {\left| x\right| }^{2}\right) }^{-1} \) ,从而成立
\begin{equation*}
    \begin{aligned}
        {x}^{\alpha}\left| {{\partial }^{\beta }\left( {{f}_{\nu }\left( x\right)  - f\left( x\right) }\right) }\right| 
\leqslant  \mathop{\sum }\limits_{{\beta  = {\alpha }_{1} + {\alpha }_{2}}}\frac{\alpha !}{{\alpha }_{1}!{\alpha }_{2}!}\left| {{x}^{\alpha}{\partial }^{{\alpha }_{1}}\left( {{\beta }_{\nu }\left( x\right)  - 1}\right) {\partial }^{{\alpha }_{2}}f\left( x\right) }\right| \\
 \leqslant  \mathop{\sum }\limits_{{{\alpha }^{\prime} \leqslant  \beta }}{C}_{{\alpha }^{\prime }}\left| {{\partial }^{{\alpha }^{\prime }}\left( {{\beta }_{\nu }\left( x\right)  - 1}\right) }\right| {\left( {\left| x\right| }^{2}\right) }^{-1}. 
    \end{aligned}
\end{equation*}

又因为 \( {\partial }^{{\alpha }^{\prime}}\left( {{\beta }_{\nu }\left( x\right)  - 1}\right) \) 在全空间有界,而且当 \( \left| x\right|  \leqslant  \nu  - 1 \) 时, \( {\beta }_{\nu }\left( x\right)  - 1 = 0\). 故
\( \mathop{\lim }\limits_{{\nu  \rightarrow  \infty }}\mathop{\sup }\limits_{{x \in  {\mathbb{R}}^{n}}}\left|  {x}^{\alpha}\left| {{\partial }^{\beta }\left( {{f}_{\nu }\left( x\right)  - f\left( x\right) }\right) }\right| \right|  \leqslant  \mathop{\lim }\limits_{{\nu  \rightarrow  \infty }}\frac{C}{{\left( \nu  - 1\right) }^{2}} = 0. \)
命题证毕.
\end{proof}
下面来证明Fourier逆定理, 其证明思想来源于Fourier本人.
\begin{theorem}
  Fourier变换 \(\mathcal{F}\) 是从 \(\mathcal{S}\left( {\mathbb{R}}^{n}\right)\) 到 \(\mathcal{S}\left( {\mathbb{R}}^{n}\right)\) 的线性可逆连续双射,其逆映射为Fourier逆变换. 即 \(\overline{\mathcal{F}}\left\lbrack  {\mathcal{F}\left\lbrack  \varphi \right\rbrack  }\right\rbrack   = \varphi  = \mathcal{F}\left\lbrack  {\overline{\mathcal{F}}\left\lbrack  \varphi \right\rbrack  }\right\rbrack\) .
\end{theorem}
\begin{proof}
  我们只证 \(\overline{\mathcal{F}}\left\lbrack  {\mathcal{F}\left\lbrack  \varphi \right\rbrack  }\right\rbrack   = \varphi\) . 由于 \(\mathcal{F}\) 及 \(\overline{\mathcal{F}}\) 都是从 \(\mathcal{S}\left( {\mathbb{R}}^{n}\right)\) 到 \(\mathcal{S}\left( {\mathbb{R}}^{n}\right)\) 的连续映射, 故只需证明对稠密子集 \({C}_{0}^{\infty }\left( {\mathbb{R}}^{n}\right)\) 中的函数 \(\varphi\) 成立 \(\overline{\mathcal{F}}\left\lbrack  {\mathcal{F}\left\lbrack  \varphi \right\rbrack  }\right\rbrack   = \varphi\) 即可. 用 \({C}_{\varepsilon }\) 记 \({\mathbb{R}}^{n}\) 中以原点为中心体积为 \({\left( \frac{2}{\varepsilon }\right) }^{n}\) 的方体. 选择 \(\varepsilon\) 足够小使得 \(\varphi\) 的支集包含在 \({C}_{\varepsilon }\) . 令
\begin{equation*}
{K}_{\varepsilon } = \left\{  {k = \left( {{k}_{1},{k}_{2},\cdots ,{k}_{n}}\right)  \in  {\mathbb{R}}^{n} \mid  {k}_{i}/{\pi \varepsilon }\text{ 为整数, }i = 1,2,\cdots ,n}\right\}  .
\end{equation*}

则 \({\left\{  {\left( \frac{1}{2}\varepsilon \right) }^{\frac{n}{2}}{e}^{{ik} \cdot  x}\right\}  }_{k \in  {K}_{\varepsilon }}\) 为 \({L}^{2}\left( {C}_{\varepsilon }\right)\) 的正交规范基. 由于 \(\varphi\) 连续可微,故 \(\varphi\) 的Fourier级数
\begin{equation}
\varphi \left( x\right)  = \mathop{\sum }\limits_{{k \in  {K}_{\varepsilon }}}\left( {\varphi ,{\left( \frac{1}{2}\varepsilon \right) }^{\frac{n}{2}}{e}^{{ik} \cdot  x}}\right) {\left( \frac{1}{2}\varepsilon \right) }^{\frac{n}{2}}{e}^{{ik} \cdot  x} \label{3.16}
\end{equation}

在 \({C}_{\varepsilon }\) 中一致收敛到 \(\varphi\) . 因此,
\begin{equation}
\varphi \left( x\right)  = \mathop{\sum }\limits_{{k \in  {K}_{\varepsilon }}}\frac{\widehat{\varphi }{e}^{{ik} \cdot  x}}{{\left( 2\pi \right) }^{\frac{n}{2}}}{\left( \pi \varepsilon \right) }^{n}. \label{3.17}
\end{equation}

由于 \({\mathbb{R}}^{n}\) 可表示成中心位于 \({K}_{\varepsilon }\) 中点体积为 \({\left( \pi \varepsilon \right) }^{n}\) 的方体的无交并,而且(3.17)的右边恰为函数 \(\widehat{\varphi }\left( k\right) {e}^{{ik} \cdot  x}/{\left( 2\pi \right) }^{\frac{n}{2}}\) 的黎曼和. 有引理知, \(\widehat{\varphi }{e}^{{ik} \cdot  x} \in  \mathcal{S}\left( {\mathbb{R}}^{n}\right)\) . 从而令 \(\varepsilon\) 趋于零得, \(\overline{\mathcal{F}}\left\lbrack  {\mathcal{F}\left\lbrack  \varphi \right\rbrack  }\right\rbrack   = \varphi\) . 同理可证 \(\mathcal{F}\left\lbrack  {\overline{\mathcal{F}}\left\lbrack  \varphi \right\rbrack  }\right\rbrack   = \varphi\) .

最后,由 \(\overline{\mathcal{F}}\left\lbrack  {\mathcal{F}\left\lbrack  \varphi \right\rbrack  }\right\rbrack   = \varphi\) 知 \(\mathcal{F}\) 为单射; 由 \(\mathcal{F}\left\lbrack  {\overline{\mathcal{F}}\left\lbrack  \varphi \right\rbrack  }\right\rbrack   = \varphi\) 知 \(\mathcal{F}\) 为满射. 至此,定理得证.
\end{proof}
\begin{corollary}
  若 \(\varphi  \in  \mathcal{S}\left( {\mathbb{R}}^{n}\right)\) ,则
\begin{equation*}
{\int }_{{\mathbb{R}}^{n}}{\left| \varphi \left( x\right) \right| }^{2}{\dif x} = {\int }_{{\mathbb{R}}^{n}}{\left| \widehat{\varphi }\left( k\right) \right| }^{2}{\dif k}.
\end{equation*}
\end{corollary}
\begin{proof}
  若 \(\varphi\) 有紧支集,由 (3.16) 知,对足够小的 \(\varepsilon\) 有

\begin{equation*}
\varphi \left( x\right)  = \mathop{\sum }\limits_{{k \in  {K}_{\varepsilon }}}\left( {\varphi ,{\left( \frac{1}{2}\varepsilon \right) }^{\frac{n}{2}}{e}^{{ik} \cdot  x}}\right) {\left( \frac{1}{2}\varepsilon \right) }^{\frac{n}{2}}{e}^{{ik} \cdot  x}.
\end{equation*}

因此,
\begin{equation*}
{\int }_{{\mathbb{R}}^{n}}{\left| \varphi \left( x\right) \right| }^{2}{\dif x} = {\int }_{{C}_{\varepsilon }}{\left| \varphi \left( x\right) \right| }^{2}{\dif x}
= \mathop{\sum }\limits_{{k \in  {K}_{\varepsilon }}}{\left| \left( \varphi ,{\left( \frac{1}{2}\varepsilon \right) }^{\frac{n}{2}}{e}^{{ik} \cdot  x}\right) \right| }^{2}
= \mathop{\sum }\limits_{{k \in  {K}_{\varepsilon }}}{\left| \widehat{\varphi }\left( k\right) \right| }^{2}{\left( \pi \varepsilon \right) }^{n}
\rightarrow  {\int }_{{\mathbb{R}}^{n}}{\left| \widehat{\varphi }\left( k\right) \right| }^{2}{\dif k},\;\left( {\varepsilon  \rightarrow  0}\right) .
\end{equation*}

这意味着推论对 \(\varphi  \in  {C}_{0}^{\infty }\left( {\mathbb{R}}^{n}\right)\) 成立. 由于Fourier变换及 \(\parallel  \cdot  {\parallel }_{2}\) 在 \(\mathcal{S}\left( {\mathbb{R}}^{n}\right)\) 上连续, 且 \({C}_{0}^{\infty }\left( {\mathbb{R}}^{n}\right)\) 在 \(\mathcal{Y}\) 中稠密，从而推论结论对所有 \(\varphi  \in\)  \(\mathcal{S}\left( {\mathbb{R}}^{n}\right)\) 都成立.
\end{proof}

\begin{example}
  设 \(\alpha  > 0\) ,求 \(\varphi \left( x\right)  = {e}^{-\frac{\alpha {x}^{2}}{2}} \in  \mathcal{S}\left( \mathbb{R}\right)\) 的Fourier变换.
\end{example}
\begin{solution}
\begin{equation*}
  \begin{aligned}
\widehat{\varphi }\left( \xi \right)  = \frac{1}{\sqrt{2\pi }}{\int }_{\mathbb{R}}{e}^{-\frac{\alpha {x}^{2}}{2}}{e}^{-{\ci\xi x}}{\dif x}
= \frac{1}{\sqrt{2\pi }}{\int }_{\mathbb{R}}\sqrt{\frac{2}{\alpha }}\exp \left( {-{t}^{2} - {it\xi }\sqrt{\frac{2}{\alpha }}}\right) {\dif t}
= \frac{{e}^{-{\xi }^{2}/{2\alpha }}}{\sqrt{\alpha \pi }}{\int }_{\mathbb{R}}\exp \left( {-{\left( t + \ci\frac{\xi }{\sqrt{\alpha \pi }}\right) }^{2}}\right) {\dif t}\\
= \frac{{e}^{-{\xi }^{2}/{2\alpha }}}{\sqrt{\alpha \pi }}{\int }_{\mathbb{R}}{e}^{-{t}^{2}}{\dif t}
= \frac{{e}^{-{\xi }^{2}/{2\alpha }}}{\sqrt{\alpha }}.
  \end{aligned}
\end{equation*}

\end{solution}

从 (3.10) 可以看出,对 \(\varphi \left( x\right)  \in  {L}^{1}\left( {\mathbb{R}}^{n}\right)\) ,我们都可以定义Fourier变换. 人们自然会问,哪些函数存在Fourier 变换呢? 由于一般的 \({L}^{p}\left( {\mathbb{R}}^{n}\right) \left( {p > 1}\right)\) 函数不一定属于 \({L}^{1}\left( {\mathbb{R}}^{n}\right)\) ,从而也就不一定能定义Fourier 变换. 实际上,最简单的函数 \(\varphi \left( x\right)  \equiv\) 1(或更一般的多项式函数)就不能按(3.10)定义Fourier变换. 因此,为了能使Fourier变换这一重要工具得到更好地应用, 我们有必要扩大Fourier变换的定义范围. 在引进广义函数的Fourier变换之前, 我们先研究 \(\mathcal{D}\) 中的函数的Fourier变换的如下性质.
\begin{proposition}
  若 \(\varphi ,\psi  \in  \mathcal{S}\) ,则
\begin{equation*}
\langle \mathcal{F}\left\lbrack  \varphi \right\rbrack  ,\psi \rangle  = \langle \varphi ,\mathcal{F}\left\lbrack  \psi \right\rbrack  \rangle ,
\end{equation*}
\begin{equation*}
\langle \overline{\mathcal{F}}\left\lbrack  \varphi \right\rbrack  ,\psi \rangle  = \langle \varphi ,\overline{\mathcal{F}}\left\lbrack  \psi \right\rbrack  \rangle .
\end{equation*}
\end{proposition}
\begin{proof}
  我们只证第一个等式, 第二个等式可同理证得. 实际上,
\begin{equation*}
\langle \mathcal{F}\left\lbrack  \varphi \right\rbrack  ,\psi \rangle  = {\int }_{{\mathbb{R}}^{n}}\left\lbrack  {{\int }_{{\mathbb{R}}^{n}}\varphi \left( x\right) {e}^{-{\ci x} \cdot  \xi }{\dif x}}\right\rbrack  \psi \left( \xi \right) {\dif\xi }\overset{\text{ Fubini定理 }}{=}
{\int }_{{\mathbb{R}}^{n}}\varphi \left( x\right) \left\lbrack  {{\int }_{{\mathbb{R}}^{n}}\psi \left( \xi \right) {e}^{-{\ci x} \cdot  \xi }{\dif\xi }}\right\rbrack  {\dif x} = \langle \varphi ,\mathcal{F}\left\lbrack  \psi \right\rbrack  \rangle \text{ . }
\end{equation*}
\end{proof}

\begin{definition}
  在 \({\mathcal{S}}^{\prime }\) 上定义 \(\widetilde{\mathcal{F}} = {\mathcal{F}}^{ * },\left( {\widetilde{\overline{\mathcal{F}}} = {\left( \overline{\mathcal{F}}\right) }^{ * }}\right)\) 为广义Fourier (逆) 变换. 在不会混淆时,记号 \({}^{ \sim  }\) 可以略去. 有时也记为 \(\mathcal{F}\left\lbrack  \varphi \right\rbrack   = \widehat{\varphi }\) .

\end{definition}
\begin{note}
  \begin{enumerate}
    \item Fourier变换 \(\mathcal{F}\) 及Fourier 逆变换 \(\overline{\mathcal{F}}\) 都是从 \({\mathcal{I}}^{\prime }\) 到 \({\mathcal{I}}^{\prime }\) 的线性映射.
    \item 由命题3.2知,当限制在 \(\mathcal{Y}\) 上时, \({\mathcal{S}}^{\prime }\) 上的Fourier 变换及其逆变换分别与普通的Fourier变换, Fourier逆变换一致.

  \end{enumerate}
\end{note}
\begin{proposition}
  Fourier变换 \(\mathcal{F}\) 和Fourier逆变换 \(\overline{\mathcal{F}}\) 都是从 \({\mathcal{S}}^{\prime }\) 到 \({\mathcal{S}}^{\prime }\) 的弱连续线性算子.
\end{proposition}
\begin{proof}
  我们只证 \(\mathcal{F}\) 的情形. 若 \({f}_{n} \rightharpoonup  f\) ,则由引理3.1知,对任意 \(\varphi  \in  \mathcal{S}\) 都有 \({f}_{n}\left( {\mathcal{F}\left\lbrack  \varphi \right\rbrack  }\right)  \rightarrow  f\left( {\mathcal{F}\left\lbrack  \varphi \right\rbrack  }\right)\) . 这意味着 \(\mathcal{F}\left\lbrack  {f}_{n}\right\rbrack  \left( \varphi \right)  \rightarrow  \mathcal{F}\left\lbrack  f\right\rbrack  \left( \varphi \right)\) ,即 \(\mathcal{F}\left\lbrack  {f}_{n}\right\rbrack   \rightharpoonup  \mathcal{F}\left\lbrack  f\right\rbrack\) . 因此 \(F\) 为弱连续算子. 定理剩余证明留给读者.

\end{proof}
\begin{theorem}
  \begin{enumerate}
    \item  \(\mathcal{F}\left\lbrack  {{\partial }^{\alpha }f}\right\rbrack   = {\left( \ci\xi \right) }^{\alpha }\mathcal{F}\left\lbrack  f\right\rbrack\) ,
    \item \(\mathcal{F}\left\lbrack  {{\left( -\ci x\right) }^{\alpha }f}\right\rbrack   = {\partial }^{\alpha }\mathcal{F}\left\lbrack  f\right\rbrack\) ,
    \item \(\mathcal{F}\left\lbrack  {{\widetilde{\tau }}_{{x}_{0}}f}\right\rbrack   = {e}^{-\ci{x}_{0} \cdot  \xi }\mathcal{F}\left\lbrack  f\right\rbrack\) ,
    \item  \(\mathcal{F}\left\lbrack  {{e}^{{x}_{0} \cdot  x}f}\right\rbrack   = {\widetilde{\tau }}_{{x}_{0}}\mathcal{F}\left\lbrack  f\right\rbrack\)

  \end{enumerate}
\end{theorem}
\begin{proof}
  这里我们只证明结论( 1 ),其他结论的证明留给读者作为习题. 由定义知,

对任意 \(\varphi  \in  \mathcal{S}\) ,有
\begin{equation*}
\left\langle  {\mathcal{F}\left\lbrack  {{\partial }^{\alpha }f}\right\rbrack  ,\varphi }\right\rangle   = \left\langle  {{\partial }^{\alpha }f,\mathcal{F}\left\lbrack  \varphi \right\rbrack  }\right\rangle   = {\left( -1\right) }^{\left| \alpha \right| }\left\langle  {f,{\partial }^{\alpha }\mathcal{F}\left\lbrack  \varphi \right\rbrack  }\right\rangle
\end{equation*}
\begin{equation*}
= {\left( -1\right) }^{\left| \alpha \right| }\left\langle  {f,\mathcal{F}\left\lbrack  {{\left( -\ci\xi \right) }^{\alpha }\varphi }\right\rbrack  }\right\rangle   = {\left( -1\right) }^{\left| \alpha \right| }\left\langle  {\mathcal{F}\left\lbrack  f\right\rbrack  ,{\left( -\ci\xi \right) }^{\alpha }\varphi }\right\rangle
\end{equation*}
\begin{equation*}
= \left\langle  {\mathcal{F}\left\lbrack  f\right\rbrack  ,{\left( \ci\xi \right) }^{\alpha }\varphi }\right\rangle   = \left\langle  {{\left( \ci\xi \right) }^{\alpha }\mathcal{F}\left\lbrack  f\right\rbrack  ,\varphi }\right\rangle  ,
\end{equation*}

这里第三个等式我们用到等式\ref{3.14}.
\end{proof}

\begin{example}
  在 \({\mathcal{S}}^{\prime }\) 中,证明 \(\mathcal{F}\left\lbrack  \delta \right\rbrack   = \frac{1}{\sqrt{{\left( 2\pi \right) }^{n}}},\overline{\mathcal{F}}\left\lbrack  1\right\rbrack   = \sqrt{{\left( 2\pi \right) }^{n}}\delta\) .
\end{example}
\begin{proof}
  由定义知,对任意 \(\varphi  \in  \mathcal{S}\) ,有
\begin{equation*}
\langle \mathcal{F}\left\lbrack  \delta \right\rbrack  ,\varphi \rangle  = \langle \delta ,\mathcal{F}\left\lbrack  \varphi \right\rbrack  \rangle  = \mathcal{F}\left\lbrack  \varphi \right\rbrack  \left( 0\right)  = \frac{1}{\sqrt{{\left( 2\pi \right) }^{n}}}{\int }_{{\mathbb{R}}^{n}}\varphi \left( x\right) {\dif x} = \left\langle  {\frac{1}{\sqrt{{\left( 2\pi \right) }^{n}}},\varphi }\right\rangle  .
\end{equation*}


同理可证 \(\overline{\mathcal{F}}\left\lbrack  1\right\rbrack   = \sqrt{{\left( 2\pi \right) }^{n}}\delta\) .
\end{proof}

\begin{example}
  在 \({\mathcal{S}}^{\prime }\) 中,证明 \(\mathcal{F}\left\lbrack  1\right\rbrack   = \sqrt{{\left( 2\pi \right) }^{n}}\delta ,\overline{\mathcal{F}}\left\lbrack  \delta \right\rbrack   = \frac{1}{\sqrt{{\left( 2\pi \right) }^{n}}}\) .
\end{example}

 \begin{proof}
  对任意 \(\varphi  \in  \mathcal{S}\) ,有  \begin{equation*}
  \langle \mathcal{F}\left\lbrack  1\right\rbrack  ,\varphi \rangle  = \langle 1,\mathcal{F}\left\lbrack  \varphi \right\rbrack  \rangle  = {\int }_{{\mathbb{R}}^{n}}\mathcal{F}\left\lbrack  \varphi \right\rbrack  \left( x\right) {\dif x} = \frac{1}{\sqrt{{\left( 2\pi \right) }^{n}}}{\int }_{{\mathbb{R}}^{n}}\int _{\mathbb{R}^n}e^{-\ci yx}\varphi(y)\dif y {\dif x}
  \end{equation*}\begin{equation*}
 \frac{(2\pi)^n}{\sqrt{{\left( 2\pi \right) }^{n}}}{\int }_{{\mathbb{R}}^{n}}\delta(y)\varphi(y)\dif y = \sqrt{{\left( 2\pi \right) }^{n}}\varphi \left( 0\right)  = \sqrt{{\left( 2\pi \right) }^{n}}\langle \delta ,\varphi \rangle .
  \end{equation*}
  
  另外,\begin{equation*}
  \langle \overline{\mathcal{F}}\left\lbrack  \delta \right\rbrack  ,\varphi \rangle  = \langle \delta ,\overline{\mathcal{F}}\left\lbrack  \varphi \right\rbrack  \rangle  = \overline{\mathcal{F}}\left\lbrack  \varphi \right\rbrack  \left( 0\right)  = \frac{1}{\sqrt{{\left( 2\pi \right) }^{n}}}{\int }_{{\mathbb{R}}^{n}}\varphi \left( x\right) {e}^{{\ci x} \cdot  0}{\dif x} = \left\langle  {\frac{1}{\sqrt{{\left( 2\pi \right) }^{n}}},\varphi }\right\rangle  .
  \end{equation*}
  
  故 \(\overline{\mathcal{F}}\left\lbrack  \delta \right\rbrack   = \frac{1}{\sqrt{{\left( 2\pi \right) }^{n}}}\) .
\end{proof}
\begin{example}
    下面证明几个有用的结论
    \begin{enumerate}
        \item $\int_{\mathbb{R}^n}e^{\ci x\xi}\dif \xi=(2\pi)^n\delta(x)$.
        \item $\int_{\mathbb{R}}\frac{e^{\ci x\xi}}{x}\dif x=(2\pi)^n\delta(\xi)$.
    \end{enumerate}
\end{example}
\begin{proof}
1.\[
\overline{\mathcal{F}}\left\lbrack \mathcal{F}\left\lbrack \delta\right\rbrack \right\rbrack =\overline{\mathcal{F}}\left\lbrack  \frac{1}{\sqrt{(2\pi)^n}}\right\rbrack  =\int_{\mathbb{R}^n}\frac{1}{(2\pi)^n}e^{\ci x\xi}\dif \xi=\delta(x)
\]

    
\end{proof}
 \begin{example}
   在 \({\mathcal{S}}^{\prime }\) 中,证明 \(\mathcal{F}\left\lbrack  {e}^{-\frac{{\left| x\right| }^{2}}{2}}\right\rbrack   = {e}^{-\frac{{\left| \xi \right| }^{2}}{2}}\) .
  
\end{example}
 \begin{proof}
  当 \(n = 1\) 时,令 \(\varphi \left( x\right)  = {e}^{-\frac{{x}^{2}}{2}}\) ,则  \begin{equation}
  {\varphi }^{\prime }\left( x\right)  + {x\varphi }\left( x\right)  = 0,\;\varphi \left( 0\right)  = 1. \label{3.18}
  \end{equation}
  
  在方程两端求Fourier变换, 并利用定理3.5可得\begin{equation*}
  {\ci\xi }\widehat{\varphi }\left( \xi \right)  + \ci{\widehat{\varphi }}^{\prime }\left( \xi \right)  = 0.
  \end{equation*}
  
  故\begin{equation}
  {\widehat{\varphi }}^{\prime }\left( \xi \right)  + \xi \widehat{\varphi }\left( \xi \right)  = 0,\;\widehat{\varphi }\left( 0\right)  = 1. \label{3.19}
  \end{equation}
  
  比较\ref{3.18}与\ref{3.19}知, \(\varphi\) 与 \(\widehat{\varphi }\) 满足同一个微分方程和相同的初值. 由常微分方程解的唯一性知,  \begin{equation*}
  \widehat{\varphi }\left( \xi \right)  = {e}^{-\frac{{\xi }^{2}}{2}}.
  \end{equation*}
  
  当 \(n > 1\) 时,即 \(x,\xi  \in  {\mathbb{R}}^{n}\) ,有
\begin{equation*}
  \mathcal{F}\left\lbrack  {e}^{-\frac{{\left| x\right| }^{2}}{2}}\right\rbrack   = \frac{1}{\sqrt{{\left( 2\pi \right) }^{n}}}{\int }_{{\mathbb{R}}^{n}}{e}^{-\frac{{\left| x\right| }^{2}}{2}}{e}^{-{\ci x} \cdot  \xi }{\dif x} = \frac{1}{\sqrt{{\left( 2\pi \right) }^{n}}}\mathop{\prod }\limits_{{j = 1}}^{n}{\int }_{\mathbb{R}}{e}^{-{x}_{j}{\xi }_{j}}{e}^{-\frac{{x}_{j}^{2}}{2}}d{x}_{j}
  = \mathop{\prod }\limits_{{j = 1}}^{n}{e}^{-\frac{{\xi }^{2}}{2}} = {e}^{-\frac{{\left| \xi \right| }^{2}}{2}}
  \end{equation*}
  
\end{proof}
\subsection{Fourier变换的应用}

Fourier变换是研究常系数线性偏微分方程的非常重要的工具之一. 下面,我们给出几个例子来说明.
 \begin{example}
  (热方程的基本解) 考虑初值问题
  \begin{equation}
  \left\{  \begin{matrix} {u}_{t} - \bigtriangleup u = 0, & \left( {x,t}\right)  \in  {\mathbb{R}}^{n} \times  \left( {0,\infty }\right) \\  u = g, & \left( {x,t}\right)  \in  {\mathbb{R}}^{n} \times  \{ 0\}  \end{matrix}\right.  \label{3.20}
  \end{equation}
\end{example}\begin{solution}
  在方程\ref{3.20}中对函数 \(u\) 仅关于空间变元作Fourier变换,则 \(u\) 的Fourier变换 \(\widehat{u}\) 满足
\begin{equation}
\left\{  \begin{array}{r} {\widehat{u}}_{t} + {\left| y\right| }^{2}\widehat{u} = 0,\;t > 0 \\  \widehat{u} = \widehat{g},\;t = 0 \end{array}\right.  \label{3.21}
\end{equation}

解得
\begin{equation*}
\widehat{u} = {e}^{-t{\left| y\right| }^{2}}\widehat{g}.
\end{equation*}

作逆Fourier变换得
\begin{equation*}
u\left( {x,t}\right)  = \frac{1}{\sqrt{{\left( 4\pi t\right) }^{n}}}{\int }_{{\mathbb{R}}^{n}}{e}^{-\frac{{\left| x - y\right| }^{2}}{4t}}g\left( y\right) {\dif y},\;x \in  {\mathbb{R}}^{n},t > 0.
\end{equation*}
\end{solution}

\begin{example}[波动方程的基本解] 考虑初值问题
\begin{equation}
\left\{  \begin{array}{ll} {u}_{tt} - \bigtriangleup u = 0, & \left( {x,t}\right)  \in  {\mathbb{R}}^{n} \times  \left( {0,\infty }\right) \\  u = g,{u}_{t} = 0, & \left( {x,t}\right)  \in  {\mathbb{R}}^{n} \times  \{ t = 0\}  \end{array}\right.  \label{3.22}
\end{equation}

其中 \(u\) 和 \(g\) 为复值函数.
\end{example}
\begin{solution}
    作 \(u\) 关于空间变元 \(x \in  {\mathbb{R}}^{n}\) 的Fourier变换,则有
\begin{equation}
\left\{  \begin{matrix} {\widehat{u}}_{tt} + {\left| y\right| }^{2}\widehat{u} = 0,\;t > 0 \\  \widehat{u} = \widehat{g},{\widehat{u}}_{t} = 0,\;t = 0 \end{matrix}\right.  \label{3.23}
\end{equation}

求解上述常微分方程得
\begin{equation*}
\widehat{u} = \frac{\widehat{g}}{2}\left( {{e}^{{it}\left| y\right| } + {e}^{-{it}\left| y\right| }}\right) .
\end{equation*}

再对 \(\widehat{u}\) 求逆Fourier变换可得\ref{3.22}的基本解为
\begin{equation*}
u\left( {x,t}\right)  = \frac{1}{\sqrt{{\left( 2\pi \right) }^{n}}}{\int }_{{\mathbb{R}}^{n}}\frac{\widehat{g}\left( y\right) }{2}\left( {{e}^{\ci\left( {x \cdot  y + t\left| y\right| }\right) } + {e}^{\ci\left( {x \cdot  y - t\left| y\right| }\right) }}\right) {\dif y},x \in  {\mathbb{R}}^{n},t \geq  0.
\end{equation*}

\end{solution}
\subsection{习 题}
\begin{problemset}
  \item 求下列函数的Fourier变换:\begin{enumerate}
  \item \(\sin {ax}\) ,
  \item  \({e}^{-a\left| x\right| }\left( {a > 0}\right)\) ,
  \item \(\frac{x}{{\left( {a}^{2} + {x}^{2}\right) }^{k}},\frac{1}{{\left( {a}^{2} + {x}^{2}\right) }^{k}}\) ( \(a > 0,k\) 为自然数).
  \item \(u\left( x\right)  = \left\{  \begin{matrix} 1, & \left| x\right|  < 1 \\  0, & \left| x\right|  > 1, \end{matrix}\right.\)
  \item  \(u\left( x\right)  = \left\{  \begin{array}{ll} 1 - \left| x\right| , & \left| x\right|  < 1 \\  0, & \left| x\right|  > 1. \end{array}\right.\)
\end{enumerate}
\begin{solution}
    1. \[
    \mathcal{F}[\sin ax]=\frac{1}{\sqrt{2\pi}}\int_{-\infty}^\infty e^{-\ci x\xi }\sin ax\dif x=\frac{1}{\sqrt{2\pi}}\int_{-\infty}^\infty  e^{-\ci x\xi }\frac{e^{\operatorname{i}ax}-e^{-\operatorname{i}ax}}{2\ci}\dif x=\frac{1}{\sqrt{2\pi}}2\pi \frac{1}{2\ci}(\delta(\xi-a)-\delta(\xi+a))
    \]
    2.
    \[
    \mathcal{F}[{e}^{-a\left| x\right| }]=\frac{1}{\sqrt{2\pi}}\int_{-\infty}^\infty e^{-\ci x\xi }{e}^{-a\left| x\right| }\dif x=\frac{1}{\sqrt{2\pi}}\int_{0}^\infty  e^{-(\ci \xi+a)x }+e^{(\ci \xi-a)x }\dif x=\frac{2a}{\sqrt{2\pi}(a^2+\xi^2)}
    \]
    3.


\end{solution}
\item 证明当 \(f\left( x\right)\) 在 \(\left( {-\infty , + \infty }\right)\) 内绝对可积时, \(\mathcal{F}\left\lbrack  f\right\rbrack\) 为连续函数.
\begin{proof}
    由Leabesgue控制收敛定理知
    \[\left|\mathcal{F}(f)(\xi_0)-\mathcal{F}(f)(\xi_0+\xi)\right|=\left|\frac{1}{\sqrt{2\pi}}\int_{-\infty}^\infty f(x)\left( e^{-\ci x\xi_0}- e^{-\ci x(\xi_0+\xi)}\right)\dif x \right|\to 0,\xi\to 0\]
    
\end{proof}
\item 证明下列公式\begin{enumerate}
  \item  \(\mathcal{F}\left\lbrack  {x}^{-{2m}}\right\rbrack   = \frac{{\left( -1\right) }^{m}\pi }{\left( {{2m} - 1}\right) !}{\left| \xi \right| }^{{2m} - 1},x \in  \mathbb{R},m = 1,2,\cdots\) .
  \item \(\mathcal{F}\left\lbrack  {{x}^{m}H\left( {a - \left| x\right| }\right) }\right\rbrack   = 2{\ci}^{m}\frac{{d}^{m}}{d{\xi }^{m}}\frac{\sin {a\xi }}{\xi },x \in  \mathbb{R},m = 1,2,\cdots\) .
  \item \(\mathcal{F}\left\lbrack  {\left( {x}^{2} + {a}^{2}\right) }^{m}\right\rbrack   = {2\pi }{\left( {a}^{2} - \frac{{d}^{2}}{d{\xi }^{2}}\right) }^{m}\delta \left( \xi \right) ,x \in  \mathbb{R},m = 1,2,\cdots\) .
  \item \(\mathcal{F}\left\lbrack  {\ln \left| x\right| }\right\rbrack   =  - \pi \left\lbrack  {\frac{1}{\left| \xi \right| } + {2\gamma \delta }\left( \xi \right) }\right\rbrack  ,x \in  \mathbb{R},\gamma\) 为Euler常数.
  \item \(\mathcal{F}\left\lbrack  {{e}^{iax}H\left( x\right) }\right\rbrack   = {\pi \delta }\left( {\xi  - a}\right)  + \frac{1}{\ci\left( {\xi  - a}\right) },x \in  \mathbb{R}\) .
\end{enumerate}\begin{proof}
    1.\[\mathcal{F}\left\lbrack  {x}^{-{2m}}\right\rbrack =\mathcal{F}\left\lbrack  \frac{(-1)^{2m-1}}{(2m-1)!}\frac{\dif^{2m-1}{x}^{-1}}{\dif x^{2m-1}}\right\rbrack    = \frac{{\left( -1\right) }^{2m-1}(\ci \xi)^{2m-1} }{\left( {{2m} - 1}\right) !}\mathcal{F}\left\lbrack\frac{1}{x}\right\rbrack=\frac{{\left( -1\right) }^{m} }{\left( {{2m} - 1}\right) !}{\left| \xi \right| }^{{2m} - 1}\sqrt{\frac{\pi}{2}}\]
    \begin{equation*}
        \begin{aligned}
        \mathcal{F}\left\lbrack \frac{1}{x}\right\rbrack=\frac{1}{\sqrt{2\pi}}\lim\limits_{\varepsilon\to 0^+}\int_{|x|>\varepsilon}\frac{e^{-\ci x\xi}}{x} \dif x =\frac{1}{\sqrt{2\pi}}\lim\limits_{\varepsilon\to 0^+}\int_{|x|>\varepsilon}\frac{-\ci\sin x\xi}{x} \dif x=-\ci \sqrt{\frac{\pi}{2}}\operatorname{sgn}(\xi)
        \end{aligned}
    \end{equation*}
    
\end{proof}
\item 设 \(f\left( x\right)  \in  {L}^{1}\left( {\mathbb{R}}^{n}\right)\) ,证明 \(f\left( x\right)\) 按 \({\mathcal{S}}^{\prime }\) 的Fourier变换与普通的Fourier变换一致.
\begin{proof}
    对 $\forall \varphi \in S\left(\mathbb{R}^n\right)$ ，
\[\left\langle\widetilde{\mathcal{F}}[f],\varphi\right\rangle=\left\langle f,\mathcal{F}[\varphi]\right\rangle=\int_{\mathbb{R}^n}f(\xi)\int_{\mathbb{R}^n}e^{-\ci x\xi}\varphi(x)\dif x\dif \xi=\int_{\mathbb{R}^n}\varphi(x)\int_{\mathbb{R}^n}e^{-\ci x\xi}f(\xi)\dif \xi\dif x=\left\langle \mathcal{F}[f],\varphi \right\rangle \]

从上式即知 $\widetilde{\mathcal{F}} f=\mathcal{F} f$ ．
\end{proof}
\item 计算下列广义函数的Fourier变换.

(1) \({PV}\left( \frac{1}{x}\right)\) ,

(2) \({\delta }^{\prime }\left( {x - a}\right)\) .
\begin{solution}
    \[\left\langle \mathcal{F
    }\left\lbrack {PV}\left( \frac{1}{x}\right)\right\rbrack ,\varphi(x)\right\rangle=\int_\mathbb{R}\varphi(\xi)\lim\limits_{\varepsilon\to 0}\int_{|x|>\varepsilon}\frac{e^{-\ci x\xi}}{x}\dif x\dif \xi =\int_\mathbb{R}\varphi(\xi)\left(-\ci \pi \operatorname{sgn}(\xi)\right)\dif \xi= \left\langle -\ci \pi \operatorname{sgn}(\xi),\varphi(x)\right\rangle\]

    \[
    \mathcal{F}\left\lbrack \delta'(x-a)\right\rbrack=\ci \xi \mathcal{F}\left\lbrack \delta(x-a)\right\rbrack=\ci \xi \frac{1}{\sqrt{2\pi}}e^{-\ci a \xi}
    \]
\end{solution}
\item 证明定理3.5的结论 \(\left( 2\right) ,\left( 3\right) ,\left( 4\right)\) .
\begin{proof}
    由引理3.3.1知\begin{equation*}
        \begin{aligned}
           {\partial }^{\alpha }  \mathcal{F}\left\lbrack  {f}\right\rbrack =  \mathcal{F}\left\lbrack  {(-\ci x)^\alpha f}\right\rbrack 
        \end{aligned}
    \end{equation*}
    \begin{equation*}
        \begin{aligned}
            \left\langle \mathcal{F}\left\lbrack \widetilde{\tau}_{x_0}f\right\rbrack,\varphi\right \rangle= \left\langle\widetilde{\tau}_{x_0}f, \mathcal{F}\left\lbrack\varphi\right\rbrack\right \rangle= \left\langle f, {\tau}_{-x_0}\mathcal{F}\left\lbrack\varphi\right\rbrack\right \rangle&= \int_{\mathbb{R}^n} f(x) \frac{1}{\sqrt{(2\pi)^n}}\int_{\mathbb{R}^n}\varphi(y)e^{-\ci y (x+x_0)}\dif y\dif x\\&=\int_{\mathbb{R}^n} \varphi(y) \frac{1}{\sqrt{(2\pi)^n}}\int_{\mathbb{R}^n}f(x) e^{-\ci y (x+x_0)}\dif x\dif y\\&=\left\langle e^{-\ci x_0\xi}\mathcal{F}\left\lbrack f\right\rbrack,\varphi\right \rangle
        \end{aligned}
    \end{equation*}
      \begin{equation*}
        \begin{aligned}
            \left\langle  \widetilde{\tau}_{x_0}\mathcal{F}\left\lbrack f\right\rbrack,\varphi\right \rangle= \left\langle \mathcal{F}\left\lbrack f\right\rbrack, \varphi(x+x_0)\right \rangle= \left\langle f,\mathcal{F} \left\lbrack \varphi(x+x_0)\right\rbrack\right \rangle &= \int_{\mathbb{R}^n} f(\xi) \frac{1}{\sqrt{(2\pi)^n}}\int_{\mathbb{R}^n}\varphi(x+x_0)e^{-\ci x\xi}\dif x\dif \xi\\&=\int_{\mathbb{R}^n} \varphi(x) \frac{1}{\sqrt{(2\pi)^n}}\int_{\mathbb{R}^n}f(\xi) e^{-\ci \xi(x-x_0)}\dif \xi\dif x \\&=\left\langle e^{x_0 x} \mathcal{F}\left\lbrack f\right\rbrack,\varphi\right \rangle
        \end{aligned}
    \end{equation*}
\end{proof}
\item (Parseval 公式) 若 \(u \in  {L}^{2}\left( {\mathbb{R}}^{n}\right)\) ,则 \(\widehat{u} \in  {L}^{2}\left( {\mathbb{R}}^{n}\right)\) 且 \(\parallel \widehat{u}{\parallel }_{{L}^{2}} = \parallel u{\parallel }_{{L}^{2}}\) .

\begin{proof}
    由于$C_0^\infty(\mathbb{R}^n)$在$L^2(\mathbb{R}^n)$中稠密,而上述公式对$C_0^\infty(\mathbb{R}^n)$成立,进而由Fourier变换及$\|\cdot\|_2$的连续性可知上述公式对$L^2(\mathbb{R}^n)$成立
\end{proof}
\item  若 \(f,g \in  {L}^{2}\left( {\mathbb{R}}^{n}\right)\) ,证明
\begin{equation*}
{\left( f,g\right) }_{{L}^{2}} = {\left( \mathcal{F}\left\lbrack  f\right\rbrack  ,\mathcal{F}\left\lbrack  g\right\rbrack  \right) }_{{L}^{2}},
\end{equation*}

其中 \({\left( \cdot , \cdot  \right) }_{{L}^{2}}\) 记 \({L}^{2}\left( {\mathbb{R}}^{n}\right)\) 中内积.

\begin{proof}
    由于Parseval公式成立,上式显然成立.
\end{proof}
\item 设 \(u,v \in  \mathcal{S}\left( {\mathbb{R}}^{n}\right)\) ,定义其卷积如下:

\begin{equation*}
\left( {u * v}\right) \left( x\right)  \mathrel{\text{ := }} {\int }_{{\mathbb{R}}^{n}}u\left( y\right) v\left( {x - y}\right) {\dif y}.
\end{equation*}

证明

(1) \(\mathcal{F}\left\lbrack  {u * v}\right\rbrack   = \sqrt{{\left( 2\pi \right) }^{n}}\mathcal{F}\left\lbrack  u\right\rbrack  \mathcal{F}\left\lbrack  v\right\rbrack\) ,

(2) \(\mathcal{F}\left\lbrack  {uv}\right\rbrack   = \frac{1}{\sqrt{{\left( 2\pi \right) }^{n}}}\mathcal{F}\left\lbrack  u\right\rbrack   * \mathcal{F}\left\lbrack  v\right\rbrack\) .
\begin{proof}
    (1)\begin{equation*}
        \begin{aligned}
\mathcal{F}\left\lbrack  {u * v}\right\rbrack&=\frac{1}{\sqrt{(2\pi)^n}}\int_{\mathbb{R}^n }e^{-\ci x\xi}\int_{\mathbb{R}^n}u(y)v(x-y)\dif y\dif x  \\
&=\frac{1}{\sqrt{(2\pi)^n}}\int_{\mathbb{R}^n }u(y)\int_{\mathbb{R}^n}e^{-\ci x\xi}v(x-y)\dif x\dif y  \\
&=\frac{1}{\sqrt{(2\pi)^n}}\int_{\mathbb{R}^n }u(y)\int_{\mathbb{R}^n}e^{-\ci( x+y)\xi}v(x)\dif x\dif y  
\\
&= \sqrt{{\left( 2\pi \right) }^{n}}\mathcal{F}\left\lbrack  u\right\rbrack  \mathcal{F}\left\lbrack  v\right\rbrack\
        \end{aligned}
    \end{equation*}
\end{proof}
\item 
\(\left( 1\right)\) 记 \(R\) 为 \({\mathbb{R}}^{n}\) 上的旋转变换, \({R}^{t}\) 为其转置. 若 \(\varphi  \in  \mathcal{S}\) ,证明
\begin{equation*}
\mathcal{F}\left\lbrack  {f \circ  R}\right\rbrack   = \mathcal{F}\left\lbrack  f\right\rbrack   \circ  {R}^{t}.
\end{equation*}

(2) 定义 \({D}_{\lambda } : {\mathbb{R}}^{n} \rightarrow  {\mathbb{R}}^{n}\) 为 \({D}_{\lambda }x = {\lambda x}\) . 若 \(\varphi  \in  \mathcal{S}\) ,证明

\begin{equation*}
\mathcal{F}\left\lbrack  {f \circ  {D}_{\lambda }}\right\rbrack   = {\lambda }^{-n}\mathcal{F}\left\lbrack  f\right\rbrack   \circ  {D}_{{\lambda }^{-1}}.
\end{equation*}

(3) 若 \(f \in  {\mathcal{S}}^{\prime }\) ,证明

\begin{equation*}
\mathcal{F}\left\lbrack  {f \circ  R}\right\rbrack   = \mathcal{F}\left\lbrack  f\right\rbrack   \circ  {R}^{t},\;\mathcal{F}\left\lbrack  {f \circ  {D}_{\lambda }}\right\rbrack   = {\lambda }^{-n}\mathcal{F}\left\lbrack  f\right\rbrack   \circ  {D}_{{\lambda }^{-1}}.
\end{equation*}
\item 设 \({H}^{m}\left( {\mathbb{R}}^{n}\right)  = \left\{  {u \in  {\mathcal{S}}^{\prime }\left( {\mathbb{R}}^{n}\right)  \mid  {\partial }^{\alpha }u \in  {L}^{2}\left( {\mathbb{R}}^{n}\right) ,\left| \alpha \right|  \leqslant  m}\right\}\) ,其中范数定义为

\begin{equation*}
\parallel u{\parallel }_{m} \mathrel{\text{ := }} {\left( \mathop{\sum }\limits_{{\left| \alpha \right|  \leqslant  m}}{\begin{Vmatrix}{\partial }^{\alpha }u\end{Vmatrix}}_{{L}^{2}}^{2}\right) }^{\frac{1}{2}}.
\end{equation*}

同时对 \(u \in  {H}^{m}\left( {\mathbb{R}}^{n}\right)\) ,定义
\begin{equation*}
\parallel u{\parallel }_{m}^{\prime } \mathrel{\text{ := }} {\left( {\int }_{{\mathbb{R}}^{n}}{\left( 1 + {\left| \xi \right| }^{2}\right) }^{m}\left| \mathcal{F}\left\lbrack  u\right\rbrack  \left( \xi \right) \right| {}^{2}\dif\xi \right) }^{\frac{1}{2}}.
\end{equation*}

证明

(1) \(\parallel u{\parallel }_{m}^{\prime } < \infty\) ;

(2) \(\parallel  \cdot  {\parallel }_{m}^{\prime }\) 是 \({H}^{m}\left( {\mathbb{R}}^{n}\right)\) 的等价范数;

(3) \({H}^{m}\left( {\mathbb{R}}^{n}\right)\) 是完备的.
\begin{proof}
    由Plancherel定理知\[
\|\partial^\alpha  u\|_{L^2}^2= \int_{\mathbb{R}^n} \xi ^{2\alpha} |\hat{u}|^2\dif \xi
    \]
存在正常数$c_1,c_2$使得
\[
c_1\left( 1 + {\left| \xi \right| }^{2}\right)^{m}\leqslant \sum\limits_{|\alpha|\leqslant m} x^{2\alpha}\leqslant c_2\left( 1 + {\left| \xi \right| }^{2}\right)^m
\]
由此可得(1)和(2).

(3)事实上，设 $\left\{u_n\right\}$ 是基本列，即
$$
\left\|u_{n+p}-u_n\right\|_m^{\prime} \rightarrow 0
$$


也就是
$$
\int_{\mathbb{R}^n}\left(1+|x|^2\right)^m\left|\hat{u}_{n+p}-\hat{u}_n\right|^2 \mathrm{~d} x \rightarrow 0 \quad(n \rightarrow \infty)
$$


对 $\forall p$ 一致成立．由此可见
$$
\left\{\left(1+|x|^2\right)^{\frac{m}{2}} \hat{u}_n(x)\right\}
$$


是 $L^2\left(\mathbb{R}^n\right)$ 的基本列，从而 $\left\{\hat{u}_n(x)\right\}$ 也为 $L^2\left(\mathbb{R}^n\right)$ 的基本列．因为 $L^2\left(\mathbb{R}^n\right)$完备，所以设
$$
\left\{\begin{array}{l}
\hat{u}_n \xrightarrow{L^2} v, \\
\left(1+|x|^2\right)^{\frac{m}{2}} \hat{u}_n(x) \xrightarrow{L^2} w,
\end{array}\right.
$$


则存在 a.e.收敛子列
$$
\left\{\hat{u}_{n_k}\right\}, \quad\left\{\left(1+|x|^2\right)^{\frac{m}{2}} \hat{u}_{n_{k_j}}(x)\right\},
$$


使得
$$
\hat{u}_{n_{k_j}} \xrightarrow{\text { a. e. }} v, \quad\left(1+|x|^2\right)^{\frac{m}{2}} \hat{u}_{n_{k_j}} \xrightarrow{\text { a. e. }} w=\left(1+|x|^2\right)^{\frac{m}{2}} v,
$$

故有
$$
\left(1+|x|^2\right)^{\frac{m}{2}} \hat{u}_n(x) \xrightarrow{L^2}\left(1+|x|^2\right)^{\frac{m}{2}} v .
$$

由此即知
$$
\begin{aligned}
& \left(1+|x|^2\right)^{\frac{m}{2}} v \in L^2\left(\mathbb{R}^n\right), \\
& \left(1+|x|^2\right)^{\frac{m}{2}} \hat{u}_n(x) \xrightarrow{L^2}\left(1+|x|^2\right)^{\frac{m}{2}} \hat{u},
\end{aligned}
$$

其中 $\hat{u}=v$ 。故有 $\left\|u_n-u\right\|_m^{\prime} \rightarrow 0(n \rightarrow \infty)$ ．于是 $\left(H^m\left(\mathbb{R}^n\right),\|u\|_m^{\prime}\right)$ 完备．
\end{proof}
\item 对任意非负实数 \(s\) ,设
\begin{equation*}
{H}^{s}\left( {\mathbb{R}}^{n}\right)  \mathrel{\text{ := }} \left\{  {u \in  {L}^{2}\left( {\mathbb{R}}^{n}\right)  \mid  {\left( 1 + {\left| \xi \right| }^{2}\right) }^{\frac{s}{2}}\mathcal{F}\left\lbrack  u\right\rbrack  \left( \xi \right)  \in  {L}^{2}\left( {\mathbb{R}}^{n}\right) }\right\}  ,
\end{equation*}

其中范数定义为
\begin{equation*}
\parallel u{\parallel }_{s} \mathrel{\text{ := }} \parallel {\left( 1 + {\left| \xi \right| }^{2}\right) }^{\frac{s}{2}}\mathcal{F}\left\lbrack  u\right\rbrack  \left( \xi \right) {\parallel }_{{L}^{2}}.
\end{equation*}

证明

(1) 当 \(s = m \in  \mathbb{N}\) 时, \({H}^{s}\left( {\mathbb{R}}^{n}\right)\) 与上题中 \({H}^{m}\left( {\mathbb{R}}^{n}\right)\) 的定义等价;

(2) \({H}^{s}\left( {\mathbb{R}}^{n}\right)\) 中可引进内积 \(\left( {\cdot , \cdot  }\right)\) 使得 \(\parallel u{\parallel }_{s} = {\left( u,u\right) }^{\frac{1}{2}}\) ;

(3) 对 \(u \in  {H}^{s}{\left( {\mathbb{R}}^{n}\right) }^{\prime }\) ,存在 \(\widetilde{u} \in  {L}_{loc}^{1}\left( {\mathbb{R}}^{n}\right)\) 使得
\begin{equation*}
\widetilde{u}\left( \xi \right) {\left( 1 + {\left| \xi \right| }^{2}\right) }^{-\frac{s}{2}} \in  {L}^{2}\left( {\mathbb{R}}^{n}\right) ,
\end{equation*}

且
\begin{equation*}
\langle u,\mathcal{F}\left\lbrack  \varphi \right\rbrack  \rangle  = {\int }_{{\mathbb{R}}^{n}}\varphi  \cdot  \widetilde{u}\left( \xi \right) {\dif\xi },\;\forall \varphi  \in  \mathcal{S}\left( {\mathbb{R}}^{n}\right) .
\end{equation*}

\item 证明方程 \({\Delta f} = f\) 在 \({\mathcal{S}}\left( {\mathbb{R}}^{n}\right)\) 中无非零解.
\begin{proof}
    由Fourier变化性质知\[
    \mathcal{F}[\partial^2_{x_i}f ]=(\ci\xi_i )^2\hat{f}
    \]
    那么由\[
    \Delta f=f\Rightarrow -|\xi|^2\hat{f}=\hat{f}\Rightarrow \hat{f}=0\]
    所以 \({\Delta f} = f\) 在 \({\mathcal{S}}\left( {\mathbb{R}}^{n}\right)\) 中无非零解.
\end{proof}
\item 记 \(f\left( x\right)  = \mathop{\sum }\limits_{{\nu  =  - \infty }}^{\infty }{e}^{\ci\nu x} = 1 + {g}^{\prime \prime }\left( x\right) ,g\left( x\right)  = \mathop{\sum }\limits_{{\nu  \neq  0}}\frac{-1}{{\nu }^{2}}{e}^{\ci\nu x} = \frac{{\pi }^{2}}{6} - \frac{1}{2}{\left( x - \pi \right) }^{2}\) ,

\(0 < x < {2\pi }\) . 证明
\begin{equation*}
{g}^{\prime \prime }\left( x\right)  =  - 1 + \mathop{\sum }\limits_{{\nu  =  - \infty }}^{\infty }{2\pi \delta }\left( {x - {2\pi \nu }}\right) \text{ (在 }{\mathcal{D}}^{\prime }\left( \mathbb{R}\right) \text{ 中), }
\end{equation*}

由此得出 \(f\left( x\right)  = \mathop{\sum }\limits_{\nu }{2\pi \delta }\left( {x - {2\pi \nu }}\right)\) . 利用这个结果证明Poisson求和公式
\begin{equation*}
\mathop{\sum }\limits_{{\nu  =  - \infty }}^{\infty }\varphi \left( \nu \right)  = \mathop{\sum }\limits_{{\nu  =  - \infty }}^{\infty }\widehat{\varphi }\left( {2\pi \nu }\right) ,\;\forall \varphi  \in  \mathcal{S}\left( \mathbb{R}\right) .
\end{equation*}


\end{problemset}


\section{Sobolev空间与嵌入定理}

\subsection{Sobolev空间概念}

设 \(\Omega  \subset  {\mathbb{R}}^{n}\) 为带光滑边界的区域, \(k\) 是非负整数, \(1 \leqslant  p < \infty\) . 定义光滑函数 \(u : \Omega  \rightarrow  \mathbb{R}\) 的 \({W}^{k,p}\) -范数为
\begin{equation}
\parallel u{\parallel }_{k,p} \mathrel{\text{ := }} {\left( \mathop{\sum }\limits_{{\left| \alpha \right|  \leqslant  k}}{\int }_{\Omega }{\left| {\partial }^{\alpha }u\left( x\right) \right| }^{p}\dif x\right) }^{\frac{1}{p}}, \label{3.24}
\end{equation}

其中 \(\alpha  = \left( {{\alpha }_{1},\cdots ,{\alpha }_{n}}\right)\) 为多重指标且 \(\left| \alpha \right|  = {\alpha }_{1} + \cdots  + {\alpha }_{n}\) .
\begin{definition}
    称集合
\begin{equation*}
{W}^{k,p}\left( \Omega \right)  = \left\{  {u \in  {L}^{p}\left( \Omega \right)  \mid  {\partial }^{\alpha }u \in  {L}^{p}\left( \Omega \right) ,\left| \alpha \right|  \leqslant  k}\right\}
\end{equation*}

按范数\ref{3.24}构成的空间为Sobolev空间,记为 \({W}^{k,p}\left( \Omega \right)\) .
\end{definition}
\begin{theorem}
     空间 \({W}^{k,p}\left( \Omega \right)\) 是完备的.
\end{theorem}

\begin{proof}
    设 \({\left\{  {u}_{j}\right\}  }_{1}^{\infty }\) 是 \({W}^{k,p}\left( \Omega \right)\) 中的基本列,那么 \(\left\{  {{\partial }^{\alpha }{u}_{j}}\right\}\) 是 \({L}^{p}\left( \Omega \right)\) 中的基本列. 对任意 \(\alpha \left( {\left| \alpha \right|  \leqslant  k}\right)\) ,由 \({L}^{p}\left( \Omega \right)\) 的完备性知,存在 \({g}_{\alpha } \in  {L}^{p}\left( \Omega \right)\) 使得对 \(\left| \alpha \right|  \leqslant  k\) 有
\begin{equation*}
{\begin{Vmatrix}{\partial }^{\alpha }{u}_{j} - {g}_{\alpha }\end{Vmatrix}}_{{L}^{p}\left( \Omega \right) } \rightarrow  0,\;j \rightarrow  \infty .
\end{equation*}

从而对任意 \(\varphi  \in  \mathcal{D}\left( \Omega \right)\) ,当 \(j \rightarrow  \infty\) 时,
\begin{equation*}
\left\langle  {{\partial }^{\alpha }{u}_{j},\varphi }\right\rangle\to 
\left\langle  {{g}_{\alpha },\varphi }\right\rangle  \;   {\left( -1\right) }^{\left| \alpha \right| }\left\langle  {{u}_{j},{\partial }^{\alpha }\varphi }\right\rangle\to{\left( -1\right) }^{\left| \alpha \right| }\left\langle  {{g}_{0},{\partial }^{\alpha }\varphi }\right\rangle  
\end{equation*}

即得 \({g}_{\alpha } = {\partial }^{\alpha }{g}_{0}\) . 由此可得 \({g}_{0} \in  {W}^{k,p}\left( \Omega \right)\) ,且
\begin{equation*}
{\begin{Vmatrix}{u}_{j} - {g}_{0}\end{Vmatrix}}_{k,p} \rightarrow  0,\;j \rightarrow  \infty .
\end{equation*}

\end{proof}
\begin{proposition}
    [Young不等式]
    设$f\in L^p(\mathbb{R}),1\leqslant p\leqslant \infty$,$K\in L^1(\mathbb{R})$,那么有
    \begin{equation*}
        \|K*f\|_p\leqslant \|K\|_1\cdot \|f\|_p
    \end{equation*}
\end{proposition}
\begin{theorem}
    \({C}_{0}^{\infty }\left( {\mathbb{R}}^{n}\right)\) 在 \({W}^{k,p}\left( {\mathbb{R}}^{n}\right)\) 中稠密.
\end{theorem}
\begin{proof}
    首先,证明 \({W}^{k,p}\left( {\mathbb{R}}^{n}\right)\) 中函数可以用一个具紧支集的 \({W}^{k,p}\left( {\mathbb{R}}^{n}\right)\) 的函数列来逼近.

取 \(g \in  {C}_{0}^{\infty }\left( {\mathbb{R}}^{n}\right)\) 满足

(i)
\begin{equation*}
g\left( x\right)  = \left\{  \begin{array}{ll} 1, & x \in  B\left( {0,1}\right) \\  0, & x \notin  B\left( {0,2}\right)  \end{array}\right.
\end{equation*}

(ii) 对任意 \(x \in  {\mathbb{R}}^{n}\) 及 \(0 \leqslant  \alpha  \leqslant  k,\left| {{\partial }^{\alpha }g\left( x\right) }\right|  \leqslant  c\) (常数).

取 \({g}_{m}\left( x\right)  = g\left( \frac{x}{m}\right) \left( {m > 1}\right)\) . 则 \({g}_{m} \in  {C}_{0}^{\infty }\left( {\mathbb{R}}^{n}\right)\) 满足
\begin{equation*}
{g}_{m}\left( x\right)  = \left\{  \begin{array}{ll} 1, & x \in  B\left( {0,m}\right) \\  0, & x \notin  B\left( {0,{2m}}\right)  \end{array}\right.
\end{equation*}

且
\begin{equation*}
\left| {{\partial }^{\alpha }{g}_{m}\left( x\right) }\right|  \leqslant  c\frac{1}{{m}^{\left| \alpha \right| }}.
\end{equation*}

设 \(u \in  {W}^{k,p}\left( {\mathbb{R}}^{n}\right)\) . 令 \({u}_{m} = {g}_{m}u\) ,则 \(\operatorname{supp}{u}_{m} \subset  B\left( {0,{2m}}\right)\) 且在 \(B\left( {0,m}\right)\) 上 \({u}_{m} = u\) .

下证 \(\left\{  {u}_{m}\right\}\) 在 \({W}^{k,p}\left( {\mathbb{R}}^{n}\right)\) 中收敛到 \(u\) ,即对任意 \(\left| \alpha \right|  \leqslant  k,{\partial }^{\alpha }{u}_{m}\overset{{L}^{p}\left( {\mathbb{R}}^{n}\right) }{ \longrightarrow  }{\partial }^{\alpha }u\) .

若 \(\alpha  = 0\) ,则有
\begin{equation*}
{\begin{Vmatrix}{u}_{m} - u\end{Vmatrix}}_{{L}^{p}\left( {\mathbb{R}}^{n}\right) }^{p} = {\int }_{{\mathbb{R}}^{n}}{\left| {u}_{m} - u\right| }^{p}{\dif x} = {\int }_{\left| x\right|  > m}{\left| {u}_{m} - u\right| }^{p}{\dif x}
\leqslant  {\left( c + 1\right) }^{p}{\int }_{\left| x\right|  > m}{\left| u\left( x\right) \right| }^{p}{\dif x} \rightarrow  0,\;\left( {m \rightarrow  \infty }\right) .
\end{equation*}

若 \(\left| \alpha \right|  = 1\) ,则对 \(1 \leqslant  j \leqslant  n\) 有 \({\partial }_{j}{u}_{m} = {g}_{m}\left( {{\partial }_{j}u}\right)  + \left( {{\partial }_{j}{g}_{m}}\right) u\) . 如上可知
\begin{equation*}
{g}_{m}\left( {{\partial }_{j}u}\right) \overset{{L}^{p}\left( {\mathbb{R}}^{n}\right) }{ \longrightarrow  }{\partial }_{j}u,\;\left( {m \rightarrow  \infty }\right) .
\end{equation*}
又
\begin{equation*}
{\begin{Vmatrix}\left( {\partial }_{j}{g}_{m}\right) u\end{Vmatrix}}_{{L}^{p}\left( {\mathbb{R}}^{n}\right) }^{p} = {\int }_{{\mathbb{R}}^{n}}{\left| \left( {\partial }_{j}{g}_{m}\right) u\right| }^{p}{\dif x}={\int }_{m < \left| x\right|  < {2m}}{\left| \left( {\partial }_{j}{g}_{m}\right) u\right| }^{p}{\dif x}
\leqslant  {c}^{p}{\int }_{m < \left| x\right|  < {2m}}{\left| u\left( x\right) \right| }^{p}{\dif x} \rightarrow  0,\;\left( {m \rightarrow  \infty }\right) .
\end{equation*}

故当 \(m \rightarrow  \infty\) 时有
\begin{equation*}
{\partial }_{j}{u}_{m}\overset{{L}^{p}\left( {\mathbb{R}}^{n}\right) }{ \longrightarrow  }{\partial }_{j}u
\end{equation*}

用数学归纳法即可证,对任意 \(\left| \alpha \right|  \leqslant  k,{\partial }^{\alpha }{u}_{m}\overset{{L}^{p}\left( {\mathbb{R}}^{n}\right) }{ \longrightarrow  }{\partial }^{\alpha }u\) .

其次,再证 \({W}^{k,p}\left( {\mathbb{R}}^{n}\right)\) 中每个具紧支集的函数可用 \({C}_{0}^{\infty }\left( {\mathbb{R}}^{n}\right)\) 中函数列来逼近.

设 \(u \in  {W}^{k,p}\left( {\mathbb{R}}^{n}\right)\) 且 \(\operatorname{supp}u\) 为紧集,即对 \(\left| \alpha \right|  \leqslant  k\) 有 \({\partial }^{\alpha }u \in  {L}^{p}\left( {\mathbb{R}}^{n}\right)\) . 对任意 \(\delta  >\) 0,令
\begin{equation}
{u}_{\delta }\left( x\right)  = {\int }_{{\mathbb{R}}^{n}}u\left( y\right) {j}_{\delta }\left( {x - y}\right) {\dif y} \label{3.25}
\end{equation}

其中 \({j}_{\delta }\left( {x - y}\right)\) 是例3.2中定义的磨光核. 从而有 \({u}_{\delta } \in  {C}_{0}^{\infty }\left( {\mathbb{R}}^{n}\right)\) 且
\begin{equation*}
{\partial }_{x}^{\alpha }{u}_{\delta }\left( x\right)  = {\int }_{{\mathbb{R}}^{n}}u\left( y\right) {\partial }_{x}^{\alpha }{j}_{\delta }\left( {x - y}\right) {\dif y}
= {\left( -1\right) }^{\left| \alpha \right| }{\int }_{{\mathbb{R}}^{n}}u\left( y\right) {\partial }_{y}^{\alpha }{j}_{\delta }\left( {x - y}\right) {\dif y}
= {\int }_{{\mathbb{R}}^{n}}{\partial }^{\alpha }u\left( y\right) {j}_{\delta }\left( {x - y}\right) {\dif y} = {\left( {\partial }^{\alpha }u\right) }_{\delta }.
\end{equation*}

即 \({\partial }^{\alpha }{u}_{\delta }\) 是 \({\partial }^{\alpha }u\) 的光滑化. 进一步,利用Young不等式可证对任意 \(v \in  {L}^{p}\left( {\mathbb{R}}^{n}\right)\) 有 \({\begin{Vmatrix}{\left( v\right) }_{\delta }\end{Vmatrix}}_{{L}^{p}} \leqslant  \parallel v{\parallel }_{{L}^{p}}\) . 由于 \({L}^{p}\left( {\mathbb{R}}^{n}\right)\) 中函数可以由 \({C}_{0}^{\infty }\left( {\mathbb{R}}^{n}\right)\) 中函数任意逼近,因此,对 \({\partial }^{\alpha }u \in  {L}^{p}\left( {\mathbb{R}}^{n}\right)\) 及任意 \(\varepsilon  > 0\) ,都存在 \({v}_{\alpha } \in  {C}_{0}^{\infty }\left( {\mathbb{R}}^{n}\right)\) 使得 \({\begin{Vmatrix}{v}_{\alpha } - {\partial }^{\alpha }u\end{Vmatrix}}_{{L}^{p}} < \frac{\varepsilon }{3}\) . 从而由Young不等式有
\begin{equation*}
{\begin{Vmatrix}{\left( {v}_{\alpha }\right) }_{\delta } - {\left( {\partial }^{\alpha }u\right) }_{\delta }\end{Vmatrix}}_{{L}^{p}} < \frac{\varepsilon }{3}.
\end{equation*}

再由定理 3.1 后注知,存在 \({\delta }_{0} = {\delta }_{0}\left( {\varepsilon ,\alpha }\right)  > 0\) 使得当 \(0 < \delta  < {\delta }_{0}\) 时有
\begin{equation*}
{\begin{Vmatrix}{\left( {v}_{\alpha }\right) }_{\delta } - {v}_{\alpha }\end{Vmatrix}}_{{L}^{p}} < \frac{\varepsilon }{3}.
\end{equation*}

从而 \({\begin{Vmatrix}{\left( {\partial }^{\alpha }u\right) }_{\delta } - {\partial }^{\alpha }u\end{Vmatrix}}_{{L}^{p}} < \varepsilon\) . 即 \({\begin{Vmatrix}{\partial }^{\alpha }{u}_{\delta } - {\partial }^{\alpha }u\end{Vmatrix}}_{{L}^{p}} < \varepsilon \left( {\left| \alpha \right|  < k}\right)\) . 亦即
\begin{equation*}
{\partial }^{\alpha }{u}_{\delta }\overset{{L}^{p}\left( {\mathbb{R}}^{n}\right) }{ \longrightarrow  }{\partial }^{\alpha }u
\end{equation*}
\end{proof}
\begin{note}
    如果用 \({W}_{0}^{k,p}\left( {\mathbb{R}}^{n}\right)\) 表示 \({C}_{0}^{\infty }\left( {\mathbb{R}}^{n}\right)\) 按 \({W}^{k,p}\) -范数\ref{3.24}完备化产生的空间, 则由该定理和定理3.4.1可知
\begin{equation*}
{W}_{0}^{k,p}\left( {\mathbb{R}}^{n}\right)  = {W}^{k,p}\left( {\mathbb{R}}^{n}\right) .
\end{equation*}

但对一般的开集 \(\Omega ,{W}_{0}^{k,p}\left( \Omega \right)  = {W}^{k,p}\left( \Omega \right)\) 未必成立. 事实上,当 \(\Omega\) 为有界开集且 \(k >\) 0时, \({W}_{0}^{k,p}\left( \Omega \right)\) 是 \({W}^{k,p}\left( \Omega \right)\) 的真子集,因为在 \(\bar{\Omega }\) 上恒等于 1 的函数属于 \({W}^{k,p}\left( \Omega \right)\) ,但它却不属于 \({W}_{0}^{1,1}\left( \Omega \right)\) ,从而也不属于任一 \({W}_{0}^{k,p}\left( \Omega \right)\) ,参见习题3.4.3.
\end{note}

至此,我们自然会问另外一个问题: 如果我们考虑 \({C}^{\infty }\left( \Omega \right)\) 在范数\ref{3.24}下的完备化,情况又如何?亦即,若记 \({H}^{k,p}\left( \Omega \right)\) 为集合
\begin{equation*}
M \mathrel{\text{ := }} \left\{  {u \in  {C}^{\infty }\left( \Omega \right)  \mid  \parallel u{\parallel }_{k,p} < \infty }\right\}
\end{equation*}

在范数\ref{3.24}下的完备化,则 \({H}^{k,p}\left( \Omega \right)\) 与 \({W}^{k,p}\left( \Omega \right)\) 有何关系? 对此,我们有
\begin{theorem}
    若 \(1 \leqslant  p < \infty\) ,则 \({H}^{k,p}\left( \Omega \right)  = {W}^{k,p}\left( \Omega \right)\) .
\end{theorem}
\begin{proof}
     由定理 3.4.1知, \({H}^{k,p}\left( \Omega \right)  \subset  {W}^{k,p}\left( \Omega \right)\) . 我们只需证集合 \(M\) 在 \({W}^{k,p}\left( \Omega \right)\) 中稠密即可. 记
\begin{equation*}
{\Omega }_{m} \mathrel{\text{ := }} \left\{  {x \in  \Omega \left| \right| x \mid   < m\text{ 且 }\operatorname{dist}\left( {x,\partial \Omega }\right)  > \frac{1}{m}}\right\}  ,\;m = 1,2,\cdots ,
\end{equation*}

并记 \({\Omega }_{0} = {\Omega }_{-1} = \varnothing\) . 那末
\begin{equation*}
\mathcal{U} \mathrel{\text{ := }} \left\{  {{U}_{m} \mid  {U}_{m} = {\Omega }_{m + 1}\bigcap {\bar{\Omega }}_{m - 1}^{c},m = 1,2,\cdots }\right\}
\end{equation*}

是 \(\Omega\) 的一族开覆盖. 由单位分解定理知,存在从属于开覆盖 \(\mathcal{U}\) 的一族 \({C}^{\infty }\) 函数 \(\mathcal{F}\) . 由于 \({\bar{U}}_{m}\) 为紧集,故由单位分解定理的结论 (1) 知, \(\mathcal{F}\) 中只有有限多个函数 \(\psi\) ,使得 \(\operatorname{supp}\psi  \subset  {U}_{m}\) . 记 \({\varphi }_{m}\) 为这些函数之和,则有 \({\varphi }_{m} \in  {C}_{0}^{\infty }\left( {U}_{m}\right)\) 而且
\begin{equation*}
\mathop{\sum }\limits_{{m = 1}}^{\infty }{\varphi }_{m}\left( x\right)  \equiv  1,\;\forall x \in  \Omega .
\end{equation*}

设 \(u \in  {W}^{k,p}\left( \Omega \right)\) ,下证对任意 \(\varepsilon  > 0\) ,存在 \(\varphi  \in  M\) 使得 \(\parallel u - \varphi {\parallel }_{k,p} \leqslant  \varepsilon\) . 由定义式 (3.25)可知,当 \(0 < \delta  < \frac{1}{\left( {m + 1}\right) \left( {m + 2}\right) }\) 时,有
\begin{equation*}
\operatorname{supp}{\left( {\varphi }_{m}u\right) }_{\delta } \subset  {\Omega }_{m + 2}\bigcap {\bar{\Omega }}_{m - 2}^{c}.
\end{equation*}

由定理3.7的证明方法可知,存在 \({\delta }_{m} \in  \left( {0,\frac{1}{\left( {m + 1}\right) \left( {m + 2}\right) }}\right)\) 满足
\begin{equation*}
{\begin{Vmatrix}{\varphi }_{m}u - {\left( {\varphi }_{m}u\right) }_{{\delta }_{m}}\end{Vmatrix}}_{k,p} < \frac{\varepsilon }{{2}^{m}}.
\end{equation*}

令 \(\varphi  = \mathop{\sum }\limits_{{m = 1}}^{\infty }{\left( {\varphi }_{m}u\right) }_{{\delta }_{m}}\) ,易知 \(\varphi  \in  {C}^{\infty }\left( \Omega \right)\) . 同时对任意正整数 \(m\) ,在 \({\Omega }_{m}\) 上成立有
\begin{equation*}
u\left( x\right)  = \mathop{\sum }\limits_{{j = 1}}^{{m + 2}}{\varphi }_{j}\left( x\right) u\left( x\right) ,
\end{equation*}
\begin{equation*}
\varphi \left( x\right)  = \mathop{\sum }\limits_{{j = 1}}^{{m + 2}}{\left( {\varphi }_{j}u\right) }_{{\delta }_{j}}\left( x\right) .
\end{equation*}

因此,
\begin{equation*}
\parallel u - \varphi {\parallel }_{{W}^{k,p}\left( {\Omega }_{m}\right) } \leqslant  \mathop{\sum }\limits_{{j = 1}}^{{m + 2}}{\begin{Vmatrix}{\left( {\varphi }_{j}u\right) }_{{\delta }_{j}} - {\varphi }_{j}u\end{Vmatrix}}_{k,p} < \varepsilon .
\end{equation*}

令 \(m \rightarrow  \infty\) ,即可得 \(\parallel u - \varphi {\parallel }_{k,p} \leqslant  \varepsilon\) . 故定理得证.

\end{proof}

\subsection{Sobolev不等式}

我们应该注意到具有弱导数的函数未必连续. 例如,考虑 \(\Omega  = {B}_{1} = \left\{  {x \in  {\mathbb{R}}^{n} \mid  }\right.\)  \(\left| x\right|  < 1\}\) 上函数
\begin{equation*}
u\left( x\right)  = {\left| x\right| }^{-\alpha },
\end{equation*}

其中 \(\alpha  \in  \mathbb{R}\) . 则
\begin{equation*}
\frac{\partial u}{\partial {x}_{j}} =  - \alpha {x}_{j}{\left| x\right| }^{-\alpha  - 2}.
\end{equation*}

由归纳法知,
\begin{equation*}
\left| {{\partial }^{\nu }u\left( x\right) }\right|  \leqslant  {c}_{\nu }{\left| x\right| }^{-\alpha  - \left| \nu \right| }.
\end{equation*}

由于函数 \({\left| x\right| }^{-\beta }\) 在 \({B}_{1}\) 上可积当且仅当 \(\beta  < n\) . 因此仅当
\begin{equation*}
{\alpha p} + {kp} < n
\end{equation*}

时才有 \(u \in  {W}^{k,p}\left( {B}_{1}\right)\) . 因此,若 \({kp} < n\) ,则可选择 \(0 < \alpha  < n/p - k\) 使得 \(u \in\)  \({W}^{k,p}\left( {B}_{1}\right)\) 但在原点不连续. 我们自然会问,当 \({kp} > n\) 时,情况又如何呢? 由于 \({W}^{k,p}\left( \Omega \right)\) 中函数 \(f\) 的弱导数也属于 \({L}^{p}\left( \Omega \right)\) ,我们自然希望,当 \(f\) 的充分高阶弱导数属于 \({L}^{p}\left( \Omega \right)\) 时, \(f\) 应为连续或适当光滑的函数. 这恰好是Sobolev嵌入定理的结论. 但在叙述和证明Sobolev嵌入定理之前, 我们先证明一些Sobolev 不等式.
\begin{lemma}\label{3.4.2}
     设 \(\Omega  \subset  {\mathbb{R}}^{n}\) 为有界区域, \(1 \leqslant  p < n\) . 则存在常数 \(C = C\left( {p,n}\right)\) 使得对任意 \(u \in  {W}_{0}^{1,p}\left( \Omega \right)\) 都有 \(u \in  {L}^{\frac{np}{n - p}}\left( \Omega \right)\) 且
\begin{equation}
\parallel u{\parallel }_{{L}^{\frac{np}{n - p}}} \leqslant  C\parallel \nabla u{\parallel }_{{L}^{p}\left( \Omega \right) }. \label{3.26}
\end{equation}

\end{lemma}
\begin{proof}
  首先假定 \(p = 1\) 且 \(u \in  {C}_{0}^{1}\left( \Omega \right)\) . 记 \(x = \left( {{x}_{1},\cdots ,{x}_{n}}\right)\) . 此时,对每个 \(i \in  \{ 1,2,\cdots ,n\}\) 有
\begin{equation*}
\left| {u\left( x\right) }\right|  \leqslant  {\int }_{-\infty }^{{x}_{i}}\left| {{\nabla }_{i}u\left( x\right) }\right| \dif x_i,
\end{equation*}

这里我们用到\(u\) 具有紧支集. 从而
\begin{equation*}
{\left| u\left( x\right) \right| }^{n} \leqslant  \mathop{\prod }\limits_{{i = 1}}^{n}{\int }_{-\infty }^{\infty }\left| {{\nabla }_{i}u}\right| \dif x_i
\end{equation*}

且
\begin{equation*}
{\left| u\left( x\right) \right| }^{\frac{n}{n - 1}} \leqslant  {\left( \mathop{\prod }\limits_{{i = 1}}^{n}{\int }_{-\infty }^{\infty }\left| {\nabla }_{i}u\right| \dif x_i\right) }^{\frac{1}{n - 1}}.
\end{equation*}

不等式两边关于第一个变量积分得
\begin{equation*}
{\int }_{-\infty }^{\infty }{\left| u\left( x\right) \right| }^{\frac{n}{n - 1}}\dif x_1 \leqslant  {\left( {\int }_{-\infty }^{\infty }\left| {\nabla }_{1}u\right| \dif x_1\right) }^{\frac{1}{n - 1}}{\int }_{-\infty }^{\infty }{\left( \mathop{\prod }\limits_{{i = 2}}^{n}{\int }_{-\infty }^{\infty }\left| {\nabla }_{i}u\right| \dif x_i\right) }^{\frac{1}{n - 1}}\dif x_1.
\end{equation*}

因此,由 \({p}_{1} = \cdots  = {p}_{n - 1} = n - 1\) 的广义Hölder不等式得
\begin{equation*}
{\int }_{-\infty }^{\infty }{\left| u\left( x\right) \right| }^{\frac{n}{n - 1}}\dif x_1 \leqslant  {\left( {\int }_{-\infty }^{\infty }\left| {\nabla }_{1}u\right| \dif x_1\right) }^{\frac{1}{n - 1}}{\left( \mathop{\prod }\limits_{{i \neq  1}}{\int }_{-\infty }^{\infty }{\int }_{-\infty }^{\infty }\left| {\nabla }_{i}u\right| \dif x_i\dif x_1\right) }^{\frac{1}{n - 1}}.
\end{equation*}

同理可得
\begin{equation*}
{\int }_{-\infty }^{\infty }{\int }_{-\infty }^{\infty }{\left| u\left( x\right) \right| }^{\frac{n}{n - 1}}\dif x_1\dif x_2 \leqslant  {\left( {\int }_{-\infty }^{\infty }{\int }_{-\infty }^{\infty }\left| {\nabla }_{1}u\right| \dif x_1\dif x_2\right) }^{\frac{1}{n - 1}}
\end{equation*}
\begin{equation*}
\cdot  {\left( {\int }_{-\infty }^{\infty }{\int }_{-\infty }^{\infty }\left| {\nabla }_{2}u\right| \dif x_1\dif x_2\right) }^{\frac{1}{n - 1}}
\end{equation*}
\begin{equation*}
\cdot  {\left( \mathop{\prod }\limits_{{i = 3}}^{n}{\int }_{-\infty }^{\infty }{\int }_{-\infty }^{\infty }{\int }_{-\infty }^{\infty }\left| {\nabla }_{i}u\right| \dif x_i\dif x_1\dif x_2\right) }^{\frac{1}{n - 1}}.
\end{equation*}

依次迭代可得
\begin{equation*}
{\int }_{\Omega }{\left| u\left( x\right) \right| }^{\frac{n}{n - 1}}{\dif x} \leqslant  {\left( \mathop{\prod }\limits_{{i = 1}}^{n}{\int }_{\Omega }\left| {\nabla }_{i}u\right| \dif x\right) }^{\frac{1}{n - 1}}.
\end{equation*}

进而有
\begin{equation*}
\parallel u{\parallel }_{\frac{n}{n - 1}} \leqslant  {\left( \mathop{\prod }\limits_{{i = 1}}^{n}{\int }_{\Omega }\left| {\nabla }_{i}u\right| \dif x\right) }^{\frac{1}{n}} \leqslant  \frac{1}{n}{\int }_{\Omega }\mathop{\sum }\limits_{{i = 1}}^{n}\left| {{\nabla }_{i}u}\right| {\dif x},
\end{equation*}

这里我们使用了几何平均值不大于算术平均值. 因此,由 \({\mathbb{R}}^{n}\) 上范数等价知,存在常数 \({C}^{\prime } = {C}^{\prime }\left( n\right)\) 使得
\begin{equation}
\parallel u{\parallel }_{\frac{n}{n - 1}} \leqslant  \frac{{C}^{\prime }}{n}\parallel \nabla u{\parallel }_{1}. \label{3.27}
\end{equation}

至此,我们证明了 \(p = 1\) 时引理结论对 \(u \in  {C}_{0}^{1}\left( \Omega \right)\) 成立.

对 \(\gamma  > 1\) ,我们应用\ref{3.27}到 \({\left| u\right| }^{\gamma }\) ,由链式法则可得
\begin{equation}
{\begin{Vmatrix}{\left| u\right| }^{\gamma }\end{Vmatrix}}_{\frac{n}{n - 1}} \leqslant  \frac{{C}^{\prime }\gamma }{n}{\int }_{\Omega }{\left| u\right| }^{\gamma  - 1}\left| {\nabla u}\right| {\dif x}
\leqslant  \frac{{C}^{\prime }\gamma }{n}{\begin{Vmatrix}{\left| u\right| }^{\gamma  - 1}\end{Vmatrix}}_{q} \cdot  \parallel \nabla u{\parallel }_{p}, \label{3.28}
\end{equation}

其中 \(\frac{1}{p} + \frac{1}{q} = 1\) ,在第二个不等式中使用了Hölder不等式.
由于 \(p < n\) ,取 \(\gamma  = \frac{\left( {n - 1}\right) p}{n - p}\) ,则有 \(\frac{np}{n - p} = \frac{\gamma n}{n - 1} = \frac{\left( {\gamma  - 1}\right) p}{p - 1}\) . 令 \(q = \frac{p}{p - 1}\) ,则由

可得
\begin{equation*}
\parallel u{\parallel }_{\frac{\gamma n}{n - 1}}^{\gamma } \leqslant  \frac{{C}^{\prime }\gamma }{n}\parallel u{\parallel }_{\frac{\gamma n}{n - 1}}^{\gamma  - 1} \cdot  \parallel \nabla u{\parallel }_{p}.
\end{equation*}

因此,
\begin{equation*}
\parallel u{\parallel }_{\frac{\gamma n}{n - 1}} \leqslant  \frac{{C}^{\prime }\gamma }{n}\parallel \nabla u{\parallel }_{p}.
\end{equation*}

取 \(C = \frac{{C}^{\prime }\gamma }{n}\) 可得引理结论对 \(u \in  {C}_{0}^{1}\left( \Omega \right)\) 成立.

最后,若 \(u \in  {W}_{0}^{1,p}\left( \Omega \right)\) ,则由定义知,存在 \({C}_{0}^{\infty }\left( \Omega \right)\) 函数列 \({\left\{  {u}_{n}\right\}  }_{n \in  \mathbb{N}}\) 按 \({W}^{1,p}\) -范数收敛到 \(u\) . 由于 \({\left\{  \nabla {u}_{n}\right\}  }_{n \in  \mathbb{N}}\) 为 \({L}^{p}\left( \Omega \right)\) 中的Cauchy列,故由前面所证结论知, \({\left\{  {u}_{n}\right\}  }_{n \in  \mathbb{N}}\) 为 \({L}^{\frac{np}{n - p}}\left( \Omega \right)\) 中的Cauchy 列. 从而 \(u\) 也属于 \({L}^{\frac{np}{n - p}}\left( \Omega \right)\) 且不等式\ref{3.26}成立.
\end{proof}
\begin{corollary}
    设 \(\Omega  \subset  {\mathbb{R}}^{n}\) 为具有光滑边界的有界区域, \(1 \leqslant  p < n,{p}^{ * } = \frac{np}{n - p}\) . 若 \(u \in  {W}_{0}^{1,p}\left( \Omega \right)\) ,则存在仅依赖于 \(n,p,q\) 及 \(\Omega\) 的常数 \(C\) 使得对每个 \(q \in  \left\lbrack  {1,{p}^{ * }}\right\rbrack\) 有
\begin{equation*}
\parallel u{\parallel }_{{L}^{q}\left( \Omega \right) } \leqslant  C\parallel \nabla u{\parallel }_{{L}^{p}\left( \Omega \right) }.
\end{equation*}

\end{corollary}
\begin{note}
    由推论3.2 知,当 \(p \rightarrow  n\) 时有 \({p}^{ * } = \frac{np}{n - p} \rightarrow   + \infty\) . 因此,我们也许会期望当 \(u \in  {W}^{1,n}\left( \Omega \right)\) 时会有 \(u \in  {L}^{\infty }\left( \Omega \right)\) . 但这一般不对,例如,若 \(\Omega  = \left\{  {x \in  {\mathbb{R}}^{n} \mid  \left| x\right|  < }\right.\)  \(\left. \frac{1}{2}\right\}\) ,则有
\begin{equation*}
u\left( x\right)  = {\left( \ln \frac{1}{\left| x\right| }\right) }^{\alpha },\;0 < \alpha  < 1 - \frac{1}{n}
\end{equation*}

属于 \({W}^{1,n}\left( \Omega \right)\) ,但 \(u\left( x\right)\) 在 \(x = 0\) 处有奇性,从而 \(u \notin  {L}^{\infty }\left( \Omega \right)\) .

\end{note}

为证明 \(p > n\) 情形的Sobolev不等式,我们先证明如下引理:

\begin{lemma}
     设 \(\Omega  \subset  {\mathbb{R}}^{n}\) 为有界开集, \(1 \leqslant  p < \infty\) . 对 \(u\left( x\right)  \in  {L}^{p}\left( \Omega \right)\) ,定义
\begin{equation*}
{\Phi }_{p}\left( u\right)  \mathrel{\text{ := }} {\left( \frac{1}{\mu \left( \Omega \right) }{\int }_{\Omega }{\left| u\left( x\right) \right| }^{p}\dif x\right) }^{\frac{1}{p}};
\end{equation*}

对满足 \({\left| u\left( x\right) \right| }^{p}\) 可测但不可积的可测函数 \(u\left( x\right)\) ,则定义 \({\Phi }_{p}\left( u\right)  \mathrel{\text{ := }} \infty\) . 则对可测函数 \(u : \Omega  \rightarrow  \mathbb{R} \cup  \{  + \infty \}\) 有
\begin{equation*}
{\left. \operatorname{ess}\mathop{\sup }\limits_{{x \in  \Omega }}\left| u\left( x\right) \right|  = \mathop{\lim }\limits_{{p \rightarrow  \infty }}{\left( \frac{1}{\mu \left( \Omega \right) }{\int }_{\Omega }{\left| u\left( x\right) \right| }^{p}\dif x\right) }^{\frac{1}{p}}\right) }^{\frac{1}{p}}
\end{equation*}

其中 \(\mu \left( \Omega \right)\) 为 \(\Omega\) 的体积.
\end{lemma}
\begin{proof}
    由Hölder不等式可知, \({\Phi }_{p}\left( u\right)\) 关于 \(p\) 单调增加且对 \(1 \leqslant  p < \infty\) 有
\begin{equation*}
{\Phi }_{p}\left( u\right)  \leqslant  \operatorname{ess}\mathop{\sup }\limits_{{x \in  \Omega }}\left| {u\left( x\right) }\right| .
\end{equation*}

由单调有界原理知, \(\mathop{\lim }\limits_{{p \rightarrow  \infty }}{\Phi }_{p}\left( u\right)  \in  \mathbb{R} \cup  \{  + \infty \}\) . 因此,只需证明
\begin{equation*}
{\operatorname{esssup}}_{x \in  \Omega }\left| {u\left( x\right) }\right|  \leqslant  \mathop{\lim }\limits_{{p \rightarrow  \infty }}{\Phi }_{p}\left( u\right) .
\end{equation*}

对 \(K \in  \mathbb{R}\) ,令
\begin{equation*}
{A}_{K} \mathrel{\text{ := }} \{ x \in  \Omega \left| \right| u\left( x\right)  \mid   \geq  K\} .
\end{equation*}

由 \(u\) 为可测函数知, \({A}_{K}\) 为可测集,而且对 \(K < \operatorname{ess}\mathop{\sup }\limits_{{x \in  \Omega }}\left| {u\left( x\right) }\right|\) 有 \(\mu \left( {A}_{K}\right)  > 0\) . 从而,
\begin{equation*}
{\Phi }_{p}\left( u\right)  \geq  \mu {\left( \Omega \right) }^{-\frac{1}{p}}{\left( {\int }_{{A}_{K}}{\left| u\left( x\right) \right| }^{p}\dif x\right) }^{\frac{1}{p}} \geq  \mu {\left( \Omega \right) }^{-\frac{1}{p}}\mu {\left( {A}_{K}\right) }^{\frac{1}{p}}K.
\end{equation*}

令 \(p \rightarrow  \infty\) 得
\begin{equation*}
\mathop{\lim }\limits_{{p \rightarrow  \infty }}{\Phi }_{p}\left( u\right)  \geq  K
\end{equation*}

由于上式对所有的 \(K < \operatorname{ess}\mathop{\sup }\limits_{{x \in  \Omega }}\left| {u\left( x\right) }\right|\) 都成立,故有
\begin{equation*}
\mathop{\lim }\limits_{{p \rightarrow  \infty }}{\Phi }_{p}\left( u\right)  \geq  \operatorname{ess}\mathop{\sup }\limits_{{x \in  \Omega }}\left| {u\left( x\right) }\right| .
\end{equation*}

引理得证.

\end{proof}

接下来,我们可以证明 \(p > n\) 情形的模估计.
\begin{lemma}\label{3.4.4}
    设 \(\Omega  \subset  {\mathbb{R}}^{n}\) 为有界区域, \(p > n \geq  1\) . 则存在常数 \(C = C\left( {n,p,\Omega }\right)\) 使得对任意 \(u \in  {W}_{0}^{1,p}\left( \Omega \right)  \cap  {L}^{\infty }\left( \Omega \right)\) 都有 \(u \in  {C}^{0}\left( \bar{\Omega }\right)\) 且
\begin{equation}
\mathop{\sup }\limits_{{x \in  \Omega }}\left| {u\left( x\right) }\right|  \leqslant  C\parallel \nabla u{\parallel }_{{L}^{p}\left( \Omega \right) }. \label{3.29}
\end{equation}

\end{lemma}
\begin{proof}
    首先,考虑 \(n = 1\) 的情形,即 \(\Omega\) 为开区间 \(I = \left( {c,d}\right)  \subset  \mathbb{R}\) . 若 \(u \in  {W}_{0}^{1,1}\left( I\right)\) , 则由 \({W}_{0}^{1,1}\left( I\right)\) 的定义知,存在函数列 \({\left\{  {u}_{n}\right\}  }_{n = 1}^{\infty } \subset  {C}_{0}^{\infty }\left( I\right)\) 使得 \({u}_{n}\) 按 \({W}^{1,1}\) 范数收敛到 \(u\) , 即 \(\mathop{\lim }\limits_{{n \rightarrow  \infty }}{\begin{Vmatrix}{u}_{n} - u\end{Vmatrix}}_{1,1} = 0\) . 由于 \({u}_{n}\) 及 \(u\) 都是 \(I\) 上局部可积函数,从而对 \(I\) 中任意有界闭集 \({I}^{\prime }\) 有, \(\nabla {u}_{n}\) 在 \({L}^{1}\left( {I}^{\prime }\right)\) 中收敛到 \(\nabla u\) ,即有
\begin{equation*}
\mathop{\lim }\limits_{{n \rightarrow  \infty }}\left\{  {{\int }_{{I}^{\prime }}\left| {{u}_{n}\left( x\right)  - u\left( x\right) }\right| {\dif x} + {\int }_{{I}^{\prime }}\left| {{u}_{n}^{\prime }\left( x\right)  - {u}^{\prime }\left( x\right) }\right| {\dif x}}\right\}   = 0.
\end{equation*}

由于 \({u}_{n}\) 光滑,故对任意 \(a,b \in  I\) 及 \(\varepsilon  > 0\) ,当 \(n\) 充分大时有
\begin{equation}
\left| {{u}_{n}\left( b\right)  - {u}_{n}\left( a\right) }\right|  = \left| {{\int }_{a}^{b}{u}_{n}^{\prime }\left( x\right) {\dif x}}\right|  \leqslant  {\int }_{a}^{b}\left| {{u}_{n}^{\prime }\left( x\right) }\right| {\dif x}
\leqslant  {\int }_{a}^{b}\left| {{u}^{\prime }\left( x\right) }\right| {\dif x} + \frac{\varepsilon }{2}. \label{3.30}
\end{equation}

令 \(a \rightarrow  c\) ,我们有
\begin{equation*}
\left| {{u}_{n}\left( b\right) }\right|  \leqslant  {\int }_{c}^{b}\left| {{u}^{\prime }\left( x\right) }\right| {\dif x} + \frac{\varepsilon }{2} \leqslant  {\int }_{I}\left| {{u}^{\prime }\left( x\right) }\right| {\dif x} + \frac{\varepsilon }{2}.
\end{equation*}

这意味着 \(\left\{  {u}_{n}\right\}\) 一致有界.

此外,由于 \(\nabla u \in  {L}^{1}\left( I\right)\) ,故存在 \(\delta  > 0\) 使得当 \(\left| {a - b}\right|  < \delta\) 时有
\begin{equation*}
{\int }_{a}^{b}\left| {\nabla u\left( x\right) }\right| {\dif x} < \frac{\varepsilon }{2}.
\end{equation*}

因此,我们有当 \(n\) 充分大且 \(\left| {a - b}\right|  < \delta\) 时有
\begin{equation}
\left| {{u}_{n}\left( b\right)  - {u}_{n}\left( a\right) }\right|  < \varepsilon . \label{3.31}
\end{equation}

这意味着 \({u}_{n}\) 等度连续. 从而由Arzela-Ascoli定理知, \({u}_{n}\) 一致收敛到极限 \(u\) .

在\ref{3.30}中令 \(n\) 趋于无穷大取极限得
\begin{equation}
\left| {u\left( b\right)  - u\left( a\right) }\right|  \leqslant  {\int }_{a}^{b}\left| {\nabla u\left( x\right) }\right| {\dif x},\;\forall a,b \in  I,a < b. \label{3.32}
\end{equation}

因此, \(u\) 为一致连续函数. 由 \(u\) 的任意性知, \({W}_{0}^{1,1}\left( I\right)\) 中所有函数都是一致连续的. 由Hölder不等式知 \({W}_{0}^{1,p}\left( I\right)  \subset  {W}_{0}^{1,1}\left( I\right)\) 中所有函数也都是一致连续的. 从而由\ref{3.30}可得引理结论成立.

下面考虑 \(n > 1\) 的情形. 首先假定 \(u \in  {C}_{0}^{\infty }\left( \Omega \right)\) . 令 \(x = \left( {{x}_{1},\cdots ,{x}_{n}}\right)\) . 若假定
\begin{equation}
\mu \left( \Omega \right)  = 1 \label{3.33}
\end{equation}

且
\begin{equation}
\parallel \nabla u{\parallel }_{{L}^{p}\left( \Omega \right) } = 1 \label{3.34}
\end{equation}

则由\ref{3.30}知,对 \(\gamma  > 1\) 有
\begin{equation*}
{\begin{Vmatrix}{\left| u\right| }^{\gamma }\end{Vmatrix}}_{\frac{n}{n - 1}} \leqslant  \frac{\gamma }{n}{\begin{Vmatrix}{\left| u\right| }^{\gamma  - 1}\end{Vmatrix}}_{\frac{p}{p - 1}}.
\end{equation*}

因此,
\begin{equation*}
\parallel u{\parallel }_{\frac{\gamma n}{n - 1}} \leqslant  {\left( \frac{\gamma }{n}\right) }^{\frac{1}{\gamma }}\parallel u{\parallel }_{\frac{\left( {\gamma  - 1}\right) p}{p - 1}}^{\frac{\gamma  - 1}{\gamma }}.
\end{equation*}

从而由 \(\mu \left( \Omega \right)  = 1\) 及引理3.3证明中的单调性结论有
\begin{equation}
\parallel u{\parallel }_{\frac{\gamma n}{n - 1}} \leqslant  {\left( \frac{\gamma }{n}\right) }^{\frac{1}{\gamma }}\parallel u{\parallel }_{\frac{\gamma p}{p - 1}}^{\frac{\gamma  - 1}{\gamma }}. \label{3.35}
\end{equation}

由 \(p > n\) 知,
\begin{equation*}
\gamma  = \frac{n}{n - 1} \cdot  \frac{p - 1}{p} > 1
\end{equation*}

在\ref{3.35}中用 \({\gamma }^{k}\) 代替 \(\gamma\) 得
\begin{equation}
\parallel u{\parallel }_{\frac{{\gamma }^{k}n}{n - 1}} \leqslant  {\left( \frac{{\gamma }^{k}}{n}\right) }^{\frac{1}{{\gamma }^{k}}}\parallel u{\parallel }_{\frac{{\gamma }^{k} - {1}_{n}}{n - 1}}^{\frac{{\gamma }^{k} - 1}{{\gamma }^{k}}}, \label{3.36}
\end{equation}

这里用到 \({\gamma }^{k}\frac{p}{p - 1} = {\gamma }^{k - 1}\frac{n}{n - 1}\) .

由引理3.3知,对可测函数 \(u\left( x\right)\)
\begin{equation}
{\left. \operatorname{ess}\mathop{\sup }\limits_{{x \in  \Omega }}\left| u\left( x\right) \right|  = \mathop{\lim }\limits_{{p \rightarrow  \infty }}{\left( \frac{1}{\mu \left( \Omega \right) }{\int }_{\Omega }{\left| u\left( x\right) \right| }^{p}\dif x\right) }^{\frac{1}{p}}\right) }^{\frac{1}{p}}. \label{3.37}
\end{equation}

下面分两种情形来考虑:

(i)若对无穷多个 \(k \in  \mathbb{N}\) 成立
\begin{equation*}
\parallel u{\parallel }_{{\gamma }^{k}\frac{n}{n - 1}} \leqslant  1
\end{equation*}

则有
\begin{equation*}
\text{ ess }\mathop{\sup }\limits_{{x \in  \Omega }}\left| {u\left( x\right) }\right|  \leqslant  1\text{ , }
\end{equation*}

其中我们用到 \(\mu \left( \Omega \right)  = 1\) .

(ii) 若存在 \({k}_{0}\) 使得当 \(k \geq  {k}_{0}\) 时有
\begin{equation*}
\parallel u{\parallel }_{{\gamma }^{k}\frac{n}{n - 1}} \geq  1
\end{equation*}

则由\ref{3.36}知,当 \(k \geq  {k}_{0} + 1\) 时有
\begin{equation}
\parallel u{\parallel }_{{\gamma }^{k}\frac{n}{n - 1}} \leqslant  {\left( \frac{{\gamma }^{k}}{n}\right) }^{\frac{1}{{\gamma }^{k}}}\parallel u{\parallel }_{{\gamma }^{k - 1}\frac{n}{n - 1}}. \label{3.38}
\end{equation}

通过迭代由 (3.38) 知
\begin{equation}
\parallel u{\parallel }_{{\gamma }^{k}\frac{n}{n - 1}} \leqslant  \frac{{\gamma }^{\mathop{\sum }\limits_{{\nu  = {k}_{0} + 1}}^{k}\nu \frac{1}{{\gamma }^{\nu }}}}{{n}^{\mathop{\sum }\limits_{{\nu  = {k}_{0} + 1}}^{k}\frac{1}{{\gamma }^{\nu }}}}\parallel u{\parallel }_{{\gamma }^{{k}_{0}}\frac{n}{n - 1}}. \label{3.39}
\end{equation}

在 \ref{3.36} 中令 \(k = {k}_{0}\) ,则右边得到 \(\parallel u{\parallel }_{{\gamma }^{{k}_{0} - 1}\frac{n}{n - 1}}\) . 不妨假设 \(\parallel u{\parallel }_{{\gamma }^{{k}_{0} - 1}\frac{n}{n - 1}} \leqslant  1\) 或者 \({k}_{0} = 1\) .

由于\ref{3.27}及 \(\parallel \nabla u{\parallel }_{p} = 1\) ,从而由\ref{3.37}令 \(k \rightarrow  \infty\) 取极限得
\begin{equation}
{\operatorname{esssup}}_{x \in  \Omega }\left| {u\left( x\right) }\right|  \leqslant  C\left( {p,n}\right) . \label{3.40}
\end{equation}

接下来,我们需要去掉规范化条件 (3.33) 和 (3.34). 若 \(\parallel \nabla u{\parallel }_{p} \neq  0\) ,令 \(v \mathrel{\text{ := }}\)  \(\frac{u}{\parallel \nabla u{\parallel }_{p}}\) . 则 \(v\) 满足 (3.34),从而由 (3.40) 知
\begin{equation}
{\operatorname{esssup}}_{x \in  \Omega }\left| {u\left( x\right) }\right|  \leqslant  C\left( {p,n}\right) \parallel \nabla u{\parallel }_{p}. \label{3.41}
\end{equation}

若 \({C}_{0}^{1}\) -函数 \(u\) 满足 \(\parallel \nabla u{\parallel }_{p} = 0\) ,则 \(u \equiv  0\) . 此时,结论自然成立.

最后, 为去掉条件 (3.33), 我们考虑坐标变换
\begin{equation*}
y = y\left( x\right)  = \mu {\left( \Omega \right) }^{-\frac{1}{n}}x
\end{equation*}

及
\begin{equation*}
\widetilde{\Omega } \mathrel{\text{ := }} \{ y\left( x\right)  \mid  x \in  \Omega \} ,\;\widetilde{u}\left( y\right)  \mathrel{\text{ := }} u\left( x\right) .
\end{equation*}

则有
\begin{equation*}
\mu \left( \widetilde{\Omega }\right)  = 1
\end{equation*}

且
\begin{equation*}
{\nabla }_{i}u\left( x\right)  = {\nabla }_{i}\widetilde{u}\left( y\right) \mu {\left( \Omega \right) }^{-\frac{1}{n}}.
\end{equation*}

从而有
\begin{equation*}
{\left( {\int }_{\Omega }{\left| {\nabla }_{i}u\left( x\right) \right| }^{p}\dif x\right) }^{\frac{1}{p}} = {\left( {\int }_{\widetilde{\Omega }}{\left| {\nabla }_{i}\widetilde{u}\left( y\right) \right| }^{p}\dif y\right) }^{\frac{1}{p}} \cdot  \mu {\left( \Omega \right) }^{\frac{1}{p} - \frac{1}{n}}.
\end{equation*}

应用(3.41)到 \(\widetilde{u}\) 可得
\begin{equation*}
{\operatorname{esssup}}_{x \in  \Omega }\left| {u\left( x\right) }\right|  \leqslant  C\left( {p,n}\right) \mu {\left( \Omega \right) }^{\frac{1}{n} - \frac{1}{p}}\parallel \nabla u{\parallel }_{p}.
\end{equation*}

至此,我们证明了对 \({C}_{0}^{1}\left( \Omega \right)\) 中所有函数 \(u\) 引理结论成立.
\end{proof}

若 \(u \in  {W}_{0}^{1,p}\left( \Omega \right)\) ,则由定义知,存在 \({C}_{0}^{\infty }\left( \Omega \right)\) 函数列 \({\left\{  {u}_{n}\right\}  }_{n \in  \mathbb{N}}\) 按 \({W}^{1,p}\) -范数收敛到 \(u\) . 由于 \({\left\{  \nabla {u}_{n}\right\}  }_{n \in  \mathbb{N}}\) 为 \({L}^{p}\left( \Omega \right)\) 中的Cauchy列,故由 (3.29) 知, \({\left\{  {u}_{n}\right\}  }_{n \in  \mathbb{N}}\) 为 \({C}^{0}\left( \bar{\Omega }\right)\) 中的Cauchy 列. 从而 \(u\) 也属于 \({C}^{0}\left( \bar{\Omega }\right)\) 且不等式 (3.29) 成立.

至此, 我们可以叙述著名的Poincaré不等式如下:
\begin{corollary}
    设 \(\Omega  \subset  {\mathbb{R}}^{n}\) 为有界区域, \(1 \leqslant  p <  + \infty\) . 则存在常数 \(C = C\left( {n,p,\Omega }\right)\) 使得对任意 \(u \in  {W}_{0}^{1,p}\left( \Omega \right)\) 都有
\begin{equation*}
\parallel u{\parallel }_{{L}^{p}\left( \Omega \right) } \leqslant  C\parallel \nabla u{\parallel }_{{L}^{p}\left( \Omega \right) }.
\end{equation*}

\end{corollary}
\begin{proof}
     由引理\ref{3.4.2}和引理\ref{3.4.4}知,我们只需证明结论对 \(p = n > 1\) 情形也成立即可. 实际上,由 (3.28) 知,对 \(\gamma  > 1\) ,有
\begin{equation*}
\parallel u{\parallel }_{\frac{\gamma n}{n - 1}}^{\gamma } \leqslant  \frac{\gamma {C}^{\prime }}{n}\parallel u{\parallel }_{\frac{\left( {\gamma  - 1}\right) n}{n - 1}}^{\gamma  - 1}\parallel \nabla u{\parallel }_{n}
\leqslant  \frac{\gamma {C}^{\prime \prime }}{n}\parallel u{\parallel }_{\frac{\gamma n}{n - 1}}^{\gamma  - 1}\parallel \nabla u{\parallel }_{n}
\end{equation*}

其中在第二个不等式中使用了Hölder 不等式. 从而有
\begin{equation*}
\parallel u{\parallel }_{\frac{\gamma n}{n - 1}} \leqslant  \frac{\gamma {C}^{\prime \prime }}{n}\parallel \nabla u{\parallel }_{n}
\end{equation*}

取 \(\gamma  = n - 1\) 即得所需结论.
\end{proof}
\begin{note}
    从Poincaré不等式可以看出,我们可以在Sobolev空间 \({W}_{0}^{k,p}\left( \Omega \right)\) 上选取与范数 \ref{3.24} 等价的范数
\begin{equation*}
{\left( {\int }_{\Omega }\mathop{\sum }\limits_{{\left| \alpha \right|  = k}}{\left| {\partial }^{\alpha }u\left( x\right) \right| }^{p}\dif x\right) }^{\frac{1}{p}}.
\end{equation*}

从数学分析中, 我们知道, 若光滑函数的所有导数都为零, 则该函数必为常值函数. 实际上,著名的Poincaré不等式告诉我们该结论对于 \(u \in  {W}^{1,p}\left( \Omega \right)\) 也对.
\end{note}

\begin{corollary}
    设 \(\Omega  \subset  {\mathbb{R}}^{n}\) 为有界区域, \(u \in  {W}^{1,p}\left( \Omega \right)\) 的弱导数 \(\frac{\partial u}{\partial {x}_{j}} \equiv  0,j =\)  \(1,2,\cdots ,n\) . 则 \(u\) 在 \(\Omega\) 上为常数. 若进一步有 \(u \in  {W}_{0}^{1,p}\left( \Omega \right)\) ,则 \(u \equiv  0\) .
\end{corollary}
\begin{proof}
    记函数 \(u\) 的平均值为
\begin{equation*}
{\left( u\right) }_{\Omega } \mathrel{\text{ := }} \frac{{\int }_{\Omega }u\left( x\right) {\dif x}}{\mu \left( \Omega \right) }.
\end{equation*}

不妨假设 \({\left( u\right) }_{\Omega } = 0\) ,否则,考虑函数 \(v = u - {\left( u\right) }_{\Omega }\) 即可. 由定理 3.4.3 知,我们只需对 \(u \in  {C}^{\infty }\left( \bar{\Omega }\right)\) 证明结论即可. 为此,我们现在

断言: 当 \(\Omega\) 为方体时,存在常数 \(C\) 使得对 \({C}^{\infty }\left( \Omega \right)\) 中平均值为零的函数 \(u\left( x\right)\) 成立
\begin{equation*}
\parallel u{\parallel }_{{L}^{p}\left( \Omega \right) } \leqslant  C\parallel \nabla u{\parallel }_{{L}^{p}\left( \Omega \right) }.
\end{equation*}

下面将关于 \(n\) 用归纳法来证明结论. 不失一般性,我们可以假设 \(\Omega\) 为单位立方体 \({Q}^{n} = \left\{  {x \in  {\mathbb{R}}^{n} \mid  0 < {x}_{j} < 1}\right\}\) .

当 \(n = 1\) 时,我们有
\begin{equation*}
u\left( x\right)  = u\left( 0\right)  + {\int }_{0}^{x}{u}^{\prime }\left( t\right) {\dif t} = u\left( 0\right) .
\end{equation*}

由 \({\left( u\right) }_{\Omega } = 0\) 知 \(u\left( x\right)  \equiv  0\) . 故断言成立.

假设断言对 \(n \geq  1\) 成立. 设 \(u \in  {C}^{\infty }\left( {Q}^{n + 1}\right)\) 具有平均值零. 定义 \(v\left( t\right)  \in  {C}^{\infty }\left( {Q}^{1}\right)\) 为
\begin{equation*}
v\left( t\right)  = {\int }_{{Q}^{n}}u\left( {{x}_{1},\cdots ,{x}_{n},t}\right) \dif x_1\cdots \dif{x}_{n}
\end{equation*}

且
\begin{equation*}
b\left( t\right)  = {\int }_{0}^{t}{v}^{\prime }\left( s\right) {\dif s}.
\end{equation*}

由于 \(v\) 具有平均值零,故
\begin{equation*}
{\int }_{0}^{1}{\left| v\left( t\right) \right| }^{p}{\dif t} \leqslant  {\int }_{0}^{1}{\left( {\int }_{0}^{t}\left| {v}^{\prime }\left( s\right) \right| \dif s\right) }^{p}{\dif t} \leqslant  {C}_{1}{\int }_{0}^{1}{\left| {v^\prime}\left( t\right) \right| }^{p}{\dif t} \leqslant  {C}_{1}{\int }_{{Q}^{n + 1}}{\left| \frac{\partial u}{\partial {x}_{n + 1}}\right| }^{p}{\dif x},
\end{equation*}

其中第二个不等式利用了Hölder不等式. 由归纳假设知,
\begin{equation*}
{\int }_{{Q}^{n}}{\left| u\left( x,t\right)  - v\left( t\right) \right| }^{p}{\dif x} \leqslant  {C}_{2}{\int }_{{Q}^{n}}{\left| \nabla u\left( x,t\right) \right| }^{p}{\dif x}.
\end{equation*}

上式两边关于变量 \(t\) 积分并利用前一不等式则可得
\begin{equation*}
\parallel u{\parallel }_{{L}^{P}\left( {Q}^{n + 1}\right) } \leqslant  C\parallel \nabla u{\parallel }_{{L}^{p}\left( {Q}^{n + 1}\right) }.
\end{equation*}

至此, 由归纳法断言成立.

由断言知,如果 \(u\) 的弱导数 \(\frac{\partial u}{\partial {x}_{j}} \equiv  0,j = 1,2,\cdots ,n\) ,则有 \(u\) 在 \(\Omega\) 中的任意方体上为常数,由 \(\Omega\) 的连通性知, \(u\) 在 \(\Omega\) 上为常数. 由于$C^\infty$在$W^{1,p}$上稠密,设$\{u_k\}\subset{ C^\infty},u\in W^{1,p},u_k\to u,u_{x_i}=0$,那么
\[
\|u-(u)_\Omega\|_p\leqslant \|u-u_k\|_p+\|u_k-(u_k)_\Omega\|_p+\|(u)_\Omega -(u_k)_\Omega\|_p\leqslant \|u-u_k\|_p+C\|\nabla u_k\|_p+\|(u)_\Omega -(u_k)_\Omega\|_p\to 0
\]

 \(u\) 在 \(\Omega\) 上为常数.


特别,若 \(u \in  {W}_{0}^{1,p}\left( \Omega \right)\) ,则有 \(u \equiv  0\) . 至此, 命题得证.
\end{proof}


\subsection{Sobolev嵌入定理}

当 \(\Omega\) 为有界区域时,由于 \({L}^{p} \subset  {L}^{{p}^{\prime }}\left( {p \geq  {p}^{\prime }}\right)\) ,故存在常数 \(C\left( {p,\Omega }\right)  > 0\) 有
\begin{equation*}
\left\{  \begin{array}{l} {W}^{k,p}\left( \Omega \right)  \subset  {W}^{{k}^{\prime },{p}^{\prime }}\left( \Omega \right) ,\;{k}^{\prime } < k,\;{p}^{\prime } \leqslant  p, \\  \parallel u{\parallel }_{{k}^{\prime },{p}^{\prime }} \leqslant  C\left( {p,\Omega }\right) \parallel u{\parallel }_{k,p},\;u \in  {W}^{k,p}. \end{array}\right.
\end{equation*}

其中范数不等式表明嵌入映射 \(i : {W}^{k,p}\left( \Omega \right)  \hookrightarrow  {W}^{{k}^{\prime }{p}^{\prime }}\left( \Omega \right)\) 为连续映射. 特别有连续嵌入
\begin{equation*}
{W}_{0}^{k,p}\left( \Omega \right)  \hookrightarrow  {W}_{0}^{k,p}\left( {\mathbb{R}}^{n}\right)  \equiv  {W}^{k,p}\left( {\mathbb{R}}^{n}\right) .
\end{equation*}

实际上, 引理\ref{3.4.2}和引理\ref{3.4.4}蕴含下面的Sobolev嵌入定理.
\begin{theorem}\label{3.9}
    设 \(\Omega  \subset  {\mathbb{R}}^{n}\) 为带光滑边界的有界区域, \(u \in  {W}_{0}^{1,p}\left( \Omega \right)\) ,则当 \(p < n\) 时有 \(u \in  {L}^{\frac{np}{n - p}}\left( \Omega \right)\) ; 当 \(p > n\) 时有 \(u \in  {C}^{0}\left( \bar{\Omega }\right)\) .
\end{theorem}
\begin{theorem}
     [Sobolev 嵌入定理] 设 \(\Omega  \subset  {\mathbb{R}}^{n}\) 为带光滑边界的有界区域,则
\begin{equation*}
{W}_{0}^{k,p} \hookrightarrow  \left\{  \begin{array}{ll} {L}^{\frac{np}{n - {kp}}}\left( \Omega \right) , & \text{ 当 }{kp} < n\text{ 时 } \\  {C}^{\ell }\left( \bar{\Omega }\right) , & \text{ 当 }0 \leqslant  \ell  < k - \frac{n}{p}\text{ 时. } \end{array}\right.
\end{equation*}

而且存在常数 \(C = C\left( {n,p,\Omega }\right)\) 使得对任意 \(u \in  {W}_{0}^{k,p}\left( \Omega \right)\) ,有
\begin{equation*}
\left\{  \begin{array}{ll} \parallel u{\parallel }_{{L}^{\frac{np}{n - {kp}}}\left( \Omega \right) } \leqslant  C\parallel u{\parallel }_{k,p}, & \text{ 当 }{kp} < n\text{ 时 } \\  \parallel u{\parallel }_{{C}^{\ell }\left( \bar{\Omega }\right) } \leqslant  C\parallel u{\parallel }_{k,p}, & \text{ 当 }0 \leqslant  \ell  < k - \frac{n}{p}\text{ 时. } \end{array}\right.
\end{equation*}

\end{theorem}
\begin{proof}
    由于 \({kp} < n\) ,故对任意 \(1 \leqslant  l \leqslant  k - 1\) 都有 \(\left( {k - l}\right) \frac{np}{n - {lp}} < n\) . 由引理\ref{3.4.2}得
\begin{equation*}
{W}_{0}^{k,p}\left( \Omega \right)  \hookrightarrow  {W}_{0}^{k - 1,\frac{np}{n - p}}\left( \Omega \right)  \hookrightarrow  \cdots  \hookrightarrow  {W}_{0}^{1,\frac{np}{n - \left( {k - 1}\right) p}}\left( \Omega \right)  \hookrightarrow  {L}^{\frac{np}{n - {kp}}}\left( \Omega \right) .
\end{equation*}

第一个嵌入得证.

当 \(n = 1\) 时,第二个嵌入可以由引理\ref{3.4.4}的证明及定理\ref{3.9}得到. 因此,下面我们假设 \(n \geq  2\) . 此时,第二个嵌入将可以由第一个嵌入及引理 \ref{3.4.2}得到. 下面我们将就 \(k = 2\) 的情形给出证明.

首先,当 \(p < n\) 时有
\begin{equation*}
{W}_{0}^{2,p}\left( \Omega \right)  \hookrightarrow  {W}_{0}^{1,\frac{np}{n - p}}\left( \Omega \right) .
\end{equation*}

实际上,设 \(u \in  {W}_{0}^{2,p}\left( \Omega \right)\) ,则存在 \({\left\{  {u}_{n}\right\}  }_{n \in  \mathbb{N}} \subset  {C}_{0}^{\infty }\left( \Omega \right)\) 按 \({W}^{2,p}\) -范数收敛到 \(u\) . 由引理 \ref{3.4.2}知,存在常数 \({C}_{0}\) 使得
\begin{equation*}
{\begin{Vmatrix}{u}_{n} - u\end{Vmatrix}}_{{W}^{1,\frac{np}{n - p}}} \leqslant  {C}_{0}{\begin{Vmatrix}{u}_{n} - u\end{Vmatrix}}_{{W}^{2,p}}.
\end{equation*}

因此, \(u \in  {W}_{0}^{1,\frac{np}{n - p}}\left( \Omega \right)\) . 从而当 \(p < n\) 时有 \({W}_{0}^{2,p}\left( \Omega \right)  \hookrightarrow  {W}_{0}^{1,\frac{np}{n - p}}\left( \Omega \right)\) .

特别,由 \(p < n\) 及 \({2p} > n\) 知 \(\frac{np}{n - p} > n\) . 因此,由定理\ref{3.9} 可得
\begin{equation*}
{W}_{0}^{2,p}\left( \Omega \right)  \hookrightarrow  {W}_{0}^{1,\frac{np}{n - p}}\left( \Omega \right)  \hookrightarrow  {C}^{0}\left( \bar{\Omega }\right) .
\end{equation*}

其次,若 \(p > n\) 且 \(u \in  {W}_{0}^{2,p}\left( \Omega \right)\) ,则由嵌入定理\ref{3.9}知, \(u\) 及 \(\nabla u\) 都在 \(\bar{\Omega }\) 上连续. 从而 \(\nabla u\) 为 \(u\) 的古典意义下的导数,即 \(u \in  {C}^{1}\left( \bar{\Omega }\right)\) . 因此,
\begin{equation*}
{W}_{0}^{2,p}\left( \Omega \right)  \hookrightarrow  {C}^{1}\left( \bar{\Omega }\right) .
\end{equation*}

最后,若 \(p = n\) ,由 \({L}^{r}\left( \Omega \right)  \subset  {L}^{q}\left( \Omega \right) \left( {1 \leqslant  q \leqslant  r \leqslant  \infty }\right)\) 可得对任意 \(q \leqslant  n\) 有
\begin{equation*}
{W}_{0}^{2,n}\left( \Omega \right)  \subset  {W}_{0}^{1,q}\left( \Omega \right) ,\;\forall q < n.
\end{equation*}

特别,当 \(0 < \varepsilon  < 1\) 时,
\begin{equation*}
{W}_{0}^{2,n}\left( \Omega \right)  \hookrightarrow  {W}_{0}^{2,n - \varepsilon }\left( \Omega \right)  \hookrightarrow  {W}_{0}^{1,\frac{n\left( {n - \varepsilon }\right) }{\varepsilon }}\left( \Omega \right) ,
\end{equation*}

其中第二个嵌入使用了引理\ref{3.4.2}. 由于对任意 \(q \in  \lbrack 1, + \infty \) ,存在 \({\varepsilon }_{0} \in  \left( {0,1}\right)\) 使得 \(q <\)  \(\frac{n\left( {n - {\varepsilon }_{0}}\right) }{{\varepsilon }_{0}}\) ,从而有
\begin{equation*}
{W}_{0}^{2,n}\left( \Omega \right)  \hookrightarrow  {W}_{0}^{2,n - {\varepsilon }_{0}}\left( \Omega \right)  \hookrightarrow  {W}_{0}^{1,\frac{n\left( {n - {\varepsilon }_{0}}\right) }{{\varepsilon }_{0}}}\left( \Omega \right)  \hookrightarrow  {W}_{0}^{1,q}\left( \Omega \right) .
\end{equation*}

取 \({q}_{0} > n\) ,则有
\begin{equation*}
{W}_{0}^{2,n}\left( \Omega \right)  \hookrightarrow  {W}_{0}^{1,{q}_{0}}\left( \Omega \right)  \hookrightarrow  {C}^{0}\left( \bar{\Omega }\right) .
\end{equation*}

至此,我们已经证明,若 \(p > n\) ,则有 \({W}_{0}^{2,p}\left( \Omega \right)  \hookrightarrow  {C}^{1}\left( \bar{\Omega }\right)\) ; 若 \(p < n\) 且 \(\frac{np}{n - p} >\)  \(n\) (即 \({2p} > n\) ),则有 \({W}_{0}^{2,p}\left( \Omega \right)  \subset  {C}^{0}\left( \bar{\Omega }\right)\) ; 若 \(p = n\) ,则有 \({W}_{0}^{2,n}\left( \Omega \right)  \subset  {C}^{0}\left( \bar{\Omega }\right)\) . 即第二个嵌入对 \(k = 2\) 成立.同理可证对任意 \(k\) 的情形.
\end{proof}
\begin{theorem}
     [Sobolev 嵌入定理] 设 \(\Omega  \subset  {\mathbb{R}}^{n}\) 为区域, \({kp} > n\) ,则 \({W}^{k,p}\left( \Omega \right)\) 中任意函数 \(u\) 有 \(\ell  = k - \left\lbrack  \frac{n}{p}\right\rbrack   - 1\) 阶以下连续导数,即 \(u \in  {C}^{\ell }\left( \Omega \right)\) ,且嵌入映射
\begin{equation}
{W}^{k,p}\left( \Omega \right)  \hookrightarrow  {C}^{\ell }\left( \Omega \right)  \label{3.42}
\end{equation}

为连续映射,即对 \(\Omega\) 的任意紧子集 \(K\) ,存在常数 \(C\) 使得
\begin{equation*}
\mathop{\sup }\limits_{\substack{{x \in  K} \\  {\left| \alpha \right|  \leqslant  \ell } }}\left| {{\partial }^{\alpha }u\left( x\right) }\right|  \leqslant  C\parallel u{\parallel }_{k,p}.
\end{equation*}

\end{theorem}
\begin{proof}
    若 \(\mathrm{K}\) 为 \(\Omega\) 中紧子集,则存在常数 \(C > 0\) ,在用适当函数 \(\varphi  \in  {C}_{0}^{\infty }\left( \Omega \right)\) 乘 \(u\) 后有 \({u}_{j} \in  {C}_{0}^{\infty }\left( \Omega \right) ,{\begin{Vmatrix}{u}_{j} - u\end{Vmatrix}}_{k,p} \rightarrow  0\) ,且可以认为 \(\Omega\) 有界,使得当 \(\left| \alpha \right|  < k - \frac{n}{p}\) 时有
\begin{equation*}
\mathop{\sup }\limits_{{x \in  K}}\left| {{\partial }^{\alpha }\left( {{u}_{i} - {u}_{j}}\right) }\right|  \leqslant  C{\begin{Vmatrix}{\partial }^{\alpha }\left( {u}_{i} - {u}_{j}\right) \end{Vmatrix}}_{k - \left| \alpha \right| ,p}
\leqslant  C{\begin{Vmatrix}{u}_{i} - {u}_{j}\end{Vmatrix}}_{k,p} \rightarrow  0,\;i,j \rightarrow  \infty .
\end{equation*}

因此,当 \(\left| \alpha \right|  < k - \frac{n}{p}\) 时, \({\partial }^{\alpha }{u}_{j}\) 在 \(\Omega\) 的任意紧子集上一致收敛. 定理得证.
\end{proof}
\begin{note}
    我们知道, \({L}^{p}\left( \Omega \right)\) 中的元素并不是处处有定义的函数. 若 \(u \in  {L}^{p}\left( \Omega \right)\) ,则在零测集上任意改变它的值并不影响其范数 \(\parallel u{\parallel }_{{L}^{p}}\) . 所以 \({L}^{p}\left( \Omega \right)\) 中所有除去 \(\Omega\) 内一个零测集外有意义且相等的函数都被看成是同一个元素,或者说, \(u \in  {L}^{p}\left( \Omega \right)\) 是代表了一个等价类,这个类中的元素除去 \(\Omega\) 中的一个零测集外互相相等, \({W}^{k,p}\left( \Omega \right)\) 中的元素也是如此. 那末 (3.42) 中的嵌入 \({W}^{k,p}\left( \Omega \right)  \hookrightarrow  {C}^{\ell }\left( \Omega \right) \left( {0 \leqslant  \ell  < k - \frac{n}{p}}\right)\) 是什么意思？准确地讲,它应该理解为 “等价类” \(u \in  {W}^{k,p}\left( \Omega \right)\) 中应包含一个函数, 它属于 \({C}^{\ell }\left( \Omega \right)\) ,而且其在 \({C}^{\ell }\left( \Omega \right)\) 中的范数可以由 \(\parallel u{\parallel }_{k,p}\) 来控制. 换言之,就是每个 \(u \in  {W}^{k,p}\left( \Omega \right)\) ,当我们把它视为一个函数时,可以在 \(\Omega\) 中的一个零测集上重新定义它的值使得修改后的函数 \(\widetilde{u}\) (它在 \({W}^{k,p}\left( \Omega \right)\) 中等于 \(u\) )属于 \({C}^{\ell }\left( \Omega \right)\) ,而且 \(\parallel \widetilde{u}{\parallel }_{{C}^{\ell }\left( \Omega \right) } \leqslant\)  \(C\parallel u{\parallel }_{k,p}\) (C与 \(u\) 无关).

 Sobolev在1938年在 \(\Omega\) 上而非在任一紧子集 \(\mathrm{K}\) 上证明了一些更精确地估计式,但要求 \(\Omega\) 满足锥条件. 从而将嵌入定理推广到更一般的区域,参见[AF].

\end{note}

\subsection{Sobolev紧嵌入定理}
\begin{definition}
    设 \(X,Y\) 为Banach空间, \(X \subset  Y\) . 称 \(X\) 紧嵌入到 \(Y\) 中,记为 \(X \subset   \subset  Y\) ,如果

(i) 存在某个常数 \(C\) 使得对任意 \(x \in  X\) 都有 \(\parallel x{\parallel }_{Y} \leqslant  C\parallel x{\parallel }_{X}\) .

(ii) \(X\) 中的有界序列在 \(Y\) 中列紧.
\end{definition}

定理 3.10 表明,对 \({\mathbb{R}}^{n}\) 中有界开集 \(\Omega\) 及 \(p < n,{W}_{0}^{1,p}\left( \Omega \right)\) 可以连续嵌入到 \({L}^{\frac{np}{n - p}}\left( \Omega \right)\) , 即存在连续 (等价于有界) 线性单射
\begin{equation}
i : {W}_{0}^{1,p}\left( \Omega \right)  \rightarrow  {L}^{\frac{np}{n - p}}\left( \Omega \right) . \label{3.43}
\end{equation}

此外,注意到对有界开集 \(\Omega\) ,存在连续嵌入
\begin{equation}
j : {L}^{r}\left( \Omega \right)  \rightarrow  {L}^{q}\left( \Omega \right) ,\;1 \leqslant  q \leqslant  r \leqslant  \infty . \label{3.44}
\end{equation}

这意味着我们有连续嵌入
\begin{equation*}
j \circ  i : {W}_{0}^{1,p}\left( \Omega \right)  \rightarrow  {L}^{q}\left( \Omega \right) ,\;1 \leqslant  q \leqslant  \frac{np}{n - p}.
\end{equation*}

我们自然会问,上述嵌入 \(j \circ  i\) 是否为紧致嵌入? Rellich-Kondrachov紧嵌入定理对此给出了肯定回答.
\begin{theorem}
     [Rellich-Kondrachov紧嵌入定理] 设 \(\Omega  \subset  {\mathbb{R}}^{n}\) 为具有光滑边界 \(\partial \Omega\) 的有界区域. 则对任意 \(1 \leqslant  p < n\) 及 \(1 \leqslant  q < {p}^{ * } = \frac{np}{n - p}\) ,和 \(p \geq  n\) 及 \(1 \leqslant  q < \infty\) 都有紧嵌入
\begin{equation*}
{W}_{0}^{1,p}\left( \Omega \right)  \subset   \subset  {L}^{q}\left( \Omega \right) .
\end{equation*}

\end{theorem}
\begin{proof}
     对 \(p \geq  n\) 及 \(1 \leqslant  q < \infty\) 的情形,我们总可以找到 \(r < n\) 使得 \(q < \frac{nr}{n - r}\) . 由于对 \(1 \leqslant  r \leqslant  p < \infty\) 有连续嵌入 \({L}^{p}\left( \Omega \right)  \hookrightarrow  {L}^{r}\left( \Omega \right)\) ,从而有连续嵌入
\begin{equation*}
{W}_{0}^{1,p}\left( \Omega \right)  \hookrightarrow  {W}_{0}^{1,r}\left( \Omega \right)
\end{equation*}

如果我们能证明 \(1 \leqslant  p < n\) 及 \(1 \leqslant  q < {p}^{ * } = \frac{np}{n - p}\) 情形结论成立,则连续嵌 \(\lambda {W}_{0}^{1,p}\left( \Omega \right)  \hookrightarrow  {W}_{0}^{1,r}\left( \Omega \right)\) 与紧嵌入 \({W}_{0}^{1,r}\left( \Omega \right)  \subset   \subset  {L}^{q}\left( \Omega \right)\) 的复合就给出我们需要的紧嵌入.

因此,我们只需证明紧嵌入结论对 \(1 \leqslant  p < n\) 及 \(1 \leqslant  q < {p}^{ * } = \frac{np}{n - p}\) 情形成立即可. 实际上,只需证明若 \({\left\{  {u}_{k}\right\}  }_{k = 1}^{\infty }\) 为 \({W}_{0}^{1,p}\left( \Omega \right)\) 中有界列,则存在子列 \({\left\{  {u}_{{k}_{j}}\right\}  }_{j = 1}^{\infty }\) 在 \({L}^{q}\left( \Omega \right)\) 中收敛.

我们先证 \(q = 1\) 的情形. 设 \({\left\{  {u}_{k}\right\}  }_{k \in  \mathbb{N}} \subset  {W}_{0}^{1,p}\left( \Omega \right)\) 为有界点列,即存在 \(K > 0\) 使得
\begin{equation}
{\begin{Vmatrix}{u}_{k}\end{Vmatrix}}_{{W}^{1,p}\left( \Omega \right) } \leqslant  K. \label{3.45}
\end{equation}

下证 \({\left\{  {u}_{k}\right\}  }_{k \in  \mathbb{N}}\) 在 \({L}^{1}\left( \Omega \right)\) 中完全有界,即需证对任给的 \(\varepsilon  > 0\) ,存在有限多个 \({f}_{1},\cdots ,{f}_{m}\)  \(\in  {L}^{1}\left( \Omega \right)\) 使得对每个 \({u}_{k}\) 存在 \(i \in  \{ 1,2,\cdots ,m\}\) 使得
\begin{equation}
{\begin{Vmatrix}{u}_{k} - {f}_{i}\end{Vmatrix}}_{{L}^{1}\left( \Omega \right) } < \varepsilon . \label{3.46}
\end{equation}

首先,由 \({W}_{0}^{1,p}\left( \Omega \right)\) 的定义知,对每个 \(k \in  \mathbb{N}\) 及 \(\varepsilon  > 0\) ,存在 \({g}_{k} \in  {C}_{0}^{1}\left( \Omega \right)\) 及仅依赖于 \(\Omega\) 的常数 \({C}_{0}\) 使得
\begin{equation}
{\begin{Vmatrix}{u}_{k} - {g}_{k}\end{Vmatrix}}_{{L}^{1}\left( \Omega \right) } \leqslant  {C}_{0}{\begin{Vmatrix}{u}_{k} - {g}_{k}\end{Vmatrix}}_{{W}^{1,p}\left( \Omega \right) } < \frac{\varepsilon }{3}. \label{3.47}
\end{equation}

考虑 \({\left\{  {g}_{k}\right\}  }_{k = 1}^{\infty }\) 的光滑化
\begin{equation*}
{g}_{k}^{\delta } \mathrel{\text{ := }} {\int }_{\left| {x - y}\right|  \leqslant  \delta }{g}_{k}\left( y\right) {j}_{\delta }\left( {x - y}\right) {\dif y}
\end{equation*}

其中 \({j}_{\delta }\left( x\right)\) 为例3.3中的磨光核. 则对每个 \({g}_{k} \in  {C}_{0}^{1}\left( \Omega \right)\) 都有
\begin{equation*}
\begin{aligned}
    {g}_{k}^{\delta }\left( x\right)  - {g}_{k}\left( x\right)  &= {\int }_{B\left( {0,1}\right) }{j}_{\delta }\left( y\right) \left( {{g}_{k}\left( {x - {\delta y}}\right)  - {g}_{k}\left( x\right) }\right) {\dif y}\\
&= {\int }_{B\left( {0,1}\right) }{j}_{\delta }\left( y\right) {\int }_{0}^{1}\frac{d}{\dif t}\left( {{g}_{k}\left( {x - {\delta ty}}\right) }\right) {\dif t\dif y}\\&
=  - \delta {\int }_{B\left( {0,1}\right) }{j}_{\delta }\left( y\right) {\int }_{0}^{1}\nabla {g}_{k}\left( {x - {\delta ty}}\right) {y\dif t\dif y}.
\end{aligned}
\end{equation*}

因此

\begin{equation}
{\int }_{\Omega }\left| {{g}_{k}^{\delta }\left( x\right)  - {g}_{k}\left( x\right) }\right| {\dif x} \leqslant  \delta {\int }_{B\left( {0,1}\right) }{j}_{\delta }\left( y\right) {\int }_{0}^{1}{\int }_{\Omega }\left| {\nabla {g}_{k}\left( {x - {\delta ty}}\right) }\right| {\dif x\dif t\dif y}
\leqslant  \delta {\int }_{\Omega }\left| {\nabla {g}_{k}\left( x\right) }\right| {\dif x}. \label{3.48}
\end{equation}

由(3.45)与(3.47)知,存在常数 \({C}_{1}\) 和 \({K}^{\prime }\) 使得
\begin{equation*}
{\begin{Vmatrix}{g}_{k}\end{Vmatrix}}_{{W}^{1,1}\left( \Omega \right) } \leqslant  {C}_{1}{\begin{Vmatrix}{g}_{k}\end{Vmatrix}}_{{W}^{1,p}\left( \Omega \right) } \leqslant  {K}^{\prime },\;\forall k \in  \mathbb{N}.
\end{equation*}

从 (3.48) 知,我们可以选择足够小的 \(\delta  > 0\) 使得
\begin{equation}
{\begin{Vmatrix}{g}_{k} - {g}_{k}^{\delta }\end{Vmatrix}}_{{L}^{1}\left( \Omega \right) } < \frac{\varepsilon }{3},\;\forall n \in  \mathbb{N}. \label{3.49}
\end{equation}

进一步,对 \(x \in  {\mathbb{R}}^{n}\) ,我们有
\begin{equation}
\left| {{g}_{k}^{\delta }\left( x\right) }\right|  \leqslant  {\int }_{B\left( {x,\delta }\right) }\left| {{g}_{k}\left( y\right) }\right| {j}_{\delta }\left( {x - y}\right) {\dif y}
\leqslant  {\begin{Vmatrix}{j}_{\delta }\end{Vmatrix}}_{{L}^{\infty }\left( {\mathbb{R}}^{n}\right) }{\begin{Vmatrix}{g}_{k}\end{Vmatrix}}_{{L}^{1}\left( \Omega \right) } \leqslant  \frac{{C}^{\prime }}{{\delta }^{n}} < \infty , \label{3.50}
\end{equation}

对任意 \(k\) 都成立. 类似地,

\begin{equation}
\left| {\nabla {g}_{k}^{\delta }\left( x\right) }\right|  \leqslant  {\int }_{B\left( {x,\delta }\right) }\left| {\nabla {j}_{\delta }\left( {x - y}\right) }\right| \left| {{g}_{k}\left( y\right) }\right| {\dif y}
\leqslant  {\begin{Vmatrix}\nabla {j}_{\delta }\end{Vmatrix}}_{{L}^{\infty }\left( {\mathbb{R}}^{n}\right) }{\begin{Vmatrix}{g}_{k}\end{Vmatrix}}_{{L}^{1}\left( \Omega \right) } \leqslant  \frac{{C}^{\prime \prime }}{{\delta }^{n + 1}} < \infty , \label{3.51}
\end{equation}

对任意 \(k\) 都成立.

从(3.50)和(3.51)可知,存在常数 \({C}_{2}\) 使得
\begin{equation}
{\begin{Vmatrix}{g}_{k}^{\delta }\end{Vmatrix}}_{{C}^{1}\left( \Omega \right) } \leqslant  {C}_{2}. \label{3.52}
\end{equation}

这里常数 \({C}_{2}\) 依赖 \(\delta\) ,但我们可以选择一个使得 (3.49) 成立的固定 \(\delta\) . 因此, \({\left\{  {g}_{k}^{\delta }\right\}  }_{k \in  \mathbb{N}}\) 在 \({C}^{1}\left( \Omega \right)\) 中一致有界且等度连续. 从而由Arzela-Ascoli定理知, \({\left\{  {g}_{k}^{\delta }\right\}  }_{k \in  \mathbb{N}}\) 包含一致收敛子列. 因此, \(\overline{{\left\{  {g}_{k}^{\delta }\right\}  }_{k \in  \mathbb{N}}}\) 在 \({L}^{1}\left( \Omega \right)\) 中紧致,从而完全有界. 故存在有限多个 \({f}_{1},{f}_{2},\cdots ,{f}_{m} \in  {L}^{1}\left( \Omega \right)\) 使得对任意 \({g}_{k}^{\delta }\) 都存在 \({f}_{i}\) 使得
\begin{equation}
{\begin{Vmatrix}{g}_{k}^{\delta } - {f}_{i}\end{Vmatrix}}_{{L}^{1}\left( \Omega \right) } < \frac{\varepsilon }{3}. \label{3.53}
\end{equation}

由(3.47)(3.49)和(3.53)可知,对任意 \({u}_{k}\) 都存在 \({f}_{i}\) 使得
\begin{equation*}
{\begin{Vmatrix}{u}_{k} - {f}_{i}\end{Vmatrix}}_{{L}^{1}\left( \Omega \right) } < \varepsilon .
\end{equation*}

因此, \({\left\{  {u}_{k}\right\}  }_{k \in  \mathbb{N}}\) 可以被有限多个球
\begin{equation*}
{U}_{i} \mathrel{\text{ := }} \left\{  {h \in  {L}^{1}\left( \Omega \right)  \mid  {\begin{Vmatrix}h - {f}_{i}\end{Vmatrix}}_{{L}^{1}} < \varepsilon }\right\}
\end{equation*}

覆盖. 从而由定理 1.2 得, \(\overline{{\left\{  {u}_{k}\right\}  }_{k \in  \mathbb{N}}}\) 在 \({L}^{1}\left( \Omega \right)\) 中紧致. 因此, \({\left\{  {u}_{k}\right\}  }_{k \in  \mathbb{N}}\) 在 \({L}^{1}\left( \Omega \right)\) 中有收敛子列. 这表明定理结论对 \(q = 1\) 的情形成立.

接下来,考虑任意 \(1 < q < \frac{np}{n - p}\) 的情形. 由插值不等式 (参见附录B中不等式f) 及引理\ref{3.4.2}可得,存在常数 \({C}_{3}\) 使得对任意 \(u \in  {W}_{0}^{1,p}\left( \Omega \right)\) 都有
\begin{equation}
\parallel u{\parallel }_{{L}^{q}\left( \Omega \right) } \leqslant  \parallel u{\parallel }_{{L}^{1}\left( \Omega \right) }^{\theta }\parallel u{\parallel }_{{L}^{\frac{np}{n - p}}\left( \Omega \right) }^{1 - \theta } \leqslant  {C}_{3}\parallel u{\parallel }_{{L}^{1}\left( \Omega \right) }^{\theta }\parallel \nabla u{\parallel }_{{L}^{p}\left( \Omega \right) }^{1 - \theta }, \label{3.54}
\end{equation}

其中 \(\frac{1}{q} = \theta  + \left( {1 - \theta }\right) \left( {\frac{1}{p} - \frac{1}{n}}\right) ,0 < \theta  < 1\) .

我们前面已经证明了 \({\left\{  {u}_{k}\right\}  }_{k \in  \mathbb{N}}\) 在 \({L}^{1}\left( \Omega \right)\) 中有收敛子列,不妨仍记为 \({\left\{  {u}_{k}\right\}  }_{k \in  \mathbb{N}}\) . 因此,对任意 \(\varepsilon  > 0\) ,存在 \(K \in  \mathbb{N}\) 使得当 \(k,l \geq  K\) 时有
\begin{equation*}
{\begin{Vmatrix}{u}_{k} - {u}_{l}\end{Vmatrix}}_{{L}^{1}\left( \Omega \right) } < \varepsilon .
\end{equation*}

从(3.45)和(3.54)可得,存在常数 \({C}_{4}\) 使得
\begin{equation*}
{\begin{Vmatrix}{u}_{k} - {u}_{l}\end{Vmatrix}}_{{L}^{q}\left( \Omega \right) } \leqslant  {C}_{4} \cdot  {\varepsilon }^{\theta }.
\end{equation*}

这意味着 \({\left\{  {u}_{k}\right\}  }_{k \in  \mathbb{N}}\) 为 \({L}^{q}\left( \Omega \right)\) 中Cauchy列. 由 \({L}^{q}\left( \Omega \right)\) 的完备性知,定理结论成立.
\end{proof}

\subsection{习题}

\begin{problemset}

\item 证明当 \(1 \leqslant p < +\infty\) 时，\({W}^{k,p}(\Omega)\) 是可分和自反的.
\begin{proof}
对每个多重指标 $\alpha$（满足 $|\alpha| \le k$），定义线性映射
\[
T_\alpha : W^{k,p}(\Omega) \to L^p(\Omega), \qquad T_\alpha(u) = D^\alpha u.
\]
将所有这些映射组合成
\[
T : W^{k,p}(\Omega) \to \prod_{|\alpha| \le k} L^p(\Omega) \cong \big(L^p(\Omega)\big)^N, 
\qquad T(u) = (D^\alpha u)_{|\alpha| \le k},
\]
其中 $N = \#\{\alpha : |\alpha| \le k\}$ 是有限的.

对任意 $u \in W^{k,p}(\Omega)$，有
\[
\|T(u)\|_{\ell^p(L^p)} 
= \Big( \sum_{|\alpha| \le k} \| D^\alpha u \|_{L^p}^p \Big)^{1/p} 
= \|u\|_{W^{k,p}},
\]
因此 $T$ 是一个等距线性同构到其像 $T(W^{k,p}(\Omega))$ 的映射，而 $T(W^{k,p}(\Omega))$ 是 $\big(L^p(\Omega)\big)^N$ 的闭子空间.

由于当 $1 \le p < \infty$ 时，$L^p(\Omega)$ 是可分的（例如由有理系数的简单函数生成可数稠密子集），有限直积 $\big(L^p(\Omega)\big)^N$ 也是可分的，而闭子空间的可分性由父空间继承.  
因此 $T(W^{k,p}(\Omega))$ 可分，从而 $W^{k,p}(\Omega)$ 可分.

当 $1 < p < \infty$ 时，$L^p(\Omega)$ 是自反的.有限直积 $\big(L^p(\Omega)\big)^N$ 仍是自反的，因为有限个自反 Banach 空间的直积仍然自反.

已知：若 $X$ 是自反 Banach 空间，则任意闭子空间 $Y \subset X$ 也是自反的（例如由双对偶映射的限制可得）.  
因此 $T(W^{k,p}(\Omega))$ 作为自反空间 $\big(L^p(\Omega)\big)^N$ 的闭子空间是自反的.

由于 $W^{k,p}(\Omega)$ 与 $T(W^{k,p}(\Omega))$ 等距同构，故 $W^{k,p}(\Omega)$ 亦自反.
\end{proof}
\item 当 \(\Omega\) 为 \({\mathbb{R}}^{n}\) 中有界区域时，\({W}_{0}^{k,p}(\Omega) \neq {W}^{k,p}(\Omega)\).

\textbf{提示:} 考虑 \(f \equiv 1\) 并利用 Poincaré 不等式.
\begin{proof}
  首先注意常函数 $f\equiv 1$ 在 $\Omega$ 上平凡地属于 $W^{k,p}(\Omega)$，因为对任一多重指标 $\alpha$，若 $|\alpha|\ge1$ 则 $D^\alpha f\equiv0$，而 $f\in L^p(\Omega)$ 当 $\Omega$ 有界且 $1\le p<\infty$.

假设相反，$1\in W^{k,p}_0(\Omega)$. 取 $k\ge1$ 的情形，考虑 $W^{1,p}_0(\Omega)$ 的 Poincaré 不等式：存在常数 $C>0$（仅依赖于 $\Omega,p$）使得对任意 $u\in W^{1,p}_0(\Omega)$ 有
\[
\|u\|_{L^p(\Omega)} \le C\|\nabla u\|_{L^p(\Omega)}.
\]
因为 $W^{k,p}_0(\Omega)\subset W^{1,p}_0(\Omega)$（当 $k\ge1$ 时成立），若 $1\in W^{k,p}_0(\Omega)$ 则 $1\in W^{1,p}_0(\Omega)$，从而
\[
\|1\|_{L^p(\Omega)} \le C\|\nabla 1\|_{L^p(\Omega)} = C\cdot 0 = 0.
\]
然而 $\|1\|_{L^p(\Omega)} = |\Omega|^{1/p}>0$，这与上式矛盾.因此 $1\notin W^{k,p}_0(\Omega)$，从而
\[
W^{k,p}_0(\Omega)\neq W^{k,p}(\Omega).
\]
\end{proof}
\item 设
\[
{H}^{k}({\mathbb{R}}^{n}) = \left\{\, u \in {\mathcal{S}}^{\prime} \;\middle|\; {\partial}^{\alpha}u \in {L}^{2}({\mathbb{R}}^{n}),~|\alpha| \leqslant k \,\right\},
\]
其中范数定义为
\[
\|u\|_{k} := \left( \sum_{|\alpha| \leqslant k} \|{\partial}^{\alpha}u\|_{L^{2}}^{2} \right)^{1/2}.
\]
证明 \({H}^{k}({\mathbb{R}}^{n})\) 是完备的.
\begin{proof}
  设$u_j\in H^k(\mathbb{R}^n)$是柯西列,设\[
  u_j\to g_0,\partial^\alpha u_j \to g_\alpha(L^p(\mathbb{R^n}))  \]
从而对任意 \(\varphi  \in  \mathcal{S}\left( \mathbb{R}^n \right)\) ,当 \(j \rightarrow  \infty\) 时,
\begin{equation*}
\left\langle  {{\partial }^{\alpha }{u}_{j},\varphi }\right\rangle\to 
\left\langle  {{g}_{\alpha },\varphi }\right\rangle  \;   {\left( -1\right) }^{\left| \alpha \right| }\left\langle  {{u}_{j},{\partial }^{\alpha }\varphi }\right\rangle\to{\left( -1\right) }^{\left| \alpha \right| }\left\langle  {{g}_{0},{\partial }^{\alpha }\varphi }\right\rangle  
\end{equation*}

即得 \({g}_{\alpha } = {\partial }^{\alpha }{g}_{0}\) . 由此可得 \({g}_{0} \in  {H}^{k}\left( \Omega \right)\) ,且
\begin{equation*}
{\begin{Vmatrix}{u}_{j} - {g}_{0}\end{Vmatrix}}_{k} \rightarrow  0,\;j \rightarrow  \infty .
\end{equation*}
\end{proof}
\item 证明若 \({kp} > n\)，则存在常数 \(c = c(k,p) > 0\)，使得对任意 \(u,v \in {C}^{\infty}(\bar{\Omega})\) 和 \(f \in {C}^{k}(\mathbb{R})\) 有
\[
\|uv\|_{k,p} \leqslant c\|u\|_{k,p}\|v\|_{L^{\infty}} + \|u\|_{L^{\infty}}\|v\|_{k,p},
\]
以及
\[
\|f \circ u\|_{k,p} \leqslant c\left(\|f\|_{C^{k}} + 1\right)\|u\|_{k,p}.
\]
\begin{proof}
  由Sobolev嵌入定理知,存在$C>0$使得
  \[
  \mathop{\sup }\limits_{\substack{{x \in  \bar{\Omega}} \\  {\left| \alpha \right|  < k-\frac{n}{p} } }}\left| {{\partial }^{\alpha }u\left( x\right) }\right|  \leqslant  C\parallel u{\parallel }_{k,p},  \mathop{\sup }\limits_{\substack{{x \in  \bar{\Omega}} \\  {\left| \alpha \right|  < k-\frac{n}{p} } }}\left| {{\partial }^{\alpha }v\left( x\right) }\right|  \leqslant  C\parallel v{\parallel }_{k,p}.
  \]
  \begin{equation*}
    \begin{aligned}
      \|uv\|_{k,p}=\left(\int_\Omega \sum_{|\alpha|\leqslant k}\left|\partial^\alpha (uv)\right|^p\dif x \right)^{\frac{1}{p}}=\left(\int_\Omega \sum_{|\alpha|\leqslant k}\left|\sum\limits_{\beta\leqslant \alpha}\partial^\beta u\partial^{\alpha-\beta }v\right|^p\dif x \right)^{\frac{1}{p}}
    \end{aligned}
  \end{equation*}
\end{proof}
\item 设 \(\Omega \subset {\mathbb{R}}^{n}\) 为有界区域，\(u \in {W}_{0}^{1,n}\).证明对任意 \(p \in [1,+\infty)\)，有 \(u \in {L}^{p}(\Omega)\).
\begin{proof}
    取$q=\frac{np}{n+p}$,由Sobolev嵌入定理知,
    \[W_0^{1,n}\hookrightarrow W_0^{1,q}\hookrightarrow L^p \]
    
\end{proof}
\item 设 \(\Omega \subset {\mathbb{R}}^{n}\) 为有界区域.证明 \(u \in {W}^{k+1,p}(\Omega)\) 当且仅当 \(u \in {W}^{1,p}(\Omega)\) 且 \({\partial}_{i}u \in {W}^{k,p}(\Omega)\)，其中 \(i = 1,2,\dots,n\).
\begin{proof}
    充分性,显然成立.

    必要性,由Holder不等式可知.
\end{proof}
\item 设 \(\Omega \subset {\mathbb{R}}^{n}\) 为具有光滑边界的有界区域，\(V\) 为 \({\mathbb{R}}^{n}\) 中的有界区域，且 \(\Omega \subsetneqq V\).证明存在有界线性算子
\[
E : {W}^{1,p}(\Omega) \to {W}^{1,p}({\mathbb{R}}^{n})
\]
使得对每个 \(u \in {W}^{1,p}(\Omega)\) 有：

(i) \({Eu} = u\) 在 \(\Omega\) 中几乎处处成立;

(ii) \(\mathrm{supp}(Eu) \subset V\);

(iii) 存在仅依赖于 \(p,\Omega,V\) 的常数 \(C\)，使得
\[
\|Eu\|_{{W}^{1,p}({\mathbb{R}}^{n})} \leqslant C\|u\|_{{W}^{1,p}(\Omega)}.
\]
\begin{proof}

\end{proof}
\item (\textbf{Gagliardo–Nirenberg–Sobolev 不等式})  
设 \(\Omega \subset {\mathbb{R}}^{n}\) 为具有光滑边界的有界区域，\(1 \leqslant p < n\)，\({p}^{*} = \frac{np}{n - p}\).  
则存在仅依赖 \(p,n\) 的常数 \(C\)，使得对任意 \(u \in {C}_{0}^{1}({\mathbb{R}}^{n})\) 有
\[
\|u\|_{L^{p^{*}}({\mathbb{R}}^{n})} \leqslant C\|\nabla u\|_{L^{p}({\mathbb{R}}^{n})}.
\]
\begin{proof}
 首先假定 \(p = 1\) 且 \(u \in  {C}_{0}^{1}\left( \Omega \right)\) . 记 \(x = \left( {{x}_{1},\cdots ,{x}_{n}}\right)\) . 此时,对每个 \(i \in  \{ 1,2,\cdots ,n\}\) 有
\begin{equation*}
\left| {u\left( x\right) }\right|  \leqslant  {\int }_{-\infty }^{{x}_{i}}\left| {{\nabla }_{i}u\left( x\right) }\right| \dif x_i,
\end{equation*}

这里我们用到\(u\) 具有紧支集. 从而
\begin{equation*}
{\left| u\left( x\right) \right| }^{n} \leqslant  \mathop{\prod }\limits_{{i = 1}}^{n}{\int }_{-\infty }^{\infty }\left| {{\nabla }_{i}u}\right| \dif x_i
\end{equation*}

且
\begin{equation*}
{\left| u\left( x\right) \right| }^{\frac{n}{n - 1}} \leqslant  {\left( \mathop{\prod }\limits_{{i = 1}}^{n}{\int }_{-\infty }^{\infty }\left| {\nabla }_{i}u\right| \dif x_i\right) }^{\frac{1}{n - 1}}.
\end{equation*}

不等式两边关于第一个变量积分得
\begin{equation*}
{\int }_{-\infty }^{\infty }{\left| u\left( x\right) \right| }^{\frac{n}{n - 1}}\dif x_1 \leqslant  {\left( {\int }_{-\infty }^{\infty }\left| {\nabla }_{1}u\right| \dif x_1\right) }^{\frac{1}{n - 1}}{\int }_{-\infty }^{\infty }{\left( \mathop{\prod }\limits_{{i = 2}}^{n}{\int }_{-\infty }^{\infty }\left| {\nabla }_{i}u\right| \dif x_i\right) }^{\frac{1}{n - 1}}\dif x_1.
\end{equation*}

因此,由 \({p}_{1} = \cdots  = {p}_{n - 1} = n - 1\) 的广义Hölder不等式得
\begin{equation*}
{\int }_{-\infty }^{\infty }{\left| u\left( x\right) \right| }^{\frac{n}{n - 1}}\dif x_1 \leqslant  {\left( {\int }_{-\infty }^{\infty }\left| {\nabla }_{1}u\right| \dif x_1\right) }^{\frac{1}{n - 1}}{\left( \mathop{\prod }\limits_{{i \neq  1}}{\int }_{-\infty }^{\infty }{\int }_{-\infty }^{\infty }\left| {\nabla }_{i}u\right| \dif x_i\dif x_1\right) }^{\frac{1}{n - 1}}.
\end{equation*}

同理可得
\begin{equation*}
{\int }_{-\infty }^{\infty }{\int }_{-\infty }^{\infty }{\left| u\left( x\right) \right| }^{\frac{n}{n - 1}}\dif x_1\dif x_2 \leqslant  {\left( {\int }_{-\infty }^{\infty }{\int }_{-\infty }^{\infty }\left| {\nabla }_{1}u\right| \dif x_1\dif x_2\right) }^{\frac{1}{n - 1}}
\end{equation*}
\begin{equation*}
\cdot  {\left( {\int }_{-\infty }^{\infty }{\int }_{-\infty }^{\infty }\left| {\nabla }_{2}u\right| \dif x_1\dif x_2\right) }^{\frac{1}{n - 1}}
\end{equation*}
\begin{equation*}
\cdot  {\left( \mathop{\prod }\limits_{{i = 3}}^{n}{\int }_{-\infty }^{\infty }{\int }_{-\infty }^{\infty }{\int }_{-\infty }^{\infty }\left| {\nabla }_{i}u\right| \dif x_i\dif x_1\dif x_2\right) }^{\frac{1}{n - 1}}.
\end{equation*}

依次迭代可得
\begin{equation*}
{\int }_{\Omega }{\left| u\left( x\right) \right| }^{\frac{n}{n - 1}}{\dif x} \leqslant  {\left( \mathop{\prod }\limits_{{i = 1}}^{n}{\int }_{\Omega }\left| {\nabla }_{i}u\right| \dif x\right) }^{\frac{1}{n - 1}}.
\end{equation*}

进而有
\begin{equation*}
\parallel u{\parallel }_{\frac{n}{n - 1}} \leqslant  {\left( \mathop{\prod }\limits_{{i = 1}}^{n}{\int }_{\Omega }\left| {\nabla }_{i}u\right| \dif x\right) }^{\frac{1}{n}} \leqslant  \frac{1}{n}{\int }_{\Omega }\mathop{\sum }\limits_{{i = 1}}^{n}\left| {{\nabla }_{i}u}\right| {\dif x},
\end{equation*}

这里我们使用了几何平均值不大于算术平均值. 因此,由 \({\mathbb{R}}^{n}\) 上范数等价知,存在常数 \({C}^{\prime } = {C}^{\prime }\left( n\right)\) 使得
\begin{equation*}
\parallel u{\parallel }_{\frac{n}{n - 1}} \leqslant  \frac{{C}^{\prime }}{n}\parallel \nabla u{\parallel }_{1}.
\end{equation*}

至此,我们证明了 \(p = 1\) 时引理结论对 \(u \in  {C}_{0}^{1}\left( \Omega \right)\) 成立.

对 \(\gamma  > 1\) ,我们应用\ref{3.27}到 \({\left| u\right| }^{\gamma }\) ,由链式法则可得
\begin{equation*}
{\begin{Vmatrix}{\left| u\right| }^{\gamma }\end{Vmatrix}}_{\frac{n}{n - 1}} \leqslant  \frac{{C}^{\prime }\gamma }{n}{\int }_{\Omega }{\left| u\right| }^{\gamma  - 1}\left| {\nabla u}\right| {\dif x}
\leqslant  \frac{{C}^{\prime }\gamma }{n}{\begin{Vmatrix}{\left| u\right| }^{\gamma  - 1}\end{Vmatrix}}_{q} \cdot  \parallel \nabla u{\parallel }_{p}, 
\end{equation*}

其中 \(\frac{1}{p} + \frac{1}{q} = 1\) ,在第二个不等式中使用了Hölder不等式.
由于 \(p < n\) ,取 \(\gamma  = \frac{\left( {n - 1}\right) p}{n - p}\) ,则有 \(\frac{np}{n - p} = \frac{\gamma n}{n - 1} = \frac{\left( {\gamma  - 1}\right) p}{p - 1}\) . 令 \(q = \frac{p}{p - 1}\) ,则由

可得
\begin{equation*}
\parallel u{\parallel }_{\frac{\gamma n}{n - 1}}^{\gamma } \leqslant  \frac{{C}^{\prime }\gamma }{n}\parallel u{\parallel }_{\frac{\gamma n}{n - 1}}^{\gamma  - 1} \cdot  \parallel \nabla u{\parallel }_{p}.
\end{equation*}

因此,
\begin{equation*}
\parallel u{\parallel }_{\frac{\gamma n}{n - 1}} \leqslant  \frac{{C}^{\prime }\gamma }{n}\parallel \nabla u{\parallel }_{p}.
\end{equation*}

取 \(C = \frac{{C}^{\prime }\gamma }{n}\) 可得引理结论对 \(u \in  {C}_{0}^{1}\left( \Omega \right)\) 成立.

\end{proof}
\begin{note}
    $C_0^1$包含于$W_0^{1,p}$.
\end{note}
\item (\({W}^{1,p}\) 估计)  
设 \(\Omega \subset {\mathbb{R}}^{n}\) 为具有光滑边界的有界区域，\(1 \leqslant p < n\)，\({p}^{*} = \frac{np}{n - p}\).  
若 \(u \in {W}^{1,p}(\Omega)\)，则 \(u \in L^{p^{*}}(\Omega)\)，且存在仅依赖于 \(n,p,\Omega\) 的常数 \(C\)，使得
\[
\|u\|_{L^{p^{*}}(\Omega)} \leqslant C\|u\|_{W^{1,p}(\Omega)}.
\]
\begin{proof}
    由第7题可设$Eu=\widetilde{u},u\in W_0^{1,p}(\mathbb{R}^n)$,由$C_0^\infty$的稠密可知$\exists u_n\in C_0^\infty(\mathbb{R}^n)$,由引理3.2知\[
    \|u_i-u_j\|_{p*}\leqslant C\|\nabla(u_i-u_j)\|_p\to 0(i,j\to \infty)\Rightarrow u_n\to \widetilde{u} \quad \textbf{in} \quad L^{p*}(\mathbb{R}^n)
    \]
    同样有\[\|u_n\|_{p*}\leqslant C \|\nabla u_n\|_p\]
    取极限得到
\[\|\widetilde{u}\|_{p*}\leqslant C\|\nabla \widetilde{u}\|_p\leqslant C'\|u\|_{1,p}\]
而从$Eu$的性质知\[
\|u\|_{L^{p^{*}}(\Omega)} =\|Eu\|_{p*}\leqslant C'\|Eu\|_{1,p}\leqslant C\|u\|_{W^{1,p}(\Omega)}
\]

\end{proof}
\item 设 \(\Omega \subset {\mathbb{R}}^{n}\) 为区域，且 \({kp} > n\).则对每个紧子集 \(K \subset \Omega\)，存在常数 \(C > 0\)，使得对任意 \(u \in {C}_{0}^{\infty}(\Omega)\) 且 \(\mathrm{supp}\,u \subset K\)，有
\[
\sup_{x \in K} |u(x)| \leqslant C\|u\|_{k,p}.
\]
\begin{proof}
    由定理3.10易证.
\end{proof}
\item 设 \(\Omega \subset {\mathbb{R}}^{n}\) 为具有 \({C}^{1}\) 边界的有界区域，且 \(1 < p < \infty\).则存在常数 \(C = C(n,p,\Omega) > 0\)，使得对任意 \(u \in {W}_{0}^{1,p}(\Omega)\) 有
\[
\|u - (u)_{\Omega}\|_{L^{p}(\Omega)} \leqslant C\|\nabla u\|_{L^{p}(\Omega)},
\]
其中 \((u)_{\Omega} = \frac{\int_{\Omega} u(x)\,dx}{\mu(\Omega)}\).
\begin{proof}
假设不成立, 那么对于每一个 \( k = 1,\ldots \) 存在 \( {u}_{k} \in  {W}^{1,p}\left( \Omega\right) \) 满足 \[ {\begin{Vmatrix}{u}_{k} - {\left( {u}_{k}\right) }_{\Omega}\end{Vmatrix}}_{{L}^{p}\left( \Omega\right) } > k{\begin{Vmatrix}D{u}_{k}\end{Vmatrix}}_{{L}^{p}\left( \Omega\right) }. \]
我们定义
\[ {v}_{k} \mathrel{\text{:=}} \frac{{u}_{k} - {\left( {u}_{k}\right) }_{\Omega}}{{\begin{Vmatrix}{u}_{k} - {\left( {u}_{k}\right) }_{\Omega}\end{Vmatrix}}_{{L}^{p}\left( \Omega\right) }}\;\left( {k = 1,\ldots }\right) . \]
那么
\[ {\left( {v}_{k}\right) }_{\Omega} = 0,{\begin{Vmatrix}{v}_{k}\end{Vmatrix}}_{{L}^{p}\left( \Omega\right) } = 1,{\begin{Vmatrix}D{v}_{k}\end{Vmatrix}}_{{L}^{p}\left( \Omega\right) } < \frac{1}{k}\;\left( {k = 1,2,\ldots }\right) . \]
由于 \( {\left\{  {v}_{k}\right\}  }_{k = 1}^{\infty } \) 在 \( {W}^{1,p}\left( \Omega\right) \) 中有界,那么由Rellich-Kondrachov 紧嵌入定理知,
存在 \( {\left\{  {v}_{{k}_{j}}\right\}  }_{j = 1}^{\infty } \subset  {\left\{  {v}_{k}\right\}  }_{k = 1}^{\infty } \) 和 \( v \in  {L}^{p}\left( \Omega\right) \)使得 \( {v}_{{k}_{j}} \rightarrow  v\; \) in \( {L}^{p}\left( \Omega\right) . \)
同时\( {\left( v\right) }_{\Omega} = 0,\parallel v{\parallel }_{{L}^{p}\left( \Omega\right) } = 1. \)

另一方面,对于任意 \( i = 1,\ldots ,n \) 和\( \phi  \in  {C}_{0}^{\infty }\left( \Omega\right) \)有
\[ {\int }_{\Omega}v{\phi }_{{x}_{i}}{dx} = \mathop{\lim }\limits_{{{k}_{j} \rightarrow  \infty }}{\int }_{\Omega}{v}_{{k}_{j}}{\phi }_{{x}_{i}}{dx} =  - \mathop{\lim }\limits_{{{k}_{j} \rightarrow  \infty }}{\int }_{\Omega}{v}_{{k}_{j},{x}_{i}}{\phi dx} = 0. \]
因此 \( v \in  {W}^{1,p}\left( \Omega\right) \) , \( {Dv} = 0 \) 几乎处处. 由推论3.4知\( v \) 是常数.由\( v \) 是常数且 \( {\left( v\right) }_{\Omega} = 0 \) 知\( v \equiv  0 \) ,进而
\( \parallel v{\parallel }_{{L}^{p}\left( \Omega\right) } = 0 \) . 这与$\|v\|_{L^p(\Omega)}=1$矛盾.
\end{proof}
\item 设 \(B(x,r)\) 为以 \(x\) 为心、\(r\) 为半径的球，\(1 \leqslant p \leqslant \infty\).则存在仅依赖于 \(n,p\) 的常数 \(C\)，使得对任意 \(u \in {W}_{0}^{1,p}(B^{0}(x,r))\) 都有
\[
\|u - (u)_{B(x,r)}\|_{L^{p}(B(x,r))} \leqslant Cr\|\nabla u\|_{L^{p}(B(x,r))}.
\]
\begin{proof}
    取上题中的$\Omega=B^0(x,1),v(y)=u(x+ry),y\in B^0(x,1)$,那么$v\in {W}_{0}^{1,p}(B^{0}(x,r))$,有
    \[
    \|v-(v)_\Omega\|_{L^p(\Omega)}\leqslant C\|\nabla v\|_{L^p(\Omega)}\Rightarrow \|u - (u)_{B(x,r)}\|_{L^{p}(B(x,r))} \leqslant Cr\|\nabla u\|_{L^{p}(B(x,r))}.
    \]
    
\end{proof}
\end{problemset}
\chapter{紧算子与Fredholm算子理论}

\section{紧算子的定义及基本性质}
\begin{definition}
    设 \(\mathcal{X},\mathcal{Y}\) 是Banach空间,称线性算子 \(T \in  \mathcal{L}\left( {\mathcal{X},\mathcal{Y}}\right)\) 为紧算子,如果 \(T\) 映 \(\mathcal{X}\) 中有界集为 \(\mathcal{Y}\) 中相对紧集,即若 \(A\) 为 \(\mathcal{X}\) 中任意有界集,则 \(T\left( A\right)\) 中任意点列必有在 \(\mathcal{Y}\) 中收敛的子列. 记从 \(\mathcal{X}\) 到 \(\mathcal{Y}\) 的紧算子的全体为 \(\mathcal{C}\left( {\mathcal{X},\mathcal{Y}}\right)\) . 特别, 当 \(\mathcal{X} = \mathcal{Y}\) 时,记为 \(\mathcal{C}\left( \mathcal{X}\right)\) .
\end{definition}
\begin{note}
    \(T \in  \mathcal{C}\left( {\mathcal{X},\mathcal{Y}}\right)\) 当且仅当 \(T\) 映 \(\mathcal{X}\) 中单位球为 \(\mathcal{Y}\) 中相对紧集.

     \(T \in  \mathcal{C}\left( {\mathcal{X},\mathcal{Y}}\right)\) 当且仅当对 \(\mathcal{X}\) 中任意有界点列 \(\left\{  {x}_{n}\right\}  ,\left\{  {T{x}_{n}}\right\}\) 有收敛子列.

     \(\mathcal{C}\left( {\mathcal{X},\mathcal{Y}}\right)  \subset  \mathcal{L}\left( {\mathcal{X},\mathcal{Y}}\right)\) .
\end{note}
由于任意有限维Banach空间的有界集都是相对紧集, 故有
\begin{example}
     设 \(\mathcal{X}\) 为Banach空间,值域为有限维的有界线性算子是紧算子,通常称为有限秩算子.
\end{example}
\begin{theorem}
     设 \(\mathcal{X},\mathcal{Y},\mathcal{Z}\) 为Banach 空间,则

(i) 有限个紧算子的和还是紧算子,即若 \({T}_{1},{T}_{2},\cdots ,{T}_{n} \in  \mathcal{C}\left( {\mathcal{X},\mathcal{Y}}\right)\) ,则 \({T}_{1} + {T}_{2} +\)  \(\cdots  + {T}_{n} \in  \mathcal{C}\left( {\mathcal{X},\mathcal{Y}}\right)\) ;

(ii) 紧算子与有界线性算子的复合仍为紧算子,即若 \(T \in  \mathcal{L}\left( {\mathcal{X},\mathcal{Y}}\right) ,S \in  \mathcal{L}\left( {\mathcal{Y},\mathcal{Z}}\right)\) 中有一个为紧算子,则有 \(S \circ  T \in  \mathcal{C}\left( {\mathcal{X},\mathcal{Z}}\right)\) ;

(iii) 若 \(\left\{  {T}_{n}\right\}   \subset  \mathcal{C}\left( {\mathcal{X},\mathcal{Y}}\right)\) 且 \({\begin{Vmatrix}{T}_{n} - T\end{Vmatrix}}_{\mathcal{L}\left( {\mathcal{X},\mathcal{Y}}\right) } \rightarrow  0\left( {n \rightarrow  \infty }\right)\) ,则 \(T\) 也是紧算子.

\end{theorem}

\begin{proof}
    (i) 和 (ii) 都是显然的.

(iii) 要证 \(T\) 为紧算子只需证明 \(T\) 将闭单位球映成相对紧集即可. 记 \({B}_{1} = \{ x \in\)  \(\mathcal{X} \mid  \parallel x\parallel  \leqslant  1\}\) 为 \(\mathcal{X}\) 中单位球. 由于 \(\begin{Vmatrix}{{T}_{n} - T}\end{Vmatrix} \rightarrow  0\left( {n \rightarrow  \infty }\right)\) 知,对任意 \(\varepsilon  > 0\) ,存在 \(N > 0\) 使得对任意 \(x \in  {B}_{1}\) 及 \(n > N\) 有 \(\begin{Vmatrix}{{T}_{n}x - {Tx}}\end{Vmatrix} < \varepsilon /4\) . 由于 \(\left\{  {T}_{n}\right\}   \subset  \mathcal{C}\left( {\mathcal{X},\mathcal{Y}}\right)\) , 故取 \({n}_{0} \geq  N\) ,则有 \(\overline{{T}_{{n}_{0}}\left( {B}_{1}\right) }\) 为完全有界集. 记它的有穷 \(\frac{\varepsilon }{4}\) - 网为 \(\left\{  {{y}_{1},{y}_{2},\cdots ,{y}_{m}}\right\}\) ,则有
\begin{equation*}
\overline{T\left( {B}_{1}\right) } \subset  \mathop{\bigcup }\limits_{{i = 1}}^{m}B\left( {{y}_{i},\varepsilon/2 }\right) .
\end{equation*}

注意这里的$y_i$不一定属于$\overline{T\left( {B}_{1}\right) } $,不妨设$z_i \in \overline{T\left( {B}_{1}\right) } \bigcap B\left( {{y}_{i},\varepsilon/2 }\right)  $,那么\begin{equation*}
\overline{T\left( {B}_{1}\right) } \subset  \mathop{\bigcup }\limits_{{i = 1}}^{m}B\left( {{z}_{i},\varepsilon/2 }\right) .
\end{equation*}

即 \(\overline{T\left( {B}_{1}\right) }\) 有有穷的 \(\varepsilon\) -网,故 \(\overline{T\left( {B}_{1}\right) }\) 为紧集.
\end{proof}
\begin{note}
    从 (i) 和 (ii) 不难得知, \(\mathcal{C}\left( {\mathcal{X},\mathcal{Y}}\right)\) 为线性空间.
\end{note}
\begin{definition}
    设 \(\mathcal{X},\mathcal{Y}\) 是Banach空间,称 \(T \in  \mathcal{L}\left( {\mathcal{X},\mathcal{Y}}\right)\) 为全连续的,如果
\begin{equation*}
{x}_{n} \rightharpoonup  x \Rightarrow  T{x}_{n} \rightarrow  {Tx}.
\end{equation*}

\end{definition}
\begin{theorem}
    设 \(\mathcal{X},\mathcal{Y}\) 是Banach空间. 若 \(T \in  \mathcal{C}\left( {\mathcal{X},\mathcal{Y}}\right)\) ,则 \(T\) 是全连续的; 反之, 若 \(\mathcal{X}\) 是自反的且 \(T\) 是全连续的,则 \(T \in  \mathcal{C}\left( {\mathcal{X},\mathcal{Y}}\right)\) .

\end{theorem}
\begin{proof}
    首先,设 \({x}_{n} \rightharpoonup  x\) ,往证 \(T{x}_{n} \rightarrow  {Tx} = y\) . 用反证法,若不然,则存在 \({\varepsilon }_{0} > 0\) 及 \(\left\{  {n}_{i}\right\}\) 使得
\begin{equation*}
\begin{Vmatrix}{T{x}_{{n}_{i}} - y}\end{Vmatrix} \geq  {\varepsilon }_{0}.
\end{equation*}

由共鸣定理知, \(\left\{  {x}_{n}\right\}\) 有界($x_n^{**}(f)\to x^{**}(f)\Rightarrow \sup\limits_n\|x_n^{**}\|<\infty$); 又由 \(T\) 紧,从而可从 \(\left\{  {x}_{{n}_{i}}\right\}\) 中抽取子列,不妨仍记为 \(\left\{  {x}_{{n}_{i}}\right\}\) ,使得 \(T{x}_{{n}_{i}} \rightarrow  z\) ,但
\begin{equation*}
\left\langle  {{y}^{ * },T{x}_{{n}_{i}} - y}\right\rangle   = \left\langle  {{T}^{ * }{y}^{ * },{x}_{{n}_{i}} - x}\right\rangle   \rightarrow  0,\;\forall {y}^{ * } \in  {\mathcal{Y}}^{ * },
\end{equation*}

即 \(T{x}_{{n}_{i}} \rightharpoonup  y\) . 从而 \(y = z\) ,矛盾,故定理第一个结论得证.

反之,由于自反空间的单位闭球是弱自列紧的,因此,若 \(\left\{  {x}_{n}\right\}\) 有界,则必有子列 \({x}_{{n}_{i}} \rightharpoonup  x\) . 由 \(T\) 全连续得 \(T{x}_{{n}_{i}} \rightarrow  {Tx}\) ,故 \(T \in  \mathcal{C}\left( {\mathcal{X},\mathcal{Y}}\right)\) .

\end{proof}
\begin{theorem}
    设 \(\mathcal{X},\mathcal{Y}\) 是Banach空间,则 \(T \in  \mathcal{C}\left( {\mathcal{X},\mathcal{Y}}\right)\) 当且仅当 \({T}^{ * } \in  \mathcal{C}\left( {{\mathcal{Y}}^{ * },{\mathcal{X}}^{ * }}\right)\) .
\end{theorem}
\begin{proof}
    必要性. 记 \({B}_{1}^{ * }\) 为 \({\mathcal{Y}}^{ * }\) 中单位球. 则要证: 若 \({y}_{n}^{ * } \in  {B}_{1}^{ * }\) ,则 \(\left\{  {{T}^{ * }{y}_{n}^{ * }}\right\}\) 中有收敛子列.\[
\|T^*y_n^*\|\leqslant \sup\limits_{\|x\|\leqslant 1}\left|\left\langle T^*y_n^*, x\right\rangle\right|\leqslant \sup\limits_{y\in T(B_1)}\left|\left\langle y_n^*, y\right\rangle\right|  
    \]
    
    对任意 \(n \in  \mathbb{N}\) ,令
\begin{equation*}
{\varphi }_{n}\left( y\right)  \mathrel{\text{ := }} \left\langle  {{y}_{n}^{ * },y}\right\rangle  ,\;\forall y \in  \overline{T\left( {B}_{1}\right) }.
\end{equation*}

显然 \({\varphi }_{n} \in  C\left( \overline{T\left( {B}_{1}\right) }\right) \left( {\forall n \in  \mathbb{N}}\right)\) . 因此,我们只需证明 \(\left\{  {\varphi }_{n}\right\}\) 作为 \(C\left( \overline{T\left( {B}_{1}\right) }\right)\) 上的函数列有一致收敛的子列就够了. 实际上, 我们有
\begin{equation*}
\left| {{\varphi }_{n}\left( y\right) }\right|  \leqslant  \begin{Vmatrix}{y}_{n}^{ * }\end{Vmatrix}\parallel y\parallel  \leqslant  \parallel T\parallel ,\;\forall n \in  \mathbb{N},\forall y \in  \overline{T\left( {B}_{1}\right) }
\end{equation*}

及
\begin{equation*}
\left| {{\varphi }_{n}\left( y\right)  - {\varphi }_{n}\left( z\right) }\right|  \leqslant  \begin{Vmatrix}{y}_{n}^{ * }\end{Vmatrix}\parallel y - z\parallel  \leqslant  \parallel y - z\parallel ,\;\forall n \in  \mathbb{N},\forall y,z \in  \overline{T\left( {B}_{1}\right) }.
\end{equation*}

上述二式表明 \(\left\{  {\varphi }_{n}\right\}\) 是一致有界和等度连续的. 由Arzela-Ascoli定理知, \(\left\{  {\varphi }_{n}\right\}\) 中有子列在 \(C\left( \overline{T\left( {B}_{1}\right) }\right)\) 中收敛,即得 \(\left\{  {{T}^{ * }{y}_{n}^{ * }}\right\}\) 有收敛子列.

充分性. 由必要性的结论知, \({T}^{* * } \in  \mathcal{C}\left( {{\mathcal{X}}^{* * },{\mathcal{Y}}^{* * }}\right)\) ,但 \(T = {\left. {T}^{* * }\right| }_{\mathcal{X}}\) . 由于 \(\mathcal{X}\) 可视为 \({\mathcal{X}}^{* * }\) 的闭线性子空间,故有 \(T \in  \mathcal{C}\left( {\mathcal{X},\mathcal{Y}}\right)\) .

\end{proof}
\begin{example}
     设 \(\Omega  \subset  {\mathbb{R}}^{n}\) 是一个有界闭集, \(K \in  C\left( {\Omega  \times  \Omega }\right)\) ,取 \(\mathcal{X} = C\left( \Omega \right)\) ; 若令
\begin{equation*}
\begin{aligned} T : C\left( \Omega \right) &  \rightarrow  C\left( \Omega \right) \\  u &  \mapsto  {\int }_{\Omega }K\left( {x,y}\right) u\left( y\right) {\dif y}, \end{aligned}
\end{equation*}

则 \(T \in  \mathcal{C}\left( \mathcal{X}\right)\) .

\end{example}
\begin{proof}
    只需证 \(\overline{T\left( {B}_{1}\right) }\) 为紧集即可. 为此应用Arzela-Ascoli定理. 若记
\begin{equation*}
M \mathrel{\text{ := }} \mathop{\max }\limits_{{x,y \in  \Omega }}\left| {K\left( {x,y}\right) }\right|
\end{equation*}

则有 \(\parallel {Tu}\parallel  \leqslant  M\parallel u{\parallel }_{{C}^{0}\left( \bar{\Omega }\right) }\mu \left( \Omega \right)\) . 由 \(K\left( {x,y}\right)\) 在 \(\Omega  \times  \Omega\) 中的一致连续性,故对任意 \(\varepsilon  > 0\) ,存在 \(\delta  > 0\) ,使得对任意 \(y \in  \Omega\) 有当 \(\left| {x - {x}^{\prime }}\right|  < \delta\) 时
\begin{equation*}
\left| {K\left( {x,y}\right)  - K\left( {{x}^{\prime },y}\right) }\right|  < \varepsilon ,
\end{equation*}

从而
\begin{equation*}
\left| {\left( {Tu}\right) \left( x\right)  - \left( {Tu}\right) \left( {x}^{\prime }\right) }\right|  \leqslant  {\int }_{\Omega }\left| {K\left( {x,y}\right)  - K\left( {{x}^{\prime },y}\right) }\right| \left| {u\left( y\right) }\right| {\dif y} \leqslant  \varepsilon \parallel u{\parallel }_{{C}^{0}\left( \bar{\Omega }\right) }\mu \left( \Omega \right) .
\end{equation*}
\end{proof}
\begin{note}
Banach空间不需要有内积.
\end{note}
\begin{example}
    设 \(\Omega  \subset  {\mathbb{R}}^{n}\) 是一个有界开区域,又设 \(T \in  \mathcal{L}\left( {{L}^{2}\left( \Omega \right) ,{W}_{0}^{1,2}\left( \Omega \right) }\right)\) 满足
\begin{equation*}
\parallel {Tu}{\parallel }_{1,2} \leqslant  C\parallel u{\parallel }_{{L}^{2}\left( \Omega \right) },
\end{equation*}

其中 \(C\) 是一个常数. \(\iota  : {W}_{0}^{1,2}\left( \Omega \right)  \hookrightarrow  {L}^{2}\left( \Omega \right)\) 为Sobolev嵌入. 那么 \({T}^{\prime } \mathrel{\text{ := }} T \circ  \iota  \in\)  \(\mathcal{C}\left( {{W}_{0}^{1,2}\left( \Omega \right) }\right)\) .

\end{example}
\begin{proof}
    由Sobolev紧嵌入定理3.12知, \(\iota  : {W}_{0}^{1,2}\left( \Omega \right)  \rightarrow  {L}^{2}\left( \Omega \right)\) 是紧嵌入,又 \({T}^{\prime }\) : \(\iota \left( {{W}_{0}^{1,2}\left( \Omega \right) }\right)  \subset  {L}^{2}\left( \Omega \right)  \rightarrow  {W}_{0}^{1,2}\left( \Omega \right)\) 连续,从而由定理 4.1 知 \({T}^{\prime } = T \circ  \iota  \in  \mathcal{C}\left( {X,Y}\right)\) .

\end{proof}
\begin{example}
     设 \(\Omega  \subset  {\mathbb{R}}^{n}\) 为有界集,函数 \(k : \bar{\Omega } \times  \bar{\Omega } \rightarrow  \mathbb{C}\) 满足条件: 当 \(x \neq\)  \(y\) 时 \(k\left( {x,y}\right)\) 是 \(\bar{\Omega } \times  \bar{\Omega }\) 上二元连续函数; 当 \(x = y\) 时, \(k\) 有弱奇性: 存在常数 \(m > 0\) 使得 \(\left| {k\left( {x,y}\right) }\right|  \leqslant  m/{\left| x - y\right| }^{\alpha },\alpha  \in  \left( {0,n}\right)\) . 则积分算子
\begin{equation}
{Kf} = {\int }_{\Omega }k\left( {x,y}\right) f\left( y\right) {\dif y} \label{4.1}
\end{equation}

是 \(C\left( \bar{\Omega }\right)\) 中紧算子.
\end{example}
\begin{proof}
    首先,由前面例题知,当 \(k : \bar{\Omega } \times  \bar{\Omega } \rightarrow  \mathbb{C}\) 是连续函数时, \(K\) 为紧算子.

当 \(k\) 有弱奇性时,令
\begin{equation*}
{k}_{j}\left( {x,y}\right)  \mathrel{\text{ := }} \left\{  \begin{array}{ll} k\left( {x,y}\right) , & \text{ 当 }\left| {x - y}\right|  > 1/j \\  {j}^{\beta }{\left| x - y\right| }^{\beta }k\left( {x,y}\right) , & \text{ 当 }\left| {x - y}\right|  \leqslant  1/j,\beta  \in  \left( {\alpha ,n}\right) . \end{array}\right.
\end{equation*}

记以 \({k}_{j}\left( {x,y}\right)\) 为核产生的积分算子为 \({K}_{j}\) ,则算子列 \({K}_{j}\) 按 \(\mathcal{L}\left( {C\left( \bar{\Omega }\right) }\right)\) 范数收敛到算子 \(K\) . 实际上, 我们有
\begin{align*}
&\left| {{\int }_{\Omega }\left\lbrack  {k\left( {x,y}\right)  - {k}_{j}\left( {x,y}\right) }\right\rbrack  {\dif y}}\right|  \leqslant  {\int }_{\Omega }\left| {k\left( {x,y}\right)  - {k}_{j}\left( {x,y}\right) }\right| {\dif y}
\\=& {\int }_{\left| {x - y}\right|  > 1/j}\left| {k\left( {x,y}\right)  - {k}_{j}\left( {x,y}\right) }\right|  + {\int }_{\left| {x - y}\right|  \leqslant  1/j}\left| {k\left( {x,y}\right)  - {k}_{j}\left( {x,y}\right) }\right|\\= &{\int }_{\left| {x - y}\right|  \leqslant  1/j}\left| {k\left( {x,y}\right)  - {k}_{j}\left( {x,y}\right) }\right|\leqslant  c{\int }_{\left| {x - y}\right|  \leqslant  1/j}\frac{\dif y}{{\left| x - y\right| }^{\alpha }} = c{\int }_{\left| z\right|  \leqslant  1/j}\frac{\dif z}{{\left| z\right| }^{\alpha }}\leqslant  c{\int }_{0}^{1/j}{r}^{n - 1 - \alpha }{\dif r}
\end{align*}


其中 \(c\) 表示常数. 因为 \({\dif z} = {r}^{n - 1}{\dif r\dif S}\) ,当 \(j \rightarrow  \infty\) 时,上式右端趋于零. 从而有
\begin{equation*}
\mathop{\lim }\limits_{{j \rightarrow  \infty }}\mathop{\sup }\limits_{\bar{\Omega }}{\int }_{\Omega }\left| {k\left( {x,y}\right)  - {k}_{j}\left( {x,y}\right) }\right| {\dif y} = 0,
\end{equation*}

因此,当 \(j \rightarrow  \infty\) 时有 \(\begin{Vmatrix}{K - {K}_{j}}\end{Vmatrix} \rightarrow  0\) .
\end{proof}
\begin{note}
    这里$\|K\|=\sup\limits_{x\in\bar{\Omega}}\int_\Omega |k(x,y)|\dif y$.
\end{note}
\begin{note}
    当 \(\alpha  = n\) 时由 (4.1) 定义的积分算子称为奇异积分算子,它是连续算子, 但不是紧算子.
\end{note}
\begin{example}
    设 \(\Omega  \subset  {\mathbb{R}}^{n}\) 为带光滑边界的有界区域. \(\bigtriangleup  = \frac{{\partial }^{2}}{\partial {x}_{1}^{2}} + \cdots  + \frac{{\partial }^{2}}{\partial {x}_{n}^{2}}\) . 由偏微分方程理论知,对每个 \(f \in  {C}^{\infty }\left( \Omega \right)\) ,边值问题
\begin{equation}
\left\{  \begin{array}{ll} \bigtriangleup u\left( x\right)  = f\left( x\right) , & x \in  \Omega , \\  u = 0, & x \in  \partial \Omega  \end{array}\right.  \label{4.2}
\end{equation}

有唯一解 \(u\) . 若记 \(u = {Sf}\) . 则 \(S : {L}^{2}\left( \Omega \right)  \rightarrow  {L}^{2}\left( \Omega \right)\) 为紧算子.

\end{example}
\begin{proof}
    在\ref{4.2}中方程两边乘以 \(u\) ,并在 \(\Omega\) 上积分,由分部积分法得
\begin{equation}
- {\int }_{\Omega }\sum {\left| \frac{\partial u}{\partial {x}_{j}}\right| }^{2}{\dif x} = {\int }_{\Omega }{fu\dif x}. \label{4.3}
\end{equation}

由推论4.3(Poincaré不等式)知,存在常数 \(C = C\left( \Omega \right)  > 0\) 使得
\begin{equation}
\parallel u{\parallel }_{{L}^{2}\left( \Omega \right) } \leqslant  C\parallel \nabla u{\parallel }_{{L}^{2}\left( \Omega \right) }. \label{4.4}
\end{equation}

在\ref{4.3}右边使用Schwarz不等式, 并由\ref{4.4}得
\begin{equation*}
\parallel \nabla u{\parallel }_{{L}^{2}\left( \Omega \right) }^{2} \leqslant  \parallel f{\parallel }_{{L}^{2}\left( \Omega \right) }\parallel u{\parallel }_{{L}^{2}\left( \Omega \right) } \leqslant  C\parallel f{\parallel }_{{L}^{2}\left( \Omega \right) }\parallel \nabla u{\parallel }_{{L}^{2}\left( \Omega \right) }.
\end{equation*}

因此,
\begin{equation}
\parallel \nabla u{\parallel }_{{L}^{2}\left( \Omega \right) } \leqslant  C\parallel f{\parallel }_{{L}^{2}\left( \Omega \right) },\;\parallel u{\parallel }_{{L}^{2}\left( \Omega \right) } \leqslant  {C}^{2}\parallel f{\parallel }_{{L}^{2}\left( \Omega \right) }. \label{4.5}
\end{equation}

用 \({B}_{1}\) 记 \({L}^{2}\left( \Omega \right)\) 中单位球,则
\begin{equation*}
S\left( {B}_{1}\right)  = \left\{  {u \in  {L}^{2}\left( \Omega \right)  \mid  \bigtriangleup u = f\text{ 且 }\parallel f{\parallel }_{{L}^{2}\left( \Omega \right) } \leqslant  1}\right\}  .
\end{equation*}

因此,对任意 \(u \in  S\left( {B}_{1}\right)\) 有
\begin{equation}
\parallel \nabla u{\parallel }_{{L}^{2}\left( \Omega \right) } \leqslant  C,\;\parallel u{\parallel }_{{L}^{2}\left( \Omega \right) } \leqslant  {C}^{2}. \label{4.6}
\end{equation}

由定理 1.11 知, \({L}^{2}\left( \Omega \right)\) 中满足 (4.6) 的函数集合为相对紧集,即有 \(S\) 为紧算子.
\end{proof}
\begin{example}
    设 \(\Omega  \subset  {\mathbb{R}}^{n}\) 为有界区域. \(\bigtriangleup  = \frac{{\partial }^{2}}{\partial {x}_{1}^{2}} + \cdots  + \frac{{\partial }^{2}}{\partial {x}_{n}^{2}}\) . 从偏微分方程理论知,抛物方程的初边值问题
\begin{equation}
\left\{  \begin{array}{ll} {u}_{t}\left( {x,t}\right)  = \bigtriangleup u\left( {x,t}\right) , & \left( {x,t}\right)  \in  \Omega  \times  \lbrack 0, + \infty ) \\  u\left( {x,0}\right)  = g\left( x\right) , & x \in  \Omega \\  u\left( {x,t}\right)  = 0, & \left( {x,t}\right)  \in  \partial \Omega  \times  \lbrack 0, + \infty ) \end{array}\right.  \label{4.7}
\end{equation}

具有唯一解. 对每个 \(T \leqslant  0\) ,定义算子 \(S\left( T\right)  : {L}^{2}\left( \Omega \right)  \rightarrow  {L}^{2}\left( \Omega \right)\) 如下:
\begin{equation}
S\left( T\right) \left( {u\left( {x,0}\right) }\right)  \mathrel{\text{ := }} u\left( {x,T}\right) , \label{4.8}
\end{equation}

其中 \(u\left( {x,t}\right)\) 为初边值问题\ref{4.7}的唯一解. 证明对任意 \(T > 0\) ,算子 \(S\left( T\right)\) 为紧算子.
\end{example}
\begin{proof}
    在方程\ref{4.7}两边乘以 \(u\) ,关于变量 \(x\) 在 \(\Omega\) 上积分,关于变量 \(t\) 在 \(\left\lbrack  {0,T}\right\rbrack\) 上积分, 由分部积分法得
\begin{equation}
{\left. \frac{1}{2}\left( {\int }_{\Omega }{u}^{2}\dif x\right) \right| }_{0}^{T} = {\int }_{0}^{T}{\int }_{\Omega }u\bigtriangleup {u\dif x\dif t} =  - {\int }_{0}^{T}{\int }_{\Omega }\mathop{\sum }\limits_{{j = 1}}^{n}{\left| \frac{\partial u}{\partial {x}_{j}}\right| }^{2}{\dif x\dif t} \leqslant  0. \label{4.9}
\end{equation}

从而有
\begin{equation*}
{\int }_{\Omega }{u}^{2}\left( {x,T}\right) {\dif x} \leqslant  {\int }_{\Omega }{u}^{2}\left( {x,0}\right) {\dif x}
\end{equation*}

及 \(\parallel u\left( t\right) {\parallel }_{{L}^{2}\left( \Omega \right) }\) 为 \(t\) 的单调递减函数. 由 \(S\left( T\right)\) 的定义知
\begin{equation}
\parallel S\left( T\right) \parallel  \leqslant  1\text{ . } \label{4.10}
\end{equation}

在方程\ref{4.7}两边乘以 \(t\bigtriangleup u\) ,关于变量 \(x\) 在 \(\Omega\) 上积分,关于变量 \(t\) 在 \(\left\lbrack  {0,T}\right\rbrack\) 上积分得
\begin{equation}
{\int }_{0}^{T}{\int }_{\Omega }t{u}_{t}\bigtriangleup {u\dif x\dif t} = {\int }_{0}^{T}{\int }_{\Omega }t{\left( \bigtriangleup u\right) }^{2}{\dif x\dif t}. \label{4.11}
\end{equation}

由于上式右边非负,在左边使用分部积分法可得
\begin{equation*}
\frac{T}{2}{\int }_{\Omega }\mathop{\sum }\limits_{{j = 1}}^{n}{\left| \frac{\partial u}{\partial {x}_{j}}\left( T\right) \right| }^{2}{\dif x} \leqslant  \frac{1}{2}{\int }_{0}^{T}{\int }_{\Omega }\mathop{\sum }\limits_{{j = 1}}^{n}{\left| \frac{\partial u}{\partial {x}_{j}}\right| }^{2}{\dif x\dif t}.
\end{equation*}

由 \ref{4.9}及上式可得
\begin{equation*}
T{\int }_{\Omega }\mathop{\sum }\limits_{{j = 1}}^{n}{\left| \frac{\partial u}{\partial {x}_{j}}\left( T\right) \right| }^{2}{\dif x} \leqslant  \frac{1}{2}{\int }_{\Omega }{u}^{2}\left( 0\right) {\dif x}.
\end{equation*}

由定理 1.6 知, 所证结论成立.
\end{proof}

从定理4.1 (iii) 知, 有穷秩算子的极限必为紧算子. Banach 和Grothendick 曾考虑过该结论的逆,即著名的 “逼近问题” : 任给一个 \(T \in  \mathcal{C}\left( {\mathcal{X},\mathcal{Y}}\right)\) ,是否存在有穷秩算子列 \({\left\{  {T}_{n}\right\}  }_{1}^{\infty }\) 使得 \({\begin{Vmatrix}{T}_{n} - T\end{Vmatrix}}_{\mathcal{L}\left( {\mathcal{X},\mathcal{Y}}\right) } \rightarrow  0\left( {n \rightarrow  \infty }\right)\) . P. Enflo \(\left\lbrack  \mathrm{{En}}\right\rbrack\) 在1972年找到一个反例. 关于逼近问题的更详细资料可参见[LT1, LT2]. 这里, 我们将在Banach 空间 \(\mathcal{X},\mathcal{Y}\) 满足一定条件时的逼近问题.
\begin{definition}
    设 \(T \in  \mathcal{L}\left( {\mathcal{X},\mathcal{Y}}\right)\) . 若 \(\dim R\left( T\right)  < \infty\) ,则称 \(T\) 是有穷秩算子. 记有穷秩算子的全体为 \(F\left( {\mathcal{X},\mathcal{Y}}\right)\) . 显然有
\begin{equation*}
F\left( {\mathcal{X},\mathcal{Y}}\right)  \subset  \mathcal{C}\left( {\mathcal{X},\mathcal{Y}}\right) .
\end{equation*}
\end{definition}
\begin{definition}
    设 \(f \in  {\mathcal{X}}^{ * },y \in  \mathcal{Y}\) ,用 \(y \otimes  f\) 表示下列算子:
\begin{equation*}
x \mapsto  \langle f,x\rangle y,\;\forall x \in  \mathcal{X},
\end{equation*}

不难看出, 上述算子为秩 1 算子.
\end{definition}
\begin{theorem}
     \(T \in  F\left( {\mathcal{X},\mathcal{Y}}\right)\) 当且仅当存在 \({y}_{i} \in  \mathcal{Y}\) 及 \({f}_{i} \in  {\mathcal{X}}^{ * },\left( {i = 1,2,\cdots ,n}\right)\) 使得
\begin{equation*}
T = \mathop{\sum }\limits_{{i = 1}}^{n}{y}_{i} \otimes  {f}_{i}
\end{equation*}
\end{theorem}
\begin{proof}
    充分性是显然的,因为 \(R\left( T\right)  = \operatorname{span}\left\{  {{y}_{1},{y}_{2},\cdots ,{y}_{n}}\right\}\) . 下证必要性. 在 \(R\left( T\right)\) 上取基 \(\left\{  {{y}_{1},{y}_{2},\cdots ,{y}_{n}}\right\}\) ,则 \(\forall x \in  \mathcal{X}\) ,存在唯一 \({\left\{  {l}_{i}\left( x\right) \right\}  }_{i = 1}^{n}\) 使得
\begin{equation*}
{Tx} = \mathop{\sum }\limits_{{i = 1}}^{n}{l}_{i}\left( x\right) {y}_{i}.
\end{equation*}

只需证 \({l}_{i}\left( x\right)\) 为 \(\mathcal{X}\) 上有界线性泛函即可,事实上,

(1) \({l}_{i}\left( {i = 1,2,\cdots ,n}\right)\) 是线性的,因为 \(T\) 是线性的且 \({l}_{i}\) 的表示是唯一的;

(2) \({l}_{i}\) 是有界的.事实上,注意到 \(\parallel {Tx}\parallel\) 及 \(\mathop{\sum }\limits_{{i = 1}}^{n}\left| {{l}_{i}\left( x\right) }\right|\) 都是 \(R\left( T\right)\) 上的模,而且 \(\dim R\left( T\right)  < \infty\) ,所以它们必须是等价范数. 于是存在 \(M > 0\) 使得
\begin{equation*}
\left| {{l}_{i}\left( x\right) }\right|  \leqslant  \mathop{\sum }\limits_{{i = 1}}^{n}\left| {{l}_{i}\left( x\right) }\right|  \leqslant  M\parallel {Tx}\parallel  \leqslant  M\parallel T\parallel \parallel x\parallel ,\forall x \in  \mathcal{X}.
\end{equation*}

因此,存在 \({f}_{i} \in  {X}^{ * }\left( {i = 1,2,\cdots ,n}\right)\) 使得对任意 \(i \in  \{ 1,2,\cdots ,n\}\) 有
\begin{equation*}
\left\langle  {{f}_{i},x}\right\rangle   = {l}_{i}\left( x\right) ,\;\forall x \in  \mathcal{X},i = 1,2,\cdots ,n.
\end{equation*}

于是有
\begin{equation*}
{Tx} = \mathop{\sum }\limits_{{i = 1}}^{n}{y}_{i}\left\langle  {{f}_{i},x}\right\rangle   = \left( {\mathop{\sum }\limits_{{i = 1}}^{n}{y}_{i} \otimes  {f}_{i}}\right) \left( x\right) ,\;\forall x \in  \mathcal{X}.
\end{equation*}


\end{proof}
下面我们来探讨 \(\mathcal{C}\left( {\mathcal{X},\mathcal{Y}}\right)\) 的构造. 由定理4.1知, \(\overline{F\left( {\mathcal{X},\mathcal{Y}}\right) } \subset  \mathcal{C}\left( {\mathcal{X},\mathcal{Y}}\right)\) . 一个自然的问题是, \(\overline{F\left( {\mathcal{X},\mathcal{Y}}\right) } = \mathcal{C}\left( {\mathcal{X},\mathcal{Y}}\right)\) 是否成立? 为简单起见,我们假设 \(\mathcal{X} = \mathcal{Y}\) .
\begin{proposition}
    若 \(\mathcal{X}\) 是一个Hilbert空间,则有 \(\overline{F\left( {\mathcal{X},\mathcal{X}}\right) } = \mathcal{C}\left( \mathcal{X}\right)\) . 
\end{proposition}
\begin{proof}
    对任意 \(T \in\)  \(\mathcal{C}\left( \mathcal{X}\right)\) ,有 \(\overline{T\left( {B}_{1}\right) }\) 紧. 因此,对任意 \(\varepsilon  > 0,\overline{T\left( {B}_{1}\right) }\) 有有限的 \(\varepsilon /2\) 网 \(\left\{  {{y}_{1},{y}_{2},\cdots ,{y}_{n}}\right\}\) ,即
\begin{equation*}
\overline{T\left( {B}_{1}\right) } \subset  \mathop{\bigcup }\limits_{{i = 1}}^{n}B\left( {{y}_{i},\frac{\varepsilon }{2}}\right) .
\end{equation*}

令 \({E}_{n} = \operatorname{span}\left\{  {{y}_{1},{y}_{2},\cdots ,{y}_{n}}\right\}\) ,并令 \({P}_{n}\) 为 \({E}_{n}\) 上的正交投影,那么 \({P}_{n}T \in  F\left( \mathcal{X}\right)\) ,并且对任意 \(x \in  {B}_{1}\) ,存在 \({y}_{i}\left( {1 \leqslant  i \leqslant  n}\right)\) 使得
\begin{equation*}
\begin{Vmatrix}{{Tx} - {y}_{i}}\end{Vmatrix} < \frac{\varepsilon }{2}
\end{equation*}

从而
\begin{equation*}
\begin{Vmatrix}{{P}_{n}{Tx} - {y}_{i}}\end{Vmatrix} = \begin{Vmatrix}{{P}_{n}\left( {{Tx} - {y}_{i}}\right) }\end{Vmatrix} \leqslant \|P_n\|\cdot \|Tx-y_i\|< \frac{\varepsilon }{2}.
\end{equation*}

由此可得
\begin{equation*}
\begin{Vmatrix}{{Tx} - {P}_{n}{Tx}}\end{Vmatrix} < \varepsilon ,\;\forall x \in  {B}_{1},
\end{equation*}

亦即 \(\begin{Vmatrix}{T - {P}_{n}T}\end{Vmatrix} \leqslant  \varepsilon\) . 至此,我们证明了算子 \(T\) 的任意 \(\varepsilon\) 邻域内都有有限秩算子 \({P}_{n}T\) 存在. 从而 \(T \in  \overline{F\left( \mathcal{X}\right) }\) .
\end{proof}
 若 \(\mathcal{X}\) 是Banach 空间,由习题4.1.1可知,我们只需考虑可分Banach 空间即可.

\begin{definition}
     设 \(\mathcal{X}\) 是一个可分的Banach 空间,称 \({\left\{  {e}_{i}\right\}  }_{1}^{\infty }\) 为 \(\mathcal{X}\) 的一组Schauder 基,如果对任意 \(x \in  \mathcal{X}\) ,存在唯一的序列 \(\left\{  {{C}_{n}\left( x\right) }\right\}\) ,使得
\begin{equation*}
x = \mathop{\lim }\limits_{{N \rightarrow  \infty }}\mathop{\sum }\limits_{{n = 1}}^{N}{C}_{n}\left( x\right) {e}_{n}.
\end{equation*}
\end{definition}


从定义中可以看出 \({C}_{n}\left( x\right)\) 是 \(\mathcal{X}\) 上的线性函数,实际上,我们有
\begin{lemma}
     \({C}_{n}\left( x\right)\) 是 \(\mathcal{X}\) 上的连续泛函.
\end{lemma}
\begin{proof}
    在 \(\mathcal{X}\) 上引入另一范数 \(\parallel  \cdot  {\parallel }_{1}\) 如下:
\begin{equation*}
\parallel x{\parallel }_{1} = \mathop{\sup }\limits_{{N \in  \mathbb{N}}}\begin{Vmatrix}{{S}_{N}x}\end{Vmatrix},
\end{equation*}

其中 \({S}_{N}x = \mathop{\sum }\limits_{{n = 1}}^{N}{C}_{n}\left( x\right) {e}_{n}\left( {\forall x \in  \mathcal{X}}\right)\) . 容易验证 \(X\) 按范数 \(\parallel  \cdot  {\parallel }_{1}\) 形成一个完备Banach空间, 并且
\begin{equation*}
\parallel x\parallel  = \mathop{\lim }\limits_{{N \rightarrow  \infty }}\begin{Vmatrix}{{S}_{N}x}\end{Vmatrix} \leqslant  \parallel x{\parallel }_{1},\;\forall x \in  \mathcal{X}.
\end{equation*}

由Banach空间等价范数定理知,存在 \({M}_{1} > 0\) ,使得
\begin{equation*}
\parallel x{\parallel }_{1} \leqslant  {M}_{1}\parallel x\parallel ,\;\forall x \in  \mathcal{X}.
\end{equation*}

于是对任意 \(n \in  \mathbb{N}\) ,有
\begin{equation*}
\begin{Vmatrix}{{C}_{n}\left( x\right) {e}_{n}}\end{Vmatrix} = \begin{Vmatrix}{{S}_{n}x - {S}_{n - 1}x}\end{Vmatrix} \leqslant  2\parallel x{\parallel }_{1} \leqslant  2{M}_{1}\parallel x\parallel ,\;\forall x \in  \mathcal{X}.
\end{equation*}

由此可得,对任意 \(n \in  \mathbb{N}\) ,
\begin{equation*}
\left| {{C}_{n}\left( x\right) }\right|  \leqslant  2{M}_{1}{\begin{Vmatrix}{e}_{n}\end{Vmatrix}}^{-1}\parallel x\parallel ,\;\forall x \in  \mathcal{X}.
\end{equation*}

故 \({C}_{n}\left( x\right) ,\left( {n = 1,2,\cdots }\right)\) 是连续的.

\end{proof}

现在, 我们可以证明
\begin{theorem}
     若可分Banach空间 \(\mathcal{X}\) 有一组Schauder基 \({\left\{  {e}_{n}\right\}  }_{1}^{\infty }\) ,则
\begin{equation*}
\overline{F\left( \mathcal{X}\right) } = \mathcal{C}\left( \mathcal{X}\right) .
\end{equation*}
\end{theorem}
\begin{proof}
    由于 \({\left\{  {e}_{n}\right\}  }_{1}^{\infty }\) 为Schauder 基,故对任意 \(x \in  \mathcal{X}\) ,存在唯一序列 \({\left\{  {C}_{n}\left( x\right) \right\}  }_{1}^{\infty }\) 使得
\begin{equation*}
x = \mathop{\lim }\limits_{{N \rightarrow  \infty }}\mathop{\sum }\limits_{{n = 1}}^{N}{C}_{n}\left( x\right) {e}_{n}.
\end{equation*}

(1)对任意 \(N \in  \mathbb{N}\) ,令
\begin{equation*}
{S}_{N}x = \mathop{\sum }\limits_{{n = 1}}^{N}{C}_{n}\left( x\right) {e}_{n},\;\forall x \in  \mathcal{X},
\end{equation*}

并记 \({R}_{N} = I - {S}_{N}\) ; 则由共鸣定理知,存在 \(M > 0\) 使得
\begin{equation*}
\begin{Vmatrix}{S}_{N}\end{Vmatrix} \leqslant  M,\;\begin{Vmatrix}{R}_{N}\end{Vmatrix} \leqslant  1 + M.
\end{equation*}

(2)若 \(T \in  \mathcal{C}\left( \mathcal{X}\right)\) ,则需证对任意 \(\varepsilon  > 0\) ,都存在 \({T}_{\varepsilon } \in  F\left( \mathcal{X}\right)\) 使得 \(\begin{Vmatrix}{T - {T}_{\varepsilon }}\end{Vmatrix} <\)  \(\varepsilon\) . 由 \(T\) 为紧算子知, \(\overline{T\left( {B}_{1}\right) }\) 紧. 从而存在有穷的 \(\varepsilon /\left\lbrack  {3\left( {M + 1}\right) }\right\rbrack\) -网 \(\left\{  {{y}_{1},{y}_{2},\cdots ,{y}_{m}}\right\}\) . 即对任意 \(x \in  {B}_{1}\) ,存在 \({y}_{i}\left( {1 \leqslant  i \leqslant  m}\right)\) 使得
\begin{equation*}
\begin{Vmatrix}{{Tx} - {y}_{i}}\end{Vmatrix} < \frac{\varepsilon }{3\left( {M + 1}\right) }.
\end{equation*}

由Schauder基的定义知,存在 \(N \in  \mathbb{N}\) 使得
\begin{equation*}
\begin{Vmatrix}{{y}_{j} - {S}_{N}{y}_{j}}\end{Vmatrix} < \frac{\varepsilon }{3},\;j = 1,2,\cdots ,m.
\end{equation*}

但由于 \(\begin{Vmatrix}{S}_{N}\end{Vmatrix} \leqslant  M\) ,故有
\begin{equation*}
\begin{Vmatrix}{{S}_{N}\left( {Tx}\right)  - {S}_{N}{y}_{i}}\end{Vmatrix} < \frac{M}{3\left( {M + 1}\right) }\varepsilon .
\end{equation*}

进而可得
\begin{equation*}
\begin{Vmatrix}{{Tx} - \left( {{S}_{N}T}\right) x}\end{Vmatrix} < \varepsilon ,\;\forall x \in  {B}_{1}.
\end{equation*}

取 \({T}_{\varepsilon } = {S}_{N}T\) 即得所求.

\end{proof}

\subsection{习 题}

\begin{problemset}
  \item 若 \(T \in  \mathcal{C}\left( {\mathcal{X},\mathcal{Y}}\right)\)，则 \(R\left( T\right)\) 可分.
\begin{proof}
    考虑$B_n,T(B_n)$是相对紧集,那么完全有界,进而$T(B_n)$可分,而$\mathcal{R}(T)=\bigcup\limits_{n=1}^\infty T(B_n)$也可分.
\end{proof}
  \item 设 \(\mathcal{X}\) 为无穷维 Banach 空间，证明：若 \(T \in  \mathcal{C}\left( \mathcal{X}\right)\)，则 \(T\) 没有有界逆.
\begin{proof}
    若 \(T\) 有有界逆$T^{-1}$,那么$I=TT^{-1}$是紧算子,那么$B_1$是相对紧集,但是有Riesz引理知无穷维Banach空间中$B_1$不是列紧的,进而与$B_1$是相对紧集矛盾.
\end{proof}
  \item 设 \(\mathcal{X}\) 为 Banach 空间，\(T \in  \mathcal{C}\left( \mathcal{X}\right)\) 满足 \(\|Tx\| \geq \alpha \|x\|\)（\(\forall x \in  \mathcal{X}\)），其中 \(\alpha\) 是正常数.  
  证明：\(T \in  \mathcal{C}\left( \mathcal{X}\right)\) 当且仅当 \(\mathcal{X}\) 是有限维的.
\begin{proof}
    必要性. 如果$\mathcal{X}$是无限维的,那么$B_1$不是列紧的,存在$x_n\in B_1,n\in \mathbb{N}^+,\|x_n-x_m\|>\frac{1}{2}$,进而$\|Tx_n-Tx_m\|\geq \frac{1}{2}\alpha$,那么$T(B_1)$也不是列紧的,进而$T(B_1)$不是相对紧集,与\(T \in  \mathcal{C}\left( \mathcal{X}\right)\)矛盾.

    充分性. 由于$\mathcal{X}$是有限维的,那么$T(B_1)$列紧,进而\(T \in  \mathcal{C}\left( \mathcal{X}\right)\).
\end{proof}
  \item 设 \(\mathcal{X}\) 为 Banach 空间.证明：若恒等映射 \(I : \mathcal{X} \rightarrow  \mathcal{X}\) 为紧映射，则 \(\mathcal{X}\) 为有限维.
\begin{proof}
    $I$是紧算子,那么$B_1$是相对紧集,但是有Riesz引理知无穷维Banach空间中$B_1$不是列紧的,进而与$B_1$是相对紧集矛盾.
\end{proof}
  \item 设 \(\mathcal{X} = C^{k}\left[ a,b\right], \mathcal{Y} = C^{j}\left[ a,b\right]\).则  
  (1) 若 \(j \leqslant k\)，则嵌入映射 \(\iota : \mathcal{X} \hookrightarrow  \mathcal{Y}\) 连续;  
  (2) 若 \(j < k\)，则嵌入映射 \(\iota : \mathcal{X} \hookrightarrow  \mathcal{Y}\) 为紧映射.
\begin{proof}
    (1) 显然成立.

    (2) 任取序列$\{f_n\}\subset B_1$,$\{f_n\}$显然一致有界,下面证明等度连续.

    由于$j<k$,所以
    \[|f^{(i)}_n(x)-f^{(i)}_n(y)|=|f^{(i+1)}(\xi)||x-y|\leqslant \|f\|_k\cdot|x-y|=|x-y|,\forall i\leqslant j<k\]

    进而$\{f_n\}$等度连续,那么$\{f_n\}$列紧.

    
\end{proof}
  \item 设 \(a,b \in  \mathbb{R}, a < b\).定义算子
  \[
  (Tx)(s) := \int_{a}^{s} x(t)\,\dif t.
  \]
  证明：\(T : C[a,b] \to C[a,b]\) 为紧算子.
\begin{proof}
    同样易证一致有界和等度连续.
\end{proof}
  \item 设 \(\Omega  \subset  \mathbb{R}^{n}\) 为有界闭集.证明：  
  
  (1) 若 \(k(x,y) \in  C(\Omega  \times  \Omega)\)，则式\ref{4.1}中定义的算子 \(K : L(\Omega) \to C(\Omega)\) 为紧算子;  
  
  (2) 若 \(k(x,y) \in  L^{2}(\Omega  \times  \Omega)\)，则式\ref{4.1}中定义的算子 \(K : L^{2}(\Omega) \to L^{2}(\Omega)\) 为紧算子.
\begin{proof}
    (1) 先证$T(B_1)$一致有界.
    \[\left\|K f\right\|=\sup\limits_{x\in \Omega}\left|\int_\Omega k(x,y)f(y)\dif y\right|\leqslant \|k(x,y)\|\cdot \|f\|_1\leqslant\|k(x,y)\| \]

    下面证明等度连续,由于$k(x,y)$在$\Omega\times \Omega$上一致连续,那么$\forall \varepsilon>0,\exists \delta >0,|x_1-x_2|<\delta,|k(x_1,y)-k(x_2,y)|<\varepsilon $.
    \begin{equation*}
            \begin{aligned}
            \left|Kf(x_1)-Kf(x_2)\right|&=\left|\int_\Omega k(x_1,y)f(y)\dif y -\int_\Omega k(x_2,y)f(y)\dif  y\right|\\
            &\leqslant \left|\int_\Omega \left(k(x_1,y)-k(x_2,y)\right)f(y)\dif  y\right|\\
            &\leqslant \left| k(x_1,y)-k(x_2,y)\right|\cdot \|f\|_1<\varepsilon
    \end{aligned}
\end{equation*}
    (2) 由于$L^2(\Omega\times \Omega)$是自反的,那么只需要证明$K$是全连续的,设$f_n\in L^2(\Omega),f_n \rightharpoonup f_0$.

    由Fubini定理知,对于几乎所有的$x,K(x,\cdot)\in L^2(\Omega)$.固定$x$,由Holder不等式知$K$是连续线性泛函.进而由$f_n$弱收敛知
\[\left|Kf_n -Kf_0\right|=\left|\int_\Omega k(x,y)(f_n(y)-f_0(y))\dif y\right|\to 0,\text{a.s.}x\in \Omega\]

进而
\begin{equation*}
        \begin{aligned}
\|Kf_n -Kf_0\|_2^2=\int_\Omega \left|\int_\Omega k(x,y)(f_n(y)-f_0(y))\dif y\right|^2\dif x\to 0
        \end{aligned}
    \end{equation*}



\end{proof}
\item 设 \( \mathcal{X} \) 是 \( B \) 空间, \( A \in  \mathcal{C}\left( \mathcal{X}\right) ,{\mathcal{X}}_{0} \) 是 \( \mathcal{X} \) 的闭子空间并使得
\( A{\mathcal{X}}_{0} \subset  {\mathcal{X}}_{0} \) . 求证: 映射 \( T : \left\lbrack  x\right\rbrack   \mapsto  \left\lbrack  {Ax}\right\rbrack \) 是商空间 \( \mathcal{X}/{\mathcal{X}}_{0} \) 上的紧算
子.
\begin{proof}
    只要证映射 \( T \) 映 \( \mathcal{X}/{\mathcal{X}}_{0} \) 上的有界集为列紧集.
事实上,设 \( \begin{Vmatrix}\left\lbrack  {x}_{n}\right\rbrack  \end{Vmatrix} \leqslant  C \) ,易证\(\exists {\widetilde{x}}_{n} \in  \left\lbrack  {x}_{n}\right\rbrack \) ,使
得 \( \begin{Vmatrix}{\widetilde{x}}_{n}\end{Vmatrix} \leqslant  {2C} \) . 又因为 \( A{\mathcal{X}}_{0} \subset  {\mathcal{X}}_{0} \) ,所以
\[ {\widetilde{x}}_{n} \in  \left\lbrack  {x}_{n}\right\rbrack \Rightarrow {\widetilde{x}}_{n} - {x}_{n} \in  {\mathcal{X}}_{0} \Rightarrow  A{\widetilde{x}}_{n} - A{x}_{n} \in  A{\mathcal{X}}_{0} \subset  {\mathcal{X}}_{0} 
\Rightarrow  \left\lbrack  {A{x}_{n}}\right\rbrack   = \left\lbrack  {A{\widetilde{x}}_{n}}\right\rbrack \]

进而
\[ {\widetilde{x}}_{n} \in  \left\lbrack  {x}_{n}\right\rbrack   \Rightarrow  \left\lbrack  {A{x}_{n}}\right\rbrack   = \left\lbrack  {A{\widetilde{x}}_{n}}\right\rbrack  . \]

因此
\( T\left\lbrack  {x}_{n}\right\rbrack   = \left\lbrack  {A{x}_{n}}\right\rbrack   = \left\lbrack  {A{\widetilde{x}}_{n}}\right\rbrack  . \)
进一步,因为 \( \begin{Vmatrix}{\widetilde{x}}_{n}\end{Vmatrix} \leqslant  {2C},A \in  \mathcal{C}\left( \mathcal{X}\right) \) ,所以 \( \exists A{\widetilde{x}}_{{n}_{k}} \rightarrow  y \) . 又 \( \left\lbrack  \cdot \right\rbrack \) 是线性连
续映射, 故有
\[ T\left\lbrack  {x}_{{n}_{k}}\right\rbrack   = \left\lbrack  {A{\widetilde{x}}_{{n}_{k}}}\right\rbrack   \rightarrow  \left\lbrack  y\right\rbrack  . \]

由此可见,映射 \( T \) 映 \( \mathcal{X}/{\mathcal{X}}_{0} \) 上的有界集为列紧集,故
\( T \in  \mathcal{C}\left( {\mathcal{X}/{\mathcal{X}}_{0}}\right) \) .
\end{proof}


  \item 设 \(\mathcal{X}, \mathcal{Y}\) 为 Banach 空间，\(A \in  \mathcal{L}(\mathcal{X},\mathcal{Y})\)，\(K \in  \mathcal{C}(\mathcal{X},\mathcal{Y})\).  
  证明：如果 \(R(A) \subset R(K)\)，则 \(A \in  \mathcal{C}(\mathcal{X},\mathcal{Y})\).
\begin{proof}
 对 \( \forall \left\lbrack  x\right\rbrack   \in  \mathcal{X}/N\left( K\right) ,\widetilde{K}\left\lbrack  x\right\rbrack  \overset{\text{ def }}{ = }{Kx} \) , 令 \( S \) 为 \( \mathcal{X} \) 空间的闭单位球,则 \( (S + N\left( K\right) )/N(K)\) 是
\( \mathcal{X}/N\left( K\right) \) 中的单位球,且 \( \widetilde{K}\left( {S + N\left( K\right) }\right)  = K\left( S\right) \) 是列紧的,从而 \( \widetilde{K} \) 是
紧算子. 又 \( \widetilde{K} \) 是单射连续算子,故 \( {\widetilde{K}}^{-1} \) 闭,且
\( D\left( {\widetilde{K}}^{-1}\right)  = R\left( K\right)  \supset  R\left( A\right) , \)

从而 \( {\widetilde{K}}^{-1}A : \mathcal{X} \rightarrow  \mathcal{X}/N\left( K\right) \) 是闭算子,且定义域是全空间 (当然闭). 根
据闭图像定理, \( {\widetilde{K}}^{-1}A \) 有界,因此,根据命题 2(6),有
\( A = {\widetilde{K}}{\left( {\widetilde{K}}^{-1}A\right) } \in  \mathcal{C}\left( {\mathcal{X},\mathcal{Y}}\right) . \)
\end{proof}
  \item 设 \(H\) 是 Hilbert 空间，\(T : H \to H\) 是紧算子.证明：若 \({x}_{n} \rightharpoonup  {x}_{0}\)，\({y}_{n} \rightharpoonup  {y}_{0}\)，则\[
  (x_{n}, T y_{n}) \rightarrow (x_{0}, T y_{0}), \quad n \rightarrow \infty.
  \]
\begin{proof}
    由于Hilbert空间中紧算子是全连续的,\begin{equation*}
        \begin{aligned}
            \left| (x_{n}, T y_{n}) - (x_{0}, T y_{0})\right|&\leqslant   \left| (x_{n}, T y_{n}) - (x_{n}, T y_{0})\right|+ \left| (x_{n}, T y_{0}) - (x_{0}, T y_{0})\right|\\&\leqslant   \|x_n\|\|Ty_n-Ty_0\|+ \left| (x_{n}, T y_{0}) - (x_{0}, T y_{0})\right|\\
       &\leqslant   M\|Ty_n-Ty_0\|+ \left| (x_{n}, T y_{0}) - (x_{0}, T y_{0})\right|\\&\to 0
        \end{aligned}
    \end{equation*}
\end{proof}
  \item 证明例 4.5 中的结论对边界不光滑的有界区域仍成立.
  \item  设 $M$ 是 Banach 空间 $\mathcal{X}$ 的闭线性子空间，称满足 $P^2=P$（幂等性）的由 $\mathcal{X}$ 到 $M$ 上的一个有界线性算子 $P$ 为由 $\mathcal{X}$ 到 $M$ 上的投影算子．求证：
  
（1）若 $M$ 是 $\mathcal{X}$ 的有穷维线性子空间，则必存在由 $\mathcal{X}$ 到 $M$ 上的投影算子;

（2）若 $P$ 是由 $\mathcal{X}$ 到闭线性子空间 $M$ 上的投影算子，则 $I-P$ 是由 $\mathcal{X}$ 到 $R(I-P)$ 上的投影算子;

（3）若 $P$ 是由 $\mathcal{X}$ 到 $M$ 上的投影算子，则 $\mathcal{X}=M \oplus N$ ，其中 $N= R(I-P)$ ．\begin{proof}
    设 $\left\{x_1, x_2, \cdots, x_n\right\}$ 是 $M$ 的一个基，根据Hanh-Banach定理， $\exists f_1, \cdots, f_n \in \mathcal{X}^*$ ，使得
$$
\left\langle f_i, x_j\right\rangle=\delta_{i j} \quad(i, j=1,2, \cdots, n) .
$$
令 $P: x \mapsto P x=\sum_{k=1}^n\left\langle f_k, x\right\rangle x_k$ ，则有
$$
\begin{aligned}
& P x_j=x_j \quad(j=1,2, \cdots, n), \\
& \|P x\| \leqslant\left(\sum_{k=1}^n\left\|f_k\right\|\left\|x_k\right\|\right)\|x\| .
\end{aligned}
$$
由此可见，$P \in \mathcal{L}(\mathcal{X})$ .

进一步要证明 $P \mathcal{X}=M$.一方面，$P \mathcal{X} \subset M$ 是显然的．另一方面，还需证明 $M \subset P \mathcal{X}$ ．事实上，
$$
\begin{aligned}
 \forall x \in M  \Rightarrow x=\sum_{k=1}^n \alpha_k x_k \Rightarrow P x=\sum_{k=1}^n \alpha_k P\left(x_k\right)=\sum_{k=1}^n \alpha_k x_k=x\Rightarrow x=P x \in P \mathcal{X} \\
\end{aligned}
$$

现在，有了 $P \mathcal{X}=M, P$ 的幂等性就容易证明了.事实上，$\forall x \in \mathcal{X}$ ，因为 $P x \in M$ ，所以 $P x=\sum_{k=1}^n \beta_k x_k$ ，于是
$$
\begin{gathered}
P^2 x =
P P x=\sum_{k=1}^n \beta_k P\left(x_k\right)=\sum_{k=1}^n \beta_k x_k
\end{gathered}
$$
综上所述即知 $P$ 是由 $\mathcal{X}$ 到 $M$ 上的投影算子．

(2)已知 $P^2=P, P \mathcal{X}=M$ ，要证 $I-P$ 是由 $\mathcal{X}$ 到 $R(I-P)$ 上的投影算子，就是要证
$$
(I-P)^2=I-P, \quad(I-P) \mathcal{X}=R(I-P) .
$$
  首先 $(I-P)^2=I-2 P+P^2 = I-P$ ;
  
其次证 $(I-P) \mathcal{X}=R(I-P)$.一方面，
$$
(I-P) \mathcal{X} \subset R(I-P)
$$
是显然的．另一方面，
$\forall x \in R(I-P),\exists x_0$ ，使得 $x=(I-P) x_0 \Longrightarrow(I-P) x=(I-P)^2 x_0=(I-P) x_0=x\Rightarrow x=(I-P) x$从此即知 $R(I-P) \subset(I-P) \mathcal{X}$.

综合以上两方面即知
$$
R(I-P)=(I-P) \mathcal{X} .
$$

(3)分解 $x=P x+(I-P) x$ ，其中 $P x \in M$ ，且
$$
(I-P) x \in R(I-P)=N .
$$

又如果 $x \in M \cap N$ ，则
$$
\left.\begin{array}{r}
x \in M \Rightarrow P x=x \\
x \in N \Longrightarrow(I-P) x=x
\end{array}\right\} \Rightarrow x=\theta .
$$

故有 $\mathcal{X}=M \oplus N$ ，其中 $N=R(I-P)$ ．
\end{proof}
\item  若 $\mathcal{X}$ 是 $B$ 空间，$M \subset \mathcal{X}$ 是闭线性子空间，求证：

（1）若 $\operatorname{dim}(\mathcal{X} / M)<\infty$ ，则 $\exists N \subset \mathcal{X}, \operatorname{dim} N<\infty$ ，使得 $N \cong \mathcal{X} / M$ ，且 $\mathcal{X}=M \oplus N$ ;

（2）若 $\operatorname{dim} M<\infty$ ，则
$$
\mathcal{X} / M \cong R\left(I-P_M\right) \quad \text { 且 } \quad \mathcal{X}=M \oplus R\left(I-P_M\right) \text {. }
$$

\begin{proof}
(1)若 $\operatorname{dim}(\mathcal{X} / M)=n$ ，设 $\left[x_1\right], \cdots,\left[x_n\right]$ 是 $\mathcal{X} / M$ 上的一组基，任取
$$
x_k \in\left[x_k\right] \quad(k=1,2, \cdots, n),
$$
则 $\left\{x_k\right\}$ 线性无关．事实上，若
$$
c_1 x_1+c_2 x_2+\cdots+c_n x_n=\theta,
$$

则有
$$
\left[c_1 x_1+c_2 x_2+\cdots+c_n x_n\right]=[\theta],
$$
即
$$
c_1\left[x_1\right]+c_2\left[x_2\right]+\cdots+c_n\left[x_n\right]=[\theta] .
$$
因为 $\left[x_1\right], \cdots,\left[x_n\right]$ 是 $\mathcal{X} / M$ 上的一组基，所以$$
c_1=c_2=\cdots=c_n=0
$$
又因为 $\left[x_k\right] \neq[\theta]$ ，所以
$$
x_k \notin M \quad(k=1,2, \cdots, n) .
$$
于是，若令 $N=\operatorname{span}\left\{x_1, x_2, \cdots, x_n\right\}$ ，则有 $N \cap M=\{\theta\}$ ．

下证 $N \cong \mathcal{X} / M$ 。令
$$
\tau: x=\sum_{k=1}^n \alpha_k x_k \mapsto \sum_{k=1}^n \alpha_k\left[x_k\right],
$$
则有

$\tau$ 的连续性：由 $\|\tau x\|=\|[x]\| \leqslant\|x\|$ 得出 $\tau$ 是连续算子．

$\tau$ 是单射：
$$
\tau x=0 \Longrightarrow[x]=[\theta] \Longrightarrow x \in N \cap M \Longrightarrow x=\theta .
$$

$\tau$ 是满射：$\forall[x]=\sum_{k=1}^n \alpha_k\left[x_k\right] \in \mathcal{X} / M$ ，
$$
\tau(\underbrace{\sum_{k=1}^n \alpha_k x_k}_{\in N})=[x] .
$$
根据逆算子定理，$\tau^{-1}$ 存在连续，故 $N \cong \mathcal{X} / M$ .

进一步，令 $P x=\tau^{-1}[x](\forall x \in \mathcal{X})$ ，则有
$$
P x_k=\tau^{-1}\left[x_k\right]=x_k .
$$
对 $\forall x \in \mathcal{X},[x]=\sum_{k=1}^n \alpha_k\left[x_k\right]$ ，
$$
P x=\tau^{-1}[x]=\sum_{k=1}^n \alpha_k x_k \in N,
$$
由上式，有
\[
P^2x = \sum_{k=1}^n \alpha_kP( x_k)= \sum_{k=1}^n \alpha_k x_k=Px  
\]
所以$P$ 是幂等算子．又显然 $R(P)=N$ ，故 $P$ 
为由 $\mathcal{X}$ 到 $N$ 上的投影算子．下面改写 $P=P_{N^*}$ 于是 $N=R\left(P_N\right)$ ．又
$$x \in M\Leftrightarrow 
[x]=[\theta] \Longleftrightarrow P_N x=\tau^{-1}[x]=\theta \Longleftrightarrow\left(I-P_N\right) x=x\Leftrightarrow
x \in R\left(I-P_N\right)$$
从此即知 $M=R\left(I-P_N\right)$ 。于是，由上题知
$$
\mathcal{X}=R\left(P_N\right) \oplus R\left(I-P_N\right)=N \oplus M .
$$
(2)由上题知$\mathcal{X}=R(P_M)\oplus R(I-P_M)$,记$N=R(I-P_M)$.定义映射\[
\tau: \mathcal{X}/M\to N,[x]\mapsto (I-P_M)x
\]
则有

$\tau$ 的连续性：$\forall [x]\in \mathcal{X}/M,\exists \bar{x}\in [x],\|\bar{x}\|\leqslant 2\|[x]\|$由 $\|\tau[ x]\|=\|(I-P_M)\bar{x}\| \leqslant2(1+\|P_M\|)\|[x]\|$ 得出 $\tau$ 是连续算子．

$\tau$ 是单射：
$$
\tau[x]=0 \Longrightarrow(I-P_M)x=\theta \Longrightarrow x \in  M \Longrightarrow [x]=[\theta] .
$$

$\tau$ 是满射: 显然成立.

而$N=R(I-P_M)=\Ker P_M$是闭空间,根据Banach逆算子定理，$\tau^{-1}$ 存在连续，故 $N \cong \mathcal{X} / M$ .

\end{proof}

\end{problemset}

\section{Riesz-Schauder理论}

在线性代数中已经对矩阵的特征值、不变子空间及Jordan标准型理论有系统的研究并有非常明确的结论. 从大学的泛函分析课程中,我们知道,有界线性算子是矩阵概念到无穷维的推广. 自然会问:对无限维的线性算子,相关结论又如何呢？ 实际上, 这些问题是算子谱论的中心问题. 从定理1.16知, 一般有界线性算子的谱集非常复杂. 但对于紧算子而言, 结果则理想的多. 这将是本节的中心问题.

本节假设 \(\mathcal{X}\) 为无穷维复Banach空间.

\subsection{几个引理}

在叙述和证明本节的主要定理之前, 我们先证明一些引理.
\begin{definition}
    对任意 \(T \in  \mathcal{L}\left( \mathcal{X}\right)\) ,分别定义其像集和核空间为
\begin{equation*}
R\left( T\right)  \mathrel{\text{ := }} T\left( \mathcal{X}\right) ,\;\text{ 及 }\;N\left( T\right)  \mathrel{\text{ := }} \{ x \in  \mathcal{X} \mid  {Tx} = 0\} .
\end{equation*}

对任意的 \(M \subset  \mathcal{X},N \subset  {\mathcal{X}}^{ * }\) ,定义
\begin{equation*}
{}^{ \bot  }M \mathrel{\text{ := }} \left\{  {f \in  {\mathcal{X}}^{ * }\mid \langle f,x\rangle  = 0,\forall x \in  M}\right\}  ,
\end{equation*}
\begin{equation*}
{N}^{ \bot  } \mathrel{\text{ := }} \{ x \in  \mathcal{X} \mid  \langle f,x\rangle  = 0,\forall f \in  N\} .
\end{equation*}

若 \(f \in  {\mathcal{X}}^{ * },x \in  \mathcal{X}\) 满足 \(\langle f,x\rangle  = 0\) ,有时简记为 \(f \bot  x\) .
\end{definition}
\begin{lemma}
    设 \(X\) 为Banach空间, \(T \in  \mathcal{C}\left( \mathcal{X}\right)\) .  若 \(\lambda  \neq  0\) ,则 \(N\left( {{\lambda I} - T}\right)\) 和 \(N\left( {{\lambda I} - {T}^{ * }}\right)\) 都是有限维的.

\end{lemma}
\begin{proof}
设 \({x}_{n}\) 为子空间 \(N\left( {{\lambda I} - T}\right)\) 中有界列,从 \(\lambda {x}_{n} = T{x}_{n}\) 和算子 \(T\) 紧知, \(\left\{  {x}_{n}\right\}\) 有收敛子列. 即有 \(N\left( {{\lambda I} - T}\right)\) 中任意有界集都是相对紧集. 又由于Banach空间的任意有界集为相对紧集当且仅当其为有限维,故有 \(\mathrm{{dim}}N({\lambda I} -\)  \(T) < \infty\) . 同理可得 \(\dim N\left( {{\lambda I} - {T}^{ * }}\right)  < \infty\) .

\end{proof}



\begin{lemma}
     设 \(\mathcal{X}\) 为Banach空间, \(T \in  \mathcal{C}\left( \mathcal{X}\right) ,\lambda  \neq  0\) . 则存在整数 \(p \geq  1\) 使得
\begin{equation*}
N\left( {\left( \lambda I - T\right) }^{p}\right)  = N\left( {\left( \lambda I - T\right) }^{p + 1}\right) .
\end{equation*}
\end{lemma}
\begin{proof}
    对 \(k \geq  1\) ,令 \({N}_{k} = N\left( {\left( \lambda I - T\right) }^{k}\right)\) . 易知 \({\left( \lambda I - T\right) }^{k} = {\lambda }^{k}I - {T}_{1}\) ,其中 \({T}_{1}\) 是 \(T\) 的多项式,因而仍是紧算子,于是 \(\dim {N}_{k} < \infty\) . 下面用反证法来证明引理结论. 若不然,即对任意 \(k \geq  1\) 都有 \({N}_{k} \subsetneqq  {N}_{k + 1}\) ,由Riesz引理知,存在序列 \({x}_{k} \in  {N}_{k + 1}\) ,使得 \(\begin{Vmatrix}{x}_{k}\end{Vmatrix} = 1\) 且
\begin{equation*}
\begin{Vmatrix}{{x}_{k} - x}\end{Vmatrix} > \frac{1}{2},\;\forall x \in  {N}_{k}.
\end{equation*}

当 \(k > m\) 时有
\begin{equation*}
{\left( \lambda I - T\right) }^{k}\left\lbrack  {\left( {{\lambda I} - T}\right) {x}_{k} + T{x}_{m}}\right\rbrack   = {\left( \lambda I - T\right) }^{k + 1}{x}_{k} + T{\left( \lambda I - T\right) }^{k}{x}_{m} = 0.
\end{equation*}

从而有 \(\left( {{\lambda I} - T}\right) {x}_{k} + T{x}_{m} \in  {N}_{k}\) . 因此
\begin{equation*}
\frac{1}{\left| \lambda \right| }\begin{Vmatrix}{T{x}_{k} - T{x}_{m}}\end{Vmatrix} = \begin{Vmatrix}{{x}_{k} - \frac{1}{\lambda }\left\lbrack  {\left( {{\lambda I} - T}\right) {x}_{k} + T{x}_{m}}\right\rbrack  }\end{Vmatrix} > \frac{1}{2}.
\end{equation*}

这与 \(T\) 的紧性矛盾,因此存在有限的 \(p\) 使得 \({N}_{p} = {N}_{p + 1}\) .

\end{proof}


\begin{lemma}
    设 \(\mathcal{X}\) 为Banach空间, \(T \in  \mathcal{C}\left( \mathcal{X}\right) ,\lambda  \neq  0\) . 则 \(R\left( {{\lambda I} - T}\right)  = \mathcal{X}\) 当且仅当 \(N\left( {{\lambda I} - T}\right)  = 0\) .
\end{lemma}
\begin{proof}
     设 \(R\left( {{\lambda I} - T}\right)  = \mathcal{X}\) . 若 \(N\left( {{\lambda I} - T}\right)  \neq  0\) ,存在非零元 \({x}_{1} \in  N\left( {{\lambda I} - T}\right)\) . 从而由 \(R\left( {{\lambda I} - T}\right)  = \mathcal{X}\) 知,存在序列 \(\left\{  {x}_{n}\right\}\) 使得 \(\left( {{\lambda I} - T}\right) {x}_{n} = {x}_{n - 1}\left( {n \geq  2}\right)\) . 易知 \({\left( \lambda I - T\right) }^{n}{x}_{n + 1} = {x}_{1},{\left( \lambda I - T\right) }^{n + 1}{x}_{n + 1} = 0\) ,于是对任意 \(n \geq  1\) 有, \(N(({\lambda I} -\)  \(\left. {T{)}^{n}}\right)  \subsetneqq  N\left( {\left( \lambda I - T\right) }^{n + 1}\right)\) . 这与引理4.2.2矛盾.

反之,如果 \(N\left( {{\lambda I} - T}\right)  = 0\) ,由引理 4.2.1 知, \(R\left( {\bar{\lambda }I - {T}^{ * }}\right)  = {\mathcal{X}}^{ * }\) . 由于 \({T}^{ * }\) 也是紧算子,故 \(N\left( {\bar{\lambda }I - {T}^{ * }}\right)  = 0\) . 同理由引理4.2.1知, \(R\left( {{\lambda I} - T}\right)  = N{\left( \bar{\lambda }I - {T}^{ * }\right) }^{ \bot  } =\)  \(\mathcal{X}\) .

\end{proof}

\begin{lemma}
 若 \(T \in  \mathcal{C}\left( X\right) ,\lambda  \neq  0,{T}_{\lambda } = {\lambda I} - T\) . 设 \( {x}_{1},{x}_{2},\cdots ,{x}_{n} \in  N\left(T_\lambda\right) \) 为 \( N\left(T_\lambda\right) \) 的一组基,存在闭线性 \( {\mathcal{X}}_{1} \subset  \mathcal{X} \) ,使得
\( \mathcal{X} = \operatorname{span}\left\{  {{x}_{1},{x}_{2},\cdots ,{x}_{n}}\right\}   \oplus  {\mathcal{X}}_{1}. \)
\end{lemma}

\begin{proof}
    由 Hahn-Banach 定理 (定理 2.4.4),存在 \( {g}_{1},{g}_{2},\cdots ,{g}_{n} \in \mathcal{X}^*\),满足
\( {g}_{i}\left( {x}_{j}\right)  = {\delta }_{ij},\;1 \leq  i,j \leq  n. \)

令 \( {\mathcal{X}}_{1} = \mathop{\bigcap }\limits^{n}N\left( {g}_{i}\right) \) ,其中 \( N\left( g\right)  = \{ x \in  \mathcal{X} \mid  g\left( x\right)  = 0\} \) ,则 \( {\mathcal{X}}_{1} \) 是闭线
性的,满足

(1) \( \operatorname{span}\left\{  {{x}_{1},{x}_{2},\cdots ,{x}_{n}}\right\}   \cap  {\mathcal{X}}_{1} = \{ \theta \} \) ;

(2) \( \forall x \in  \mathcal{X} \) ,取 \( {c}_{i} = {g}_{i}\left( x\right) \) ,有 \( x - \mathop{\sum }\limits_{{i = 1}}^{n}{c}_{i}{x}_{i} \in  {\mathcal{X}}_{1} \) .

从而有 \( \mathcal{X} = \operatorname{span}\left\{  {{x}_{1},{x}_{2},\cdots ,{x}_{n}}\right\}   \oplus  {\mathcal{X}}_{1} \) .
\end{proof}
\begin{lemma}
   若 \(T \in  \mathcal{C}\left( X\right) ,\lambda  \neq  0,{T}_{\lambda } = {\lambda I} - T\) . 设 \( {f}_{1},{f}_{2},\cdots ,{f}_{m} \in \)\( N\left( {T_\lambda}^{ * }\right) \) 为 \( N\left( {T_\lambda}^{ * }\right) \) 的一组基,   存在 \( {y}_{1},{y}_{2},\cdots ,{y}_{m} \in  \mathcal{X} \) ,使得
\( {f}_{i}\left( {y}_{j}\right)  = {\delta }_{ij},\;1 \leq  i,j \leq  m. \)
\end{lemma}
\begin{proof}
    考察线性连续映射 \( V : \mathcal{X} \rightarrow  {\mathbb{K}}^{m} \) 如下:
\[V : x \mapsto  \left( {\left\langle  {{f}_{1},x}\right\rangle  ,\left\langle  {{f}_{2},x}\right\rangle  ,\cdots ,\left\langle  {{f}_{m},x}\right\rangle  }\right) . \]
只要证明它是满射就够了. 如其不然, \( V\left( \mathcal{X}\right) \) 是 \( {\mathbb{K}}^{m} \) 的一个真子空
间. 由 Hahn-Banach 定理 (定理 2.4.4),存在 \( \alpha  = \left( {{\alpha }_{1},{\alpha }_{2},\cdots ,{\alpha }_{m}}\right)  \in \)
\( {\mathbb{K}}^{m} \smallsetminus  \{ 0\} \) ,使得 \( \alpha \) 有 \( V\left( \mathcal{X}\right) \) 上为 0,即
\( {\left( \alpha ,V\left( x\right) \right) }_{{\mathbb{K}}^{m}} = 0,\;\forall x \in  \mathcal{X}, \)
亦即 \( \left\langle  {\mathop{\sum }\limits_{{j = 1}}^{m}{\alpha }_{j}{f}_{j},x}\right\rangle   = 0,\;\forall x \in  \mathcal{X}. \)
从而 \( \mathop{\sum }\limits_{{j = 1}}^{m}{\alpha }_{j}{f}_{j} = \theta \) ,这与 \( {\left\{  {f}_{j}\right\}  }_{j = 1}^{m} \) 是 \( N\left( {T_\lambda}^{ * }\right) \) 的一组基矛盾.
\end{proof}
\begin{lemma}
    \( \dim N\left(T_\lambda\right)  = \dim N\left(  T_{\bar{\lambda}}^{ * }\right) \) .
\end{lemma}
\begin{proof}
   先证明$ \operatorname{span}\left\{  {{y}_{1},{y}_{2},\cdots ,{y}_{m}}\right\}  \bigcap  R\left( T_\lambda\right)=\{0\},设$$T_\lambda x=\sum\limits_{i=1}^mc_iy_i$,那么
\[c_i=\left\langle f_i,T_\lambda x\right\rangle=\left\langle T_\lambda^* f_i,x\right\rangle=0\Rightarrow T_\lambda x=0\]
   
   因为$N(T_\lambda)\subset N(T_\lambda^{**})$.所以设 \( n < m \) ,考虑
\[ \widehat{T} : \mathcal{X} = \operatorname{span}\left\{  {{x}_{1},{x}_{2},\cdots ,{x}_{n}}\right\}   \oplus  {\mathcal{X}}_{1} \rightarrow  \operatorname{span}\left\{  {{y}_{1},{y}_{2},\cdots ,{y}_{n}}\right\}   \oplus  R\left( T\right)  \hookrightarrow  \mathcal{X},  \widehat{T}\left( {\mathop{\sum }\limits_{{i = 1}}^{n}{c}_{i}{x}_{i} + y}\right)  = \mathop{\sum }\limits_{{i = 1}}^{n}{c}_{i}{y}_{i} + {Ty}.\]

根据引理4.4, \( \widehat{T} \) 是满射. 但显然有 \( {y}_{m} \in  R\left( \widehat{T}\right) \) ,矛盾,所以\( \dim N\left( T_{\bar{\lambda}}^{ * }\right)  \leq  \dim N\left( T_\lambda\right) . \)
同样.
\( \dim N\left( {T_\lambda}^{* * }\right)  \leq  \dim N\left( T_{\bar{\lambda}}^{ * }\right) . \)
因
\( \dim N\left( T_\lambda\right)  \leq  \dim N\left( {T_\lambda}^{* * }\right) , \)
所以有
\( \dim N\left( T_\lambda\right)  = \dim N\left( T_{\bar{\lambda}}^{ * }\right) . \)
\end{proof}
\begin{corollary}
          若 \(T \in  \mathcal{C}\left( X\right) ,\lambda  \neq  0,{T}_{\lambda } = {\lambda I} - T\) . 则存在闭线性子空间$\mathcal{X}_1$,和有限维子空间$\mathcal{Y}_1$,$\dim \mathcal{Y}_1=\dim N(T_\lambda)$,使得\[
\mathcal{X}=\mathcal{X}_1\oplus N(T_\lambda) =\mathcal{Y}_1\oplus R(T_\lambda)
      \]
\end{corollary}
\begin{corollary}
        若 \(T \in  \mathcal{C}\left( X\right) ,\lambda  \neq  0,{T}_{\lambda } = {\lambda I} - T\) . 则 \(\operatorname{codim}R\left( {T}_{\lambda }\right)  = \dim N\left( {T}_{\bar{\lambda }}^{ * }\right)\) ;
\end{corollary}

\begin{lemma}
        设 \(X\) 为Banach空间, \(T \in  \mathcal{C}\left( \mathcal{X}\right)\) . 若 \(\lambda  \neq  0\) ,则 \(R\left( {{\lambda I} - T}\right)\) 和 \(R\left( {{\lambda I} - {T}^{ * }}\right)\) 都是闭子空间,且
\begin{equation*}
R\left( {{\lambda I} - T}\right)  = N{\left( \bar{\lambda }I - {T}^{ * }\right) }^{ \bot  },\;R\left( {\bar{\lambda }I - {T}^{ * }}\right) { = }^{ \bot  }N\left( {{\lambda I} - T}\right) .
\end{equation*}
\end{lemma}
\begin{proof}
    $N(T_\lambda)$是$\mathcal{X}$的闭子空间,那么定义映射\[
    \widetilde{T}:\mathcal{X}/N(T_\lambda)\to R(T_\lambda),[x]\to T_\lambda x
    \]
    是既单又满的映射,考虑$\widetilde{T}^{-1}$,如果它是连续的,那么$R(T_\lambda)$是闭的,假设它不是连续的,那么存在$[x_n]\in \mathcal{X}/N(T_\lambda),[x_n]\not\to 0,T_\lambda x_n\to 0$,那么存在$\varepsilon_0>0$以及$[y_m]=[x_{n_{m}}]$
使得\[
\|[y_m]\|\geq \varepsilon_0,\textbf{取}[z_m]=\frac{[y_m]}{\|[y_m]\|},\widetilde{T}[z_m]\to 0
\]
又存在$w_m\in[z_m]$,满足\[
\|w_m\|\leqslant2 \|[z_m]\|\leqslant 2
\]
同时有\[
\widetilde{T}[w_m]=\widetilde{T}[z_m]\to 0
\]
由于$T$是紧的,那么存在$Tw_{m_k}\to z$,进而有
\[\lambda w_{m_k}=T_\lambda w_{m_k}+Tw_{m_k}\to z\]
于是\[T_\lambda z=\lambda\lim\limits_{k\to \infty}T_\lambda w_{m_k}=\lambda\lim\limits_{k\to \infty}\widetilde{T}[w_{m_k}]=0\]
得到$z\in N(T_\lambda)$,那么\[
\|[\lambda w_{m_k}]\|=\|[\lambda w_{m_k}-z]\|\leqslant \|\lambda w_{m_k}-z\|\to 0
\]
与\[\|[w_{m_k}]\|=1\]
矛盾,所以是闭的.

下证 \(R\left( {{\lambda I} - T}\right)  = N{\left( \bar{\lambda }I - {T}^{ * }\right) }^{ \bot  }\) . 显然有 \(R\left( {{\lambda I} - T}\right)  \subset  N{\left( \bar{\lambda }I - {T}^{ * }\right) }^{ \bot  }\) .假设是真包含,那么存在$x\in  N{\left( \bar{\lambda }I - {T}^{ * }\right) }^{ \bot  },x\notin R\left( {{\lambda I} - T}\right)  $,由于$R\left( {{\lambda I} - T}\right) $是闭集,由Hahn-Banach定理知存在$f\in \mathcal{X}^*$,满足\[
f(x)=\|x\|,f(R\left( {{\lambda I} - T}\right) )=0,\|f\|=1
\]
那么\[(\overline{\lambda}I-T^*)f=0\Rightarrow f\in N{\left( \bar{\lambda }I - {T}^{ * }\right) }^{ \bot  }\Rightarrow f(x)=0\Rightarrow x=0\]
与$x\notin R\left( {{\lambda I} - T}\right) $矛盾.

由于 \({T}^{ * } \in  \mathcal{C}\left( \mathcal{X}\right)\) ,由上面结论可知\begin{equation*}
\;R\left( {\bar{\lambda }I - {T}^{ * }}\right) = N{\left( {\lambda }I - {T}^{ ** }\right) }^{ \bot  }
\end{equation*}

由引理5.2.6知$\dim  N{\left( {\lambda }I - {T}^{ ** }\right) }=\dim  N{\left( {\lambda }I - {T}\right) }$,而$ N{\left( {\lambda }I - {T}\right) }\subset  N{\left( {\lambda }I - {T}^{ ** }\right) }$,所以$ N{\left( {\lambda }I - {T}\right) }=  N{\left( {\lambda }I - {T}^{ ** }\right) }$,下面证明$N{\left( {\lambda }I - {T}^{ ** }\right) }^{ \bot  }={}^\bot N(\lambda I-T)$.
\begin{equation*}
    \begin{aligned}
        f\in N{\left( {\lambda }I - {T}^{ ** }\right) }^{ \bot  }&\Leftrightarrow \left\langle x^{**},f\right\rangle=0,\forall x^{**}\in N{\left( {\lambda }I - {T}^{ ** }\right) }\\
        &\Leftrightarrow \left\langle f,x\right\rangle=0,\forall x\in N{\left( {\lambda }I - {T}\right) }
        \\
        &\Leftrightarrow f\in {}^\bot N{\left( {\lambda }I - {T}\right) }
    \end{aligned}
\end{equation*}

进而可得\begin{equation*}
R\left( {\bar{\lambda }I - {T}^{ * }}\right) { = }^{ \bot  }N\left( {{\lambda I} - T}\right) .
\end{equation*}

\end{proof}




\subsection{紧算子的谱}

关于紧算子谱集的结构, 我们有
\begin{theorem}
    设 \(\mathcal{X}\) 为Banach 空间. 若 \(T \in  \mathcal{C}\left( \mathcal{X}\right)\) ,则

(1) \(0 \in  \sigma \left( T\right)\) ,除非 \(\dim \mathcal{X} < \infty\) ;

(2) \(\sigma \left( T\right)  \smallsetminus  \{ 0\}  = {\sigma }_{p}\left( T\right)  \smallsetminus  \{ 0\}\) ;

(3) \({\sigma }_{p}\left( T\right)\) 至多以 0 为聚点.
\end{theorem}
\begin{proof}
    (1) 用反证法. 若 \(\dim \mathcal{X} =  + \infty\) 且 \(0 \notin  \sigma \left( T\right)\) ,即 \({T}^{-1}\) 存在,由于 \(T \in\)  \(\mathcal{C}\left( \mathcal{X}\right)\) ,从而 \(I = T \circ  {T}^{-1} \in  \mathcal{C}\left( \mathcal{X}\right)\) . 故得到矛盾.

(2) 若 \(\lambda  \in  \sigma \left( T\right)  \smallsetminus  \{ 0\}\) ,必有$N(T_\lambda )\neq \{0\}$,否则有 \(R\left( {T}_{\lambda }\right)  \neq  \mathcal{X}\) ,其中 \({T}_{\lambda } = {\lambda I} - T\) . 从而由推论5.2.2知, \(N\left( {T}_{\lambda }\right)  \mathrel{\text{ := }} \left\{  {x \in  \mathcal{X} \mid  {T}_{\lambda }x = 0}\right\}   \neq  0\) . 因此 \(\left( {{\lambda I} - T}\right) x = 0\) 有非零解. 故 \(\lambda\) 为 \(T\) 的特征值,即 \(\lambda  \in  {\sigma }_{p}\left( T\right)  \smallsetminus  \{ 0\}\) .

(3)用反证法. 若有 \({\lambda }_{n} \in  {\sigma }_{p}\left( T\right)  \smallsetminus  \{ 0\} \left( {n = 1,2,\cdots }\right) ,{\lambda }_{n} \neq  {\lambda }_{m}\) (当 \(n \neq  m\) ),并且 \({\lambda }_{n} \rightarrow  \lambda  \neq  0\left( {n \rightarrow  \infty }\right)\) ; 那末
\begin{equation*}
\exists {x}_{n} \in  N\left( {{\lambda }_{n}I - T}\right)  \smallsetminus  \{ 0\} ,\;n = 1,2,\cdots .
\end{equation*}

我们有:

\({1}^{ \circ  }\text{ 、 }\left\{  {{x}_{1},{x}_{2},\cdots ,{x}_{n}}\right\}\) 是线性无关的. 事实上,可用数学归纳法证明,设结论对 \(n\) 已成立. 若有
\begin{equation*}
{x}_{n + 1} \in  N\left( {{\lambda }_{n + 1}I - T}\right)  \smallsetminus  \{ 0\} ,
\end{equation*}

使得 \({x}_{n + 1} = \mathop{\sum }\limits_{{i = 1}}^{n}{a}_{i}{x}_{i}\) ,则有
\begin{equation*}
{\lambda }_{n + 1}{x}_{n + 1} = T{x}_{n + 1} = \mathop{\sum }\limits_{{i = 1}}^{n}{a}_{i}{\lambda }_{i}{x}_{i}
\end{equation*}

从而
\begin{equation*}
\mathop{\sum }\limits_{{i = 1}}^{n}{a}_{i}\left( {{\lambda }_{n + 1} - {\lambda }_{i}}\right) {x}_{i} = 0.
\end{equation*}

由归纳法假设知 \(\left\{  {{x}_{1},\cdots ,{x}_{n}}\right\}\) 是线性无关的,因而有
\begin{equation*}
\left( {{\lambda }_{n + 1} - {\lambda }_{i}}\right) {a}_{i} = 0 \Rightarrow  {a}_{i} = 0,\;i = 1,2,\cdots ,n.
\end{equation*}

这与 \({x}_{n + 1} \neq  0\) 矛盾. 因此 \(\left\{  {{x}_{1},{x}_{2},\cdots ,{x}_{n + 1}}\right\}\) 是线性无关的.

\({2}^{ \circ  }\) 、若令 \({E}_{n} = \operatorname{span}\left\{  {{x}_{1},\cdots ,{x}_{n}}\right\}\) ,则 \({E}_{n} \subsetneqq  {E}_{n + 1}\) ,应用Riesz引理可得
\begin{equation*}
\exists {y}_{n + 1} \in  {E}_{n + 1}\text{ ,使得 }\begin{Vmatrix}{y}_{n + 1}\end{Vmatrix} = 1\text{ ,且 }\operatorname{dist}\left( {{y}_{n + 1},{E}_{n}}\right)  \geq  \frac{1}{2}\text{ . }
\end{equation*}

从而对任意 \(n,p \in  \mathbb{N}\) ,有
\begin{equation*}
\begin{Vmatrix}{\frac{1}{{\lambda }_{n + p}}T{y}_{n + p} - \frac{1}{{\lambda }_{n}}T{y}_{n}}\end{Vmatrix} = \begin{Vmatrix}{{y}_{n + p} - \left( {{y}_{n + p} - \frac{1}{{\lambda }_{n + p}}T{y}_{n + p} + \frac{1}{{\lambda }_{n}}T{y}_{n}}\right) }\end{Vmatrix} \geq  \frac{1}{2},
\end{equation*}

这是因为
\begin{equation*}
{y}_{n + p} - \frac{1}{{\lambda }_{n + p}}T{y}_{n + p} + \frac{1}{{\lambda }_{n}}T{y}_{n} \in  {E}_{n + p - 1}.
\end{equation*}

这与 \(T\) 的紧性矛盾.

\end{proof}
\begin{note}
    该定理表明: 对于无穷维空间上的紧算子 \(T\) ,只有三种可能情形:

(1) \(\sigma \left( T\right)  = \{ 0\}\) ;

(2) \(\sigma \left( T\right)  = \left\{  {0,{\lambda }_{1},\cdots ,{\lambda }_{n}}\right\}\) ;

(3) \(\sigma \left( T\right)  = \left\{  {{\lambda }_{1},{\lambda }_{2},\cdots ,{\lambda }_{n},\cdots ,0}\right\}\) ,其中 \({\lambda }_{n} \rightarrow  0\) .

\end{note}
\begin{example}
    证明 (4.12) 中定义的紧算子 \(T\) 仅有零谱.
\end{example}
\begin{proof}
 \(T\) 为紧算子,从而其谱集由零和特征值组成. 下证 \(T\) 没有非零特征值. 用反证法,若不然,则存在 \(\lambda  \neq  0\) 及非零函数 \(x \in  C\left\lbrack  {a,b}\right\rbrack\) 使得
\begin{equation}
\left( {Tx}\right) \left( s\right)  = {\int }_{a}^{s}x\left( t\right) {\dif t} = {\lambda x}\left( s\right) . \label{4.13}
\end{equation}

由于方程左边为 \(s\) 的可微函数,故右边也是 \(s\) 的可微函数. 因此,在方程 (4.13) 两边求导得
\begin{equation*}
x\left( s\right)  = \lambda {x}^{\prime }\left( s\right) .
\end{equation*}

由于该方程的所有非零解形如 \(x\left( s\right)  = c{e}^{s/\lambda },c \neq  0\) . 特别地, \(x\left( a\right)  = c{e}^{a/\lambda } \neq  0\) . 在 (4.13)中令 \(s = a\) 得 \(0 = {\lambda x}\left( a\right)\) . 与 \(\lambda  \neq  0\) 及 \(c \neq  0\) 矛盾. 从而 \(T\) 仅有零谱.
\end{proof}


\subsection{不变子空间}

对于紧算子的不变子空间, 我们有
\begin{definition}
    设 \(\mathcal{X}\) 为Banach空间, \(M \subset  \mathcal{X}\) 称为算子 \(T \in  \mathcal{L}\left( \mathcal{X}\right)\) 的不变子空间, 如果 \({TM} \subset  M\) .

\end{definition}

由定义立即可得
\begin{proposition}
    设 \(\mathcal{X}\) 为Banach空间, \(T \in  \mathcal{L}\left( \mathcal{X}\right)\) ,那末

(1) \(\{ 0\} ,\mathcal{X}\) 都是 \(T\) 的不变子空间;

(2)若 \(M\) 是 \(T\) 的不变子空间,则 \(\bar{M}\) 也是 \(T\) 的不变子空间;

( 3 )若 \(\lambda  \in  {\sigma }_{p}\left( T\right)\) ,即 \(\lambda\) 是 \(T\) 的本征值,则 \(N\left( {{\lambda I} - T}\right)\) 是 \(T\) 的不变子空间;

(4)对任意 \(y \in  \mathcal{X}\) ,若记 \({L}_{y} \mathrel{\text{ := }} \{ P\left( T\right) y \mid  P\) 是任意多项式 \(\}\) ,则 \({L}_{y}\) 是 \(T\) 的不变子空间.

\end{proposition}

当 \(\dim \mathcal{X} = \infty\) 时,对任意 \(T \in  \mathcal{L}\left( \mathcal{X}\right)\) ,问是否存在 \(T\) 的非平凡的闭不变子空间? 这里,非平凡是指 \(M \neq  \{ 0\}\) 或 \(\mathcal{X}\) . 这是一个长期未解决的公开问题. 直到1984年才由C.J.Read给出否定回答:存在无穷维Banach空间 \(\mathcal{X}\) 及线性算子 \(T \in\)  \(\mathcal{L}\left( \mathcal{X}\right)\) 使得 \(T\) 没有非平凡的闭不变子空间[Read]. 接下来,我们自然会问,当 \(\mathcal{X}\) 或算子 \(T \in  \mathcal{L}\left( \mathcal{X}\right)\) 满足一定条件时上述公开问题答案如何？对此,我们有
\begin{theorem}
    若 \(\dim \mathcal{X} \geq  2\) ,则对任意 \(T \in  \mathcal{C}\left( \mathcal{X}\right) ,T\) 必有非平凡的闭不变子空间.
\end{theorem}
\begin{proof}
    我们不妨设 \(\dim \mathcal{X} = \infty ,T \neq  0\) ,并且 \({\sigma }_{p}\left( T\right)  \smallsetminus  \{ 0\}  = \varnothing\) . 于是由定理5.2.1有 \(\sigma \left( T\right)  = \{ 0\}\) . 用反证法,若不然,即 \(T\) 没有非平凡的闭不变子空间,则对任意 \(y \in  X \smallsetminus  \{ 0\}\) ,由命题5.2.1中定义的 \({L}_{y}\) 蕴含着
\begin{equation*}
\overline{{L}_{y}} = \mathcal{X}.
\end{equation*}

不妨设 \(\parallel T\parallel  = 1\) ,那末存在 \({x}_{0} \in  \mathcal{X}\) 使得 \(\begin{Vmatrix}{T{x}_{0}}\end{Vmatrix} > 1\) 且 \(\begin{Vmatrix}{x}_{0}\end{Vmatrix} > 1\) . 取 \(C \mathrel{\text{ := }}\)  \(\overline{T\left( {B\left( {{x}_{0},1}\right) }\right) }\) ,则有 \(C\) 为紧集,并且
\begin{equation*}
0 \notin  C\text{ . }
\end{equation*}

同时对任意 \({y}_{0} \in  C\) ,存在多项式 \({P}_{{y}_{0}}\left( T\right)\) ,使得
\begin{equation*}
\begin{Vmatrix}{{P}_{{y}_{0}}\left( T\right) {y}_{0} - {x}_{0}}\end{Vmatrix} < 1,
\end{equation*}

从而有 \({\delta }_{{y}_{0}} > 0\) ,使得
\begin{equation*}
\begin{Vmatrix}{{P}_{{y}_{0}}\left( T\right) y - {x}_{0}}\end{Vmatrix} < 1,\;\forall y \in  B\left( {{y}_{0},{\delta }_{{y}_{0}}}\right) .
\end{equation*}

由于 \(C\) 紧致,故有有限覆盖
\begin{equation*}
\mathop{\bigcup }\limits_{{i = 1}}^{n}B\left( {{y}_{i},{\delta }_{i}}\right)  \supset  C
\end{equation*}

其中 \({\delta }_{i} \mathrel{\text{ := }} {\delta }_{{y}_{i}}\left( {i = 1,2,\cdots ,n}\right)\) . 从而对任意 \(y \in  C\) ,存在 \({i}_{1}\left( {1 \leqslant  {i}_{1} \leqslant  n}\right)\) ,使得
\begin{equation}
\begin{Vmatrix}{{P}_{{i}_{1}}\left( T\right) y - {x}_{0}}\end{Vmatrix} < 1, \label{4.14}
\end{equation}

这里记 \({P}_{i}\left( T\right)  \mathrel{\text{ := }} {P}_{{y}_{i}}\left( T\right) \left( {i = 1,2,\cdots ,n}\right)\) .

由 上式知, \({P}_{{i}_{1}}\left( T\right) y \in  B\left( {{x}_{0},1}\right)\) ,所以 \(T{P}_{{i}_{1}}\left( T\right) y \in  C\) ,又存在 \({i}_{2}\left( {1 \leqslant  {i}_{2} \leqslant  }\right.\)  \(n)\) 使得
\begin{equation*}
\begin{Vmatrix}{{P}_{{i}_{2}}\left( T\right) T{P}_{{i}_{1}}\left( T\right) y - {x}_{0}}\end{Vmatrix} < 1.
\end{equation*}

由于 \({P}_{{i}_{1}}\) 是与 \(T\) 可交换的多项式,故有
\begin{equation*}
\begin{Vmatrix}{{P}_{{i}_{2}}\left( T\right) {P}_{{i}_{1}}\left( T\right) {Ty} - {x}_{0}}\end{Vmatrix} < 1.
\end{equation*}

依此下去,存在 \({i}_{1},{i}_{2},\cdots ,{i}_{k},\cdots\) 使得
\begin{equation*}
\begin{Vmatrix}{\mathop{\prod }\limits_{{j = 1}}^{{k + 1}}{P}_{{i}_{j}}\left( T\right) \left( {{T}^{k}y}\right)  - {x}_{0}}\end{Vmatrix} < 1.
\end{equation*}

从而
\begin{equation*}
\left\|\left( {\mathop{\prod }\limits_{{j = 1}}^{{k + 1}}{P}_{{i}_{j}}\left( T\right) }\right)\left( {{T}^{k}y}\right) \right\|> \begin{Vmatrix}{x}_{0}\end{Vmatrix} - 1.
\end{equation*}

设 \(\mu  = \mathop{\max }\limits_{{1 \leqslant  i \leqslant  n}}\begin{Vmatrix}{{P}_{i}\left( T\right) }\end{Vmatrix}\) ,则有
\begin{equation*}
\begin{Vmatrix}{x}_{0}\end{Vmatrix} - 1 \leqslant  {\mu }^{k + 1}\begin{Vmatrix}{{T}^{k}y}\end{Vmatrix},\;\mu  > 0,k \in  \mathbb{N}.
\end{equation*}

因此,
\begin{equation*}
\frac{1}{\mu }{\left( \frac{\begin{Vmatrix}{x}_{0}\end{Vmatrix} - 1}{\mu \parallel y\parallel }\right) }^{\frac{1}{k}} \leqslant  {\left( \frac{\begin{Vmatrix}{T}^{k}y\end{Vmatrix}}{\parallel y\parallel }\right) }^{\frac{1}{k}} \leqslant  {\begin{Vmatrix}{T}^{k}\end{Vmatrix}}^{\frac{1}{k}}.
\end{equation*}

当 \(k \rightarrow  \infty\) 时,上式左端极限为 \(\frac{1}{\mu }\) ,而右端极限是 0 (此处用到 \(T\) 的谱半径 \({r}_{\sigma }\left( T\right)  =\) 0及 \({r}_{\sigma }\left( T\right)  = \mathop{\lim }\limits_{{k \rightarrow  \infty }}{\begin{Vmatrix}{T}^{k}\end{Vmatrix}}^{\frac{1}{k}}\) ). 得到矛盾.

\end{proof}

\subsection{紧算子的结构}

在线性代数中, 从有关线性变换的结构的研究已经知道, 任何矩阵都可以分解为Jordan标准型. 一个自然地问题是, 一般Banach空间上的有界线性算子, 是否也有类似结果? 对紧算子, 我们有
\begin{theorem}
    设 \(T \in  \mathcal{C}\left( \mathcal{X}\right) ,\lambda  \neq  0,{T}_{\lambda } = {\lambda I} - T\) . 则存在非负整数 \(p\) 使得
\begin{equation*}
\mathcal{X} = N\left( {T}_{\lambda }^{p}\right)  \oplus  R\left( {T}_{\lambda }^{p}\right) ,
\end{equation*}

且 \({T}^{\prime } \mathrel{\text{ := }} {\left. {T}_{\lambda }\right| }_{R\left( {T}_{\lambda }^{p}\right) }\) 有线性有界逆算子.

\end{theorem}
\begin{proof}
    用 \(p\) 记使得 \(N\left( {T}_{\lambda }^{k}\right)  = N\left( {T}_{\lambda }^{k + 1}\right)\) 成立的最小整数 \(p\) ,由引理4.3知这样的 \(p\) 存在,通常称为算子 \({T}_{\lambda }\) 的零链长. 下面分四步来证明定理.

(1)存在$p<\infty$, \(R\left( {T}_{\lambda }^{p}\right)  = R\left( {T}_{\lambda }^{p + 1}\right)\) .事实上,我们有下列链包含关系:
\begin{equation*}
\mathcal{X} \supseteq  R\left( {T}_{\lambda }\right)  \supseteq  R\left( {T}_{\lambda }^{2}\right)  \supseteq  \cdots
\end{equation*}

假设都是真包含,那么由Reisz引理知存在$y_n\in R(T_\lambda ^n)$,满足\[
\|y_n\|=1,\operatorname{dist}(y_n,R(T_\lambda^{n+1}))\geqslant \frac{1}{2}
\]
下面考虑
\[
\|Ty_n-Ty_{n+p}\|=\|\lambda y_n-T_\lambda y_n-\lambda y_{n+p}+T_\lambda y_{n+p}\|\geqslant |\lambda |\operatorname{dist}(y_n,R(T_\lambda^{n+1}))\geqslant \frac{|\lambda|}{2}
\]
与$T$是紧的矛盾.

 记使得 \(R\left( {T}_{\lambda }^{k}\right)  = R\left( {T}_{\lambda }^{k + 1}\right)\) 成立的最小整数为 \(q\) ,称为 \({T}_{\lambda }\) 的像链长. 由引理4.2 及 \(R\left( {T}_{\lambda }^{q}\right)  = R\left( {T}_{\lambda }^{q + 1}\right)\) 知,
\begin{equation*}
\dim N\left( {T}_{\lambda }^{q}\right)  = \operatorname{codim}R\left( {T}_{\lambda }^{q}\right)  = \operatorname{codim}R\left( {T}_{\lambda }^{q + 1}\right)  = \dim N\left( {T}_{\lambda }^{q + 1}\right) .
\end{equation*}

由于 \(\dim N\left( {T}_{\lambda }^{q}\right)  < \infty\) ,故 \(N\left( {T}_{\lambda }^{q}\right)  = N\left( {T}_{\lambda }^{q + 1}\right)\) . 由零链长的定义知, \(p \leqslant  q\) . 类似可证 \(p \geq  q\) . 从而 \(p = q\) .

(2) \(N\left( {T}_{\lambda }^{p}\right) \bigcap R\left( {T}_{\lambda }^{p}\right)  = 0\) . 事实上,若有 \(y \in  N\left( {T}_{\lambda }^{p}\right) \bigcap R\left( {T}_{\lambda }^{p}\right)\) ,则必存在 \(x \in\)  \(\mathcal{X}\) 使得 \(y = {T}_{\lambda }^{p}x\) ,且有 \({T}_{\lambda }^{p}y = 0\) . 从而
\begin{equation*}
x \in  N\left( {T}_{\lambda }^{2p}\right)  = N\left( {T}_{\lambda }^{p}\right)  \Rightarrow  y = {T}_{\lambda }^{p}x = 0.
\end{equation*}

(3) \(\mathcal{X} = N\left( {T}_{\lambda }^{p}\right)  \oplus  R\left( {T}_{\lambda }^{p}\right)\) . 事实上,对任意 \(x \in  \mathcal{X}\) 有
\begin{equation*}
{T}_{\lambda }^{p}x \in  R\left( {T}_{\lambda }^{p}\right)  = R\left( {T}_{\lambda }^{2p}\right) ,
\end{equation*}

因此存在 \(u \in  \mathcal{X}\) 使得 \({T}_{\lambda }^{2p}u = {T}_{\lambda }^{p}x\) . 令 \(y \mathrel{\text{ := }} {T}_{\lambda }^{p}u \in  R\left( {T}_{\lambda }^{p}\right)\) ,则由 \({T}_{\lambda }^{p}y = {T}_{\lambda }^{p}x\) 有 \(z \mathrel{\text{ := }}\)  \(x - y \in  N\left( {T}_{\lambda }^{p}\right)\) . 于是可得
\begin{equation*}
x = y + z,\;y \in  R\left( {T}_{\lambda }^{p}\right) ,z \in  N\left( {T}_{\lambda }^{p}\right) .
\end{equation*}

(4) \({T}^{\prime } \mathrel{\text{ := }} {\left. {T}_{\lambda }\right| }_{R\left( {T}_{\lambda }^{p}\right) }\) 有线性有界逆算子. 事实上,由前面结论(1)知
\begin{equation*}
R\left( {T}_{\lambda }^{p}\right)  = R\left( {T}_{\lambda }^{p + 1}\right) .
\end{equation*}

由此可知 \({T}^{\prime }\) 为满射. 下证 \({T}^{\prime }\) 是单射. 实际上,若 \(y \in  R\left( {T}_{\lambda }^{p}\right)\) 且 \({T}_{\lambda }y = 0\) ,则存在 \(x \in\)  \(\mathcal{X}\) 使得 \(y = {T}_{\lambda }^{p}x\) ,从而
\begin{equation*}
x \in  N\left( {T}_{\lambda }^{p + 1}\right)  = N\left( {T}_{\lambda }^{p}\right) ,
\end{equation*}

即有 \(y = 0\) . 于是由Banach逆算子定理知结论成立.

\end{proof}

\subsection{习 题}
\begin{problemset}
  \item 给定数列 \(\{a_n\}\)，在空间 \({\ell }^{2}\) 上定义算子
  \[
  T : (x_{1},x_{2},\cdots) \mapsto (a_{1}x_{1},a_{2}x_{2},\cdots).
  \]
  (1) 证明：\(T \in  \mathcal{L}\left( {\ell }^{2}\right)\) 当且仅当存在 \(M > 0\) 使得 \(|a_n| \leqslant M\).  
  (2) 若 \(T \in  \mathcal{L}\left( {\ell }^{2}\right)\)，求 \(\sigma(T)\) 并判别谱点类型.
\begin{proof}
    (1)若$|a_n|\leqslant M,\forall x\in \ell^2$,有\[
    \|Tx\|\leqslant M\|x\|
    \]
    若$\{a_n\}$无界,那么$\forall k\in \mathbb{N}^+,\exists |a_{n_k}|>k,取x_k=(\delta_i^{n_k})\in \ell^2$,那么有\[ 
    \|Tx_k\|=|a_{n_k}|>k\]
    进而$T$不是有界算子.
    
    (2)记$M=\sup\{|a_n|\}$,由(1)可知$\|T\|=M,\|T_k\|=M^k$,由Gelfand定理知$r_\sigma(T)=M$.

    若$|\lambda|=M$
\end{proof}
  \item 在 \(C[a,b]\) 中，考虑映射
  \[
  T : x(t) \mapsto \int_{a}^{t}x(s)\,\dif s,\quad \forall x(t) \in C[a,b].
  \]
  (1) 证明 \(T\) 是紧算子.  
  (2) 证明 \({C}_{a} := \{u \in C[0,1] \mid u(x) = 0,\ \forall x \in [0,a]\}\) 为算子 \(T\) 的非平凡不变子空间.
\begin{proof}
    (1)由习题5.1.6可知.

    (2)$\forall  u\in C_a$
    \[
    Tu(t)=\int_0^t u(s)\dif s,Tu(t)=0,\forall t\in [0,a]
    \]
    
\end{proof}
  \item 设 \(A,B \in L(X)\)，且 \(AB = BA\).证明：  
  (1) \(R(A)\) 和 \(N(A)\) 都是 \(B\) 的不变子空间;  
  (2) 对任意 \(n \in \mathbb{N}\)，\(R(B^{n})\) 和 \(N(B^{n})\) 都是 \(B\) 的不变子空间.
\begin{proof}
    (1) $\forall y\in R(A),y=Ax,By=ABx\in R(A)$
    同理可证.
\end{proof}
  \item 设 \([a,b]\) 为 \(\mathbb{R}\) 中有界闭区间，\(G \in L^{2}([a,b]\times[a,b])\)，\(G \neq 0\).定义算子$K: L^2[a,b]\to L^2[a,b]$为
  \[
  (Ku)(x) := \int_{a}^{b} G(x,y)u(y)\,\dif y,\quad x \in [a,b].
  \]
  证明：  
  
  (1) \(K\) 为紧算子;  
  
  (2) \(\|K\| \le \|G\|_{L^{2}}\);  
  
  (3) 伴随算子 \(K^{*}\) 的核为 \(G^{*}(x,y) = \overline{G(y,x)}\);  
  
  (4) 若 \(K_{1}, K_{2}\) 分别由核 \(G_{1}, G_{2}\) 定义，则
  \[
  (K_{2}K_{1})(u)(x) = \int_{a}^{b} G(x,y)u(y)\,\dif y,
  \quad G(x,y) = \int_{a}^{b} G_{2}(x,z)G_{1}(z,y)\,\dif z.
  \]
\begin{proof}
    (1) 由习题5.1.7知.

    (2)由Holder不等式得
\begin{equation*}
        \begin{aligned}
              \|(Ku)(x)\| &= \left(\int_a^b\left|\int_{a}^{b} G(x,y)u(y)\dif y\right|^2\dif x\right)^{\frac{1}{2}}\leqslant  \left(\int_a^b\left(\int_{a}^{b} G(x,y)^2\dif y\int _a^bu(y)^2\dif y\right)\dif x\right)^{\frac{1}{2}}\\
              &= \left(\int_a^b\left(\int_{a}^{b} G(x,y)^2\dif y\|u\|^2\right)\dif x\right)^{\frac{1}{2}}= \|u\|\|G\|
        \end{aligned}
    \end{equation*}
    (3) 由于$(L^2)^*=L^2$,所以设$G^*(x,y)\in L\left(\left(L^2\right)^*\right)$,那么有\begin{equation*}
        \begin{aligned}
            \left\langle K^*u,v\right\rangle=      \left\langle u ,Kv\right\rangle  ,\forall u \in (L^2[a,b])^*,v\in L^2[a,b]
        \end{aligned}
    \end{equation*}
   由Fubini定理得
\begin{equation*}
        \begin{aligned}
             \left\langle u ,Kv\right\rangle=\int_a^b u(x)\overline{\int_{a}^{b} G(x,y)v(y)\,\dif y}\dif x= \int_a^b \overline{v(x)}{\int_{a}^{b} \overline{G(y,x)}u(y)\,\dif y}\dif x
        \end{aligned}
    \end{equation*}
    所以 \(K^{*}\) 的核为 \(G^{*}(x,y) = \overline{G(y,x)}\).

    (4)\begin{equation*}
        \begin{aligned}
              (K_{2}K_{1})(u)(x)& = \int_a^b G_2(x,y)K_1u(y)\dif y= \int_a^b G_2(x,y)\int_a^b G_1(y,z)u(z)\dif z \dif y\\
              &= \int_a^bu(z) \int_a^b G_2(x,y)G_1(y,z)\dif y \dif z = \int_a^bu(y) \int_a^b G_2(x,z)G_1(z,y)\dif z \dif y
        \end{aligned}
    \end{equation*}
\end{proof}
  \item* 设 \(H = L^{2}[0,1]\).对每个 \(a \in [0,1]\)，定义
  \[
  H_{a} := \{u \in L^{2}[0,1] \mid u(x)=0,\ \forall x \in [0,a]\}.
  \]
  证明：对任意 \(a \in [0,1]\)，\(H_{a}\) 为 (4.12) 中所定义算子 \(T\) 的不变子空间.反之，证明 \(H_{a}\) 是 \(T\) 唯一的不变子空间.

  \item 设 \(\mathcal{X}\) 为 Banach 空间，\(S \in \mathcal{C}(\mathcal{X})\).  
  (1) 对任意满足 \(\operatorname{Ker}(S - \lambda I) = \{0\}\) 的非零 \(\lambda \neq 0\)，存在 \(A_{j} \in \mathcal{L}(\mathcal{X},\mathcal{X})\)，使得当 \(z\) 很小时，
  \[
  (S - (\lambda + z)I)^{-1} = \sum_{j=-n}^{\infty} A_{j} z^{j}
  \]
  为收敛幂级数展开.  
  (2) 对不满足 \(\operatorname{Ker}(S - \lambda I) = \{0\}\) 的非零 \(\lambda \neq 0\)，存在 \(A_{j} \in \mathcal{L}(X,X)\)，使得当 \(z\) 很小时，
  \[
  (S - (\lambda + z)I)^{-1} = \sum_{j=-n}^{\infty} A_{j} z^{j}
  \]
  为 Laurent 展式，其中 \(-A_{-1}\) 为沿
\[
  F_{\lambda} = \bigcap_{k=1}^{\infty}(S - \lambda I)^{k}\mathcal{X}
  \quad \text{到} \quad
  N_{\lambda} = \bigcup_{k=1}^{\infty}\operatorname{Ker}(S - \lambda I)^{k}
  \]
  的投影.
\begin{proof}
    
\end{proof}
  \item* 设 \(Q\) 为紧致 Hausdorff 空间，\(\mathcal{X} = C(Q)\) 为 \(Q\) 上连续实值函数的集合.称线性算子 \(K : \mathcal{X} \to \mathcal{X}\) 为严格正的，若对任意非负且非零 \(\rho \in \mathcal{X}\)，有 \(K\rho\) 为正函数.证明：  
  (1) 若 \(K\) 为严格正的紧算子，则 \(K\) 具有重数为 1 的正本征值 \(\sigma\)，且相应本征函数为正函数;  
  (2) \(K\) 的其它本征值 \(\mu\) 都满足 \(|\mu| < \sigma\).

  \item* 证明二阶椭圆偏微分算子
  \[
  -\Delta = \sum_{i=1}^{n} \frac{\partial^{2}}{\partial x_{i}^{2}}
  \]
  具有最小正本征值，且相应的本征函数为正函数.
  \item 如果 $M$ 是 $B^*$ 空间 $\mathcal{X}$ 的闭子空间，求证：
$$
\left({ }^{\perp}M\right)^{\perp} =M .
$$

\begin{proof}
    根据 ${ }^{\perp}\left(M^{\perp}\right)$ 的定义，
$$
x \in\left({ }^{\perp}M\right)^{\perp} \Leftrightarrow\langle f, x\rangle=0 \quad\left(\forall f \in{}^{\perp} M\right) .
$$
一方面，$\forall x \in M, f \in M^{\perp}$ ，都有 $\langle f, x\rangle=0$ ，因此 $\left({ }^{\perp}M\right)^{\perp} $.

另一方面，要证 $\left({ }^{\perp}M\right)^{\perp} \subseteq M$ ，只要证
$$
\forall x_1 \in \mathcal{X} \backslash M \Rightarrow x_1 \notin\left({ }^{\perp}M\right)^{\perp}  .
$$
事实上，任意给定 $x_1 \in \mathcal{X} \backslash M$ 。因为 $M$ 是 $\mathcal{X}$ 的闭子空间，所以
$$
d\left(x_1, M\right) \xlongequal{\text { def }} \inf _{v \in M}\left\|x_1-z\right\|>0 .
$$
根据Hanh-Banach定理 ，存在 $f \in \mathcal{X}^*$ ，使得
$$
\left\{\begin{array}{l}
f\left(x_1\right)=d\left(x_1, M\right), \\
\|f\|=1, \\
f(x)=0(\forall x \in M), \text { 即 } f \in M^{\perp} .
\end{array}\right.
$$
因为 $f\left(x_1\right) \neq 0$ ，所以 $x_1 \notin^{\perp}\left(M^{\perp}\right)$ ．

综合以上两方面得到
$$
\left({ }^{\perp}M\right)^{\perp} =M .
$$


\end{proof}
\item  如果 $S \subset \mathcal{X}, M$ 是由 $S$ 张成的闭子空间，求证：$M^{\perp}=S^{\perp}$ ，并且 $M={ }^{\perp}\left(S^{\perp}\right)$ 。
\begin{proof}
    因为 $M$ 闭，据上题第二个论断 $M={ }^{\perp}\left(S^{\perp}\right)$ ，是第一个论断 $M^{\perp}=S^{\perp}$ 的结果。

至于第一个论断，一方面，因为 $S \subseteq M$ ，所以显然有
$$
M^{\perp} \subseteq S^{\perp} .
$$
另一方面，如果 $f \in S^{\perp}$ ，那么
$$
\langle f, x\rangle=0, \quad \forall x \in S .
$$
对 $\forall x \in M$ ，因为 $M$ 是由 $S$ 张成的闭子空间，所以 $\exists x_n \in S$ ，使得 $x_n \rightarrow x$ ，于是，由 $f$ 的连续性，令 $n \rightarrow \infty$ ，
$$
\left\langle f, x_n\right\rangle=0 \Longrightarrow\langle f, x\rangle=0 \Longrightarrow f \in M^{\perp} .
$$
故有 $S^{\perp} \subseteq M^{\perp}$ ．

综合以上两方面得到 $M^{\perp}=S^{\perp}$ 。
\end{proof}
\item 设 $\mathcal{X}, \mathcal{Y}$ 是 $B^*$ 空间，$T \in \mathcal{L}(\mathcal{X}, \mathcal{Y})$ ，求证：$\mathcal{R}(T)= \mathcal{N}\left(T^*\right)^{\perp}$ 的充分必要条件是 $\mathcal{R}(T)$ 在 $\mathcal{Y}$ 中闭．
\begin{proof}
    首先我们证明 ${ }^{\perp} \mathcal{R}(T)=\mathcal{N}\left(T^*\right)$ ，事实上，
\[f\in {}^\perp \mathcal{R}(T)\Leftrightarrow f (Tx)=0,\forall x\in \mathcal{X}\Leftrightarrow T^*f(x)=0,\forall x\in \mathcal{X}\Leftrightarrow f\in \mathcal{N}(T^*)\]
可知 ${ }^\perp \mathcal{R}(T)=\mathcal{N}\left(T^*\right)$ ．

进一步，若 $\mathcal{R}(T)$ 在 $\mathcal{Y}$ 中闭，在上上例中取 $M=\mathcal{R}(T)$ ，则有
$$
\mathcal{N}\left(T^*\right)^{\perp}=\left[{ }^{\perp} \mathcal{R}(T)\right]^{\perp}=\mathcal{R}(T) .
$$


反过来，若 $\mathcal{R}(T)=\mathcal{N}\left(T^*\right)^{\perp}$ ，则有
$$
\overline{\mathcal{R}(T)}=\left[{ }^{\perp} \mathcal{R}(T)\right]^{\perp}=\mathcal{N}\left(T^*\right)^{\perp} .
$$


故有 $\mathcal{R}(T)=\overline{\mathcal{R}(T)}$ ，即 $\mathcal{R}(T)$ 在 $\mathcal{Y}$ 中闭．
\end{proof}
\item 设 $\mathcal{X}, \mathcal{Y}$ 是 Banach 空间，$T \in \mathcal{L}(\mathcal{X}, \mathcal{Y}), \mathcal{R}(T)$ 在 $\mathcal{Y}$ 中闭．求证：$\mathcal{R}\left(T^*\right)={ }^{\perp} \mathcal{N}(T)$ ．
\begin{proof}
    首先证明 $\mathcal{R}\left(T^*\right) \subset^{\perp} \mathcal{N}(T)$ ．这就是说，对任意给定的 $f \in \mathcal{R}\left(T^*\right)$ ，要证明该 $f \in{ }^{\perp} \mathcal{N}(T)$ ，就是要证
$$
\langle f, x\rangle=0, \quad \forall x \in \mathcal{N}(T) .
$$
事实上，对该 $f \in \mathcal{R}\left(T^*\right)$ ，按定义，$\exists y^* \in \mathcal{Y}^*$ ，使得 $T^* y^*=f$ ．所以对 $\forall x \in \mathcal{N}(T)$ ，有 $T x=0$ ，从而
$$
f(x)=\left\langle T^* y^*, x\right\rangle=\left\langle y^*, T x\right\rangle=0
$$

反过来，要证明 ${ }^{\perp} \mathcal{N}(T) \subset \mathcal{R}\left(T^*\right)$ ．为此，任取 $x^* \in{ }^{\perp} \mathcal{N}(T)$ ，作 $\mathcal{R}(T)$ 上的线性泛函
$$
r^*(y)=x^*(x) \quad \text { (当 } y=T x \text { 时), }
$$


则 $r^*(y)$ 的值是唯一确定的．事实上，当 $y=\theta$ 时，因为 $y=T x$ ，所以 $x= \theta$ ，即 $x \in \mathcal{N}(T)$ 。而 $x^* \in{ }^{\perp} \mathcal{N}(T)$ ，故 $x^*(x)=0$ 。 这表明 $r^*(y)$ 在 $y=\theta$ 的值是唯一确定的，因此在任何点处的值也唯一确定。

下面证明 $r^*$ 在 $\mathcal{R}(T)$ 上是有界的．事实上，$\forall z \in \mathcal{N}(T)$ ，
$$
r^*(y)=x^*(x)=x^*(x-z) \Longrightarrow\left|r^*(y)\right| \leqslant\left\|x^*\right\|\|x-z\|,
$$
两边对 $z \in \mathcal{N}(T)$ 取下确界，即知
$$
\left|r^*(y)\right| \leqslant\left\|x^*\right\| d(x, \mathcal{N}(T)) .
$$
进一步，因为 $\mathcal{R}(T)$ 闭，根据Hanh-Banach定理$\exists a>0$ ，使得
$$
d(x, \mathcal{N}(T)) \leqslant a\|T x\|=a\|y\| .
$$
联合上面两式即得
$$
\left|r^*(y)\right| \leqslant a\left\|x^*\right\|\|y\| .
$$
这样，我们便证明了 $r^*$ 是子空间 $\mathcal{R}(T)$ 上的有界线性泛函．

根据 Hahn－Banach 定理，可将 $r^*$ 延拓到整个空间 $\mathcal{Y}$ 上．用 $y^*$ 表示延拓后的泛函，则 $y^* \in \mathcal{Y}^*$ ，且满足
$$
\left\langle y^*, T x\right\rangle=x^*(x), \quad \forall x \in \mathcal{X} .
$$
由共轭算子的定义，有
$$
\left\langle y^*, T x\right\rangle=\left\langle T^* y^*, x\right\rangle, \quad \forall x \in \mathcal{X} .
$$
联合上式即得
$$
\left\langle T^* y^*, x\right\rangle=x^*(x), \quad \forall x \in \mathcal{X} .
$$
故 $x^*=T^* y^*$ ，也就是 $x^* \in \mathcal{R}\left(T^*\right)$ ．
\end{proof}

\end{problemset}



\section{Fredholm算子理论}

\subsection{Fredholm二择一性质}

我们知道,在有限维空间 \({\mathbb{R}}^{n}\) 的情形,每个有界线性算子都可以用一个矩阵 \(T\) 来表示, 在不同的正交基下它们彼此酉相似, 因此算子的谱可用线性代数方程组
\begin{equation*}
\left( {{\lambda I} - T}\right) x = y \label{4.15}
\end{equation*}

的可解性来刻画: 或者对任意 \(y \in  {\mathbb{R}}^{n}\) 上述方程组有唯一解,此时 \(\lambda  \in  \rho \left( T\right)\) ; 或者当 \(y = 0\) 时方程组有非零解存在,此时 \(\lambda  \in  \sigma \left( T\right)  = {\sigma }_{p}\left( T\right)\) . 通常称这个结果为Fredholm二择一定理.
\begin{definition}[Fredholm择一性]
    在赋范线性空间 \(\mathcal{X}\) 上的有界线性算子 \(T : \mathcal{X} \rightarrow\)  \(\mathcal{X}\) 称为满足Fredholm择一性,如果 \(T\) 满足1或者满足2 .
\begin{enumerate}
    \item 非齐次方程
\begin{equation*}
{Tx} = y,\;{T}^{ * }f = g
\end{equation*}

对任意给定的 \(y \in  \mathcal{X}\) 和 \(g \in  {\mathcal{X}}^{ * }\) 均各有解 \(x\) 和 \(f\) ,且解唯一;

对应的齐次方程
\begin{equation*}
{Tx} = 0,\;{T}^{ * }f = 0
\end{equation*}

各只有平凡解 \(x = 0\) 和 \(f = 0\) .
\item 齐次方程
\begin{equation*}
{Tx} = 0,\;{T}^{ * }f = 0
\end{equation*}

各有相同数目的线性无关解
\begin{equation*}
{x}_{1},\cdots ,{x}_{n}\text{ 与 }{f}_{1},\cdots ,{f}_{n},\;\left( {n \geq  1}\right) \text{ . }
\end{equation*}

非齐次方程
\begin{equation*}
{Tx} = y,\;{T}^{ * }f = g
\end{equation*}

各对任意的 \(y\) 与 \(g\) 并不总有解,其各有解的充分必要条件是 \(y\) 与 \(g\) 各满足
\begin{equation*}
{f}_{k}\left( y\right)  = 0,\;g\left( {x}_{k}\right)  = 0,\;k = 1,2,\cdots ,n.
\end{equation*}

\end{enumerate}



\end{definition}在无限维空间中对紧算子我们有相应的结果.
\begin{theorem}
    设 \(\mathcal{X}\) 为Banach空间, \(T \in  \mathcal{C}\left( \mathcal{X}\right)\) . 设 \(\lambda  \neq  0\) ,那末或者对任意 \(y \in  \mathcal{X}\) ,非齐次方程
\begin{equation*}
\left( {{\lambda I} - T}\right) x = y
\end{equation*}

有唯一解,此时, \(\lambda  \in  \rho \left( T\right)\) ; 或者齐次方程有非零解,此时 \(\lambda  \in  {\sigma }_{p}\left( T\right)\) ,解空间的维数有限. 相应的非齐次方程有解当且仅当对任意 \(f \in  N\left( {\bar{\lambda }I - {T}^{ * }}\right)\) 有 \(f\left( y\right)  = 0\) .
\end{theorem}
\begin{example}
    设 \(K \in  C\left( {\left\lbrack  {0,1}\right\rbrack   \times  \left\lbrack  {0,1}\right\rbrack  }\right)\) . 考察方程
\begin{equation*}
x\left( t\right)  = {\int }_{0}^{1}K\left( {s,t}\right) x\left( s\right) {\dif s} + y\left( t\right)  \label{4.16}
\end{equation*}

及其共轭方程
\begin{equation*}
f\left( t\right)  = {\int }_{0}^{1}K\left( {t,s}\right) f\left( s\right) {\dif s} + g\left( t\right)  \label{4.17}
\end{equation*}

其中 \(x,y,f,g \in  {L}^{2}\left( \left\lbrack  {0,1}\right\rbrack  \right)\) ,则关于方程 (4.16) 只有两种情形:

结论1: (1) 或者对任意 \(y \in  {L}^{2}\left( \left\lbrack  {0,1}\right\rbrack  \right) ,\left( {4.16}\right)\) 存在唯一解 \(x \in  {L}^{2}\left( \left\lbrack  {0,1}\right\rbrack  \right)\) ;

( 2 )或者当 \(y = 0\) 时,( 4.16 )有非零解.

结论2: 方程 (4.17) 与方程 (4.16) 的情形一样.
\end{example}.


\subsection{Fredholm算子}

设 \(\mathcal{X},\mathcal{Y}\) 为有限维实向量空间, \(T : \mathcal{X} \rightarrow  \mathcal{Y}\) 为线性算子,则 \(T\) 的值域 \(R\left( T\right)\) 的维数等于其核 \(\mathrm{{Ker}}T\) 的余维,即
\begin{equation*}
\dim \mathcal{X} - \dim \operatorname{Ker}T = \dim \mathcal{Y} - \dim \operatorname{Coker}T,
\end{equation*}

其中 \(\operatorname{Coker}T = \mathcal{Y}/T\left( \mathcal{X}\right)\) . 这意味着
\begin{equation*}
\dim \operatorname{Ker}T - \dim \operatorname{Coker}T = \dim \mathcal{X} - \dim \mathcal{Y}
\end{equation*}

不依赖算子 \(T\) . 称上式左边为 \(T\) 的指标. 我们自然会问: 在无限维情形线性算子的指标是否仍是稳定的?

若 \(\mathcal{X},\mathcal{Y}\) 都是Banach空间且 \(T \in  L\left( {\mathcal{X},\mathcal{Y}}\right)\) ,则
\begin{equation*}
\operatorname{Ker}T = \{ x \in  \mathcal{X} \mid  {Tx} = 0\}
\end{equation*}

为 \(\mathcal{X}\) 的闭子空间,但未必是有限维的. 值域 \(T\left( \mathcal{X}\right)\) 也不一定是闭的,但我们有
\begin{lemma}
   若 \(T \in  \mathcal{L}\left( {\mathcal{X},\mathcal{Y}}\right)\) 且值域 \(T\left( \mathcal{X}\right)\) 在 \(\mathcal{Y}\) 中具有有限的余维数,则 \(T\left( \mathcal{X}\right)\) 为闭子空间.  \
\end{lemma}
\begin{proof}
    不妨假设 \(T\) 为单射. 否则,考虑由 \(T\) 诱导的从 \(\mathcal{X}/\operatorname{Ker}T\) 到 \(\mathcal{Y}\) 的映射即可. 若 \(T\left( \mathcal{X}\right)\) 的余维为 \(n\) ,则存在线性映射
\begin{equation*}
S : {\mathbb{C}}^{n} \rightarrow  \mathcal{Y}
\end{equation*}

使得 \(S\left( {\mathbb{C}}^{n}\right)\) 为 \(T\left( \mathcal{X}\right)\) 在 \(\mathcal{Y}\) 中的补,即映射
\begin{equation*}
\begin{matrix} {T}_{1} : \mathcal{X} \oplus  {\mathbb{C}}^{n} &  \rightarrow  & \mathcal{Y} \\  \left( {x,y}\right) &  \mapsto  & {Tx} + {Sy} \end{matrix}
\end{equation*}

为双射. 由Banach逆算子定理知, \({T}_{1}^{-1} \in  \mathcal{L}\left( {\mathcal{Y},\mathcal{X} \oplus  {\mathbb{C}}^{n}}\right)\) ,即 \({T}_{1}\) 为同胚映射. 故 \(T\left( \mathcal{X}\right)  = {T}_{1}\left( {\mathcal{X}\oplus \{ 0\} }\right)\) 为闭的.
\end{proof}
\begin{note}
同胚映射将闭集映成闭集.

    利用 \(T\) 的图像可以将该结果推广到 \(T\) 为闭算子但不必连续的情形.
\end{note}



如果我们要考虑一个算子的指标, 我们往往需要判断其何时具有闭值域和有限维的核空间. 这里我们介绍几个引理.
\begin{lemma}
     设 \(\mathcal{X},\mathcal{Y},\mathcal{Z}\) 为Banach空间, \(T \in  \mathcal{L}\left( {\mathcal{X},\mathcal{Y}}\right)\) 且 \(K \in  \mathcal{C}\left( {\mathcal{X},\mathcal{Z}}\right)\) . 若存在常数 \(c > 0\) 使得
\begin{equation*}
\parallel x{\parallel }_{\mathcal{X}} \leqslant  c\left( {\parallel {Tx}{\parallel }_{\mathcal{Y}} + \parallel {Kx}{\parallel }_{\mathcal{Z}}}\right) ,x \in  \mathcal{X}. \label{4.18}
\end{equation*}

则 \(T\) 具有闭值域和有限维的核空间.
\end{lemma}
\begin{proof}
     为证 \(\dim\Ker T <  + \infty\) ,我们只需证\(\Ker T\) 中单位球为紧集即可. 为此,选择点列 \({x}_{\nu } \in  \mathcal{X}\) 满足
\begin{equation*}
\begin{Vmatrix}{x}_{\nu }\end{Vmatrix} \leqslant  1,\;T{x}_{\nu } = 0.
\end{equation*}

由于 \({x}_{\nu }\) 有界及 \(K\) 为紧算子,故存在子列,仍记为 \({x}_{\nu }\) ,使得 \(K{x}_{\nu }\) 收敛. 由 \({x}_{\nu } \in\)  \(\operatorname{Ker}T\) 知 \({x}_{\nu }\) 为Cauchy列. 由 \(\mathcal{X}\) 的完备性知, \({x}_{\nu }\) 收敛.

下证 \(T\) 的值域 \(T\left( \mathcal{X}\right)\) 为闭的.  $N(T)$是$\mathcal{X}$的闭子空间,那么定义映射\[
    \widetilde{T}:\mathcal{X}/N(T)\to R(T),[x]\to T x
    \]
    是既单又满的映射,考虑$\widetilde{T}^{-1}$,如果它是连续的,那么$R(T)$是闭的,假设它不是连续的,那么存在$[x_n]\in \mathcal{X}/N(T),[x_n]\not\to 0,T x_n\to 0$,那么存在$\varepsilon_0>0$以及$[y_m]=[x_{n_{m}}]$
使得\[
\|[y_m]\|\geq \varepsilon_0,\textbf{取}[z_m]=\frac{[y_m]}{\|[y_m]\|},\widetilde{T}[z_m]\to 0
\]
又存在$w_m\in[z_m]$,满足\[
\|w_m\|\leqslant2 \|[z_m]\|\leqslant 2
\]
同时有\[
\widetilde{T}[w_m]=\widetilde{T}[z_m]\to 0
\]
由于$K$是紧的,那么存在$Kw_{m_k}$是柯西列,由条件知$w_{m_k}$也是柯西列,设其收敛到$w$,那么
\[T w=\lim\limits_{k\to \infty}T w_{m_k}=\lim\limits_{k\to \infty}\widetilde{T}[w_{m_k}]=0\]
得到$z\in N(T)$,那么\[
\|[ w_{m_k}]\|=\|[ w_{m_k}-z]\|\leqslant \|w_{m_k}-z\|\to 0
\]
与\[\|[w_{m_k}]\|=1\]
矛盾,所以是闭的.

\end{proof}

\begin{corollary}
    设 \(\mathcal{X},\mathcal{Y}\) 为Banach空间, \(T \in  \mathcal{L}\left( {\mathcal{X},\mathcal{Y}}\right)\) 具有闭值域和有限维核空间. 则

(i)对任意紧算子 \(K \in  \mathcal{C}\left( {\mathcal{X},\mathcal{Y}}\right)\) ,算子 \(T + K\) 也具有闭值域及有限维核空间.

(ii) 存在 \(\varepsilon  > 0\) 使得若 \(P \in  \mathcal{L}\left( {\mathcal{X},\mathcal{Y}}\right)\) 满足 \(\parallel P\parallel  < \varepsilon\) ,则 \(T + P\) 具有闭值域及有限维核空间.

\end{corollary}
\begin{proof}
记 \(\dim\Ker T = n\) . 选择有界线性算子 \({K}_{0} : \mathcal{X} \rightarrow  {\mathbb{R}}^{n}\) 使得 \({K}_{0}\) 到 \(\Ker
T\) 上的限制为向量空间同构. 则算子
\begin{equation*}
\begin{aligned} \left( {T,{K}_{0}}\right)  : \mathcal{X} &  \rightarrow  T\mathcal{X} \oplus  {\mathbb{R}}^{n} \\  x &  \mapsto  \left( {{Tx},{K}_{0}x}\right)  \end{aligned}
\end{equation*}

既为单射又为满射. 因此,由Banach逆算子定理知, \(\left( {T,{K}_{0}}\right)\) 有有界逆. 从而存在常数 \(c > 0\) 使得
\begin{equation*}
\parallel x{\parallel }_{\mathcal{X}} \leqslant  c\left( {\parallel {Tx}{\parallel }_{\mathcal{Y}} + {\begin{Vmatrix}{K}_{0}x\end{Vmatrix}}_{{\mathbb{R}}^{n}}}\right) .
\end{equation*}

因此,对任意紧算子 \(K \in  \mathcal{C}\left( {\mathcal{X},\mathcal{Y}}\right)\) 有
\begin{equation*}
\begin{aligned}
\widetilde{K}:\mathcal{ X}\to \mathcal{Y}\oplus \mathbb{R}^n,x\mapsto (Kx,K_0x)\textbf{ 为紧算子}\\
    \parallel x{\parallel }_{\mathcal{X}} \leqslant  c\left( {\parallel \left( {T + K}\right) x{\parallel }_{\mathcal{Y}} + \parallel {Kx}{\parallel }_{\mathcal{Y}} + {\begin{Vmatrix}{K}_{0}x\end{Vmatrix}}_{{\mathbb{R}}^{n}}}\right) .
\end{aligned}
\end{equation*}

同理,若 \(\parallel P\parallel  < \frac{1}{c}\) ,则
\begin{equation*}
\parallel x{\parallel }_{\mathcal{X}} \leqslant  \frac{c}{1 - c\parallel P\parallel }\left( {\parallel \left( {T + P}\right) x{\parallel }_{\mathcal{Y}} + {\begin{Vmatrix}{K}_{0}x\end{Vmatrix}}_{{\mathbb{R}}^{n}}}\right) .
\end{equation*}

因此, 由引理4.10立即可得推论.
\end{proof}
\begin{corollary}
    设 \(\mathcal{X},\mathcal{Y}\) 为Banach空间, \(T \in  \mathcal{L}\left( {\mathcal{X},\mathcal{Y}}\right)\) 具有闭值域和有限维上核空间 \(\Coker T\) . 则

(i)对任意紧算子 \(K \in  \mathcal{C}\left( {\mathcal{X},\mathcal{Y}}\right)\) ,算子 \(T + K\) 也具有闭值域和有限维上核空间 \(\Coker\left( {T + K}\right)\) .

(ii) 存在 \(\varepsilon  > 0\) 使得若 \(P \in  \mathcal{L}\left( {\mathcal{X},\mathcal{Y}}\right)\) 满足 \(\parallel P\parallel  < \varepsilon\) ,则 \(T + P\) 具有闭值域和有限维上核空间Coker \(\left( {T + P}\right)\) .
\end{corollary}
\begin{proof}
    由于算子 \(T\) 有闭值域当且仅当 \({T}^{ * }\) 有闭值域,而且 \(T\) 有有限维上核空间当且仅当 \({T}^{ * }\) 有有限维核空间. 因此,本推论的结果可直接由推论得出.
\end{proof}
\begin{proposition}\label{pro4.2}
         若 \(T \in  \mathcal{L}\left( {\mathcal{X},\mathcal{Y}}\right)\) ,则下列条件等价:

(i) \(\dim \operatorname{Ker}T < \infty\) 且 \(T\left( \mathcal{X}\right)\) 为闭的.

(ii) \(\mathcal{X}\) 中任意使得 \(T{x}_{j}\) 收敛的有界点列 \({x}_{j}\) 都有收敛子列.
\end{proposition}

\begin{proof}
    首先,若(ii)成立,则有 \(\operatorname{Ker}T\) 中单位球为紧集,从而由Riesz引理知, \(\operatorname{Ker}T\) 为有限维. 因此,下面我们假设 \(\operatorname{Ker}T\) 为有限维.

由Hahn-Banach定理知,有限维空间 \(\operatorname{Ker}T\) 具有闭补空间 \({\mathcal{X}}_{0} \subset  \mathcal{X}\) ,即 \(\mathcal{X} =\) Ker \(T \oplus  {\mathcal{X}}_{0}\) . 我们断言有结论 (ii) 等价于
\begin{equation}
\text{ 存在常数 }C > 0\text{ 使得 }\parallel x{\parallel }_{\mathcal{X}} \leqslant  C\parallel {Tx}{\parallel }_{\mathcal{Y}},x \in  {\mathcal{X}}_{0}\text{ . } \label{4.19}
\end{equation}

(ii) \(\Rightarrow\) \ref{4.19} : 用反证法,若不然,则存在点列 \({x}_{j} \in  {X}_{0}\) 使得 \({\begin{Vmatrix}{x}_{j}\end{Vmatrix}}_{\mathcal{X}} = 1\) 且 \({\begin{Vmatrix}T{x}_{j}\end{Vmatrix}}_{\mathcal{Y}} \leqslant  \frac{1}{j}\) . 因此由结论 (ii) 知, \({x}_{j}\) 有收敛子列,记其极限为 \(x\) ,则有 \(x \in\)  \({\mathcal{X}}_{0},\parallel x{\parallel }_{\mathcal{X}} = 1,{Tx} = 0\) . 这与 \({\mathcal{X}}_{0}\) 的定义相矛盾,故 (4.19) 成立.

\ref{4.19} \( \Rightarrow  \) (ii): 设\ref{4.19}成立. 若 \({x}_{j}\) 是 \(X\) 中使得 \(T{x}_{j}\) 收敛的有界点列,则有 \({x}_{j} = {g}_{j} + {h}_{j}\) ,由习题4.1.12 \({g}_{j} \in  \operatorname{Ker}T\) 和 \({h}_{j} \in  {\mathcal{X}}_{0}\) 均有界. 由于 \(T{h}_{j} = T{x}_{j}\) ,故由\ref{4.19}可得, \({h}_{j}\) 收敛. 又由 \(\Ker T\) 为有限维知,有界点列 \({g}_{j}\) 有收敛子列. 因此,(ii) 成立.

最后,由引理4.10知,\ref{4.19}蕴含(i)成立. 反之,由(i)知, \(T{ \mid  }_{{\mathcal{X}}_{0}} : {\mathcal{X}}_{0} \rightarrow\)  \(T{\mathcal{X}}_{0}\) 为双射. 从而由Banach逆算子定理知,\ref{4.19}成立. 故定理得证.
\end{proof}
\begin{definition}
    设 \(\mathcal{X},\mathcal{Y}\) 为Banach空间. \(T \in  \mathcal{L}\left( {\mathcal{X},\mathcal{Y}}\right)\) 称为Fredholm 算子如果 \(\dim\mathrm{{Ker}}T\) 有限且 \(T\left( \mathcal{X}\right)\) 为闭并有有限余维. 对Fredholm算子 \(T\) ,我们定义其指标为
\begin{equation*}
\text{ ind }T = \dim \operatorname{Ker}T - \dim \operatorname{Coker}T,\Coker T= \mathcal{Y}/\mathcal{R}(T).
\end{equation*}

\end{definition}
\begin{note}
    在仅有 \(\dim\operatorname{Ker}T\) 有限且 \(T\left( \mathcal{X}\right)\) 闭时,若容许指标为 \(- \infty\) ,有时也沿用上述定义.
\end{note}

关于Fredholm算子指标的稳定性质, 我们有
\begin{theorem}
设 $\mathcal{X},\mathcal{Y}$ 为 Banach 空间，$T\in \mathcal{L}(\mathcal{X},\mathcal{Y})$ 满足
\[
\dim \ker T < \infty, \qquad R(T) \text{ 闭}.
\]
则存在 $\varepsilon>0$，使得对任意 $S\in \mathcal{L}(\mathcal{X},\mathcal{Y})$ 且 $\|S\|<\varepsilon$，有：
\begin{enumerate}[(1)]
\item $\dim \ker (T+S) \leq \dim \ker T$；
\item $T+S$ 具有闭值域；
\item $\operatorname{ind}(T+S)=\operatorname{ind}T$。
\end{enumerate}
\end{theorem}
\begin{proof}
首先，假设 $T$ 可逆。则有
\[
T + S = T \bigl(I + T^{-1} S \bigr).
\]
如果 $\|T^{-1}\| \cdot \|S\| < 1$，则 $I + T^{-1}S$ 可由 Neumann 级数求逆，因此 $T + S$ 可逆，从而 $\operatorname{Ker}(T+S) = \{0\}$，值域闭合，指标不变：
\[
\operatorname{ind}(T + S) = \operatorname{ind}T = 0.
\]

其次，假设 $T$ 仅是单射。由条件
\[
\|f\|_1 \le C \|Tf\|_2 \le C \|(T+S)f\|_2 + C \|S\| \|f\|_1.
\]
若 $C\|S\| < 1/2$，则
\[
\|f\|_1 \le 2C \|(T+S)f\|_2,
\]
因此 $T + S$ 是单射且值域闭合。

若 $\dim \operatorname{Coker}T < \infty$，选取 $T(B_1)$ 的有限维补空间 $W \subset B_2$，并令自然投影为
\[
q: B_2 \to B_2 / W.
\]
则 $qT: B_1 \to B_2 / W$ 可逆。当 $\|S\|$ 足够小时，$qT + qS$ 也可逆。因此：
\[
\operatorname{ind}(T + S) = -\dim W = \operatorname{ind}T.
\]

如果只知道 $\dim \operatorname{Coker}T > \nu$，可以选择 $\dim W = \nu + 1$ 且 $W \cap T(B_1) = \{0\}$，则当 $\|S\|$ 足够小时有
\[
\operatorname{ind}(T + S) < -\nu.
\]
将 $T$ 替换为 $T + S$ 可得集合
\[
\{S : C\|S\| < 1/2, \operatorname{ind}(T+S) = -\nu\}, \quad
\{S : C\|S\| < 1/2, \operatorname{ind}(T+S) < -\nu\}
\]
都是开集。由于 $\{S : C\|S\| < 1/2\}$ 连通，故 $\operatorname{ind}(T + S)$ 必为常数。

最后，若 $N = \operatorname{Ker}T \neq \{0\}$，选取 $B_0 \subset B_1$ 为 $N$ 的拓扑补空间，则 $B_1 = B_0 \oplus N$。由上面结果，$T + S$ 限制在 $B_0$ 上是单射且值域闭合，并且
\[
\operatorname{codim}((T + S)B_0) = \operatorname{codim} T B_0 = \operatorname{codim} T B_1
\]
当 $\|S\|$ 足够小时成立。

设 $N' = \operatorname{Ker}(T + S)$，则 $N' \cap B_0 = \{0\}$，并存在有限维 $W$ 使得
\[
B_1 = N' \oplus W \oplus B_0.
\]
因此：
\[
\operatorname{codim}((T+S)B_1) = \operatorname{codim}((T+S)W \oplus (T+S)B_0) = \operatorname{codim} T B_1 - \dim W,
\]
于是：
\[
\operatorname{ind}(T + S) = \dim N' + \dim W - \operatorname{codim} T B_1 = \dim N - \dim \operatorname{Coker} T = \operatorname{ind}T.
\]
证毕。
\end{proof}

\begin{corollary}
    记\(\operatorname{Fred}\left( {\mathcal{X},\mathcal{Y}}\right)  = \{ T \in  \mathcal{L}\left( {\mathcal{X},\mathcal{Y}}\right)  \mid  T\) 为Fredholm 算子 \(\}\) ,则集合 \(\operatorname{Fred}\left( {\mathcal{X},\mathcal{Y}}\right)\) 为 \(\mathcal{L}\left( {\mathcal{X},\mathcal{Y}}\right)\) 中的开集且ind \(T\) 在每个连通分支上取常值.
\end{corollary}

\begin{corollary}
    若 \({T}_{1} \in  \operatorname{Fred}\left( {\mathcal{X},\mathcal{Y}}\right) ,{T}_{2} \in  \operatorname{Fred}\left( {\mathcal{Y},\mathcal{Z}}\right)\) ,则 \({T}_{2}{T}_{1} \in  \operatorname{Fred}\left( {\mathcal{X},\mathcal{Z}}\right)\) 且
\begin{equation*}
\operatorname{ind}\left( {{T}_{2}{T}_{1}}\right)  = \operatorname{ind}{T}_{1} + \operatorname{ind}{T}_{2}.
\end{equation*}

\end{corollary}
\begin{proof}
    由于 \({T}_{1}\) 映 \(\operatorname{Ker}{T}_{2}{T}_{1}\) 到 \(\operatorname{Ker}{T}_{2}\) 且具有核 \(\operatorname{Ker}{T}_{1}\) ,因此,
\begin{equation*}
\dim \operatorname{Ker}{T}_{2}{T}_{1} \leqslant  \dim \operatorname{Ker}{T}_{1} + \dim \operatorname{Ker}{T}_{2}.
\end{equation*}

同理可得,
\begin{equation*}
\dim \operatorname{Coker}{T}_{2}{T}_{1} \leqslant  \dim \operatorname{Coker}{T}_{1} + \dim \operatorname{Coker}{T}_{2}.
\end{equation*}

从而 \({T}_{2}{T}_{1} \in  \operatorname{Fred}\left( {\mathcal{X},\mathcal{Z}}\right)\) . 若用 \({I}_{2}\) 记 \(\mathcal{Y}\) 上的恒等算子,用分块矩阵的记号则有对任意 \(t\) ,
\begin{equation*}
{F}_{t} \mathrel{\text{ := }} \left( \begin{matrix} {I}_{2} & 0 \\  0 & {T}_{2} \end{matrix}\right) \left( \begin{matrix} {I}_{2}\cos t & {I}_{2}\sin t \\   - {I}_{2}\sin t & {I}_{2}\cos t \end{matrix}\right) \left( \begin{matrix} {T}_{1} & 0 \\  0 & {I}_{2} \end{matrix}\right)
\end{equation*}

为从 \(\mathcal{X} \oplus  \mathcal{Y}\) 到 \(\mathcal{Y} \oplus  \mathcal{Z}\) 的Fredholm 算子. 由于 \({F}_{0} = {T}_{1} \oplus  {T}_{2}\) 且 \({F}_{-\frac{\pi }{2}} : \left( {{x}_{1},{x}_{2}}\right)  \mapsto\)  \(\left( {-{x}_{2},{T}_{2}{T}_{1}{x}_{1}}\right)\) ,从而有 \(\operatorname{ind}{F}_{0} = \operatorname{ind}{T}_{1} + \operatorname{ind}{T}_{2}\) 及 \(\operatorname{ind}{F}_{-\frac{\pi }{2}} = \operatorname{ind}\left( {{T}_{2}{T}_{1}}\right)\) . 故
\begin{equation*}
\operatorname{ind}\left( {{T}_{2}{T}_{1}}\right)  = \operatorname{ind}{T}_{1} + \operatorname{ind}{T}_{2}.
\end{equation*}

\end{proof}
\begin{corollary}
     设 \(T \in  \operatorname{Fred}\left( {\mathcal{X},\mathcal{Y}}\right) ,K \in  \mathcal{C}\left( {\mathcal{X},\mathcal{Y}}\right)\) ,则 \(T + K \in  \operatorname{Fred}\left( {\mathcal{X},\mathcal{Y}}\right)\) 且 \(\operatorname{ind}\left( {T + K}\right)  = \operatorname{ind}T\) .
\end{corollary}
\begin{proof}
    任取 \(t \in  \mathbb{R}\) . 若 \({x}_{j}\) 为 \(\mathcal{X}\) 中使得 \(\left( {T + {tK}}\right) {x}_{j}\) 收敛的有界点列,则存在子列 \({x}_{{j}_{k}}\) 使得 \(K{x}_{{j}_{k}}\) 收敛. 从而 \(T{x}_{{j}_{k}}\) 也收敛. 由 \(T\) 为Fredholm算子及命题4.2 知 \({x}_{{j}_{k}}\) 有收敛子列. 再次对算子 \(T + {tK}\) 应用命题4.2可得 \(T + {tK}\) 具有有限维核及闭值域. 由定理4.10知 \(\operatorname{ind}\left( {T + {tK}}\right)\) 为 \(t \in  \mathbb{R}\) 的局部常值函数,从而不依赖 \(t\) . 因此, \(T + {tK}\) 总是Fredholm算子且 \(\operatorname{ind}\left( {T + {tK}}\right)  = \operatorname{ind}T\) . 特别地,取 \(t = 1\) 则得推论结论.
\end{proof}

\begin{corollary}
    设 \(T \in  \mathcal{L}\left( {\mathcal{X},\mathcal{Y}}\right)\) 且存在 \({S}_{j} \in  \mathcal{L}\left( {Y,\mathcal{X}}\right) ,i = 1,2\) 使得 \(T{S}_{2} = {I}_{2} +\)  \({K}_{2},{S}_{1}T = {I}_{1} + {K}_{1}\) ,其中 \({K}_{j},j = 1,2\) 为紧算子. 则 \(T,{S}_{1}\) 及 \({S}_{2}\) 都是Fredholm算子且 \(\operatorname{ind}T =  - \operatorname{ind}{S}_{j},j = 1,2\) .
\end{corollary}

\begin{proof}
    由于
\begin{equation*}
\dim \operatorname{Ker}T \leqslant  \dim \operatorname{Ker}\left( {{I}_{1} + {K}_{1}}\right) ,
\end{equation*}
\begin{equation*}
\dim \operatorname{Coker}T \leqslant  \dim \operatorname{Coker}\left( {{I}_{2} + {K}_{2}}\right) ,
\end{equation*}

故由引理4.6知, \(T\) 为Fredholm 算子. 由于
\begin{equation*}
{S}_{1}T{S}_{2} = {S}_{1} + {S}_{1}{K}_{2} = {S}_{2} + {K}_{1}{S}_{2},
\end{equation*}

故 \({S}_{2} - {S}_{1}\) 为紧算子. 因此 \({S}_{2}T - {I}_{1}\) 及 \(T{S}_{1} - {I}_{2}\) 也是紧算子. 由前面的证明知, \({S}_{1}\) 及 \({S}_{2}\) 也是Fredholm算子且由推论4.3及推论4.5知, ind \(T +\) ind \({S}_{j} = \operatorname{ind}\left( {{I}_{j} + {K}_{j}}\right)  =\) 0. 至此, 推论得证.
\end{proof}\begin{theorem}
设 $V$ 是一个在域$\mathbb{K}$上的有限维线性空间,$\dim V = n$. 它的对偶空间$V^* := \{ f : V \to \mathbb{K} \mid f \text{线性} \}$ 也有有限维数 $n$, 即
\[
\dim V^* = \dim V = n.
\]
\end{theorem}

\begin{proof}
取 $\{ e_1, \dots, e_n \}$ 是 $V$的一组基,定义对偶基$\{ e_1^*, \dots, e_n^* \} \subset V^*$ by
\[
e_i^*(e_j) := \delta_{ij} = 
\begin{cases} 
1, & i=j \\ 
0, & i \neq j 
\end{cases}.
\]
 对于任意 $f \in V^*$, 有
\[
f = \sum_{i=1}^n f(e_i) e_i^*.
\]
实际上, 对于任意 $v = \sum_{i=1}^n v_i e_i \in V$,
\[
\sum_{i=1}^n f(e_i) e_i^*(v) = \sum_{i=1}^n f(e_i) v_i = f\Big(\sum_{i=1}^n v_i e_i\Big) = f(v).
\]
因此$\operatorname{span}\{ e_1^*, \dots, e_n^* \}=V^*$.

假设$\sum_{i=1}^n \alpha_i e_i^* = 0$. 那么对于任意 $v \in V$,
\[
\sum_{i=1}^n \alpha_i e_i^*(v) = 0.
\]
取 $v = e_j$ 得到 $\alpha_j = 0$  $j=1,\dots,n$. 因此 $\{ e_1^*, \dots, e_n^* \}$ 是线性无关的.

因此 $\{ e_1^*, \dots, e_n^* \}$ 是 $V^*$的基,那么
\[
\dim V^* = n = \dim V.
\]
\end{proof}

\begin{theorem}[Rudin, Theorem 4.12]
设 $\mathcal{X}$ 是Banach空间 $M \subset \mathcal{X}$ 是闭子集. 那么
\[
(\mathcal{X} / M)^* \cong {}^\perp M.
\]
特别地,如果$\dim(\mathcal{X}/M) < \infty$, 那么
\[
\dim {}^\perp M = \dim(\mathcal{X}/M) = \operatorname{codim} M.
\]
\end{theorem}

\begin{proof}
考虑商映射
\[
\pi: \mathcal{X} \to \mathcal{X}/M, \quad \pi(x) = x + M.
\]
定义对偶映射
\[
\pi^* : (\mathcal{X}/M)^* \to \mathcal{X}^*, \quad \pi^*(f) = f \circ \pi.
\]
任意 $f \in (\mathcal{X}/M)^*$, 显然 $\pi^*(f)(m) = f(\pi(m)) = f(0) = 0$ $\forall m \in M$, 所以$\pi^*(f) \in{}^\perp M$.

设 $x^* \in M^\perp \subset \mathcal{X}^*$. 定义 $f : \mathcal{X}/M \to \mathbb{K}$ 
\[
f(x+M) := x^*(x).
\]
这是良定义的,因为$x+M = y+M$, 那么 $x-y \in M$, 所以 $x^*(x-y) = 0$. 此外, $f$ 是线性有界且 $\pi^*(f) = x^*$.

如果 $\pi^*(f) = f \circ \pi = 0$, 那么 $f(x+M) = 0$ $,\forall x \in \mathcal{\mathcal{X}}$, 因此 $f = 0$ in $(\mathcal{X}/M)^*$. 因此 $\pi^*$是单射.

所以
\[
(\mathcal{X}/M)^* \cong {}^\perp M.
\]
特别地, 如果 $\dim(\mathcal{X}/M) < \infty$, 那么
\[
\dim{}^\perp M =\dim (\mathcal{X}/M)^*= \dim(\mathcal{X}/M) = \operatorname{codim} M.
\]

\end{proof}


\subsection{习 题}

\begin{problemset}

  \item 设 \({V}_{1},{V}_{2}\) 为向量空间（维数不必有限）.若线性算子 \(T : {V}_{1} \rightarrow  {V}_{2}\) 满足 \(\dim \operatorname{Ker}T\) 或 \(\dim \operatorname{Coker}T\) 有限，则定义 \(T\) 的指标为
\[
  \operatorname{ind}T = \dim \operatorname{Ker}T - \dim \operatorname{Coker}T.
  \]
  证明：
  \begin{enumerate}
    \item 若 \(T : {V}_{1} \rightarrow  {V}_{2}\) 的指标有定义，且 \(S \in  F\left( {{V}_{1},{V}_{2}}\right)\)，则 \(T + S\) 的指标也可定义且 \(\operatorname{ind}T = \operatorname{ind}\left( {T + S}\right)\).
    \item 若 \({T}_{1} : {V}_{1} \rightarrow  {V}_{2}\) 和 \({T}_{2} : {V}_{2} \rightarrow  {V}_{3}\) 的指标都有定义，则
\[
    \operatorname{ind}\left( {{T}_{2}{T}_{1}}\right)  = \operatorname{ind}{T}_{2} + \operatorname{ind}{T}_{1}.
    \]
  \end{enumerate}
\begin{proof}
    
\end{proof}
  \item 设 \(\mathcal{X},\mathcal{Y}\) 为 Banach 空间.证明有界线性算子 \(T : \mathcal{X} \rightarrow  \mathcal{Y}\) 为 Fredholm 算子当且仅当存在有界线性算子 \(S : \mathcal{Y} \rightarrow  \mathcal{X}\)，使得 \({TS} - {I}_{\mathcal{Y}}\) 及 \({ST} - {I}_{\mathcal{X}}\) 都是紧算子.
\begin{proof}
  令  $$
\widetilde{T}[x]=T x \quad(\forall x \in[x]) \quad(\forall[x] \in \mathcal{X} / \mathcal{N}(T)) .
$$

由假设 $T \in \mathcal{F}(\mathcal{X}, \mathcal{Y})$ ，从而 $\widetilde{T}: \mathcal{X} / \mathcal{N}(T) \rightarrow \mathcal{R}(T)$ 有连续逆 $\widetilde{T}^{-1}$ ，并且存在投影算子（参看习题4.1）：
$$
\begin{aligned}
& A_1: \mathcal{X} \rightarrow \mathcal{N}(T), \\
& A_2: \mathcal{Y} \rightarrow \mathcal{Y} / \mathcal{R}(T) .
\end{aligned}
$$

显然 $A_1$ 和 $A_2$ 都是有穷秩的，从而是紧的.令
$$
S \triangleq f^{-1} \widetilde{T}^{-1}\left(I_y-A_2\right)
$$

由于$I_y-A_2$为值域为$\mathcal{R}(T)$的幂等算子,便有
$$
\left(I_y-A_2\right) T x=Tx= \widetilde{T}f(I_x-A_1) x=\widetilde{T}([x])=Tx\Rightarrow ST=I_x-A_1
$$


以及
$$
T S=Tf^{-1} \widetilde{T}^{-1}\left(I_y-A_2\right)=I_y-A_2 .
$$
\begin{center}
\begin{tikzcd}[column sep=6em]
  && &\mathcal{N}(T)&\\
  &&&\bigoplus& \\
  \mathcal{X}\arrow[rrr,"I_x-A_1" ] 
      \arrow[rrruu,"A_1"]      
      \arrow[rrrd, "T"]                        
  & & & R(I_x-A_1)\arrow[r,"f"]&\mathcal{X}/\mathcal{N}(T)\arrow[l,"f^{-1}"] \arrow[ld, "\widetilde{T}"]\\[2.5em]
  \mathcal{Y} \arrow[rrr,"I_y-A_2" ]    
      \arrow[rrrdd,"A_2"]         
      \arrow[rrru, "S", dashed]               
  & & & {\mathcal{R}(T)} \arrow[ru, "\widetilde{T}^{-1}"]& \\
  &&&\bigoplus&\\
  & & & R(A_2)   \arrow[r,"g"]               &{\mathcal{Y}/\mathcal{R}(T)}  \arrow[l,"g^{-1}"]                      
\end{tikzcd}
\end{center}
如果存在 $S$ 及 $A_1, A_2$ ，使得上式成立，那么由推论5.2.2知
$$
\begin{aligned}
\mathcal{N}(T) & \subset N\left(S T\right)=N\left(I_x-A_1\right) 
\Longrightarrow  \operatorname{dim} \mathcal{N}(T)<\infty
\end{aligned}
$$

以及
$$
\begin{aligned}
\mathcal{R}(T) \supset R\left(T S\right)=R\left(I_y-A_2\right)  \Longrightarrow \operatorname{codim} \mathcal{R}(T) \leqslant \operatorname{codim} R\left(I_y-A_2\right)<\infty
\end{aligned}
$$

由于 $\mathcal{N}(T)$ 是 $\mathcal X$ 的闭子空间，考虑商算子
\[
\widetilde T:\mathcal X/\mathcal{N}(T)\to \mathcal Y,\qquad 
\widetilde T([x]) = Tx .
\]
显然 $N(\widetilde T)=\{[0]\}$ 且 $R(\widetilde T)=\mathcal{R}(T)$。为证明 $\mathcal{R}(T)$ 闭，只需证明
$\widetilde T^{-1}:\mathcal{R}(T)\to\mathcal X/\mathcal{N}(T)$ 连续。

采用反证法。若 $\widetilde T^{-1}$ 不连续，则存在序列 $[w_m]\in \mathcal X/\mathcal{N}(T)$ 满足
\[
\widetilde T ([w_m]) = T w_m \to 0\quad \text{于 }\mathcal Y,
\]
但 $[w_m]\not\to [0]$。于是可取子列（仍记为 $[w_m]$），满足
\[
\|[w_m]\| \ge \varepsilon > 0\qquad (\forall m).
\]
令
\[
[x_m] = \frac{[w_m]}{\|[w_m]\|},
\]
则 $\|[x_m]\| = 1$ 对每个 $m$ 成立。可从商类中选取代表元 $x_m\in[x_m]$ 使得
\[
\|x_m\| < 2\qquad (\forall m),
\]
并且
\[
T x_m = \widetilde T([x_m]) = \frac{T w_m}{\|[w_m]\|} \longrightarrow 0.
\]

利用 $S$ 的有界性及 $K_1 = ST - I_{\mathcal X}$ 的定义，有
\[
x_m = STx_m - K_1 x_m .
\]
由于 $T x_m\to 0$，故
\[
ST x_m \to 0,
\]
从而
\[
x_m = - K_1 x_m + o(1).
\]

因为 $\{x_m\}$ 有界且 $K_1$ 为紧算子，存在子列（仍记为 $x_m$）使得
\[
K_1 x_m \longrightarrow y \qquad (\text{于 }\mathcal X).
\]
由上式得
\[
x_m = -K_1 x_m + o(1) \longrightarrow -y.
\]
设 $z = -y$，则 $x_m\to z$ 于 $\mathcal X$ 中。

由 $T x_m \to 0$ 以及 $T$ 的连续性可得
\[
T z = 0,
\]
即 $z\in \mathcal{N}(T)$，故 $[z]=[0]$ 于商空间 $\mathcal X/\mathcal{N}(T)$。

最后注意
\[
\|[x_m]\| 
= \inf_{n\in \mathcal{N}(T)}\|x_m - n\|
\le \|x_m - z\| \longrightarrow 0,
\]
这与 $\|[x_m]\|=1$ 矛盾。因此假设不成立，$\widetilde T^{-1}$ 必连续。

故 $R(\widetilde T)=\mathcal{R}(T)$ 为闭子空间。证毕。
\end{proof}
  \item 设 \(\mathcal{X},\mathcal{Y}\) 为 Banach 空间.证明有界线性算子 \(T : \mathcal{X} \rightarrow  \mathcal{Y}\) 为 Fredholm 算子当且仅当其共轭算子 \({T}^{ * } : {\mathcal{Y}}^{ * } \rightarrow  {\mathcal{X}}^{ * }\) 也是 Fredholm 算子，且 \[
  \operatorname{ind}{T}^{ * } =  - \operatorname{ind}T.
  \]
\begin{proof}
    由上题可知存在有界线性算子$S:\mathcal{Y}\to \mathcal{X}$和紧算子$K_1,K_2$使得\[
    K_1=I_\mathcal{X}-ST,K_2=I_\mathcal{Y}-TS
    \]
    而$K_1^*,K_2^*$也是紧算子,所以$T^*$也是Fredholm算子.

    由$T$是Fredholm算子知$\mathcal{R}(T)$闭,由习题知\[ {}^\perp\mathcal{R}(T) = \mathcal{N}(T^*)\]
由定理知$\dim \mathcal{N}(T^*)=\operatorname{codim}\mathcal{R}(T)$
    
由习题知\[{}^\perp \mathcal{R}(T^*) = \mathcal{N}(T)\]
由定理知$\dim \mathcal{N}(T)=\operatorname{codim}\mathcal{R}(T^*)$.所以$\operatorname{ind}(T)=-\operatorname{ind}(T^*)$.
    
\end{proof}
  \item 设 \(\mathcal{X},\mathcal{Y}\) 为 Banach 空间，\(T \in  \mathcal{L}\left( {\mathcal{X},\mathcal{Y}}\right)\) 具有有限维核空间和闭值域，\(S \in  \mathcal{C}\left( {\mathcal{X},\mathcal{Y}}\right)\).若 \(\{x_{j}\} \subset  \mathcal{X}\) 为有界点列且 \((T + S)x_{j} = y\) 在 \(\mathcal{Y}\) 中成立，则 \(\{x_{j}\}\) 有极限点 \(x \in  \mathcal{X}\) 使得 \((T + S)x = y\).
\begin{proof}
    由推论4.3知$T+K$也有有限维核空间和闭值域,由命题4.3知$x_j$有收敛点列.
\end{proof}
  \item (F. Riesz) 设 \(\mathcal{X}\) 为 Banach 空间，\(S \in  \mathcal{C}\left( \mathcal{X}\right)\).则\[
  N := \{ x \in  \mathcal{X} \mid  \exists\, k > 0,\ {\left( I + S\right)}^{k}x = 0\}
  \]
  为 \(\mathcal{X}\) 的有限维子空间，  \[
  F := \bigcap_{k = 1}^{\infty }{\left( I + S\right)}^{k}\left( \mathcal{X}\right)
  \]
  为 \(\mathcal{X}\) 的闭子空间，且
\[
  \mathcal{X} = N \oplus  F,\quad
  (I + S)|_{N}:N \to N \text{ 为幂零算子},\quad
  (I + S)|_{F}:F \to F \text{ 为双射}.
  \]
\begin{proof}
    由引理4.3知,存在$p\in \mathbb{N}$,使得$N((I+S)^p)=N((I+S)^k),R((I+S)^p)=R((I+S)^k)\forall k>p$,
    那么\[N=\bigcup\limits_{n=1}^pN\left(\left(I+S\right)^n\right)=N((I+S)^p)\]
    为有限维的.
\[  F = \bigcap_{n = 1}^{p }{\left( I + S\right)}^{p}\left( \mathcal{X}\right)=R((I+S)^p)\]
也是闭的且\[\mathcal{X}=N\oplus F\]

      \begin{itemize}
    \item 对任意 $x\in N$，由定义存在 $k$ 使得 $(I+S)^k x=0$。又 $N=N((I+S)^p)$，因此对一切 $x\in N$ 有 $(I+S)^p x=0$，故 $(I+S)|_N$ 的 $p$ 次幂为零算子，即 $(I+S)|_N$ 是幂零算子（nilpotent）。
    \item 对 $F=R((I+S)^p)$，任取 $z\in F$，写 $z=(I+S)^p x$，则把 $y=(I+S)^{p-1}x\in F$，有 $(I+S)y=(I+S)^p x=z$，所以 $(I+S)|_F$ 是满射。若 $y\in F$ 且 $(I+S)y=0$，则 $y\in N\cap F=\{0\}$，从而 $(I+S)|_F$ 也为单射。故 $(I+S)|_F$ 为双射。
  \end{itemize}
\end{proof}
  \item 设 \(\mathcal{X},\mathcal{Y}\) 为 Banach 空间，\(T \in  \operatorname{Fred}\left( {\mathcal{X},\mathcal{Y}}\right)\)，\(S \in  \mathcal{C}\left( {\mathcal{X},\mathcal{Y}}\right)\).令
  \[
  k = \min_{\lambda  \in  \mathbb{K}} \dim \operatorname{Ker}\left( {T + {\lambda S}}\right), \quad
  w = \min_{\lambda  \in  \mathbb{K}} \dim \operatorname{Coker}\left( {T + {\lambda S}}\right),
  \]
  则除非 \(\lambda\) 属于仅有孤立点的闭集，我们都有
  \[
  \dim \operatorname{Ker}\left( {T + {\lambda S}}\right)  = k,\quad
  \dim \operatorname{Coker}\left( {T + {\lambda S}}\right)  = w.
  \]
  \item 设下图为向量空间及线性映射组成的交换图，即 \({jF} = {F}^{\prime }i,\ q{F}^{\prime } = {F}^{\prime \prime }p\)：
  \[
  \begin{matrix}
  0 & \rightarrow & {H}_{1} & \xrightarrow{i} & {H}_{1}^{\prime } & \xrightarrow{p} & {H}_{1}^{\prime \prime } & \rightarrow & 0 \\
  & & F \downarrow & & {F}^{\prime } \downarrow & & {F}^{\prime \prime } \downarrow \\
    0& \rightarrow  &{H}_{2} &\xrightarrow{j} &{H}_{2}^{\prime } &\xrightarrow{q} &{H}_{2}^{\prime \prime } &\rightarrow & 0
  \end{matrix}
  \]

  若垂直序列为正合序列且 \(F,{F}^{\prime },{F}^{\prime \prime }\) 均为 Fredholm 算子，则
  \[
  \operatorname{ind}F - \operatorname{ind}{F}^{\prime } + \operatorname{ind}{F}^{\prime \prime } = 0.
  \]

  \item 设 \(\mathcal{H},{\mathcal{H}}^{\prime }\) 为 Hilbert 空间，\(F : \mathcal{H} \rightarrow  \mathcal{H}\) 和 \(G : {\mathcal{H}}^{\prime } \rightarrow  {\mathcal{H}}^{\prime }\) 均为 Fredholm 算子.证明直和算子
  \[
  F \oplus  G : \mathcal{H} \oplus  {\mathcal{H}}^{\prime } \rightarrow  \mathcal{H} \oplus  {\mathcal{H}}^{\prime }
  \]
  为 Fredholm 算子，且
  \[
  \operatorname{ind}\left( {F \oplus  G}\right)  = \operatorname{ind}F + \operatorname{ind}G.
  \]

\begin{proof}
先说明 $F\oplus G$ 是有界的（显然）并计算其核与像的结构。

1. \emph{核的分解。} 若 $(x,y)\in\mathcal H\oplus\mathcal H'$ 在核中，则
\[
(F\oplus G)(x,y)=(Fx,Gy)=(0,0),
\]
所以 $x\in\ker F$ 且 $y\in\ker G$。反之显然。因此
\[
\ker(F\oplus G)=\ker F\oplus\ker G,
\]
于是
\[
\dim\ker(F\oplus G)=\dim\ker F+\dim\ker G<\infty,
\]
因为 $F,G$ 为 Fredholm，核维数有限。

2. \emph{像的表达与闭性。} 对于直和算子，有
\[
\operatorname{Ran}(F\oplus G)=\operatorname{Ran}F\oplus\operatorname{Ran}G
\subset \mathcal H\oplus\mathcal H'.
\]
既然 $\operatorname{Ran}F$ 与 $\operatorname{Ran}G$ 在各自的空间中为闭（Fredholm 的性质），它们的直和在 $\mathcal H\oplus\mathcal H'$ 中也是闭的。因此 $\operatorname{Ran}(F\oplus G)$ 闭。  

3. \emph{余核（或商空间）维数相加。} 在 Hilbert 空间中可以用正交补来表示余核：
\[
\operatorname{Coker}(F)\cong (\operatorname{Ran}F)^\perp,\qquad
\operatorname{Coker}(G)\cong (\operatorname{Ran}G)^\perp,
\]
并且
\[
\operatorname{Coker}(F\oplus G)\cong \big(\operatorname{Ran}F\oplus\operatorname{Ran}G\big)^\perp.
\]
由于对 Hilbert 空间的直和有正交补的分解
\[
\big(\operatorname{Ran}F\oplus\operatorname{Ran}G\big)^\perp
=(\operatorname{Ran}F)^\perp\oplus(\operatorname{Ran}G)^\perp,
\]
因此
\[
\dim\operatorname{Coker}(F\oplus G)
=\dim(\operatorname{Ran}F)^\perp+\dim(\operatorname{Ran}G)^\perp
=\dim\operatorname{Coker}(F)+\dim\operatorname{Coker}(G)<\infty.
\]

由上述 (1)--(3) 可知 $\ker(F\oplus G)$ 维数有限且 $\operatorname{Ran}(F\oplus G)$ 闭且余维有限，故 $F\oplus G$ 为 Fredholm 算子。

4. \emph{指标相加。} 指数定义为
\[
\operatorname{ind}A=\dim\ker A-\dim\operatorname{Coker}A.
\]
因此
\[
\begin{aligned}
\operatorname{ind}(F\oplus G)
&=\dim\ker(F\oplus G)-\dim\operatorname{Coker}(F\oplus G)\\
&=(\dim\ker F+\dim\ker G)
-(\dim\operatorname{Coker}F+\dim\operatorname{Coker}G)\\
&=(\dim\ker F-\dim\operatorname{Coker}F)+(\dim\ker G-\dim\operatorname{Coker}G)\\
&=\operatorname{ind}F+\operatorname{ind}G.
\end{aligned}
\]

综上所述，$F\oplus G$ 是 Fredholm，且指数可加，证毕。
\end{proof}
\end{problemset}
\chapter{无界算子}

我们前面详细讨论了有界线性算子的有关理论, 但在数学物理中出现的许多重要算子都是无界的,如 \({L}^{2}\left( \Omega \right)\) 上的微分算子及量子力学中的Schrödinger算子等. 因此, 我们有必要系统地认识和掌握无界算子的理论. 前面已经看到, 在Banach空间上有界线性算子的理论中, 算子范数起到非常重要的作用, 但对Banach 空间上的无界线性算子, 由于不存在算子范数, 因此, 我们必须引进另外的手段来研究, 即通过考察算子的图或算子在内积下的性质来研究.

\section{闭算子及其共轭算子}

设 \(\mathcal{X},\mathcal{Y}\) 为Banach 空间,若在乘积空间 \(\mathcal{X} \times  \mathcal{Y}\) 上对 \(\left( {x,y}\right)  \in  \mathcal{X} \times  \mathcal{Y}\) 定义范数 \(\parallel \left( {x,y}\right) \parallel  = \parallel x{\parallel }_{\mathcal{X}} + \parallel y{\parallel }_{\mathcal{Y}}\) ,则 \(\mathcal{X} \times  \mathcal{Y}\) 为Banach空间.
\begin{definition}
    设 \(\mathcal{X},\mathcal{Y}\) 是Banach 空间, \(T\) 是一个线性算子,其定义域 \(\mathcal{D}(T)  \subset  \mathcal{X}\) 为 \(\mathcal{X}\) 的线性子空间,其值域 \(R\left( T\right)  \subset  \mathcal{Y}\) . 称乘积空间 \(\mathcal{X} \times  \mathcal{Y}\) 的线性子空间
\begin{equation*}
\Gamma \left( T\right)  \mathrel{\text{ := }} \{ \left( {x,{Tx}}\right)  \mid  x \in  \mathcal{D}(T) \}
\end{equation*}

为线性算子 \(T\) 的图. 称算子 \(T\) 为闭的如果其图 \(\Gamma \left( T\right)\) 在 \(\mathcal{X} \times  \mathcal{Y}\) 中是闭的.

\end{definition}
\begin{definition}
    设 \(\mathcal{X},\mathcal{Y}\) 为Banach空间, \(T,{T}_{1}\) 是从 \(\mathcal{X}\) 到 \(\mathcal{Y}\) 的线性算子,称 \({T}_{1}\) 为 \(T\) 的一个扩张算子,如果 \(\Gamma \left( T\right)  \subset  \Gamma \left( {T}_{1}\right)\) ,记作 \(T \subset  {T}_{1}\) . 称算子 \(T\) 为可闭的,如果存在扩张算子 \(S \supset  T\) 使得 \(\overline{\Gamma \left( T\right) } = \Gamma \left( S\right)\) . 此时称 \(S\) 为 \(T\) 的闭包,记为 \(S = \bar{T}\) .
\end{definition}
\begin{note}
    从 \(\mathcal{X}\) 到 \(\mathcal{Y}\) 的线性算子 \(T\) 为闭的当且仅当
若 \({x}_{n} \in  \mathcal{D}(T)\) ,且在 \(\mathcal{X} \textbf{中} {x}_{n} \rightarrow  x\) ,在 \(\mathcal{Y} \text{中}T{x}_{n} \rightarrow  y\) ,则 \(x \in  \mathcal{D}(T)\) 且 \(y = {Tx}\) .

 直观来看,获得线性算子 \(T\) 的闭包的最自然地途径就是取其图 \(\Gamma \left( T\right)\) 在乘积空间 \(\mathcal{X} \times  \mathcal{Y}\) 中的闭包,但事实上 \(\overline{\Gamma \left( T\right) }\) 并不一定是某个线性算子的图. 也就是说,并不是每个线性算子都可闭化. 线性算子 \(T\) 可闭化当且仅当若 \(\left( {0,y}\right)  \in  \overline{\Gamma \left( T\right) }\) ,则 \(y = 0\) .

\end{note}
\begin{example}
     设 \(\mathcal{X} = \mathcal{Y} = {L}^{2}\left( \mathbb{R}\right) ,\mathcal{D}(T)  = {C}_{0}^{\infty }\left( \mathbb{R}\right) ,D\left( {T}_{1}\right)  = {C}_{0}^{1}\left( \mathbb{R}\right)\) . 定义
\begin{equation*}
{Tf} \mathrel{\text{ := }} \ci{f}^{\prime }\left( x\right) ,\;\forall f \in  \mathcal{D}(T)
\end{equation*}

及
\begin{equation*}
{T}_{1}f \mathrel{\text{ := }} \ci{f}^{\prime }\left( x\right) ,\;\forall f \in  D\left( {T}_{1}\right) .
\end{equation*}

则 \({T}_{1}\) 为 \(T\) 的一个扩张. 下证 \(\overline{\Gamma \left( T\right) } \supset  \Gamma \left( {T}_{1}\right)\) .

对任意 \(\phi  \in  D\left( {T}_{1}\right)\) ,记其光滑化为
\begin{equation*}
{\phi }_{\delta }\left( x\right)  = {\int }_{-\infty }^{\infty }\phi \left( y\right) {j}_{\delta }\left( {x - y}\right) {\dif y},
\end{equation*}

其中 \({j}_{\delta }\left( x\right)\) 为例磨光核. 由前面章节知,当 \(\delta  \rightarrow  0\) 时,在 \({L}^{2}\left( \mathbb{R}\right)\) 中 \({\phi }_{\delta } \rightarrow  \phi\) . 类似地,
\begin{equation*}
\ci\frac{\dif}{\dif x}{\phi }_{\delta }\left( x\right)  = {\int }_{-\infty }^{\infty }\ci\frac{\dif}{\dif x}{j}_{\delta }\left( {x - t}\right) \phi \left( t\right) {\dif t}
= {\int }_{-\infty }^{\infty } - \ci\left( {\frac{\dif }{\dif t}{j}_{\delta }\left( {x - t}\right) }\right) \phi \left( t\right) {\dif t}
\end{equation*}
\begin{equation*}
= {\int }_{-\infty }^{\infty }{j}_{\delta }\left( {x - t}\right) \ci\frac{\dif }{\dif t}\phi \left( t\right) {\dif t}
\overset{{L}^{2}\left( \mathbb{R}\right) }{ \rightarrow  }\ci\frac{\dif }{\dif x}\phi \left( x\right) .
\end{equation*}

因此,对任意 \(\phi  \in  D\left( {T}_{1}\right)\) ,都有 \({\phi }_{\delta }\overset{{L}^{2}\left( \mathbb{R}\right) }{ \rightarrow  }\phi\) 和 \({T}_{1}{\phi }_{\delta }\overset{{L}^{2}\left( \mathbb{R}\right) }{ \rightarrow  }{T}_{1}\phi\) . 由于对每个 \(\delta  > 0\) , \({\phi }_{\delta } \in  \mathcal{D}(T)\) ,故 \(\overline{\Gamma \left( T\right) } \supset  \Gamma \left( {T}_{1}\right)\) . 从而有 \(\bar{T}\) 为 \({T}_{1}\) 的扩张.

\end{example}

由有限维空间中有界闭集的紧性知, 定义在有限维欧氏空间全空间上的线性算子必连续. 因而, 线性算子的定义范围和它的有界性, 即连续性之间存在一定联系. 由闭图象定理知,定义在全空间 \(\mathcal{X}\) 上的闭算子一定有界. 因此,对无界闭算子 \(T\) 而言,我们主要关心 \(\mathcal{D}(T)  \neq  \mathcal{X}\) 的情形. 由于我们总可以将闭算子的定义域所在的空间 \(\mathcal{X}\) 缩小到其定义域的闭包 \(\overline{\mathcal{D}(T) }\) 来考虑,因此,我们引进
\begin{definition}
    称线性算子 \(T\) 为稠定算子,如果 \(\mathcal{D}(T)\) 在 \(\mathcal{X}\) 中稠密,即 \(\overline{\mathcal{D}(T) } = \mathcal{X}\) .
\end{definition}


前面我们经常用有界算子的共轭算子来描述其性质, 类似地, 我们也可以给出无界算子的共轭算子的概念.
\begin{definition}
    设 \(\mathcal{X},\mathcal{Y}\) 为Banach空间, \({\mathcal{X}}^{ * },{\mathcal{Y}}^{ * }\) 分别表示 \(\mathcal{X},\mathcal{Y}\) 的共轭空间, \(T\) 是从 \(\mathcal{X}\) 到 \(\mathcal{Y}\) 上的稠定算子. 记
\begin{equation*}
D\left( {T}^{ * }\right)  \mathrel{\text{ := }} \left\{  {{y}^{ * } \in  {\mathcal{Y}}^{ * } \mid  \text{ 存在 }x^*\in \mathcal{X}^*\text{ 使得 }\left\langle  {{y}^{ * },{Tx}}\right\rangle=\langle x^*,x\rangle,\forall x \in  \mathcal{D}(T) }\right\}  . \label{5.1}
\end{equation*}
令
\begin{equation*}
{T}^{ * } : {y}^{ * } \mapsto  {x}^{ * },\forall {y}^{ * } \in  D\left( {T}^{ * }\right) ,
\end{equation*}
则称 \({T}^{ * }\) 为 \(T\) 的共轭算子, \(D\left( {T}^{ * }\right)\) 为 \({T}^{ * }\) 的定义域.

\end{definition}
\begin{note}
    $T$必须是稠定算子才是良定义.
    
    若 \(S \subset  T\) ,则 \({T}^{ * } \subset  {S}^{ * }\) .
\end{note}

\begin{example}
     设 \(\mathcal{X} = {L}^{2}\left( \left( {-1,1}\right) \right)\) . 对 \(u \in  {C}_{0}^{2}\left( \left( {-1,1}\right) \right)\) ,定义 \({Tu} = {u}^{\prime \prime }\) . 从而 \(T\) 为稠定算子. 若 \({u}_{n} \in  {C}_{0}^{2}\left( \left( {-1,1}\right) \right)\) ,且在 \({L}^{2}\left( \left( {-1,1}\right) \right)\) 中成立 \({u}_{n} \rightarrow  u,{u}_{n}^{\prime \prime } \rightarrow  f\) ,则
\begin{equation*}
{u}_{n}^{\prime }\left( t\right)  = {\int }_{-1}^{t}{u}_{n}^{\prime \prime }\left( s\right) {\dif s} \rightarrow  {\int }_{-1}^{t}f\left( s\right) {\dif s}
\end{equation*}

在$(-1,1)$上一致收敛. 因此
\begin{equation}
\left\{  \begin{array}{l} u \in  {C}^{1}\left( \left\lbrack  {-1,1}\right\rbrack  \right) , \\  u\left( {\pm 1}\right)  = {u}^{\prime }\left( {\pm 1}\right)  = 0, \\  \text{ 在分布意义下有 }{u}^{\prime \prime } = f \in  {L}^{2}\left( \left( {-1,1}\right) \right)  \end{array}\right.  \label{5.2}
\end{equation}

反之,条件\ref{5.2}意味着 \(\left( {u,f}\right)  \in  \overline{\Gamma \left( T\right) }\) . 实际上,若将 \(u,f\) 零延拓到 \(\left\lbrack  {-1,1}\right\rbrack\) 之外,则在分布意义下 \({u}^{\prime \prime } = f\) 在 \(\mathbb{R}\) 上成立. 定义
\begin{equation*}
{u}_{n}\left( t\right)  \mathrel{\text{ := }} n{\int }_{\left| s\right|  < \frac{1}{2n}}u\left( \frac{n\left( {t - s}\right) }{n - 1}\right) {\dif s} \in  {C}_{0}^{2}\left( \left( {-1,1}\right) \right) ,
\end{equation*}
\begin{equation*}
{u}_{n}^{\prime \prime }\left( t\right)  \mathrel{\text{ := }} \frac{{n}^{3}}{{\left( n - 1\right) }^{2}}{\int }_{\left| s\right|  < \frac{1}{2n}}u\left( \frac{n\left( {t - s}\right) }{n - 1}\right) {\dif s},
\end{equation*}

则当 \(n \rightarrow  \infty\) 时,在 \({L}^{2}\left( \left( {-1,1}\right) \right)\) 中有 \({u}_{n} \rightarrow  u,{u}_{n}^{\prime \prime } \rightarrow  f\) 成立.

另外, \({T}^{ * }u = f\) 意味着 \(u,f \in  {L}^{2}\left( \left( {-1,1}\right) \right)\) 且
\begin{equation*}
\left( {u,{Tv}}\right)  = \left( {f,v}\right) ,\;\forall v \in  {C}_{0}^{2}\left( \left( {-1,1}\right) \right) ,
\end{equation*}

即在分布意义下有 \({u}^{\prime \prime } = f\) 在$(-1,1)$上成立. 这意味着 \(u \in  {C}^{1}\left( \left\lbrack  {-1,1}\right\rbrack  \right)\) .

\end{example}

由上可以看出,算子 \(T\) 与 \({T}^{ * }\) 的差别在于算子 \(T\) 的定义域中要求有\ref{5.2}中边界条件的限制,但算子 \({T}^{ * }\) 的定义域则不要求.

我们自然会问, \(D\left( {T}^{ * }\right)\) 在 \({\mathcal{Y}}^{ * }\) 中是否稠密? 一般来说, \(D\left( {T}^{ * }\right)\) 不一定在 \({\mathcal{Y}}^{ * }\) 中稠密, \(D\left( {T}^{ * }\right)\) 甚至可以为零空间,如
\begin{example}
    设 \(f \notin  {L}^{2}\left( \mathbb{R}\right)\) 为有界可测函数,令
\begin{equation*}
\mathcal{D}(T)  = \left\{  {\psi  \in  {L}^{2}\left( \mathbb{R}\right) | \int \left| f\left( x\right) \psi \left( x\right)  \right|  {\dif x} < \infty }\right\}  .
\end{equation*}

由于 \(\mathcal{D}(T)\) 包含所有具有紧支集的 \({L}^{2}\) 函数,因此, \(\mathcal{D}(T)\) 在 \({L}^{2}\left( \mathbb{R}\right)\) 中稠密. 取 \({\psi }_{0} \in\)  \({L}^{2}\left( \mathbb{R}\right)\) 并对每个 \(\psi  \in  \mathcal{D}(T)\) 定义 \({T\psi } = \left( {f,\psi }\right) {\psi }_{0}\) ,其中 \(\left( {f,\psi }\right)  = \int \left| {f\left( x\right) \psi \left( x\right) }\right| {\dif x}\) . 若 \(\phi  \in\)  \(D\left( {T}^{ * }\right)\) ,则对任意 \(\psi  \in  \mathcal{D}(T)\) 有
\begin{equation*}
\left\langle  {\psi ,{T}^{ * }\phi }\right\rangle   = \langle {T\psi },\phi \rangle  = \left\langle  {\left( {f,\psi }\right) {\psi }_{0},\phi }\right\rangle
= \left( {f,\psi }\right) \left\langle  {{\psi }_{0},\phi }\right\rangle   = \left\langle  {\psi ,\left( {{\psi }_{0},\phi }\right) f}\right\rangle  .
\end{equation*}

因此, \({T}^{ * }\phi  = \left( {{\psi }_{0},\phi }\right) f\) . 由于 \(f \notin  {L}^{2}\left( \mathbb{R}\right)\) ,故必须有 \(\left( {{\psi }_{0},\phi }\right)  = 0\) . 即 \(D\left( {T}^{ * }\right)\) 中任意元素 \(\phi\) 都与 \({\psi }_{0}\) 正交. 因此, \(D\left( {T}^{ * }\right)\) 在 \({L}^{2}\left( \mathbb{R}\right)\) 中不稠密且 \({T}^{ * } = 0\) .

\end{example}

若 \({T}^{ * }\) 是稠定算子,则可进一步定义 \({T}^{* * } = {\left( {T}^{ * }\right) }^{ * }\) .
\begin{theorem}
    任何从 \(\mathcal{X}\) 到 \(\mathcal{Y}\) 的稠定线性算子 \(T\) 的共轭算子 \({T}^{ * }\) 总是闭的.
\end{theorem}
\begin{proof}
    定义线性等距映射
\begin{equation*}
V : \mathcal{X} \times  \mathcal{Y} \rightarrow  \mathcal{Y} \times  \mathcal{X},
\left( {x,y}\right)  \mapsto  \left( {-y,x}\right) .
\end{equation*}

记空间 \({\mathcal{Y}}^{ * } \times  {\mathcal{X}}^{ * }\) 与 \(\mathcal{Y} \times  \mathcal{X}\) 的对偶为
\begin{equation*}
\left\langle  {\left( {{y}^{ * },{x}^{ * }}\right) ,\left( {y,x}\right) }\right\rangle   = \left\langle  {{y}^{ * },y}\right\rangle   + \left\langle  {{x}^{ * },x}\right\rangle  .
\end{equation*}

图 \(\Gamma \left( {T}^{ * }\right)\) 是 \({\mathcal{Y}}^{ * } \times  {\mathcal{X}}^{ * }\) 的线性子空间, \(\Gamma \left( T\right)\) 是 \(\mathcal{X} \times  \mathcal{Y}\) 的线性子空间,因此, \(V\left( {\Gamma \left( T\right) }\right)\) 是 \(\mathcal{Y} \times  \mathcal{X}\) 的线性子空间. 于是考察 \(\Gamma \left( {T}^{ * }\right)\) 中点与 \(V\left( {\Gamma \left( T\right) }\right)\) 中点的对偶知, \(\left( {\phi ,\eta }\right)  \in\)  \(V{\left( \Gamma \left( T\right) \right) }^{ \bot  }\) 当且仅当
\begin{equation*}
\langle \left( {\phi ,\eta }\right) ,\left( {-{T\psi },\psi }\right) \rangle  = 0,\;\forall \psi  \in  \mathcal{D}(T) ,
\end{equation*}

当且仅当
\begin{equation*}
\langle \phi ,{T\psi }\rangle  = \langle \eta ,\psi \rangle ,\;\forall \psi  \in  \mathcal{D}(T) ,
\end{equation*}

当且仅当
\begin{equation*}
\left( {\phi ,\eta }\right)  \in  \Gamma \left( {T}^{ * }\right) .
\end{equation*}

因此, \(\Gamma \left( {T}^{ * }\right)  = V{\left( \Gamma \left( T\right) \right) }^{ \bot  }\) . 由于 \(V{\left( \Gamma \left( T\right) \right) }^{ \bot  }\) 总是 \({\mathcal{Y}}^{ * } \times  {\mathcal{X}}^{ * }\) 的闭子空间,故 \({T}^{ * }\) 为闭算子.
\end{proof}
\begin{theorem}
     (i) 设 \(\mathcal{H}\) 为Hilbert空间, \(T\) 是 \(\mathcal{H}\) 到自身的稠定线性算子,则 \(T\) 为可闭化的当且仅当 \({T}^{ * }\) 是稠定算子. 此时, \(\bar{T} = {T}^{* * }\) .

(ii) 若 \(T\) 是可闭化的,则 \({\left( \bar{T}\right) }^{ * } = {T}^{ * }\) .

\end{theorem}
\begin{proof}
    (i)首先,若 \({T}^{ * }\) 稠定,则由定理5.1知, \({T}^{* * }\) 是 \(H\) 到其自身的闭算子,并且由定理5.1证明知
\begin{equation*}
\Gamma \left( {T}^{* * }\right)  = V{\left( \Gamma \left( {T}^{ * }\right) \right) }^{ \bot  } = V{\left\lbrack  V{\left( \Gamma \left( T\right) \right) }^{ \bot  }\right\rbrack  }^{ \bot  }
= {\left( {V}^{2}{\left( \Gamma \left( T\right) \right) }^{ \bot  }\right) }^{ \bot  } = {\left( \Gamma {\left( T\right) }^{ \bot  }\right) }^{ \bot  } = \overline{\Gamma \left( T\right) }.
\end{equation*}

其中 \(V\) 为定理 5.1 证明中定义的等距线性映射. 由定义知, \(T\) 是可闭化算子.

其次,若 \(D\left( {T}^{ * }\right)\) 在 \(\mathcal{H}\) 中不稠密,则必存在 \({y}_{0} \in  \mathcal{H},{y}_{0} \neq  0\) ,使得 \({y}_{0} \in  D{\left( {T}^{ * }\right) }^{ \bot  }\) . 从而 \(\left( {{y}_{0},0}\right)  \in  \Gamma {\left( {T}^{ * }\right) }^{ \bot  }\) . 显然 \(\left( {0,{y}_{0}}\right)  \in  V{\left( \Gamma \left( {T}^{ * }\right) \right) }^{ \bot  }\) . 这说明 \(V{\left( \Gamma \left( {T}^{ * }\right) \right) }^{ \bot  }\) 不可能是某个线性算子的图. 但 \(\Gamma \left( \bar{T}\right)  = \overline{\Gamma \left( T\right) } = V{\left( \Gamma \left( {T}^{ * }\right) \right) }^{ \bot  }\) . 矛盾,故 \({T}^{ * }\) 是稠定算子.

(ii) 若 \(T\) 可闭化,则由定理5.1知
\begin{equation*}
{T}^{ * } = \overline{\left( {T}^{ * }\right) } = {T}^{* *  * } = {\left( \bar{T}\right) }^{ * }.
\end{equation*}

证毕.

\end{proof}
\begin{definition}
    设 \({\mathcal{H}}_{1},{\mathcal{H}}_{2}\) 为Hilbert空间, \(T\) 为一个线性算子, \(\mathcal{D}(T)  \subset  {\mathcal{H}}_{1},R\left( T\right)  \subset\)  \({\mathcal{H}}_{2}\) . 称 \(T\) 为等距如果
\begin{equation*}
{\left( Tx,Ty\right) }_{2} = {\left( x,y\right) }_{1},\;\forall x,y \in  \mathcal{D}(T) .
\end{equation*}

特别,若还有 \(\mathcal{D}(T)  = {\mathcal{H}}_{1},R\left( T\right)  = {\mathcal{H}}_{2}\) ,则称 \(T\) 为酉算子.
\end{definition}

\begin{theorem}
    设 \({\mathcal{H}}_{1},{\mathcal{H}}_{2}\) 为可分Hilbert空间. 若 \(T\) 为一个线性等距且 \(\mathcal{D}(T)  \subset  {\mathcal{H}}_{1}\) , \(R\left( T\right)  \subset  {\mathcal{H}}_{2}\) ,则其闭包 \(\bar{T}\) 为具有闭定义域及闭值域的等距. \(T\) 能扩张成从 \({\mathcal{H}}_{1}\) 到 \({\mathcal{H}}_{2}\) 的酉算子当且仅当 \(\operatorname{codim} \mathcal{D}\left( \bar{T}\right)  = \operatorname{codim} \mathcal{R}\left( \bar{T}\right)\) .
\end{theorem}
\begin{proof}
    若 \({u}_{n} \in  \mathcal{D}(T)\) 且 \({f}_{n} = T{u}_{n}\) ,则 \(\begin{Vmatrix}{{u}_{n} - {u}_{m}}\end{Vmatrix} = \begin{Vmatrix}{{f}_{n} - {f}_{m}}\end{Vmatrix}\) . 因此, \({u}_{n}\) 收敛到 \(u\) 当且仅当 \({f}_{n}\) 收敛到 \(f\) . 从而有 \(\parallel u\parallel  = \parallel f\parallel\) . 至此,定理第一个结论得证. 若 \(U\) 为 \(T\) 的一个等距扩张, 则
\begin{equation*}
{\left( Ux,Ty\right) }_{2} = {\left( Ux,Uy\right) }_{2} = {\left( x,y\right) }_{1} = 0,\;\forall x \in  {\left( \mathcal{D}(T) \right) }^{ \bot  },y \in  \mathcal{D}(T) .
\end{equation*}

这意味着 \(U\) 映 \({\left( \mathcal{D}(T) \right) }^{ \bot  }\) 到 \(\mathcal{R}{\left( T\right) }^{ \bot  }\) . 若 \(U\) 为酉算子,则 \({\left. U\right| }_{{\left( \mathcal{D}(T) \right) }^{ \bot  }}\) 也是酉算子. 这意味着 \(\dim \mathcal{D}{\left( \bar{T}\right) }^{ \bot  } = \dim R{\left( \bar{T}\right) }^{ \bot  }\) . 反之,若 \(\dim \mathcal{D}{\left( \bar{T}\right) }^{ \bot  } = \dim R{\left( \bar{T}\right) }^{ \bot  }\) ,由 \({H}_{1},{H}_{2}\) 为可分的知,存在酉算子 \({U}_{1} : \mathcal{D}{\left( T\right) }^{ \bot  } \rightarrow  \mathcal{R}{\left( T\right) }^{ \bot  }\) 使得算子
\begin{equation*}
\bar{T} \oplus  {U}_{1} : {\mathcal{H}}_{1} =  \mathcal{D}\left( \bar{T}\right)  \oplus  \mathcal{D}{\left( T\right) }^{ \bot  } \rightarrow   \mathcal{R}\left( \bar{T}\right)  \oplus  \mathcal{R}{\left( T\right) }^{ \bot  } = {\mathcal{H}}_{2}
\end{equation*}

也为酉算子.
\end{proof}

\begin{definition}
    设 \(T\) 是 Hilbert空间 \(\mathcal{H}\) 上的闭算子,称复数 \(\lambda\) 为 \(T\) 的正则值,如果 \({\lambda I} -\)  \(T\) 是 \(\mathcal{D}(T)\) 到 \(\mathcal{H}\) 上的双射且有有界逆,记正则值的全体为 \(\rho \left( T\right)\) . 若 \(\lambda  \in  \rho \left( T\right)\) ,则称 \({R}_{\lambda }\left( T\right)  = {\left( \lambda I - T\right) }^{-1}\) 为 \(T\) 的在 \(\lambda\) 处的预解式.
\end{definition}

\begin{theorem}
    设 \(T\) 是从Hilbert空间 \(\mathcal{H}\) 到其自身的闭稠定线性算子,则 \(T\) 的正则值集 \(\rho \left( T\right)\) 是复平面上的开子集且相应的预解式 \({R}_{\lambda }\left( T\right)\) 是 \(\rho \left( T\right)\) 上的解析算子值函数. 而且
\begin{equation*}
{R}_{\lambda }\left( T\right)  - {R}_{\mu }\left( T\right)  = \left( {\mu  - \lambda }\right) {R}_{\mu }\left( T\right) {R}_{\lambda }\left( T\right) .
\end{equation*}

定理的证明和有界算子情形的相同.
\end{theorem}

\subsection{习 题}

\begin{problemset}

  \item 证明：若闭算子 \(S\) 是稠定对称算子 \(T\) 的一个扩张，则 \(S\) 为 \(\bar{T}\) 的一个扩张。
\begin{proof}
    先证明对称稠定算子 \(T\) 可闭化,只需证明$(0,y)\in \overline{\Gamma(T)}\Rightarrow y=0$,设$x_n\to 0,Tx_n\to  g$,下面证明$g=0$.
    \[\left\langle Tx_n,y\right\rangle=\left\langle x_n,Ty\right\rangle\to 0,\forall y\in \mathcal{D}(T)\]
    由于$T$是稠定算子,那么\[\left\langle g,y\right\rangle=0,\forall y\in\mathcal{D}(T)\Rightarrow g=0\]
    所以$T$可闭化.而\[\Gamma(T)\subset \Gamma(S)=\overline{\Gamma(S)}\Rightarrow \Gamma(\bar{T})\subset \Gamma(S)\]
    所以\(S\) 为 \(\bar{T}\) 的一个扩张.
\end{proof}
  \item 设 \(T\) 为 Hilbert 空间 \(\mathcal{H}\) 上的稠定对称算子，\(\bar{T}\) 为其闭包。  
  证明：对任何非实复数 \(z\)，算子 \({zI} - \bar{T}\) 将 \(D(\bar{T})\) 映到 \(\mathcal{H}\) 的一个闭子空间。
\begin{proof}
    设$z=u+\ci v,v\neq 0$,设$\{zx_n-\bar{T}x_n\},x_n\in \mathcal{D}(\bar{T})$是柯西列.由对称算子的性质知\[\|zx_n-\bar{T}x_n\|^2=\|ux_n-\bar{T}x_n\|^2+|v|^2\|x_n\|\]
    而$v$不等于$0$,所以$\{x_n\}$也是$\mathcal{D}(\bar{T})$上的柯西列,进而$\{\bar{T}x_n\}$也是$\mathcal{R}(\bar{T})$上的柯西列,由于$T$是可闭化的,所以设$x_n\to x$,则有$x\in \mathcal{D}(\bar{T}),y=\bar{T}x$,进而有$zx_n-\bar{T}x_n\to zx-\bar{T}x$,所以算子 \({zI} - \bar{T}\) 将 \(D(\bar{T})\) 映到 \(\mathcal{H}\) 的一个闭子空间。
\end{proof}
  \item 证明：若自伴算子 \(S\) 为稠定对称算子 \(T\) 的一个扩张，则 \(S\) 亦为 \(\bar{T}\) 的一个扩张。
  \begin{proof}
      由于$S$是稠定的,由定理5.1知$S^*=S$是闭算子,进而由第一题知$S$亦为 \(\bar{T}\) 的一个扩张。
      \end{proof}
  \item 设 \({AC}\left[0,1\right]\) 为区间 \(\left[0,1\right]\) 上具有 \({L}^{2}\left[0,1\right]\) 导数的绝对连续函数的集合。  
  设 \({T}_{1}\) 和 \({T}_{2}\) 为 \({L}^{2}\left[0,1\right]\) 上的稠定算子 \(i\dfrac{d}{dx}\)，其定义域分别为
  \[
  D({T}_{1}) = \{ \varphi \mid \varphi \in {AC}[0,1] \},
  \quad
  D({T}_{2}) = \{ \varphi \mid \varphi \in {AC}[0,1],\ \varphi(0) = 0 \}.
  \]
  试证：
  \begin{enumerate}[(i)]
    \item \({T}_{1},{T}_{2}\) 都是闭算子；
    \item \({T}_{1}\) 的谱集为复平面 \(\mathbb{C}\)；
    \item \({T}_{2}\) 的谱集为空集。
  \end{enumerate}
\begin{proof}
    
\end{proof}
  \item 设 \(S\) 是从 \(\mathcal{D}(S)\) 到 Hilbert 空间 \(\mathcal{H}\) 的单射算子。考虑以下关于 \(S\) 的条件：
  \begin{enumerate}[(1)]
    \item \(S\) 为闭算子；
    \item \(\mathcal{R}(S)\) 在 \(\mathcal{H}\) 中稠密；
    \item \(\mathcal{R}(S)\) 为闭集；
    \item 存在常数 \(C > 0\)，使得对所有 \(\psi \in D(S)\) 有 \(\|S\psi\| \ge C\|\psi\|\)。
  \end{enumerate}
  试证明：
  \begin{enumerate}[(i)]
    \item 条件（1）–（3）意味着（4）；
    \item 条件（2）–（4）意味着（1）；
    \item 条件（1）及（4）意味着（3）。
  \end{enumerate}

\end{problemset}

\section{对称算子及自伴算子}

在本节中,我们假设 \(\mathcal{H}\) 为Hilbert 空间, \(T : \mathcal{H} \rightarrow  \mathcal{H}\) 为稠定线性算子.
\begin{definition}
    设 \(\mathcal{H}\) 为Hilbert空间, \(T\) 是 \(\mathcal{H}\) 到其自身的线性稠定算子. 称 \(T\) 为对称的如果 \(T \subset  {T}^{ * }\) . 等价地, \(T\) 为对称算子当且仅当
\begin{equation*}
\left( {{T\phi },\psi }\right)  = \left( {\phi ,{T\psi }}\right) ,\;\forall \phi ,\psi  \in  \mathcal{D}(T) .
\end{equation*}

称 \(T\) 为自伴的如果 \(T = {T}^{ * }\) ,即当且仅当 \(T\) 为对称的且 \(\mathcal{D}(T)  = \mathcal{D}\left( {T}^{ * }\right)\) .

\end{definition}

由于 \(\mathcal{D}\left( {T}^{ * }\right)  \supset  \mathcal{D}(T)\) 在 \(\mathcal{H}\) 中稠密,故对称算子总是可闭化的. 由定理 5.1 知, 若 \(T\) 是对称算子,则 \({T}^{ * }\) 为 \(T\) 的闭扩张,从而由定理 \({5.2}\left( \mathrm{i}\right)\) 知, \(T\) 的最小闭扩张 \({T}^{* * }\) 必须包含在 \({T}^{ * }\) 内. 因此,对于对称算子 \(T\) 有
\begin{equation*}
T \subset  {T}^{* * } \subset  {T}^{ * }\text{ . }
\end{equation*}

对闭对称算子 \(T\) 有
\begin{equation*}
T = {T}^{* * } \subset  {T}^{ * }\text{ . }
\end{equation*}

对自伴算子 \(T\) 有
\begin{equation*}
T = {T}^{* * } = {T}^{ * }.
\end{equation*}

由此可知,闭对称算子 \(T\) 为自伴的当且仅当 \({T}^{ * }\) 为对称的. 闭对称算子和自伴算子有着本质不同, 在此, 我们先看看自伴算子的若干性质.
\begin{proposition}
     设 \(T\) 为闭对称算子. 若 \(\lambda  \in  \mathbb{C} \smallsetminus  \mathbb{R}\) ,那末 \({\lambda I} - T\) 为单射, \(R\left( {{\lambda I} - T}\right)\) 为闭的,而且
\begin{equation*}
\left| {\Im\lambda }\right| \parallel u\parallel  \leq  \parallel \left( {{\lambda I} - T}\right) u\parallel ,\;\forall u \in  \mathcal{D}(T) . \label{5.3}
\end{equation*}

\(R\left( {{\lambda I} - T}\right)\) 的余维仅依赖于 \(\Im\lambda\) 的符号.
\end{proposition}
\begin{proof}
    设 \(\lambda  = x + {iy}\) ,则由 \(T\) 的对称性知,
\begin{equation*}
\parallel \left( {{\lambda I} - T}\right) u{\parallel }^{2} = \parallel \left( {{xI} - T}\right) u{\parallel }^{2} + {y}^{2}\parallel u{\parallel }^{2} \geq  {y}^{2}\parallel u{\parallel }^{2}, \label{5.4}
\end{equation*}

即上式成立. 同时意味着 \({\lambda I} - T\) 为单射.

若 \({u}_{n} \in  \mathcal{D}(T)\) 且 \(\left( {{\lambda I} - T}\right) {u}_{n} \rightarrow  f\) ,对 \({u}_{n} - {u}_{m}\) 应用不等式得, \({u}_{n} \rightarrow  u\) ,因而, \(u \in  \mathcal{D}(T)\) 且 \(\left( {{\lambda I} - T}\right) u = f\) ,至此, \(R\left( {{\lambda I} - T}\right)\) 为闭的得证.

对 \(\left( {u,{Tu}}\right)  \in  \Gamma \left( T\right)\) ,定义
\begin{equation*}
{T}_{1}\left( {u,{Tu}}\right)  \mathrel{\text{ := }} \left( {{\lambda I} - T}\right) u,\;S\left( {u,{Tu}}\right)  \mathrel{\text{ := }} u.
\end{equation*}

不难验证算子 \({T}_{1} : \Gamma \left( T\right)  \rightarrow  \mathcal{H},S : \Gamma \left( T\right)  \rightarrow  \mathcal{H}\) 满足定理的条件且 \(\operatorname{Ker}{T}_{1} =\)  \(\{ 0\}\) ,从而有当 \(\left| \zeta \right|\) 充分小时, \(R\left( {\left( {{\lambda } + \zeta }\right) I - T}\right)\) 与 \(R\left( {{\lambda I} - T}\right)\) 具有相同的余维. 由于半平面 \(\{ \lambda  \mid   \pm  \Im\lambda  > 0\}\) 连通知 \(R\left( {{\lambda I} - T}\right)\) 的余维仅依赖于 \(\Im\lambda\) 的符号. 定理得证.
\end{proof}
\begin{definition}
    如果 \({\lambda I} - T\) 为从 \(\mathcal{D}(T)\) 到 \(\mathcal{H}\) 的双射,则称复数 \(\lambda\) 属于算子 \(T\) 的预解集.
\end{definition}
\begin{note}
    复数 \(\lambda\) 属于算子 \(T\) 的预解集意味着 \({\lambda I} - T\) 有逆算子 \({R}_{\lambda }\left( T\right)  \mathrel{\text{ := }} {\left( \lambda I - T\right) }^{-1}\) : \(H \rightarrow  \mathcal{D}(T)\) . 由命题 5.1 知 \(\begin{Vmatrix}{{R}_{\lambda }\left( T\right) }\end{Vmatrix} \leq  \frac{1}{\left| \Im\lambda \right| }\) .

\end{note}

\begin{theorem}
    闭对称算子 \(T\) 为自伴的当且仅当所有 \(\lambda  \in  \mathbb{C} \smallsetminus  \mathbb{R}\) 都属于算子 \(T\) 的预解集.
\end{theorem}
\begin{proof}
    由命题5.1知, \(R\left( {{\lambda I} - T}\right)\) 为 \(\mathcal{H}\) 的闭子空间. 下证 \(R\left( {{\lambda I} - T}\right)  = \mathcal{H}\) . 若不然,则存在非零向量 \(k \in  \mathcal{H}\) 与 \(R\left( {{\lambda I} - T}\right)\) 正交,即
\begin{equation*}
\begin{aligned}
    \left( {{\lambda v} - {Tv},k}\right)  = \left( {v,\bar{\lambda }k}\right)  - \left( {{Tv},k}\right)  = 0,\;\forall v \in  \mathcal{D}(T) .\\
\left( {{Tv},k}\right)  =     \left( {v,\bar{\lambda }k}\right)  ,\;\forall v \in  \mathcal{D}(T) \Rightarrow k\in\mathcal{D}(T^*),T^*(k)=\bar{\lambda}k
\end{aligned}
\end{equation*}

由 \(T\) 的自伴性知, \(k \in  \mathcal{D}(T)\) 且 \({Tk} = \bar{\lambda }k\) . 从而 \(\left( {k,{Tk}}\right)  = \lambda \left( {k,k}\right)\) 为非实数. 这与 \(T\) 为对称算子矛盾. 因此, \(R\left( {{\lambda I} - T}\right)  = \mathcal{H}\) . 又由命题5.1知 \({\lambda I} - T\) 为单射. 因此, \({\lambda I} - T\) 将 \(\mathcal{D}(T)\) 一一对应映到 \(\mathcal{H}\) 上,即 \(\lambda\) 属于 \(T\) 的预解集.

反之,任取 \(\lambda  \in  \mathbb{C} \smallsetminus  \mathbb{R}\) ,则 \(\lambda\) 及 \(\bar{\lambda }\) 都属于 \(T\) 的预解集. 下面先证 \({R}_{\lambda }{\left( T\right) }^{ * } = {R}_{\bar{\lambda }}\left( T\right)\) .

实际上,由 \(T\) 为对称算子知,
\begin{equation*}
\left( {x,\left( {\bar{\lambda }I - T}\right) y}\right)  = \left( {\left( {{\lambda I} - T}\right) x,y}\right) ,\;\forall x,y \in  \mathcal{D}(T) . \label{5.5}
\end{equation*}

若记
\begin{equation*}
{\left( \lambda I - T\right) }^{-1}f \mathrel{\text{ := }} x,\;{\left( \bar{\lambda }I - T\right) }^{-1}g \mathrel{\text{ := }} y, \label{5.6}
\end{equation*}

由于 \({\lambda I} - T\) 及 \(\bar{\lambda }I - T\) 将 \(\mathcal{D}(T)\) 满射映到 \(\mathcal{H}\) 上,故由上上式知
\begin{equation*}
\left( {{R}_{\lambda }\left( T\right) f,g}\right)  = \left( {f,{R}_{\bar{\lambda }}\left( T\right) g}\right) ,\;\forall f,g \in  \mathcal{H}. \label{5.7}
\end{equation*}

这意味着 \({R}_{\lambda }{\left( T\right) }^{ * } = {R}_{\bar{\lambda }}\left( T\right),\mathcal{D}(R_\lambda(T)^*)=\mathcal{H}\) .

最后,我们来证 \(T\) 为自伴的. 为此,只需证若 \(v \in  D\left( {T}^{ * }\right)\) ,则 \(v \in  \mathcal{D}(T)\) 且 \({T}^{ * }v =\)  \({Tv}\) . 任取 \(v \in  D\left( {T}^{ * }\right)\) ,记 \({T}^{ * }v = w\) . 从而
\begin{equation*}
\left( {{Tx},v}\right)  = \left( {x,w}\right) ,\;\forall x \in  \mathcal{D}(T) .
\end{equation*}

上式两边同时减 \(\lambda \left( {x,v}\right)\) 得
\begin{equation*}
\left( {\left( {{\lambda I} - T}\right) x,v}\right)  = \left( {x,\bar{\lambda }v - w}\right) ,\;\forall x \in  \mathcal{D}(T) .
\end{equation*}

令 \(g = \bar{\lambda }v - w\) 得
\begin{equation*}
\left( {f,v}\right)  = \left( {{R}_{\lambda }\left( T\right) f,\bar{\lambda }v - w}\right)  = \left( {f,{R}_{\bar{\lambda }}\left( T\right) \left( {\bar{\lambda }v - w}\right) }\right) ,\;\forall f \in  \mathcal{H}.
\end{equation*}

因此,
\begin{equation*}
v = {R}_{\bar{\lambda }}\left( T\right) \left( {\bar{\lambda }v - w}\right) .
\end{equation*}

由于 \(R\left( {\left( \bar{\lambda }I - T\right) }^{-1}\right)  = \mathcal{D}(T)\) ,故 \(v \in  \mathcal{D}(T)\) . 上式两边同时用 \(\bar{\lambda }I - T\) 作用得 \({Tv} = w\) , 即 \({Tv} = {T}^{ * }v\) .

\end{proof}
\begin{example}
    设 \(\mathcal{H} = {L}^{2}\left( \mathbb{R}\right)\) . 在 \(\mathcal{D}(T)  = {C}_{0}^{1}\left( \mathbb{R}\right)\) 上定义算子 \(T = \ci\frac{d}{\dif x}\) . 则 \(T\) 为对称算子且 \(\bar{T}\) 为自伴的.
\end{example}

\begin{proof}
由分部积分知, \(T\) 为对称算子. 对 \(\lambda  \in  \mathbb{C}\) ,则 \({\lambda I} - T\) 的值域为
\begin{equation*}
R\left( {{\lambda I} - T}\right)  = \left\{  {f \mid  \text{ 存在 }u \in  {C}_{0}^{1}\left( \mathbb{R}\right) \text{ 使得 }\ci\frac{\dif}{\dif x}u - {\lambda u} = f}\right\}  .
\end{equation*}

在 \(f\) 的定义式两边同乘 \({e}^{\ci\lambda x}\) 得
\begin{equation*}
\frac{\dif}{\dif x}\left( {{e}^{\ci\lambda x}u}\right)  = {e}^{\ci\lambda x}f.
\end{equation*}

由于 \(u\) 具有紧支集,两边积分得
\begin{equation*}
0 = {\int }_{-\infty }^{\infty }{e}^{\ci\lambda x}{f\dif x} \label{5.8}
\end{equation*}

即 \(R\left( {{\lambda I} - T}\right)\) 中每个函数 \(f\) 都满足上述条件. 反之,我们来证明 \(\mathbb{R}\) 上满足上述条件的具有紧支集的连续函数 \(f\) 都属于 \(R\left( {{\lambda I} - T}\right)\) . 为此,定义
\begin{equation*}
u\left( x\right)  \mathrel{\text{ := }}  - \ci{\int }_{-\infty }^{x}{e}^{{\ci\lambda }\left( {y - x}\right) }f\left( y\right) {\dif y}. \label{5.9}
\end{equation*}

显然, \(u\) 具有连续导数,且若 \(f\) 在紧区间 \(S\) 外为零,由 上式知, \(u\) 在 \(S\) 之外也为零.

由于 \({e}^{\ci\lambda x}\) 在 \(\mathbb{R}\) 上不是平方可积的,因此,满足条件的具有紧支集的连续函数 \(f\) 在 \({L}^{2}\left( \mathbb{R}\right)\) 中稠密. 由命题5.1知,对任意 \(\lambda  \in  \mathbb{C} \smallsetminus  \mathbb{R},R\left( {{\lambda I} - \bar{T}}\right)\) 为闭的. 因此, \(R\left( {{\lambda I} - \bar{T}}\right)  = \mathcal{H}\) . 又由 \({\lambda I} - \bar{T}\) 为单射,故 \(\lambda\) 属于 \(\bar{T}\) 的预解集. 由定理5.1知, \(\bar{T}\) 为自伴的.
\end{proof}



关于自伴算子的预解式, 我们有
\begin{theorem}
     若 \(T\) 为自伴算子,则对任意 \(\lambda ,\mu  \in  \mathbb{C} \smallsetminus  \mathbb{R}\) 有
\begin{equation*}
{R}_{\lambda }\left( T\right)  - {R}_{\mu }\left( T\right)  = \left( {\mu  - \lambda }\right) {R}_{\mu }\left( T\right) {R}_{\lambda }\left( T\right)  = \left( {\mu  - \lambda }\right) {R}_{\lambda }\left( T\right) {R}_{\mu }\left( T\right) . \label{5.10}
\end{equation*}

\end{theorem}
\begin{proof}
    由命题5.1知,只需证上. 由于 \({R}_{\mu }\left( T\right)\) 的值域包含在 \(T\) 的定义域内, 故
    \begin{equation*}
        \begin{aligned}
R_\lambda-R_\mu & =(T-\lambda I)^{-1}-(T-\mu I)^{-1} \\
& =(T-\lambda I)^{-1}[(T-\mu I)-(T-\lambda I)](T-\mu I)^{-1} \\
& =(\mu-\lambda)(T-\lambda I)^{-1}(T-\mu I)^{-1} \\
& =(\mu-\lambda) R_\lambda R_\mu,
\end{aligned}
    \end{equation*}



\end{proof}


由命题5.1知, \(R\left( {{\lambda I} - T}\right)\) 的余维仅依赖于 \({\Im\lambda }\) 的符号. 因此,我们可以定义对称算子的亏损指标如下:
\begin{definition}
    当 \(\pm  \Im\lambda  > 0\) 时,称 \(\overline{R\left( {{\lambda I} - T}\right) }\) 的余维数 \({n}_{ \pm  }\) 为对称算子 \(T\) 的亏损指标.
\end{definition}
\begin{proposition}
    设 \(T\) 为闭稠定对称算子. 则 \(T\) 为自伴的当且仅当 \(T\) 的亏损指标 \({n}_{ \pm  } = 0\) .
\end{proposition}
\begin{proof}
    由于 \(\left( {u,\left( {{\lambda I} - T}\right) v}\right)  = 0,v \in  \mathcal{D}(T)\) 等价于 \(u \in  D\left( {T}^{ * }\right)\) 且 \(\left( {\bar{\lambda }I - {T}^{ * }}\right) u = 0\) , 故由命题5.1得
\begin{equation*}
{n}_{ \pm  } = \dim \operatorname{Ker}\left( {{\lambda I} - {T}^{ * }}\right) , \pm  \Im\lambda  < 0. \label{5.11}
\end{equation*}

若 \(T\) 为自伴算子,则由命题5.1知,对任意 \(\lambda  \in  \mathbb{C} \smallsetminus  \mathbb{R}\) 有 \({\lambda I} - {T}^{ * }\) 为单射. 故 \({n}_{ + } =\)  \({n}_{ - } = 0\) . 反之,若 \(u \in  D\left( {T}^{ * }\right) ,\Im\lambda  > 0\) ,由 \({n}_{ + } = 0\) 知,存在 \(v \in  \mathcal{D}(T)\) 使得 \(({\lambda I} - {T})v = \left( {{\lambda I} - {T}^{ * }}\right) u\text{ . 从而有 }\left( {{\lambda I} - {T}^{ * }}\right) \left( {u - v}\right)  = 0\) 故由$ \dim \operatorname{Ker}\left( {{\lambda I} - {T}^{ * }}\right) =0$及 \({n}_{ - } = 0\text{ 知 }u - v = 0\). 因此, \(u \in  \mathcal{D}(T)\) ,即 \(T\) 为自伴的.
\end{proof}


\subsection{习 题}
\begin{problemset}
    \item 若 \(T\) 为自伴算子,则对任意非负整数 \(k,{T}^{k}\) 亦为自伴算子.
    \begin{proof}
        由于$T$是自伴算子,那么\[
        (T^k)^*=(T^*)^k=T^k
        \]
    
    \end{proof}
    \item 设 \(T\) 为Hilbert 空间 \(\mathcal{H}\) 上的自伴算子. 则

(i) \({T}^{2} + I\) 有有界逆.

(ii) \({\left( {T}^{2} + I\right) }^{-n}\) 映 \(H\) 到 \(D\left( {T}^{2n}\right)\) .

\begin{proof}
    (i)由于$T$是自伴的,那么\[T-\ci I,T+\ci\]有有界逆.进而\[
    T^2+I=(T-\ci I)(T+\ci I)
    \]
    也有有界逆.

    (ii) 
\end{proof}
\item  设 \(T\) 为Hilbert空间 \(\mathcal{H}\) 上的有界自伴算子. 则下列条件等价

(i) \(T\) 具有非负谱;

(ii) 存在算子 \(S\) 使得 \(T = {S}^{ * }S\) ;

(iii) 对任意 \(v \in  \mathcal{H}\) 都有 \(\left( {v,{Tv}}\right)  \geq  0\) .

若算子 \(T\) 满足上述三条中的一条,则称 \(T\) 为正算子,记为 \(T \geq  0\) .
\begin{proof}
(iii)$\Rightarrow$(i) 由习题5.2.6可知.

(ii)$\Rightarrow$(iii) 显然成立.


\end{proof}
\item (i) 证明例题5.1.2中的算子 \({T}^{ * }\) ,且 \({T}^{ * }\) 的共轭算子 \({\left( {T}^{ * }\right) }^{ * }\) 为 \(T\) 的闭包.

(ii) 若 \({a}_{ + },{b}_{ + },{a}_{ - },{b}_{ - } \in  \mathbb{R}\) 满足 \({a}_{ - }^{2} + {b}_{ - }^{2} \neq  0,{a}_{ + }^{2} + {b}_{ + }^{2} \neq  0\) . 设 \(D\left( \widetilde{T}\right) ,D\left( {T}^{ * }\right)\) 为例题5.1.2中算子. 令
\begin{equation*}
D\left( \widetilde{T}\right)  \mathrel{\text{ := }} \left\{  {u \in  D\left( {T}^{ * }\right)  \mid  {a}_{ + }u\left( 1\right)  + {b}_{ + }{u}^{\prime }\left( 1\right)  = 0,{a}_{ - }u( - 1) + {b}_{ - }{u}^{\prime }\left( {-1}\right)  = 0}\right\}  .
\end{equation*}

证明 \({T}^{ * }\) 在 \(D\left( \widetilde{T}\right)\) 上的限制算子 \(\widetilde{T}\) 为自伴的.

\begin{proof}
   (i) 由于$L^2$是Hilbert空间,那么$T^*$也是稠定的,由定理知$\overline{T}=T^{**}$.

   (ii)
\end{proof}\item 证明有界对称算子 \(T\) 的谱半径等于其范数,即
\begin{equation*}
\left| {r\left( T\right) }\right|  = \parallel T\parallel .
\end{equation*}
\begin{proof}
    先证明$\|T^{2^n}\|=\|T\|^{2^n},n\in \mathbb{N}$,显然$\|T^k\|\leqslant \|T\|^k$,当$n=1$时有
    \[\|Tx\|^2=\left\langle T^2x,x\right\rangle\leqslant \|T^2\|\|x\|^2\Rightarrow \|T^2\|\geqslant\|T\|^2\]
    那么有$\|T^2\|=\|T\|^2$,由数学归纳法知$\|T^{2^n}\|=\|T\|^{2^n},n\in \mathbb{N}$.
    那么由Gelfand定理知\[|r(T)|=\lim\limits_{n\to \infty}\|T^{2^n}\|^{\frac{1}{2^n}}=\|T\|\]
    
\end{proof}
\item  设 \(T\) 为有界对称算子. 令
\begin{equation*}
a \mathrel{\text{ := }} \mathop{\inf }\limits_{{\parallel x\parallel  = 1}}\left( {x,{Tx}}\right) ,\;b \mathrel{\text{ := }} \mathop{\sup }\limits_{{\parallel x\parallel  = 1}}\left( {x,{Tx}}\right) .
\end{equation*}

证明 \(T\) 的谱包含在闭区间 \(\left\lbrack  {a,b}\right\rbrack\) 且端点 \(a,b\) 属于 \(T\) 的谱集.
\begin{proof}
    若 $\lambda>b$,对所有单位向量 $\|x\|=1$ 有：
$$
(x, T x) \leq b \quad \Rightarrow \quad(x,(T-\lambda I) x)=(x, T x)-\lambda<b-\lambda<0 .
$$

对任意 $x \neq 0$ ，令 $u=x /\|x\|$ ，则
$$
((T-\lambda I) x, x)=\|x\|^2((T-\lambda I) u, u) \leq\|x\|^2(b-\lambda)<0 .
$$

因此存在常数
$$
\delta:=\lambda-b>0
$$

使对所有 $x$ ，
$$
\|(T-\lambda I) x\|\|x\| \geq|((T-\lambda I) x, x)| \geq \delta\|x\|^2 .
$$

故
$$
\|(T-\lambda I) x\| \geq \delta\|x\| .
$$


这说明 $T-\lambda I$ 单射且有闭值域$\mathcal{H}=\Ker (T-\lambda I)^\perp$，并且下界大于零,那么逆算子存在且有界。因此 $T-\lambda I$ 可逆，即 $\lambda \notin \sigma(T)$ 。

若 $\lambda<a$完全类似。所以
$$
\sigma(T) \subset[a, b] .
$$

要证明 $a$ 属于谱，只需证明：$T-a I$ 不可逆．假设 $T-a I$ 可逆。则存在常数 $\delta>0$ 使
$$
\|(T-a I) x\| \geq \delta\|x\|, \quad \forall x
$$

同样利用 Cauchy 不等式：
$$
|((T-a I) x, x)| \leq\|(T-a I) x\|\|x\| \quad \Rightarrow \quad((T-a I) x, x) \geq \delta\|x\|^2 .
$$

归一化 $\|x\|=1$ ，可得
$$
(x, T x)-a=((T-a I) x, x) \geq \delta .
$$

即
$$
(x, T x) \geq a+\delta, \quad \forall\|x\|=1
$$

但这是不可能的，故假设矛盾。同样证明$b$的情况.
\end{proof}
\item 设 \(T\) 为紧致自伴算子, \(\lambda  \in  \mathbb{R} \smallsetminus  \{ 0\}\) . 则 \(\operatorname{Ker}\left( {{\lambda I} - T}\right)\) 为有限维且 \({\lambda I} - T\) 限制到 \(\operatorname{Ker}\left( {{\lambda I} - T}\right)\) 的正交补上为双射.
\end{problemset}
\begin{proof}
    由$T$为紧算子知\(\operatorname{Ker}\left( {{\lambda I} - T}\right)\) 为有限维,设$\mathcal{H}=\operatorname{Ker}\left( {{\lambda I} - T}\right)\oplus M$,下面证明$\lambda I-T$是单射\[
    (\lambda I-T)x=0\Rightarrow x\in \Ker(\lambda I-T)\bigcap M=\{0\}\Rightarrow x=0 
    \]
    
    再证$\lambda I-T$是满射,设$x\in M,y\in \Ker(\lambda I-T)$\[
    \left\langle (\lambda I-T)x,y\right\rangle=  \left\langle x,(\lambda I-T)y\right\rangle =0\Rightarrow  (\lambda I-T)x\in M
    \]
  
    由于$M$是闭的,那么$ \lambda I-T$在$M$上限制是闭算子,假设不是满的,那么存在非零$z\in M,z$与$\mathcal{R} (\lambda I-T|_M)$正交,那么有\[
      \forall  x\in M, \left\langle (\lambda I-T)x,z\right\rangle=  \left\langle x,(\lambda I-T)z\right\rangle =0\Rightarrow  (\lambda I-T)z\in M^\perp=\Ker(\lambda I-T)
    \]
    
    进而由\[(\lambda I-T)z\in M\bigcap M^\perp\Rightarrow (\lambda I-T)z=0\Rightarrow z\in M\bigcap M^\perp \Rightarrow z=0 \]
    
    那么与$z$非零矛盾,所以是满射.
\end{proof}

\section{对称算子的自伴扩张}

定理5.3给出了线性等距算子能扩张成酉算子的充分必要条件. 实际上,自伴算子的扩张问题总是可以约化到等距算子的扩张问题.

\begin{proposition}
     若 \(T\) 为稠定对称算子,则算子
\begin{equation}
A : \left( {{iI} + T}\right) u \mapsto  \left( {-{iI} + T}\right) u,\;u \in  \mathcal{D}(T) , \label{5.12}
\end{equation}

为等距, \(\operatorname{Ker}\left( {I - A}\right)  = \{ 0\} ,\overline{R\left( {I - A}\right) } = \mathcal{H}\) 且
\begin{equation}
\mathcal{D}(T)  = \{ v - {Av} \mid  v \in  D\left( A\right) \} ,\;T\left( {v - {Av}}\right)  = i\left( {v + {Av}}\right) . \label{5.13}
\end{equation}

反之,若 \(A\) 满足 \(\overline{R\left( {I - A}\right) } = \mathcal{H}\) 的等距算子,则 \(\operatorname{Ker}\left( {I - A}\right)  = \{ 0\}\) 且由上式定义的算子 \(A\) 为稠定对称算子. \(A\) 为闭算子当且仅当 \(T\) 为闭算子.
\end{proposition}
\begin{proof}
    若 \(T\) 为对称算子,则
\begin{equation*}
\parallel \left( {{iI} + T}\right) u\parallel  = \parallel \left( {-{iI} + T}\right) u\parallel ,\forall u \in  \mathcal{D}(T) ,
\end{equation*}

即\ref{5.12}定义的算子 \(A\) 为等距. 记 \(w = {Tu} + {iu},{Aw} = {Tu} - {iu}\) . 则有
\begin{equation*}
{2iu} = w - {Aw},\;{2Tu} = w + {Aw}.
\end{equation*}

若令 \(w = {2iv}\) ,则可得 \(u = v - {Av}\) 且 \({Tu} = i\left( {v + {Av}}\right)\) . 因此,若 \(v \neq  0\) ,则 \(v - {Av} \neq  0\) , 即 \(\operatorname{Ker}\left( {I - A}\right)  = 0\) 且 \(R\left( {I - A}\right)  = \mathcal{D}(T)\) .

反之,假设等距算子 \(A\) 满足 \(\overline{R\left( {I - A}\right) } = \mathcal{H}\) . 若 \(v \in  D\left( A\right)\) 且 \({Av} = v\) ,则对任意 \(w \in  D\left( A\right)\) ,
\begin{equation*}
\left( {w - {Aw},v}\right)  = \left( {w,v}\right)  - \left( {{Aw},v}\right)  = \left( {{Aw},{Av}}\right)  - \left( {{Aw},v}\right)  = \left( {{Aw},{Av} - v}\right)  = 0.
\end{equation*}

这意味着 \(v = 0\) ,即 \(\operatorname{Ker}\left( {I - A}\right)  = 0\) . 因此, \(v - {Av}\) 决定 \(v\) . 从而,可以由\ref{5.13}式定义算子 \(T\) . 由于 \(A\) 为等距,即 \(\left( {v,w}\right)  = \left( {{Av},{Aw}}\right)\) ,故对任意 \(v,w \in  D\left( A\right)\) 有
\begin{equation*}
\left( {T\left( {v - {Av}}\right) ,w - {Aw}}\right)  = \left( {v - {Av},T\left( {w - {Aw}}\right) }\right) .
\end{equation*}

从而 \(T\) 为稠定对称算子. 由于 \(A\) 与 \(T\) 的图可以通过连续线性同构来连接,从而定理得证.

\end{proof}
\begin{proposition}
        稠定闭对称算子 \(T\) 有自伴扩张当且仅当其亏损指标相等.
\end{proposition}

\begin{proof}
    由命题5.3知, \(T\) 的自伴扩张的存在性等价于相应等距算子 \(A\) 的酉扩张的存在性. 又由于亏损指标是 \(R\left( {\mp {iI} + T}\right)\) 的余维,即 \(D\left( A\right)\) 和 \(R\left( A\right)\) 的余维,从而由定理5.3知定理结论成立.
\end{proof}

接下来, 我们给出一个自伴扩张存在的充分条件
\begin{theorem}
    设稠定对称算子 \(T\) 满足存在 \(\lambda  \in  \mathbb{R}\) 及常数 \(C\) 使得
\begin{equation*}
\parallel u\parallel  \leq  C\parallel \left( {{\lambda I} - T}\right) u\parallel ,\;u \in  \mathcal{D}(T) .
\end{equation*}

则 \(T\) 的亏损指标相等,从而 \(T\) 有自伴扩张.
\end{theorem}



\begin{proof}
不失一般性，假设 $T$ 为闭算子。由给定的不等式可知
\[
\|(\lambda I - T)u\|\ge \frac{1}{C}\|u\|,
\]

因此 $(\lambda I - T)$ 是单射，并且其逆算子在 $R(\lambda I - T)$ 上有界。由此可知
\[
R(\lambda I - T)
\quad \text{是闭的子空间}.
\]

取 $z\in\mathbb{C}$ 使得 $|z-\lambda|$ 足够小。将算子写为
\[
zI - T
= (\lambda I - T) - (\lambda - z)I
= (\lambda I - T)\bigl(I - (\lambda - z)(\lambda I - T)^{-1}\bigr).
\]
由于 $(\lambda I - T)^{-1}$ 有界，当
\[
|z-\lambda|\;\|(\lambda I - T)^{-1}\|<1
\]
时，算子
\[
I - (\lambda - z)(\lambda I - T)^{-1}
\]

可由 Neumann 级数倒数，从而是可逆的。因此在该范围内，
\[
R(zI - T)
= R(\lambda I - T).
\]

由此得到：当 $z$ 足够接近 $\lambda$ 时，
\[
\operatorname{codim} R(zI - T)
= \operatorname{codim} R(\lambda I - T).
\]

特别地，取 $z = \lambda \pm i\varepsilon$，其中 $\varepsilon>0$ 足够小，则根据 
\[
\operatorname{codim} R\bigl((\lambda+i\varepsilon)I - T \bigr)
=
\operatorname{codim} R\bigl((\lambda-i\varepsilon)I - T \bigr).
\]

另一方面，根据亏损指标的定义，
\[
n_{+}(T)
= \dim \ker (T^{*}-iI)
= \operatorname{codim} R(iI - T),
\]
\[
n_{-}(T)
= \dim \ker (T^{*}+iI)
= \operatorname{codim} R(-iI - T).
\]

令 $\varepsilon\downarrow 0$，由上式得到
\[
\operatorname{codim}R(iI - T)
=
\operatorname{codim}R(-iI - T),
\]
即
\[
n_{+}(T)=n_{-}(T).
\]

亏损指标相等意味着 $T$ 可自伴扩张（按照 von Neumann 扩张理论）。
\end{proof}

\begin{theorem}
    设 \(T\) 为稠定对称算子满足
\begin{equation*}
C = \mathop{\inf }\limits_{{u \in  \mathcal{D}(T) ,\parallel u\parallel  = 1}}\left( {{Tu},u}\right)  >  - \infty .
\end{equation*}

则存在 \(T\) 的自伴扩张 \({T}_{F}\) 满足 \(\left( {{T}_{F}u,u}\right)  \geq  C\parallel u{\parallel }^{2},u \in  D\left( {T}_{F}\right)\) .

\end{theorem}
\begin{proof}
    不妨假设 \(C = 1\) ,否则,用 \(T + \left( {1 - C}\right) I\) 来代替 \(T\) . 令
\begin{equation*}
{\left( u,v\right) }_{T} = \left( {u,{Tv}}\right) ,\;\parallel u{\parallel }_{T}^{2} = {\left( u,u\right) }_{T},\;u,v \in  \mathcal{D}(T) .
\end{equation*}

用 \({\mathcal{H}}_{T}\) 记 \(\mathcal{D}(T)\) 关于范数 \(\parallel  \cdot  {\parallel }_{T}\) 的完备化. 由于对任意 \(u \in  \mathcal{D}(T)\) 有 \(\parallel u\parallel  \leq  \parallel u{\parallel }_{T}\) , 故存在从 \({\mathcal{H}}_{T}\) 到 \(\mathcal{H}\) 的范数小于 1 的自然映射 \(\iota  : {\mathcal{H}}_{T} \rightarrow  \mathcal{H}\) . 下证该映射 \(\iota\) 为单射. 设 \({u}_{j} \in  \mathcal{D}(T)\) 为 \({\mathcal{H}}_{T}\) 中柯西列使得 \(\left\{  {u}_{j}\right\}\) 在 \(\mathcal{H}\) 中收敛到 0,即
\begin{equation*}
\left( {{u}_{j} - {u}_{k},T\left( {{u}_{j} - {u}_{k}}\right) }\right)  \rightarrow  0\left( {j,k \rightarrow  \infty }\right) ,\;\begin{Vmatrix}{u}_{j}\end{Vmatrix} \rightarrow  0\left( {j \rightarrow  \infty }\right) .
\end{equation*}

我们需证当 \(j \rightarrow  \infty\) 时有 \(\left( {T{u}_{j},{u}_{j}}\right)  \rightarrow  0\) . 实际上,
\begin{equation*}
\left( {{u}_{j} - {u}_{k},T\left( {{u}_{j} - {u}_{k}}\right) }\right)  = \left( {{u}_{j},T{u}_{j}}\right)  + \left( {{u}_{k},T{u}_{k}}\right)  - 2\left( {{u}_{k},T{u}_{j}}\right)
\geq  \left( {{u}_{j},T{u}_{j}}\right)  - 2\left( {{u}_{k},T{u}_{j}}\right)\rightarrow  \left( {T{u}_{j},{u}_{j}}\right) ,\;k \rightarrow  \infty .
\end{equation*}

因此,当 \(j \rightarrow  \infty\) 时有 \(\left( {{u}_{j},T{u}_{j}}\right)  \rightarrow  0\) . 从而我们可以视 \({\mathcal{H}}_{T}\) 为 \(\mathcal{H}\) 的一个子空间.

\end{proof}


\begin{definition}[Friedrichs 扩张]
设 $T$ 为半有界稠定对称算子，定义其形式域为
\[
\mathcal H_T=\overline{D(T)}^{\|\cdot\|_T},
\qquad 
\|u\|_T^2 = \|u\|^2 + (Tu,u).
\]
下面构造 $T$ 的 Friedrichs 扩张，记为 $T_F$。
\end{definition}

\begin{proof}[Friedrichs 扩张的构造]
任取 $g\in\mathcal H$，定义线性泛函
\begin{equation}\label{eq:l-def}
\ell(v):=(v,g), \qquad v\in\mathcal H .
\end{equation}
显然，
\[
|\ell(v)|\le \|v\|\,\|g\|.
\]
因此对 $v\in\mathcal H_T$ 亦有
\begin{equation}\label{eq:l-bound}
|\ell(v)|\le \|v\|_T\,\|g\|.
\end{equation}
故 $\ell$ 是 $\mathcal H_T$ 上的有界线性泛函。由 Riesz--Fréchet 表示定理，存在唯一的
\[
w\in \mathcal H_T
\]
使得
\begin{equation}\label{eq:riesz}
\ell(v) = (v,w)_T, \qquad \forall v\in \mathcal H_T .
\end{equation}

我们定义 Friedrichs 扩张的定义域为所有此类 $w$ 的集合：
\[
D(T_F):=\{ w\in\mathcal H_T : \exists g\in \mathcal H \text{ 满足 \eqref{eq:l-def} 与 \eqref{eq:riesz}}\}.
\]
并定义
\begin{equation}\label{eq:TF-def}
T_F w := g .
\end{equation}

由构造知 $g$ 与 $w$ 一一对应，故 $T_F$ 为单值算子。


\textbf{Friedrichs 扩张包含原算子 $T$。}
在 \eqref{eq:l-def} 中取 $g=Tu$（$u\in D(T)$），则
\[
\ell(v)=(v,Tu)=(v,u)_T,\qquad v\in D(T).
\]
由 Riesz 定理知 $w=u$，且 $T_F u = Tu$。故
\[
D(T)\subset D(T_F),\qquad T_F|_{D(T)} = T.
\]

\textbf{$T_F$ 为对称的。}
由 \eqref{eq:riesz} 与 \eqref{eq:TF-def}，
\[
(v,w)_T = (v,T_F w),\qquad v\in\mathcal H_T.
\]
若 $v,w\in D(T_F)$，交换二者得
\[
(w,v)_T = (w,T_F v).
\]
由内积的反对称性可得
\[
(v,w)_T = (T_F v,w).
\]
与前式比较得
\[
(T_F v,w) = (v,T_F w),
\]
即 $T_F$ 为对称算子。


\textbf{$T_F$ 的逆 $M=T_F^{-1}$ 为有界对称闭算子。}
$T_F$ 将 $D(T_F)$ 双射到 $\mathcal H$，故其逆 $M$ 定义在整个 $\mathcal H$ 上。对任意 $x_n\to x$ 且 $Mx_n\to u$，由对称性，
\[
(Mx_n,y)=(x_n,My),\qquad \forall y\in\mathcal H.
\]
取极限得
\[
(u,y)=(x,My)=(Mx,y),\qquad \forall y.
\]
故 $u=Mx$。因此 $M$ 为闭算子。由闭图像定理知 $M$ 为有界算子。


\textbf{$T_F$ 为自伴的。}
由命题 5.1 可知：对每个 $\lambda\in\mathbb C\setminus\mathbb R$，有
\[
\lambda\in\rho(M),
\]
即 $\lambda I - M$ 可逆。换言之，
\[
\lambda^{-1} \in \rho(M^{-1}) = \rho(T_F).
\]
由定理 5.5（“若算子有非实预解点则自伴”）可得 $T_F$ 为自伴算子。

因此，Friedrichs 扩张 $T_F$ 是 $T$ 的自伴扩张。
\end{proof}
\begin{note}
     实际上, 上述证明中所构造的扩张不依赖正下界是如何选择的, 在某种意义上来讲该扩张是最大可能的自伴扩张. 我们称该扩张为Friedrichs 扩张.
\end{note}
\begin{example}
    若 \(u \in  {C}^{1}\left( \bar{I}\right)\) 中满足 \(u\left( {\pm 1}\right)  = {u}^{\prime }\left( {\pm 1}\right)  = 0\) 及 \({u}^{\prime \prime } \in  {L}^{2}\left( I\right)\) ,则例题5.2中算子 \(T\) 的闭包 \(\bar{T}\) 可以定义. 从而我们有
\begin{equation*}
\left( {-{Tu},u}\right)  = {\int }_{0}^{1}{\left| {u}^{\prime }\right| }^{2}{dt},\;u \in  \mathcal{D}(T)
\end{equation*}

且对任意 \(u \in  \mathcal{D}(T)\) 有
\begin{equation*}
{\int }_{0}^{1}{\left| u\right| }^{2}{dt} =  - {\int }_{0}^{1}\left( {{u}^{\prime }\bar{u} + u{\bar{u}}^{\prime }}\right) {tdt} \leq  2\parallel u\parallel \begin{Vmatrix}{u}^{\prime }\end{Vmatrix}.
\end{equation*}

从而有 \(\parallel u\parallel  \leq  2\begin{Vmatrix}{u}^{\prime }\end{Vmatrix}\) . 这意味着算子 \(- {4T}\) 满足定理5.9证明中的假设,且
\begin{equation*}
D = \left\{  {u \in  {C}^{1}\left( \bar{I}\right)  \mid  u\left( {\pm 1}\right)  = {u}^{\prime }\left( {\pm 1}\right)  = 0,{u}^{\prime } \in  {L}^{2}\left( I\right) }\right\}  .
\end{equation*}

因此, Friedrichs 扩张对应于Dirichlet边界条件 \({\left. u\right| }_{\partial I} = 0\) .

\end{example}


在应用中, 我们常常需要考虑自伴算子的扰动算子, 我们自然会关心自伴算子的小扰动是否还是自伴的。对此, Rellich证明了如下扰动定理.
\begin{theorem}[Rellich]
     设 \(T\) 为 \(\mathcal{H}\) 上的自伴算子, \(V\) 为对称算子且满足 \(\mathcal{D}(T)  \subset\)  \(\mathcal{D}\left( V\right)\) 及
\begin{equation*}
\parallel {Vu}\parallel  \leq  a\parallel {Tu}\parallel  + b\parallel u\parallel ,\;u \in  \mathcal{D}(T) . \label{5.20}
\end{equation*}

若 \(a < 1\) ,则 \(T + V\) 在 \(\mathcal{D}(T)\) 上的限制为自伴算子.

\end{theorem}
\begin{proof}
     \(T + V\) 为对称算子是显然的. 当 \(u \in  \mathcal{D}(T)\) 时有
\begin{equation*}
\parallel \left( {{Tu} + {Vu}}\right) \parallel  \geq  \parallel {Tu}\parallel  - \parallel {Vu}\parallel  \geq  \left( {1 - a}\right) \parallel {Tu}\parallel  - b\parallel u\parallel .
\end{equation*}

由 \(a < 1\) 知,
\begin{equation*}
\parallel {Tu}\parallel  \leq  {\left( 1 - a\right) }^{-1}\parallel \left( {T + V}\right) u\parallel  + {\left( 1 - a\right) }^{-1}b\parallel u\parallel .
\end{equation*}

这意味着 \(\mathcal{D}\left( \overline{T + V}\right)  = \mathcal{D}(T)  = \mathcal{D}\left( {T + V}\right)\) . 从而 \(T + V\) 为闭算子. 因此,要证 \(T + V\) 自伴只需证其亏损指标为零. 由于对任意 \(t \in  \left\lbrack  {0,1}\right\rbrack\) ,算子 \(T + {tV}\) 都是闭对称算子,由定理3.10及命题5.1的证明知, \(R\left( {\pm {iI} + T + {tV}}\right)\) 的余维不依赖于 \(t\) . 因此,由 \(T\) 的亏损指标为零知, \(T + {tV}\) 的亏损指标为零.
\end{proof}


在应用中, 如何证明算子为自伴的往往显得极为重要. 有时会遇到算子本身并不自伴, 但其有自伴扩张. 此时, 我们往往会关心其自伴扩张的个数. 为此, 我们先引进本质自伴的概念.
\begin{definition}
    对称算子 \(T\) 称为本质自伴的如果其闭包 \(\bar{T}\) 为自伴的.

\end{definition}
\begin{note}
 若 \(T\) 为本质自伴的,则 \(T\) 有且仅有一个自伴扩张. 事实上,若 \(S\) 为 \(T\) 的一个自伴扩张,则 \(S\) 为闭算子. 由 \(S \supset  T,S \supset  {T}^{* * }\) ,因此, \(S = {S}^{ * } \subset  {\left( {T}^{* * }\right) }^{ * } = {T}^{* * }\) ,从而 \(S = {T}^{* * }\) .

若 \(T\) 有且仅有一个自伴扩张,则 \(T\) 为本质自伴的,参见习题5.3.5.

由于 \({T}^{ * } = {\bar{T}}^{ * } = {T}^{* *  * }\) ,故 \(T\) 为本质自伴当且仅当

\begin{equation*}
T \subset  {T}^{* * } = {T}^{ * }\text{ . }
\end{equation*}

\end{note}


\subsection{习 题}
\begin{problemset}
    \item  (i) 设 \(S\) 为对称算子, \(S \supset  T\) . 若 \(R\left( {T + i}\right)  = R\left( {S + i}\right)\) ,则 \(T = S\) .

(ii) 设 \(T\) 为对称算子且 \(R\left( {T + i}\right)  = \mathcal{H}\) ,但 \(R\left( {T - i}\right)  \neq  \mathcal{H}\) ,则 \(T\) 没有自伴扩张.
\begin{proof}
    
\end{proof}
\item  设 \(\mathcal{H} = {L}^{2}\left( {\mathbb{R}}_{ + }\right)\) . 在 \(\mathcal{D}(T)  = {C}_{0}^{1}\left( {\mathbb{R}}_{ + }\right)\) 上定义算子 \(T = i\frac{d}{\dif x}\) . 则 \(T\) 为对称算子且 \(\bar{T}\) 不是自伴的. 进而证明 \(T\) 没有自伴扩张.

    
\begin{proof}
\textbf{1. $T$ 是对称的。}

对任意 $f,g\in C_0^1(\mathbb R_+)$，两者在 $[0,\infty)$ 上具有紧支集，积分分部得
\[
\langle Tf,g\rangle
= \int_0^\infty i f'(x)\overline{g(x)}\,dx
= i\left[f(x)\overline{g(x)}\right]_{0}^{\infty}
- \int_0^\infty i f(x)\overline{g'(x)}\,dx.
\]
由于 $f,g$ 有紧支集且在足够靠近 $0$ 的邻域上为 $0$，边界项为 $0$，于是
\[
\langle Tf,g\rangle = \int_0^\infty f(x)\overline{i g'(x)}\,dx
= \langle f,Tg\rangle.
\]
因此 $T$ 为对称算子。

\textbf{2. 计算伴随算子 $T^*$ 和亏损方程。}

首先描述 $T^*$ 的定义域与作用。由定义 $u\in D(T^*)$ 当且仅当存在 $v\in L^2(\mathbb R_+)$ 使得
\[
\langle Tf,u\rangle=\langle f,v\rangle\qquad\forall f\in C_0^1(\mathbb R_+).
\]
对任意 $f\in C_0^1(\mathbb R_+)$ 积分分部（注意 $f$ 在 $0$ 的一个邻域上为 $0$，因此边界项在 $0$ 处消失）：
\[
\int_0^\infty i f'(x)\overline{u(x)}\,dx
= -\int_0^\infty i f(x)\overline{u'(x)}\,dx,
\]
形式上可写为 $\langle Tf,u\rangle=\langle f,\,i u'\rangle$（若 $u$ 是绝对连续且 $u'\in L^2$）。于是必要且充分条件是：  
\(u\) 在 $(0,\infty)$ 上为绝对连续函数，$u'\in L^2(\mathbb R_+)$，并且可以取 $v=i u'$。 因此
\[
T^* u = i u',\qquad
D(T^*)=\{u\in L^2(\mathbb R_+): u\in AC_{\mathrm{loc}}(0,\infty),\ u'\in L^2(0,\infty)\}.
\]
（注意：由于测试函数在 $0$ 的邻域内为 $0$，在这一步并未得到任何关于 $u(0)$ 的条件，故 $u(0)$ 可任意。）

\medskip

\textbf{3. 求解亏损方程并计算亏损指标。}

按 von Neumann 理论，亏损空间为 $\ker(T^*\mp iI)$，其维数为亏损指标
\[
n_+(T)=\dim\ker(T^*-iI),\qquad n_-(T)=\dim\ker(T^*+iI).
\]

先解方程 $(T^*-iI)u=0$，即
\[
i u'(x) - i u(x)=0 \quad\Longleftrightarrow\quad u'(x)-u(x)=0.
\]
通解为 $u(x)=C e^{x}$. 在 $(0,\infty)$ 上 $e^{x}\not\in L^2(0,\infty)$，故此方程在 $L^2$ 中仅有平凡解，得
\[
n_+(T)=0.
\]

再解 $(T^*+iI)u=0$，即
\[
i u'(x) + i u(x)=0 \quad\Longleftrightarrow\quad u'(x)+u(x)=0.
\]
通解为 $u(x)=C e^{-x}$，且 $e^{-x}\in L^2(0,\infty)$，因此该空间为一维，得
\[
n_-(T)=1.
\]

于是亏损指标为 $(n_+,n_-)=(0,1)$，二者不相等。

\medskip

\textbf{4. 结论：$\overline T$ 不是自伴的，且 $T$ 无自伴扩张。}

由于 $T$ 的两个亏损指标不相等，按照 von Neumann 的自伴扩张判据（存在自伴扩张当且仅当 $n_+=n_-$）可得：$T$ 无自伴扩张。特别地，$\overline T$ 也不是自伴的（若闭包自伴则 $T$ 就有自伴扩张），因此 $\overline T\neq T^*$。

这完成证明。
\end{proof}
\item 设 \(H = {L}^{2}\left( {0,1}\right)\) . 在 \(\mathcal{D}(T)  = \left\{  {u \in  {C}_{0}^{1}\left( \left\lbrack  {0,1}\right\rbrack  \right)  \mid  u\left( 0\right)  = u\left( 1\right)  = 0}\right\}\) 上定义算子 \(T = i\frac{d}{\dif x}\) . 则 \(T\) 为对称算子且 \(\bar{T}\) 不是自伴的,但 \(T\) 有自伴扩张.

\begin{proof}
\textbf{(1) $T$ 对称。} 

对任意 $u,v\in \mathcal D(T)$，积分分部得：
\[
\langle Tu, v \rangle = \int_0^1 i u'(x) \overline{v(x)} \, dx
= i \big[ u(x)\overline{v(x)} \big]_0^1 - \int_0^1 i u(x) \overline{v'(x)} \, dx.
\]
由于 $u(0)=u(1)=v(0)=v(1)=0$，边界项为零，因此
\[
\langle Tu, v \rangle = \int_0^1 u(x) \overline{i v'(x)} \, dx = \langle u, Tv \rangle,
\]
故 $T$ 对称。

\textbf{(2) 闭包与亏损指标。}

由定义，$u \in \mathcal D(T^*)$ 当且仅当存在 $w\in L^2(0,1)$ 使得
\[
\langle Tu, \phi \rangle = \langle u, w \rangle, \quad \forall \phi \in \mathcal D(T).
\]
积分分部可得 $\mathcal D(T^*) = \{ u \in L^2(0,1) : u \in AC[0,1],\ u' \in L^2(0,1) \}$，且 $T^* u = i u'$。

亏损空间定义为
\[
\mathcal N_\pm = \ker(T^* \mp i I).
\]

方程 $(T^* - i I)u = 0$ 通解为 $u(x) = C e^x$，且 $e^x \in L^2(0,1)$，因此 $\dim \mathcal N_+ = 1$。  
方程 $(T^* + i I)u = 0$ 通解为 $u(x) = C e^{-x}$，因此 $\dim \mathcal N_- = 1$。  
由 von Neumann 理论，亏损指标 $(n_+, n_-) = (1,1)$，非零，故 $\overline T$ 不是自伴算子，但存在自伴扩张。

\textbf{(3) 自伴扩张。} 

自伴扩张由从 $\mathcal N_+$ 到 $\mathcal N_-$ 的酉映射一一对应。由于 $\dim \mathcal N_\pm = 1$，任意酉映射为乘以 $e^{i\theta}$，$\theta \in [0,2\pi)$。

因此每个自伴扩张 $T_\theta$ 可定义为：
\[
\mathcal D(T_\theta) = \{ u \in AC[0,1] : u' \in L^2(0,1),\ u(1) = e^{i\theta} u(0) \}, \quad
T_\theta u = i u'.
\]

验证对称性：对任意 $u,v \in \mathcal D(T_\theta)$，
\[
\langle T_\theta u, v \rangle - \langle u, T_\theta v \rangle
= i[u(x) \overline{v(x)}]_0^1 = i(u(1)\overline{v(1)} - u(0)\overline{v(0)}) = 0,
\]
因为 $u(1)=e^{i\theta} u(0)$，$v(1)=e^{i\theta} v(0)$。由此可验证 $T_\theta$ 对称，且亏损指标为零，从而自伴。
\end{proof}
\item  若 \(T\) 为下方有界的稠定对称算子,即存在常数 \(C\) 使得
\begin{equation*}
\left( {{Tu},u}\right)  \leq  C\left( {u,u}\right) ,\;u \in  \mathcal{D}(T) ,
\end{equation*}

则 \(T\) 有自伴扩张.
\item  若 \(T\) 有一个且只有一个自伴扩张,则 \(T\) 为本质自伴的.

\end{problemset}




\end{document}